% This is part of Agregation : modélisation
% Copyright (c) 2011
%   Laurent Claessens
% See the file fdl-1.3.txt for copying conditions.

% This is part of Mes notes de mathématique
% Copyright (c) 2011-2013
%   Laurent Claessens
% See the file fdl-1.3.txt for copying conditions.

\documentclass[a4paper,oneside,11pt]{book}
%\documentclass[a4paper,12pt,draft]{book}

\usepackage{etex}

\usepackage{ifthen}
%\usepackage{pdfsync}

\usepackage{latexsym}
\usepackage{amsfonts}
\usepackage{amsmath}
\usepackage{amsthm}
\usepackage{amssymb}
\usepackage{mathrsfs}           
\usepackage{mathabx}           % Pour \divides
\usepackage{bbm}


\usepackage{pstricks,pst-eucl,pstricks-add,calc,pst-math,catchfile}   % Les dépendances de phystricks.
\usepackage{graphicx}                   % Pour l'inclusion d'image en pfd.

\newcommand{\EpsOrPdfincludegraphics}[2][]{%
        \ifpdf
            \includegraphics[#1]{#2.png}
        \else
            \includegraphics[#1]{#2.eps}
        \fi
        }

\usepackage{subfigure}

\usepackage{fancyvrb}
\usepackage{stmaryrd}       % Pour le \obslash
\usepackage{xstring}        % Utilisé pour les références vers wikipédia
\usepackage{cases}
\usepackage{lscape}         % pour l'environnement landscape, utilisé dans la correction corr0076.tex
\usepackage{multicol}
\usepackage{import}         % Pour le hack qui sert à inclure GeomAnal

\usepackage[normalem]{ulem}		% Pour le barré, commande \sout
%\usepackage{soul}		% Pour le barré, commande \st           Celui-ci est enlevé parce que l'autre permet de barrer des é

\usepackage[all]{xy}

\let\second\undefined      % le paquet amthabx définit \second
\let\degree\undefined       % le paquet amthabx définit \degree
\usepackage[cdot,thinqspace,amssymb]{SIunits} 
 % L'option amssymb sert à éviter un conflit avec la commande \square de amssymb. Note qu'elle n'est plus accessible. Si tu en as besoin, faudra RTFM
%ftp://ftp.belnet.be/packages/ctan/macros/latex/contrib/SIunits/SIunits.pdf

%%%%%%%%%%%%%%%%%%%%%%%%%%
%
%   Trucs mathématiques
%
%%%%%%%%%%%%%%%%%%%%%%%%

% ENSEMBLES DE NOMBRES
\newcommand{\eA}{\mathbbm{A}}
\newcommand{\eC}{\mathbbm{C}}
\newcommand{\eD}{\mathbbm{D}}
\newcommand{\eE}{\mathbbm{E}}
\newcommand{\eF}{\mathbbm{F}}
\newcommand{\eG}{\mathbbm{G}}
\newcommand{\eH}{\mathbbm{H}}
\newcommand{\eK}{\mathbbm{K}}
\newcommand{\eL}{\mathbbm{L}}
\newcommand{\eM}{\mathbbm{M}}
\newcommand{\eN}{\mathbbm{N}}
\newcommand{\eP}{\mathbbm{P}}
\newcommand{\eQ}{\mathbbm{Q}}
\newcommand{\eR}{\mathbbm{R}}
\newcommand{\eZ}{\mathbbm{Z}}

% ENSEMBLES de fonctions
\newcommand{\aL}{\mathcal{L}}       % Les applications linéaires
\newcommand{\aC}{\mathcal{C}}       % Les fonctions C^1, C^2 etc

% AUTRES
\newcommand{\sdS}{\mathcal{S}}      % L'ensemble des subdivisions d'un intervalle.



\newcommand{\mF}{\mathcal{F}}
\newcommand{\mG}{\mathcal{G}}
\newcommand{\mI}{\mathcal{I}}
\newcommand{\mL}{\mathcal{L}}
\newcommand{\mS}{\mathcal{S}}   % Utilisé pour l'espace des fonctions Schwartz
\newcommand{\mZ}{\mathcal{Z}}


\newcommand{\mtu}{\mathbbm{1}}              % La matrice unité
\newcommand{\caract}{\mathbbm{1}}    % Characteristic function of a set

\DeclareMathOperator{\val}{val}     % valuation d'un polynôme
\DeclareMathOperator{\res}{res}     % Le résultant de deux polynômes
\DeclareMathOperator{\Inv}{Inv}     % L'application inverse


%\newcommand{\efrac}[2]{\frac{ \displaystyle #1 }{\displaystyle #2 }}
%%%%%%%%%%%%%%%%%%%%%%%%%%
%
%   Numérotations en tout genre
%
%%%%%%%%%%%%%%%%%%%%%%%%

\setcounter{tocdepth}{2}        % Profondeur de la table des matières
\setcounter{secnumdepth}{2}     % Profondeur dans le texte

%%%%%%%%%%%%%%%%%%%%%%%%%%
%
%   Les lignes magiques pour le texte en français.
%
%%%%%%%%%%%%%%%%%%%%%%%%

\usepackage[utf8]{inputenc}
\usepackage[T1]{fontenc}

\usepackage{listingsutf8}
\lstset{language=python,basicstyle=\footnotesize,tabsize=3,numbers=left,numberstyle=\tiny,frame=single,commentstyle=\ttfamily\color[rgb]{0,0,0.5},stringstyle=\color[rgb]{0,0.5,0},title=\lstname,inputencoding=utf8/latin1}
% Il y a du code Octave dans le chapitre et les corrections de Matlab.

\usepackage{textcomp}
\usepackage{lmodern}
\usepackage[a4paper,margin=2cm]{geometry} 
\usepackage[english,frenchb]{babel}


%\usepackage[ps2pdf]{hyperref}                  %Doit êre appelé en dernier.
\usepackage{hyperref}                           %Doit êre appelé en dernier.
    % Le hypersetup doit être dans mes_notes.tex pour pouvoir être changé automatiquement par le script 
    % pour faire la version agreg.

\usepackage[fr]{exocorr}
%%%%%%%%%%%%%%%%%%%%%%%%%%
%
%   Les théorèmes et choses attenantes
%
%%%%%%%%%%%%%%%%%%%%%%%%


\newcounter{numtho}
\newcounter{numprob}

\makeatletter
\@addtoreset{numtho}{chapter}
\@addtoreset{CountExercice}{chapter}
\makeatother

\newlength{\EnvSpace}
\setlength{\EnvSpace}{9pt}      % C'est la distance que je veux mettre avant et après les théorèmes, remarques, \ldots

\newtheoremstyle{MyTheorems}%
        {\EnvSpace}{\EnvSpace}%
        {\itshape}%
        {}%
        {\bfseries}{.}%
        {\newline}%
        {}%
\newtheoremstyle{MyExamples}%
        {\EnvSpace}{\EnvSpace}%
        {}%
        {}%
        {\bfseries}{.}%
        {\newline}%
        {}%
\newtheoremstyle{MyRemarks}%
        {\EnvSpace}{\EnvSpace}%
        {}%
        {}%
        {\bfseries}{.}%
        {\newline}%
        {}%

%\theoremstyle{MyExamples}   %\newtheorem{exemple}[numtho]{Exemple}      % Pour unification, ne plus utiliser
%                            \newtheorem{example}[numtho]{Exemple}
\newcounter{CounterExample}
\renewcommand{\theCounterExample}{\thechapter.\arabic{CounterExample}}

\newenvironment{example}{\vspace{\EnvSpace}\refstepcounter{numtho}\noindent{\bf Exemple \thenumtho}\newline}{\phantom{a}\hfill $\triangle$\vspace{\EnvSpace}}
%TODO : example est un environnement et non un théorème seulement pour cette histoire de triangle. Il faudrait trouver quelque chose.

\theoremstyle{MyRemarks}    \newtheorem{remark}[numtho]{Remarque}

                \newtheorem{amusement}[numtho]{Amusement}
                \newtheorem{erreur}[numtho]{Error}
                \newtheorem{probleme}[numprob]{\fbox{\bf Problèmes et choses à faire}}

\theoremstyle{MyTheorems}
            \newtheorem{lemma}[numtho]{Lemme}
            \newtheorem{corollary}[numtho]{Corollaire}
            \newtheorem{theorem}[numtho]{Théorème}      
            \newtheorem{definition}[numtho]{Définition}      
            \newtheorem{proposition}[numtho]{Proposition}      

            \newtheorem{exo}[CountExercice]{Exercice}       % C'est provisoire, pour Chafaï

\renewcommand{\thenumtho}{\thechapter.\arabic{numtho}}
% La numérotation des équations change dans les corrigés
\renewcommand{\theequation}{\thechapter.\arabic{equation}}
\renewcommand{\theCountExercice}{\arabic{CountExercice}}       % Ce compteur est défini dans SystemeCorr.sty
\newcommand{\defe}[2]{\textbf{#1}\index{#2}}

\renewcommand{\labelenumi}{\theenumi}
\renewcommand{\theenumi}{(\arabic{enumi})}


%%%%%%%%%%%%%%%%%%%%%%%%%%
%
%   Les macros qui font des choses
%
%%%%%%%%%%%%%%%%%%%%%%%%

\newcommand{\mA}{\mathcal{A}}
\newcommand{\mB}{\mathcal{B}}
\newcommand{\mO}{\mathcal{O}}
\newcommand{\mR}{\mathcal{R}}
\newcommand{\mT}{\mathcal{T}}
\newcommand{\mU}{\mathcal{U}}

\newcommand{\scal}[2]{ \langle #1,#2\rangle }
\newcommand{\vect}[1]{\overrightarrow{#1}}

\newcommand{\tq}{\text{ tel que }}
\newcommand{\tqs}{\text{ tels que }}
\newcommand{\quext}[1]{ \footnote{\textsf{#1}}  }
\newcommand{\info}[1]{\texttt{#1}}

\newcommand{\normal}{\lhd}
\newcommand{\swS}{\mathscr{S}}          % L'ensemble des fonctions Schwartz

\newcommand{\Borelien}{\mathcal{B}\text{or}}       % Les boréliens
\newcommand{\tribA}{\mathcal{A}}            % Une tribu A
\newcommand{\tribB}{\mathcal{B}}            
\newcommand{\tribF}{\mathcal{F}}            % Une tribu F

\newcommand{\affE}{\mathcal{E}}            % Un espace affine E
\newcommand{\affF}{\mathcal{F}}            
\newcommand{\affG}{\mathcal{G}}            

\newcommand{\statS}{\mathcal{S}}            % Un modèle statistique
\newcommand{\partP}{\mathcal{P}}            % L'ensemble des parties d'un ensemble

\newcommand{\polyP}{\mathcal{P}}            % L'ensemble des polynômes

\newcommand{\dB}{\mathscr{B}}       % la distribution de Bernoulli
\newcommand{\dE}{\mathscr{E}}       % la distribution exponentielle
\newcommand{\dG}{\mathscr{G}}       % la distribution géométrique.
\newcommand{\dM}{\mathscr{M}}       % la distribution multinomiale
\newcommand{\dN}{\mathscr{N}}       % la distribution normale.
\newcommand{\dP}{\mathscr{P}}       % la distribution de Poisson.
\newcommand{\dT}{\mathscr{T}}       % la distribution de Student
\newcommand{\dU}{\mathscr{U}}       % la distribution uniforme

\newcommand{\hL}{\mathscr{L}}       
\newcommand{\cL}{\hL}           % Pour la partie Chafai

\newcommand{\modE}{\mathcal{E}}         % Le E des modules
\newcommand{\modF}{\mathcal{F}}         % Le F des modules
\newcommand{\hH}{\mathscr{H}}           % Le H des espaces de Hilbert

%%%%%%%%%%%%%%%%%%%%%%%%%%
%
%   Bibliographie, index et liste des notations
%
%%%%%%%%%%%%%%%%%%%%%%%%

\usepackage{makeidx}
\usepackage[nottoc]{tocbibind}      % Le paquetage qui fait en sorte que la biblio soit inclue correctement dans la table des matières.
\usepackage[numbers]{natbib}                 % Pour avoir le champ URL dans le fichier bibtex
\usepackage[refpage]{nomencl}
\renewcommand{\nomname}{Liste des notations}
%
%   Comment introduire des éléments dans l'index des notations.
%
% La liste des tags à mettre pour bien classer mes notations est :
% T     pour la topologie et théorie des ensembles
%
% La syntaxe est facile, par exemple 
%       $\SL(2,\eR)$\nomenclature[G]{$\SL(2,\eR)$}{Le groupe de matrices deux par deux de déterminant 1.}
\renewcommand{\nomgroup}[1]{%
    \ifthenelse{\equal{#1}{A}}{\item[\textbf{Algèbre}]}{}%
    \ifthenelse{\equal{#1}{G}}{\item[\textbf{Géométrie}]}{}%
    \ifthenelse{\equal{#1}{R}}{\item[\textbf{Théorie des groupes}]}{}%
    \ifthenelse{\equal{#1}{P}}{\item[\textbf{Probabilités et statistique}]}{}%
    \ifthenelse{\equal{#1}{Y}}{\item[\textbf{Analyse}]}{}%
    \ifthenelse{\equal{#1}{M}}{\item[\textbf{Chaînes de Markov}]}{}%
}

%%%%%%%%%%%%%%%%%%%%%%%%%%
%
%   DeclareMathOperator
%
%%%%%%%%%%%%%%%%%%%%%%%%

\DeclareMathOperator{\signe}{sgn}
\DeclareMathOperator{\Vol}{Vol}
\DeclareMathOperator{\Int}{Int}     % Intérieur d'un ensemble.
\DeclareMathOperator{\Ind}{Ind}     % l'indice d'un chemin en analyse complexe
\DeclareMathOperator{\Diam}{Diam}
\DeclareMathOperator{\id}{Id}
\DeclareMathOperator{\Graph}{Graph}
\DeclareMathOperator{\Conv}{\mathcal{C}onv}
\DeclareMathOperator{\pr}{\texttt{proj}}
\DeclareMathOperator{\dom}{dom}

\DeclareMathOperator{\Graphe}{Gr}
\DeclareMathOperator{\Spec}{Spec}   % spectre d'un opérateur
\DeclareMathOperator{\arctg}{arctg}
\DeclareMathOperator{\cotg}{cotg}
\DeclareMathOperator{\cosec}{cosec}
\DeclareMathOperator{\arcsinh}{arcsinh}

%\DeclareMathOperator{\GL}{GL}   % le groupe linéaire
\newcommand*{\Sp}{\textup{Sp}}
\newcommand*{\Gl}{\textup{GL}}
\newcommand*{\GL}{\textup{GL}}      % Le groupe linéaire; je crois qu'ontologiquement c'est mieux que le DeclareMathOperator
\DeclareMathOperator{\PGL}{PGL}   % le groupe projectif
\DeclareMathOperator{\SO}{SO}
\DeclareMathOperator{\SL}{SL}
\DeclareMathOperator{\PSL}{PSL}   % Le groupe modulaire SL(2,Z)/Z2
\DeclareMathOperator{\gO}{O}
\DeclareMathOperator{\gS}{O}        % le groupe des matrices symétriques
\DeclareMathOperator{\SU}{SU}
\DeclareMathOperator{\gU}{U}

\DeclareMathOperator{\Reel}{Re}        % La partie réelle d'un nombre complexe

\DeclareMathOperator{\Image}{Image}        % ... avec \Image qui donne l'image d'une fonction ou d'un opérateur.
\DeclareMathOperator{\rang}{rg}
\DeclareMathOperator{\Kernel}{Ker}
\DeclareMathOperator{\Domaine}{Dom}
\DeclareMathOperator{\Span}{Span}
\DeclareMathOperator{\diag}{diag}
\DeclareMathOperator{\Hom}{Hom}
\DeclareMathOperator{\End}{End}     % L'ensemble des endomorphismes
\DeclareMathOperator{\tr}{Tr}       % la trace
\DeclareMathOperator{\Majorant}{Maj}
\DeclareMathOperator{\codim}{codim} % pour la codimension.
\DeclareMathOperator{\diam}{diam} % le diamètre d'un ensemble.

\DeclareMathOperator{\Var}{Var}     % Variance d'une variable aléatoire.
\DeclareMathOperator{\Fun}{\texttt{Fun}}     % Ensemble des applications d'un ensemble vers l'autre.
\DeclareMathOperator{\Cov}{Cov}     % la covariance.
\DeclareMathOperator{\gr}{gr}     % le groupe engendré
\DeclareMathOperator{\pgcd}{pgcd}     
\DeclareMathOperator{\ppcm}{ppcm}     
\DeclareMathOperator{\Frob}{Frob}     
\DeclareMathOperator{\Card}{Card}       % Le cardinal d'un ensemble.
\DeclareMathOperator{\Stab}{Stab}       % Le stabilisateur d'un point sous l'action d'un groupe.

\DeclareMathOperator{\Frac}{Frac}       % le corps des fractions d'un anneau
\DeclareMathOperator{\Aff}{Aff}         %  l'espace affine engendré

\newenvironment{subproof}{\begin{description}}{\end{description}}

\newcommand{\Dsddb}[4]{\frac{d}{d#2}\Big[#1\Big]_{#3=#4}}
\newcommand{\Dsdd}[3]{ \Dsddb{#1}{#2}{#2}{#3}   }

%%%%%%%%%%%%% TRUCS DE YVIK POUR FAIRE FONCTIONNER CdI1 %%%%%%%%%%%%%%%%%%%%%%
%

%\newcommand{\proofend}{\hspace*{\fill} $\Box$\\}
%\newcommand{\diam}{\hspace*{\fill} $\Diamond$\\}
%\def\s{\smallskip}
%\def\m{\medskip}
%\def\my{\bf}
\newcommand{\eps}{\varepsilon}
\newcommand{\Ker}{\operatorname{Ker}}
\newcommand{\IM}{\operatorname {Im}}
\newcommand{\cat}{\operatorname{cat}}
\newcommand{\crit}{\operatorname{crit}}
\newcommand{\Crit}{\operatorname{Crit}}
\newcommand{\Rest}{\operatorname{Rest}}
\newcommand{\grad}{\operatorname{grad}}
\newcommand{\sgrad}{\operatorname{sgrad}}
\newcommand{\Fix}{\operatorname{Fix}}
\newcommand{\pt}{\operatorname{pt}}
\newcommand{\cl}{\operatorname{cl}}
\newcommand{\B}{\operatorname {B}}
\newcommand{\C}{\operatorname {C}}
%\newcommand{\S}{\operatorname {S}}
\newcommand{\Gr}{\operatorname {Gr\;\!}}
%\def\dim{\operatorname {dim}}
\newcommand{\inj}{\operatorname {inj}}
%\newcommand{\Vol}{\operatorname {Vol}\:\!}
%\newcommand{\Int}{\operatorname {Int}\:\!}
\newcommand{\dist}{\operatorname {dist}}
%\def\inter{\operatorname {int}}
\newcommand{\ext}{\operatorname {ext}}
%\newcommand{\diameter}{\operatorname {diam}\:\!}
\newcommand{\Emb}{\operatorname {Emb}}
\newcommand{\can}{\operatorname {can}}
\newcommand{\euler}{\mbox{\rm e}}
\newcommand{\sii}{\mbox{\rm \scriptsize i}}
\newcommand{\VB}{\mbox{V}_{\!\!B}}   
\newcommand{\VC}{\mbox{V}_{\!\!C}}   
\newcommand{\VS}{\mbox{V}_{\!\!S}}   
\newcommand{\f}{\frac}
\newcommand{\ga}{\alpha}
\newcommand{\gb}{\beta}
%\newcommand{\gg}{\gamma}
\newcommand{\gd}{\delta}
\newcommand{\gve}{\varepsilon}
\newcommand{\gf}{\varphi}
\newcommand{\gk}{\kappa}
\newcommand{\gkk}{\varkappa}
\newcommand{\gl}{\lambda}
\newcommand{\go}{\omega}
\newcommand{\gs}{\sigma}
\newcommand{\gt}{\vartheta}
\newcommand{\gy}{\upsilon}
\newcommand{\gv}{\varrho}
\newcommand{\gz}{\zeta}
\newcommand{\gD}{\Delta}
\newcommand{\gF}{\Phi}
\newcommand{\gG}{\Gamma}
\newcommand{\gL}{\Lambda}
%\newcommand{\gO}{\Omega}
%\newcommand{\gS}{\Sigma}

%\long\def\forget#1\forgotten{} %
%\def\end{center}{{\mathfrak C}}
%\def\ea{{\mathfrak A}}

\newcommand{\ca}{{\mathcal A}}
\newcommand{\cb}{{\mathcal B}}
\newcommand{\cc}{{\mathcal C}}
\newcommand{\cd}{{\mathcal D}}
\newcommand{\ce}{{\mathcal E}}
\newcommand{\cf}{{\mathcal F}}
\newcommand{\cg}{{\mathcal G}}
\newcommand{\ch}{{\mathcal H}}
\newcommand{\cj}{{\mathcal J}}
\newcommand{\ck}{{\mathcal K}}
\newcommand{\cn}{{\mathcal N}}
\newcommand{\co}{{\mathcal O}}
\newcommand{\cp}{{\mathcal P}}
\newcommand{\cq}{{\mathcal Q}}
\newcommand{\cs}{{\mathcal S}}
\newcommand{\ct}{{\mathcal T}}
\newcommand{\cu}{{\mathcal U}}
\newcommand{\cv}{{\mathcal V}}
\newcommand{\cw}{{\mathcal W}}
\newcommand{\eb}{{\mathfrak B}}
\newcommand{\ed}{{\mathfrak D}}
\newcommand{\ee}{{\mathfrak E}}
\newcommand{\ef}{{\mathfrak F}}
\newcommand{\eg}{{\mathfrak G}}
\newcommand{\ej}{{\mathfrak J}}
\newcommand{\eh}{{\mathfrak H}}
\newcommand{\en}{{\mathfrak N}}
\newcommand{\eo}{{\mathfrak O}}
\newcommand{\ep}{{\mathfrak P}}
\newcommand{\eq}{{\mathfrak Q}}
\newcommand{\es}{{\mathfrak S}}
\newcommand{\et}{{\mathfrak T}}
\newcommand{\eu}{{\mathfrak U}}
\newcommand{\ev}{{\mathfrak V}}
\newcommand{\ew}{{\mathfrak W}}



%\def\NN{\mathbbm{N}}
%\def\QQ{\mathbbm{Q}}
%\def\RR{\mathbbm{R}}
%\def\SS{\mathbbm{S}}
%\def\11{\mathbbm{1}}
%\def\ZZ{\mathbbm{Z}}
%\def\TT{\mathbbm{T}}
\newcommand{\RR}{\eR}
\newcommand{\DD}{\mathbbm{D}}
\newcommand{\HH}{\mathbbm{H}}
\newcommand{\II}{\mathbbm{I}}
\newcommand{\N}{\mathbbm{N}}
\newcommand{\PP}{\mathbbm{P}}
\newcommand{\Q}{\mathbbm{Q}}
\newcommand{\RRR}{\mathbbm{R}_+}
\newcommand{\Z}{\mathbbm{Z}}
\newcommand{\RP}{{\RR\PP}} 
%\newcommand{\CP}{{\CC\PP}} 
\newcommand{\pp}{\partial}
\newcommand{\ww}{\wedge}
%\newcommand{\dc}{d^\CC}
\newcommand{\sym}{Sp(n;\RR)}
\newcommand{\ha}{\hookrightarrow}
\newcommand{\Ra}{\Rightarrow}
\newcommand{\Lra}{\Leftrightarrow} 

%\def\ni{\noindent}
%\def\b{\bigskip}
%\def\m{\medskip}
%\def\im{\mbox{Im}\,}

\newcommand{\de}{\stackrel{\mbox{\scriptsize{def}}}{=}}
%\newcommand{\id}{\mbox{id}}

%\def\sq{\square}
%\def\tr{\triangle}
%\def\trd{\bigtriangledown}
%\def\proof{\noindent {\it Proof. \;}}


%	La num\'erotation des exercices


\newcounter{exoNico}
\setcounter{exoNico}{1}
\newcommand{\exerNico}{\stepcounter{exoNico}{\bf Exercice }\arabic{exoNico}. }


%++++++++++ACCENTS++++++++++++++++++
\newcommand{\e}{\'{e}}
%\newcommand{\esp}{\'{e }}
%\newcommand{\eg}{\`{e}}
\newcommand{\ac}{\`{a} }
%\newcommand{\meme}{m\^{e}me }
\newcommand{\ou}{o\`{u} }

%+++++++++++NEWCOMMANDS+++++++++++
\newcommand{\dst}{\displaystyle}
\newcommand{\ba}{\begin{array}}
%\newcommand{\ea}{\end{array}}
%++++++++++FORMULAS+++++++++++++
\newcommand{\hs}{\hspace{0.3cm}}
%\newcommand{\eps}{\epsilon}
%\newcommand{\f}{\frac}
\newcommand{\arcth}{{\rm arctanh}}
\newcommand{\arcsh}{{\rm arcsinh}}
\newcommand{\arcch}{{\rm arccosh}}
\newcommand{\csec}{{\rm cosec}}
\newcommand{\cotan}{{\rm cotg}}
\newcommand{\cis}{(\cos+i\sin)( }
%\newcommand{\ra}{\rightarrow}
\newcommand{\lra}{\longrightarrow}
\newcommand{\ceil}{\rm plafond(}
\newcommand{\dfdu}{\frac{\partial f}{\partial u}}
\newcommand{\dfdw}{\frac{\partial f}{\partial w}}
\newcommand{\dfdx}{\frac{\partial f}{\partial x}}
\newcommand{\dfdy}{\frac{\partial f}{\partial y}}
\newcommand{\dudx}{\frac{\partial u}{\partial x}}
\newcommand{\dvdx}{\frac{\partial v}{\partial x}}
\newcommand{\dUdx}{\dfrac{\partial U}{\partial x}}
\newcommand{\dVdx}{\dfrac{\partial V}{\partial x}}
\newcommand{\dhdx}{\frac{\partial h}{\partial x}}
\newcommand{\dhdy}{\frac{\partial h}{\partial y}}
\newcommand{\dgdu}{\frac{\partial g}{\partial u}}
\newcommand{\dgdv}{\frac{\partial g}{\partial v}}
\newcommand{\dgudu}{\frac{\partial g_1}{\partial u}}
\newcommand{\dgudv}{\frac{\partial g_1}{\partial v}}
\newcommand{\dgddu}{\frac{\partial g_2}{\partial u}}
\newcommand{\dgddv}{\frac{\partial g_2}{\partial v}}
\newcommand{\dhdu}{\frac{\partial h}{\partial u}}
\newcommand{\dhdv}{\frac{\partial h}{\partial v}}
\newcommand{\dldu}{\frac{\partial l}{\partial u}}
\newcommand{\dldv}{\frac{\partial l}{\partial v}}
\newcommand{\dgudr}{\frac{\partial g_1}{\partial r}}
\newcommand{\dgudth}{\frac{\partial g_1}{\partial \theta}}
\newcommand{\dgddr}{\frac{\partial g_2}{\partial r}}
\newcommand{\dgddth}{\frac{\partial g_2}{\partial \theta}}
\newcommand{\dfdv}{\frac{\partial f}{\partial v}}
\newcommand{\dfdr}{\frac{\partial f}{\partial r}}

\newcommand{\dfdth}{\frac{\partial f}{\partial \theta}}
\newcommand{\ddfdx}{\frac{\partial^2 f}{\partial x^2}}
\newcommand{\ddfdy}{\frac{\partial^2 f}{\partial y^2}}
\newcommand{\ddfdxy}{\frac{\partial^2 f}{\partial y\partial x}}
\newcommand{\ddfdt}{\frac{\partial^2 f}{\partial^2 t}}

\newcommand{\ud}{\underline}

 % *** Blackboard math symbols ***
 %\newcommand{\N}{\mathbb{N}}
 %\newcommand{\Z}{\mathbb{Z}}
 %\newcommand{\Q}{\mathbb{R}}
 %\newcommand{\K}{\mathbb{K}}
 %\newcommand{\R}{\mathbb{R}}
 %\newcommand{\C}{\mathbb{C}}
 %\newcommand{\F}{\mathbb{F}}
 %\newcommand{\J}{\mathbb{J}}
\newcommand{\Qn}{\mathbb{Q}}

\newcommand{\Rn}{\eR} 
\newcommand{\Nn}{\eN}


\newtheorem{theo}{Th{\'e}or{\`e}me}[section]
\newtheorem{defn}{D{\'e}finition}
\newtheorem{prop}{Proposition}     % redef encore dans Chafaï
%\newtheorem{rem}{Remarque}[section]
\newtheorem{lem}{Lemme}[section]
\newcommand{\R}{\mathbb{R}}
\newcommand{\dem}{\textbf{D{\'e}monstration.}}
\newcommand{\vc}[1]{\boldsymbol{#1}}
\newcommand{\p}{\textrm{P}}
%\newcommand{\e}{\textrm{E}}
\newcommand{\mbt}{arbre binaire markovien}
\newcommand{\mbts}{arbres binaires markoviens}

\newcommand{\ea}{\end{array}}


%%%%%%%%%%%%%%%%%%%%%%%%%%%%%%%%%%%%%%
%
% les petis yeux 
%
%%%%%%%%%%%%%%%%%%%%%%%%%%%%%%%%%%%%%%%%%%%%%

\newcommand{\coolexo}{$\circledast\circledast$}
\newcommand{\boringexo}{$\circleddash\circleddash$}
\newcommand{\minsyndical}{$\odot\odot$}
\newcommand{\mortelexo}{$\obslash\oslash$}


%%%%%%%%%%%%%% FIN TRUCS DE YVIK %%%%%%%%%%%%%%%%%%%%%%

%%%%%%%%%%%%%% TRUCS DE PIERRE %%%%%%%%%%%%%%%%%%%%%%


% Le paquet array est là pour faire fonctionner l'environement arrowcases dans les trucs de Pierre.
\usepackage{array}

%\documentclass[11pt,a4paper,openany]{book}
%\usepackage[ansinew]{inputenc}
%\usepackage{pstricks, pst-node, array, ifpdf, comment, pst-plot}
%\usepackage[marginparwidth=2cm]{geometry}
%\usepackage[dvips,colorlinks]{hyperref}
%\usepackage[frenchb]{entetes}

%\usepackage{bigcenter}
%%%% debut macro %%%%
%%% ----------debut de bigcenter.sty--------------

%%% nouvel environnement bigcenter
%%% pour centrer sur toute la page (sans overfull)
%\makeatletter
%\newskip\@bigflushglue \@bigflushglue = -100pt plus 1fil

%\def\bigcenter{\trivlist \bigcentering\item\relax}
%\def\bigcentering{\let\\\@centercr\rightskip\@bigflushglue%
%\def\endbigcenter{\endtrivlist}

%\leftskip\@bigflushglue
%\parindent\z@\parfillskip\z@skip}
%\makeatother

%%% ----------fin de bigcenter.sty--------------
%%%% fin macro %%%%

%\input{mfpic}

% À régler par l'utilisateur
\newlength{\arrowsep}\setlength{\arrowsep}{3pt}
\newlength{\arrowlength}\setlength{\arrowlength}{1cm}

% Cet environnement est sympa, mais il dépend trop de ps; en tout cas il ne passe pas dans pdflatex
% 15 mars 1012
%\newenvironment{arrowcases}	{%
%			\pnode(\arrowsep,0.5ex){A}%
%			\hspace{\arrowlength}%
%			\begin{array}{>{\displaystyle\pnode(-\arrowsep,0.5ex){B}}l<{\ncline{A}{B}}@{}}
%				}
%			{
%			  \end{array}
%			}

\newenvironment{arrowcases}%
{\begin{cases}}
{\end{cases}}



\makeatletter %% \limite[condition]x x_0
\newcommand*{\limite}[3][\@empty]{\lim_{\substack{#2\rightarrow#3\\#1}}}
\makeatother

% \newenvironment{split+justif}{%
% \begin{split}%
% \let\ampori&
% \def&#1&#2\\{}
% }{%

% \end{split}}%

%\def\ncov{\tilde\nabla} % Nouvelle dérivée covariante
\newcommand*\sev{<} % 

%\newcommand{\hgot}{\mathfrak{h}} % h gothique (ss algebre de Lie)
%\def\var#1{{\mathbf #1}} % \var <-> Une variété
%\def\pardef{\stackrel{def}{=}} % = par définition.
%\newcommand{\bbar#1}{\bar{\bar{#1}}}
%\def\cov{\nabla} % Derivee covariante / connexion
%\newcommand{\gl}{\mathfrak{gl}} % algèbre linéaire
%\def\doubleprime{{\prime\prime}} % Isomorphique 
%\def\scal(#1,#2){\langle #1,#2\rangle}
%\def\agit(#1,#2){\langle #1,#2^\vee\rangle}
%\let\phiori\phi
%\let\phi\varphi
\let\ssi\iff
%\def\iddc{\mathcal I}
\newcommand*{\ideal}[1]{\{#1\}}
\newcommand*{\fleche}[1]{\stackrel{#1}\longrightarrow}

%\newcounter{exercice}
%\setcounter{exercice}{0}
\setcounter{CountExercice}{0}

% \newenvironment{exo}[1][\relax]{%
% \stepcounter{exercice}%
% \par\medskip%
% #1{\textbf{Exercice}~\arabic{exercice}.}\quad}%
% {\par}

% \newenvironment{rep}{\hspace{1em}\par\textbf{Solution
%     proposée.\quad}}{\par\noindent\hrulefill\par}

%\date{}
\newcommand{\Acplx}{A_\cdot}
\newcommand{\Bcplx}{B_\cdot}
\newcommand{\toisom}{\fleche\simeq}
\newcommand{\D}{\partial}
%\newcommand{\cat}[1]{{\bf #1}}
%\newcommand{\donc}{\Rightarrow}
%\newcommand{\im}{\text{im}}
%\newcommand{\coker}{\text{coker}}
\newcommand{\lied}{\mathcal L}
\newcommand*{\nom}[1]{\textsc{#1}}
\makeatletter
\newcommand*{\attention}[1]{\@latex@warning{#1}{!\small\bf #1!}\marginpar{Warning}}
\makeatother
\newcommand*{\inner}{\imath}
\newcommand*{\newexo}{}
\newcommand*{\principe}{}
\newcommand*{\etape}{}
\newcommand*{\preuve}{}
\newcommand*{\exr}{\item}
%\def\prim#1\expandafter\d#2 {\int #1\d#2}

\newcommand*{\crochets}[1]{\Bigl[ #1 \Bigr]}
\newcommand*{\llbrack}[1]{\left\lbrack #1 \right\lbrack}
\newcommand*{\rlbrack}[1]{\left\rbrack #1 \right\lbrack}
\newcommand*{\lrbrack}[1]{\left\lbrack #1 \right\rbrack}
\newcommand*{\rrbrack}[1]{\left\rbrack #1 \right\rbrack}
\newcommand*{\vecteur}[1]{\mathbf{#1}}


%% Maths : Les ensembles
\newcommand*{\ens}[1]{\mathbb{#1}} % Ensemble de nombres
\newcommand*{\var}[1]{\mathbf{#1}} % Variété
\newcommand*{\alg}[1]{\mathcal{#1}} % Algèbre
%\newcommand*{\RR}{\ens R}%
\newcommand*{\TT}{\ens T}% Tore !
%\expandafter\show\csname SS \endcsname
%\renewcommand*{\SS}{\var S}% 
%\newcommand*{\CC}{\ens C}%
\newcommand*{\ZZ}{\ens Z}%
\newcommand*{\QQ}{\ens Q}%
\newcommand*{\NN}{\ens N}%
\newcommand{\schwartz}{\mathcal S} % Espace de Schwartz
\newcommand*{\topologie}{\mathscr{T}}
\newcommand*{\Topologie}{\textcursive{T}}
\newcommand{\LL}{\text{\textup{L}}} %% Espace de Lebesgue droit
\newcommand{\Ll}{\mathcal{L}} %% Lebesgue ronde
\newcommand{\fronde}{\mathcal{F}} %% Transformée de Fourier.
\newcommand{\sigmaalgebre}[1]{\mathcal{#1}} %% Une sigma algèbre...
\DeclareMathOperator{\SymMatrix}{Sym}
\DeclareMathOperator{\ASymMatrix}{ASym}
\newcommand{\Sym}{\SymMatrix}
\newcommand{\ASym}{\ASymMatrix}
\newcommand{\transpose}[1]{{\vphantom{#1}}^{\mathit t}{\/#1}}
%\renewcommand*{\sp}{\textup{sp}}
\newcommand*{\dprime}{{\prime\prime}}
%\show\span
%\newcommand*{\Span}[1]{\mathopen> #1 \mathclose<}

%% Maths : Symboles divers
%\newcommand{\pp}{\text{\textup{~p.p.}}} %% Presque partout
%\PackageWarning{entetes}{Redefining command \d}
%\renewcommand{\d}{\mbox{$\,$\textrm{d}}}
\newcommand{\surj}{\vers}
\newcommand{\isom}{\simeq}
\newcommand*{\Tau}{\alg T}
\newcommand{\cdv}{\mathfrak{X}} % Champs de vecteurs


%% Maths : Constructions

%\let\Exp\exp
%\renewcommand{\exp}[1]{e^{#1}} % On préfère e^{} que exp{}

%\renewcommand{\exp}[1]{e^{#1}} % On préfère e^{} que exp{}
%\renewcommand{\vec}[1]{\mathbf{#1}} % Désigner un vecteur
\newcommand{\set}[1]{\left\{#1\right\}} % Un ensemble { }
\newcommand*{\abs}[1]{\left\vert#1\right\vert} % Valeur absolue.
\newcommand*{\module}[1]{\left\vert#1\right\vert} % Valeur absolue.
\newcommand*{\norme}[1]{\left\Vert#1\right\Vert} % norme
\newcommand*{\ordre}[1]{\left\vert#1\right\vert} % L'ordre d'un élément.
%\def\scal(#1,#2){% Produit scalaire.
%  \PackageWarning{entetes}{Obsolete command \string\scal}%
%  \scalprod{#1}{#2}%
%}
\newcommand*{\scalprod}[2]{\left\langle #1,#2\right\rangle}
\let\dual\ast

\newcommand*{\pardef}{\stackrel{\text{def}}{=}} % Par définition.
\newcommand*{\iffdefn}{\stackrel{\text{def}}{\iff}} % Par définition.
%\newcommand*{\telque}{\mbox{~\entetes@name@telque~}} % tel que, dans un ensemble.
\newcommand*{\Defn}[1]{\emph{#1}} %
\newcommand*{\tensor}{\otimes}
\newcommand*{\pder}[2]{\frac{\partial #1}{\partial #2}}

%{{{ Fraction in-line plus jolie
% \DeclareRobustCommand\sfrac[1]{\@ifnextchar/{\@sfrac{#1}}%
%                                             {\@sfrac{#1}/}}
% \def\@sfrac#1/#2{\leavevmode\kern.1em\raise.5ex
%          \hbox{$\m@th{\fontsize\sf@size\z@
%                            \selectfont#1}$}\kern-.1em
%          /\kern-.15em\lower.25ex
%           \hbox{$\m@th{\fontsize\sf@size\z@
%                             \selectfont#2}$}}
%}}} 

\DeclareRobustCommand{\sfrac}[3][\mathrm]{\hspace{0.1em}%
  \raisebox{0.4ex}{$#1{\scriptstyle
#2}$}\hspace{-0.1em}/\hspace{-0.07em}%
  \mbox{$#1{\scriptstyle #3}$}}



%% Maths : Opérateurs
%\DeclareMathOperator{\tr}{Tr}
%\DeclareMathOperator{\pr}{\texttt{pr}}
\DeclareMathOperator{\supp}{supp}
\DeclareMathOperator{\adh}{adh}
\DeclareMathOperator{\interior}{int}
\DeclareMathOperator{\im}{Im}
\DeclareMathOperator{\Id}{Id}
\DeclareMathOperator{\Aut}{Aut}
\DeclareMathOperator{\Iso}{Iso}
\DeclareMathOperator{\Jac}{Jac} % jacobienne
\DeclareMathOperator{\coker}{coker}
\DeclareMathOperator{\interieur}{int}
\DeclareMathOperator{\Tor}{Tor}
\DeclareMathOperator{\divg}{div}
\DeclareMathOperator{\rot}{rot}
%\DeclareMathOperator{\cosec}{cosec}


%% Pour obtenir le \Sha cyrillique...
% \RequirePackage[OT2,T1]{fontenc}
% \DeclareSymbolFont{cyrletters}{OT2}{wncyr}{m}{n}
% \DeclareMathSymbol{\Sha}{\mathalpha}{cyrletters}{"58}


% \newcounter{@institute}
% \let\authorori\author
% \def\@institute{}\def\@auteurs{}
% %\newcommand{\institute}[2]{\refstepcounter{@institute}\label{#1}\def\@institute{\@institute\small
% %#2}\set@authors}
% \newcommand{\institute}[2]{\refstepcounter{@institute}\label{#1}%
%   \let\maketitleori\maketitle%
%   \renewcommand\maketitle{\footnote{#2}\maketitleori}%
% }%
% \renewcommand{\author}[1]{\def\@auteurs{#1}\set@authors}
% \def\the@institute{${}^{(\roman{@institute})}$}
% \newcommand{\inst}[1]{\ref{#1}}

% \newcommand{\set@authors}{\authorori{\@auteurs}}% \\ \@institute}}


% \newcounter{@institute}
% \newcommand{\institute}[1]{
%   \let\labelori\label
%   \renewcommand{\label}[1]{%
%     \refstepcounter{@institute}\labelori{##1}
%     \begin{tabular}{cc}%format ?!
%     \begin{minipage}[t]
      


%   }

% }%

\newcommand*{\conclusion}{\emph{Conclusion~:~}}
\newcommand{\hint}{\par\emph{Aide~:~}\hspace{1em}}
\newcommand{\rappel}{\par\emph{Rappels~:~}\hspace{1em}}

%\newcounter{enumarray} 
%\newenvironment{enumarray}[1]{% Merci Ulrike Fischer
% \setcounter{enumarray}{0}%
% \begin{array}{% motif
%     >{% Au début de chaque ligne
%       \stepcounter{enumarray}%
%       (\alph{enumarray})%
%       \hspace{2em}
%     }% 
%     #1%
%   }% fin motif
% }{%
% \end{array}%
%}
\newenvironment{displayinline}{% displaystyle + inline.
  $\displaystyle%
}{%
  $%
}

\newcommand{\telque}{\vert\,}
\newcommand{\donc}{\Rightarrow}
 
%\newcounter{coroNico}
%\newcommand{\corrNico}{\refstepcounter{coroNico}\paragraph{Correction \arabic{coroNico}.}}


%

% Les commandes suivantes sont pour Chafaï:

%%% Macros Dj. 

%%% Pour les notes
\newcommand{\NB}[1]{{\large\textbf{*** #1 ***}}}

%%% Pour passer en mode mathématiques automatiquement 
\newcommand{\EM}{\ensuremath}

%%% Pour théorèmes etc...  Nécessite le paquetage amsthm ou bien classe amsart
%\newcommand{\THMEN}{%
% \theoremstyle{plain}%
% \newtheorem{thm}{Theorem}[section]%
% \newtheorem{cor}[thm]{Corollary}%
% \newtheorem{prop}[thm]{Proposition}%
% \newtheorem{lem}[thm]{Lemma}%
% \theoremstyle{definition}%
% \newtheorem{defi}[thm]{Definition}%
% \theoremstyle{remark}%
% \newtheorem{rem}[thm]{Remark}%
% \newtheorem{xpl}[thm]{Example}%
% \newtheorem{exe}[thm]{Exercise}%
% \newtheorem{hyp}[thm]{Hypothesis}%
%}

%\theoremstyle{plain}%
%\newtheorem{thm}{Théorème}[section]%
%\newtheorem{cor}[thm]{Corollaire}%
%\newtheorem{prop}[thm]{Proposition}%
%\newtheorem{lem}[thm]{Lemme}%
%\theoremstyle{definition}%
%\newtheorem{defi}[thm]{Définition}%
%\theoremstyle{remark}%
%%\newtheorem{rem}[thm]{Remarque}%
%\newtheorem{xpl}[thm]{Exemple}%
%%\newtheorem{exo}[thm]{Exercice}%
%\newtheorem{hyp}[thm]{Hypothèse}%
\newtheorem{eur}[numtho]{Heuristique}%
%\newtheorem{pro}[thm]{Problème}%

%%% Titre des *sections en + petit
\makeatletter
\newcommand{\SMALLSECS}{%
%\renewcommand{\section}{\@startsection%
 %{section}%                           % name
 %{1}%                                 % level
 %{0em}%                               % indent
 %{\baselineskip}%                     % beforeskip
 %{0.5\baselineskip}%                  % afterskip
 %{\normalfont\large\bfseries}}%       % style
%\renewcommand{\subsection}{\@startsection%
 %{subsection}%                        % name
 %{2}%                                 % level
 %{0em}%                               % indent
 %{\baselineskip}%                     % beforeskip
 %{0.25\baselineskip}%                 % afterskip
 %{\normalfont\bfseries}}%             % style
}
\makeatother

%%% Limites
\newcommand{\limLeb}[2]
{\underset{#1\to+\infty}{\overset{\bL^{#2}}{\longrightarrow}}}
\newcommand{\limE}[1]
{\underset{#1\to+\infty}{\overset{\text{étr.}}{\longrightarrow}}}
\newcommand{\limP}[1]
{\underset{#1\to+\infty}{\overset{\dP}{\longrightarrow}}}
\newcommand{\limL}[1]
{\underset{#1\to+\infty}{\overset{\cL}{\longrightarrow}}}
\newcommand{\limn}[1]
{\underset{#1\to+\infty}{\longrightarrow}}
\newcommand{\mylim}[2]
{\underset{#1\to#2}{\longrightarrow}}
\newcommand{\limPS}[1]
{\underset{#1\to+\infty}{\overset{\text{p.s}}{\longrightarrow}}}

%%% Heure 
\makeatletter
\providecommand{\timenow}{\@tempcnta\time
\@tempcntb\@tempcnta
\divide\@tempcntb60
\ifnum10>\@tempcntb0\fi\number\@tempcntb
\multiply\@tempcntb60
\advance\@tempcnta-\@tempcntb
:\ifnum10>\@tempcnta0\fi\number\@tempcnta}
\makeatother

% Pour le draft-stamping dans les bas de page
%\makeatletter
%\newcommand{\versiondetravail}{%
% \renewcommand{\@evenfoot}{%
% \hfil{\tiny\texttt{%
%   Version préliminaire, compilée le \today{} à \timenow.}\hfill}}%
% \renewcommand{\@oddfoot}{\@evenfoot}%
%}
%\makeatother

%%% Doubles lettres 
\newcommand{\dA}{\EM{\mathbb{A}}}
%\newcommand{\dB}{\EM{\mathbb{B}}}
\newcommand{\dC}{\EM{\mathbb{C}}}
\newcommand{\dD}{\EM{\mathbb{D}}}
%\newcommand{\dE}{\EM{\mathbb{E}}}
\newcommand{\dF}{\EM{\mathbb{F}}}
%\newcommand{\dG}{\EM{\mathbb{G}}}
\newcommand{\dH}{\EM{\mathbb{H}}}
\newcommand{\dI}{\EM{\mathbb{I}}}
\newcommand{\dJ}{\EM{\mathbb{J}}}
\newcommand{\dK}{\EM{\mathbb{K}}}
\newcommand{\dL}{\EM{\mathbb{L}}}
%\newcommand{\dM}{\EM{\mathbb{M}}}
%\newcommand{\dN}{\EM{\mathbb{N}}}
\newcommand{\dO}{\EM{\mathbb{O}}}
%\newcommand{\dP}{\EM{\mathbb{P}}}
\newcommand{\dQ}{\EM{\mathbb{Q}}}
\newcommand{\dR}{\EM{\mathbb{R}}}
\newcommand{\dS}{\EM{\mathbb{S}}}
%\newcommand{\dT}{\EM{\mathbb{T}}}
%\newcommand{\dU}{\EM{\mathbb{U}}}
\newcommand{\dV}{\EM{\mathbb{V}}}
\newcommand{\dW}{\EM{\mathbb{W}}}
\newcommand{\dX}{\EM{\mathbb{X}}}
\newcommand{\dY}{\EM{\mathbb{Y}}}
\newcommand{\dZ}{\EM{\mathbb{Z}}}

%%% Lettres droites
\newcommand{\rA}{\EM{\mathrm{A}}}
\newcommand{\rB}{\EM{\mathrm{B}}}
\newcommand{\rC}{\EM{\mathrm{C}}}
\newcommand{\rD}{\EM{\mathrm{D}}}
\newcommand{\rE}{\EM{\mathrm{E}}}
\newcommand{\rF}{\EM{\mathrm{F}}}
\newcommand{\rG}{\EM{\mathrm{G}}}
\newcommand{\rH}{\EM{\mathrm{H}}}
\newcommand{\rI}{\EM{\mathrm{I}}}
\newcommand{\rJ}{\EM{\mathrm{J}}}
\newcommand{\rK}{\EM{\mathrm{K}}}
\newcommand{\rL}{\EM{\mathrm{L}}}
\newcommand{\rM}{\EM{\mathrm{M}}}
\newcommand{\rN}{\EM{\mathrm{N}}}
\newcommand{\rO}{\EM{\mathrm{O}}}
\newcommand{\rP}{\EM{\mathrm{P}}}
\newcommand{\rQ}{\EM{\mathrm{Q}}}
\newcommand{\rR}{\EM{\mathrm{R}}}
\newcommand{\rS}{\EM{\mathrm{S}}}
\newcommand{\rT}{\EM{\mathrm{T}}}
\newcommand{\rU}{\EM{\mathrm{U}}}
\newcommand{\rV}{\EM{\mathrm{V}}}
\newcommand{\rW}{\EM{\mathrm{W}}}
\newcommand{\rX}{\EM{\mathrm{X}}}
\newcommand{\rY}{\EM{\mathrm{Y}}}
\newcommand{\rZ}{\EM{\mathrm{Z}}}

% Lettres caligraphiques
\newcommand{\cA}{\EM{\mathcal{A}}}
\newcommand{\cB}{\EM{\mathcal{B}}}
\newcommand{\cC}{\EM{\mathcal{C}}}
\newcommand{\cD}{\EM{\mathcal{D}}}
\newcommand{\cE}{\EM{\mathcal{E}}}
\newcommand{\cF}{\EM{\mathcal{F}}}
\newcommand{\cG}{\EM{\mathcal{G}}}
\newcommand{\cH}{\EM{\mathcal{H}}}
\newcommand{\cI}{\EM{\mathcal{I}}}
\newcommand{\cJ}{\EM{\mathcal{J}}}
\newcommand{\cK}{\EM{\mathcal{K}}}
%\newcommand{\cL}{\EM{\mathcal{L}}}
\newcommand{\cM}{\EM{\mathcal{M}}}
\newcommand{\cN}{\EM{\mathcal{N}}}
\newcommand{\cO}{\EM{\mathcal{O}}}
\newcommand{\cP}{\EM{\mathcal{P}}}
\newcommand{\cQ}{\EM{\mathcal{Q}}}
\newcommand{\cR}{\EM{\mathcal{R}}}
\newcommand{\cS}{\EM{\mathcal{S}}}
\newcommand{\cT}{\EM{\mathcal{T}}}
\newcommand{\cU}{\EM{\mathcal{U}}}
\newcommand{\cV}{\EM{\mathcal{V}}}
\newcommand{\cW}{\EM{\mathcal{W}}}
\newcommand{\cX}{\EM{\mathcal{X}}}
\newcommand{\cY}{\EM{\mathcal{Y}}}
\newcommand{\cZ}{\EM{\mathcal{Z}}}

%%% Lettres grasses 
\newcommand{\bA}{\EM{\mathbf{A}}}
\newcommand{\bB}{\EM{\mathbf{B}}}
\newcommand{\bC}{\EM{\mathbf{C}}}
\newcommand{\bD}{\EM{\mathbf{D}}}
\newcommand{\bE}{\EM{\mathbf{E}}}
\newcommand{\bF}{\EM{\mathbf{F}}}
\newcommand{\bG}{\EM{\mathbf{G}}}
\newcommand{\bH}{\EM{\mathbf{H}}}
\newcommand{\bI}{\EM{\mathbf{I}}}
\newcommand{\bJ}{\EM{\mathbf{J}}}
\newcommand{\bK}{\EM{\mathbf{K}}}
\newcommand{\bL}{\EM{\mathbf{L}}}
\newcommand{\bM}{\EM{\mathbf{M}}}
\newcommand{\bN}{\EM{\mathbf{N}}}
\newcommand{\bO}{\EM{\mathbf{O}}}
\newcommand{\bP}{\EM{\mathbf{P}}}
\newcommand{\bQ}{\EM{\mathbf{Q}}}
\newcommand{\bR}{\EM{\mathbf{R}}}
\newcommand{\bS}{\EM{\mathbf{S}}}
\newcommand{\bT}{\EM{\mathbf{T}}}
\newcommand{\bU}{\EM{\mathbf{U}}}
\newcommand{\bV}{\EM{\mathbf{V}}}
\newcommand{\bW}{\EM{\mathbf{W}}}
\newcommand{\bX}{\EM{\mathbf{X}}}
\newcommand{\bY}{\EM{\mathbf{Y}}}
\newcommand{\bZ}{\EM{\mathbf{Z}}}

%%% Quelques lettres grecques et symboles
\newcommand{\al}{\alpha}
\newcommand{\be}{\beta}
\newcommand{\De}{\Delta}
%\newcommand{\de}{\delta}
%\newcommand{\ga}{\gamma}
\newcommand{\Ga}{\Gamma}
\newcommand{\g}{\gamma}
\newcommand{\gn}{\g_n}
%\newcommand{\gt}[1]{\g^{\otimes #1}}
\newcommand{\la}{\lambda}
\newcommand{\La}{\Lambda}
\newcommand{\lan}{\la_n}
\newcommand{\na}{\nabla}
\newcommand{\Om}{\Omega}
\newcommand{\om}{\omega}
\newcommand{\ph}{\Phi}
\newcommand{\Si}{\Sigma}
\newcommand{\si}{\sigma}
\newcommand{\Te}{\Theta}
\newcommand{\te}{\theta}
\newcommand{\ta}{\tau}
\newcommand{\veps}{\varepsilon}
\newcommand{\vphi}{\varphi}
\newcommand{\bul}{\EM{\bullet}}

%%% Prototype de fonction avec dimensionnement
%\newcommand{\p}[4]{{#3}\!\left#1{#4}\right#2} 

%%% Normes et assimilées 
\newcommand{\ABS}[1]{\EM{{\left| #1 \right|}}} % |1|
\newcommand{\BRA}[1]{\EM{{\left\{#1\right\}}}} % {1}
\newcommand{\DP}[1]{\EM{{\left<#1\right>}}} % <1>
\newcommand{\NRM}[1]{\EM{{\left\| #1\right\|}}} % ||1||
\newcommand{\NI}[1]{\EM{{\left\| #1\right\|}_\infty}} % norme infinie
\newcommand{\OSC}[1]{\EM{{\p(){\mathrm{osc}}{#1}}}} % oscillation
\newcommand{\PAR}[1]{\EM{{\left(#1\right)}}} % (1)
\newcommand{\BPAR}[1]{\EM{{\biggl(#1\biggr)}}} % (1)
\newcommand{\BABS}[1]{\EM{{\biggl|#1\biggr|}}} % (1)
\newcommand{\pd}{\EM{\partial}} % dérivée partielle
\newcommand{\PD}[2]{\EM{{\frac{\partial #1}{\partial #2}}}} % dérivée partielle en fraction
\newcommand{\SBRA}[1]{\EM{{\left[#1\right]}}} % [1]
\newcommand{\VT}[1]{\EM{{\| #1\|}_{\mbox{{\scriptsize VT}}}}} % variation totale
\newcommand{\LIP}[1]{\EM{\|#1\|_{\mathrm{Lip}}}} % norme de lipschitz

%%% Fonctions et fonctionnelles 
\newcommand{\sentf}{\bH}
\newcommand{\sent}[1]{\p(){\sentf}{#1}}
\newcommand{\bentf}[1]{\bH_{#1}}
\newcommand{\bent}[2]{\p(){\bentf{#1}}{#2}}
\newcommand{\eentf}{\bN}
\newcommand{\eent}[1]{\p(){\eentf}{#1}}
\newcommand{\entf}[1]{\mathbf{Ent}_{#1}}
\newcommand{\ent}[2]{\p(){\entf{#1}}{#2}}
\newcommand{\ientf}[1]{\bH^{#1}}
\newcommand{\ient}[2]{\p(){\ientf{#1}}{#2}}
\newcommand{\rentf}{\mathbf{Ent}}
\newcommand{\rent}[2]{\p(){\rentf}{#1\,\vert\,#2}}
\newcommand{\entr}{\mathbf{Ent}_r}
\newcommand{\enef}[1]{\boldsymbol{\mathcal{E}}_{#1}}
\newcommand{\ene}[2]{\p(){\enef{#1}}{#2}}
\newcommand{\imutf}{\bI}
\newcommand{\imut}[1]{\p(){\imutf}{#1}}
\newcommand{\fcrf}{\dI}
\newcommand{\fcr}[1]{\p(){\fcrf}{#1}}
\newcommand{\fishf}{\bJ}
\newcommand{\fish}[1]{\p(){\fishf}{#1}}
\newcommand{\fishmf}{\dJ}
\newcommand{\fishm}[1]{\p(){\fishmf}{#1}}
\newcommand{\moyf}[1]{\bE_{#1}}
\newcommand{\moyp}[2]{{\moyf{#1}}{#2}}
\newcommand{\bmoy}[2]{\moyf{#1}\!\biggl(#2\biggr)}
\newcommand{\corrf}[1]{\mathbf{Cor}_{#1}}
\newcommand{\corr}[3]{\p(){\corrf{#1}}{#2,#3}}
\newcommand{\covf}[1]{\mathbf{Cov}_{#1}}
\newcommand{\cov}[3]{\p(){\covf{#1}}{#2,#3}}
\newcommand{\bcov}[3]{\covf{#1}\!\biggl(#2,#3\biggr)}
%\newcommand{\suppf}{\mathrm{supp}}
%\newcommand{\supp}[1]{\suppf\PAR{#1}}
\newcommand{\varf}[1]{\mathbf{Var}_{#1}}
%\newcommand{\var}[2]{\p(){\varf{#1}}{#2}}
\newcommand{\varp}[2]{{\varf{#1}}{#2}}
\newcommand{\bvar}[2]{\varf{#1}\!\biggl(#2\biggr)}
\newcommand{\Kf}{\bK}
\newcommand{\K}[1]{\p(){\Kf}{#1}}
\newcommand{\vrs}[1]{\mathbf{L}_{#1}}

%%% Ensembles, espaces de fonctions... 
%\newcommand{\C}[1]{\p(){\cC}{#1}}
\newcommand{\Cb}[1]{\p(){\cC_b}{#1}}
\newcommand{\Cc}[1]{\p(){\cC_c}{#1}}
\newcommand{\Cn}[2]{\p(){\cC^{#1}}{#2}}
\newcommand{\Ci}[1]{\Cn{\infty}{#1}}
\newcommand{\Cic}[1]{\p(){\cC_c^\infty}{#1}}
\newcommand{\Cnc}[2]{\p(){\cC_c^{#1}}{#2}}
\newcommand{\Cnb}[2]{\p(){\cC_b^{#1}}{#2}}
\newcommand{\Cib}[1]{\p(){\cC_b^\infty}{#1}}
\newcommand{\leb}[2]{\p(){\bL^{#1}}{#2}}
\newcommand{\lebb}[1]{\bL^{#1}}

%%% Determinant, trace...
\newcommand{\Tr}{\mathrm{Tr\,}}
\newcommand{\Det}[1]{\mathrm{Det}\,}
\newcommand{\TR}[1]{\p(){\mathrm{Tr}}{#1}}
\newcommand{\DIAG}[1]{\p(){\mathrm{Diag}}{#1}}
\newcommand{\DET}[1]{\p(){\mathrm{Det}}{#1}}
\newcommand{\SIG}[1]{\p(){\mathrm{Sign}}{#1}}
\newcommand{\ID}{\mathbf{Id}}
%\newcommand{\Id}{\mathrm{Id}}
\newcommand{\Mo}{\mathbf{1}}
\newcommand{\Vo}{\mathrm{1}}

%%% Semi-groupes, générateurs, carré-du-champ... 
\newcommand{\CD}{CD(\rho ,\infty)}
\newcommand{\GA}{\boldsymbol{\Gamma}}
\newcommand{\GD}{\GA_{\!\!{\mathbf 2}}}
\newcommand{\GI}{\bL}
\newcommand{\GR}{\nabla}
\newcommand{\GIV}{\EM{\overrightarrow{\GI}}}
\newcommand{\GIB}{\EM{\overline{\GI}}}
\newcommand{\LA}{\boldsymbol{\Delta}}
\newcommand{\ROT}{\mathbf{rot}\,}
\newcommand{\DIV}{\mathbf{div}\,}
\newcommand{\PT}[1][t]{\mathbf{P}_{\!#1}}
\newcommand{\SGf}[1]{{\mathbf P}_{#1}}
\newcommand{\SG}[2]{\p(){\SGf{\!#1}}{#2}}
\newcommand{\SGQf}[1]{{\mathbf Q}_{#1}}
\newcommand{\SGQ}[2]{\p(){\SGQf{#1}}{#2}}
\newcommand{\isopf}{\mathbf{I}}
\newcommand{\isop}[1]{\p(){\isopf}{#1}}
\newcommand{\disop}[1]{\p(){\isopf'}{#1}}
\newcommand{\ddisop}[1]{\p(){\isopf''}{#1}}
\newcommand{\HESS}[1]{\GR^2\!{#1}}
\newcommand{\Hess}[1]{{\p(){\mathrm{Hess}}{#1}}}
\newcommand{\HAM}{\bH}
\newcommand{\POT}{\bV}
\newcommand{\PF}{\bZ}
\newcommand{\DOM}{\EM{\cD_2(\GI)}}
\newcommand{\laV}{\EM{\overrightarrow{\la}}}
\newcommand{\DOML}{\EM{\cD_2\PAR{\GI}}}
\newcommand{\DOMB}{\EM{\cD_2(\GIB)}}
\newcommand{\DOMV}{\EM{\cD_2(\GIV)}}
\newcommand{\GAV}{\EM{\overrightarrow{\GA}}}
%\newcommand{\D}{\mathbf{D}}

%%% Topo
\newcommand{\inter}[1]{{\overset{\circ}{#1}}} % interior
\newcommand{\adher}[1]{{\overline{#1}}} % adherence
\newcommand{\ADH}[1]{\mathbf{adh}(#1)}
\newcommand{\INT}[1]{\mathbf{int}(#1)}

%%% Racourcis vers des noms propres 
\newcommand{\bern}{\textsc{Bernoulli}}
\newcommand{\bob}{\textsc{Bobkov}}
\newcommand{\boc}{\textsc{Bochner}}
%\newcommand{\cs}{\textsc{Cauchy}-\textsc{Schwarz}}
\newcommand{\fub}{\textsc{Fubini}}
\newcommand{\gaus}{\textsc{Gauss}}
\newcommand{\hol}{\textsc{Hölder}}
\newcommand{\jen}{\textsc{Jensen}}
\newcommand{\lap}{\textsc{Laplace}}
\newcommand{\ls}{\sobo{} logarithmique}
\newcommand{\mink}{\textsc{Minkowski}}
\newcommand{\mrkv}{\textsc{Markov}}
%\newcommand{\ou}{\textsc{Ornstein-Uhlenbeck}}
\newcommand{\poin}{\textsc{Poincaré}}
\newcommand{\pois}{\textsc{Poisson}}
\newcommand{\ric}{courbure de \textsc{Ricci}}
\newcommand{\shan}{\textsc{Shannon}}
\newcommand{\sobo}{\textsc{Sobolev}}

%%% Divers
\newcommand{\SSK}[1]{\substack{#1}}
\renewcommand{\leq}{\leqslant}
\renewcommand{\geq}{\geqslant}
\newcommand{\bs}{\EM{\backslash}} 
\newcommand{\Inf}{\boldsymbol{\inf}}
\newcommand{\Sup}{\boldsymbol{\sup}}
\newcommand{\ds}[1]{\EM{\displaystyle{#1}}}
%\newcommand{\eg}{\overset{\Delta}{=}}
\newcommand{\fdefeq}{\overset{\text{déf.}}{=}}
\newcommand{\edefeq}{\overset{\text{def.}}{=}}
\newcommand{\defeq}{:=}
\newcommand{\equ}{\; \Leftrightarrow \;}
\newcommand{\ex}{\exists \,}
\newcommand{\imp}{\Rightarrow  \;}
%\newcommand{\ssi}{{\it ssi}}
\newcommand{\tout}{\forall \,}
%\newcommand{\tq}{\,|\,}
%\newcommand{\1}{\hbox{1}\!\!\hbox{I}} %Bug avec paquet xy
\newcommand{\DSFRAC}[2]{\EM{\frac{\ds{#1}}{\ds{#2}}}}
\newcommand{\SR}[2]{\EM{\strackrel{#1}{#2}}}

%Ajout à Babel
\newcommand{\Ieme}{\up{\lowercase{ième}}\xspace}

% Here are two LaTeX macros for the word "C++". They prevent line breaks
% between the C and "++", and the first packs the two "+"s close to each
% other but the second does not. Try them both and see which one you
% like best.
%\newcommand{\CC}{C\nolinebreak\hspace{-.05em}\raisebox{.4ex}{\tiny\bf +}%
%\nolinebreak\hspace{-.10em}\raisebox{.4ex}{\tiny\bf +}}
%\def\CC{{C\nolinebreak[4]\hspace{-.05em}\raisebox{.4ex}{\tiny\bf ++}}}
% Here are two more LaTeX macros for the word "C++". They allow line
% breaks between the C and "++", which may not be desirable, but they're
% included here just in case.
%\def\CC{C\raise.22ex\hbox{{\footnotesize +}}\raise.22ex\hbox{\footnotesize +}}
%\def\CC{{C\hspace{-.05em}\raisebox{.4ex}{\tiny\bf ++}}}


\usepackage{amsmath,amsfonts,amssymb,amsthm}
%\usepackage[a4paper,portrait]{geometry}
%\geometry{left=2.5cm,right=2.5cm}
\usepackage{graphicx}
\usepackage{rotating}
\usepackage{multicol}
\usepackage{moreverb}
%\usepackage[latin1]{inputenc}
\usepackage{xspace}
%\usepackage[francais]{babel} 
\usepackage{color} 
\usepackage{listingsutf8}

\lstset{extendedchars=true,%
       basicstyle=\small,%
       commentstyle=\ttfamily\color[rgb]{0,0,0.5},%
       stringstyle=\color[rgb]{0,0.5,0},%
       frame=tblr,frameround=tttt,%
        extendedchars=\true,       %   Ces deux lignes sont les miennes pour faire fonctionner le UTF8
       inputencoding=utf8/latin1   %
       %labelstyle=\tiny, labelstep=1,labelsep=10pt
       }

%\ifx\pdfoutput\undefined
% % no problemo.
%\else % compatibilité listings/hyperref.
% \makeatletter
% \providecommand*{\toclevel@lstlisting}{0}
% \makeatother
%\fi

%\input{macros}
%\THMFR
%
%\newcommand{\MFILE}[1]{
%\begin{lstinputlisting}[language=Matlab]{#1.m}%
%\label{co:#1}% 
%\end{lstinputlisting}}
%
\newcommand{\OC}{Octave}
%\newcommand{\SL}{Scilab}
%\newcommand{\MP}{Maple}
\newcommand{\ML}{Matlab}
%\newcommand{\MU}{Mupad}
\newcommand{\SB}{Stixbox}
%
\newcommand{\FIG}[4]{%  fname scale floatparams caption
 \begin{figure}[#3]%
  \begin{center}%
  \includegraphics[scale=#2]{#1}%
  \caption{#4}%
  \label{fi:#1}%
  \end{center}%
 \end{figure}
}
%
%% Parce que j'ai l'impression que ces figures font mourir la conversion en pdf
%\renewcommand{\FIG}[4]{}


\newcommand{\TextePourISBN}{

\begin{center}

        \ifpdf
            \includegraphics{CreativeCommons.png}
        \else
            \includegraphics{CreativeCommons.eps}
        \fi

        \emph{Mes notes de mathématique} de \href{http://student.ulb.ac.be/~lclaesse/}{Laurent Claessens} est mis à disposition selon les termes de la licence \href{http://creativecommons.org/licenses/by-nc-nd/2.0/fr/}{licence Creative Commons Attribution - Pas d’Utilisation Commerciale - Pas de Modification 2.0 France.}.

        Une version plus à jour de ce texte, publiée sous la \href{http://fr.wikipedia.org/wiki/FDL}{licence GNU-FDL} est disponible sur \href{http://student.ulb.ac.be/~lclaesse/}{le site de l'auteur}.

        ISBN : Y-Y-YYYYYYY-Y-Y
\end{center}

}



\makeindex
\makenomenclature


\begin{document}
\corrPosition{1}


    \thispagestyle{empty}
\begin{center}
  \begin{minipage}{15cm}
    \hrule\par
    \vspace{2mm}
    \begin{center}
    \Huge \bfseries Mes notes de mathématique \par
    \end{center}
    \hrule\par
  \end{minipage}
\end{center}

\vspace{2cm}

\begin{center}
    Laurent \textsc{Claessens}\\
    Université de Franche-Comté\\
    \today\\
    \url{http://student.ulb.ac.be/~lclaesse/mes_notes.pdf}\\
    \url{https://www.gitorious.org/mes-notes-de-math-matique}
\end{center}

\vfill

\begin{center}

           \ifpdf
            \includegraphics[width=1cm]{Copyleft.svg}
        \else
            \includegraphics[width=1cm]{Copyleft.eps}
        \fi

Copyright (c) 2011  Laurent Claessens,

Permission is granted to copy, distribute and/or modify this document under the terms of the \href{http://www.gnu.org/licenses/fdl-1.3.html}{GNU Free Documentation License}, Version 1.3 or any later version published by the Free Software Foundation; with no Invariant Sections, no Front-Cover Texts, and no Back-Cover Texts. A copy of the license is included in the section entitled ``GNU Free Documentation License''.

% Ajouter ici l'ISBN. Pour les révisions, mettre un nouvel ISBN et indiquer que c'est une révision.
% Pour l'ISBN:
% Copier tout dans un nouveau répertoire
% Créer une nouvelle branche git
% Coder en dur la date (càd enlever \today)
% Changer la licence vers non modifiable, copiable, et ajouter un lien vers la
%   version modifiable.

% http://www.bnf.fr/fr/professionnels/s_informer_obtenir_isbn/s.qu_est_ce_que_isbn.html

% TODO: écrire le README.txt

\end{center}

\clearpage


    %+++++++++++++++++++++++++++++++++++++++++++++++++++++++++++++++++++++++++++++++++++++++++++++++++++++++++++++++++++++++++++
\section{Avertissement}
%+++++++++++++++++++++++++++++++++++++++++++++++++++++++++++++++++++++++++++++++++++++++++++++++++++++++++++++++++++++++++++

Ceci sont des notes «prises au vol» du cours de modélisation pour l'agrégation. Aucune garantie. Merci de me signaler toute faute ou remarque.


\tableofcontents

\part{Mes notes d'agrégation}

\chapter{Théorie des groupes}
Nous allons suivre dans un premier temps \cite{Kropholler}.


\begin{definition}
    Soit \( G\) un groupe. Le \defe{centralisateur}{centralisateur} de \( H\subset G\) est 
    \begin{equation}
        \mZ_G(H)=\{ g\in G\tq hg=gh\,\forall h\in h\}
    \end{equation}
    Si \( H\) est un sous groupe, son \defe{centralisateur}{centralisateur} est
    \begin{equation}
        N_G(H)=\{ g\in G\tq gH=Hg \}.
    \end{equation}
\end{definition}

\begin{definition}
    Un sous groupe \( N\) de \( G\) est \defe{normal}{normal!sous groupe} ou \defe{distingué}{distingué} si pour tout \( g\in G\) et pour tout \( n\in N\), \( gng^{-1}\in N\). Autrement dit lorsque \( gNg^{-1}\subset N\). Lorsque \( N\) est normal dans \( G\) nous noterons \( N\normal G\)\nomenclature[]{\(N \normal G\)}{\( N\) est normal dans \( G\)}.
\end{definition}

\begin{proposition}
    Soit \( N\) un sous groupe de \( G\). Les propriétés suivantes sont équivalentes :
    \begin{enumerate}
        \item
            \( gNg^{-1}\subseteq N\) pour tout \( g\in G\),
        \item
            \( gNg^{-1}= N\) pour tout \( g\in G\),
        \item
            \( gN=Ng\) pour tout \( g\in G\),
        \item
            \( N\) est une union de classes de conjugaison de \( G\),
        \item
            \( N\) est normal.
    \end{enumerate}
\end{proposition}

\begin{definition}
    Soit \( g\in G\) et \( n\in \eZ\). Nous définissons \( g^n\) par
    \begin{enumerate}
        \item
            \( g^0=e\) et \( g^n=gg^{n-1}\) si \( n\) est positif.
        \item
            si \( n<0\), nous posons \( g^n=(g^{-1})^{-n}\).
    \end{enumerate}
\end{definition}

\begin{definition}
    Si \( G\) est un groupe, l'\defe{ordre}{ordre!d'un groupe} est la cardinalité de \( G\) et est noté \( | G |\). L'\defe{ordre}{ordre!élément} d'un élément \( g\) de \( G\) est le naturel
    \begin{equation}
        \min\{ n\in\eN\tq g^n=e \}.
    \end{equation}
    Si le minimum n'existe pas, nous disons que l'ordre de \( g\) est infini.
\end{definition}

%+++++++++++++++++++++++++++++++++++++++++++++++++++++++++++++++++++++++++++++++++++++++++++++++++++++++++++++++++++++++++++
\section{Théorèmes d'isomorphismes}
%+++++++++++++++++++++++++++++++++++++++++++++++++++++++++++++++++++++++++++++++++++++++++++++++++++++++++++++++++++++++++++

Si \( G\) est un groupe et si \( N\) est un sous groupe normal, alors l'ensemble \( G/N\) a une structure de groupe et la projection canonique \( \pi\colon G\to G/N\) est un homomorphisme surjectif de noyau~\( N\).

\begin{theorem}[premier théorème d'isomorphisme]        \label{ThoPremierthoisomo}
    Soit \( \theta\colon G\to H\) un homomorphisme de groupe. Alors
    \begin{enumerate}
        \item
            \( \Kernel\theta\) est normal dans \( G\),
        \item
            \( \Image \theta\) est un sous groupe de \( H\)
        \item
            nous avons un isomorphisme naturel
            \begin{equation}
                G/\Kernel\theta\simeq \Image\theta
            \end{equation}
    \end{enumerate}
\end{theorem}

\begin{proof}
    \begin{enumerate}
        \item
        \item
        \item
            Si \( [g]\) représente la classe de \( g\) dans \( G/\Kernel\theta\), l'isomorphisme est donné par \( \varphi[g]=\theta(g)\).
    \end{enumerate}
\end{proof}


\begin{theorem}[Deuxième théorème d'isomorphisme]
    Soient \( H\) et \( N\) deux sous groupes de \( G\) et supposons que \( N\) soit normal. Alors
    \begin{enumerate}
        \item
            \( NH=HN\) est un sous groupe
        \item
            \( N\normal NH\)
        \item
            \( N\cap H\normal H\)
        \item
            nous avons un isomorphisme
            \begin{equation}
                NH/N\simeq H/H\cap N.
            \end{equation}
    \end{enumerate}
\end{theorem}

\begin{proof}
    \begin{enumerate}
        \item
        \item
        \item
        \item
            Il faut d'abord remarquer que \( H\) et \( N\) étant des groupes et le produit \( NH\) étant un groupe, nous avons \( NH=HN\). Soit le morphisme injectif
            \begin{equation}
                \begin{aligned}
                    j\colon H&\to HN \\
                    h&\mapsto h
                \end{aligned}
            \end{equation}
            et la surjection canonique
            \begin{equation}
                \sigma\colon HN\to HN/N 
            \end{equation}
            Nous considérons ensuite l'application composée
            \begin{equation}
                \begin{aligned}
                    f\colon H&\to HN/N \\
                    h&\mapsto hN. 
                \end{aligned}
            \end{equation}
            L'application \( f\) est surjective parce que l'élément \( hnN\in HN/N\) est l'image de \( h\), étant donné que \( hnN=hN\).

            Le noyau de \( f\) est \( \Kernel f=H\cap N\). En effet si \( a\in H\cap N\), nous avons \( f(a)=\sigma(a)\in K\). Par conséquent \( H\cap N\subset \Kernel f\). D'autre part si \( h\in H\) vérifie \( h\in\Kernel f\), alors \( f(h)=hN=N\), ce qui est uniquement possible si \( h\in N\).

            Le premier théorème d'isomorphisme implique alors que \( H/\Kernel f\simeq \Image f\), c'est à dire
            \begin{equation}
                H/N\cap H\simeq HN/N.
            \end{equation}
    \end{enumerate}
\end{proof}

\begin{theorem}[Troisième théorème d'isomorphisme]
    Soient \( N\) et \( M\) deux sous groupes normaux de \( G\) avec \( M\subset N\). Alors \( N/M\) est normal dans \( G/M\) et
    \begin{equation}
        (G/M)/(N/M)\simeq G/N.
    \end{equation}
\end{theorem}

\begin{proof}
    Afin de montrer que \( N/M\) est normal dans \( G/M\), nous considérons \( g\in G\), \( nM\in N/M\) et nous calculons
    \begin{equation}
        gnMg^{-1}=gng^{-1}\underbrace{gMg^{-1}}_{=M}=\underbrace{gng^{-1}}_{\in N}M\in N/M.
    \end{equation}

    Pour prouver l'isomorphisme nous considérons le morphisme
    \begin{equation}
        \begin{aligned}
            \varphi\colon G/M&\to G/N \\
            gM&\mapsto gN. 
        \end{aligned}
    \end{equation}
    C'est surjectif et le noyau est \( N/M\) parce que \( \varphi(gM)=N\) uniquement si \( g\in N\). Nous pouvons appliquer le premier théorème d'isomorphisme à \( \varphi\) en écrivant
    \begin{equation}
        (G/M)/\Kernel \varphi\simeq\Image \varphi,
    \end{equation}
    c'est à dire
    \begin{equation}
        (G/M)/(N/M)\simeq G/N.
    \end{equation}
\end{proof}

%+++++++++++++++++++++++++++++++++++++++++++++++++++++++++++++++++++++++++++++++++++++++++++++++++++++++++++++++++++++++++++
\section{Indice d'un sous groupe}
%+++++++++++++++++++++++++++++++++++++++++++++++++++++++++++++++++++++++++++++++++++++++++++++++++++++++++++++++++++++++++++

Soit \( G\) un groupe fini et \( H\), un sous groupe. L'\defe{indice}{indice} de \( H\) dans \( G\) est le nombre \( | G |/| H |\), souvent noté \( | G:H |\). Le théorème de Lagrange dira en particulier que l'indice est toujours un nombre entier.

\begin{theorem}[Théorème de Lagrange]\index{théorème!Lagrange}
    Soit \( H\) un sous groupe du groupe fini \( G\). Alors \( | H |\) divise \( | G |\) et les trois nombres suivants sont égaux :
    \begin{enumerate}
        \item
            le nombre de classes de \( H\) à gauche,
        \item
            le nombre de classes de \( H\) à droite,
        \item
            l'indice de \( H\) dans \( G\).
    \end{enumerate}
\end{theorem}

\begin{proof}
    Nous commençons par montrer que les classes de \( H\) ont toutes les même nombre d'éléments que \( H\). En effet pour chaque \( g\in G\) nous avons la bijection
    \begin{equation}
        \begin{aligned}
            \varphi\colon H&\to gH \\
            h&\mapsto gh. 
        \end{aligned}
    \end{equation}
    L'injectivité de \( \varphi\) est le fait que \( gh=gh'\) implique \( h=h'\). La surjectivité est par définition de la classe. 

    Les classes à gauche formant une partition de \( G\), le cardinal de \( G\) est le produit de la taille des classes par le nombre de classes :
    \begin{equation}
        | G |=| H |\cdot\text{nombre de classes}.
    \end{equation}
    En particulier nous voyons que \( | H |\) divise \( | G |\).
\end{proof}

%+++++++++++++++++++++++++++++++++++++++++++++++++++++++++++++++++++++++++++++++++++++++++++++++++++++++++++++++++++++++++++
\section{Le groupe des entiers}
%+++++++++++++++++++++++++++++++++++++++++++++++++++++++++++++++++++++++++++++++++++++++++++++++++++++++++++++++++++++++++++

\begin{definition}
    Soit \( G\), un groupe et \( A\) une partie de \( G\). Nous notons \( \gr(A)\)\nomenclature[A]{\( \gr\)}{groupe engendré} l'intersection de tous les sous groupes de \( G\) contenant \( A\). C'est le plus petit (pour l'inclusion) groupe de \( G\) contenant \( A\). Si \( \gr(A)=G\), alors nous disons que \( A\) est une \defe{partie génératrice}{partie génératrice} le groupe \( G\).

    Un groupe est \defe{monogène}{monogène} si il a une partie génératrice réduite à un seul élément.

    Un élément \( a\in G\) est un \defe{générateur}{générateur} de \( G\) si tous les éléments de \( G\) s'écrivent sous la forme \( a^n\) pour un certain \( n\in\eZ\). Un groupe fini et monogène est dit \defe{cyclique}{cyclique}.
\end{definition}

\begin{example}
    Soit \( G\) le groupe des rotations d'angle \( 2k\pi/10\). La rotation d'angle \( \pi/2\) n'est pas génératrice parce qu'elle n'engendre que \( \pi/2\), \( \pi\),\( 3\pi/2\) et l'identité. La rotation d'angle \( 2\pi/10\) par contre est génératrice.
\end{example}


Nous notons \( \eN_{k}=\{ 0,1,\ldots, k \}\subset\eN\)\nomenclature[A]{\( \eN_k\)}{Un sous ensemble des naturels}.
 
\begin{proposition}
    Une partie \( H\) du groupe \( \eZ\) est un sous groupe si et seulement si il existe \( n\in\eN\) tel que \( H=n\eZ\).
\end{proposition}

\begin{proof}
    Soit \( H\neq\{ 0 \}\) un sous groupe de \( \eZ\). L'ensemble \( H\cap\eN^*\) contient un élément minimum que nous notons \( n\). Nous avons certainement \( n\eZ\subset H\) parce que \( H\) est un groupe (donc \( n+n\) et \( -n\) sont dans \( H\) dès que \( n\) est dans \( H\)). Nous devons prouver que \( H\subset n\eZ\).

    Si \( x\in H\), il existe \( q\in\eZ\) et \( r\in\eN_{n-1}\) tels que \( x=nq+r\).

    \begin{probleme}
        Justification par la division euclidienne à venir.
    \end{probleme}
    Nous savons déjà que \( nq\in H\), donc \( r\in H\) parce que \( x\) est également dans \( H\). Mais nous avions décidé que \( n\) serait le plus petit naturel de \( H\cap\eN^*\). Par conséquent \( r=0\) et \( x=nq\in n\eZ\).

\end{proof}

Notons que si un sous groupe \( H\) de \( \eZ\) est donné, alors le nombre \( n\) tel que \( H=n\eZ\) est unique. En effet si \( n\eZ=m\eZ\) nous avons que \( n\) divise \( m\) (parce que \( m\in m\eZ\subset n\eZ\)) et que \( m\) divise \( n\) parce que \( n\in m\eZ\). Par conséquent \( n=m\).

%---------------------------------------------------------------------------------------------------------------------------
\subsection{Division euclidienne}
%---------------------------------------------------------------------------------------------------------------------------

\begin{theorem}[Division euclidienne]     \label{ThoDivisEuclide}
    Soient \( a\in\eZ\) et \( b\in\eN^*\). Il existe un unique couple \( (q,r)\in\eZ\times\eN\) tel que
    \begin{equation}
        a=bq+r
    \end{equation}
    avec \( 0\leq r<b\).
\end{theorem}

L'opération \( (a,b)\mapsto(a,r)\) donnée par le théorème \ref{ThoDivisEuclide} est la \defe{division euclidienne}{division!euclidienne}. Le nombre \( q\) est le \defe{quotient}{quotient} et \( r\) est le \defe{reste}{reste} de la division de \( a\) par \( b\).

%---------------------------------------------------------------------------------------------------------------------------
\subsection{pgcd et ppcm}
%---------------------------------------------------------------------------------------------------------------------------

Soient \( p,q\in\eZ^*\). Les ensembles \( p\eZ\cap q\eZ\) et \( p\eZ+q\eZ\) sont des sous groupes de \( \eZ\). Par conséquent il existe des entiers \( \ppcm(p,q)\) et \( \pgcd(p,q)\) tels que
\begin{subequations}
    \begin{align}
        p\eZ\cap q\eZ&=\ppcm(p,q)\eZ\\
        p\eZ + q\eZ&=\pgcd(p,q)\eZ.
    \end{align}
\end{subequations}

Si \( \pgcd(p,q)=1\), nous disons que \( p\) et \( q\) sont \defe{premiers entre eux}{premiers!deux nombres}.

\begin{theorem}[Théorème de Bezout]
    Deux entiers non nuls \( p,q\in\eZ^*\) sont premiers entre eux si et seulement si il existe \( m,n\in\eZ\) tels que 
    \begin{equation}
        pm+qn=1.
    \end{equation}
\end{theorem}

\begin{proposition}
    soient \( a,b\in\eZ^*\). Si
    \begin{equation}
        \begin{aligned}[]
            a&=\epsilon\prod_{\text{\( p\) premiers}}p^{\mu(p)}&b&=\epsilon'\prod_{\text{\( p\) premier}}p^{\nu(p)},
        \end{aligned}
    \end{equation}
    alors
    \begin{subequations}
        \begin{align}
            \pgcd(a,b)&=\prod p^{\min\{ \mu(p),\nu(p) \}}\\
            \ppcm(a,b)&=\prod p^{\max\{ \mu(p),\nu(p) \}}
        \end{align}
    \end{subequations}    
\end{proposition}

Cette proposition implique que \( m\leq p^n\) et \( \pgcd(m,p^n)\) si et seulement si \( m=qp\) avec \( q\leq p^{n-1}\).

%---------------------------------------------------------------------------------------------------------------------------
\subsection{Quotients}
%---------------------------------------------------------------------------------------------------------------------------

Nous notons \( \eZ_n=\eZ/n\eZ\)\nomenclature[A]{\( \eZ_n\)}{groupe des entiers modulo \( n\)} le groupe des entiers modulo \( n\).

\begin{proposition}
    Soit \( n\in\eN\). Le groupe \( \eZ_n\) est monogène. Si \( n\neq 0\), le groupe \( \eZ_n\) est cyclique d'ordre \( n\).
\end{proposition}

\begin{proof}
    Nous considérons la surjection canonique \( \mu\colon \eZ\to \eZ_n\). Si \( a\in\eZ\), alors \( \mu(a)=a\mu(1)\). Par conséquent \( \eZ_n=\gr\big( \mu(1) \big)\) parce que tout groupe contenant \( \mu(1)\) contient tous les multiples de \( \mu(1)\), et par conséquent contient \( \mu(\eZ)=\eZ_n\).

    Soit \( x\in\eZ_n\) et considérons \( x_0\), le plus petit naturel représentant \( x\). Nous notons \( x=[x_0]\). Le théorème de la division euclidienne \ref{ThoDivisEuclide} donne l'existence de \( q\) et \( r\) avec \( 0\leq r<n\) et \( q\geq 0\) tels que
    \begin{equation}
        x_0=nq+r.
    \end{equation}
    Nous avons \( [x_0]=[r]=\mu(r)\) parce que \( x_0-r\) est un multiple de \( n\). Nous avons donc \( [x_0]\in\mu(\eN_{n-1})\). Par conséquent
    \begin{equation}
        \eZ_n=\mu(\eZ)=\mu(\eN_{n-1}).
    \end{equation}
    La restriction \( \mu\colon \eN_{n-1}\to \eZ_n\) est donc surjective. Montrons qu'elle est également injective. Si \( \mu(x_0)=\mu(x_1)\), alors \( x_1=x_0+kn\). Si nous supposons que \( x_1>x_0\), alors \( k>0\) et si \( x_0\in\eN_{n-1}\), alors \( x_1>n-1\).

    L'ordre de \( \eZ_n\) est donc le même que le cardinal de \( \eN_{n-1}\), c'est à dire \( n\). Le groupe \( \eZ_n\) est donc fini, d'ordre \( n\) et monogène (parce que \( \eZ_n=\gr(\mu(1))\)). Il est donc cyclique.
\end{proof}
 

Pour \( n\in\eN^+\) nous introduisons l'ensemble
\begin{equation}
    P_n=\{ x\in\{ 1,\ldots,n \}\tq\pgcd(x,n)=1 \}.
\end{equation}
C'est l'ensemble des entiers inférieurs à \( n\), premiers avec \( n\). La fonction \( \varphi\) donnée par
\begin{equation}
    \varphi(n)=\Card(P_n)
\end{equation}
est l'\defe{indicatrice d'Euler}{indicatrice d'Euler}\index{Euler!indicatrice}. Nous avons par exemple
\begin{subequations}
    \begin{align}
        P_8&=\{ 1,3,5,7 \}\\
        P_7&=\{ 1,2,3,4,5,6 \}.
    \end{align}
\end{subequations}
Si \( p\) est un nombre premier, alors \( \varphi(p)=p-1\).

\begin{lemma}
    Si \( p\) est un nombre premier, alors \( \varphi(p^n)=p^n-p^{n-1}\).
\end{lemma}

\begin{proof}
    Les éléments de \( \{ 1,\ldots,p^n \}\) qui ont un \( \pgcd\) différent de \( 1\) avec \( p^n\) sont des nombres qui s'écrivent sous la forme \( qp\) avec \( q\leq p^{n-1}\). Il y a évidemment \( p^{n-1}\) tels nombres.

    Par conséquent le cardinal de \( P_{p^n}\) est \( \varphi(p^{n})=p^n-p^{n-1}\).
\end{proof}

\begin{proposition}     \label{PropZnmuphiGensn}<++>
    Soit \( n\in\eN^*\) et \( \mu\colon \eZ\to \eZ_n\) la projection canonique. Les générateurs de \( \eZ_n\) sont les \( \mu(x)\) avec \( x\in P_n\). En particulier \( \eZ_n\) est un groupe contenant \( \varphi(n)\) générateurs.
\end{proposition}

\begin{proof}
    La classe de \( 1\) dans \( \eZ_n\) est \( \mu(1)\), par conséquent nous avons \( \gr\big( \mu(1) \big)=\eZ_n\). L'élément \( \mu(x)\) sera générateur si et seulement si il génère \( \mu(1)\), c'est à dire si il existe \( p\) tel que \( p\mu(x)=\mu(1)\). Cette dernière égalité étant une égalité de classes dans \( \eZ_n\), elle sera vraie si et seulement si il existe \( q\) tel que
    \begin{equation}
        px+qn=1.
    \end{equation}
    Cela signifie entre autres que \( x\eZ+n\eZ=\eZ\), c'est à dire que \( \pgcd(x,n)=1\) et que \( x\in P_n\).
\end{proof}


\begin{theorem}
    Un groupe monogène est abélien. Plus précisément,
    \begin{enumerate}
        \item
            un groupe monogène infini est isomorphe à \( \eZ\),
        \item
            un groupe monogène fini est isomorphe à \( \eZ_n\) pour un certain \( n\).
    \end{enumerate}
    Un groupe monogène d'ordre \( n\) possède \( \varphi(n)\) générateurs.
\end{theorem}

\begin{proof}
    Le groupe est abélien parce que $g=a^n$, \( g'=a^{n'}\) implique \( gg'=q^{n+n'}=g'g\). Nous considérons un générateur \( a\) de \( G\) (qui existe parce que $G$ est monogène) et le morphisme surjectif
    \begin{equation}
        \begin{aligned}
            f\colon \eZ&\to G \\
            p&\mapsto a^p. 
        \end{aligned}
    \end{equation}
    Si \( G\) est infini, alors \( f\) est injective parce que si \( a^n=a^{n'}\), alors \( a^{n-n'}=e\), ce qui rendrait \( G\) cyclique et par conséquent non infini. Nous concluons que si \( G\) est infini, alors \( f\) est une bijection et donc un isomorphisme \( \eZ\simeq G\).

    Si \( G\) est fini, alors \( f\) n'est pas injective et a un noyau \( \ker f\). Étant donné que \( \ker f\) est un sous groupe de \( G\), il existe un (unique) \( n\) tel que \( \ker f=n\eZ\) et le premier théorème d'isomorphisme (théorème \ref{ThoPremierthoisomo}) nous indique que
    \begin{equation}
        \eZ/\ker f=\eZ/n\eZ=\Image f=G.
    \end{equation}
    Dans ce cas, le fait qu'un groupe monogène d'ordre \( n\) possède \( \varphi(n)\) générateurs est le contenu de la proposition \ref{PropZnmuphiGensn}.
\end{proof}

%+++++++++++++++++++++++++++++++++++++++++++++++++++++++++++++++++++++++++++++++++++++++++++++++++++++++++++++++++++++++++++
\section{Théorèmes de Sylow}
%+++++++++++++++++++++++++++++++++++++++++++++++++++++++++++++++++++++++++++++++++++++++++++++++++++++++++++++++++++++++++++

\begin{definition}
    Soit \( p\) un nombre premier. Un \( p\)-groupe est un groupe dont tous les éléments sont d'ordre \( p^m\) pour un certain \( m\) (dépendant de l'élément).

    Soit \( G\) un groupe fini de cardinal \( N\) et \( p\), un diviseur premier de $| G |$. Un \( p\)-Sylow dans \( G\) est une \( p\)-sous groupe d'ordre \( p^n\) où \( p^n\) est la plus grande puissance de \( p\) divisant \( | G |\).
\end{definition}

Si \( S\) est un \( p\)-Sylow, alors \( p\) ne divise pas \( | G:S |=| G |/| S |\).


\chapter{Anneaux et corps}
%+++++++++++++++++++++++++++++++++++++++++++++++++++++++++++++++++++++++++++++++++++++++++++++++++++++++++++++++++++++++++++
\section{Généralités}
%+++++++++++++++++++++++++++++++++++++++++++++++++++++++++++++++++++++++++++++++++++++++++++++++++++++++++++++++++++++++++++

Source : \cite{Tauvel}.

\begin{definition}
    Un \defe{anneau}{anneau} est un triple \( (A,+,\cdot)\) avec les conditions
    \begin{enumerate}
        \item
            \( (A,+)\) est un groupe abélien. Nous notons \( 0\) le neutre.
        \item
            La multiplication est associative et nous notons \( 1\) le neutre
        \item
            La multiplication est distributive par rapport à l'addition.
    \end{enumerate}
\end{definition}

\begin{remark}
    Un anneau est ce qu'on appelle «\emph{ring}» en anglais.
\end{remark}

Soit \( X\) un ensemble et $A$ un anneau. Nous considérons \( \Fun(X,A)\)\nomenclature[A]{\( \Fun(X,Y)\)}{les applications de \( X\) vers \( Y\)} l'ensemble des applications \( X\to A\). Cet ensemble devient un anneau avec les définitions
\begin{subequations}
    \begin{align}
        (f+g)(x)=f(x)+g(x)\\
        (fg)(x)=f(x)g(x).
    \end{align}
\end{subequations}
Cela est la \defe{structure canonique}{structure d'anneau canonique} d'anneau sur \( \Fun(X,A)\).

Le \defe{centralisateur}{centralisateur} de \( x\in A\) dans \( A\) est l'ensemble
\begin{equation}
    \{ y\in A\tq xy=yx \},
\end{equation}
le \defe{centre}{centre} de \( A\) est
\begin{equation}
    \{ y\in A\tq xy=yx\forall x\in A \}.
\end{equation}
Un élément \( a\neq 0\) est un \defe{diviseur de zéro à gauche}{diviseur!de zéro} si il existe \( x\neq 0\) tel que $xa=0$. L'élément \( a\) est un diviseur de zéro \defe{à droite}{diviseur!de zéro à droite} si il existe \( b\) tel que \( ab=0\). Un anneau est \defe{intègre}{intègre!anneau}\index{anneau!intègre} si il est non nul et ne possède pas de diviseurs de zéro.

\begin{example}
    L'ensemble \( \eZ\) avec les opérations usuelles est un anneau intègre.
\end{example}

Un élément \( a\in A\) est \defe{régulier à droite}{régulier à droite} \( ba=0\) implique \( b=0\). Il est régulier ) gauche si \( ab=0\) implique \( b=0\).

L'ensemble \( U(A)\)\nomenclature[A]{\( U(A)\)}{ensemble des inversibles} des éléments inversibles de \( A\) est un groupe pour la multiplication. Nous notons \( A^*=A\setminus\{ 0 \}\).

\begin{lemma}
    Si \( a\) et \( b\) commutent, nous avons la formule
    \begin{equation}        \label{Eqarpurmkbk}
        a^{r+1}-b^{r+1}=(a-b)(\sum_{k=0}^ra^{r-k}b^k).
    \end{equation}
\end{lemma}

\begin{proposition}
    Si \( a\) est un élément nilpotent de l'anneau \( A\), alors \( 1-a\) est inversible. Si \( a\) est nilpotent non nul, alors il est diviseur de zéro.
\end{proposition}

\begin{proof}
    Soit \( n\) le minimum tel que \( a^n=0\). En vertu de la formule \eqref{Eqarpurmkbk} nous avons
    \begin{equation}
        1=1-a^n=(1-a)(1+a+\ldots+a^{n-1})=(1+a+\ldots+a^{n-1})(1-a).
    \end{equation}
    La somme \( 1+a+\ldots+a^{n-1}\) est donc un inverse de \( (1-a)\).
\end{proof}

\begin{definition}
    Si \( A\) et \( B\) sont des anneaux, un \defe{morphisme}{morphisme!d'anneaux} est une application \( f\colon A\to B\) telle que pour tout \( x,y\in A\) nous ayons
    \begin{enumerate}
        \item
            \( f(x+y)=f(x)+f(y)\)
        \item
            \( f(xy)=f(x)f(y)\)
        \item
            \( f(1)=1\)
    \end{enumerate}
\end{definition}

Si \( f\) est un morphisme, nous avons \( f(0)=0\) et \( f(x)^{-1}=f(x^{-1})\).

%+++++++++++++++++++++++++++++++++++++++++++++++++++++++++++++++++++++++++++++++++++++++++++++++++++++++++++++++++++++++++++
\section{Idéaux dans des anneaux}
%+++++++++++++++++++++++++++++++++++++++++++++++++++++++++++++++++++++++++++++++++++++++++++++++++++++++++++++++++++++++++++

\begin{definition}
    Un sous ensemble \( B\subset A\) d'un anneau est un \defe{sous anneau}{sous anneau} si
    \begin{enumerate}
        \item
            \( 1\in B\)
        \item
            \( B\) est un sous groupe pour l'addition
        \item
            \( B\) est stable pour la multiplication.
    \end{enumerate}
    Un sous ensemble \( I\subset A\) est un \defe{idéal}{idéal!dans un anneau} à gauche si
    \begin{enumerate}
        \item
            \( I\) est un sous groupe pour l'addition
        \item
            si \( x\in I\) et \( a\in A\), alors \( ax\in I\).
    \end{enumerate}
\end{definition}

Lorsqu'un ensemble est idéal à gauche et à droite, nous disons que c'est un \defe{idéal bilatère}{idéal!bilatère}. Lorsque nous parlons d'idéal sans précisions, nous parlons d'idéal bilatère.

\begin{remark}
    Un idéal n'est pas toujours un anneau parce que l'identité pourrait manquer. Un idéal qui contient l'identité est l'anneau complet.
\end{remark}

\begin{example}
    L'ensemble \( 2\eZ\) est un idéal de \( \eZ\). Tous les idéaux de \( \eZ\) sont de la forme \( n\eZ\). En effet en vertu de la proposition \ref{PropSsgpZestnZ}, les seule sous groupes de \( \eZ\) (en tant que groupe additif) sont les \( n\eZ\).
\end{example}

Soit \( A\), un anneau, \( I\) un idéal bilatère de \( A\). Nous considérons la relation d'équivalence \( x\sim y\) si et seulement si \( x-y\in I\). Dans ce cas, le quotient
\begin{equation}
    A/\sim=A/I
\end{equation}
est un anneau appelé \defe{anneau quotient}{anneau!quotient par un idéal}. La surjection \( A\to A/I\) est un morphisme.

\begin{proposition}
    Soient \( A\) et \( B\) des anneaux et un homomorphisme \( f\colon A\to B\). Nous considérons l'injection canonique \( j\colon f(A)\to B\) et la surjection canonique \( \phi\colon A\to A/\ker f\). Alors il existe un unique isomorphisme
    \begin{equation}
        \tilde f \colon A/\ker f\to f(A)
    \end{equation}
    tel que \( f=j\circ\tilde f\circ\phi\).

    \begin{equation}
        \xymatrix{%
        A \ar[r]^{f}\ar[d]_{\phi}        &   B\ar[d]^{j}\\
           A/\ker f \ar[r]_{\tilde f}   &   f(A)\subset B
           }
    \end{equation}
\end{proposition}

\begin{proposition}     \label{PropIJJIdsousphi}
    Soit \( I\), un idéal de \( A\) et \( \phi\colon A\to A/I\) la surjection canonique. Les idéaux de \( A/I\) sont les \( \phi(J)\) où \( J\) est un idéal de \( A\) contenant \( I\)
\end{proposition}

\begin{proof}
    Si \( I\subset J\) et si \( J \) est un idéal de \( A\), alors \( \phi(J)\) est un idéal dans \( A/I\). En effet un élément de \( \phi(J)\) est de la forme \( \phi(j)\) et un élément de \( A/I\) est de la forme \( \phi(i)\). Leur produit vaut
    \begin{equation}
        \phi(i)\phi(j)=\phi(ij)\in\phi(J).
    \end{equation}
    
    Soit maintenant \( K\), un idéal dans \( A/I\). Soit \( J=\phi^{-1}(K)\). Étant donné qu'un idéal doit contenir \( 0\) (parce qu'un idéal est un groupe pour l'addition), \( [0]\in K\) et par conséquent \( I\subset\phi^{-1}(K)\).
\end{proof}

\begin{corollary}
    Les quotients de \( \eZ\) sont \( \eZ_n\)
\end{corollary}

\begin{proof}
    Nous avons déjà vu que les seuls idéaux de \( \eZ\) sont les \( n\eZ\).
\end{proof}

\begin{proposition}     \label{PropZpintssiprempUzn}
    Soit \( n\geq 2\) un entier et \( \phi\colon \eZ\to \eZ_n\), la surjection canonique. Nous noterons \( \tilde a=\phi(a)\). Alors
    \begin{equation}
        U(\eZ_n)=\phi(P_n)=\{ \tilde x\tq 0\leq x\leq n\tq\pgcd(x,n)=1 \}.
    \end{equation}
    où \( P_n\) est l'ensemble décrit par l'équation \eqref{EqDefPnEntierldeost}. En particulier, \( \Card\big( U(\eZ_n) \big)=\varphi(n)\).

\end{proposition}

\begin{proof}
    Soit \( 0\leq x\leq n\) tel que \( \pgcd(x,n)=1\). Il existe donc \( p,q\in\eZ\) tels que \( px+qn=1\). En passant aux classes,
    \begin{equation}
        \tilde p\tilde x=\tilde 1,
    \end{equation}
    donc \( \tilde p\) est l'inverse de \( \tilde x\). Cela prouve que \( \phi(P_n)\subset U(\eZ_n)\).

    Nous prouvons maintenant l'inclusion inverse. Soit \( \tilde x\) et \( \tilde y\) inverses l'un de l'autre : $\tilde x\tilde y=\tilde 1$. Il existe donc \( q\in\eZ\) tel que \( xy-qn=1\), ce qui prouve que \( \pgcd(x,n)=1\).

\end{proof}

\begin{corollary}   \label{CorZnInternprem}
    L'anneau \( \eZ_n\) est intègre si et seulement si \( n\) est premier.
\end{corollary}

\begin{proof}
    Si \( n\) est premier, tous les éléments de \( \eZ_n\) sont inversibles parce que tous les éléments rentrent dans \( \phi(P_n)\). Donc \( \eZ_n\) est intègre.

    Si \( n\) n'est pas premier, il existe \( p,q\in\eN^*\) tels que \( pq=n\). Dans ce cas au niveau des classes nous avons \( \tilde p\tilde q=0\) avec \( \tilde p\neq 0\neq\tilde q\), ce qui montre que \( \eZ_n\) a des diviseurs de zéro et n'est pas intègre.
\end{proof}


%---------------------------------------------------------------------------------------------------------------------------
\subsection{Caractéristique}
%---------------------------------------------------------------------------------------------------------------------------

L'application 
\begin{equation}
    \begin{aligned}
        \mu\colon \eZ&\to A \\
        n&\mapsto n\cdot 1_A 
    \end{aligned}
\end{equation}
est un morphisme d'anneaux. Le noyau de \( \mu\) étant un sous groupe de \( \eZ\), il existe un et un seul \( p\in\eZ\) tel que \( \ker\mu=p\eZ\). Ce \( p\) est la \defe{caractéristique}{caractéristique!d'un anneau} de \( A\).

\begin{lemma}
    Si \( A\) est de caractéristique nulle, alors \( A\) est infini.
\end{lemma}

\begin{proof}
    En effet, \( \ker\mu=0\) implique que \( n1_A\neq  m1_A\) et par conséquent \( A\) est infini.
\end{proof}

\begin{lemma}       \label{LemHmDaYH}
    Si \( p\) est la caractéristique de l'anneau \( A\), alors nous avons l'isomorphisme d'anneaux
    \begin{equation}
         \eZ 1_A\simeq\eZ/p\eZ.
    \end{equation}
\end{lemma}

\begin{proof}
    L'isomorphisme est donné par l'application \( n1_A\mapsto \phi(n)\) si \( \phi\) est la projection canonique \( \eZ\to \eZ_p\).
\end{proof}

\begin{lemma}       \label{LemCaractIntergernbrcartpre}
    La caractéristique d'un anneau intègre est zéro ou un nombre premier.
\end{lemma}

\begin{proof}
    Si \( A\) est intègre, alors \( \eZ 1_A\) est intègre (a fortiori), et \( \eZ_p\) est intègre parce qu'il est isomorphe à \( \eZ A_A\). Mais nous savons que \( \eZ_p\) est intègre si et seulement si \( p\) est premier (proposition \ref{CorZnInternprem}).
\end{proof}

\begin{theorem}[Théorème chinois]\index{théorème!chinois}
    Soit \( A\) un anneau commutatif, \( n\geq 2\), des éléments \( x_1,\ldots,x_n\) dans \( A\) et des idéaux \( I_1,\ldots,I_n\) tels que \( I_i+I_j=A\) pour tout \( i\neq j\).

    Alors il existe un \( x\in A\) tel que \( x-x_i\in I_i\) pour tout \( 1\leq i\leq n\).
\end{theorem}

\begin{proof}
    Pour \( i\in\{ 1,\ldots,n \}\) nous notons \( J_i\) le produit \( J_i=\prod_{k\neq i}I_k\). Étant donné que chaque \( I_i\) est un idéal, nous avons \( I_k\in J_i\) lorsque \( i\neq k\).

    Soit \( i\) fixé, et considérons \( j\neq i\). Nous pouvons trouver \( a_j\in I_i\) et \( b_j\in I_j\) tel que \( a_j+b_j=1\). Nous avons alors
    \begin{equation}
        1=\prod_{j\neq i}(a_j+b_j).
    \end{equation}
    Par ailleurs \( I_i+J_i=A\) parce que \( J_i\) contient \( I_k\) avec \( k\neq i\) et \( I_i+I_k=A\). Nous pouvons donc prendre \( \alpha_i\in I_i\) et \( \beta_i\in J_i\) tels que
    \begin{equation}
        \prod_{j\neq i}(a_j+b_j)=\alpha_i+\beta_i.
    \end{equation}
    Nous considérons alors l'élément \( x=\beta_1x_1+\ldots+\beta_nx_n\) et nous avons
    \begin{subequations}
        \begin{align}
            x-x_1&=(\beta_1-1)x_1+\beta_2x_2+\ldots+\beta_nx_n\\
            &=-\alpha_1x_1+\beta_2x_2+\ldots+\beta_nx_n
        \end{align}
    \end{subequations}
    Mais \( \alpha_1\in I_1\) et tous les autres termes sont dans les \( J_i\) avec \( i\neq 1\). Par conséquent le tout est dans \( I_1\). Ici nous utilisons par exemple le fait que \( \beta_2\in J_2\subset I_1\) parce que les éléments de \( J_2\) sont des produits d'éléments dont un facteur est dans \( I_1\).
\end{proof}

%---------------------------------------------------------------------------------------------------------------------------
\subsection{Anneaux factoriels}
%---------------------------------------------------------------------------------------------------------------------------

On dit que les éléments \( a\) et \( b\) d'un anneau sont \defe{associés}{associé}\index{éléments!associés} si il existe un élément \( u\) inversible dans \( \eA\) tel que \( a=ub\).

\begin{definition}
    Un anneau commutatif unitaire \( \eA\) est \defe{factoriel}{factoriel!anneau}\index{anneau!factoriel} si il vérifie les propriétés suivantes.
    \begin{enumerate}
        \item
            L'anneau \( \eA\) est intègre.
        \item
            Si \( a\in \eA\) est non nul et non inversible alors il admet une décomposition \( a=p_1\ldots p_k\) où les \( p_i\) sont irréductibles.
        \item
            Si \( a=q_1\ldots q_m\) est une autre décomposition de \( a\) en irréductibles, alors \( m=k\) et il existe une permutation \( \sigma\in S_k\) telle que \( p_i\) et \( q_{\sigma(i)}\) soient associés.
    \end{enumerate}
    <++>
\end{definition}
<++>

%+++++++++++++++++++++++++++++++++++++++++++++++++++++++++++++++++++++++++++++++++++++++++++++++++++++++++++++++++++++++++++
\section{Corps}
%+++++++++++++++++++++++++++++++++++++++++++++++++++++++++++++++++++++++++++++++++++++++++++++++++++++++++++++++++++++++++++

\begin{definition}
    Un \defe{corps}{corps} est un anneau non nul dans lequel tout élément non nul est inversible.
\end{definition}

\begin{lemma}       \label{LemAnnCorpsnonInterdivzer}
    En tant que anneau, un corps n'a pas de diviseurs zéro.
\end{lemma}

\begin{proof}
    En effet si \( a\) est un diviseur de zéro, alors \( ax=0\) pour un certain \( x\neq 0\). Si \( a\) était inversible, nous aurions \( x=a^{-1}ax=0\), ce qui est impossible.
\end{proof}

Par le lemme \ref{LemCaractIntergernbrcartpre}, un anneau sans diviseur de zéro est de caractéristique zéro ou première. Le fait d'être intègre pour un anneau n'assure cependant pas le fait d'être un corps. Nous avons cependant ce résultat pour les anneaux finis.

\begin{proposition}     \label{PropanfinintimpCorp}
    Un anneau fini intègre est un corps.
\end{proposition}

\begin{proof}
    Soit \( A\) un tel anneau. Soit \( a\neq 0\). Les applications 
    \begin{subequations}
        \begin{align}
            l_a\colon x\to ax\\
            r_a\colon x\to xa
        \end{align}
    \end{subequations}
    sont injectives. En tant que applications injectives entre ensembles finis, elles sont surjectives. Il existe donc \( b\) et \( c\) tels que \( 1=ba=ac\). Il se fait que \( b\) et \( c\) sont égaux parce que
    \begin{equation}
        b=b(ac)=(ba)c=c.
    \end{equation}
    Par conséquent \( b\) est un inverse de \( a\).
\end{proof}

\begin{proposition}     \label{PropzhFgNJ}
    Soit \( n\in\eN^*\). Les conditions suivantes sont équivalentes :
    \begin{enumerate}
        \item
            \( n\) est premier.
        \item
            \( \eZ_n\) est un anneau intègre.
        \item
            \( \eZ_n\) est un corps.
    \end{enumerate}
\end{proposition}

\begin{proof}
    L'équivalence entre les deux premiers points est le contenu du corollaire \ref{CorZnInternprem}. Le fait que \( \eZ_n\) soit un corps lorsque \( \eZ_n\) est intègre est la proposition \ref{PropanfinintimpCorp}. Le fait que \( \eZ_n\) soit intègre lorsque \( \eZ_n\) est un corps est une propriété générale des corps : ce sont en particulier des anneaux intègres (lemme \ref{LemAnnCorpsnonInterdivzer}).
\end{proof}

\begin{proposition}     \label{AnnCorpsIdeal}
    Si \( A\) est un anneau, nous avons les équivalences
    \begin{enumerate}
        \item
            \( A\) est un corps.
        \item
            \( A\) est non nul et ses seuls idéaux à gauche sont \( \{ 0 \}\) et \( A\).
        \item
            \( A\) est non nul et ses seuls idéaux à droite sont \( \{ 0 \}\) et \( A\).
    \end{enumerate}
\end{proposition}

\begin{proposition}
    Soit \( A\), un anneau commutatif non nul et \( I\), un idéal dans \( A\). L'ensemble \( I\) est un idéal maximum de \( A\) si et seulement si \( A/I\) est un corps.
\end{proposition}

\begin{proof}
    Par la proposition \ref{AnnCorpsIdeal}, le fait pour \( A/I\) d'être un corps signifie qu'il n'a pas d'idéaux non triviaux. Si \( \phi\) est la projection canonique \( A\mapsto A/I\), nous savons que les idéaux de \( A/I\) sont les \( \phi(J)\) où \( J\) est un idéal de \( A\) contenant \( I\) (proposition \ref{PropIJJIdsousphi}). Si \( I\) est un idéal maximum de \( A\), un tel \( J\) n'existe pas. Inversement si \( A/I\) n'a pas d'autres idéaux que \( A/I\) et \( \phi(0)\), c'est que \( I\) est un idéal maximum.
\end{proof}

%---------------------------------------------------------------------------------------------------------------------------
\subsection{Morphisme de Frobenius}
%---------------------------------------------------------------------------------------------------------------------------

\begin{proposition}     \label{Propqrrdem}
    Soit \( \eA\) un anneau commutatif de caractéristique première \( p\). Alors \( \sigma(x)=x^p\) est un automorphisme de l'anneau \( \eA\). Nous avons la formule
    \begin{equation}
        (a+b)^p=a^p+b^p
    \end{equation}
    pour tout \( a,b\in \eA\).
\end{proposition}

\begin{proof}
    Nous utilisons la formule du binôme est le fait que les coefficients binomiaux non extrêmes sont divisibles par \( p\) et donc nuls.
\end{proof}

\begin{proposition} \label{PropFrobHAMkTY}
    Soit \( \eA\) un anneau commutatif unitaire de caractéristique \( p\). L'application
    \begin{equation}
        \begin{aligned}
            \Frob_\eA\colon \eA&\to \eA \\
            x&\mapsto x^p 
        \end{aligned}
    \end{equation}
    est un automorphisme d'anneau unitaire.
\end{proposition}
Nous le nommons le \defe{morphisme de Frobenius}{morphisme!Frobenius}\index{Frobenius!morphisme}. Nous utiliserons aussi les itérés du morphisme de Frobenius : \( \Frob^k\colon x\mapsto x^{p^k}\).

%---------------------------------------------------------------------------------------------------------------------------
\subsection{Corps des fractions}
%---------------------------------------------------------------------------------------------------------------------------

\begin{theorem}     \label{ThogbhWgo}
    Soit \( \eA\) un anneau commutatif intègre. Il existe un corps commutatif \( \eK\) et un morphisme injectif (d'anneaux) \( \epsilon\colon \eA\to \eK\) tels que pour tout \( \lambda\in\eK\), il existe \( (a,b)\in \eA\times \eA^*\) tels que
    \begin{equation}
        \lambda=\epsilon(a)\big( \epsilon(b) \big)^{-1}
    \end{equation}

    De plus si \( (\eK',\epsilon')\) est un autre couple qui vérifie la propriété, les corps \( \eK\) et \( \eK'\) sont isomorphes.
\end{theorem}
Le corps \( \eK\) associé à l'anneau \( \eA\) est le \defe{corps des fractions}{corps!des fractions}\index{fractions (corps)} de \( \eA\).

Notons que l'application \( \eA\times \eA^*\to \eK\) donnée par \( (a,b)\mapsto \epsilon(a)\big( \epsilon(b) \big)^{-1}\) envoie \( (xa,xb)\) sur le même que \( (a,b)\).

L'ensemble \( \eQ\)\nomenclature[A]{\( \eQ\)}{corps des fractions sur \( \eZ\)} est le corps des fractions de \( \eZ\).

%---------------------------------------------------------------------------------------------------------------------------
\subsection{Corps premier}
%---------------------------------------------------------------------------------------------------------------------------
\label{subseccorpspremhBlYIv}

\begin{definition}
    Un corps est \defe{premier}{corps!premier}\index{premier!corps} si il est son seul sous corps. Le \defe{sous corps premier}{premier!sous corps} d'un corps est l'intersection de tous ses sous corps.
\end{definition}

\begin{lemma}
    Un corps premier est commutatif
\end{lemma}

\begin{proof}
    Le centre d'un corps est certainement un sous corps. Par conséquent un corps premier doit être contenu dans son propre centre, c'est à dire être commutatif.
\end{proof}

Soit \( p\) un nombre premier. Nous notons \( \eF_p=\eZ_p=\eZ/p\eZ\)\nomenclature[A]{\( \eF_p\)}{lorsque \( p\) est premier}. Nous verrons plus loi (section \ref{SecCorpsFinizkAcbS}) comment nous pouvons définir \( \eF_{p^n}\) lorsque \( p\) est premier.

Nous avons par exemple 
\begin{equation}
    \eF_2=\eZ/2\eZ=\{ 0,1 \}
\end{equation}
avec la loi \( 2=0\).

Notons que \( \eF_p\) est un corps possédant \( p\) éléments. L'ensemble \( \eF_p^*\) est un groupe d'ordre \( p-1\).

\begin{lemma}
    Les corps \( \eQ\) et \( \eF_p\) (avec \( p\) premier) sont premiers.
\end{lemma}

\begin{proof}
    Tout sous corps de \( \eQ\) doit contenir \( 1\), et par conséquent \( \eZ\). Devant également contenir tous les inverses, il contient \( \eQ\).

    Tout sous corps de \(\eF_p \) doit contenir \( 1\) et donc \( \eF_p\) en entier. Par ailleurs nous savons de la proposition \ref{PropzhFgNJ} que \( \eF_p\) est un corps lorsque \( p\) est premier.
\end{proof}

\begin{proposition}
    Soit \( \eK\) un corps de caractéristique \( p\) et \( \eP\) son sous corps premier. Si \( p=0\) alors \( \eP=\eQ\). Si \( p>0\), alors \( \eP=\eF_p\).
\end{proposition}

\begin{proof}
    Notons d'abord que la caractéristique d'un corps est toujours un nombre premier parce qu'un corps est en particulier un anneau intègre (proposition \ref{LemCaractIntergernbrcartpre}).

    Étant donné que \( 1\) est dans tout sous corps, nous devons avoir \( \eZ 1\subseteq \eP\). Si \( p=0\), alors \( \eZ 1\simeq \eZ\), et nous avons
    \begin{equation}
        \eZ 1_{\eA}\subset \eP\subset \eK.
    \end{equation}
    Pour chaque \( (n,m)\in \eZ 1_{\eA}\times (\eZ 1_{\eA})^*\) l'élément \( nm^{-1}\in \eK\) est dans \( \eP\) parce que \( \eP\) est un corps. Nous en déduisons que le corps des fractions de \( \eZ\) est contenu dans \( \eP\) par conséquent \( \eP=\eQ\) (théorème \ref{ThogbhWgo}). 

    Si par contre la caractéristique de \( \eK\) est \( p\neq 0\), nous avons \( \eZ 1_{\eA}\simeq\eZ/p\eZ=\eF_p\) par le lemme \ref{LemHmDaYH}. L'ensemble \( \eF_p\) étant un corps, c'est le corps premier de \( \eK\).
\end{proof}

\begin{proposition}     \label{PropqPPrgJ}
    Soit \( \eK\) un corps et \( \eP\) sont sous corps premier. Si \( \sigma\in\Aut(\eK)\) alors \( \sigma|_{\eP}=\id\), c'est à dire que
    \begin{equation}
        \sigma(x)=x
    \end{equation}
    pour tout \( x\in \eP\).
\end{proposition}

\begin{theorem}[Petit théorème de Fermat]
    Soit \( p\) un nombre premier. Si \( x\in \eF_p\) alors \( x^p=x\). Si \( x\in(\eF_p)^*\), alors \( x^{p-1}=1\).

    En particulier si \( x\in \eF_p^*\) alors \( x^{-1}=x^{p-2}\).
\end{theorem}

\begin{proof}
    Étant donné que \( \eF_p\) est un corps commutatif et que \( p\) est premier, la proposition \ref{Propqrrdem} nous indique que \( \sigma(x)=x^p\) est un automorphisme. La proposition \ref{PropqPPrgJ} nous indique alors que
    \begin{equation}
        a^p=a.
    \end{equation}
    Si \( a\) est inversible alors \( a^{p-1}=a^pa^{-1}=1\).
\end{proof}

\begin{example}
    Soit \( \eK=\eF_29^*\). Le nombre \( 29\) étant premier, \( \eK\) est un corps premier. C'est le corps des entier modulo \( 29\). Nous avons donc
    \begin{equation}
            -142=-113=-84=-55=-26=3=32=61=90=119.
    \end{equation}
    Le petit théorème de Fermat nous permet aussi de calculer des exposants et des inverses. Nous avons
    \begin{equation}
        x^{-a}=(x^a)^{27}=x^{27a}.
    \end{equation}
    Le nombre \( 27 a\) peut être grand par rapport à \( 29\). Étant donné que \( 1=x^{28}\) nous avons
    \begin{equation}
        x^{-a}=x^{27 a\mod 28}.
    \end{equation}
    Dans les exposants nous calculons modulo \( 28\).
\end{example}

%+++++++++++++++++++++++++++++++++++++++++++++++++++++++++++++++++++++++++++++++++++++++++++++++++++++++++++++++++++++++++++
\section{Modules}
%+++++++++++++++++++++++++++++++++++++++++++++++++++++++++++++++++++++++++++++++++++++++++++++++++++++++++++++++++++++++++++

Soit \( \modE\) un \( A\)-module et \( x=(x_i)_{i\in I}\) une famille d'éléments de \( \modE\), paramétrée par l'ensemble \( I\). Nous considérons l'application
\begin{equation}
    \begin{aligned}
        \mu_x\colon A^{(I)}&\to \modE \\
        (a_i)_{i\in I}&\mapsto \sum_{i\in I}a_ix_i.
    \end{aligned}
\end{equation}
Ici \( A^{(I)}\) désigne l'ensemble de toutes les applications \( I\to A\) de support fini.  

\begin{definition}      \label{DefBasePouyKj}
    À l'instar des espaces vectoriels, les modules ont une notion de partie libre, génératrice et de bases :
    \begin{enumerate}
        \item
            Si \( \mu_x\) est surjective, nous disons que \( x\) est une partie \defe{génératrice}{génératrice!partir d'un module}.
        \item
            Si \( \mu_x\) est injective, nous disons que la partie \( x\) est \defe{libre}{libre!partie d'un module}.
        \item
            Si \( \mu_x\) est bijective, nous disons que la partie \( x\) est une \defe{base}{base!d'un module}.
    \end{enumerate}
\end{definition}

\begin{definition}
    Soit \( \modE\) un module sur un anneau commutatif \( A\). Un \defe{projecteur}{projecteur!dans un module} est une application linéaire \( p\colon \modE\to \modE\) telle que \( p^2=p\).

    Une famille \( (p_i)_{i\in I}\) sur \( \modE\) est \defe{orthogonale}{orthogonal!famille de projecteurs} si \( p_i\circ p_j=0\) pour tout \( i\neq j\). La famille est \defe{complète}{complète!famille de projecteurs} si \( \sum_{i\in I}p_i=\mtu\).
\end{definition}

\begin{theorem}     \label{ThoProjModpAlsUR}
    Soient des sous modules \( \modE_1,\ldots,\modE_n\) du module \( \modE\) tels que \( \modE=\modE_1\oplus\ldots\oplus\modE_n\). Les applications \( p_i\) définies par
    \begin{equation}
        p_i(x_1+x_n)=x_i
    \end{equation}
    forment une famille orthogonale de projecteurs et \( p_1+\ldots +p_n=\id\).

    Inversement, si \( (p_1,\ldots, p_n)\) est une famille orthogonale de projecteurs dans un module \( \modE\) tel que \( \sum_{i=1}^np_i=\id\), alors
    \begin{equation}
        \modE=\bigoplus_{i=1}^np_i(\modE).
    \end{equation}
\end{theorem}

%+++++++++++++++++++++++++++++++++++++++++++++++++++++++++++++++++++++++++++++++++++++++++++++++++++++++++++++++++++++++++++
\section{Polynômes}
%+++++++++++++++++++++++++++++++++++++++++++++++++++++++++++++++++++++++++++++++++++++++++++++++++++++++++++++++++++++++++++

Soit \( A\) un anneau commutatif. Nous considérons \( \polyP\) l'ensemble des suites presque nulles d'éléments de \( A\), ce sont les suites \( (a_n)_{n\in\eN}\) telles que il existe \( N\) tel que \( a_i=0\) pour tout \( i>N\).

Cela est un \( A\)-module libre de base (définition \ref{DefBasePouyKj})
\begin{equation}
    (e_n)_k=\delta_{nk}.
\end{equation}
Si \( (a_n)_{n\in \eN}\) et \( (b_n)_{n\in\eN}\) sont des éléments de \( \polyP\), nous définissons le produit \( ab\) par
\begin{equation}
    (ab)_n=\sum_{p+q=n}a_pb_q.
\end{equation}
Cela est bien un élément de \( \polyP\) parce qu'il existe \( N\in\eN\) tel que \( a_n=b_n=0\) pour tout \( n\geq N\). Avec la somme et le produit par un scalaire, le module \( \polyP\) devient une \( A\)-algèbre commutative unitaire. L'unité est 
\begin{equation}
    e_0=(1,0,\ldots).
\end{equation}

\begin{definition}
    En tant que \( A\)-algèbre, l'ensemble \( \polyP\) est l'\defe{algèbre des polynômes en une indéterminée}{algèbre!polynômes} à coefficients dans \( A\).
\end{definition}

Si nous posons que \( X=e_1\), et que nous prenons la convention \( X^0=1\), alors nous avons \( e_k=X^k\) et nous notons \( A[X]\) l'anneau \( \polyP\) exprimé avec \( X\). Les éléments de la forme \( \lambda X^k\) avec \( \lambda\in A\) et \( k\in\eN\) sont des \defe{monômes}{monôme}. Nous allons aussi considérer
\begin{equation}\nomenclature[A]{\( A_n[X]\)}{les polynômes à coefficients dans \( A\) et de degré inférieur à \( n\)}
    A_n[X]=\{ P\in A[X]\tq \deg(P)\leq n \}.
\end{equation}
Cela est un sous module libre.

\begin{theorem}
    L'anneau \( A\) est intègre si et seulement si \( A[X]\) est intègre.
\end{theorem}

\begin{proof}
    Soient \( P\) et \( Q\) des éléments non nuls de \( A[X]\). Vu que l'anneau \( A\) est intègre, nous avons
    \begin{equation}
        \deg(PQ)=\deg(P)+\deg(Q)
    \end{equation}
    et le produit ne peut pas être nul. L'anneau \( A[X]\) est donc intègre.

    Si \( A[X]\) est intègre, \( A\) est intègre parce qu'il peut être vu comme sous anneau.
\end{proof}

\begin{remark}
    Si \( A\) n'est pas intègre, soit \( \alpha\beta=0\), alors \( (\alpha X)(\beta x)=0\) et le degré du produit n'est pas la somme des degrés.
\end{remark}

\begin{corollary}
    Si \( A\) est intègre, les inversibles de \( A[X]\) sont les éléments de \( U(A)\).
\end{corollary}

\begin{proof}
    Pour que \( Q\) soit inversible, il faut un \( P\) tel que \( PQ=1\). Mais l'anneau \( A\) étant intègre, les degrés s'additionnent. Par conséquent ils doivent être de degrés zéro et il faut que \( P,Q\in A\). Enfin pour qu'ils soient inversibles, ils doivent être dans \( U(A)\).
\end{proof}

La \defe{valuation}{valuation} de \( P\) du polynôme \( P=\sum_n a_nX^n\), notée \( \val(P)\), est 
\begin{equation}
    \val(P)=\min\{ n\tq a_n\neq 0 \}.
\end{equation}
Nous avons \( \val(P)\leq \deg(P)\) et \( \val(P)=\deg(P)\) si et seulement si \( P\) est un monôme. Si \( P=0\), nous convenons que \( \val(0)=\infty\) et \( \deg(0)=-\infty\).

%---------------------------------------------------------------------------------------------------------------------------
\subsection{Irréductibilité}
%---------------------------------------------------------------------------------------------------------------------------

\begin{theorem}[d'Alembert-Gauss]\index{théorème!d'Alembert-Gauss}      \label{ThovgyUuA}
    Tout polynôme non constant à coefficients complexes possède au moins une racine complexe.
\end{theorem}


\begin{definition}      \label{DefIrredfIqydS}
    Un polynôme est \defe{irréductible}{irréductible!polynôme} lorsqu'il ne peut pas être écrit sous la forme de produits de polynômes de degré supérieurs à \( 1\).
\end{definition}

\begin{proposition}
    Un polynôme irréductible à coefficients réels est soit de degré un soit de degré \( 2\) avec un discriminant négatif.
\end{proposition}

\begin{proof}
    Soit un polynôme \( P\) à coefficients réels de degré plus grand que \( 1\). Alors le théorème de d'Alembert-Gauss (théorème \ref{ThovgyUuA}) implique l'existence d'une racine \( \alpha\). Il est facile de montrer que le conjugué complexe \( \bar \alpha\) est également racine. Par conséquent les polynômes \( (X-\alpha)\) et \( (X-\bar \alpha)\) divisent \( P\).

    Ces deux polynômes sont premiers entre eux parce que
    \begin{equation}
        a(X-\alpha)+b(X-\bar \alpha)=0
    \end{equation}
    implique \( a=b=0\). Par conséquent le produit 
    \begin{equation}
        X^2-(\alpha+\bar \alpha)X+\alpha\bar\alpha
    \end{equation}
    divise également \( P\). Ce dernier est un polynôme à coefficients réels de degré \( 2\). Donc tout polynôme de degré \( 3\) ou plus est réductible.
\end{proof}

Nous disons que \( P\in\eK[X]\setminus\eK\) est \defe{scindé}{polynôme!scindé} sur \(\eK\) si il est produit dans \(\eK[X]\) de polynômes de degré \( 1\).

%---------------------------------------------------------------------------------------------------------------------------
\subsection{Division euclidienne}
%---------------------------------------------------------------------------------------------------------------------------

Le théorème suivant établit la \defe{division euclidienne}{division!euclidienne} dans \( \eA[X]\) du polynôme \( A\) par \( B\).
\begin{theorem}     \label{ThodivEuclPsFexf}
    Soit \( B\neq 0\) dans \( \eA[X]\) de coefficient dominant inversible dans \( \eA\). Pour tout \( A\in\eA[X]\), il existe \( P,Q\in \eA[X]\) tels que
    \begin{equation}
        A=BQ+R
    \end{equation}
    avec \( \deg(R)<\deg(B)\).

    Les polynômes \( Q\) et \( R\) sont déterminés de façon univoque par cette condition. Le polynôme \( Q\) est le \defe{quotient}{quotient} et \( R\) est le \defe{reste}{reste} de la division euclidienne de \( A\) par \( B\).
\end{theorem}

Deux polynômes \( P\) et \( Q\) sont dits \defe{étrangers}{étrangers!polynômes} entre eux si \( 1\) est un \( \pgcd\) de \( P\) et \( Q\). Un ensemble de polynômes \( (P_i)_{i\in I}\) est étranger \defe{dans leur ensemble}{étranger!dans leur ensemble} si \( 1\) est un \( \pgcd\) des \( P_i\).

\begin{theorem}[Bezout] \label{ThoBezoutOuGmLB}
    Les polynômes \( P_1,\ldots,P_n\) dans \( \eK[X]\) sont étrangers entre eux si et seulement si il existe des polynômes \( Q_1,\ldots,Q_n\in\eK[X]\) tels que
    \begin{equation}
        P_1Q_1+\ldots+P_nQ_n=1.
    \end{equation}
\end{theorem}

\begin{lemma}       \label{LemuALZHn}
    Soient \( (P_i)_{i=1,\ldots,n}\in \eK[X]\) des polynômes étrangers deux à deux. Alors les polynômes
    \begin{equation}
        Q_i=\prod_{j\neq i}P_j
    \end{equation}
    sont étrangers entre eux\footnote{Et non juste deux à deux.}.
\end{lemma}

%---------------------------------------------------------------------------------------------------------------------------
\subsection{Idéaux}
%---------------------------------------------------------------------------------------------------------------------------

Soit \( P\in \eK[X]\) un polynôme. Nous notons \( (P)\) l'idéal engendré par \( P\) :
\begin{equation}        \label{EqDefxMkDtW}
    (P)=\{ PR\tq R\in\eK[X] \}.
\end{equation}

\begin{lemma}
    Nous avons
    \begin{enumerate}
        \item
            \( (P)\subset (Q)\) si et seulement si \( Q\) divise \( P\),
        \item
            \( (P)=(Q)\) si et seulement si \( P\) et \( Q\) sont multiples (non nuls) l'un de l'autre.
    \end{enumerate}
\end{lemma}

\begin{proof}
    Si \( (P)\subset (Q)\), en particulier \( P\in(Q)\) et il existe \( R\in\eK[X]\) tel que \( P=QR\), ce qui signifie que \( Q\) divise \( P\).

    Si les idéaux de \( P\) et de \( Q\) sont identiques, l'un divise l'autre et l'autre divise l'un. Ils sont donc multiples l'un de l'autre.
\end{proof}

\begin{theorem}     \label{ThoCCHkoU}
    Soit \( I\) un idéal dans \( \eK[X]\). Alors il existe un polynôme \( P\) tel que \( I=(P)\). Plus précisément, si \( P\) est de degré minimal, alors \( (P)=I\).

    De plus si \( I\neq \{  0\}\), il existe un unique polynôme unitaire \( U\) tel que \( I=(U)\).
\end{theorem}

\begin{proof}
    Si \( I=\{ 0 \}\), le résultat est évident. Nous supposons donc \( I\) non nul. Soit \( P\) de degré minimum parmi les éléments de \( I\). Évidemment \( (P)\subset I\). Nous allons démontrer qu'en réalité \( (P)=I\).

    Soit \( A\in I\). Par le théorème \ref{ThodivEuclPsFexf} de la division euclidienne\footnote{Ici \( \eK\) est un corps et donc l'hypothèse d'inversibilité est automatiquement vérifiée.}, il existe \( Q\) et \( R\) dans \( \eK[X]\) tels que \( A=PQ+R\) avec \( \deg(R)<\deg(P)\). Étant donné que \( R=A-PQ\) nous avons \( R\in I\) et par conséquent \( R=0\) parce que \( P\) a été choisit de degré minimum dans \( I\). Nous avons donc \( A=PQ\) et \( I\subset (P)\).

    L'existence d'un polynôme unitaire qui génère \( I\) est obtenu en choisissant \( U=P/a_n\) où \( a_n\) est le coefficient du terme de plus haut degré.
\end{proof}
Nous voyons que n'importe quel polynôme de degré minimum dans un idéal génère l'idéal. Une importante conséquence du théorème \ref{ThoCCHkoU} que nous verrons plus bas est que tout polynôme annulateur d'un endomorphisme est divisé par le polynôme minimal (proposition \ref{PropAnnncEcCxj}).

\begin{corollary}       \label{CorvlvOWA}
    Soit \( \eK\) un corps commutatif. Alors l'anneau des polynômes \( \eK[X]\) est principal.
\end{corollary}

\begin{proof}
    En effet \( \eK[X]\) est un anneau commutatif intègre (pas de diviseurs de zéro), et d'après le théorème \ref{ThoCCHkoU}, tous les idéaux sont principaux en tant que pour tout idéal \( I\), il existe \( P\in I\) tel que \( I=\eK[X]P=(P)\).
\end{proof}

%+++++++++++++++++++++++++++++++++++++++++++++++++++++++++++++++++++++++++++++++++++++++++++++++++++++++++++++++++++++++++++
\section{Racines de polynômes}
%+++++++++++++++++++++++++++++++++++++++++++++++++++++++++++++++++++++++++++++++++++++++++++++++++++++++++++++++++++++++++++

Soit \( \eA\) un anneau et \( P\in \eA[X]\) un polynôme et \( \alpha\in \eA\). Le \defe{degré}{degré!d'une racine} ou la \defe{multiplicité}{multiplicité!d'une racine} de \( \alpha\) par rapport à \( P\) est l'entier \( h\) tel que \( P\) est divisible par \( (X-\alpha)^h\) mais pas divisible par \( (X-\alpha)^{h+1}\).

Nous noterons \( \theta_{\alpha}(P)\)\nomenclature[A]{\( \theta_{\alpha}(P)\)}{l'ordre de \( \alpha\) par rapport à \( P\)} l'ordre de \( \alpha\) par rapport à \( P\).

\begin{proposition}     \label{PropahQQpA}
    L'élément \( \alpha\in \eA\) est d'ordre \( h\) par rapport à \( \) si et seulement si il existe \( Q\in\eA[X]\) tel que \( P(X-\alpha)^hQ\) avec \( Q(\alpha)\neq 0\).
\end{proposition}

\begin{lemma}       \label{LemIeLhpc}
    Soient \( P\) et \( Q\) des polynômes non nuls de \( \eA[X]\) et \( \alpha\in \eA\) d'ordre \( p\) pour \( P\) et d'ordre \( q\) pour \( Q\). Alors
    \begin{enumerate}
        \item
            \( \theta_{\alpha}(P+Q)\geq\ln\{ \theta_{\alpha}(P),\theta_{\alpha}(Q) \}\)
        \item
            si \( \theta_{\alpha}(P)\neq \theta_{\alpha}(Q)\), alors \( \theta_{\alpha}(P+Q)=\min\{ \theta_{\alpha}(P),\theta_{\alpha}(Q) \}\)
        \item
            \( \theta_{\alpha}(PQ)\geq \theta_{\alpha(P)}+\theta_{\alpha}(Q)\);
        \item       \label{ItemIeLhpciv}
            si \(\eA \) est intègre alors \( \theta_{\alpha}(PQ)= \theta_{\alpha}(P)+\theta_{\alpha}(Q)\);
    \end{enumerate}
\end{lemma}

\begin{theorem}
    Soit \( \eA\) un anneau intègre et \( P\in \eA[X]\setminus\{ 0 \}\), un polynôme de degré \( n\). Si \( \alpha_1,\ldots, \alpha_p\in\eA\) sont des racines deux à deux distinctes d'ordres \( k_1,\ldots, k_p\), alors il existe \( Q\in \eA[X]\) tel que
    \begin{enumerate}
        \item
            \( Q(\alpha_i)\neq 0\);
        \item
            \( P=Q\prod_{i=1}^p(X-\alpha_i)\);
    \end{enumerate}
    De plus la sommes des ordres des racines de \( P\) est au plus \( \deg(P)\).
\end{theorem}

\begin{proof}
    Si \( p=1\), alors le résultat est la proposition \ref{PropahQQpA}. Nous supposons que \( p\geq 2\) et nous effectuons une récurrence sur \( P\). Nous considérons donc pas \( p-1\) premières racines \( \alpha_1,\ldots, \alpha_{p-1}\) et un polynôme \( R\in\eA[X]\) tel que \( R(\alpha_i)\neq 0\) pour \( i=1,\ldots, p-1\) et
    \begin{equation}
        P=\underbrace{(X-\alpha_1)^{k_1}\ldots (X-\alpha_{p-1})^{k_{p-1}}}_SR.
    \end{equation}
    Par hypothèse \( P(\alpha_p)=S(\alpha_p)R(\alpha_p)=0\). L'anneau \( \eA\) étant intègre, \( S(\alpha_p)\neq 0\) parce que \( \alpha_i\neq \alpha_p\) pour \( i\neq p\). Par conséquent, \( R(\alpha_p)=0\).
    
    Nous devons encore vérifier que l'ordre de \( \alpha_p\) est \( k_p\) par rapport à \( R\). Pour cela nous utilisons le point \ref{ItemIeLhpciv} du lemme \ref{LemIeLhpc} affin de dire que le degré de \( \alpha_p\) pour \( P=SR\) est \( k_p\). Par conséquent
    \begin{equation}
        R=(X-\alpha_p)^{k_p}T
    \end{equation}
    avec \( T(\alpha_p)\neq 0\) et enfin
    \begin{equation}
        P=\prod_{i=1}^p(X-\alpha_i)T.
    \end{equation}
    De plus \( T(\alpha_i)\neq 0\), sinon \( R(\alpha_i)\) serait nul.
\end{proof}

%+++++++++++++++++++++++++++++++++++++++++++++++++++++++++++++++++++++++++++++++++++++++++++++++++++++++++++++++++++++++++++
\section{Anneaux principaux}
%+++++++++++++++++++++++++++++++++++++++++++++++++++++++++++++++++++++++++++++++++++++++++++++++++++++++++++++++++++++++++++

Nous avons parlé de l'idéal des polynôme annulateurs dans le théorème \ref{ThoCCHkoU}.

\begin{definition}      \label{DefIdPrinpuMrbOq}
    Un idéal \( I\) dans \(\eA\) est \defe{principal à gauche}{idéal!principal!à gauche} si il existe \( a\in I\) tel que \( I=\eA a\). Il est \defe{principal à droite}{idéal!principal!à droite} si il existe \( a\in I\) tel que \( I=a\eA\). Nous disons qu'il est \defe{principal}{principal!idéal} si il est principal à gauche et à droite.

    Un anneau commutatif intègre est \defe{principal}{principal!anneau} si tous ses idéaux sont principaux.
\end{definition}

Un idéal \( I\) dans l'anneau \( \eA\) est \defe{maximal}{maximal!idéal}\index{idéal!maximal} si les seuls idéaux de \( \eA\) contenants \( I\) sont \( I\) et \( \eA\).

\begin{definition}  \label{DeirredBDhQfA}
    Soit \( \eA\) un anneau commutatif intègre. Un élément \( a\in\eA\) est \defe{irréductible}{irréductible!dans un anneau} si \( a\) n'est pas inversible, mais si \( a=xy\), alors soit \( x\) soit \( y\) est inversible. Nous notons \( U(\eA)\) l'ensemble des éléments inversibles de \( \eA\).
\end{definition}

\begin{remark}
    Un corps n'a pas d'éléments irréductibles parce qu'à part zéro tous les éléments sont inversibles alors que \( 0\) n'est certainement pas irréductible vu que \( 0=0\cdot 0\).
\end{remark}


\begin{definition}
    Nous disons qu'un idéal \( I\) dans \( \eA\) est \defe{premier}{premier!idéal} si \( \eA\) est un anneau commutatif intègre et si \( A/I\) est intègre.
\end{definition}

\begin{proposition}
    Soit \( \eA\) un anneau principal qui n'est pas un corps. Pour un idéal \( I\subset \eA\), les conditions suivantes sont équivalentes :
    \begin{enumerate}
        \item
            \( I\) est un idéal maximum;
        \item
            \( I\) est un idéal premier non nul;
        \item
            il existe \( p\) irréductible dans \( \eA\) tel que \( I=(p)\).
    \end{enumerate}
\end{proposition}

\begin{proposition}
    Un idéal \( I\) dans \( \eA\) est premier si et seulement si \( I\) est strictement inclus dans \( \eA\) et si pour tout \( a,b\in\eA\) tels que \( ab\in I\) nous avons \( a\in I\) ou \( b\in I\).
\end{proposition}

\begin{proposition}     \label{PropoTMMXCx}
    Si \( \eA\) est un anneau principal et si \( p\) est irréductible, alors \( A/p\) est un corps.
\end{proposition}

%---------------------------------------------------------------------------------------------------------------------------
\subsection{Anneaux principaux et polynômes}
%---------------------------------------------------------------------------------------------------------------------------

Nous supposons que \( \eK\) est une corps commutatif, et nous étudions l'anneau \( \eK[X]\). Étant donné que \( \eK\) est commutatif pour tout polynôme que l'idéal engendré par \( P\) est \( (P)=\eK[X]P\), voir la notation \eqref{EqDefxMkDtW}.

\begin{remark}
    Un polynôme est irréductible dans \( \eK[X]\) au sens de la définition \ref{DeirredBDhQfA} si et seulement si il est irréductible au sens de la définition \ref{DefIrredfIqydS} parce que seules les constantes (non nulles) sont inversibles dans \( \eK[X]\).
\end{remark}

\begin{corollary}       \label{CorsLGiEN}
    Si \( \eK\) est un corps et si \( P\) est un polynôme irréductible de degré \( n\), alors l'ensemble \( \eL=\eK[X]/(P)\) est un corps. De plus \( \eL\) est un espace vectoriel de dimension \( n\).
\end{corollary}

\begin{proof}
    En effet \( \eK[X]\) est un anneau principal par le corollaire \ref{CorvlvOWA}, par conséquent la proposition \ref{PropoTMMXCx} déduit que \( \eK[X]/(P)\) est un corps.

    Une base de \( \eL\) est donnée par les projections de \( 1,X,X^2,\ldots, X^n\).
\end{proof}

%+++++++++++++++++++++++++++++++++++++++++++++++++++++++++++++++++++++++++++++++++++++++++++++++++++++++++++++++++++++++++++
\section{Polynômes cyclotomiques}
%+++++++++++++++++++++++++++++++++++++++++++++++++++++++++++++++++++++++++++++++++++++++++++++++++++++++++++++++++++++++++++

Soit \( n\in \eN^*\). Nous considérons le polynôme \( (X^n-1)\in\eC[X]\). Les racines de ce polynôme forment le groupe
\begin{equation}
    \gU_n=\{ \xi^k\tq k=0,\ldots, n-1 \}
\end{equation}
avec \( \xi= e^{2i\pi/n}\). Ce groupe est un groupe cyclique généré par \( \xi\). Les autres générateurs sont les \( \xi^p\) avec \( \pgcd(p,n)=1\). Nous notons \( \Delta_n\) l'ensemble des générateurs de \( \gU_n\) :
\begin{equation}
    \Delta_n=\{  e^{2ki\pi/n}\tq 0\leq k\leq n-1,\pgcd(k,n)=1 \}.
\end{equation}
Ces éléments sont les \defe{racines primitives}{racine!primitive de l'unité} de l'unité dans \( \eC\). Nous avons 
\begin{equation}
    \Card(\Delta_n)=\varphi(n)
\end{equation}
où \( \varphi\) est la fonction d'Euler définie par \eqref{EqEulerGqPsvi}.

Nous avons par exemple
\begin{subequations}
    \begin{align}
        \Delta_1&=\{ 1 \}\\
        \Delta_2&=\{  e^{\pi i} \}\\
        \Delta_4&=\{  e^{\pi i/2}, e^{3\pi i/2} \}.
    \end{align}
\end{subequations}

\begin{lemma}       \label{LemKcpjee}
    Nous avons
    \begin{equation}        \label{EqpZuIyL}
        \gU_n=\bigcup_{d\divides n}\Delta_d
    \end{equation}
    et l'union est disjointe. Nous avons aussi la formule
    \begin{equation}        \label{EqTPHqgJ}
        n=\sum_{d\divides n}\varphi(d).
    \end{equation}
\end{lemma}

\begin{proof}
    À l'application \( x\mapsto  e^{2i\pi x}\) près, nous pouvons considérer
    \begin{equation}
        \Delta_d=\{ \frac{ k }{ d }\tq k=0,\ldots, d-1, \pgcd(k,d)=1 \},
    \end{equation}
    c'est à dire l'ensemble des fractions irréductibles dont le dénominateur est \( d\). L'union des \( \Delta_d\) sera donc disjointe.
    
    Toujours à l'application \( x\mapsto  e^{2i\pi x}\) près, le groupe \( \gU_n\) est donné par
    \begin{equation}
        \gU_n=\{ \frac{ k }{ n }\tq k=0,\ldots, n-1 \}.
    \end{equation}
    L'égalité \eqref{EqpZuIyL} revient maintenant à dire que toute fraction de la forme \( \frac{ k }{ n }\) s'écrit de façon irréductible avec un dénominateur qui divise \( n\).

    La relation \eqref{EqTPHqgJ} consiste à prendre le cardinal des deux côtés de \eqref{EqpZuIyL}. Nous avons \( \Card(\gU_n)=n\) et l'union étant disjointe, à droite nous avons la somme des cardinaux.
\end{proof}

Le \defe{polynôme cyclotomique}{polynôme!cyclotomique} d'indice \( n\) est le polynôme
\begin{equation}
    \phi_n(X)=\prod_{a\in\Delta_n}(X-a).
\end{equation}
C'est un polynôme unitaire de degré \( \varphi(n)\). Pour former \( \Delta_n\), nous posons \( \xi= e^{2\pi i/n}\) et puis
\begin{equation}
    \Delta_n=\{ \xi^k\, 0\leq k\leq n-1\tq \pgcd(k,n)=1 \}.
\end{equation}
Nous avons par exemple
\begin{subequations}
    \begin{align}
        \Delta_1&=\{ 1 \}\\
        \Delta_2&=\{ -1 \}\\
        \Delta_3&=\{  e^{2\pi i/3, e^{4\pi i/3}} \}
    \end{align}
\end{subequations}
et les premiers polynômes cyclotomiques sont donnés par
\begin{subequations}
    \begin{align}
        \phi_1(X)&=X-1\\
        \phi_2(X)&=X+1\\
        \phi_3(X)&=X^2+X+1.
    \end{align}
\end{subequations}
Pour le dernier nous avons utilisé le fait que \(  e^{6\pi i/3}=1\) et \(  e^{4\pi i/3+ e^{2\pi i/3}}=-1\).

%+++++++++++++++++++++++++++++++++++++++++++++++++++++++++++++++++++++++++++++++++++++++++++++++++++++++++++++++++++++++++++
\section{Extensions de corps}
%+++++++++++++++++++++++++++++++++++++++++++++++++++++++++++++++++++++++++++++++++++++++++++++++++++++++++++++++++++++++++++

Soit \( \eK\) un corps commutatif. Une \defe{extension}{extension!de corps} de \( \eK\) est un corps \( \eL\) muni d'un morphisme \( i\colon \eL\to \eK\). Nous identifions le plus souvent \( \eK\) avec \( i(\eK)\subset \eL\).

Notons que \( \eL\) est un espace vectoriel sur \( \eK\) parce que si \( \lambda\in \eK\) et \( x\in \eL\) nous pouvons considérer la multiplication scalaire
\begin{equation}
    \lambda\cdot x=i(\lambda)x
\end{equation}
où la multiplication du membre de droite est celle du corps \( \eL\). Le \defe{degré}{degré!extension de corps} de \( \eL\) est la dimension de cet espace vectoriel. Il est noté \( [\eL:\eK]\)\nomenclature[A]{\( [\eL:\eK]\)}{degré d'une extension de corps}; notons qu'il peut être infini.

Soit \( \eL\) une extension de \( \eK\) et \( A\subset \eL\). Nous notons \( \eK(A)\)\nomenclature[A]{$\eK(A)$}{corps contenant \( \eK\) et \( A\)} le plus petit sous corps de \( \eL\) qui contient \( \eK\) et \( A\). Nous notons \( \eK[A]\)\nomenclature[A]{$\eK[A]$}{anneau contenant \( \eK\) et \( A\)} le plus petit sous anneau de \( \eL\) qui contienne \( \eK\) et \( A\).

Nous disons que l'extension \( \eL\) de \( \eK\) est \defe{monogène}{monogène!extension de corps} ou \defe{\wikipedia{fr}{Extension_simple}{simple}}{extension!simple}\index{simple!extension de corps} si il existe \( x\in\eL\) tel que \( \eL=\eK(x)\).

Une extension \( \eL\) de \( \eK\) est \defe{séparable}{séparable!extension de corps} si le polynôme minimum de tout élément de \( \eL\) n'admet que des racines simples.

\begin{lemma}
    Soit \( P\in\eK[X]\) un polynôme unitaire irréductible de degré \( n\). Il existe une extension \( \eL\) de \( \eK\) et \( a\in \eL\) telle que \( \eL=\eK(a)\) et \( P\) est le polynôme minimal de \( a\) dans \( \eL\).
\end{lemma}

\begin{proof}
    Nous prenons \( \eL=\eK[X]/(P)\) où \( (P)\) est l'idéal dans \( \eK[X]\) généré par \( P\). Cela est un corps par le corollaire \ref{CorsLGiEN}. Nous identifions \( \eK\) avec \( \phi(\eK)\) où
    \begin{equation}
        \phi\colon \eK[X]\to \eL 
    \end{equation}
    est la projection canonique. Nous considérons également \( a=\phi(X)\).

    Nous avons alors \( P(a)=0\) dans \( \eL\). En effet \( P(a)=P\big( \phi(X) \big)\) est à voir comme l'application du polynôme \( P\) au polynôme \( X\), le résultat étant encore un élément de \( \eL\). En l'occurrence le résultat est \( P\) qui vaut \( 0\) dans \( \eL\).

    Le polynôme \( P\) étant unitaire et irréductible, il est minimum dans \( \eL\).

    Nous devons encore montrer que \( \eL=\eK(a)\). Le fait que \( \eK(a)\subset \eL\) est une tautologie parce qu'on calcule \( \eK(a)\) dans \( \eL\). Pour l'inclusion inverse soit \( Q(X)=\sum_iQ_iX^i\) dans \( \eK[X]\). Dans \( \eL\) nous avons évidemment \( Q=\sum_iQ_ia^i\).
\end{proof}

%---------------------------------------------------------------------------------------------------------------------------
\subsection{Corps de rupture}
%---------------------------------------------------------------------------------------------------------------------------

\begin{definition}
    Soit \( P\in\eK[X]\) un polynôme irréductible. Une extension \( \eL\) de \( \eK\) est un \defe{corps de rupture}{corps!de rupture}\index{rupture!corps} pour \( P\) si il existe \( a\in \eL\) tel que \( P(a)=0\) et \( \eL=\eK(a)\).
\end{definition}

\begin{example}     \label{ExemGVxJUC}
    Soit \( \eK=\eQ\) et \( P(X)=X^2-2\). On pose \( a=\sqrt{2}\) et \( \eL=\eQ(\sqrt{2})\subset\eR\). De cette façon \( P\) est scindé :
    \begin{equation}
        P=(X-\sqrt{2})(X+\sqrt{2}).
    \end{equation}
    Le corps \( \eQ(\sqrt{2})\) est donc un corps de rupture pour \( P\)
\end{example}

\begin{example}
    Dans l'exemple \ref{ExemGVxJUC}, le polynôme \( P\) était scindé dans son corps de rupture. Il n'en est pas toujours ainsi. Prenons 
    \begin{equation}
        P(X)=X^3-2
    \end{equation}
    et \( a=\sqrt[3]{2}\). Nous avons, certes, \( P(a)=0\) dans \( \eQ(\sqrt[3]{2})\), mais \( P\) n'est pas scindé parce qu'il y a deux racines complexes.
\end{example}

%---------------------------------------------------------------------------------------------------------------------------
\subsection{Corps de décomposition}
%---------------------------------------------------------------------------------------------------------------------------

\begin{definition}
    Soit \( \eK\) un corps commutatif et \( F=(P_i)_{i\in I}\) une famille d'éléments non constants de \( \eK[X]\). Un \defe{corps de décomposition}{corps!de décomposition}\index{décomposition!corps} de \( F\) est une extension \( \eL\) de \( \eK\) telle que
    \begin{enumerate}
        \item
            les \( P_i\) sont scindés sur \( \eL\),
        \item
            \( \eL=\eK(R)\) où \( R=\bigcup_{i\in I}\{ x\in\eL\tq P_i(x)=0 \}\).
    \end{enumerate}
    C'est à dire que \( \eL\) étends \( \eK\) par toutes les racines de tous les polynômes de \( F\).
\end{definition}

L'unicité est due à la proposition suivante.
\begin{proposition}     \label{PropTMkfyM}
    Soit \( \eK\) un corps et \( P\in\eK[X]\). Soient \( \eL\) et \( \eF\) deux corps de décomposition de \( P\). Alors il existe un isomorphisme \( f\colon \eL\to \eF\) tel que \( f|_{\eK}=\id\).
\end{proposition}
Nous pouvons donc parler du corps de décomposition d'un polynôme.

%+++++++++++++++++++++++++++++++++++++++++++++++++++++++++++++++++++++++++++++++++++++++++++++++++++++++++++++++++++++++++++
\section{Corps finis}
%+++++++++++++++++++++++++++++++++++++++++++++++++++++++++++++++++++++++++++++++++++++++++++++++++++++++++++++++++++++++++++
\label{SecCorpsFinizkAcbS}

%---------------------------------------------------------------------------------------------------------------------------
\subsection{Existence, unicité}
%---------------------------------------------------------------------------------------------------------------------------

Nous avons déjà défini le corps fini \( \eF_p\) lorsque \( p\) est un nombre premier dans la section \ref{subseccorpspremhBlYIv}. Le théorème suivant sert à définir \( \eF_{p^n}\)\nomenclature[A]{\( \eF_{p^n}\)}{corps fini à \( p^n\) éléments} lorsque \( p\) est premier.
\begin{theorem}     \label{ThoOzgSfy}
    Soit \( p\) un nombre premier, soit \( n\in \eN^*\) et \( q=p^n\). Alors il existe un unique corps \( \eK\) de cardinal \( q\). Ce corps est le corps de décomposition du polynôme \( X^q-X\) sur \( \eF_p\).
\end{theorem}

\begin{proof}
    Montrons l'unicité. Soit \( \eK\) un corps fini de cardinal \( q=p^n\). Le groupe multiplicatif \( \eK^*\) est de cardinal \( q-1\), et par le corollaire \ref{CorpZItFX} tous les éléments de \( \eK^*\) vérifient \( g^{q-1}=e\), c'est à dire que dans \( \eK[X]\), les éléments de \( \eK^*\) sont des racines du polynôme
    \begin{equation}
        X^{q-1}-1
    \end{equation}
    Par conséquent \( \eK\) est un corps de décomposition pour le polynôme \( Q(X)=X^q-X=X(X^{q-1}-1)\) parce que \( Q(X)=0\) dans \( \eK\). Il est unique par la proposition \ref{PropTMkfyM}.

    Montrons maintenant que le corps de décomposition de \( P=X^q-X\) sur \( \eF_p\) est un corps de cardinal \( q\). Pour ce faire nous considérons \( \eK\) ce corps de décomposition et \(\eE\), l'ensemble des racines de \( P\) dans \( \eK\). Nous allons montrer que \( \eE=\eK\) et que \( \eE\) est un corps contenant \( q\) éléments.

    Montrons que \( \eE\) est un corps. Pour \( \alpha,\beta\in \eE\) nous avons
    \begin{equation}
        (\alpha\beta)^q=\alpha^q\beta^q=\alpha\beta
    \end{equation}
    parce que \( \alpha^q=\alpha\). Le produit \( \alpha\beta\) est donc encore dans \( \eE\). Pour la somme,
    \begin{equation}
        (\alpha+\beta)^q=(\alpha+\beta)^{p^n}=\Big( (\alpha+\beta)^p \Big)^{p^{n-1}}=(\alpha^p+\beta^p)^{p^{n-1}}=\ldots=\alpha^{p^n}+\beta^{p^n}=\alpha+\beta.
    \end{equation}
    En ce qui concerne l'inverse,
    \begin{equation}
        (\alpha^{-1})^q=(\alpha^q)^{-1}=\alpha^{-1}.
    \end{equation}
    Donc \( \eE\) est un corps. Évidemment \( \eE\) est un corps de décomposition de \( P\) au sens où \( \eE\) est une extension de \( \eF_p\) sur lequel \( P\) est scindé (parce qu'il est scindé sur \( \eK\) et \( \eE\) est le sous corps de \( \eK\) contenant les racines de \( P\)) et tel que \( \eE=\eF_p(\{ \alpha_i \})\) où les \( \alpha_i\) sont les racines de \( P\). Notons que \( \eF_p\subset \eE\) parce que dans \( \eF_p\) on a \( x^q=x\).

    Par unicité, nous avons \( \eK=\eE\). Nous devons montrer que \( P\) possède exactement \( q\) racines distinctes, affin d'avoir \( \Card(\eE)=q\). Pour cela remarquons que 
    \begin{equation}
        P'(X)=qX^{q-1}-1=-1
    \end{equation}
    dans \( \eF_p\). En effet \( P\in\eF_p\) et \( q=0\) dans \( \eF_p\). Par conséquent \( P'\) ne s'annule pas et \( P\) n'a pas de racines doubles. Toutes les racines étant simples, il y en a exactement \( q\).

\end{proof}

Le théorème \ref{ThoOzgSfy} ne permet pas de \emph{construire} le corps à \( q=p^n\) éléments. Nous allons maintenant voir un certain nombre de résultats donnant des façons de construire. Ces résultats proviennent de \cite{MichelMerlecorpsfinis,GabrielPeyre,RodierCorpsFinis} et de \wikipedia{fr}{Théorème_de_l'élément_primitif}{wikipedia} 

\begin{proposition}     \label{PropnfebjI}
    Soit \( \eK\) un corps fini. Alors le groupe multiplicatif \( \eK^*\) est cyclique.
\end{proposition}

\begin{proof}
    Soit \( \eK\) un corps ayant \( q\) éléments. Le groupe \( \eK^*\) en a \( q-1\), de telle façon à ce que l'ordre des éléments de \( \eK^*\) soient des diviseurs de \( q-1\); c'est le corollaire \ref{CorpZItFX}. Soit \( d\) un diviseur de \( q-1\) et
    \begin{subequations}
        \begin{align}
            H^*_d&=\{ x\,\text{d'ordre \( d\) dans \( \eK^*\)} \}\\
            H_d&=\{ \text{racines de \( X^d-1\) dans \( \eK\)} \}.
        \end{align}
    \end{subequations}
    Ici le polynôme \( X^d-1\) est vu dans \( \eK[X]\). Notons que nous avons automatiquement \( H^*_d\subset H_d\), mais l'inclusion inverse n'est pas assurée parce que les éléments d'ordre \( d/2\) par exemple sont aussi dans \( H_d\). Supposons \( H^*_d\neq \emptyset\) et considérons \( a\in H^*_d\). Alors l'application
    \begin{equation}
        \begin{aligned}
            \phi\colon \eZ/d\eZ&\to H_d \\
            n&\mapsto a^n 
        \end{aligned}
    \end{equation}
    est un isomorphisme d'anneaux. En effet étant donné que \( a\in H^*_d\subset H_d\), l'ensemble \( H_d\) contient le groupe cyclique engendré par \( a\). Ce dernier contient, par construction, \( d\) éléments. Mais \( \Card(H_d)\leq d\) parce que \( H_d\) est l'ensemble des racines d'un polynôme de degré \( d\). Par conséquent \( \Card(H_d)=d\) et l'ensemble \( H_d\) est bien engendré par \( a\) et \( \phi\) est bien un isomorphisme. Par conséquent tous les éléments de \( H^*_d\) sont des générateurs de \( H_d\).

    Inversement soit \( x\) un générateur de \( H_d\). L'ordre de \( H_d\) étant \( d\), l'ordre de \( x\) doit être un diviseur de \( d\). Supposons donc que \( x\) soit d'ordre \( d/k\). Dans ce cas nous devrions avoir \( \Card(H_d)=d/k\), ce qui contredit l'isomorphisme \( \phi\).

    En conclusion, \( H^*_d\) est l'ensemble des générateurs du groupe \( H_d\). Le nombre de générateurs de \( \eZ/d\eZ\) étant \( \varphi(d)\) par la proposition \ref{PropZnmuphiGensn}, et \( H_d\) étant isomorphe à \( \eZ/d\eZ\) nous avons
    \begin{equation}
        \Card(H^*_d)=\varphi(d).
    \end{equation}
    
    Par conséquent si \( H^*_d\) n'est pas vide, son cardinal est \( \varphi(d)\). Nous avons 
    \begin{subequations}
        \begin{align}
            q-1&=\Card(\eK^*)\\
            &=\Card\big( \bigcup_{d\divides q-1}H^*_d \big)\\
            &=\sum_{d\divides q-1}\Card(H^*_d)\\
            &\leq \sum_{d\divides q-1}\varphi(d)\\
            &=q-1
        \end{align}
    \end{subequations}
    où nous avons utilisé le lemme \ref{LemKcpjee}. Par conséquent pour tout \( d\) divisant \( q-1\) nous avons \( \Card(H^*_d)=\varphi(d)\) et il y a au moins un élément d'ordre \( q-1\) dans \( \eK\). Cet élément engendre \( \eK^*\) parce que \( \eK^*\) contient exactement \( q-1\) éléments. Par conséquent \( \eK\) est cyclique.
\end{proof}

Lorsque \( \eK\) est un corps les éléments du groupe \( \eK^*\) sont les \defe{éléments primitifs}{primitif!élément d'un corps} de \( \eK\).

\begin{proposition}     \label{propQRcUlq}
    Soit \( \eK\) un corps contenant \( q\) éléments. Alors
    \begin{enumerate}
        \item
            \( x^q=x\) pour tout \( x\in \eK\),
        \item
            \( X^q-X=\prod_{a\in \eK}(X-a)\).
    \end{enumerate}
\end{proposition}

\begin{proof}
    Le groupe \( \eK^*\) ayant \( q-1\) éléments, ses éléments vérifient \( a^{q-1}=1\) par le corollaire \ref{CorpZItFX} et par conséquent \( a^q=aa^{q-1}=a \).

    Soit \( a\in \eK\). Étant donné que \( a^q-a=0\), le polynôme \( (X-a)\) divise \( X^q-X\) dans \( \eK[X]\). Par conséquent 
    \begin{equation}
        \prod_{a\in \eK}(X-a)
    \end{equation}
    divise également \( X^q-X\). Les polynômes \( X^q-X\) et \( \prod_{a\in \eK}(X-a)\) étant deux polynômes unitaires de même degré, le fait que l'un divise l'autre montre qu'ils sont égaux.
\end{proof}

\begin{example}
    Soit \( \eK=\eQ\) et \( \eL=\eQ(\sqrt{2},\sqrt{3})\). Afin de montrer que \( \eL=\eQ(\alpha)\) avec \( \alpha=\sqrt{2}+\sqrt{3}\) nous devons montrer que \( \sqrt{2}\) et \( \sqrt{3}\) sont des polynômes en \( \alpha\).
\end{example}

\begin{definition}
    Soit \( \eK\) un corps. Une extension \( \eL\) de \( \eK\) est dite \defe{finie}{extension!finie} si \( \eL\) est un espace vectoriel de dimension finie sur \( \eK\).
\end{definition}
Notez que la définition d'extension finie ne suppose ni que \( \eK\) ni que \( \eL\) soient finis en tant qu'ensembles.

\begin{theorem}[de l'élément primitif]\index{théorème!élément primitif}
    Si \( \eK\) est un corps fini, toute extension finie de \( \eK\) est simple.

    Si \( \eK\) est un corps quelconque alors toute extension séparable finie est simple.
\end{theorem}

\begin{proof}
    Nous ne donnons la preuve que dans le cas où \( \eK\) est fini. Dans ce cas nous savons par la proposition \ref{PropnfebjI} que le groupe \( \eK^*\) est cyclique. Si de plus \( \eL\) est une extension finie alors \( \eL\) est fini en tant qu'ensemble. Par conséquent \( \eL^*\) est un groupe cyclique. Si \( \alpha\) est un générateur de \( \eL\) alors \( \eL=\eK(\alpha)\) et l'extension est donc simple.

    Une preuve de l'assertion dans le cas où \( \eK\) est infini peut être trouvée sur wikipédia.
\end{proof}

\begin{definition}\label{DefvBFpsY}
    Soit \( P\) un polynôme irréductible de degré \( n\) sur \( \eF_p[X]\). L'\defe{ordre}{ordre!d'un polynôme} de \( P\) est
    \begin{equation}
        \min\{ k\tq P\divides X^k-1 \}.
    \end{equation}

    Soit \( p\), un nombre premier et \( P\) un polynôme unitaire irréductible de degré $n$ dans \( \eF_p[X]\). Nous disons que \( P\) est \defe{primitif}{primitif!polynôme}\index{polynôme!primitif} si les racines de \( P\) sont d'ordre \( p^n-1\) dans \( \eF_p[X]/P\).
\end{definition}

\begin{remark}
    Si \( \eA\) est un anneau factoriel, \wikipedia{fr}{Polynôme}{il est souvent dit} qu'un polynôme \( P\in \eA[X]\) est primitif si le pgcd de ses coefficients est \( 1\). Cette notion de primitivité n'est apparemment pas reliée à celle de la définition \ref{DefvBFpsY} qui sera celle que nous retiendrons pour la suite. Notons au passage que dans un corps, tous les polynômes non nuls et unitaires sont primitifs au sens du pgcd des coefficients; cette notions de primitivité n'a donc pas d'intérêt dans notre cadre.
\end{remark}

\begin{proposition}
    L'ordre d'un polynôme \( P\) vérifie les propriétés suivantes :
    \begin{enumerate}
        \item
            L'ordre de \( P\) est l'ordre multiplicatif de ses racines
        \item
            L'ordre de \( P\) divise \( p^n-1\).
    \end{enumerate}
\end{proposition}

\begin{lemma}
    Si \( \alpha\in \eF_q\) est une racine d'ordre \( k\) de \( P\) (de degré \( n\)) alors les racines de \( X^k-1\) sont \( \{ \alpha^i\tq i=0,\ldots, k-1 \}\).
\end{lemma}

Nous serions donc intéressé à construire $\eF_{q}$ comme quotient de \( \eF_p[X]\) par un polynôme primitif. Le théorème suivant donne une description abstraite de \( \eF_q\) qui va nous servir de point de départ pour la construction.
\begin{theorem}[Théorème de l'élément primitif]\index{théorème!élément primitif}    \label{ThoqSludu}
    Soit \( p\) un nombre premier, \( n\in \eN\) et \( q=p^n\). Soit \( \eK\) un corps à \( q\) éléments. Alors
    \begin{enumerate}
        \item
            Il existe \( \alpha\in \eK\) tel que \( \eK=\eF_p[\alpha]\).
        \item
            Il existe une polynôme irréductible \( P\in\eF_p[X]\) de degré \( n\) tel que
            \begin{equation}
                \begin{aligned}
                    \phi\colon \eF_p[X]/(P)&\to \eK \\
                    \bar X&\mapsto \alpha 
                \end{aligned}
            \end{equation}
            soit un isomorphisme de corps.
    \end{enumerate}
    Soit \( \alpha\) et \( P\) choisis pour avoir les propriétés citées plus haut. Alors nous avons les propriétés suivantes.
    \begin{enumerate}
        \item
            \( P\) est primitif.
        \item
            \( P\) est scindé dans \( \eK\).
        \item
            L'ensemble des racines de \( P\) est \( \{ \alpha,\alpha^p,\ldots, \alpha^{p^{n-1}} \}\).
        \item
            Le polynôme \( P\) divise \( X^q-X\) dans \( \eF_p[X]\).
    \end{enumerate}
\end{theorem}

\begin{proof}
    Le corps \( \eK\) étant fini, il est cyclique par la proposition \ref{PropnfebjI}. Si \( \alpha\) un générateur de \( \eK^*\) alors
    \begin{equation}
        \eK=\eF_p[\alpha].
    \end{equation}
    Soit \( \ell\) le plus grand entier tel que l'ensemble
    \begin{equation}
        \{ 1,\alpha,\cdots,\alpha^{\ell-1} \}\subset\eK
    \end{equation}
    soit libre. Pour rappel \( \eK\) est une espace vectoriel sur \( \eF_p\). Il existe des \( a_i\in \eF_p\) tels que
    \begin{equation}
        \alpha^{\ell}+a_{\ell-1}\alpha^{\ell-1}+\ldots+a_0=0.
    \end{equation}
    De façon équivalente, il existe un polynôme unitaire \( P\in\eF_p[X]\) de degré \( \ell\) tel que \( P(\alpha)=0\). Étant donné que \( \alpha\) est générateur de \( \eK\),
    \begin{equation}
        \eK=\Span\{ 1,\alpha,\ldots, \alpha^{\ell-1} \}
    \end{equation}
    parce que \( \eK\) est généré par les puissances de \( \alpha\) alors que les puissances de \( \alpha\) plus hautes que \( \ell-1\) peuvent être générées par \( 1,\alpha,\ldots, \alpha^{\ell-1}\). L'espace \( \eK\) est donc un \( \eF_p\)-espace vectoriel de dimension \( \ell\); par conséquent
    \begin{equation}
        \Card(\eK)=p^n=q
    \end{equation}
    et \( \ell=n\).

    Montrons que \( P\) est irréductible dans \( \eF_p\). Si \( P\) était réductible dans \( \eF_p\), l'élément \( \alpha\in \eK\) serait une racine d'un des facteurs, c'est à dire qu'il serait racine d'un polynôme de degré inférieur à \( n\), ce qui contredirait le fait que 
    \begin{equation}
        \{ \alpha^{\ell-1},\ldots, 1 \}
    \end{equation}
    est libre.

    Montrons que l'application
    \begin{equation}
        \begin{aligned}
            \infty\colon \eF_p[X]/(P)&\to \eK \\
            \bar X&\mapsto \alpha 
        \end{aligned}
    \end{equation}
    est un isomorphisme. Pour l'injectivité, deux éléments \( Q_1,Q_2\in \eF_p[X]/(P)\) s'écrivent
    \begin{subequations}
        \begin{align}
            Q_1&=\sum_{k=0}^{n-1}a_k\bar X^k\\
            Q_2&=\sum_{k=0}^{n-1}b_k\bar X^k.
        \end{align}
    \end{subequations}
    Dans ce cas si \( \phi(Q_1)=\phi(Q_2)\) alors
    \begin{equation}
        \phi(Q_1)=\sum_{k=0}^{n-1}a_k\alpha^k=\phi(Q_2)=\sum_{k=0}^{n-1}b_k\alpha^k.
    \end{equation}
    Mais l'ensemble \( \{ 1,\alpha,\ldots, \alpha^{n-1} \}\) étant libre sur \( \eF_p\), cela implique \( a_k=b_k\). La surjectivité de \( \phi\) provient du fait que \( \alpha\) génère \( \eK\).

    Nous passons maintenant à la seconde partie de la démonstration. Soient \( \alpha\in \eK\) tel que \( \eK=\eF_p[\alpha]\) et \( P\in \eF_p[X]\) un polynôme irréductible de degré \( n\) tel que \( \alpha\mapsto \bar X\) soit un isomorphisme entre \( \eK\) et \( \eF_p[X]/(P)\).

    Le polynôme \( P\) est primitif parce que \( \alpha\) es d'ordre \( p^n\) dans \( \eK\) alors que \( \bar X\mapsto \alpha\) est un isomorphisme. Par conséquent \( \bar X\) est d'ordre \( p^n\) dans \( \eF_p[X]/P\).

    Nous commençons par prouver que l'ensemble
    \begin{equation}        \label{EqAcsQHL}
        \{ \alpha,\alpha^p,\alpha^{p^2},\ldots, \alpha^{p^{n-1}} \}
    \end{equation}
    est l'ensemble des racines distinctes de \( P\). Pour cela nous posons
    \begin{equation}
        P(X)=\sum_{k=0}^na_kX^k
    \end{equation}
    avec \( a_k\in\eF_p\). D'abord \( \alpha\) est une racine de \( P\). En effet
    \begin{equation}        \label{EqbTAmKG}
        P(\bar X)=\sum_ka_k\bar X^k=0
    \end{equation}
    parce que cette somme est calculée dans \( \eF_p[X]/(P)\). En appliquant l'isomorphisme \( \phi\) à l'égalité \eqref{EqbTAmKG} nous trouvons
    \begin{equation}
        0=\phi\big( P(\bar X) \big)=\sum_ka_k\phi(\bar X^k)=\sum_ka_k\alpha^k.
    \end{equation}
    Donc \( \alpha\) est bien une racine de \( P\) dans \( \eF_p[X]\). Nous devons montrer qu'il en est de même pour les autres puissances dans l'ensemble \eqref{EqAcsQHL}. Étant donné que pour tout \( x\) dans \( \eF_p\) nous avons \( x^p=x\), nous avons aussi
    \begin{equation}
        P(X^p)=\sum_ka_k(X^p)^k=\sum_ka_k^p(X^p)^k=\sum_k(a_kX^k)^p
    \end{equation}
    alors que nous savons que \( x\mapsto x^p\) est un automorphisme de \( \eF_p\) par la proposition \ref{PropFrobHAMkTY}. Par conséquent
    \begin{equation}
        P(X^p)=\sum_k(a_kX^k)^p=\left( \sum_k a_kX^k\right)^p=P(X)^p.
    \end{equation}
    Nous avons montré que si \( \beta\) est une racine de \( P\), alors \( \beta^p\) est également une racine de \( P\). Nous savons déjà que \( \alpha\) est une racine de \( P\), et que \( \alpha\) est également générateur de \( \eK\), c'est à dire que \( \alpha\) est d'ordre \( q-1\). Les puissances
    \begin{equation}
        \alpha,\alpha^p,\alpha^{p^2},\ldots, \alpha^{p^{n-1}}
    \end{equation}
    sont donc distinctes (\( \alpha^{p^n}=\alpha^q=1\)) et sont toutes des racines de \( P\). Étant donné que \( P\) est de degré \( n\) il ne peut pas y avoir d'autres racines. Nous concluons que l'ensemble
    \begin{equation}
        \{ \alpha,\alpha^p,\alpha^{p^2},\ldots, \alpha^{p^{n-1}}\}
    \end{equation}
    est l'ensemble des racines distinctes de \( P\) dans \( \eK\). Le polynôme \( P\) est alors scindé dans \( \eK[X]\).

    Le dernier point du théorème est de montrer que \( P\) divise \( X^q-X\). Pour cela nous allons montrer que toutes les racines de \( P\) sont des racines de \( X^q-X\). Soit \( \beta\) une racine de \( P\); il s'écrit \( \beta=\alpha^k\) pour un certain \( k\). Étant donné que \( \alpha^{q-1}=e=\alpha^{p^n-1}\),
    \begin{subequations}
        \begin{align}
            \beta^q&=(\alpha^{p^n})^k\\
            &=\left( \alpha^{p^n-1}\alpha \right)^k\\
            &=\left( \alpha^{q-1}\alpha \right)^k\\
            &=\alpha^k\\
            &=\beta.
        \end{align}
    \end{subequations}
    Cela signifie que \( \beta^q=\beta\) et donc que \( \beta\) est racine de \( X^q-X\).
\end{proof}

\begin{lemma}       \label{Lembcerei}
    Si \( \alpha\) est une racine primitive de \( P\) alors les autres racines sont également primitives.
\end{lemma}

\begin{proof}
    Soit \( \alpha\in \eF_p[X]/P\) une racine primitive de \( P\). L'élément \( \alpha^p\) est également une racine parce que si \( P=\sum_ka_kX^k\),
    \begin{equation}
        P(\alpha^p)=\sum_k(a_k\alpha^k)^p=\big( \sum_ka_k\alpha^k \big)^p=0
    \end{equation}
    où nous avons utilisé le fait que \( a_k^p=a_k\) étant donné que \( a_k\in\eF_p\). Par hypothèse \( \alpha\) est une racine primitive et les éléments \( \alpha,\alpha^p,\alpha^{p^2},\ldots,\alpha^{p^n-1}\) sont des racines distinctes de \( P\). Elles sont donc \emph{toutes} les racines de \( P\).

    Toutes les racines de \( P\) sont donc primitives parce que pour la racine \( \beta=\alpha^{p^k}\) nous avons
    \begin{equation}
        \alpha^{p^l}=\beta^{p^{l-k}}
    \end{equation}
    c'est à dire que la racine \( \alpha^l\) est bien de la forme \( \beta^{p^{r}}\) pour un certain \( r\).
\end{proof}

\begin{lemma}       \label{LemkzWjse}
    Soit \( p\) un nombre premier et \( n\), un entier. Un polynôme de degré \( d\) irréductible dans \( \eF_p[X]\) divise \( X^{p^n}-X\) si et seulement si \( d\) divise \( n\).
\end{lemma}

\begin{theorem}
    Soient \( P\) et \( Q\) deux polynômes irréductibles de degré \( n\) dans \( \eF_p[X]\). Alors les quotients \( \eF_p[X]/P\) et \( \eF_p[X]/Q\) sont isomorphes en tant que corps.
\end{theorem}
En guise de démonstration de ce théorème, nous allons démontrer la proposition suivante.
\begin{proposition}
    Si \( \eK\) et \( \eL\) sont deux corps à \( q=p^n\) éléments, alors ils sont isomorphes.
\end{proposition}

\begin{proof}
    Soit \( a\) un élément primitif de \( \eK\) et \( P\) son polynôme minimal. Nous savons que \( \eK\simeq \eF_p[X]/P\) par le théorème de l'élément primitif \ref{ThoqSludu}. L'élément \( a\) est en particulier une racine de \( X^q-X\). Par ailleurs \( P\) divise \( X^q-X\) par le lemme \ref{LemkzWjse}.

    Nous avons aussi
    \begin{equation}
        X^q-X=\prod_{b\in \eL}(X-b)
    \end{equation}
    par la proposition \ref{propQRcUlq}. Étant donné que \( P\) divise \( X^q-X\), un des éléments de \( \eL\) annule \( P\). Soit \( b\in \eL\) tel que \( P(b)=0\). Soit \( Q\) le polynôme minimal de \( b\). Par définition nous avons que \( Q\) divise \( P\), mais \( P\) étant irréductible et unitaire nous avons immédiatement \( P=Q\). En particulier nous avons
    \begin{equation}
        \eF_p[X]/P\simeq\eF_p[X]/Q\simeq \eK.
    \end{equation}
    Nous montrons maintenant que \( \eF_p[X]/Q\simeq \eL\) par l'application
    \begin{equation}
        \begin{aligned}
            \phi\colon \eF_p[X]/Q&\to \eL \\
            \bar X&\mapsto b 
        \end{aligned}
    \end{equation}
    qui se prolonge en \( R(\bar X)\mapsto R(b)\) pour tout \( R\in \eF_p[X]\). Cette application est bien définie parce que \( Q(b)=0\). Elle est injective parce que \( R(b)=0\) ne peut pas avoir lieu avec \( R\in \eF_p[X]/Q\) parce que \( Q\) est le polynôme minimal de \( b\). La surjectivité vient alors du fait que les deux corps ont le même nombre d'éléments.
\end{proof}

%---------------------------------------------------------------------------------------------------------------------------
\subsection{Construction de $\eF_q$}
%---------------------------------------------------------------------------------------------------------------------------

Soit \( p\) un nombre premier et \( n\in \eN\). Nous souhaitons construire un corps à \( q=p^n\) éléments. Nous savons déjà que ce corps est unique (théorème \ref{ThoOzgSfy}) et que nous le notons \( \eF_q\). Le théorème \ref{ThoqSludu} nous incite à l'écrire sous la forme
\begin{equation}
    \eF_q=\eF_p[X]/(P)
\end{equation}
pour un certain polynôme irréductible \( P\in\eF_p[X]\).

%///////////////////////////////////////////////////////////////////////////////////////////////////////////////////////////
\subsubsection{La version du faignant}
%///////////////////////////////////////////////////////////////////////////////////////////////////////////////////////////

Nous pouvons construire le corps à \( p^n\) éléments en prenant le quotient de \( \eF_p[X]\) par n'importe quel polynôme irréductible de degré \( n\). Le résultat est le suivant.
\begin{proposition}
    Soit \( P\) une polynôme unitaire irréductible dans \( \eF_p[X]\). Nous posons \( \eK=\eF_p[X]/(P)\). Alors
    \begin{enumerate}
        \item
            \( \eK\) est un corps à \( q\) éléments.
        \item
            \( \alpha=\bar X\) est une racine de \( P\) dans \( \eK\).
        \item
            \( \eK=\eF_p[\alpha]\).
    \end{enumerate}
\end{proposition}

\begin{proof}
    \begin{enumerate}
        \item
            En vertu du corollaire \ref{CorsLGiEN}, \( \eK\) est un corps. Il est aussi un espace vectoriel de dimension \( n\) sur \( \eF_p\), et contient donc \( p^n=q\) éléments.
        \item
            Par construction de \( \eK/(P)\).
        \item
            En tant que quotient de \( \eF_p[X]\), les éléments de \( \eK\) sont des polynômes en \( \bar X\).
    \end{enumerate}
\end{proof}

%///////////////////////////////////////////////////////////////////////////////////////////////////////////////////////////
\subsubsection{La version plus élaborée}
%///////////////////////////////////////////////////////////////////////////////////////////////////////////////////////////

Construire \( \eF_q\) comme quotient de \( \eF_p[X]\) par un polynôme irréductible quelconque ne donne pas d'informations sur les générateurs de \( \eF_q^*\), et en particulier il n'est pas toujours vrai que \( \bar X\) est générateur.

\begin{example}     \label{ExemWUdrcs}
    Cherchons à construire \( \eF_{16}\) comme quotient de \( \eF_2\) par un polynôme de degré \( 4\).
    \begin{verbatim}
----------------------------------------------------------------------
| Sage Version 4.7.1, Release Date: 2011-08-11                       |
| Type notebook() for the GUI, and license() for information.        |
----------------------------------------------------------------------
sage: x=polygen(GF(2))
sage: -x-1
x + 1
sage: Q=x**15-1
sage: Q.factor()
(x + 1) * (x^2 + x + 1) * (x^4 + x + 1) * (x^4 + x^3 + 1) * (x^4 + x^3 + x^2 + x + 1)
    \end{verbatim}
    Les polynômes candidats a avoir des racines génératrices sont donc au nombre de \( 3\):
    \begin{subequations}
        \begin{align}
            P_1&=X^4+X+1\\
            P_2&=X^4+X^3+1\\
            P_3&=X^4+X^3+X^2+X+1.
        \end{align}
    \end{subequations}
    Dans le quotient \( \eF_2[X]/P_3\), l'élément \( \bar X\) n'est pas générateur. En effet nous avons \( X^4=X^3+X^2+X+1\) et par conséquent les puissances successives de \( X\) sont
    \begin{subequations}
        \begin{align}
            X&\\
            X^2&\\
            X^3&\\
            X^4&=X^3+X^2+X+1\\
            1.
        \end{align}
    \end{subequations}
    La classe de \( X\) dans \( \eF_2[X]/P_3\) n'est donc pas génératrice du groupe \( (\eF_2[X]/P_3)^*\).

    Le polynôme \( P_1=X^4+X+1\) par contre est primitif parce que les puissances de \( X\) dans \( \eF_2[X]/P_1\) sont
    \begin{subequations}
        \begin{align}
            X\\
            X^2\\
            X^3\\
            X+1\\
            X^2+X\\
            X^3+X^2\\
            X+1+X^3\\
            X^2+1\\
            X^3+X\\
            X+1+X^2\\
            X^2+X+X^3\\
            X^3+X^2+X+1\\
            1+X^2+X^2\\
            1+X^3\\
            1
        \end{align}
    \end{subequations}
    Cela fait \( 15\) puissances distinctes, ce qui prouve que \( P_1\) est primitif. Nous verrons un peu plus loin comment alléger un peu la vérification de la primitivité de \( P_1\).
\end{example}

\begin{proposition}     \label{PropNsLqWb}
    Soit \( P\) un polynôme irréductible unitaire primitif dans \( \eF_p[X]\). Nous considérons \( \eK=\eF_p/P\) et \( \alpha=\bar X\in \eK\). Alors
    \begin{enumerate}
        \item
            \( P\) est le polynôme minimal de \( \alpha\).
        \item
            Les racines de \( P\) sont \( \{ \alpha,\alpha^p,\ldots, \alpha^{p^{n-1}} \}\) et \( \alpha^q=\alpha\).
        \item
            \( P\) est scindé dans \( \eK\).
        \item
            \( P\) divise \( X^q-X\) dans \( \eK\).
        \item
            La famille \( \{1, \alpha,\alpha^2,\ldots, \alpha^{n-1} \}\) est une base de \( \eK\) en tant qu'espace vectoriel sur \( \eF_p\).
        \item
            En tant qu'ensemble,
            \begin{equation}
                \eF_q=\{0, \alpha,\alpha^2,\alpha^3,\ldots, \alpha^{q-1} \},
            \end{equation}
            et les \( \alpha^k\) sont distincts pour \( k=1,\ldots, q-1\).
    \end{enumerate}
\end{proposition}

\begin{proof}
    La plupart des assertions sont des corollaires ou des paraphrases de résultats contenus dans les propositions précédentes.
    \begin{enumerate}
        \item
            Le fait que \( P(\bar X)=0\) est la définition du quotient de \( \eF_p[X]\) par l'idéal engendré par \( P\). Supposons avoir un polynôme \( \tilde P\) tel que \( \tilde P(\alpha)=0\) et \( \deg(\tilde P)<\deg(P)\). Par division euclidienne (théorème \ref{ThodivEuclPsFexf}) nous avons des polynômes \( Q\) et \( R\) tels que
            \begin{equation}        \label{EqPWaRUl}
                P=\tilde PQ+R
            \end{equation}
            avec \( \deg(R)<\deg(\tilde P)\). En appliquant l'égalité \( \eqref{EqPWaRUl}\) à \( \alpha\) nous trouvons \( R(\alpha)=0\). Par récurrence nous avons des polynômes annulateurs de \( \alpha\) de degré arbitrairement bas, ce qui est absurde.
        \item
            L'assertion à propos des racines de \( P\) est contenue dans le lemme \ref{Lembcerei}. Par le corollaire \ref{CorpZItFX}, dans un groupe cyclique d'ordre \( q\), tous les éléments vérifient \( x^q=x\).
        \item
            Possédant \( n\) racines distinctes dans \( \eK\), le polynôme \( P\) est scindé.
        \item
            Le lemme \ref{LemxneRcs} un polynôme irréductible de degré \( n\) divise le polynôme \( X^{p^n}-X\). Une autre façon de montrer ce point est de remarquer que le polynôme \( P\) est scindé et que toutes ses racines sont également racines de \( X^q-X\).
        \item
            Une combinaison linéaire nulle entre les éléments de \( \{ 1,\alpha,\alpha^2,\ldots, \alpha^{n-1} \}\) serait un polynôme annulateur de degré \( n-1\) de \( \alpha\). Cet ensemble est donc libre. Par ailleurs un ensemble libre de \( n\) éléments dans un espace vectoriel de dimension \( n\) est générateur.
        \item
            Si \( \alpha^l=\alpha^k\) avec \( k<l\) et \( k,l\leq q\) alors nous avons \( \alpha^r=1\) avec \( r=l-k<q\), ce qui contredirait la primitivité de \( P\). Les éléments \( 0,\alpha,\ldots, \alpha^{q-1}\) étant distincts et au nombre de \( q\), ils forment tout l'ensemble \( \eF_q\).

    \end{enumerate}
\end{proof}

%---------------------------------------------------------------------------------------------------------------------------
\subsection{Exemple : étude de \texorpdfstring{$\eF_{16}$}{F16}}
%---------------------------------------------------------------------------------------------------------------------------




%+++++++++++++++++++++++++++++++++++++++++++++++++++++++++++++++++++++++++++++++++++++++++++++++++++++++++++++++++++++++++++
\section{Mini introduction aux nombres \texorpdfstring{p}{$p$}-adiques}
%+++++++++++++++++++++++++++++++++++++++++++++++++++++++++++++++++++++++++++++++++++++++++++++++++++++++++++++++++++++++++++


\subsection{La flèche d'Achille}\label{s:un}

C'est un grand classique que je donne ici juste comme introduction pour montrer que des série infinies peuvent donner des nombres finis de manière tout à fait intuitive.

Achille tire une flèche vers un arbre situé à $\unit{10}{\meter}$ de lui. Disons que la flèche avance à une vitesse constante de $\unit{1}{\meter\per\second}$. Il est clair que la flèche mettra $\unit{10}{\second}$ pour toucher l'arbre. En $\unit{5}{\second}$, elle aura parcouru la moitié de son chemin. On le note :
\[
\text{temps}=5s+\ldots
\]
Reste \( \unit{5}{\meter}\) à faire. En $\unit{2.5}{\second}$, elle aura fait la moitié de ce chemin chemin, soit $2.5m=\frac{10}{4}m$. On le note :
\[
\text{temps}=\frac{10}{2}s+\frac{10}{4}s+
\]
Reste $2.5m$ à faire. La moitié de ce trajet, soit $\frac{10}{8}m$, est parcouru en $\frac{10}{8}s$; on le note encore, mais c'est la dernière fois !

\[
\text{temps}=\frac{10}{2}s+\frac{10}{4}s+\frac{10}{8}s+
\]
En continuant ainsi à regarder la flèche qui parcours des demi-trajets puis des demi de demi-trajets et encore des demi de demi de demi-trajets, et en sachant que le temps total est $10s$, on trouve :
\[
10\left( \frac{1}{2}+\frac{1}{4}+\frac{1}{8}+\frac{1}{16}+\ldots  \right)=10.
\]
On doit donc croire que la somme jusqu'à l'infini des inverse des puissances de deux vaut $1$ :
\[
   \sum_{n=1}^{\infty}\frac{1}{2^n}=1.
\]
Cela peut être démontré à la loyale.

\subsection{La tortue et Achille}

Maintenant qu'on est convaincu que des sommes infinies peuvent représenter des nombres tout à fait normaux, passons à un truc plus marrant.

Achille, qui marche peinard à $\unit{10}{\meter\per\hour}$, part avec $1m$ d'avance sur une tortue qui avance à $\unit{1}{\meter\per\hour}$. Le temps que la tortue arrive au point de départ d'Achille, Achille aura parcouru $10m$, et le temps que la tortue mettra pour arriver à ce point, eh bien, Achille ne sera déjà plus là : il sera à $100m$. Si la tortue tient bon pendant un temps infini, et si l'on est confiant en le genre de raisonnements faits à la section \ref{s:un}, elle rattrapera Achille dans 
\[
1m+10m+100m+1000m+\ldots
\]
Autant dire que ça ne risque pas d'arriver. Et pourtant, mettons en équations : 
\begin{subequations}
    \begin{numcases}{}
        x_{\text{Achile}}(t)=1+10t\\
        x_{\text{tortue}}(t)=t.
    \end{numcases}
\end{subequations}
La tortue rejoints Achille au temps \( t\) tel que \( x_{\text{Achille}(t)}=x_{\text{tortue}}(t)\). Un mini calcul donne $t=-1/9$. Physiquement, c'est une situation logique. Peut-on en déduire une égalité mathématique du style de 
\[
1+10+100+1000+\ldots=-\frac{1}{9}\; ???
\]
Là où les choses deviennent jolies, c'est quand on cherche à voir ce que peut bien être la valeur d'un hypothétique $x=1+10+100+1000+\ldots$. En effet, logiquement on devrait avoir
\begin{equation*}
\begin{split}
\frac{x}{10}&=\frac{1}{10}+1+10+100+\ldots\\
            &=\frac{1}{10}+x.
\end{split}
\end{equation*}
Reste à résoudre l'équation du premier degré : $\frac{x}{10}=x+\frac{1}{10}$. Ai-je besoin de donner la solution ?

%---------------------------------------------------------------------------------------------------------------------------
\subsection{Dans les nombres \texorpdfstring{p}{$ p$}-adiques, c'est vrai}
%---------------------------------------------------------------------------------------------------------------------------

Nous nous proposons d'apprendre sur les nombres \( p\)-adiques juste ce qu'il faut pour montrer que l'égalité
\begin{equation}
    \sum_{k=0}^{\infty}10^k=-\frac{1}{ 9 }
\end{equation}
est vraie dans les nombres \( 5\)-adiques. Tout ce qu'il faut est sur \wikipedia{fr}{Nombre_p-adique}{wikipedia}.

Soit \( a\in \eN\) et \( p\), un nombre premier. La \defe{valuation}{valuation!$p$-adique} \( p\)-adique de \( a\) est l'exposant de \( p\) dans la décomposition de \( a\) en nombres premiers. On la note \( v_p(a)\). Pour un rationnel on définit
\begin{equation}
    v_p\left( \frac{ a }{ b } \right)=v_p(a)-v_p(b)
\end{equation}
La \defe{valeur absolue}{valeur absolue!$p$-adique} \( p\)-adique de \( r\in \eQ\) est 
\begin{equation}
    | r |_p=p^{-v_p(r)}.
\end{equation}
Nous posons \( | 0 |_p=0\). De là nous considérons la distance
\begin{equation}
    d_p(x,y)=| x-y |_p.
\end{equation}

\begin{lemma}
    L'espace \( (\eQ,d_p)\) est un espace métrique.
\end{lemma}

Nous considérons maintenant \( p=5\). Étant donné que \( a=5\cdot 2\) nous avons \( v_5(10)=1\) et
\begin{equation}
    v_5\left( \frac{1}{ 9 } \right)=v_5(1)-v_5(9)=0.
\end{equation}
Nous avons
\begin{equation}
    \sum_{k=0}^N10^k+\frac{1}{ 9 }=\frac{ 10^{N-1} }{ 9 }
\end{equation}
mais
\begin{equation}
    v_p\left( \frac{ 10^{N-1} }{ 9 } \right)=v_5(10^{N-1})-v_5(9)=N-1.
\end{equation}
Par conséquent
\begin{equation}
    d_5\big( \sum_{k=0}^N,-\frac{1}{ 9 } \big)=| \frac{ 10^{N-1} }{ 9 } |_p=p^{-(N-1)}.
\end{equation}
En passant à la limite,
\begin{equation}
    \lim_{N\to \infty} d_5\big( \sum_{k=0}^N,-\frac{1}{ 9 } \big)=0,
\end{equation}
ce qui signifie que
\begin{equation}
    \sum_{k=0}^{\infty}10^k=-\frac{1}{ 9 }.
\end{equation}


\chapter{Espaces vectoriels}
% This is part of Mes notes de mathématique
% Copyright (c) 2011-2013
%   Laurent Claessens, Carlotta Donadello
% See the file fdl-1.3.txt for copying conditions.


%+++++++++++++++++++++++++++++++++++++++++++++++++++++++++++++++++++++++++++++++++++++++++++++++++++++++++++++++++++++++++++
\section{Espaces de matrices}
%+++++++++++++++++++++++++++++++++++++++++++++++++++++++++++++++++++++++++++++++++++++++++++++++++++++++++++++++++++++++++++

%---------------------------------------------------------------------------------------------------------------------------
\subsection{Connexité par arcs}
%---------------------------------------------------------------------------------------------------------------------------

\begin{lemma}
    Les groupes \( \gU(n)\) et \( \SU(n)\) sont connexes par arcs.
\end{lemma}

\begin{proof}
    Soit \( A\), une matrice unitaire et \( Q\) une matrice unitaire qui diagonalise \( A\). Étant donné que les valeurs propres arrivent par paires complexes conjuguées,
    \begin{equation}
        QAQ^{-1}=\begin{pmatrix}
            e^{i\theta_1}    &       &       &       &   \\  
            &    e^{-i\theta_1}    &       &       &   \\  
            &       &    \ddots    &       &   \\  
            &       &       &    e^{i\theta_r}    &   \\  
            &       &       &       &        e^{-i\theta_r}
        \end{pmatrix}.
    \end{equation}
    Le chemin \( U(t)\) obtenu en remplaçant \( \theta_i\) par \( t\theta_i\) avec \( t\in\mathopen[ 0 , 1 \mathclose]\) joint \( QAQ^{-1}\) à l'identité. Par conséquent \( Q^{-1}U(t)Q\) joint \( A\) à l'unité.
\end{proof}

\begin{theorem}
    Les matrices \wikipedia{fr}{Endomorphisme_normal}{normales} forment un espace connexe par arc.
\end{theorem}

\begin{proof}
    Soit \( A\) une matrice normale, et \( U\) une matrice unitaire qui diagonalise \( A\). Nous considérons \( U(t)\), un chemin qui joint \( \mtu\) à \( U\) dans \( \gU(n)\). Pour chaque \( t\), la matrice
    \begin{equation}
        A(t)=U(t)^{-1} AU(t)
    \end{equation}
    est normale. Nous avons donc trouvé un chemin dans les matrices normales qui joint \( A\) à une matrice diagonale. Il est à présent facile de la joindre à l'identité.

    Toutes les matrices normales étant connexes à l'identité, l'ensemble des matrices normales est connexe.
\end{proof}

%---------------------------------------------------------------------------------------------------------------------------
\subsection{Densité}
%---------------------------------------------------------------------------------------------------------------------------

\begin{proposition}     \label{PropDigDensVxzPuo}
    Les matrices diagonalisables sont denses dans \( \eM(n,\eC)\).
\end{proposition}

\begin{proof}
    D'après le lemme de Schur \ref{LemSchurComplHAftTq}, une matrice de \( \eM(n,\eC)\) est de la forme
    \begin{equation}
        A=Q\begin{pmatrix}
            \lambda_1    &   *    &   *    \\
              0  &   \ddots    &   *    \\
            0    &   0    &   \lambda_n
        \end{pmatrix}Q^{-1}.
    \end{equation}
    Les valeurs propres sont sur la diagonale. La matrice est diagonalisable si les éléments de la diagonales sont tous différents. Il suffit maintenant de considérer \( n\) suites \( (\epsilon^{(r)}_k)_{k\in\eN}\) convergentes vers zéro telles que pour chaque \( k\) les nombres \( \lambda_r+\epsilon^{(r)}_k\) soient tous différents. La suite de matrices
    \begin{equation}
        A_k=Q\begin{pmatrix}
            \lambda_1+\epsilon^{(1)}_k    &   *    &   *    \\
              0  &   \ddots    &   *    \\
              0    &   0    &   \lambda_n+\epsilon^{(n)}_k
        \end{pmatrix}Q^{-1}.
    \end{equation}
    est alors diagonalisable pour tout \( k\) et nous avons \( \lim_{k\to \infty} A_k=A\).
\end{proof}

\begin{proposition}
    Si \( A\in\eM(n,\eC)\) alors
    \begin{equation}
        e^{\tr(A)}=\det( e^{A}).
    \end{equation}
\end{proposition}

\begin{proof}
    Le résultat est un simple calcul pour les matrices diagonalisable. Si \( A\) n'est pas diagonalisable, nous considérons une suite de matrices diagonalisables \( A_k\) dont la limite est \( A\) (proposition \ref{PropDigDensVxzPuo}). La suite
    \begin{equation}
        a_k= e^{\tr(A_k)}
    \end{equation}
    converge vers \(  e^{\tr(A)}\) tandis que la suite 
    \begin{equation}
        b_k=\det( e^{A_k})
    \end{equation}
    converge vers \( \det( e^{A})\). Mais nous avons \( a_k=b_k\) pour tout \( k\); les limites sont donc égales.
\end{proof}

%---------------------------------------------------------------------------------------------------------------------------
\subsection{Racine carré d'une matrice hermitienne positive}
%---------------------------------------------------------------------------------------------------------------------------

\begin{proposition}     \label{PropVZvCWn}
    Si \( A\in \eM(n,\eC)\) est une matrice hermitienne positive, alors il existe une unique matrice hermitienne positive \( R\) telle que \( A=R^2\). De plus \( R\) est un polynôme (de \( \eR[X]\)) en \( A\).
\end{proposition}
\index{matrice!semblable}
\index{polynôme!d'endomorphisme}
\index{endomorphisme!diagonalisable}
\index{matrice!hermitienne!racine carré}
\index{racine!carré!de matrice hermitienne}

La matrice \( R\) ainsi définie est la \defe{racine carré de}{matrice!racine carré}\index{racine!carré de matrice!hermitienne positive} de \( A\), et est notée \( \sqrt{A}\)\nomenclature[A]{\( \sqrt{A}\)}{racine d'une matrice hermitienne positive}. Une des applications usuelles de cette proposition est la décomposition polaire.

\begin{proof}
    \begin{subproof}
    \item[Existence]
        Étant donné que \( A \) est hermitienne, elle est diagonalisable par une unitaire (proposition \ref{ThogammwA}), et ses valeurs propres sont réelles et positives (parce que \( A\) est positive). Soit donc \( P\) une matrice unitaire telle que
        \begin{equation}
            P^*AP=\begin{pmatrix}
                \alpha_1    &       &       \\
                    &   \ddots    &       \\
                    &       &   \alpha_n
            \end{pmatrix}
        \end{equation}
        avec \( \alpha_i>0\). Si on pose
        \begin{equation}
            R=P\begin{pmatrix}
                \sqrt{\alpha_1}    &       &       \\
                    &   \ddots    &       \\
                    &       &   \sqrt{\alpha_n}
            \end{pmatrix}P^*,
        \end{equation}
        alors \( R^2=A\) parce que \( P^*P=\mtu\).
    \item[Hermitienne positive]
        La matrice \( R\) est hermitienne parce que, avec un peu de notation raccourcie, \( R=P^*\sqrt{\alpha}P\) et \( R^*=P^*\sqrt{\alpha}P\). D'autre part, elle est positive parce que ses valeurs propres sont les \( \sqrt{\alpha_i}\) qui sont positives.
        
    \item[Polynôme]
        Nous montrons maintenant que la matrice \( R\) est un polynôme en \( A\). Pour cela nous considérons un polynôme \( Q\) tel que \( A(\alpha_i)=\sqrt{\alpha_i}\) pour tout \( i\). Soit \( \{ e_i \}\) une base de diagonalisation de \( A\) : \( Ae_i=\alpha_ie_i\). Alors c'est encore une base de diagonalisation de \( Q(A)\). En effet si \( Q=\sum_ka_kX^k\), alors
        \begin{equation}
            Q(A)e_i=(\sum_ka_kA^k)e_i=(\sum_ka_k\alpha_i^k)e_i=Q(\alpha_i)e_i=\sqrt{\alpha_i}e_i.
        \end{equation}
        Les valeurs propres de \( Q(A)\) sont donc \( \sqrt{\alpha_i}\). Nous savons maintenant que \( Q(A)\) a la même base de diagonalisation de \( A\) (et donc la même matrice unitaire \( P\) qui diagonalise), c'est à dire que
        \begin{equation}
            Q(A)=P^*\begin{pmatrix}
                \sqrt{\alpha_1}    &       &       \\
                    &   \ddots    &       \\
                    &       &   \sqrt{\alpha_n}
            \end{pmatrix}=R.
        \end{equation}
        Donc oui, \( R\) est un polynôme en \( A\).

        Notons que ce \( Q\) n'est pas du tout unique; il existe une infinité de polynômes qui envoient \( n\) nombres donnés sur \( n\) nombres donnés.

    \item[Unicité]
        Soit \( S\) une matrice hermitienne positive telle que \( R^2=S^2=A\). D'abord \( S\) commute avec \( A\) parce que
        \begin{equation}
            SA=S^3=S^2S=AS.
        \end{equation}
        Donc \( S\) commute aussi avec \( Q(A)=R\). Étant donné que \( S\) et \( R\) commutent et sont diagonalisables, ils sont simultanément diagonalisables par le corollaire \ref{CorQeVqsS}. Soient \( D_R=PRP^*\) et \( D_S=PSP^*\) les formes diagonales de \( R\) et \( S\) dans une base de simultanée diagonalisation. Les carrés des valeurs propres de \( R\) et \( S\) étant identiques (ce sont les valeurs propres de \( A\)) et les valeurs propres de \( R\) et \( S\) étant positives, nous déduisons que \( D_R=D_S\) et donc que \( R=P^*D_RP=P^*D_SP=S\).
    \end{subproof}
\end{proof}

%---------------------------------------------------------------------------------------------------------------------------
\subsection{Racine carré d'une matrice symétrique positive}
%---------------------------------------------------------------------------------------------------------------------------

\begin{lemma}[\cite{JJdQPyK}]   \label{LemTLlTAAf}
    Le groupe orthogonal \( O(n,\eR) \) est compact.
\end{lemma}

\begin{proof}
    Nous avons \( O(n)=f^{-1}\big( \{ \mtu_n \} \big)\) où \( f\) est l'application continue \( A\mapsto A^tA\). En tant qu'image inverse d'un fermé par une application continue, le groupe \( O(n)\) est fermé.

    De plus il est borné parce que tous les coefficients d'une matrice orthogonale sont \( \leq 1\), donc \( \| A \|_{\infty}\) pour tout \( A\in O(n)\).
\end{proof}

\begin{proposition} \label{PropPEMDqVT}
    Une matrice symétrique semi (ou pas) définie positive admet une unique racine carré symétrique. Le spectre de la racine carré est la racine carré du spectre de la matrice de départ.
\end{proposition}

\begin{proof}
    Ceci est une phrase pour que les titres se mettent bien.
    \begin{subproof}
        \item[Existence]
            Soit \( T\) une matrice symétrique et \( Q\) une matrice orthogonale qui diagonalise\footnote{Théorème \ref{ThoeTMXla}.} \( T\) : \( QTQ^{-1}=D\) avec \( D=\diag(\lambda_i)\) et \( \lambda_i\geq 0\). En posant \( R=Q^{-1}\sqrt{D}Q\), il est vite vérifié que \( R^2=T\) et que \( R\) est symétrique. En ce qui concerne le spectre, \( R\) a pour valeurs propres les \( \sqrt{\lambda_i}\).
        \item[Unicité]

            Soit \( R\) une matrice symétrique de \( T\) : \( R^2=T\). Du coup \( R\) et \( T\) commutent : \( RT=R^3=TR\). Par conséquent les espaces propres de \( T\) sont stables sous \( R\). Soit \( E_{\lambda} \) l'un d'eux de dimension \( d\), et \( T_F\), \( R_F\) les restrictions de \( T\) et \( R\) à \( E_{\lambda}\). L'application \( T_F\) est une homothétie et \( R_F^2=T_F=\lambda\mtu\). Mais \( R_F\) est encore une matrice symétrique définie positive, donc nous pouvons considérer une base \( \{ e_1,\ldots, e_d \}\) de \( E_{\lambda}\) qui diagonalise \( R_F\) avec les valeurs propres \( \mu_i\); nous avons donc en même temps
            \begin{subequations}
                \begin{align}
                    R_f^2(e_i)&=\mu_i^2 e_i\\
                    T_F(e_i)&=\lambda e_i,
                \end{align}
            \end{subequations}
            de telle sorte que \( \mu_i^2=\lambda\). Mais les valeurs propres de \( R_F\) sont positives, sont \( \mu_i=\sqrt{\lambda}\) pour tout \( i\). En conclusion \( R_F\) est univoquement déterminé par la donnée de \( T\). Vu que cela est valable pour tous les espaces propres de \( T\) et que ces espaces propres engendrent tout \( E\), l'opérateur \( R\) est déterminé de façon univoque par \( T\).
    \end{subproof}
\end{proof}
Notons que nous n'avons démontré l'unicité qu'au sein des matrices symétriques.

%--------------------------------------------------------------------------------------------------------------------------- 
\subsection{Décomposition polaires : cas réel}
%---------------------------------------------------------------------------------------------------------------------------

Nous nommons \( S^+(n,\eR)\) l'ensemble des matrices \( n\times n\) symétriques réelles définies positives et \( S^{++}(n,\eR)\) le sous-ensemble de \( S^+(n,\eR)\) des matrices strictement définies positives.
\nomenclature[R]{\( S^+(n,\eR)\)}{matrices symétriques définies positives}
\nomenclature[R]{\( S^{++}(n,\eR)\)}{matrices symétriques strictement définies positives}

\begin{lemma}   \label{LemZKJWqIP}
    La fermeture de l'ensemble des matrice symétriques strictement définies positives est l'ensemble des matrices définies positives : \( \overline{ S^{++}(n,\eR) }=S^+(n,\eR)\).
\end{lemma}

\begin{proof}
    Nous savons que \( S^+(n,\eR)\) est fermé et que \( S^{++}(n,\eR)\) y est inclus, donc il suffit de montrer que toute suite de \( S^{++}(n,\eR)\) convergente dans \( \eM(n,\eR)\) a une limite dans \( S^+(n,\eR)\).

    En effet si \( S_k\) est une suite de matrices symétriques convergeant dans \( \eM(n,\eR)\) vers la matrice \( A\), les suites \( (S_k)_{ij}\) et \( (S_k)_{ji}\) des composantes \( ij\) et \( ji\) sont des suites égales, et donc leurs limites sont égales\footnote{Ici nous utilisons le critère de convergence composante par composante et le fait que nous ne sommes pas trop inquiétés par la norme que nous choisissons parce que toutes les normes sont équivalentes par le théorème \ref{ThoNormesEquiv}.}. Donc la limite est symétrique.

    En ce qui concerne le spectre, nous diagonalisons : \( S_k=Q_kD_kQ_k^{-1}\) où les \( D_k\) sont des matrices diagonales remplies de nombres strictement positifs. Vu que \( O(n)\) est compact\footnote{Lemme \ref{LemTLlTAAf}.}, nous avons une sous-suite \( Q_{\varphi(k)}\) convergente : \( Q_{\varphi(k)}\to Q\). Pour chaque \( k\), nous avons
    \begin{equation}
        S_{\varphi(k)}=Q_{\varphi(k)}D_{\varphi(k)}Q^{-1}_{\varphi(k)},
    \end{equation}
    dont la limite existe et vaut \( A\). Étant donné que le membre de droite est le produit de trois facteurs dont deux sont convergents et dont le tout converge, le troisième facteur est également convergent : \( S_{\varphi(k)}\to S\). En passant à la limite :
    \begin{equation}
        A=\lim_{k\to \infty } S_{\varphi(k)}=QDQ^{-1},
    \end{equation}
    et donc le spectre de \( A\) est la limite de ceux des matrices \( D_{\varphi(k)}\). Chacun étant strictement positif, la limite est positive (mais pas spécialement strictement). Donc \( A\in S^+(n,\eR)\).
\end{proof}

\begin{theorem}[Décomposition polaire de matrices symétriques définies positives\cite{JJdQPyK,AABkVai}] \label{ThoLHebUAU}
   En ce qui concerne les matrices inversibles :
   \begin{equation}
       \begin{aligned}
           f\colon O(n,\eR)\times S^{++}(n,\eR)&\to \GL(n,\eR) \\
           (Q,S)&\mapsto SQ 
       \end{aligned}
   \end{equation}
   est un difféomorphisme.

   En ce qui concerne les matrices en général :
   \begin{equation}
       \begin{aligned}
           g\colon O(n,\eR)\times S^+(n,\eR)&\to \eM(n,\eR) \\
           (Q,S)&\mapsto SQ 
       \end{aligned}
   \end{equation}
   est une surjection mais pas une injection.

   De plus les mêmes conclusions tiennent si nous regardons \( (Q,S)\mapsto QS\) au lieu de \( SQ\).
\end{theorem}
\index{groupe!linéaire!décomposition polaire}
\index{endomorphisme!décomposition!polaire}
\index{décomposition!polaire}

%TODO : prouver le difféomorphisme.
%TODO : je crois qu'on doit pouvoir prouver que les éléments de la décomposition polaire sont des polynômes en M.

\begin{proof}
    La preuve qui suit ne démontre qu'un homéomorphisme dans le cas des matrices inversibles.
    \begin{subproof}
        \item[Existence et unicité]

            Si \( M=SQ\), alors \( MM^t=SQQ^tS^t=S^2\), donc \( S\) doit être une racine carré symétrique de la matrice définie positive \( MM^t\). La proposition \ref{PropPEMDqVT} nous dit que ça existe et que c'est unique. Donc \( S\) est univoquement déterminé par \( M\). Maintenant avoir \( Q=MS^{-1}\) est obligatoire (unicité) et fonctionne :
            \begin{equation}
                Q^tQ=(S^{-1})^tM^tMS^{-1}=S^{-1}S^2S^{-1}=\mtu,
            \end{equation}
            donc \( Q\) ainsi défini est orthogonale.

            Ceci prouve l'existence et l'unicité de la décomposition polaire d'une matrice inversible.
        
        \item[Homéomorphisme]

            Le fait que \( f\) soit continue n'est pas un problème : c'est un produit de matrice. Nous devons vérifier que \( f^{-1}\) est continue. Soit une suite convergente \( M_k\to M\) dans \( \GL(n,\eR)\). Si nous nommons \( (Q_k,S_k)\) la décomposition polaire de \( M_k\) et \( (Q,S)\) celle de \( M\), nous devons prouver que \( Q_k\to Q\) et \( S_k\to S\). En effet dans ce cas nous aurions
            \begin{equation}    \label{EqJIkoaJv}
                \lim_{k\to \infty} f^{-1}(M_k)=\lim_{k\to \infty} (Q_k,S_k)=(Q,S)=f^{-1}(M).
            \end{equation}
            
            Étant donné que \( O(n)\) est compact (lemme \ref{LemTLlTAAf}), la suite \( (Q_k)\) admet une sous-suite convergente (Bolzano-Weierstrass, théorème \ref{ThoBWFTXAZNH}) que nous nommons
            \begin{equation}
                Q_{\varphi(k)}\to F\in O(n).
            \end{equation}
            Vu que la suite \( (M_k)\) converge, sa sous-suite converge vers la même limite : \( M_{\varphi(k)}\to M\) et vu que pour tout \( k\) nous avons \( S_k=M_kQ_k^{-1}\),
            \begin{equation}
                S_{\varphi(k)}\to G=MF^{-1}.
            \end{equation}
            Vu que chacune des matrices \( S_{\varphi(k)}\) est symétrique définie positive, la limite est symétrique et semi-définie positive\footnote{Lemme \ref{LemZKJWqIP}}. Donc \( G\in S^+(n,\eR)\cap \GL(n,\eR)\) parce que de plus \( M\) et \( F\) étant inversibles, \( G\) est inversible. En ce qui concerne la sous-suite nous avons
            \begin{equation}
                M_{\varphi(k)}=S_{\varphi(k)}Q_{\varphi(k)}\to GF=M
            \end{equation}
            où \( F\in O(n)\) et \( G\in S^+(n,\eR)\). Par unicité de la décomposition polaire de \( M\) (partie déjà démontrée), nous avons \( G=S\) et \( F=Q\).

            Nous avons prouvé que toute sous-suite convergente de \( Q_k\) a \( Q\) pour limite. Donc la suite elle-même converge\footnote{Proposition \ref{PropHNylIAW}, pas difficile.} vers \( Q\). Donc \( Q_k\to Q\). Du coup vu que \( S_k=M_kQ_k^{-1}\) est un produit de suites convergentes, \( S_k\) converge également, vers \( S\) :  \( S_k\to S\).

            Au final l'application \( f^{-1}\) est bien continue parce que les égalités \eqref{EqJIkoaJv} ont bien lieu.

    \end{subproof}
\end{proof}

%---------------------------------------------------------------------------------------------------------------------------
\subsection{Enveloppe convexe du groupe orthogonal}
%---------------------------------------------------------------------------------------------------------------------------

Le théorème suivant n'est pas indispensablissime parce qu'il est le même que le théorème de la projection sur les espaces de Hilbert\footnote{Théorème \ref{ThoProjOrthuzcYkz}}. Cependant la partie existence est plus simple en se limitant au cas de dimension finie.
\begin{theorem}[Théorème de la projection]  \label{ThoWKwosrH}
    Soit \( E\) un espace vectoriel réel ou complexe de dimension finie, \( x\in E\), et \( C\) un sous ensemble fermé convexe de \(E\).
    \begin{enumerate}
        \item
            Les deux conditions suivantes sur \( y\in E\) sont équivalentes:
    \begin{enumerate}
        \item   \label{zzETsfYCSItemi}
            \( \| x-y \|=\inf\{ \| x-z \|\tq z\in C \}\),
        \item\label{zzETsfYCSItemii}
            pour tout \( z\in C\), \( \Reel\langle x-y, z-y\rangle \leq 0\).
    \end{enumerate}
\item
    Il existe un unique \( y\in E\), noté \( y=\pr_C(x)\) vérifiant ces conditions.
    \end{enumerate}
\end{theorem}
%TODO : il y a surement un endroit plus adapté pour mettre ce théorème.

\begin{proof}
    Nous commençons par prouver l'existence et l'unicité d'un élément dans \( C\) vérifiant la première condition. Ensuite nous verrons l'équivalence. 

    \begin{subproof}
        \item[Existence]
        
            Soit \( z_0\in C\) et \( r=\| x-z_0 \|\). La boule fermée \( \overline{ B(x,r) }\) est compacte\footnote{C'est ceci qui ne marche plus en dimension infinie.} et intersecte \( C\). Vu que \( C\) est fermé, l'ensemble \( C'=C\cap\overline{ B(x,r) }\) est compacte. Tous les points qui minimisent la distance entre \( x\) et \( C\) sont dans \( C'\); la fonction 
            \begin{equation}
                \begin{aligned}
                     C'&\to \eR \\
                    z&\mapsto d(x,z) 
                \end{aligned}
            \end{equation}
            est continue sur un compact et donc a un minimum qu'elle atteint. Un point \( P\) réalisant ce minimum prouve l'existence d'un point vérifiant la première condition.

        \item[Unicité]
            Soient \( y_1\) et \( y_2\), deux éléments de \( C\) minimisant la distance avec \( x\), et soit \( d\) ce minimum. Nous avons par l'identité du parallélogramme \eqref{EqYCLtWfJ} que
            \begin{equation}
                \| y_1-y_2 \|^2=-4\left\| \frac{ y_1+y_2-x }{2} \right\|^2+2\| y_1-x \|^2+2\| y_2-x \|^2\leq -4d+2d+2d=0.
            \end{equation}
            Par conséquent \( y_1=y_2\).

        \item[\ref{zzETsfYCSItemi}\( \Rightarrow\) \ref{zzETsfYCSItemii}]

            Soit \( z\in C\) et \( t\in \mathopen] 0 , 1 \mathclose[\); nous notons \( P=\pr_Cx\). Par convexité le point \( z=ty+(1-t)P\) est dans \( C\), et par conséquent,
                \begin{equation}
                    \| x-P \|^2\leq\| x-tz-(1-t)P \|^2=\| (x-P)-t(z-P) \|^2.
                \end{equation}
                Nous sommes dans un cas \( \| a \|^2\leq | a-b |^2\), qui implique \( 2\Reel\langle a, b\rangle \leq \| b \|^2\). Dans notre cas,
                \begin{equation}
                    2\Reel\langle x-P , t(z-P)\rangle \leq t^2\| z-P \|^2.
                \end{equation}
                En divisant par \( t\) et en faisant \( t\to 0\) nous trouvons l'inégalité demandée :
                \begin{equation}
                    2\Reel\langle x-P, z-P\rangle \leq 0.
                \end{equation}
                
        \item[\ref{zzETsfYCSItemii}\( \Rightarrow\) \ref{zzETsfYCSItemi}]

            Soit un point \( P\in C\) vérifiant 
            \begin{equation}
                \Reel\langle x-P, z-P\rangle \leq 0
            \end{equation}
            pour tout \( z\in C\). Alors en notant \( a=x-P\) et \( b=P-z\),
            \begin{equation}
                \begin{aligned}[]
                \| x-z \|^2=\| x-P+P-z \|^2&=\| a+b \|^2\\
                &=\| a \|^2+\| b \|^2+2\Reel\langle a, b\rangle \\
                &=\| a \|^2+\| b \|^2-2\Reel\langle x-P, z-P\rangle \\
                &\geq \| b \|^2,
                \end{aligned}
            \end{equation}
            ce qu'il fallait.

    \end{subproof}

\end{proof}

\begin{definition}
    Soit \( A\) une partie d'un espace vectoriel \( E\). L'\defe{enveloppe convexe}{enveloppe!convexe} de \( A\), notée \( \Conv(A)\)\nomenclature[G]{\( \Conv(A)\)}{enveloppe convexe} est l'intersection de tous les convexes contenant \( A\).
\end{definition}
L'enveloppe convexe est un convexe. En effet soit \( C\) un convexe contenant \( A\) et \( x,y\in\Conv(A)\); alors \( x\) et \( y \) sont dans \( C\) et par conséquent le segment \( [x,y]\) est inclus à \( C\). Ce segment étant inclus à tout convexe contenant \( A\), il est inclus à \( \Conv(A)\).

\begin{theorem}[Carathéodory\cite{KXjFWKA}] \label{ThoJLDjXLe}
    Dans un espace affine de dimension \( n\), l'enveloppe convexe de \( A\) est l'ensemble des barycentres à coefficients positifs ou nuls de familles de \( n+1\) points.
\end{theorem}
\index{barycentre!enveloppe convexe}
\index{espace!vectoriel!dimension}

\begin{proof}
    Soit \( x\in\Conv(A)\); on sait par la proposition \ref{PropYHMTmZX} que \( x\) est barycentre de points de \( A\) avec des coefficients positifs :
    \begin{equation}    \label{EqWJDwOTH}
        x=\sum_{k=1}^p\lambda_kx_k
    \end{equation}
    avec \( \sum_k\lambda_k=1\). Nous supposons que \( p>n+1\) (sinon le théorème est réglé), et nous allons faire une récurrence à l'envers en montrant qu'on peut aussi écrire \( x\) sous forme d'un barycentre de strictement moins de \( p\) points.
    
    Étant donné que \( p-1>n\), la famille \( \{ x+i-x_1 \}_{i=2,\ldots, p}\) est liée et il existe donc \( \alpha_1,\ldots, \alpha_p\in \eR\) tels que \( \sum_{i=2}^p\alpha_i(x_i-x_1)=0\), c'est à dire telle que
    \begin{equation}
        \sum_{i=2}^p\alpha_ix_i=\sum_{i=2}^p\alpha_ix_1.
    \end{equation}
    Nous posons \( \alpha_1=-\sum_{i=2}^p\alpha_1\). Remarquons qu'alors \( \sum_{i=1}^p\alpha_ix_i=0\) parce que
    \begin{equation}
        \sum_{i=1}^p\alpha_ix_i=\alpha_1x_1+\sum_{i=2}^p\alpha_ix_i=\alpha_1x_1+\sum_{i=2}^p\alpha_ix_1=\sum_{i=1}^p\alpha_ix_1=0.
    \end{equation}
    Par conséquent ça ne coûte rien de récrire \eqref{EqWJDwOTH} sous la forme
    \begin{equation}
        x=\sum_{i=1}^p(\lambda_i+t\alpha_i)x_i.
    \end{equation}
    Les \( \alpha_i\) ne sont pas tous nuls, mais leur somme est nulle, donc il y en a au moins un négatif. Nous notons
    \begin{equation}
        \tau=\min\{ -\frac{ \lambda_i }{ \alpha_i }\tq \alpha_i<0 \},
    \end{equation}
    et \( J\) l'ensemble de \( i\) pour lesquels ce minimum est atteint. Nous considérons aussi le nombres \( \mu_i=\lambda_i+\tau\alpha_i\). Plusieurs remarques. 
    \begin{enumerate}
        \item
            Si \( j\in J\), alors \( \mu_j=0\)
        \item
            Si \( \alpha_i>0\) alors \( \mu_i\geq 0\), mais si \( \alpha_i<0\) alors
            \begin{equation}
                \lambda_i+\tau\alpha_i\geq \lambda_i+(-\frac{ \lambda_i }{ \alpha_i })\alpha_i=0m
            \end{equation}
            donc \( \mu_i\geq 0\) quand même.
        \item
            \( \sum_{i=1}^p\mu_i=1\), toujours parce que \( \sum_{i=1}^p\alpha_i=0\).
    \end{enumerate}
    Avec tout ça, nous avons
    \begin{equation}
        \sum_{i\notin J}\mu_ix_i=\sum_{i=1}^p\mu_ix_i=x.
    \end{equation}
    Et voila, nous avons écrit \( x\) comme un barycentre à coefficients positifs de moins de \( p\) éléments parce que \( J\) n'est pas vide.
\end{proof}

\begin{corollary}   \label{CorOFrXzIf}
    Dans un espace affine de dimension finie, l'enveloppe convexe d'un compact est compacte.
\end{corollary}

\begin{proof}
    Soit \( A\) une partie compacte de l'espace vectoriel \( E\), et \( \Conv(A)\) son enveloppe convexe. Nous allons montrer que toute suite dans \( \Conv(A)\) admet une sous-suite convergente en écrivant un point de \( \Conv(A)\) comme le théorème de Carathéodory \ref{ThoJLDjXLe} nous le suggère. Pour cela nous considérons le simplexe
    \begin{equation}
        \Lambda=\left\{  \lambda\in \eR^{n+1}\tq \sum_{k=1}^{n+1}\lambda_k=1\text{ et } \lambda_k\geq 0\forall k   \right\}.    
    \end{equation}
    Montrons en passant que \( \Lambda\) est compact. Si \( \lambda_k\in \Lambda\) est une suite, alors chacun des \( \lambda_k\) est un \( (n+1)\)-uple de nombres dans \( \mathopen[ 0 , 1 \mathclose]\) :
    \begin{equation}
        k\mapsto (\lambda_k)_i
    \end{equation}
    est une suite qui possède une sous-suite convergente. En passant \( n+1\) fois à une sous-suite, nous tombons sur une suite convergente vers \( \lambda\in\Lambda\), grâce à la convergence composante par composante. De plus pour chaque \( k\) nous avons \( \sum_{i=1}^{n+1}(\lambda_k)_i=1\), et en passant à la limite, la somme étant une application continue, \( \sum_{i}\lambda_i=1\).

    Considérons l'application
    \begin{equation}
        \begin{aligned}
            f\colon \Lambda\times A^{n+1}&\to \Conv(A) \\
            (\lambda,x)&\mapsto \sum_{k=1}^{n+1}\lambda_kx_k. 
        \end{aligned}
    \end{equation}
    C'est une application continue parce qu'elle est bilinéaire en dimension finie; son image est contenue dans \( \Conv(A)\) par la proposition \ref{PropSVvAQzi}, et elle est surjective par le théorème de Carathéodory. Bref, \( \Conv(A)=f(\Lambda\times A^{n+1})\) est donc l'image d'un compact par une application continue; elle est donc compacte par le théorème \ref{ThoImCompCotComp}.
\end{proof}
Notons que sans le théorème de Carathéodory, peut être que le nombre de points utiles pour décomposer les différents \( a_k\) n'était pas borné; dans ce cas nous aurions du prendre une infinité de sous-suites et rien n'aurait été sûr.

\begin{definition}
    Sur \( C\) est un ensemble convexe, un point \( x\in C\) est un \defe{point extrémal}{extrémal!point dans un convexe} si \( C\setminus\{ x \}\) est encore convexe.
\end{definition}

\begin{theorem}[\cite{KXjFWKA}] \label{ThoBALmoQw}
    Soit \( E\) un espace euclidien de dimension \( n\geq 1\) et \( \aL(E)\) l'espace des opérateurs linéaires sur \( E\) sur lequel nous considérons la norme subordonnée\footnote{Voir la définition donnée dans l'exemple \ref{ExemdefnormpMrt}.} à celle sur \( E\). L'ensemble des points extrémaux de la boule unité fermée de \( \aL(E)\) est le groupe orthogonal \( O(n,\eR)\).
\end{theorem}
\index{densité!points extrémaux dans \( \aL\)}

\begin{proof}
    Nous notons \( \mB\) la boule unité fermée de \( \aL(E)\). Montrons pour commencer que les éléments de \( O(n)\) sont extrémaux dans \( \mB\). D'abord si \( A\in O(E)\) alors \( \| A \|=1\) parce que \( \| Ax \|=\| x \|\). Supposons maintenant que \( A\) n'est pas extrémal, c'est à dire qu'il est le milieu d'un segment joignant deux points (distincts) de la boule unité de \( \aL(E)\). Soient donc \( T,U\in\mB\) tels que \( A=\frac{ 1 }{2}(T+U)\). Pour tout \( x\in E\) tel que \( \| x \|=1\) nous avons 
    \begin{equation}    \label{EqKTuAIIE}
        1=\| x \|=\| Ax \|=\frac{ 1 }{2}\| Tx+Ux \|\leq \frac{ 1 }{2}\big( \| Tx \|+\| Ux \| \big)\leq\frac{ 1 }{2}\big( \| T \|+| U | \big)\leq 1
    \end{equation}
    Toutes les inégalités sont en réalité des égalités. En particulier nous avons
    \begin{equation}
        \| Tx+Ux \|=\| Tx \|+\| Ux \|,
    \end{equation}
    mais alors nous sommes dans un cas d'égalité dans l'inégalité de Cauchy-Schwartz (théorème \ref{ThoAYfEHG}) et donc il existe \( \lambda\geq 0\) tel que \( Tx=\lambda Ux\). Mais de plus les \sout{inégalité} égalités \eqref{EqKTuAIIE} nous donnent
    \begin{equation}
        \frac{ 1 }{2}\big( \| Tx \|+\| Ux \| \big)=1
    \end{equation}
    alors que nous savons que \( \| Tx \|,\| Ux \|\leq 1\), donc \( \| Tx \|=\| Ux \|=1\). La seule possibilité est d'avoir \( \lambda=1\) et donc que \( U=T\) parce que nous avons \( Tx=Ux\) pour tout \( x\) de norme \( 1\). Au final \( A\) n'est pas le milieu d'un segment dans \( \mB\).

    Nous passons donc à l'inclusion inverse : nous prouvons que les points extrémaux de \( \mB\) sont dans \( O(E)\). Pour cela nous prenons \( U\in\mB\setminus O(E)\) et nous allons montrer que \( U\) n'est pas un point extrémal : nous allons l'écrire comme milieu d'un segment dans \( \mB\).

    Par la seconde partie du théorème de décomposition polaire \ref{ThoLHebUAU}, il existe \( Q\in O(n,\eR)\) et \( S\in S^+(n,\eR)\) tels que \( U=QS\). Nous diagonalisons \( S\) à l'aide de la matrice orthogonale \( P\) :
    \begin{equation}
        S=PDP^{-1}
    \end{equation}
    avec \( D=\diag(\lambda_i)\). En termes de normes, nous avons
    \begin{equation}
        \| U \|=\| S \|=\| S \|.
    \end{equation}
    En effet vu que \( Q\) est orthogonale, \( \| Ux \|=\| QSx \|=\| Sx \|\) pour tout \( x\), donc \( \| U \|=\| S \|\). De plus pour tout \( x\) nous avons
    \begin{equation}
        \| Sx \|=\| PDP^{-1} x \|=\| DP^{-1}x \|.
    \end{equation}
    Étant donné que \( P^{-1}\) est une bijection, le supremum des \( \| Sx \|\) sera le même que celui des \( \| Dx \|\) et donc \( \| S \|=\| D \|\). Étant donné que par définition \( \| U \|\leq 1\), nous avons aussi \( \| D \|\leq 1\) et donc \( 0\leq\lambda_i\leq 1\) (pour rappel, les valeurs propres de \( D\) sont positives ou nulles parce que \( S\) est ainsi). 

    Comme \( U\notin O(E)\), au moins une des valeurs propres n'est pas \( 1\), supposons que ce soit \( \lambda_1\). Alors nous avons \( \alpha,\beta\in\mathopen[ -1 , 1 \mathclose]\) avec \( -1\leq \alpha<\beta\leq 1\) et \( \lambda_1=\frac{ 1 }{2}(\alpha+\beta)\). Nous posons alors
    \begin{subequations}
        \begin{align}
            D_1=\diag(\alpha,\lambda_2,\ldots, \lambda_n)\\
            D_2=\diag(\beta,\lambda_2,\ldots, \lambda_n).
        \end{align}
    \end{subequations}
    Nous avons bien \( D_1\neq D2\) et \( D_1+ D_2=D\). Par conséquent
    \begin{equation}
        U=\frac{ 1 }{2}\big( QPD_1P^{-1}+QPD_2P^{-1} \big)
    \end{equation}
    avec \( QPD_1P^{-1}\neq QPD_2^{-1}\). La matrice \( U \) est donc le milieu d'un segment. Reste à montrer que ce segment est dans \( \mB\). Pour ce faire, prenons \( x\in E\) et calculons :
    \begin{equation}
        \| QPD_iP^{-1}x \|=\| D_iP^{-1}x \|\leq\| P^{-1}x \|=\| x \|
    \end{equation}
    parce que \( \| D_i \|\leq 1\) et \( P^{-1}\) est orthogonale. Au final la norme de \( QPD_iP\) est plus petite que \( 1\) et donc \( U\) est bien le milieu d'un segment dans \( \mB\), et donc non extrémal.
\end{proof}

\begin{theorem}[\cite{NHXUsTa}] \label{ThoVBzqUpy}
    L'enveloppe convexe de \( O(n)\) dans \( \eM_n(\eR)\) est la boule unité pour la norme induite de \( \| . \|_2\) sur \( \eR^n\).
\end{theorem}
\index{convexité!enveloppe de $O(n)$}
\index{groupe!linéaire!enveloppe convexe de $\Omega(n)$}

\begin{proof}
    Nous notons \( \mB\) la boule unité fermée de \( \eM(n,\eR)\) et \( \Conv\big( O(n,\eR) \big)\) l'enveloppe convexe de \( O(n,\eR)\). Vu que \( \mB\) est convexe nous avons \( \Conv\big( O(n) \big)\subset\mB\).


    Maintenant nous devons prouver l'inclusion inverse. Pour ce faire nous supposons avoir un élément \( A\in \mB\setminus\Conv\big( O(n) \big)\) et nous allons dériver une contradiction.
    
    Remarquons que \( O(n)\) est compact par le lemme \ref{LemTLlTAAf} et que par conséquent \( \Conv(O(n))\) est compacte par le corollaire \ref{CorOFrXzIf} et donc fermée. Nous considérons un produit scalaire \( (X,Y)\mapsto X\cdot Y\) sur \( \eM\). Vu que \( \Conv\big( O(n) \big)\) est un fermé convexe nous pouvons considérer la projection\footnote{Le théorème de projection : théorème \ref{ThoWKwosrH}.} sur \( \Conv(A)\) relativement au produit scalaire choisis.

    Nous notons \( P=\pr_{\Conv\big( O(n) \big)}(A)\). En vertu du théorème de projection, nous avons
    \begin{equation}    \label{EqYSisLTL}
        (A-P)\cdot (M-P)\leq 0
    \end{equation}
    pour tout \( M\in\Conv O(n)\). Notons \( B=A-P\) pour alléger les notations. L'équation \eqref{EqYSisLTL} s'écrit
    \begin{equation}    \label{EqQDLZqXQ}
        B\cdot M\leq B\cdot P.
    \end{equation}
    D'autre par vu que \( B \neq 0\) nous avons \( B\cdot B> 0\), c'est à dire \( B\cdot (A-P)>0\) et donc
    \begin{equation}
        B\cdot A>B\cdot P.
    \end{equation}
    En combinant avec \eqref{EqQDLZqXQ},
    \begin{equation}        \label{EqIQNlwql}
        B\cdot M\leq B\cdot P<B\cdot A.
    \end{equation}
    Nous utilisons maintenant la décomposition polaire, théorème \ref{ThoLHebUAU}, pour écrire \( B=QS\) avec \( Q\in O(n)\) et \( S\in S^+(n,\eR)\). Vu que l'inégalité \eqref{EqIQNlwql} tient pour tout \( M\in\Conv(O(n))\), elle tient en particulier pour \( Q\in O(n)\). Donc
    \begin{equation}
        B\cdot Q=B\cdot A.
    \end{equation}
    Nous nous particularisons à présent au produit scalaire \( (X,Y)\mapsto\tr(X^tY)\) de la proposition \ref{PropMAQoKAg}. D'abord
    \begin{equation}    \label{EaHVxWdau}
        B\cdot Q=\tr(B^tQ)=\tr(S^tQ^tQ)=\tr(S^t)=\tr(S),
    \end{equation}
    et ensuite l'inégalité \eqref{EaHVxWdau} devient
    \begin{equation}
        \tr(S)<B\cdot A=\tr(S^tQ^tA).
    \end{equation}
    Nous choisissons une basse \( \{ e_i \}\) diagonalisant \( S\) : \( Se_i=\lambda_ie_i\) vérifiant automatiquement \( \lambda_i>0\) parce que \( S\) est semi-définie positive. Alors
    \begin{subequations}
        \begin{align}
            \tr(S)&<\tr(S^tQ^tA)\\
            &=\sum_i\langle S^tQ^tAe_i, e_i\rangle \\
            &=\sum_i\langle Ae_i, QSe_i\rangle \\
            &\leq \sum_i \| Ae_i \| | \lambda_i | \underbrace{\| Qe_i \|}_{=1} \\
            &\leq \sum_i\lambda_i   & A\in\mB\Rightarrow\| Ae_i \|\leq 1\\
            &=\tr(S).
        \end{align}
    \end{subequations}
    Il faut noter que la première inégalité est stricte, et donc nous avons une contradiction.
\end{proof}

%+++++++++++++++++++++++++++++++++++++++++++++++++++++++++++++++++++++++++++++++++++++++++++++++++++++++++++++++++++++++++++
\section{Sous espaces caractéristiques}
%+++++++++++++++++++++++++++++++++++++++++++++++++++++++++++++++++++++++++++++++++++++++++++++++++++++++++++++++++++++++++++

% TODO : lire le blog de Pierre Bernard; en particulier celle-ci : http://allken-bernard.org/pierre/weblog/?p=2299

Sources : \cite{MneimneReduct} et \wikipedia{fr}{Décomposition_de_Dunford}{divers articles sur Wikipédia}.
%TODO : citer mieux Wikipédia.

Lorsqu'un opérateur n'est pas diagonalisable, les valeurs propres jouent quand même un rôle important.

Soit \( E\) un \( \eK\)-espace vectoriel et \( f\in\End(E)\). Pour \( \lambda\in \eK\) nous définissons
\begin{equation}
    F_{\lambda}(f)=\{ v\in E\tq (f-\lambda\mtu)^nv=0, n\in\eN \}.
\end{equation}
C'est l'ensemble de nilpotence de l'opérateur \( f-\lambda\mtu\).

\begin{lemma}
    L'ensemble \( F_{\lambda}(f)\) est non vide si et seulement si \( \lambda\) est une valeur propre de \( f\). L'espace \( F_{\lambda}(f)\) est invariant sous \( f\).
\end{lemma}

\begin{proof}
    Si \( F_{\lambda}(f)\) est non vide, nous considérons \( v\in F_{\lambda}(f)\) et \( n\) le plus petit entier non nul tel que \( (f-\lambda)^nv=0\). Alors \( (f-\lambda)^{n-1}v\) est un vecteur propre de \( f\) pour la valeur propre \( \lambda\). Inversement si \( v\) est une valeur propre de \( f\) pour la valeur propre \( \lambda\), alors \( v\in F_{\lambda}(f)\).

    En ce qui concerne l'invariance, remarquons que \( f\) commute avec \( f-\lambda\mtu\). Si \( x\in F_{\lambda}(f)\) il existe \( n\) tel que \( (f-\lambda\mtu)^nx=0\). Nous avons aussi
    \begin{equation}
        (f-\lambda\mtu)^nf(x)=f\big( (f-\lambda\mtu)^nx \big)=0,
    \end{equation}
    par conséquent \( f(x)\in F_{\lambda}(f)\).
\end{proof}

\begin{remark}
    Toute matrice sur \( \eC\) n'est pas diagonalisable. Considérons en effet l'endomorphisme \( f\) donné par la matrice
    \begin{equation}
        \begin{pmatrix}
            a&    \alpha    &   \beta    \\
            0    &   a    &   \gamma    \\
            0    &   0    &   b
        \end{pmatrix}
    \end{equation}
    où \( a\neq b\), \( \alpha\neq 0\), \( \beta\) et \( \gamma\) sont des nombres complexes quelconques.
    Son polynôme caractéristique est 
    \begin{equation}
        \chi_f(\lambda)=(a-\lambda)^2(b-\lambda)
    \end{equation}
    de telle façon à ce que les valeurs propres soient \( a\) et \( b\). Nous trouvons les vecteurs propres pour la valeur \( a\) en résolvant
    \begin{equation}
        \begin{pmatrix}
            a    &   \alpha    &   \beta    \\
            0    &   a    &   \gamma    \\
            0    &   0    &   b
        \end{pmatrix}\begin{pmatrix}
            x    \\ 
            y    \\ 
            z    
        \end{pmatrix}=\begin{pmatrix}
            ax    \\ 
            ay    \\ 
            az    
        \end{pmatrix}.
    \end{equation}
    L'espace propre \( E_a(f)\) est réduit à une seule dimension générée par \( (1,0,0)\). De la même façon l'espace propre correspondant à la valeur propre \( b\) est donné par 
    \begin{equation}
        \begin{pmatrix}
            \frac{1}{ b-a }\left( \beta+\frac{ \alpha\gamma }{ b-a } \right)    \\ 
            \frac{ \gamma }{ b-a }    \\ 
            1    
        \end{pmatrix}.
    \end{equation}
    Il n'y a donc pas trois vecteurs propres linéairement indépendants, et l'opérateur \( f\) n'est pas diagonalisable.

    Par contre nous pouvons voir que
    \begin{equation}
        (f-\alpha\mtu)^2\begin{pmatrix}
             0   \\ 
            1    \\ 
            0    
        \end{pmatrix}=
        \begin{pmatrix}
            a    &   \alpha    &   \beta    \\
            0    &   a    &   \gamma    \\
            0    &   0    &   b
        \end{pmatrix}
        \begin{pmatrix}
            \alpha    \\ 
            0    \\ 
            0    
        \end{pmatrix}-\begin{pmatrix}
            a\alpha    \\ 
            0    \\ 
            0    
        \end{pmatrix}=\begin{pmatrix}
            0    \\ 
            0    \\ 
            0    
        \end{pmatrix},
    \end{equation}
    de telle sorte que le vecteur \( (0,1,0)\) soit également dans l'espace caractéristique \( F_a(f)\).

    Dans cet exemple, la multiplicité algébrique de la racine \( a\) du polynôme caractéristique vaut \( 2\) tandis que sa multiplicité géométrique vaut seulement \( 1\).
\end{remark}

Le théorème suivant est aussi appelé le théorème de \defe{décomposition primaire}{décomposition!primaire}.


\begin{theorem}[Théorème spectral, décomposition primaire]\index{théorème!spectral}     \label{ThoSpectraluRMLok}
    Soit \( E\) espace vectoriel de dimension finie sur le corps algébriquement clos \( \eK\) et \( f\in\End(E)\). Alors
    \begin{equation}
        E=F_{\lambda_1}(f)\oplus\ldots\oplus F_{\lambda_k}(f)
    \end{equation}
    où la somme est sur les valeurs propres distinctes de \( f\).

    Les projecteurs sur les espaces caractéristique forment un système complet et orthogonal.
\end{theorem}

\begin{proof}
    Soit \( P\) le polynôme caractéristique de \( u\) et une décomposition
    \begin{equation}
        P=(u-\lambda_1)^{\alpha_1}\ldots(u-\lambda_r)^{\alpha_r}
    \end{equation}
    en facteurs irréductibles. La le théorème de noyaux (\ref{ThoDecompNoyayzzMWod}) nous avons
    \begin{equation}        \label{EqDeFVSaYv}
        E=\ker(u-\lambda_1)^{\alpha_1}\oplus\ldots\oplus\ker(u-\lambda_r)^{\alpha_r}.
    \end{equation}
    Les projecteurs sont des polynômes en \( u\) et forment un système orthogonal. Il nous reste à prouver que \( \ker(u-\lambda_i)^{\alpha_i}=F_{\lambda_i}(u)\). L'inclusion
    \begin{equation}    \label{EqzmNxPi}
        \ker(u-\lambda_i)^{\alpha_i}\subset F_{\lambda_i}(u)
    \end{equation}
    est évidente. Nous devons montrer l'inclusion inverse. Prouvons que la somme des \( F_{\lambda_i}(u)\) est directe. Si \( v\in F_{\lambda_i}(u)\cap F_{\lambda_j}(u)\), alors il existe \( v_1=(u-\lambda_i)^nv\neq 0\) avec \( (u-\lambda_i)v_1=0\). Étant donné que \( (u-\lambda_i)\) commute avec \( (u-\lambda_j)\), ce \( v_1\) est encore dans \( F_{\lambda_j}(u)\) et par conséquent il existe \( w=(u-\lambda_j)^mv_1\) non nul tel que 
    \begin{subequations}
        \begin{numcases}{}
            (u-\lambda_i)w=0\\
            (u-\lambda_j)w=0.
        \end{numcases}
    \end{subequations}
    Ce \( w\) serait donc un vecteur propre simultané pour les valeurs propres \( \lambda_i\) et \( \lambda_j\), ce qui est impossible parce que les espaces propres sont linéairement indépendants. Les espaces \( F_{\lambda_i}\) sont donc en somme directe et
    \begin{equation}
        \sum_i\dim F_{\lambda_i}(u)\leq \dim E.
    \end{equation}
    En tenant compte de l'inclusion \eqref{EqzmNxPi} nous avons même
    \begin{equation}
        \dim E=\sum_i\dim\ker(u-\lambda_i)^{\alpha_i}\leq\sum_i F_{\lambda_i}(u)\leq \dim E.
    \end{equation}
    Par conséquent nous avons \( \dim\ker(u-\lambda_i)^{\alpha_i}=\dim F_{\lambda_i}(u)\) et l'égalité des deux espaces.
    
\end{proof}

Le théorème suivant généralise le théorème de diagonalisabilité \ref{ThoDigLEQEXR} au cas où le polynôme minimum est seulement scindé.

\begin{probleme}
    \begin{enumerate}
\item 
    Dans le cas où le corps n'est pas algébriquement clos, il paraît qu'il faut remplacer «diagonalisable» par «semi-simple».
    \end{enumerate}
\end{probleme}

\begin{definition}
    Un endomorphisme d'un espace vectoriel est \defe{semi-simple}{semi-simple!endomorphisme} si tout sous-espace stable par \( u\) possède un supplémentaire stable.
\end{definition}
Si l'espace vectoriel est sur un corps algébriquement clos, alors les endomorphismes semi-simples sont les endomorphismes diagonaux.


\begin{theorem}[Décomposition de Dunford]\index{décomposition!Dunford}\index{Dunford!décomposition} \label{ThoRURcpW}
    Soit \( E\) un espace vectoriel sur le corps algébriquement clos \( \eK\) et \( u\in\End(E)\) un endomorphisme de \( E\). Alors \( u\) se décompose de façon unique sous la forme
    \begin{equation}
        u=s+n
    \end{equation}
    où \( s\) est diagonalisable, \( n\) est nilpotent et \( [s,n]=0\).

    De plus \( s\) et \( n\) sont des polynômes en \( u\) et commutent avec \( u\).
\end{theorem}

\begin{proof}
    Le théorème spectral \ref{ThoSpectraluRMLok} nous indique que
    \begin{equation}
        E=\bigoplus_iF_{\lambda_i}(f).
    \end{equation}
    Nous considérons l'endomorphisme \( s\) de \( E\) qui consiste à dilater d'un facteur \( \lambda\) l'espace caractéristique \( F_{\lambda}(f)\) :
    \begin{equation}
        s=\sum_i\lambda_ip_i
    \end{equation}
    où \( p_i\colon E\to F_{\lambda_i}(u)\) est la projection de \( E\) sur \( F_{\lambda_i}(u)\).

    Nous allons prouver que \( [s,f]=0\) et \( n=f-s\) est nilpotent. Cela impliquera que \( [s,n]=0\).

    Si \( x\in F_{\lambda}(f)\), alors nous avons \( sf(x)=\lambda f(x)\) parce que \( f(x)\in F_{\lambda}(f)\) tandis que \( fs(x)=f(\lambda x)=\lambda f(x)\). Par conséquent \( f\) commute avec \( s\).

    Pour montrer que \( f-s\) est nilpotent, nous en considérons la restriction
    \begin{equation}
        f-s\colon F_{\lambda}(f)\to F_{\lambda}(f).
    \end{equation}
    Cet opérateur est égal à \( f-\lambda\mtu\) et est par conséquent nilpotent.

    Prouvons à présent l'unicité. Soit \( u=s'+n'\) une autre décomposition qui satisfait aux conditions : \( s'\) est diagonalisable, \( n'\) est nilpotent et \( [n',s']=0\). Commençons par prouver que \( s'\) et \( n'\) commutent avec \( u\). En multipliant \( u=s'+n'\) par \( s'\) nous avons
    \begin{equation}
        s'u=s'^2+s'n'=s'^2+n's'=(s'+n')s'=us',
    \end{equation}
    par conséquent \( [u,s']=0\). Nous faisons la même chose avec \( n'\) pour trouver \( [u,n']=0\). Notons que pour obtenir ce résultat nous avons utilisé le fait que \( n'\) et \( s'\) commutent, mais pas leur propriétés de nilpotence et de diagonalisibilité.
    
    
    Si \( s'+n'=s+n\) est une autre décomposition, \( s'\) et \( n'\) commutent avec \( u\), et par conséquent avec tous les polynômes en \( u\). Ils commutent en particulier avec \( n\) et \( s\). Les endomorphismes \( s\) et \( s'\) sont alors deux endomorphismes diagonalisables qui commutent. Par la proposition \ref{PropGqhAMei}, ils sont simultanément diagonalisables. Dans la base de simultanée diagonalisation, la matrice de l'opérateur \( s'-s=n-n'\) est donc diagonale. Mais \( n-n'\) est également nilpotent, en effet si \( A\) et \( B\) sont deux opérateurs nilpotents,
    \begin{equation}
        (A+B)^n=\sum_{k=0}^n\binom{k}{n}A^kB^{n-k}.
    \end{equation}
    Si \( n\) est assez grand, au moins un parmi \( A^k\) ou \( B^{n-k}\) est nul.

    Maintenant que \( n-n'\) est diagonal et nilpotent, il est nul et \( n=n'\). Nous avons alors immédiatement aussi \( s=s'\).
\end{proof}


%---------------------------------------------------------------------------------------------------------------------------
\subsection{Calcul de l'exponentielle d'une matrice}
%---------------------------------------------------------------------------------------------------------------------------

Nous reprenons l'exemple de \cite{MneimneReduct}. Soit \( A\) une matrice dont le polynôme minimum s'écrit
\begin{equation}
    P(X)=(X-1)^2(X-2).
\end{equation}
Par le théorème \ref{ThoDecompNoyayzzMWod} de décomposition des noyaux nous avons
\begin{equation}
    E=\ker(A-1)^2\oplus\ker(A-2).
\end{equation}
En suivant les notations de ce théorème nous avons \( P_1(X)=(X-1)^2\), \( P_2(X)=X-2\) et
\begin{subequations}
    \begin{align}
        Q_1(X)&=X-2\\
        Q_2(X)&=(X-1)^2.
    \end{align}
\end{subequations}
Les polynômes \( R_i\) dont l'existence est assurée par le théorème de Bézout sont
\begin{equation}
    \begin{aligned}[]
        R_1(X)&=-X\\
        R_2(X)&=1.
    \end{aligned}
\end{equation}
Nous avons
\begin{equation}
    R_1Q_1+R_2Q_2=1.
\end{equation}
Le projecteur \( p_i\) sur \( \ker P_i\) est \( R_iQ_i\) :
\begin{equation}
    \begin{aligned}[]
        p_1&=-A(A-2)=\pr_{\ker(u-1)^2}\\
        p_2&=(A-1)^2=\pr_{\ker(u-2)}.
    \end{aligned}
\end{equation}
Passons maintenant au calcul de l'exponentielle. Nous avons évidemment
\begin{equation}
    e^A=e^Ap_1+e^Ap_2.
\end{equation}
Étant donné que \( p_1\) est le projecteur sur le noyau de \( (A-1)^2\), nous avons
\begin{equation}
    e^Ap_1=ee^{A-1}p_1=ep_1+e(u-1)1=ep_1=-Ae(A-2).
\end{equation}
En effet \( e^{A-1}p_1=\sum_{k=0}^{\infty}(A-1)^k\circ p_1\). De la même façon nous avons
\begin{equation}
    e^Ap_2=e^2e^{A-2}p_2=e^2p_2=e^2(A-1)^2.
\end{equation}
Au final,
\begin{equation}
    e^A=-Ae(A-2)+e^2(A-1)^2.
\end{equation}

\begin{theorem}
    Soit une matrice \( A\in \eM(n,\eC)\). On a que la suite \( (A^kx)\) tends vers zéro pour tout \( x\) si et seulement si \( \rho(A)<1\) où \( \rho(A)\)\index{rayon!spectral} est le rayon spectral de $A$
\end{theorem}

\begin{proof}
    Dans le sens direct, il suffit de prendre comme \( x\), un vecteur propre de \( A\). Dans ce cas nous avons \( A^kx=\lambda^kx\). Mais \( \lambda^kx\) ne tend vers zéro que si \( \lambda<1\). Donc toute les valeurs propres de \( A\) doivent être plus petite que \( 1\) et \( \rho(A)<1\).

    Pour l'autre sens nous utilisons la décomposition de Dunford (théorème \ref{ThoRURcpW}) : il existe une matrice inversible \( P\) telle que
    \begin{equation}
        A=P^{-1}(D+N)P
    \end{equation}
    où \( D\) est diagonale, \( N\) est nilpotente et \( [D,N]=0\). Étant donné que \( D+N\) est triangulaire, son polynôme caractéristique que
    \begin{equation}
        \chi_{D+N}(\lambda)=\prod_i D_{ii}-\lambda.
    \end{equation}
    Par similitude, c'est le même polynôme caractéristique que celui de \( A\) et nous savons alors que la diagonale de \( D\) contient les valeurs propres de \( A\).

    Par ailleurs nous avons
    \begin{subequations}
        \begin{align}
            A^k&=P^{-1}(D+N)^kP\\
            &=P^{-1}\sum_{j=0}^k{j\choose k}D^{j-k}N^jP\\
            &==P^{-1}\sum_{j=0}^{n-1}{j\choose k}D^{j-k}N^jP
        \end{align}
    \end{subequations}
    où nous avons utilité le fait que \( D\) et \( N\) commutent ainsi que \( N^{n-1}=0\) parce que \( N\) est nilpotente. Nous utilisons la norme matricielle usuelle, pour laquelle \( \| D \|=\rho(D)=\rho(A)\). Nous avons alors
    \begin{equation}
        \| (D+N)^k \|\leq \sum_{j=0}^k{j\choose k}\rho(D)^{k-j}\| N \|^j.
    \end{equation}
    Du coup si \( \rho(D)<1\) alors \( \| (D+N)^k \|\to 0\) (et c'est même un si et seulement si).
\end{proof}


%---------------------------------------------------------------------------------------------------------------------------
\subsection{Valeurs singulières}
%---------------------------------------------------------------------------------------------------------------------------

\begin{definition}
    Soit \( M\) une matrice \( m\times n\) sur \( \eK\) (\( \eK\) est \( \eR\) ou \( \eC\)). Un nombre réel \( \sigma\) est une \defe{valeur singulière}{valeur!singulière} de \( M\) si il existent des vecteurs unitaires \( u\in \eK^m\), \( v\in \eK^n\) tels que
    \begin{subequations}
        \begin{align}
            Mv&=\sigma u\\
            M^*u&=\sigma v.
        \end{align}
    \end{subequations}
\end{definition}

\begin{theorem}[Décomposition en valeurs singulières]
    Soit \( M\in \eM(m\times n,\eK)\) où \( \eK=\eR,\eC\). Alors \( M\) se décompose en
    \begin{equation}
        M=ADB
    \end{equation}
    où
    il existe deux matrices unitaires \( A\in \gU(m\times m)\), \( B\in \gU(n\times n)\) et une matrice (pseudo)diagonale \( D\in \eM(m\times n)\) tels que
    \begin{enumerate}
        \item 
            \( A\in\gU(m\times m)\), \( B\in\gU(n\times n)\) sont deux matrices unitaires;,
        \item
            \( D\) est (pseudo)diagonale,
        \item
            les éléments diagonaux de \( D\) sont les valeurs singulières de \( M\),
        \item
            le nombre d'éléments non nuls sur la diagonale de \( D\) est le rang de \( M\).
    \end{enumerate}
\end{theorem}

\begin{corollary}
    Soit \( M\in \eM(n,\eC)\). Il existe un isomorphisme \( f\colon \eC^n\to \eC^n\) tel que \( fM\) soit autoadjoint.
\end{corollary}

\begin{proof}
    Si \( M=ADB\) est la décomposition de \( M\) en valeurs singulières, alors nous pouvons prendre \( f=\overline{ B }^tA^{-1}\) qui est une matrice inversible. Pour la vérification que ce \( f\) répond bien à la question, ne pas oublier que \( D\) est réelle, même si \( M\) ne l'est pas.
\end{proof}

%+++++++++++++++++++++++++++++++++++++++++++++++++++++++++++++++++++++++++++++++++++++++++++++++++++++++++++++++++++++++++++
\section{Matrice compagnon et endomorphismes cycliques}
%+++++++++++++++++++++++++++++++++++++++++++++++++++++++++++++++++++++++++++++++++++++++++++++++++++++++++++++++++++++++++++

%---------------------------------------------------------------------------------------------------------------------------
\subsection{Matrice compagnon}
%---------------------------------------------------------------------------------------------------------------------------

Soit le polynôme \( P=X^n-a_{n-1}X^{n-1}-\ldots-a_1X-a_0\) dans \( \eK[X]\). La \defe{matrice compagnon}{matrice!compagnon} de \( P\) est la matrice\nomenclature[A]{\( C(P)\)}{matrice compagnon} donnée par
\begin{equation}
    C(P)=\begin{pmatrix}
        0    &   \cdots    &   \cdots    &   0    &   a_0\\  
        1    &   0    &       &   \vdots    &   a_1\\  
        0    &   \ddots    &   \ddots    &   \vdots    &   \vdots\\  
        \vdots    &   \ddots    &   \ddots    &   0    &   a_{n-2}\\  
        0    &   \cdots    &   0    &   1    &   a_{n-1}    
    \end{pmatrix}
\end{equation}
si \( n\geq 2\) et par \( (a_0)\) si \( n=1\). Si \( f\) est l'endomorphisme associé à la matrice \( C(P)\) nous avons
\begin{equation}
    f(e_i)=\begin{cases}
        e_{i+1}    &   \text{si \( i<n\)}\\
        (a_0,\ldots, a_{n-1})    &    \text{si \( i=n\)}.
    \end{cases}
\end{equation}
Cet endomorphisme est conçu pour vérifier \( P(f)e_1=0\).

\begin{lemma}[\cite{RapportArgreg2011}] \label{LemkVNisk}
    Soit \( P\), un polynôme sur un corps commutatif \( \eK\). Si \( f\) est l'endomorphisme associé à la matrice compagnon de \( P\), alors \( P\) est la polynôme caractéristique de \( f\). En d'autres termes, \( P=\chi_f\).
\end{lemma}

\begin{proof}
    La propriété \( P(f)e_1=0\) nous indique que le polynôme minimal ponctuel de \( f\) en \( e_1\) divise \( P\). L'ensemble des puissances de \( f\) appliquées à \( e_1\), \( \big( f^i(e_1) \big)_{i=1,\ldots, n-1}\) est libre, donc le polynôme minimal ponctuel en \( e_1\) est de degré \( n\) au minimum. En reprenant les notations du théorème \ref{ThoCCHkoU}, nous avons \( I_{e_1}=(P)\) parce que \( P\) est de degré minimum dans \( I_{e_1}\) et \( \chi_f\in I_{e_1}\).

    Donc \( P\) divise \( \chi_f\) et est de degré égal à celui de \( \chi_f\). Étant donné qu'ils sont tous deux unitaires, ils sont égaux.
\end{proof}

\begin{remark}  \label{RemmQjZOA}
    Les matrices compagnons ne sont pas les seules dont le polynôme caractéristique est égal au polynôme minimal. En fait les matrices sont le polynôme caractéristique est égale au polynôme minimal sont denses dans les matrices. En effet une matrice dont le polynôme minimal n'est pas égal au polynôme caractéristique a un polynôme caractéristique avec une racine double. Il est possible, en modifiant arbitrairement peu la matrice de séparer la racine double en deux racines distinctes.
\end{remark}

\begin{definition}[Matrices, endomorphismes et vecteurs cycliques]
    Une matrice est \defe{cyclique}{cyclique!matrice}\index{matrice!cyclique} si elle est semblable à une matrice compagnon. Un endomorphisme \( f\colon E\to E\) est \defe{cyclique}{cyclique!endomorphisme}\index{endomorphisme!cyclique} si il existe un vecteur \( x\in E\) tel que \( \{ f^k(x)\tq k=1,\ldots, n-1 \}\) est une base de \( E\). Un vecteur ayant cette propriété est un \defe{vecteur cyclique}{vecteur!cyclique} pour \( f\).
\end{definition}

\begin{lemma}
    Un endomorphisme est cyclique si et seulement si sa matrice associée est cyclique.
\end{lemma}

\begin{lemma}   \label{LemSGmdnE}
    Une matrice est cyclique si et seulement si ses polynômes minimal et caractéristiques coïncident.
\end{lemma}

\begin{lemma}   \label{LemAGZNNa}
    Si \( f\colon E\to E\) est un endomorphisme cyclique et si \( y\) est un vecteur cyclique de \( f\), alors le polynôme minimal de \( f\) est égal au polynôme minimal de \( f\) au point \( y\) : \( \mu_{f}=\mu_{f,y}\).
\end{lemma}

\begin{proof}
    Montrons que \( \mu_{f,y}\) est un polynôme annulateur de \( f\), ce qui prouvera que \( \mu(f)\) divise \( \mu_{f,y}\). Étant donné que \( y\) est cyclique, tout élément de \( E\) s'écrit sous la forme \( x=Q(f)y\). Prenons un polynôme \( P\) annulateur de \( f\) en \( y\) : \( P(f)y=0\). Nous montrons que \( P\) est alors un polynôme annulateur de \( f\). En effet, nous avons
    \begin{equation}
        P(f)x=\big( P(f)\circ Q(f) \big)y=\big( Q(f)\circ P(f) \big)y=0
    \end{equation}
    où nous avons utilisé le lemme \ref{LemQWvhYb}.
\end{proof}

%---------------------------------------------------------------------------------------------------------------------------
\subsection{Réduction de Frobenius}
%---------------------------------------------------------------------------------------------------------------------------

\begin{theorem}[Réduction de Frobenius \cite{AutourFrobCompa,Vialivs}]      \index{réduction!Frobénius}\index{Frobénius!réduction}
    Soit \( E\), un \( \eK\)-espace vectoriel où \( \eK\) est \( \eR\) ou \( \eC\), et \( f\in \End(E)\). Alors il existe une suite de sous-espaces \( E_1,\ldots, E_r\) stables par \( f\) tels que
    \begin{enumerate}
        \item   \label{ItemmpwjnSs}
            \( E=\bigoplus_{i=1}^rE_i\);
        \item
            pour chaque \( E_i\), l'endomorphisme restreint \( f_i=f|_{E_i}\) est cyclique;
        \item
            si \( \mu_i\) est le polynôme minimal de \( f_i\) alors \( \mu_{i+1}\) divise \( \mu_i\);
    \end{enumerate}
    Une telle décomposition vérifie automatiquement \( \mu_1=\mu_f\) et \( \mu_1\cdots \mu_r=\chi_f\), et la suite \( (\mu_i)_{i=1,\ldots, r}\) ne dépend que de \( f\) et non du choix de la décomposition du point \ref{ItemmpwjnSs}.
\end{theorem}

Les polynômes \( \mu_i\) sont les \defe{invariants de similitude}{invariant!de similitude} de l'endomorphisme \( f\).

\begin{proof}
    Nous commençons par montrer que si une telle décomposition existe, alors
    \begin{subequations}    \label{subEqzcGouz}
        \begin{align}
            \chi_f=\prod_{i=1}^r\mu_i  \label{EqTaxsvb}\\
            \mu_f=\mu_1
        \end{align}
    \end{subequations}
    où \( \chi_f\) est le polynôme caractéristique de \( f\) et \( \mu_f\) est le polynôme minimal\footnote{Cette partie de la preuve provient de \cite{MoncetIVS}.}. D'abord le polynôme caractéristique de \( f\) devra être égal au produit des polynômes caractéristique des \( f|_{E_i}\), mais ces derniers endomorphismes étant cycliques, leurs polynôme caractéristiques sont égaux à leurs polynômes minimaux (lemme \ref{LemSGmdnE}). Cela prouve l'égalité \eqref{EqTaxsvb}. Ensuite tous les \( \mu_i\) doivent diviser le polynôme minimal, donc \( \ppcm(\mu_1,\ldots, \mu_r)\) divise \(\mu_f\). Cependant le polynôme minimal ne doit contenir une et une seule fois chacun des facteurs irréductibles du polynôme caractéristique, et chacun de ces facteurs sont dans les polynômes \( \mu_i\). Par conséquent \( \ppcm(\mu_1,\ldots, \mu_r)=\mu_f\). Mais par ailleurs \( \mu_1=\ppcm(\mu_1,\ldots, \mu_r)\), donc \( \mu_1=\mu_f\).
    
    Mais le produit des \( \mu_i\) est le polynôme caractéristique, donc tous les facteurs irréductibles du polynôme minimal sont dans les \( \mu_i\); cela signifie que \( \mu_f=\ppcm(\mu_1,\ldots, \mu_r)\).

    Soit \( d\), le degré du polynôme minimal de \( f\) et \( y\in E\) tel que \( \mu_f=\mu_{f,y}\) (voir lemme \ref{LemSYsJJj}). Le plus petit espace stable sous \( f\) contenant \( y\) est
    \begin{equation}
        E_y=\Span\{ y,f(y),\ldots, f^{d-1}(y) \}.
    \end{equation}
    Nous notons \( e_i=f^{i-1}(y)\). Notons que les vecteurs donnés forment bien une base de \( E_y\) parce que si les \( e_i\) n'était pas linéairement indépendants, alors soit \( \sum_ka_ke_k=0\). Donc ce cas nous aurions
    \begin{equation}
        \big( \sum_ka_kX^k \big)(f)y=0,
    \end{equation}
    ce qui contredirait la minimalité de \( \mu_{f,y}\).

    La difficulté du théorème est de trouver un complément de \( E_y\) qui soit également stable sous \( f\). Nous commençons par étendre\footnote{Pour autant que j'aie compris, cette extension manque dans \cite{AutourFrobCompa}. Corrigez moi si je me trompe.} \( \{ e_1,\ldots, e_d \}\) en une base \( \{ e_1,\ldots, e_n \}\) de \( E\). Ensuite nous allons montrer que
    \begin{equation}
        E=E_y\oplus F
    \end{equation}
    avec
    \begin{equation}
        F=\{ x\in E\tq  e^*_d\big( f^k(x) \big)=0\forall k\in \eN \}.
    \end{equation}
    Par construction, \( F\) est invariant sous \( f\). Montrons pour commencer que \( E_y\cap F=\{ 0 \}\). Un élément de \( E_y\) s'écrit
    \begin{equation}
        z=a_1e_1+\ldots +a_ke_k
    \end{equation}
    avec \( k\leq d\). Étant donné que \( f\) décale les vecteurs de base, nous avons \( e^*_d\big( f^{d-k}(z) \big)=a_k\). Du coup \( z\in F\) si et seulement si \( a_1=\ldots=a_d=0\), c'est à dire que \( E_y\cap F=\{ 0 \}\).

    Nous montrons maintenant que \( \dim F=n-d\). Pour cela nous considérons l'application
    \begin{equation}
        \begin{aligned}
            T\colon \eK[F]&\to E^* \\
            g&\mapsto e^*_d\circ g. 
        \end{aligned}
    \end{equation}
    Cette application est injective. En effet un élément général de \( \eK[f]\) est
    \begin{equation}
        g=a_1\id+a_2f+\ldots +a_pf^{p-1}
    \end{equation}
    avec \( p\leq d\). Si \( T(g)=0\), alors nous avons en particulier
    \begin{equation}
        0=T(g)e_{_d-p+1}=e^*_d(a_1e_{d-p+1}+a_2e_{d-p+2}+\ldots +a_pe_d)=a_p.
    \end{equation}
    Donc \( a_p=0\) et en appliquant maintenant \( T(g)\) à \( e_{d-p}\) nous obtenons \( a_{p-1}=0\). Au final nous trouvons que \( g=0\) et donc que \( T\) est injective.

    Étant donné que \( \dim\eK[f]=d\) et que \( T\) est injective, \( \dim\Image(T)=d\). Nous regardons l'orthogonal de l'image :
    \begin{subequations}
        \begin{align}
            (\Image(T))^{\perp}&=\{ x\in E\tq T(g)x=0\forall g\in\eK[f] \}\\
            &=\{ x\in E\tq e^*_d\big( g(x) \big)=0\forall g\in \eK[f] \}\\
            &=F.
        \end{align}
    \end{subequations}
    Par conséquent \( F^{\perp}=\Image(T)\). Vu que \( \dim\Image(T)=d\), nous avons donc \( \dim F=n-d\) et il est établi que \( E=E_y\oplus F\). 

    Nous avons donc trouvé \( F\), stable par \( f\) et tel que \( E=E_y\oplus F\). Nous devons maintenant nous assurer que cette décomposition tombe bien pour les polynômes minimaux. Si \( P_1\) est le polynôme minimal de \( f|_{E_yj}\), alors par le lemme \ref{LemAGZNNa} nous avons \( P_1=\mu_{f,y}=\mu_f\) parce que \( f|_{E_y}\) est cyclique sur \( E_y\). Mettons \( P_2\), le polynôme minimal de \( f|_F\). Étant attendu que \( F\) est stable par \( f\), le polynôme \( P_2\) divise \( P_1\). En recommençant la construction sur \( F\), nous construisons un nouvel espace \( F'\) stable sous \( F\) et vérifiant \( \mu_{f|_{F'}}=P_2\), etc.

    Nous passons maintenant à la partie unicité du théorème. Soient deux suites \( F_1,\ldots, F_r\) et \( G_1,\ldots, G_s\) de sous-espaces stables par \( f\) et vérifiant
    \begin{enumerate}
        \item
            \( E=\bigoplus_{i=1}^rF_i\),
        \item
            \( f|_{F_i}\) est cyclique,
        \item
            \( \mu_{f|_{F_{i+1}}}\) divise \( \mu_{f|_{F_i}}\),
    \end{enumerate}
    et, \emph{mutatis mutandis}, les mêmes conditions pour la famille \( \{ G_i \}\). Nous posons \( P_i=\mu_{f_{F_i}}\) et \( Q_i=\mu_{f|_{G_i}}\). Nous allons montrer par récurrence que \( P_i=Q_i\) et \( \dim F_i=\dim G_i\). Il ne sera cependant pas garanti que \( F_i=G_i\). D'abord, \( P_1=Q_1\) parce qu'ils sont tous deux égaux à \( \mu_f\) par les relations \eqref{subEqzcGouz}. Nous supposons que \( P_i=Q_i\) pour \( i\leq 1\leq j-1\) et nous tentons de montrer que \( P_j=Q_j\).

    Nous avons 
    \begin{equation}    \label{EqMrCtZO}
        P_j(f)=P_j(f)|_{F_1}\oplus\ldots\oplus P_j(f)|_{F_{j-1}}.
    \end{equation}
    En effet étant donné que \( P_{j+k}\) divise \( P_j\), nous avons\footnote{En vertu du lemme \ref{LemQWvhYb}.} \( P_{j}(f)=A(f)\circ P_{j+k}(f)\), mais \( P_{j+k}(f)F_{j+k}=0\), donc \( P_j(f)F_{j+k}=0\). Les espaces \( G_i\) n'ayant a priori aucun rapport avec les polynômes \( P_i\), nous écrivons
    \begin{equation}    \label{EqJreLiO}
        P_j(f)=P_j(f)|_{G_1}\oplus\ldots\oplus P_j(f)|_{G_{j-1}}\oplus P_j(f)|_{G_j}\oplus\ldots\oplus P_j(f)|_{G_s}.
    \end{equation}
    Pour \( 1\leq i\leq j-1\), nous avons supposé \( P_i=Q_i\). Étant donné que \( f|_{F_i}\) est semblable à \( C_{_i}\) et \( f|_{G_i}\) est semblable à \( C_{Q_i}\), la matrice de \( f|_{E_i}\) est semblable à la matrice de \( f|_{G_i}\). En particulier,
    \begin{equation}
        \dim P_j(f)F_i=\dim P_j(f)G_i.
    \end{equation}
    En prenant les dimensions des images dans les égalités \eqref{EqMrCtZO} et \eqref{EqJreLiO}, nous trouvons que
    \begin{equation}
        P_j(f)|_{G_j}=\ldots=P_j(f)|_{G_s}=0.
    \end{equation}
    Par conséquent \( P_j\in I_{f|G_j}\) et donc \( P_j\) divise \( Q_j\), qui est générateur de \( I_{f|_{G_j}}\). La situation étant symétrique entre \( P\) et \( Q\), nous montrons de même que \( Q_j\) divise \( P_j\) et donc que \( P_j=Q_j\).

    Ceci achève la démonstration du théorème de réduction de Frobenius.

\end{proof}


Sous forme matricielle, ce théorème dit que toute matrice est semblable à une matrice de la forme bloc-diagonale
\begin{equation}
    f=\begin{pmatrix}
        C_{\mu_1}    &       &       \\
            &   \ddots    &       \\
            &       &   C_{\mu_r}
    \end{pmatrix}
\end{equation}

\begin{remark}
    Si nous travaillons sur \( \eR\), la réduite de Frobenius restera une matrice réelle, même si les valeurs propres sont complexes. En effet le procédé de Frobenius ne regarde absolument pas les valeurs propres, mais seulement les facteurs irréductibles du polynôme caractéristique. La réduite de Frobenius ne tente pas de résoudre ces polynômes, mais se contente d'en utiliser les matrices compagnon.

    La situation sera différente dans le cas de la forme normal de Jordan.
\end{remark}

%---------------------------------------------------------------------------------------------------------------------------
\subsection{Forme normale de Jordan}
%---------------------------------------------------------------------------------------------------------------------------

Il existe une preuve directe de la réduction de Jordan ne nécessitant pas la réduction de Frobenius\cite{LecLinAlgAllen}. Cette dernière passe par les espaces caractéristiques\footnote{Aussi appelés «espaces propres généralisés».} et est à mon avis plus compliquée que la démonstration de Frobenius elle-même. Nous allons donc nous contenter de donner la réduction de Jordan comme un cas particulier de Frobenius.

\begin{theorem}[Réduction de Jordan]\index{réduction!Jordan}\index{Jordan!réduction}
    Soit \( E\) un espace vectoriel sur \( \eK\), et \( f\in\End(E)\) un endomorphisme dont le polynôme caractéristique \( \chi_f\) est scindé\footnote{C'est pour cette hypothèse que \( \eK=\eR\) n'est pas le bon cadre.}. Il existe une base de \( E\) dans laquelle la matrice de \( f\) s'écrit sous la forme
    \begin{equation}
        M=\begin{pmatrix}
            J_{n_1}(\lambda_1)    &       &       \\
                &   \ddots    &       \\
                &       &   J_{n_k}(\lambda_k)
        \end{pmatrix}
    \end{equation}
    où les \( \lambda_i\) sont les valeurs propres de \( f\) (avec éventuelle répétitions) et \( J_n(\lambda)\) représente le bloc \( n\times n\)
    \begin{equation}
        J_n(\lambda)=\begin{pmatrix}
            \lambda    &   1    &       &       &   \\  
                &   \lambda    &   1    &       &   \\  
                &       &   \lambda    &       &   \\  
                &       &       &   \ddots    &   1\\  
                &       &       &       &   \lambda    
        \end{pmatrix}.
    \end{equation}
    En d'autres termes, \( J_n(\lambda)_{ii}=\lambda\) et \( J_n(\lambda)_{i-1,i}=1\).    
\end{theorem}

\begin{proof}
    Nous commençons par le cas où \( f\) est nilpotente; nous notons \( M\) sa matrice. Dans ce cas la seule valeur propre est zéro et le polynôme caractéristique est \( X^m\) pour un certain \( m\). Nous savons par le lemme \ref{LemkVNisk} que (la matrice de) \( f\) est semblable à sa matrice compagnon. En l'occurrence pour \( f\) nous avons
    \begin{equation}
        C_{X^m}=\begin{pmatrix}
             0   &       &       &  0     \\
             1   &   \ddots    &       &   \vdots    \\
                &   \ddots    &   \ddots    &    \vdots   \\ 
                &       &   1    &   0     
         \end{pmatrix}.
    \end{equation}
    Ensuite le changement de base (qui est une similitude) \( (e_1,\ldots, e_n)\mapsto(e_n,\ldots, e_1)\) montre que \( C_{X^m}\) est semblable à un bloc de Jordan \( J_m(0)\).

    Supposons à présent que \( f\) ne soit pas nilpotente. Par l'hypothèse de polynôme caractéristique scindé, nous supposons que \( f\) a \( m\) valeurs propres distinctes et que son polynôme caractéristique est
    \begin{equation}
        \chi_f=(X-\lambda_1)^{l_1}\ldots (X-\lambda_m)^{l_m}.
    \end{equation}
    Le lemme des noyaux (théorème \ref{ThoDecompNoyayzzMWod}) nous enseigne que
    \begin{equation}
        E=\bigoplus_{i=1}^m\underbrace{\ker(f-\mu_i\mtu)^{l_i}}_{F_i}.
    \end{equation}
    La restriction de \( f-\lambda_i\mtu\) à \( F_i\) est par construction un endomorphisme nilpotent, et donc peut s'écrire comme un bloc de Jordan avec des zéros sur la diagonale. En utilisant la décomposition
    \begin{equation}
        f|_{F_i}=(f-\lambda_i\mtu)|_{F_i}+\lambda_i\mtu_{F_i},
    \end{equation}
    nous voyons que \( f|_{F_i}\) s'écrit comme un bloc de Jordan avec \( \lambda_i\) sur la diagonale.
\end{proof}

\begin{remark}
    Nous pouvons calculer la forme normale de Jordan pour une matrice complexe ou réelle, mais dans les deux cas nous devons nous attendre à obtenir une matrice complexe parce que les valeurs propres d'une matrice réelle peuvent être complexes. Cependant nous demandons que le polynôme caractéristique de \( f\) soit scindé sur \( \eK\). En pratique, la décomposition de Jordan n'est garantie que sur les corps algébriquement clos, c'est à dire sur \( \eC\).

    La suite des invariants de similitude sur laquelle repose Frobenius, elle, est disponible sur tout corps, y compris \( \eR\).
\end{remark}

Une application de la décomposition de Jordan est l'existence d'un logarithme pour les matrices.
\begin{proposition}
    Toute matrice inversible complexe est une exponentielle\index{exponentielle!de matrice}.
\end{proposition}

\begin{proof}
    Soit \( A\in \GL(n,\eC)\); nous allons donner une matrice \( B\in \eM(n,\eC)\) telle que \( A=\exp(B)\). D'abord remarquons qu'il suffit de prouver le résultat pour une matrice par classe de similitude. En effet si \( A=\exp(B)\) et si \( M\) est inversible alors 
    \begin{subequations}    \label{EqqACuGK}
        \begin{align}
            \exp(MBM^{-1})&=\sum_k\frac{1}{ k! }(MBM^{-1})^k\\
            &=\sum_k\frac{1}{ k! }MB^kM^{-1}\\
            &=M\exp(B)M^{-1}.
        \end{align}
    \end{subequations}
    Donc \( MAM^{-1}=\exp(MBM^{-1})\). Nous pouvons donc nous contenter de trouver un logarithme pour les blocs de Jordan. Nous supposons donc que \( A=(\mtu+N)\) avec \( N^m=0\). En nous inspirant de \eqref{EqweEZnV}, nous posons
    \begin{equation}
        D(t)=tN-\frac{ t^2 }{ 2 }N^2+\ldots +(-1)^m\frac{ t^{m-1} }{ m-1 }N^{m-1}
    \end{equation}
    et nous allons prouver que \(  e^{D(1)}=\mtu+N\). Notons que \( N\) étant nilpotente, cette somme ainsi que toutes celles qui viennent sont finies. Il n'y a donc pas de problèmes de convergences dans cette preuve (si ce n'est les passages des équations \eqref{EqqACuGK}).

    Nous posons \( S(t)= e^{D(t)}\) (la somme est finie), et nous avons
    \begin{equation}
        S'(t)=D'(t) e^{D(t)}
    \end{equation}
    Afin d'obtenir une expression qui donne \( S'\) en termes de \( S\), nous multiplions par \( (\mtu+tN)\) en remarquant que \( (\mtu+tN)D'(t)=N\) nous avons
    \begin{equation}
        (\mtu+tN)S'(t)=NS(t).
    \end{equation}
    En dérivant à nouveau,
    \begin{equation}    \label{EqKjccqP}
        (\mtu+tN)S''(t)=0.
    \end{equation}
    La matrice \( (\mtu+tN)\) est inversible parce que son noyau est réduit à \( \{ 0 \}\). En effet si \( (\mtu+tN)x=0\), alors \( Nx=-\frac{1}{ t }x\), ce qui est impossible parce que \( N\) est nilpotente. Ce que dit l'équation \eqref{EqKjccqP} est alors que \( S''(t)=0\). Si nous développons \( S(t)\) en puissances de \( t\) nous nous arrêtons au terme d'ordre \( 1\) et nous avons
    \begin{equation}
        S(t)=S(0)+tS'(0)=\mtu+tD'(0)=1+tN.
    \end{equation}
    En \( t=1\) nous trouvons \( S(1)=\mtu+N\). La matrice \( D(1)\) donnée est donc bien un logarithme de $\mtu+N$.
\end{proof}

%+++++++++++++++++++++++++++++++++++++++++++++++++++++++++++++++++++++++++++++++++++++++++++++++++++++++++++++++++++++++++++
\section{Mini introduction au produit tensoriel}
%+++++++++++++++++++++++++++++++++++++++++++++++++++++++++++++++++++++++++++++++++++++++++++++++++++++++++++++++++++++++++++
\label{SeOOpHsn}

%---------------------------------------------------------------------------------------------------------------------------
\subsection{Définitions}
%---------------------------------------------------------------------------------------------------------------------------

Soit \( E\), un espace vectoriel de dimension finie. Si \( \alpha\) et \( \beta\) sont deux formes linéaires sur un espace vectoriel \( E\), nous définissons \( \alpha\otimes \beta\) comme étant la \( 2\)-forme donnée par
\begin{equation}
    (\alpha\otimes \beta)(u,v)=\alpha(u)\beta(v).
\end{equation}
Si \( a\) et \( b\) sont des vecteurs de \( E\), ils sont vus comme des formes sur \( E\) via le produit scalaire et nous avons
\begin{equation}
    (a\otimes b)(u,v)=(a\cdot u)(b\cdot v).
\end{equation}
Cette dernière équation nous incite à pousser un peu plus loin la définition de \( a\otimes b\) et de simplement voir cela comme la matrice de composantes
\begin{equation}
    (a\otimes b)_{ij}=a_ib_j.
\end{equation}
Cette façon d'écrire a l'avantage de ne pas demander de se souvenir qui est une vecteur ligne, qui est un vecteur colonne et où il faut mettre la transposée. Évidemment \( (a\otimes b)\) est soit \( ab^t\) soit \( a^tb\) suivant que \( a\) et \( b\) soient ligne ou colonne.

\begin{lemma}   \label{LemMyKPzY}
    Soient \( x,y\in E\) et \( A,B\) deux opérateurs linéaires sur \( E\) vus comme matrices. Alors
    \begin{equation}        \label{EqXdxvSu}
        (Ax\otimes By)=A(x\otimes y)B^t.
    \end{equation}
\end{lemma}

\begin{proof}
    Calculons la composante \( ij\) de la matrice \( (Ax\otimes By)\). Nous avons
    \begin{subequations}
        \begin{align}
            (Ax\otimes By)_{ij}&=(Ax)_i(By)_j\\
            &=\sum_{kl}A_{ik}x_kB_{jl}y_l\\
            &=A_{ik}(x\otimes y)_{kl}B_{jl}\\
            &=\big( A(x\otimes y)B^t \big)_{ij}.
        \end{align}
    \end{subequations}
\end{proof}

%+++++++++++++++++++++++++++++++++++++++++++++++++++++++++++++++++++++++++++++++++++++++++++++++++++++++++++++++++++++++++++
\section{Espaces hermitiens}
%+++++++++++++++++++++++++++++++++++++++++++++++++++++++++++++++++++++++++++++++++++++++++++++++++++++++++++++++++++++++++++

\begin{definition}  \label{DefMZQxmQ}
Si \( E\) est un espace vectoriel sur \( \eC\), nous disons qu'une application \( \langle ., .\rangle \colon E\times E\to \eC\) est un \defe{produit scalaire hermitien}{produit!scalaire!hermitien}\index{hermitien!produit scalaire} si pour tout \( u,v\in E\) nous avons
\begin{enumerate}
    \item
        \( \langle u, v\rangle =\overline{ \langle v, u\rangle  }\)
    \item
        \( \lambda\langle u, v\rangle =\langle \lambda u, v\rangle =\langle u, \bar \lambda v\rangle \)
    \item
        \( \langle u, u\rangle \in \eR^+\) et \( \langle u, u\rangle =0\) si et seulement si \( u=0\).
\end{enumerate}
\end{definition}

%+++++++++++++++++++++++++++++++++++++++++++++++++++++++++++++++++++++++++++++++++++++++++++++++++++++++++++++++++++++++++++
\section{Formes bilinéaires et quadratiques}
%+++++++++++++++++++++++++++++++++++++++++++++++++++++++++++++++++++++++++++++++++++++++++++++++++++++++++++++++++++++++++++

%--------------------------------------------------------------------------------------------------------------------------- 
\subsection{Isométries d'espaces euclidiens}
%---------------------------------------------------------------------------------------------------------------------------

À une forme bilinéaire \( b\) nous associons la forme quadratique \( q(x)=b(x,x)\). Étant donné une forme quadratique, la forme bilinéaire peut être retrouvée grâce aux \defe{identités de polarisation}{identité!polarisation}\index{polarisation (identité)} :
\begin{equation}    \label{EqMrbsop}
    b(x,y)=\frac{ 1 }{2}\big( q(x)+q(y)-q(x-y) \big).
\end{equation}
Une forme bilinéaire est \defe{non dégénérée}{forme!bilinéaire!non dégénérée} lorsque l'unique \( x\in E\) tel que \( b(x,z)=0\) pour tout \( z\) est \( x=0\).

\begin{lemma}   \label{LemyKJpVP}
    Soit \( b\) une forme bilinéaire non dégénérée. Si \( x\) et \( y\) sont tels que \( b(x,z)=b(y,z)\) pour tout \( z\), alors \( x=y\).
\end{lemma}

\begin{proof}
    C'est immédiat du fait de la linéarité en le premier argument et de la non-dégénérescence : si \( b(x,z)-b(y,z)=0\) alors
    \begin{equation}
        b(x-y,z)=0
    \end{equation}
    pour tout \( z\), ce qui implique \( x-y=0\).
\end{proof}

\begin{proposition}
    La forme bilinéaire \( b\) est non-dénénérée si et seulement si sa matrice associée est inversible.
\end{proposition}

\begin{proof}
    Nous savons que la matrice associée est symétrique et qu'elle peut donc être diagonalisée (théorème \ref{ThoeTMXla}). En nous plaçant dans une base de diagonalisation, nous devons prouver que la forme est non-dégénérée si et seulement si les éléments diagonaux de la matrice sont tous non nuls.

    Écrivons \( b(x,z)\) en choisissant pour \( z\) le vecteur de base \( e_k\) de composantes \( (e_k)_j=\delta_{kj}\) :
    \begin{equation}
            b(x,e_k)=\sum_{ij}x_i(e_k)_j
            =\sum_i b_{ik}x_i
            =b_{kk}x_k.
    \end{equation}
    Si \( b\) est dégénérée et si \( x\) est un vecteur non nul (disons que la composante \( x_i\) est non nulle) de \( E\) tel que \( b(x,z)=0\) pour tout \( z\in E\), alors \( b_{ii}=0\), ce qui montre que la matrice de \( b\) n'est pas inversible.

    Réciproquement si la matrice de \( b\) est inversible, alors tous les \( b_{kk}\) sont différents de zéro, et le seul vecteur \( x\) tel que \( b_{kk}x_k=0\) pour tout \( k\) est le vecteur nul.
\end{proof}

\begin{example}
    La forme quadratique \( q(x)=x_1^2+x_2^2\) donne la norme euclidienne. La forme bilinéaire associée est \( b(x,y)=x_1y_1+x_2y_2\), qui est le produit scalaire usuel.
\end{example}

Il ne faudrait pas déduire trop vite que la formule \( \| x \|^2=q(x)\) donne une norme dès que \( q\) est non dégénérée. En effet \( q\) peut ne pas être définie positive. La forme \( q(x)=x_1^2-x_2^2\) prend des valeurs positives et négatives. A fortiori \( d(x,y)=q(x-y)\) ne donne pas toujours une distance.

Nous allons cependant appeler \defe{isométrie}{isométrie!de forme quadratique} pour la forme \( q\) une application bijective \( f\colon V\to V\) telle que \( q(x-y)=q\big( f(x)-f(y) \big)\). Dans les cas où \( q\) donne une distance, alors c'est une isométrie au sens usuel.

\begin{lemma}   \label{LemewGJmM}
    Pour une application bijective \( f\colon E\to E\) telle que \( f(0)=0\), les conditions suivantes sont équivalentes: 
    \begin{enumerate}
        \item
            \( b\big( f(x),f(y) \big)=b(x,y)\) pour tout \( x,y\in E\);
        \item
            \( q\big( f(x)-f(y) \big)=q(x-y)\) pour tout \( x,y\in E\).
    \end{enumerate}
\end{lemma}

\begin{proof}
    Dans le sens direct, en posant \( x=y\) nous trouvons tout de suite \( q(f(x))=q(f)\); ensuite en utilisant la distributivité de \( b\),
    \begin{subequations}
        \begin{align}
            q\big( f(x)-f(y) \big)&=b\big( f(x)-f(y),f(x)-f(y) \big)\\
            &=q\big( f(x) \big)-2b\big( f(x),f(y) \big)+q\big( f(y) \big)\\
            &=q(x)+q(y)-2b(x,y)\\
            &=q(x-y).
        \end{align}
    \end{subequations}
    
    Dans l'autre sens, nous commençons par remarquer que l'hypothèse \( f(0)=0\) donne \( q(x)=q\big( f(x) \big)\). Ensuite nous utilisons l'identité de polarisation \eqref{EqMrbsop} :
    \begin{subequations}
        \begin{align}
            b\big( f(x),f(y) \big)&=\frac{ 1 }{2}\big[ q\big( f(x) \big)+q\big( f(y) \big)-q\big( f(x-y) \big) \big]\\
            &=\frac{ 1 }{2}\big[ q(x)+q(y)-q(x-y) \big]\\
            &=b(x,y).
        \end{align}
    \end{subequations}
\end{proof}

\begin{theorem}     \label{ThoDsFErq}
    Soit \( f\colon E\to E\) une bijection telle que
    \begin{equation}
        q(x-y)=q\big( f(x)-f(y) \big)
    \end{equation}
    pour tout \( x,y\in E\). Alors
    \begin{enumerate}
        \item
            si \( f(0)=0\), alors \( f\) est linéaire;
        \item
            si \( f(0)\neq 0\) alors \( f\) est affine.
    \end{enumerate}
\end{theorem}
La rédaction la preuve a bénéficié d'un coup de main de la part de \href{http://www.ilemaths.net/forum-sujet-500814.html}{GaBuZoMeu}. Une autre preuve, utilisant un peut plus d'indices et un peu plus de mots comme «tenseurs», peut être trouvée  \href{http://physics.stackexchange.com/questions/12664/proving-that-interval-preserving-transformations-are-linear}{ici}. Le fait que la preuve donnée soit tensorielle me fait penser que le résultat peut encore être généralisé.

\begin{proof}
    Si \( f(0)=0\), nous savons par le lemme \ref{LemewGJmM} que \( b\big( f(x),f(y) \big)=b(x,y)\). Soit \( z\in E\); étant donné que \( f\) est bijective nous pouvons considérer l'élément \( f^{-1}(z)\in E\) et calculer
    \begin{subequations}
        \begin{align}
            b\big( f(x+y),z \big)&=b\big( f(x+y),f(f^{-1}(z)) \big)\\
            &=b(x+y,f^{-1}(z))\\
            &=b(x,f^{-1}(z))+b(y,f^{-1}(z))\\
            &=b(f(x),z)+b(f(y),z)\\
            &=b\big( f(x)+f(y),z \big),
        \end{align}
    \end{subequations}
    donc \( f(x+y)=f(x)+f(y)\) par le lemme \ref{LemyKJpVP}. 

    De la même façon on trouve \( b\big( f(\lambda x),z \big)=b\big( \lambda f(x),z \big)\) qui prouve que \( f(\lambda x)=\lambda f(x)\) et donc que \( f\) est linéaire.

    Si \( f(0)\neq 0\), alors nous posons \( g(x)=f(x)-f(0)\) qui vérifie \( g(0)=0\) et
    \begin{equation}
        q\big( g(x)-g(y) \big)=q\big( f(x)-f(0)-f(y)+f(0) \big)=q(x-y).
    \end{equation}
    Nous pouvons donc appliquer le premier point à \( g\), déduire que \( g\) est linéaire et donc que \( f\) est affine.
\end{proof}

Maintenant nous savons que le groupe des isométries d'un espace quadratique \( (E,q)\) est un sous-groupe de \( \GL(E)\). Dans le cas de la métrique euclidienne, il est connu que ce sont les matrices orthogonales.

Nous pouvons maintenant avoir une discussion plus détaillée des groupes d'isométries de l'espace euclidien, parce que nous savons maintenant qu'elles sont des applications linéaires. Pour en savoir plus sur le groupe des isométries, il faut lire le théorème de Cartan-Dieudonné dans \cite{JGAdTA}.

Soit \( f\), une forme bilinéaire symétrique non dégénérée  sur l'espace vectoriel \( E\) de dimension \( n\) sur \( \eK\) où \( \eK\) est un corps de caractéristique différente de \( 2\). Nous notons \( q\) la forme quadratique associée.

Un vecteur est \defe{isotrope}{isotrope (vecteur)} si il est perpendiculaire à lui-même; en d'autres termes, \( x\) est isotrope si et seulement si \( f(x,x)=0\). Un sous-espace \( W\subset E\) est \defe{totalement isotrope}{isotrope!totalement} si pour tout \( x,y\in W\), nous avons \( f(x,y)=0\).

\begin{lemma}[\cite{JGAdTA}]
    Si \( n\geq 3\), alors toute droite est intersection de deux plans non isotropes.
\end{lemma}

%--------------------------------------------------------------------------------------------------------------------------- 
\subsection{Théorème de Sylvester}
%---------------------------------------------------------------------------------------------------------------------------

% TODO : Il y a une démonstration sur wikipédia, à voir.

\begin{theorem}[de Sylvester]   \label{ThoQFVsBCk}
    Soit $Q$ une forme quadratique réelle de signature \( (p,q)\). Alors pour toute base orthonormée on a
    \begin{subequations}
        \begin{align}
            p&=\Card\{ i\tq Q(e_i)>0 \}\\
            q&=\Card\{ i\tq Q(e_i)<0 \}.
        \end{align}
    \end{subequations}
    Le rang de \( Q\) est \( p+q\).

    Si \( A\) est la matrice de \( Q\) dans une base, alors il existe une matrice inversible \( P\) telle que
    \begin{equation}
        P^tAP=\begin{pmatrix}
            -\mtu_q    &       &       \\
                &   \mtu_p    &       \\
                &       &   0
        \end{pmatrix}.
    \end{equation}
\end{theorem}
\index{théorème!Sylvester}
\index{rang}
\index{matrice!semblables}
\index{forme!quadratique}

%+++++++++++++++++++++++++++++++++++++++++++++++++++++++++++++++++++++++++++++++++++++++++++++++++++++++++++++++++++++++++++
\section{Isométries de l'espace euclidien}
%+++++++++++++++++++++++++++++++++++++++++++++++++++++++++++++++++++++++++++++++++++++++++++++++++++++++++++++++++++++++++++

Nous considérons l'espace affine euclidien \( A=\affE_n(\eR)\) modelé sur \( \eR^n\) avec sa métrique usuelle. Nous avons montré par le théorème \ref{ThoDsFErq} que les isométries de cet espaces sont des applications linéaires.

%---------------------------------------------------------------------------------------------------------------------------
\subsection{Produit semi-direct}
%---------------------------------------------------------------------------------------------------------------------------

Les isométries directes\footnote{Directes au sens où nous ne considérons pas les retournements.} de cet espace sont données d'une part par les rotations de \( \SO(n)\) et d'autre part par les translations données par les vecteurs de \( \eR^n\). Plus précisément, un couple \( (v,\Lambda)\in \eR^n\times\SO(n)\) agit sur \( x\in A\) par
\begin{equation}
    (v,\Lambda)x=\Lambda x+v.
\end{equation}
La loi de composition est donnée par
\begin{subequations}
    \begin{align}
        (v,\Lambda)\cdot(v',\Lambda')x&=(v,\Lambda)(\Lambda'x+v')\\
        &=\Lambda\Lambda'x+\Lambda v'+v\\
        &=(\Lambda v'+v,\Lambda\Lambda')x.
    \end{align}
\end{subequations}
Nous avons donc, pour tout \( v,v'\in \eR^n\), \( \Lambda,\Lambda'\in\SO(n)\) la loi de groupe
\begin{equation}    \label{EqDiHcut}
        (v,\Lambda)\cdot(v',\Lambda')=(\Lambda v'+v,\Lambda\Lambda').
\end{equation}
    
Le groupe \( \SO(n)\) agit naturellement sur \( \eR^n\) par
\begin{equation}
    \begin{aligned}
        \phi\colon \SO(n)&\to \Aut(\eR^n) \\
        \Lambda&\mapsto \phi_{\Lambda}\colon v\to \Lambda v. 
    \end{aligned}
\end{equation}
Il est à noter qu'ici, \( \eR^n\) est vu comme l'ensemble des applications \( v\colon A\to A\), \( v(x)=x+a\). Voir aussi la remarque \ref{RemAobrlX}.

Nous pourrions alors présenter le groupe de isométries de \( A\) sous la forme du produit semi-direct
\begin{equation}
    \Iso^+(A)=\eR^n\times_{\phi}\SO(n).
\end{equation}
Plusieurs choses sont à vérifier :
\begin{enumerate}
    \item
        Pour chaque \( \Lambda\), l'application \( \phi_{\Lambda}\) est un automorphisme du groupe \( \eR^n\) (en tant qu'agissant sur \( A\)). Le fait que \( \phi_{\Lambda}\) soit une bijection n'est pas un problème. Nous devons vérifier que
        \begin{equation}
            \phi_{\Lambda}(v+w)=\phi_{\Lambda}(v)\circ\phi_{\Lambda}(w)
        \end{equation}
        en tant qu'égalité dans l'ensemble des isométries de \( A\). Nous la testons donc sur un élément \( x\in A\). D'une part
        \begin{equation}
            \phi_{\Lambda}(v+w)x=x+\Lambda(v+m),
        \end{equation}
        et d'autre part,
        \begin{equation}
            \phi_{\Lambda}(v)\circ\phi_{\Lambda}(w)x=\phi_{\Lambda}(v)\big( x+\Lambda w \big)=x+\Lambda w+\Lambda v.
        \end{equation}
    \item
        L'application \( \phi\colon \SO(n)\to \Aut(\eR^n)\) est un morphisme de groupe. Nous devons vérifier l'égalité
        \begin{equation}
            \phi_{\Lambda\Lambda'}=\phi_{\Lambda}\circ\phi_{\Lambda'}
        \end{equation}
        dans \( \Aut(\eR^n)\), c'est à dire que pour tout \( v\in \eR^n\) et \( x\in A\) nous devons avoir
        \begin{equation}
            \phi_{\Lambda\Lambda'}(v)x=\big( \phi_{\Lambda}\circ\phi_{\Lambda'}\big)(v)x.
        \end{equation}
        Le membre de gauche fait immédiatement \( x+\Lambda\Lambda'v\) tandis que le membre de droite vaut
        \begin{equation}
            \big( \phi_{\Lambda}\circ\phi_{\Lambda'}\big)(v)x=\big( \phi_{\Lambda}(\Lambda'v) \big)x=(\Lambda\Lambda'v)x=x+\Lambda\Lambda'v.
        \end{equation}
    \item
        La loi de groupe donnée par \( \phi\) sur \( \SO(n)\times \eR^n\) par la définition \eqref{EqDRgbBI} est bien la loi de groupe \eqref{EqDiHcut}. Cela est encore un calcul immédiat. L'utilisation de la définition \eqref{EqDRgbBI} donne
        \begin{equation}
            (v,\Lambda)\cdot(v',\Lambda')=(v+\phi_{\Lambda}(v'),\Lambda\Lambda')=(v+\Lambda v',\Lambda\Lambda'),
        \end{equation}
        qui est bien la formule \eqref{EqDiHcut}.
\end{enumerate}

%---------------------------------------------------------------------------------------------------------------------------
\subsection{Groupe diédral}
%---------------------------------------------------------------------------------------------------------------------------
\label{subsecHibJId}

%///////////////////////////////////////////////////////////////////////////////////////////////////////////////////////////
\subsubsection{Définition et générateurs : vue géométrique}
%///////////////////////////////////////////////////////////////////////////////////////////////////////////////////////////

Le \defe{groupe diédral}{groupe!diédral} \( D_n\)\nomenclature[R]{\( D_n\)}{groupe diédral} est le groupe des isométries de \( \eR^2\) laissant invariant un polygone régulier à \( n\) côtés. Il peut être vu comme le stabilisateur de l'ensemble
\begin{equation}
    \{  e^{2ik\pi/n},k=0,\ldots, n-1 \}
\end{equation}
dans le groupe des isométries affines de \( \eC^*\).
\index{groupe!agissant sur un ensemble!diédral}
\index{groupe!en géométrie}
\index{groupe!fini!diédral}
\index{groupe!permutation!diédral}
% TODO : prouver que les racines de l'unité forment un polygone régulier.

Si \( f\in D_n\), alors \( f( e^{2ik\pi/n}) \) doit être l'un des \(  e^{2ik'\pi/n}\), et vu que \( f\) conserve les longueurs dans \( \eC\), nous devons avoir
\begin{equation}
    1=d(0, e^{2ik\pi/n})=d\big( f(0), e^{2ik'\pi/n} \big).
\end{equation}
Donc \( f(0)\) est à l'intersection de tous les cercles de rayon \( 1\) centrés en les \(  e^{2ik\pi/n}\), ce qui montre que \( f(0))0\) (dès que \( n\geq 3\)). Par conséquent notre étude du groupe diédral ne doit prendre en compte que les isométries vectorielles de \( \eR^2\). En d'autres termes
\begin{equation}
    D_n\subset O(2,\eR).
\end{equation}

\begin{proposition}[\cite{tzHydF}]
    Le groupe \( D_n\) contient un sous groupe cyclique d'ordre \( 2\) et un sous groupe cyclique d'ordre \( n\).
\end{proposition}

\begin{proof}
    Si \( s\) est la réflexion d'axe \( \eR\), alors \( s\) est d'ordre \( 2\). De plus \( s\) est bien dans \( D_n\) parce que
    \begin{equation}    \label{EqSUshknP}
        s\big(  e^{2ki\pi/n} \big)= e^{2(n-k)i\pi/n}.
    \end{equation}

    De la même façon, la rotations d'angle \(2\pi/n\), que l'on note \( r\), agit sur les racines de l'unité et engendre un le groupe d'ordre \( n\) des rotations d'angle \(2 k\pi/n\).
\end{proof}

Notons que la conjugaison complexe ne fait pas spécialement partie du groupe \( D_n\). En effet pour \( n=3\) par exemple les points fixes sont \( A_1=(1,0)\), \( A_2=(-\frac{ 1 }{2},\frac{ \sqrt{3} }{2})\) et \( A_3=(\frac{ 1 }{2},-\frac{ \sqrt{3} }{2})\). La conjugaison complexe envoie évidemment \( A_1\) sur \( A_1\), mais pas du tout \( A_2\) sur \( A_3\).
%TODO : un dessin du triangle équilatéral serait pas mal ici.

\begin{proposition}[\cite{tzHydF}]
    Nous avons \( (sr)^2=\id\).
\end{proposition}

\begin{proof}
    Si \( z^n=1\), alors
    \begin{equation}
        (srsr)z=srs e^{2 i\pi/n}z=sr\big( e^{-2\pi i/n\bar z}\big)=s\bar z=z.
    \end{equation}
\end{proof}

\begin{proposition}[\cite{tzHydF}] \label{PropLDIPoZ}
    Le groupe diédral \( D_n\) est engendré par \( s\) et \( r\). De plus tous les éléments de \( D_n\) s'écrivent sous la forme \( s\circ r^m\).
\end{proposition}
\index{groupe!diédral!générateurs (preuve)}
\index{racine!de l'unité}
\index{géométrie!avec nombres complexes}
\index{géométrie!avec des groupes}

\begin{proof}
    Nous considérons les points \( A_0=1\) et \( A_k= e^{2ki\pi/n}\) avec \( k\in\{ 1,\ldots, n-1 \}\). Par convention, \( A_n=A_0\). L'action des éléments \( s\) et \( r\) sur ces points est
    \begin{subequations}
        \begin{align}
            r(A_k)&=A_{k+1}\\
            s(A_k)&=A_{n-k}.
        \end{align}
    \end{subequations}
    Cette dernière est l'équation \eqref{EqSUshknP}.
    
    Soit \( f\in D_n\). Étant donné que c'est une isométrie de \( \eR^2\) avec un point fixe (le point \( 0\)), \( f\) est soit une rotation soit une réflexion.
    %TODO : il faut démontrer ce point et mettre un lien vers ici.

    Supposons pour commencer que un des \( A_k\) est fixé par \( f\). Dans ce cas \( f\) a deux points fixes : \( O\) et \( A_k\) et est donc la réflexion d'axe \( (OA_k)\). Dans ce cas, nous avons \( f=s\circ r^{n-2k}\). En effet
    \begin{equation}
        s\circ r^{n-2k}(A_k)=s(A_{k+n-2k})=s(A_{n-k})=A_k.
    \end{equation}
    Donc \( O\) et \( A_k\) sont deux points fixes de l'isométrie \( f\); donc \( f\) est bien la réflexion sur le bon axe.

    Nous passons à présent au cas où \( f\) ne fixe aucun des \( A_k\). 
    \begin{enumerate}
        \item
            Supposons que \( f\) soit une rotation. Si \( f(A_k)=A_m\), alors l'angle de la rotation est 
            \begin{equation}
                \frac{ 2(m-k)\pi }{ n },
            \end{equation}
            et donc \( f=r^{m-k}\), qui est de la forme demandée.
        \item
            Supposons à présent que \( f\) soit une réflexion d'axe \( \Delta\). Cette fois, \( \Delta\) ne passe par aucun des points \( A_k\), par contre \( \Delta\) passe par \( 0\). Nous commençons par montrer que \( \Delta\) doit être la médiatrice d'un des côtés \( [A_p,A_{p+1}]\) du polygone. Vu que \( \Delta\) passe par \( O\) et n'est aucune des droites \( (OA_k)\), cette droite passe par l'intérieur d'un des triangles \( OA_pA_{p+1}\) et intersecte donc le côté correspondant.

            Notre tâche est de montrer que \( \Delta\) coupe \( [A_p,A_{p+1}]\) en son milieu. Dans ce cas, \( \Delta\) sera automatiquement perpendiculaire parce que le triangle \( OA_pA_{p+1}\) est isocèle en \( O\). Nommons \( l\) la longueur des côtés du polygone, \( P=\Delta\cap[A_p,A_{p+1}]\), \( x=d(A_p,P)\) et \( \delta=d(A_p,\Delta)\). Vu que \( f\) est la symétrie d'axe \( \Delta\), nous avons aussi \( d\big( f(A_p),\Delta \big)=\delta\) et \( d\big( A_p,f(A_p) \big)=2\delta\). D'autre part, par la définition de la distance, \( \delta<x\). Si \( x<\frac{ l }{2}\), alors \( \delta<\frac{ \delta }{2}\) et donc \( d\big( A_p,f(A_p) \big)<l\). Or cela est impossible parce que le polygone ne possède aucun sommet à distance plus courte que \( l\) de \( A_p\).

            De la même manière si \( x>\frac{ l }{2}\), nous raisonnons avec \( A_{p+1}\) pour obtenir une contradiction. Nous en concluons que la seule possibilité est \( x=\frac{ l }{2}\), et donc \( f(A_p)=A_{p+1}\). Montrons alors que \( f=s\circ r^{n-2p-1}\). Il faut montrer que c'est une réflexion qui envoie \( A_p\) sur \( A_{p+1}\). D'abord c'est une réflexion parce que
            \begin{equation}
                \det(sr^{n-2p-1})=\det(s)\det(r^{n-2p-1})=-1
            \end{equation}
            parce que \( \det(s)=-1\) alors que \( \det(r^k)=1\) parce que \( r\) est une rotation dans \( \SO(2)\). Ensuite nous avons
            \begin{equation}
                s\circ r^{n-2p-1}(A_p)=s(A_{p+n-2p-1})=s(A_{n-p-1})=A_{n-(n-p-1)}=A_{p+1}.
            \end{equation}

            Donc \( s\circ r^{n-2p-1}\) est bien une réflexion qui envoie \( A_p\) sur \( A_{p+1}\).

    \end{enumerate}
\end{proof}

\begin{corollary}   \label{CorWYITsWW}
La liste des éléments de \( D_n\) est 
\begin{equation}
    D_n=\{ 1,r,\ldots, r^{n-1},s,sr,\ldots, sr^{n-1} \}
\end{equation}
et \( | D_n |=2n\).
\end{corollary}

\begin{proof}
    Nous savons par la proposition \ref{PropLDIPoZ} que tous les élément de \( D_n\) s'écrivent sous la forme \( r^k\) ou \( sr^k\). Vu que \( r\) est d'ordre \( n\), il ne faut considérer que \( k\in\{ 1,\ldots, n-1 \}\). Les éléments \( 1\), \( r\),\ldots, \( r^{n-1}\) sont tous différents, et sont (pour des raisons de déterminant) tous différents des \( sr^k\). Les isométries \( sr^k\) sont toutes différentes entre elles pour essentiellement la même raison :
    \begin{equation}
        sr^k(A_p)=s(A_{p+k})=A_{n-p+k}
    \end{equation}
    donc si \( k\neq k'\), \( sr^k(A_p)\neq sr^{k'}(A_p)\). La liste des éléments de \( D_n\) est donc
    \begin{equation}
        D_n=\{ 1,r,\ldots, r^{n-1},s,sr,\ldots, sr^{n-1} \}
    \end{equation}
    et donc \( | D_n |=2n\).
\end{proof}

\begin{example}
    Nous considérons le carré \( ABCD\) dans \( \eR^2\) et nous cherchons les isométries de \( \eR^2\) qui laissent le carré invariant. Nous nommons les points comme sur la figure \ref{LabelFigIsomCarre}. La symétrie d'axe vertical est nommée \( s\) et la rotation de \( 90\) degrés est notée \( r\).
    \newcommand{\CaptionFigIsomCarre}{Le carré dont nous étudions le groupe diédral.}
    \input{Fig_IsomCarre.pstricks}

    Il est facile de vérifier que toutes les symétries axiales peuvent être écrites sous la forme \( r^is\). De plus le groupe engendré par \( s\) agit sur le groupe engendré par \( r\) parce que
    \begin{equation}
        (srs^{-1})(A,B,C,D)=sr(B,A,D,C)=s(A,D,C,B)=(B,C,D,A),
    \end{equation}
    c'est à dire \( srs^{-1}=r^{-1}\). Nous sommes alors dans le cadre du corollaire \ref{CoroGohOZ} et nous pouvons écrire que
    \begin{equation}
        D_4=\gr(r)\times_{\sigma}\gr(s).
    \end{equation}
\end{example}

%///////////////////////////////////////////////////////////////////////////////////////////////////////////////////////////
\subsubsection{Générateurs : vue abstraite}
%///////////////////////////////////////////////////////////////////////////////////////////////////////////////////////////

Nous allons montrer que \( D_n\) peut être décrit de façon abstraite en ne parlant que de ses générateurs. Nous considérons un groupe \( G\) engendré par des éléments \( a\) et \( b\) tels que
\begin{enumerate}
    \item
        \( a\) est d'ordre \( 2\),
    \item
        \( b\) est d'ordre \( n\) avec \( n\geq 3\),
    \item
        \( abab=e\).
\end{enumerate}
Nous allons prouver que ce groupe doit avoir la même liste d'éléments que celle du corollaire \ref{CorWYITsWW}.

\begin{proposition}[\cite{tzHydF}]
    Le groupe \( G\) n'est pas abélien.
\end{proposition}

\begin{proof}
    Nous savons que \( abab=e\), donc \( abab^{-1}=b^{-2}\), mais \( b^{-2}\neq e\) parce que \( b\) est d'ordre \( n>2\). Donc \( abab^{-1}\neq e\). En manipulant un peu :
    \begin{equation}
        e\neq abab^{-1}=(ab)(ba^{-1})^{-1}=(ab)(ba)^{-1}
    \end{equation}
    parce que \( a^{-1}=a\). Donc \( ab\neq ba\).
\end{proof}

\begin{lemma}[\cite{tzHydF}]        \label{LemKKXdqdL}
    Pour tout \( k\) entre \( 1\) et \( n-1\) nous avons
    \begin{equation}
        ab^ka=b^{-k}.
    \end{equation}
\end{lemma}

\begin{proof}
    Nous faisons la démonstration par récurrence. D'abord pour \( k=1\), nous devons avoir \( aba=b^{-1}\), ce qui est correct parce que par construction de \( G\) nous avons \( abab=e\). Ensuite nous supposons que le lemme tient pour \( k\) et nous regardons ce qu'il se passe avec \( k+1\) :
    \begin{equation}
            ab^{k+1}ba=ab^kba=\underbrace{ab^ka}_{b^{-k}}\underbrace{aba}_{b^{-1}}=b^{-k}b^{-1}=b^{-(k+1)}.
    \end{equation}
\end{proof}

\begin{proposition}
    L'élément \( a\) n'est pas une puissance de \( b\).
\end{proposition}

\begin{proof}
    Supposons le contraire : \( a=b^k\). Dans ce cas nous aurions
    \begin{equation}
        e=(ab)(ab)=b^{k+1}b^{k+1}=b^{2k+2}=b^{2k}b^2=a^2b^2=b^2,
    \end{equation}
    ce qui signifierait que \( b\) est d'ordre \( 2\), ce qui est exclu par construction.
\end{proof}

\begin{proposition}[\cite{tzHydF}]
    La liste des éléments de \( G\) est donnée par
    \begin{equation}
        G=\{ 1,b,\cdots,b^{n-1},a,ab,\ldots, ab^{n-1} \}.
    \end{equation}
\end{proposition}

\begin{proof}
    Étant donné que \( a\) n'est pas une puissance de \( b\), les éléments \( 1\), \( a\), \( b\),\ldots, \( b^{n-1}\) sont distincts. De plus si \( k\) et \( m=k+p\) sont deux éléments distincts de \( \{ 1,\ldots, n-1 \}\), nous avons \( ab^k\neq ab^m\) parce que si \( ab^k=ab^{k+p}\), alors \( a=ab^p\) avec \( p<n\), ce qui est impossible. Pour la même raison, \( ab^k\neq e\), et \( ab^k\neq b^m\).

    Au final les éléments \( 1,a,b,\ldots, b^{n-1},ab,\ldots, ab^{n-1}\) sont tous différents. Nous devons encore voir qu'il n'y en a pas d'autres.

    Par définition le groupe \( G\) est engendré par \( a\) et \( b\), donc tout élément \( x\in G\) s'écrit $x=a^{m_1}b^{k_1}\ldots a^{m_r}b^{k_r}$ pour un certain \( r\) et avec pour tout \( i\), \( k_i\in\{ 1,\ldots, n-1 \}\) (sauf \( k_r\) qui peut être égal à zéro) et \( m_i=1\), sauf \( m_1\) qui peut être égal à zéro. Donc
    \begin{equation}
        x=a^mb^{k_1}ab^{k_2}a\ldots b^{k_{r-1}}ab^{k_r}
    \end{equation}
    où \( m\) et \( k_r\) peuvent éventuellement être zéro. En utilisant le lemme \ref{LemKKXdqdL} sous la forme \( b^{k_i}a=ab^{-k_i}\), quitte à changer les valeurs des exposants, nous pouvons passer tous les \( a \) à gauche et tous les \( b\) à droite pour finir sous la forme \( x=a^kb^m\). 

    Donc non, il n'existe pas d'autres éléments dans \( G\) que ceux déjà listés.

\end{proof}

\begin{theorem}
    Les groupes \( G\) et \( D_n\) sont isomorphes.
\end{theorem}

\begin{proof}
        Nous utilisons l'application
    \begin{equation}
        \begin{aligned}
            \psi\colon G&\to D_n \\
            a^kb^m&\mapsto s^kr^m. 
        \end{aligned}
    \end{equation}
    C'est évidemment bien défini et bijectif, mais c'est également un homomorphisme parce que si nous calculons \( \psi\) sur un produit, nous devons comparer
    \begin{equation}        \label{EqBULPilp}
        \psi\big( a^{k_1}b^{m_1}a^{k_2}b^{m_2} \big)
    \end{equation}
    avec
    \begin{equation}        \label{EqIVEIphI}
        \psi\big( a^{k_1}b^{m_1}\big)\psi\big(a^{k_2}b^{m_2} \big)= s^{k_1}r^{m_1}s^{k_2}r^{m_2}.
    \end{equation}
    Vu que \( D_n\) et \( G\) ont les mêmes propriétés qui permettent de permuter \( a\) et \( b\) ou \( s\) et \( r\), l'expression à l'intérieur du \( \psi\) dans \eqref{EqBULPilp} se simplifie en \( a^kb^m\) avec les même \( k\) et \( n\) que l'expression à droite dans \eqref{EqIVEIphI} ne se simplifie en \( s^kr^m\).
\end{proof}

\begin{corollary}
    Toutes les propriétés démontrées pour \( G\) sont vraies pour \( D_n\). En particulier, avec quelque redites :
    \begin{enumerate}
        \item
            Le groupe \( D_n\) peut être défini comme étant le groupe engendré par un élément \( s\) d'ordre \( 2\) et un élément \( r\) d'ordre \( n-1\) assujettis à la relation \( srsr=e\).
        \item
            Le groupe \( D_n\) n'est pas abélien.
        \item
            Pour tout \( k\in\{ 1,\ldots, n-1 \}\) nous avons \( sr^ks=r^{-k}\).
        \item
            L'élément \( s\) ne peut pas être obtenu comme une puissance de \( r\).
        \item
            La liste des éléments de \( D_n\) est
            \begin{equation}
                D_n=\{ 1,r,\ldots, r^{n-1},s,sr,\ldots, sr^{n-1} \}
            \end{equation}
        \item
            Le groupe diédral \( D_n\) est d'ordre \( 2n\).
    \end{enumerate}
\end{corollary}

%///////////////////////////////////////////////////////////////////////////////////////////////////////////////////////////
\subsubsection{Classes de conjugaison}
%///////////////////////////////////////////////////////////////////////////////////////////////////////////////////////////
\label{subsubsecZQnBcgo}

Pour les classes de conjugaison du groupe diédral nous suivons \cite{HRIMAJJ}.

D'abord pour des raisons de déterminants\footnote{Vous notez qu'ici nous utilisons un argument qui utilise la définition de \( D_n\) comme isométries de \( \eR^2\). Si nous avions voulu à tout prix nous limiter à la définition «abstraite» en termes de générateurs, il aurait fallu trouver autre chose.}, les classes des éléments de la forme \( r^k\) et de la forme \( sr^k\) ne se mélangent pas. Nous notons \( C(x)\) la classe de conjugaison de \( x\), et \( y\cdot x=yxy^{-1}\).

Les relations que nous allons utiliser sont 
\begin{subequations}
    \begin{align}
        sr^ks=r^{-k}\\
        rs=sr^{-1}=sr^{n-1}.
    \end{align}
\end{subequations}

La classe de conjugaison qui ne rate jamais est bien entendu \( C(1)={1}\). Nous commençons les vraies festivités \( C(r^{m})\). D'abord \( r^k\cdot r^m=r^m\), ensuite
\begin{equation}
    (sr^k)\cdot r^m=sr^kr^mr^{-k}s^{-1}=sr^ms^{-1}=r^{-m}.
\end{equation}
Donc
\begin{equation}    \label{EqVFfFxgi}
    C(r^m)=\{ r^m,r^{-m} \}.
\end{equation}
À ce niveau il faut faire deux remarques. D'abord si \( m>\frac{ n }{2}\), alors \( C(r^m)\) est la classe de \( C^{n-m}\) avec \( n-m<\frac{ n }{2}\). Donc les classes que nous avons trouvées sont uniquement à lister avec \( m<\frac{ n }{2}\). Ensuite si \( m=\frac{ n }{2}\) alors \( r^m=r^{-m}\) et la classe est un singleton. Cela n'arrive que si \( n\) est pair.

Nous passons ensuite à \( C(s)\). Nous avons
\begin{equation}
    r^k\cdot s=r^ksr^{-k}=ssr^ksr^{-k}=sr^{-k}r^{-k}=sr^{n-2k},
\end{equation}
et
\begin{equation}
    (sr^k)\cdot s=\underbrace{sr^ks}_{r^{-k}}r^{-k}s^{-1}=r^{-2k}s=r^{n-2k}s=sr^{(n-1)(n-2k)}=sr^{n^2-2kn-n+2k}=sr^{2k}.
\end{equation}
donc
\begin{equation}
    C(s)=\{ sr^{n-2k},sr^{2k} \}_{k=0,\ldots, n-1}.
\end{equation}
Ici aussi l'écriture n'est pas optimale : peut-être que pour certains \( k\) il y a des doublons. Nous reportons l'écriture exacte à la discussion plus bas qui distinguera \( n\) pair de \( n\) impair. Notons juste que si \( n\) est pair, l'élément \( sr\) n'est pas dans la classe \( C(s)\).

Nous en faisons donc à présent le calcul en gardant en tête le fait qu'il n'a de sens que si \( n\) est pair. D'abord
\begin{equation}
    s\cdot (sr)=ssrs=rs=sr^{n-1}.
\end{equation}
Ensuite
\begin{equation}
    (sr^k)\cdot (sr)=sr^ksrr^{-k}s=r^{-2k+1}s=sr^{2k-1}.
\end{equation}
Avec \( k=\frac{ n }{2}\), cela rend \( s\cdot (sr)\), donc pas besoin de le recopier. Nous avons
\begin{equation}
    C(sr)=\{ sr^{2k-1} \}_{k=1,\ldots, n-1}.
\end{equation}

%///////////////////////////////////////////////////////////////////////////////////////////////////////////////////////////
\subsubsection{Le compte pour $ n$ pair}
%///////////////////////////////////////////////////////////////////////////////////////////////////////////////////////////
\label{SubsubsecROVmHuM}

Si \( n\) est pair, nous avons les classes
\begin{subequations}
    \begin{align}
        C(1)&=\{ 1 \}       &&&\text{\( 1\) élément}\\
        C(r^m)&=\{ r^m,r^{m-1} \}&\text{ pour }&0<m<\frac{ n }{2}   &\text{\( \frac{ n }{2}-1\) fois \( 2\) éléments}\\
        C(r^{n/2})&=\{ r^{n/2} \}   &&& \text{\( 1\) élément}\\ 
        C(s)&=\{ sr^{2k} \}_{k=0,\ldots, \frac{ n }{2}-1} &&& \text{\( \frac{ n }{2}\) éléments}\\
        C(sr)&=\{ sr^{2k+1} \}_{k=0,\ldots, \frac{ n }{2}-1} &&& \text{\( \frac{ n }{2}\) éléments}.
    \end{align}
\end{subequations}
Au total nous avons bien listé \( 2n\) éléments comme il se doit, dans \(  \frac{ n }{2}+3\) classes différentes.

%///////////////////////////////////////////////////////////////////////////////////////////////////////////////////////////
\subsubsection{Le compte pour $ n$ impair}
%///////////////////////////////////////////////////////////////////////////////////////////////////////////////////////////
\label{Subsubsec*GJIzDEP}

Si \( n\) est impair, nous avons les classes
\begin{subequations}
    \begin{align}
        C(1)&=\{ 1 \}       &&&\text{\( 1\) élément}\\
        C(r^m)&=\{ r^m,r^{m-1} \}&\text{ pour }&0<m<\frac{ n-1 }{2}   &\text{\( \frac{ n-1 }{2}\) fois \( 2\) éléments}\\
        C(s)&=\{ sr^k \}_{k=0,\ldots, n-1} &&& \text{\( n\) éléments}
    \end{align}
\end{subequations}
Au total nous avons bien listé \( 2n\) éléments comme il se doit, dans \(  \frac{ n+3 }{2}\) classes différentes.


\chapter{Espaces de Hilbert}
% This is part of Mes notes de mathématique
% Copyright (c) 2011-2013
%   Laurent Claessens
% See the file fdl-1.3.txt for copying conditions.

%+++++++++++++++++++++++++++++++++++++++++++++++++++++++++++++++++++++++++++++++++++++++++++++++++++++++++++++++++++++++++++
\section{Sommes de familles infinies}
%+++++++++++++++++++++++++++++++++++++++++++++++++++++++++++++++++++++++++++++++++++++++++++++++++++++++++++++++++++++++++++

%---------------------------------------------------------------------------------------------------------------------------
\subsection{Convergence commutative}
%---------------------------------------------------------------------------------------------------------------------------

\begin{definition}
    Soit \( x_k\) une suite dans un espace vectoriel normé \( E\). Nous disons que la suite \defe{converge commutativement}{convergence!commutative} vers \( x\in E\) si \( \lim_{n\to \infty}\| x_n-x \| =0\) et si pour toute bijection \( \tau\colon \eN\to \eN\) nous avons aussi
    \begin{equation}
        \lim_{n\to \infty} \| x_{\tau(k)}-x \|=0.
    \end{equation}
    La notion de convergence commutative est surtout intéressante pour les séries. La somme
    \begin{equation}
        \sum_{k=0}^{\infty}x_k
    \end{equation}
    converge commutativement vers \( x\) si \( \lim_{N\to \infty} \| x-\sum_{k=0}^Nx_k \|=0\) et si pour toute bijection \( \tau\colon \eN\to \eN\) nous avons
    \begin{equation}
        \lim_{N\to \infty} \| x-\sum_{k=0}^Nx_{\tau(k)} \|=0.
    \end{equation}
\end{definition}

Nous démontrons maintenant qu'une série converge commutativement si et seulement si elle converge absolument.

Pour le sens inverse, nous avons la proposition suivante.
\begin{proposition}
    Soit \( \sum_{k=0}^{\infty}a_k\) une série réelle qui converge mais qui ne converge pas absolument. Alors pour tout \( b\in \eR\), il existe une bijection \( \tau\colon \eN\to \eN\) telle que \( \sum_{i=0}^{\infty}a_{\tau(i)}=b\).
\end{proposition}
Pour une preuve, voir \href{http://gilles.dubois10.free.fr/analyse_reelle/seriescomconv.html}{chez Gilles Dubois}.

\begin{proposition} \label{PopriXWvIY}
    Soit \( (a_i)_{i\in \eN}\) une suite dans \( \eC\) convergent absolument. Alors elle converge commutativement.
\end{proposition}

\begin{proof}
    Soit \( \epsilon>0\). Nous posons \( \sum_{i=0}^\infty a_i=a\) et nous considérons \( N\) tel que
    \begin{equation}
        | \sum_{i=0}^Na_i-a |<\epsilon.
    \end{equation}
    Étant donné que la série des \( | a_i |\) converge, il existe \( N_1\) tel que pour tout \( p,q>N_1\) nous ayons \( \sum_{i=p}^q| a_i |<\epsilon\). Nous considérons maintenant une bijection \( \tau\colon \eN\to \eN \). Prouvons que la série \( \sum_{i=0}^{\infty}| a_{\tau(i)} |\) converge. Nous choisissons \( M\) de telle sorte que pour tout \( n>M\), \( \tau(n)>N_1\); alors si \( p,q>M\) nous avons
    \begin{equation}
        \sum_{i=p}^q| a_{\tau(i)} |<\epsilon.
    \end{equation}
    Par conséquent la somme de la suite \( (a_{\tau(i)})\) converge. Nous devons montrer à présent qu'elle converge vers la même limite que la somme «usuelle» \( \lim_{N\to \infty} \sum_{i=0}^Na_i\).

    Soit \( n>\max\{ M,N \}\). Alors
    \begin{equation}
        \sum_{k=0}^na_{\tau(k)}-\sum_{k=0}^na_k=\sum_{k=0}^Ma_{\tau(k)}-\sum_{k=0}^Na_k+\underbrace{\sum_{M+1}^na_{\tau(k)}}_{<\epsilon}-\underbrace{\sum_{k=N+1}^na_k}_{<\epsilon}.
    \end{equation}
    Par construction les deux derniers termes sont plus petits que \( \epsilon\) parce que \( M\) et \( N\) sont les constantes de Cauchy pour les séries \( \sum a_{\tau(i)}\) et \( \sum a_i\). Afin de traiter les deux premiers termes, quitte à redéfinir \( M\), nous supposons que \( \{ 1,\ldots, N \}\subset \tau\{ 1,\ldots, M \}\); par conséquent tous les \( a_i\) avec \( i<N\) sont atteints par les \( a_{\tau(i)}\) avec \( i<M\). Dans ce cas, les termes qui restent dans la différence
    \begin{equation}
        \sum_{k=0}a_{\tau(k)}-\sum_{k=0}^Na_k
    \end{equation}
    sont des \( a_k\) avec \( k>N\). Cette différence est donc en valeur absolue plus petite que \( \epsilon\), et nous avons en fin de compte que
    \begin{equation}
        \left| \sum_{k=0}^na_{\tau(k)}-\sum_{k=0}^na_k \right| <\epsilon.
    \end{equation}
\end{proof}

\begin{proposition}     \label{PropyFJXpr}
    Soit \( \sum_{i=O}^{\infty}a_i\) une série qui converge mais qui ne converge pas absolument. Pour tout \( b\in \eR\), il existe une bijection \( \tau\colon \eN\to \eN\) telle que \( \sum_{i=}^{\infty}a_{\tau(i)}=b\).
\end{proposition}

Les propositions \ref{PopriXWvIY} et \ref{PropyFJXpr} disent entre autres qu'une série dans \( \eC\) est commutativement sommable si et seulement si elle est absolument sommable.

Soit \( (a_i)_{i\in I}\) une famille de nombres complexes indexée par un ensemble \( I\) quelconque. Nous allons nous intéresser à la somme \( \sum_{i\in I}a_i\).


Soit \( \{ a_i \}_{i\in I}\) des nombres positifs. Nous définissons la somme
\begin{equation}
    \sum_{i\in I}a_i=\sup_{\text{\( J\) fini}}\sum_{j\in J}a_j.
\end{equation}
Notons que cela est une définition qui ne fonctionne bien que pour les sommes de nombres positifs. Si \( a_i=(-1)^i\), alors selon la définition nous aurions \( \sum_i(-1)^i=\infty\). Nous ne voulons évidemment pas un tel résultat.

Dans le cas de familles de nombres réels positifs, nous avons une première définition de la somme. 
\begin{definition}  \label{DefHYgkkA}
Soit \( (a_i)_{i\in I}\) une famille de nombres réels positifs indexés par un ensemble quelconque \( I\). Nous définissons
\begin{equation}
    \sum_{i\in I}a_i=\sup_{\text{\( J\) fini dans \( I\)}}\sum_{j\in J}a_j.
\end{equation}
\end{definition}

\begin{definition}  \label{DefIkoheE}
    Si \( \{ v_i \}_{i\in I}\) est une famille de vecteurs dans un espace vectoriel normé indexée par un ensemble quelconque \( I\). Nous disons que cette famille est \defe{sommable}{famille!sommable} de somme \( v\) si pour tout \( \epsilon>0\), il existe un \( J_0\) fini dans \( I\) tel que pour tout ensemble fini \( K\) tel que \( J_0\subset K\) nous avons
    \begin{equation}
        \| \sum_{j\in K}v_j-v \|<\epsilon.
    \end{equation}
\end{definition}
Notons que cette définition implique la convergence commutative.

\begin{example}
    La suite \( a_i=(-1)^i\) n'est pas sommable parce que quel que soit \( J_0\) fini dans \( \eN\), nous pouvons trouver \( J\) fini contenant \( J_0\) tel que \( \sum_{j\in J}(-1)^j>10\). Pour cela il suffit d'ajouter à \( J_0\) suffisamment de termes pairs. De la même façon en ajoutant des termes impairs, on peut obtenir \( \sum_{j\in J'}(-1)^i<-10\).
\end{example}

\begin{example}
    De temps en temps, la somme peut sortir d'un espace. Si nous considérons l'espace des polynômes \( \mathopen[ 0 , 1 \mathclose]\to \eR\) muni de la norme uniforme, la somme de l'ensemble
    \begin{equation}
        \{ 1,-1,\pm\frac{ x^n }{ n! } \}_{n\in \eN}
    \end{equation}
    est zéro.

    Par contre la somme de l'ensemble \( \{ 1,\frac{ x^n }{ n! } \}_{n\in \eN}\) est l'exponentielle qui n'est pas un polynôme.
\end{example}

\begin{example}
    Au sens de la définition \ref{DefIkoheE} la famille
    \begin{equation}
        \frac{ (-1)^n }{ n }
    \end{equation}
    n'est pas sommable. En effet la somme des termes pairs est \( \infty\) alors que la somme des termes impairs est \( -\infty\). Quel que soit \( J_0\in \eN\), nous pouvons concocter, en ajoutant des termes pairs, un \( J\) avec \( J_0\subset J\) tel que \( \sum_{j\in J}(-1)^j/j\) soit arbitrairement grand. En ajoutant des termes négatifs, nous pouvons également rendre \( \sum_{j\in J}(-1)^j/j\) arbitrairement petit.
\end{example}

La définition générale de la somme \ref{DefIkoheE} est compatible avec la définition usuelle dans les cas où cette dernière s'applique.
\begin{proposition}[commutative sommabilité]\label{PropoWHdjw}
    Soit \( I\) un ensemble dénombrable et une bijection \( \tau\colon \eN\to I\). Soit \( (a_i)_{i\in I}\) une famille dans un espace vectoriel normé. Alors
    \begin{equation}
        \sum_{k=0}^{\infty}a_{\tau(k)}=\sum_{i\in I}a_i
    \end{equation}
    dès que le membre de droite existe. Le membre de gauche est définit par la limite usuelle.
\end{proposition}

\begin{proof}
    Nous posons \( a=\sum_{i\in I}a_i\). Soit \( \epsilon>0\) et \( J_0\) comme dans la définition. Nous choisissons
    \begin{equation}
        N>\max_{j\in J_0}\{ \tau^{-1}(j) \}.
    \end{equation}
    En tant que sommes sur des ensembles finis, nous avons l'égalité
    \begin{equation}
        \sum_{k=0}^Na_{\tau(k)}=\sum_{j\in J_0}a_j
    \end{equation}
    où \( J\) est un sous-ensemble de \( I\) contenant \( J_0\). Soit \( J\) fini dans \( I\) tel que \( J_0\subset J\). Nous avons alors
    \begin{equation}
        \| \sum_{k=0}^Na_{\tau(k)}-a \|=\| \sum_{j\in J}a_j-a \|<\epsilon.
    \end{equation}
    Nous avons prouvé que pour tout \( \epsilon\), il existe \( N\) tel que \( n>N\) implique \( \| \sum_{k=0}^na_{\tau(k)}-a\| <\epsilon\).
\end{proof}

\begin{corollary}
    Nous pouvons permuter une somme dénombrable et une fonction continue. C'est à dire que si \( f\) est une fonction continue sur l'espace vectoriel normé \( E\) et \( (a_i)_{i\in I}\) une famille sommable dans \( E\) alors
    \begin{equation}
        f\left( \sum_{i\in I}a_i \right)=\sum_{i\in I}f(a_i).
    \end{equation}
\end{corollary}

\begin{proof}
    En utilisant une bijection \( \tau\) entre \( I\) et \( \eN\) ainsi que le résultat connu à propos des sommes sur \( \eN\), nous avons
    \begin{equation}
        f\left( \sum_{i\in I}a_i \right)=f\left( \sum_{k=0}^{\infty}a_{\tau(k)} \right)=\sum_{k=0}^{\infty}f(a_{\tau(k)})=\sum_{i\in I}f(a_i).
    \end{equation}
\end{proof}

\begin{corollary}
    Si une série de fonctions continues converge normalement, alors la somme est continue.
\end{corollary}

\begin{proof}
    Soit \( I\) dénombrable et considérons une famille de fonctions continues \( (f_n)_{n\in I}\) telles que la famille \( (\| f_i \|_{\infty})_{i\in I}\) soit sommable. Utilisant une bijection entre \( I\) et \( \eN\), le théorème \ref{PropUEMoNF} s'applique.
\end{proof}


La proposition suivante nous enseigne que les sommes infinies peuvent être manipulée de façon usuelle.
\begin{proposition} \label{PropMpBStL}
    Soit \( I\) un ensemble dénombrable. Soient \( (a_i)_{i\in I}\) et \( (b_i)_{i\in I}\), deux familles de réels positifs telles que \( a_i<b_i\) et telles que \( (b_i)\) est sommable. Alors \( (a_i)\) est sommable.

    Si \( (a_i)_{i\in I}\) est une famille de complexes telle que \( (| a_i |)\) est sommable, alors \( (a_i)\) est sommable.
\end{proposition}

%+++++++++++++++++++++++++++++++++++++++++++++++++++++++++++++++++++++++++++++++++++++++++++++++++++++++++++++++++++++++++++
\section{Complétude}
%+++++++++++++++++++++++++++++++++++++++++++++++++++++++++++++++++++++++++++++++++++++++++++++++++++++++++++++++++++++++++++

\begin{proposition}[Inégalité de Minkowski]     \label{PropInegMinkKUpRHg}
    Si \( 1\leq p<\infty\) et si \( f,g\in L^p(\Omega,\tribA,\mu)\) alors
    \begin{enumerate}
        \item
            \( \| f+g \|_p\leq \| f \|_p+\| g_p \|\)
        \item
            Il y a égalité si et seulement si les vecteurs \( f(x)\) et \( g(x)\) sont presque partout colinéaires : il existe \( \alpha,\beta\) tels que \( \alpha f+\beta g=0\) presque partout.
    \end{enumerate}
\end{proposition}

\begin{theorem}
    Pour \( 1\leq p<\infty\), l'espace \( L^p(\Omega,\tribA,\mu)\) est complet.
\end{theorem}

\begin{proof}
    Cette preuve provient de \cite{SuquetFourierProba}, disponible aussi sous forme d'application du \wikipedia{fr}{théorème_de_Riesz-Fischer}{théorème de Riesz-Fischer} sur wikipedia. 
    
    Soit \( (f_n)_{n\in\eN}\) une suite de Cauchy dans \( L^p\). Pour tout \( i\), il existe \( N_i\in\eN\) tel que $\| f_p-f_q \|_p\leq 2^{-i}$ pour tout \( p,q\geq N_i\). Nous considérons la sous suite \( g_i=f_{N_i}\), de telle sorte qu'en particulier
    \begin{equation}    \label{EqJLoDID}
        \|g_i-g_{i-1}\|_p\leq 2^{-i}.
    \end{equation}
    Pour chaque \( j\) nous considérons la somme télescopique
    \begin{equation}
        g_j=g_0+\sum_{i=1}^j(g_i-g_{i-1})
    \end{equation}
    et l'inégalité
    \begin{equation}
        | g_j |\leq | g_0 |+\sum_{i=1}^j| g_i-g_{i-1} |.
    \end{equation}
    Nous allons noter
    \begin{equation}        \label{EqSomPaFPQOWC}
        h_j=| g_0 |+\sum_{i=1}^j| g_i-g_{i-1} |.
    \end{equation}
    La suite de fonctions \( (h_j)\) ainsi définie est une suite croissante de fonctions positive qui converge donc (ponctuellement) vers une fonction \( h\) qui peut éventuellement valoir l'infini en certains points. Par continuité de la fonction \( x\mapsto x^p\) nous avons
    \begin{equation}
        \lim_{j\to \infty} h_j^p=h^p,
    \end{equation}
    puis par le théorème de la convergence monotone (théorème \ref{ThoConvMonFtBoVh}) nous avons
    \begin{equation}
        \lim_{j\to \infty} \int_{\Omega}h_j^pd\mu=\int_{\Omega}h^pd\mu.
    \end{equation}
    Utilisant à présent la continuité de la fonction \( x\mapsto x^{1/p}\) nous trouvons
    \begin{equation}
        \lim_{j\to \infty} \left( \int h_j^p \right)^{1/p}=\left( \int | h |^p \right)^{1/p}.
    \end{equation}
    Nous avons donc déjà montré que
    \begin{equation}
        \lim_{j\to \infty} \| h_j \|_p=\left( \int | h |^p \right)^{1/p}
    \end{equation}
    où, encore une fois, rien ne garantit à ce stade que l'intégrale à droite soit un nombre fini. En utilisant l'inégalité de Minkowski (proposition \ref{PropInegMinkKUpRHg}) et l'inégalité \eqref{EqJLoDID} nous trouvons
    \begin{equation}
        \|h_j\|_p\leq \|g_0\|_p+\sum_{i=1}^j\|g_i-g_{i-1}\|_p\leq \|g_0\|_p+1.
    \end{equation}
    En passant à la limite,
    \begin{equation}
        \left( \int| h |^p \right)^{1/p}=\lim_{j\to \infty}\|h_j\|_p \leq \|g_0\|_p+1<\infty.
    \end{equation}
    Par conséquent \( \int| h |^p\) est finie et
    \begin{equation}    \label{EqgLpdUPOBP}
        h\in L^p(\Omega,\tribA,\mu).
    \end{equation}
    En particulier, l'intégrale \( \int h\) est finie (parce que \( p\geq 1\)) et donc que \( h(x)<\infty\) pour presque tout \( x\in\Omega\).

    Nous savons que \( h(x)\) est la limite des sommes partielles \eqref{EqSomPaFPQOWC}, en particulier la série
    \begin{equation}
        \sum_{j=1}^{\infty}| g_i-g_{i-1} |
    \end{equation}
    converge ponctuellement. En vertu du corollaire \ref{CorCvAbsNormwEZdRc}, la série de terme général \( g_i-g_{i-1}\) converge ponctuellement. La suite \( g_i\) converge donc vers une fonction que nous notons \( g\). Par ailleurs la suite \( g_i\) est dominée par \( h\in L^p\), le théorème de la convergence dominée (théorème \ref{ThoConvDomLebVdhsTf}) implique que
    \begin{equation}
        \lim_{j\to \infty} \|g_j-g\|_p=0.
    \end{equation}
    Nous allons maintenant prouver que \( \lim_{n\to \infty\|f_n-g\|_p} =0\). Soit \( \epsilon>0\). Pour tout \( n\) et \( i\) nous avons
    \begin{equation}
        \|f_n-g\|_p=\|f_n-f_{N_i}+f_{N_i}-g\|_p\leq\|f_n-f_{N_i}\|_p+\|f_{N_i}-g\|_p.
    \end{equation}
    Pour rappel, \( f_{N_i}=g_i\). Si \(i\) et \( n\) sont suffisamment grands nous pouvons obtenir que chacun des deux termes est plus petit que \( \epsilon/2\).

    Il nous reste à prouver que \( g\in L^p(\Omega,\tribA,\mu)\). Nous avons déjà vu (équation \eqref{EqgLpdUPOBP}) que \( h\in L^p\), mais \( | g_i |\leq h^p\), par conséquent  \( g\in L^p\).

    Nous avons donc montré que la suite de Cauchy \( (f_n)\) converge vers une fonction de \( L^p\), ce qui signifie que \( L^p\) est complet.
\end{proof}

%+++++++++++++++++++++++++++++++++++++++++++++++++++++++++++++++++++++++++++++++++++++++++++++++++++++++++++++++++++++++++++
\section{L'espace \texorpdfstring{$L^2$}{$L^2$}}
%+++++++++++++++++++++++++++++++++++++++++++++++++++++++++++++++++++++++++++++++++++++++++++++++++++++++++++++++++++++++++++

Soit \( (\Omega,\tribA,\mu)\) un espace mesuré. Nous considérons l'opération
\begin{equation}    \label{DefProdScalLubrgTj}
    \langle f, g\rangle =\int_{\Omega}f(\omega)\overline{ g(\omega)}d\mu(\omega)
\end{equation}
et la norme associée
\begin{equation}
    \| f \|_2=\sqrt{\langle f, f\rangle }.
\end{equation}
Nous considérons l'ensemble
\begin{equation}
    \mL^2(\Omega,\mu)=\{ f\colon \Omega\to \eR\tq \| f \|_2<\infty \}
\end{equation}
et la relation d'équivalence \( f\sim g\) si et seulement si \( f(x)=g(x)\) pour \( \mu\)-presque tout \( x\). L'espace que nous considérons est
\begin{equation}
    L^2=\mL^2/\sim.
\end{equation}

\begin{lemma}
    La formule \eqref{DefProdScalLubrgTj} définit un produit scalaire sur \( L^2\).
\end{lemma}

\begin{proof}
    Nous devons d'abord montrer que la formule passe au quotient. Pour cela, nous considérons des fonctions \( \alpha\) et \( \beta\) nulles presque partout et nous regardons le produit de \( f_1=f+\alpha\) par \( g_1=g+\beta\) :
    \begin{equation}
        \langle f_1, g_1\rangle =\int fg+\beta f+\alpha g+ \alpha\beta.
    \end{equation}
    Les fonction \( \beta f\), \( \alpha g\) et \( \alpha\beta\) étant nulles presque partout, leur intégrale est nulle et nous avons bien \( \langle f_1, g_1\rangle =\langle f,g \rangle \). Nous pouvons donc considérer le produit sur l'ensemble des classes.

    Pour vérifier que la formule est un produit scalaire, le seul point non évidement est de prouver que \( \langle f, f\rangle =0\) implique \( f=0\). Cela découle du fait que
    \begin{equation}
        \langle f, f\rangle =\int_{\Omega}| f |^2.
    \end{equation}
    La fonction \( x\mapsto | f(x) |^2\) vérifie les hypothèses du lemme \ref{Lemfobnwt}. Par conséquent \( | f(x) |^2\) est presque partout nulle.
\end{proof}

%+++++++++++++++++++++++++++++++++++++++++++++++++++++++++++++++++++++++++++++++++++++++++++++++++++++++++++++++++++++++++++
\section{Espaces de Hilbert}
%+++++++++++++++++++++++++++++++++++++++++++++++++++++++++++++++++++++++++++++++++++++++++++++++++++++++++++++++++++++++++++

\begin{definition}
    Un espace vectoriel muni d'un produit scalaire est une espace \defe{préhilbertien}{préhilbertien}. Si il est complet pour la norme induite par le produit scalaire alors il est de \defe{Hilbert}{Hilbert}.

    Dans les deux cas nous considérons la topologie métrique dérivant du produit scalaire.
\end{definition}

Les espaces de Hilbert sont des cas particuliers des \defe{espaces de Banach}{espace!Banach}\index{Banach!espace} qui sont les espaces vectoriels normés complets pour la topologie de la norme. La différence avec les espaces de Hilbert sont que l'espace de Banach ne demande pas que la norme dérive d'un produit scalaire.

\begin{proposition}     \label{PropTdupIG}
    Si \( \hH\) est un espace de Hilbert réel, alors
    \begin{equation}
        \| x+y \|^2=\| x \|^2+\| y \|^2+2\langle x, y\rangle .
    \end{equation}
    Si \( \hH\) est un espace de Hilbert complexe alors
    \begin{equation}        \label{EqrbBlkK}
        \| x+y \|^2=\| x \|^2+\| y \|^2+2\Reel\langle x, y\rangle .
    \end{equation}
    Dans les deux cas nous avons l'inégalité de Cauchy-Schwarz\index{inégalité!Cauchy-Schwarz}\index{Cauchy-Schwarz} :
    \begin{equation}
        | \langle x, y\rangle  |\leq \| x \|\| y \|.
    \end{equation}
\end{proposition}

%+++++++++++++++++++++++++++++++++++++++++++++++++++++++++++++++++++++++++++++++++++++++++++++++++++++++++++++++++++++++++++
\section{Théorème de la projection}
%+++++++++++++++++++++++++++++++++++++++++++++++++++++++++++++++++++++++++++++++++++++++++++++++++++++++++++++++++++++++++++

Sources : \cite{ProbCOndutetz} et \wikipedia{en}{Hilbert_projection_theorem}{théorème de projection} sur wikipédia.

\begin{theorem}[Projection sur partie fermée convexe]\index{théorème!projection!partie fermée convexe}
    \label{ThoProjOrthuzcYkz}
    Soit \( \hH\) un espace de Hilbert, et \( C\) un sous ensemble fermé convexe de \( \hH\). Pour chaque \( x\in \hH\), il existe un unique \( y\in C\) tel que
    \begin{enumerate}
        \item
            \( \| x-y \|=\inf\{ \| x-z \|\tq z\in C \}\),
        \item
            pour tout \( z\in C\), \( \langle x-y, z\rangle =0\).
    \end{enumerate}
\end{theorem}

\begin{proof}
    Nous nommons \( d\) l'infimum en question. 
    \begin{description}
        \item[Existence] 
    Soit \( (y_n)\) une suite dans \( C\) telle que 
    \begin{equation}
        \lim_{n\to \infty} \| x-y_n \|=\inf\{ \| x-y \|\tq z\in C \}=d.
    \end{equation}
    Nous allons montrer que cette suite peut être choisie de Cauchy. Elle convergera donc dans \( \hH\) parce que ce dernier est complet. Mais \( C\) étant supposé fermé dans \( \hH\), la limite appartiendra à \( C\). Soient \( r,s\in \eN\). D'abord nous avons
    \begin{subequations}
        \begin{align}
            \| y_r-y_r \|^2&=\langle y_r-y_s+x-x, y_r-y_s+x-x\rangle \\
            &=\| y_r-x \|^2+\| y_s-x \|^2-2\langle y_r-x, y_s-x\rangle .
        \end{align}
    \end{subequations}
    Ensuite,
    \begin{subequations}
        \begin{align}
            4\left\| \frac{ y_r+y_s }{2}-x \right\|^2&=\langle y_r+y_s-2x, y_r-y_s-2x\rangle \\
            &=\| y_r-x \|^2+\| y_s-x \|^2+2\langle y_r-x, y_s-x\rangle .
        \end{align}
    \end{subequations}
    Si nous égalisons les valeurs de \( 2\langle y_r-x, y_s-x\rangle \) nous trouvons
    \begin{equation}    \label{EqiqCyUa}
        \| y_r-y_s \|^2=-4\left\| \frac{ y_r+y_s }{2}-x \right\|^2+2\| y_r-x \|^2+2\| y_s-x \|^2.
    \end{equation}
    La distance infimum étant \( d\), nous pouvons choisir \( y_n\) de telle façon à avoir 
    \begin{equation}
        \| y_n-x \|\leq d+\frac{1}{ n }.
    \end{equation}
    D'autre part étant donné que \( C\) est convexe, \( (y_r+y_s)/2\) est dans \( C\) et nous avons
    \begin{equation}
        \left\| \frac{ y_r+y_s }{2}-x \right\| \leq d.
    \end{equation}
    En mettant ces majorations dans \eqref{EqiqCyUa} nous trouvons
    \begin{equation}
        \| y_r-y_s \|^2\leq -4d+2\left( d+\frac{1}{ r } \right)+2\left( d+\frac{1}{ s } \right)=\frac{1}{ r }+\frac{1}{ s }.
    \end{equation}
    La suite \( (y_n)\) est donc de Cauchy et la limite est un élément de \( C\). Prouvons que cet élément \( y\) réalise l'infimum. Pour cela nous avons les inégalités
    \begin{equation}
        d\leq \| x-y \|\leq\| x-y_n \|+\| y_n-y \|.
    \end{equation}
    En prenant le limite \( n\to\infty\) nous trouvons
    \begin{equation}
        d\leq \| x-y \|\leq d.
    \end{equation}
    
        \item[Unicité]
            Si \( y_1\) et \( y_2\) sont deux éléments de \( C\) qui minimisent la distance avec \( x\), nous avons
            \begin{equation}
                \| y_1-y_2 \|^2=-4\left\| \frac{ y_1+y_2j-x }{2} \right\|^2+2\| y_1-x \|^2+2\| y_2-x \|^2\leq -4d+2d+2d=0.
            \end{equation}
            Par conséquent \( y_1=y_2\).

    \end{description}
\end{proof}

\begin{lemma}       \label{LemGheBQz}
    Soit \( C\) une partie convexe et fermée de l'espace de Hilbert \( \hH\). Soit \( x\in \hH\). Un élément \( y\in C\) minimise la distance avec \( x\) si et seulement si
    \begin{equation}        \label{EqQCKubx}
        \Reel\langle x-y,z-y, \rangle \leq 0
    \end{equation}
    pour tout \( z\in C\).
\end{lemma}

\begin{proof}
    Considérons d'abord \( y=\pr_C(x)\). Par convexité de \( C\), si \( z\in C\) nous avons \( (1-t)y+tz\in C\) et par conséquent
    \begin{equation}
        \| x-(1-t)y-tz \|^2\geq \| x-y \|^2.
    \end{equation}
    Si nous développons complètement les deux membres de cette expression en utilisant la formule \eqref{EqrbBlkK} nous tombons sur
    \begin{equation}
        (-2t+t^2)\| y \|^2+2t\Reel\langle x, y\rangle +t^2\| y \|^2-2t\Reel\langle x, z\rangle +2t(1-t)\Reel\langle y, z\rangle \geq 0.
    \end{equation}
    En divisant par \( t\) et en faisant \( t\to 0\) nous trouvons le résultat \( \Reel\langle z-y, x-y\rangle \leq 0\).

    Réciproquement si \( y\) vérifie la relation \eqref{EqQCKubx} alors nous avons pour chaque \( z\in C\) 
    \begin{subequations}
        \begin{align}
            \| x-y \|^2&=\| (x-y)+(y-z) \|^2\\
            &=\| x-y \|^2+\| y-z \|^2+2\Reel\langle x-y,y-z \rangle \\
            \geq&\| x-y \|^2.
        \end{align}
    \end{subequations}
\end{proof}

\begin{proposition}     \label{PropAXJpCe}
    Soit \( C\) une partie convexe et fermée de l'espace de Hilbert \( \hH\). Alors pour tout \( x_1,x_2\in\hH\) nous avons
    \begin{equation}
        \| \pr_C(x_1)-\pr_C(x_2) \|\leq \| x_1-x_2 \|.
    \end{equation}
    En particulier la projection est une application continue.
\end{proposition}

\begin{proof}
    Nous posons \( y_1=\pr_C(x_1)\) et \( y_2=\pr_C(x_2)\). Par le lemme \ref{LemGheBQz} nous avons, pour tout \( z,z'\in C\) les inégalités
    \begin{subequations}
        \begin{align}
            \Reel\langle x_1-y_1, z-y_1\rangle \leq 0\\
            \Reel\langle x_2-y_2, z'-y_2\rangle \leq 0.
        \end{align}
    \end{subequations}
    En prenant \( z=y_2\) et \( z'=y_1\) et en sommant nous trouvons
    \begin{equation}
        \Reel\langle (x_1-y_1)+(y_2-x_2), y_2-y_1\rangle \leq 0.
    \end{equation}
    Nous pouvons maintenant calculer
    \begin{equation}
        \begin{aligned}[]
            \| y_1-y_2 \|^2&=\Reel\| y_1-y_2 \|\\
            &=\Reel\langle y_1-y_2, (y_1-x_1)+x_1-x_2+(x_2-y_2)\rangle \\
            &=\Reel\langle y_1-y_2, x_1-x_2\rangle +\Reel\langle y_2-y_1, (x_1-y_1)+(y_2-x_2)\rangle \\
            &\leq\langle y_1-y_2, x_1-x_2\rangle \\
            &\leq\big\| \langle y_1-y_2, x_1-x_2\rangle  \big\|\\
            &=\leq \| y_1-y_2 \|\| x_1-x_2 \|.
        \end{aligned}
    \end{equation}
    En simplifiant par \( \| y_1-y_2 \|\)\footnote{Si c'est nul, alors la preuve est évidente.} nous trouvons le résultat
    \begin{equation}
        \| y_1-y_2 \|\leq \| x_1-x_2 \|.
    \end{equation}
\end{proof}

\begin{theorem}[Projection orthogonale]\index{projection!orthogonale}   \index{théorème!projection!cas vectoriel}   \label{ThoMXwOjb}
    Soit \( \hH\) un espace de Hilbert et \( K\), un sous-espace vectoriel fermé non réduit à \( \{ 0 \}\) et \( x\in\hH\). L'élément \( \pr_K(x)\) est l'unique élément de \( K\) tel que
    \begin{equation}
        x-y\in K^{\perp}.
    \end{equation}
    De plus l'application \( x\mapsto\pr_K(x)\) est linéaire, continue et de norme \( 1\).
\end{theorem}

    L'élément \( y\) ainsi définit est la \defe{projection orthogonale}{projection!orthogonale} de \( x\) sur \( K\) et sera noté \( \pr_K(f)\).\nomenclature[Y]{\( \pr_K(x)\)}{projection orthogonale de \( x\) sur \( y\)}

\begin{proof}

    La continuité de la projection est donné par la proposition \ref{PropAXJpCe}.

            Soit \( z\) un élément de \( K\) tel que \( \langle z-x, a\rangle =0\) pour tout \( a\in K\). Nous avons
            \begin{subequations}
                \begin{align}
                    \| x-a \|^2&=\| z-x \|^2+\| a-z \|^2+2\underbrace{\langle z-x, a-z\rangle}_{=0}\\
                    &\geq \| z-x \|^2.
                \end{align}
            \end{subequations}
            Le produit scalaire est nul parce que \( a-z\in K\). La distance \( \| z-x \|\) est donc bien la plus petite distance entre \( x\) et les éléments de \( K\).

            Dans l'autre sens, nous supposons que \( y\in K\) minimise la distance à \( x\) dans \( K\). Par hypothèse pour tout \( a\) et pour tout \( \lambda\in \eR\), la différence
            \begin{equation}
                \| (y+\lambda a)-x \|^2-\| y-x \|^2
            \end{equation}
            est positive. En développant les produits scalaires nous trouvons la conditions suivante
            \begin{equation}
                \lambda^2\| a \|^2+2\lambda\langle a, y-x\rangle \geq 0
            \end{equation}
            qui doit être vraie pour tout \( \lambda\in\eR\). En tant que polynôme du second degré en \( \lambda\), cela n'aura pas deux racines réelles distinctes uniquement si \( \langle a, y-x\rangle =0\).
    
            Nous montrons maintenant la linéarité de la projection orthogonale. Soient \( x_1,x_2\in\hH\). L'élément \( y=\pr_Kx_1+\pr_Kx_2\) satisfait à la condition d'orthogonalité : pour tout \( z\in K\),
    \begin{equation}
        \langle x_1+x_2-\pr_Kx_1-\pr_Kx_2, z\rangle =\langle x_1-\pr_Kx_1, z\rangle +\langle x_2-\pr_Kx_2, z\rangle =0.
    \end{equation}
    Étant donné que \( K\) est un sous-espace vectoriel, la condition de minimalité est automatiquement vérifiée (seconde partie du théorème \ref{ThoProjOrthuzcYkz}).

    En ce qui concerne la norme opérateur de \( \pr_K\), la décomposition de \( x\in\hH\) en composantes dans \( K\) et \( K^{\perp}\) est
    \begin{equation}
        x=x+(x-\pr_Kx).
    \end{equation}
    Étant deux parties orthogonales nous avons
    \begin{equation}
        \| \pr_Kx \|^2=\| x \|^2-\| x-\pr_Kx \|^2.
    \end{equation}
    En prenant \( \| x \|=1\) nous trouvons \( \| \pr_Kx \|^2\leq 1\) et par conséquent \( \| \pr_K \|\leq 1\). Mais d'autre part en prenant \( x\in K\) nous avons automatiquement \( \| \pr_K \|\geq 1\).
\end{proof}

\begin{proposition}
    Soit \( \hH=L^2(\Omega,\tribA,\mu)\), \( \tribF\) une sous tribu de \( \tribA\) et \( K\) l'ensemble de fonctions \( \tribF\)-mesurables dans \( L^2(\Omega,\tribA,\mu)\). Si \( f\in L^2(\Omega,\tribA,\mu)\) est positive, alors \( \pr_Kf\) est positive (presque partout).
\end{proposition}

\begin{proof}
    L'ensemble \( A=\{ \pr_Kf<0 \}\) est dans \( \tribF\). En effet 
    \begin{equation}
        A=(\pr_Kf)^{-1}\big( \mathopen] -\infty , 0 \mathclose[\big)
    \end{equation}
    alors que, par construction, \( \pr_Kf\) est \( \tribF\)-mesurable. La conséquent la fonction indicatrice \( \mtu_A\) est \( \tribF\)-mesurable (c'est à dire \( 1_A\in K\)) et nous avons
    \begin{equation}
        0\leq\int_{\Omega}f\mtu_A=\int_{\Omega}\pr_Kf\mtu_A\leq 0.
    \end{equation}
    Étant donné que nous avons supposé \( f\geq 0\) nous avons alors \( \mu(A)=0\). D'où le fait que \( \pr_Kf\) est presque partout positive.
\end{proof}


%+++++++++++++++++++++++++++++++++++++++++++++++++++++++++++++++++++++++++++++++++++++++++++++++++++++++++++++++++++++++++++
\section{Systèmes orthogonaux et bases}
%+++++++++++++++++++++++++++++++++++++++++++++++++++++++++++++++++++++++++++++++++++++++++++++++++++++++++++++++++++++++++++

Sources : \cite{HilbertLi}. Dans cette partie nous noterons \( \eK\) le corps de base de l'espace \( \hH\). Seuls deux cas sont envisagés : \( \eK=\eC\) ou \( \eK=\eR\).

Pour chaque \( x\in\hH\) nous considérons l'application
\begin{equation}
    \begin{aligned}
        \Phi_x\colon \hH&\to \eK \\
        y&\mapsto \langle x, y\rangle . 
    \end{aligned}
\end{equation}

\begin{lemma}       \label{LemjYVcHE}
    L'application \( y\mapsto\Phi_y\) est une isométrie : nous avons \( \| \Phi_y \|=\| y \|\). De plus pour chaque \( y\), l'application \( \Phi_y\) est continue.
\end{lemma}

\begin{proof}
    En utilisant la définition de la norme opérateur et l'inégalité de Cauchy-Schwarz,
    \begin{equation}
        \| \Phi_y \|=\sup_{\| x \| =1}\| \Phi_y(x) \|=\sup| \langle x, y\rangle  |\leq\sup\| x \|\| y \|=\| y \|.
    \end{equation}
    Par conséquent \( \| \Phi_y \|\leq\| y \|\). Mais d'autre part le fait que \( \Phi_y(y)=\| y \|^2\) montre que \( \| \Phi_y \|\geq \| y \|\).

    En ce qui concerne la continuité de \( \Phi_y\), elle est garantie par le fait que c'est une application linéaire bornée via la proposition \ref{PropmEJjLE}.
\end{proof}

%---------------------------------------------------------------------------------------------------------------------------
\subsection{Orthogonalité}
%---------------------------------------------------------------------------------------------------------------------------

Si \( A\) est une partie de \( \hH\) alors nous notons\nomenclature[Y]{\( A^{\perp}\)}{orthogonal d'une partie}
\begin{equation}
    A^{\perp}=\{ y\in\hH\tq y\perp c\forall x\in A \}.
\end{equation}
cet ensemble est l'\defe{orthogonal}{orthogonal} de \( A\).

\begin{proposition}     \label{PropdpaMpH}
    Si \( \hH\) est une préhilbert et si \( A\subset \hH\), alors l'ensemble \( A^{\perp}\) est un sous-espace fermé de \( \hH\).
\end{proposition}

\begin{proof}
    Par définition l'application \( \Phi_x\) est continue pour chaque \( x\in\hH\) et par conséquent l'ensemble \( \ker\Phi_x=\Phi_x^{-1}(\{ 0 \})\) est fermé. L'ensemble
    \begin{equation}
        A^{\perp}=\bigcap_{x\in A}\ker\Phi_x
    \end{equation}
    est donc fermé.    
\end{proof}

\begin{proposition}[\wikipedia{fr}{Supplémentaire_topologique}{wikipédia}]     \label{PropKZDqTR}
    Soit \( V\) un espace vectoriel normé et une décomposition en somme directe \( V=F\oplus G\). Alors les trois conditions suivantes sont équivalentes.
    \begin{enumerate}
        \item
            L'isomorphisme naturel \( F\times G\to v\) est un homéomorphisme.
        \item
            \( F\) et \( G\) sont fermés et la restriction à \( G\) de la projection \( V\to V/F\) est un homéomorphisme.
        \item
            \( F\) et \( G\) sont fermés et la projection de \( E\) sur \( F\) est continue.
    \end{enumerate}
\end{proposition}
Lorsqu'une décomposition en somme directe vérifie la proposition \ref{PropKZDqTR}, nous disons que la décomposition est \defe{topologique}{topologique!somme directe}.

\begin{theorem}     \label{ThowZyaiz}
    Soit \( \hH\) un espace de Hilbert et \( F\) un sous-espace fermé de \( \hH\). Alors
    \begin{enumerate}
        \item
            Nous avons la décomposition
            \begin{equation}        \label{EqxXIvrG}
                \hH=F\oplus F^{\perp}.
            \end{equation}
        \item       \label{ItemThowZyaizii}
            La projection sur \( F\) par rapport à la somme directe \eqref{EqxXIvrG} est la projection \( \pr_F\) du théorème de projection.
        \item
            La décomposition \eqref{EqxXIvrG} est topologique.
    \end{enumerate}
\end{theorem}

\begin{proposition}
    Si \( F\) est un sous-espace de l'espace de Hilbert \( \hH\) alors on a \( F^{\perp\perp}=\bar F\).
\end{proposition}

\begin{proof}
    Nous savons par la proposition \ref{PropdpaMpH} que \( F^{\perp}\) est fermé, par conséquent le théorème \ref{ThowZyaiz} donne la somme directe
    \begin{equation}
        \hH=F^{\perp}\oplus F^{\perp\perp}.
    \end{equation}
    Mais \( \bar F\) étant également fermé nous avons la somme directe
    \begin{equation}
        \hH=\bar F\oplus(\bar F)^{\perp}.
    \end{equation}
    Montrons que \( (\bar F)^{\perp}=F^{\perp}\). En effet si \( x\in F^{\perp}\) et si \( y\in \bar F\), alors il existe une suite \( y_n\) dans \( F\) qui converge vers \( y\). Pour chaque \( n\) nous avons \( \langle x, y_n\rangle =0\) et dons \( \langle x, y\rangle =0\) par continuité du produit scalaire.

    Nous avons donc
    \begin{equation}
        \hH=F^{\perp}\oplus F^{\perp\perp}=\bar F\oplus F^{\perp}.
    \end{equation}
    Mais \( \bar F\subset F^{\perp\perp}\) (prendre une suite). Les espaces \( \bar F\) et \( F^{\perp\perp}\) étant tous deux des supplémentaires de \( F^{\perp}\), nous déduisons qu'ils doivent être égaux.
\end{proof}

\begin{proposition}     \label{PropqiWonByiBmc}
    Soit \( \hH\) un espace de Hilbert et \( F\) un sous-espace vectoriel de \( \hH\). Alors \( F\) est dense si et seulement si \( F^{\perp}=\{ 0 \}\).
\end{proposition}

\begin{proof}
    Nous savons que
    \begin{equation}
        \hH=F^{\perp}\oplus \bar F.
    \end{equation}
    Donc nous avons \( H=\bar F\) si et seulement si \( F^{\perp}=\{ 0 \}\).
\end{proof}
Pour vérifier si un sous-espace vectoriel d'un espace de Hilbert est dense, il suffit donc de montrer que son orthogonal est réduit à zéro.

\begin{lemma}[Théorème fondamental d'approximation \cite{TribuLi}]      \label{LempTBaUw}
    Soit \( \Omega\) un espace mesurable et \( f\colon \Omega\to \mathopen[ 0 , \infty \mathclose]\) une application mesurable. Alors il existe une suite croissante d'applications étagées \( \varphi_n\colon \Omega\to \eR^+\) dont la limite est \( f\).

    De plus si \( f\) est bornée, la convergence est uniforme.
\end{lemma}

\begin{theorem}[\cite{HilbertLi}]       \label{ThoJsBKir}
    Soit \( I\) un intervalle de \( \eR\). L'espace \( C_c(I)\)\nomenclature[Y]{\( C_c(I)\)}{fonctions continues à support compact dans \( I\)} des fonctions continues à support compact sur \( I\) est dense dans \( L^2(I)\).
\end{theorem}

\begin{proof}
    Soit \( g\in L^2(I)\) une fonction telle que \( g\perp f\) pour toute fonction \( f\in C_c(I)\). Nous avons donc
    \begin{equation}
        \langle f, g\rangle =\int_If\bar g=0.
    \end{equation}
    En passant éventuellement aux composantes réelles et imaginaires nous pouvons supposer que les fonctions sont toutes réelles. Nous décomposons \( g\) en parties positives et négatives : \( g=g^+-g^-\). Notre but est de montrer que \( g^+=g^-\), c'est à dire que \( g\) est nulle. La proposition \ref{PropqiWonByiBmc} conclura que \( C_c(I)\) est dense dans \( L^2(I)\).

    Soit un intervalle \( \mathopen[ a , b \mathclose]\subset I\) et une suite croissante de fonctions \( f_n\in C_c(I)\) convergent vers \( \mtu_{\mathopen[ a , b \mathclose]}\). Par hypothèse pour chaque \( n\) nous avons
    \begin{equation}
        \int_If_ng^+=\int_I f_ng^-.
    \end{equation}
    La suite étant croissante, le théorème de la convergence monotone (théorème \ref{ThoConvMonFtBoVh}) s'applique et nous avons
    \begin{equation}
        \lim_{n\to \infty} \int_I f_ng^+=\int_a^bg^+,
    \end{equation}
    de telle sorte que nous ayons, pour tout intervalle \( \mathopen[ a , b \mathclose]\subset I\) l'égalité
    \begin{equation}        \label{EqYlErAM}
        \int_a^bg^+=\int_a^bg^-.
    \end{equation}
    De plus ces intégrales sont finies parce que
    \begin{equation}
        \int_a^b g^+\leq\int_a^b| g |=\int_I| g |\mtu_{\mathopen[ a , b \mathclose]}=\langle | g |, \mtu_{\mathopen[ a , b \mathclose]}\rangle \leq \| g \|_{L^2}\sqrt{b-a}<\infty
    \end{equation}
    par l'inégalité de Cauchy-Schwarz.

    Soit maintenant un ensemble mesurable \( A\subset I\). La fonction caractéristique \( \mtu_A\) est mesurable et il existe une suite croissante de fonctions étagées \( (\varphi_n)\) convergente vers \( a\) par le lemme \ref{LempTBaUw}. À multiples près, les fonctions \( \varphi_n\) sont des sommes de fonctions caractéristiques du type \( \mtu_{\mathopen[ a , b \mathclose]}\), par conséquent, en vertu de \eqref{EqYlErAM} nous avons
    \begin{equation}
        \int_I\varphi_ng^+=\int_I\varphi_ng^-.
    \end{equation}
    Une fois de plus nous pouvons utiliser le théorème de la convergence monotone et obtenir
    \begin{equation}
        \int_Ag^+=\int_A g^-
    \end{equation}
    pour tout ensemble mesurable \( A\subset I\). Si nous notons \( dx\) la mesure de Lebesgue, les mesures \( g^+dx\) et \( g^-dx\) sont par conséquent égales et dominées par \( dx\). Par le corollaire \ref{CorZDkhwS} du théorème de Radon Nikodym, les fonctions \( g^+\) et \( g^-\) sont égales.
\end{proof}
 
%---------------------------------------------------------------------------------------------------------------------------
\subsection{Dual}
%---------------------------------------------------------------------------------------------------------------------------

Le \defe{dual}{dual} de l'espace de Hilbert \( \hH\) est l'ensemble\nomenclature[Y]{\( \hH^*\)}{dual}
\begin{equation}
    H^*=\{ \Phi\colon \hH\to \eK\text{ linéaire et continue} \}.
\end{equation}
Notons que dans le contexte des espaces linéaires nous demandons la continuité des éléments du dual parce qu'elle n'est pas automatique par la linéarité dans les cas de dimensions infinies.

\begin{theorem}[Fréchet-Riesz]\index{théorème!Fréchet-Riesz}        \label{ThoQgTovL}
    L'application
    \begin{equation}
        \begin{aligned}
            J\colon \hH&\to \hH^* \\
            y&\mapsto \Phi_y 
        \end{aligned}
    \end{equation}
    est une bijection isométrique.
\end{theorem}

\begin{proof}
    Nous savons du lemme \ref{LemjYVcHE} que \(J\) est une isométrique. Nous devons seulement montrer que \( J\) est surjective. Soit \( \Phi\in \hH^*\). Par continuité nous savons que \( F=\ker\Phi\) est fermé, donc
    \begin{equation}
        \hH=\ker\Phi\oplus(\ker\Phi)^{\perp}
    \end{equation}
    par le théorème \ref{ThowZyaiz}. Nous considérons une base de \( H\) adaptée à cette décomposition, c'est à dire
    \begin{itemize}
        \item \( (u_s)_{s\in S}\) une base de \( \ker\Phi\),
        \item
            \( (v_t)_{t\in T}\) une base de \( (\ker\Phi)^{\perp}\).
    \end{itemize}
    Il se fait que \( T\) se réduit à un seul élément parce que si \( v\in (\ker\Phi)^{\perp}\), alors \( \Phi(v)\neq 0\), et par conséquent \( \{ \Phi(v) \}\) est déjà une base de \( \Image\Phi=\eK\). Le théorème du rang (théorème \ref{ThoGkkffA}) assure alors que
    \begin{equation}
        \{ v \}\cup\{ u_s \}_{s\in S}
    \end{equation}
    est une base de \( \hH\) avec \( v\in(\ker\Phi)^{\perp}\) et \( u_s\in \ker\Phi\). Nous choisissons \( v\) pour avoir \( \| v \|=1\).

    Nous pouvons maintenant prouver l'existence de \( y\) tel que \( \Phi_y=\Phi\). En effet si nous posons \( y=\overline{ \Phi(v) }v\) nous avons \( \Phi_y=\Phi\), en effet
    \begin{subequations}
        \begin{align}
            \Phi_y(v)=\langle v, y\rangle =\Phi(v)\langle v, v\rangle =\Phi(v)\\
            \Phi_y(u_s)=\langle u_s, y\rangle =0.
        \end{align}
    \end{subequations}
    Par conséquent \( \Phi_y\) et \( \Phi\) coïncident sur une base de \( \hH\).

    En ce qui concerne l'unicité, d'abord si \( \Phi_y=\Phi\) alors nous devons avoir \( y\in(\ker\Phi)^{\perp}\) et par conséquent \( y=\lambda v\) pour un certain \( \lambda\in \eK\). Nous avons alors
    \begin{equation}
        \Phi_y(v)=\bar\lambda\langle v, v\rangle =\bar\lambda.
    \end{equation}
    Pour que cela soit égal à \( \Phi(v)\), nous fixons \( \lambda=\overline{ \Phi(v) }\).
    
\end{proof}
Notons que nous avons réellement utilisé le théorème du rang pour l'unicité. Si nous ne demandions pas l'unicité, alors nous n'avions pas besoin du fait que \( \dim(\ker\Phi)^{\perp}=1\), et nous aurons donc pu parler de formes plus générales à valeurs dans \( \eK^n\).

%---------------------------------------------------------------------------------------------------------------------------
\subsection{Séparabilité}
%---------------------------------------------------------------------------------------------------------------------------

\begin{definition}
    Un espace topologique est \defe{séparable}{séparable} si il possède une partie dénombrable dense.

    Si \( E\) est un espace vectoriel normé nous disons que \( \Delta\) est une partie \defe{totale}{totale}\index{partie!totale} de \( E\) si \( \Span\Delta\) est dense dans \( E\).
\end{definition}

\begin{proposition}
    Un espace vectoriel normé est séparable si et seulement si il possède une partie totale dénombrable.
\end{proposition}

\begin{proof}
    Si \( \Delta\) est une partie dénombrable de \( E\), alors le \( \eQ\)-espace vectoriel \( \Span_{\eQ}\Delta\) est dénombrable, et sa fermeture est la même que celle de \( \Span_{\eK}\Delta\).
\end{proof}

\begin{definition}
    Une famille \( (u_i)_{i\in I}\) d'éléments de \( \hH\) indicée par un ensemble quelconque \( I\) est un \defe{système orthonormé}{orthonormé!système}\index{système orthonormé} si
    \begin{enumerate}
        \item
            \( \| u_i \|=1\) pour tout \( i\in I\),
        \item
            \( u_i\perp u_j\) pour tout \( i\neq j\).
    \end{enumerate}
\end{definition}
Notons que si \( (u_n)_{n\in \eN}\) est un système orthonormé dénombrable alors en utilisant les formules de la proposition \ref{PropTdupIG}, nous avons
\begin{equation}    \label{EqCLQbMy}
    \big\| \sum_{k=1}^n\xi_ku_k \big\|^2=\sum_{k=1}^n| \xi_k |^2.
\end{equation}

% TODO : cohérence. Je devrais poser e_n(t)=exp(int), puis aussi faire attention aux conventions pour les intervalles non 2pi.
% Ceci est à vérifier un peu partout où on parle de série de Fourier.
\begin{example}
    Dans l'ensemble \( L^2(a,b)\) avec \( b-a=L\) l'ensemble des fonctions
    \begin{equation}        \label{EqxNXguH}
        e_n(t)=\frac{1}{ \sqrt{b-a} }\exp(\frac{ 2\pi }{ L }int)
    \end{equation}
    avec \( n\in \eZ\) forme un système orthonormé. En effet
    \begin{equation}
        \langle e_n, e_n\rangle =\frac{1}{ b-a }\int_a^b\exp(\frac{ 2\pi }{ L }int)\exp(-\frac{ 2\pi }{ L }int)=1
    \end{equation}
    et
    \begin{equation}
        \langle e_n, e_m\rangle =\frac{1}{ b-a }\int_a^b e^{\frac{ 2\pi }{ L }(n-m)t}dt=0.
    \end{equation}
    Dans cette intégrale nous utilisons le fait que \( b=a+(b-a)\) pour simplifier les expressions en cours de calcul.

    La famille \eqref{EqxNXguH} est le \defe{système trigonométrique}{système!trigonométrique} de \( L^2(a,b)\). On en parle aussi dans \cite{KuttlerTopInAl}.
\end{example}

\begin{proposition}
    Un système orthonormé est libre.
\end{proposition}

\begin{proof}
    Soit \( (u_i)_{i\in I}\) un système orthonormé de \( \hH\) et une combinaison linéaire \emph{finie} nulle :
    \begin{equation}
        \sum_{k=1}^{n}a_ku_k=0.
    \end{equation}
    Nous développons la somme en utilisant les formules de la proposition \ref{PropTdupIG} :
    \begin{equation}
        0=\| \sum_ka_ku_k \|^2=\sum_k| a_k |^2\| u_k \|^2+2\sum_{k<l}\langle a_ku_k, a_lu_l\rangle .
    \end{equation}
    Le système étant orthonormé, les choses se simplifient en
    \begin{equation}
        \sum_k| a_k |^2=0,
    \end{equation}
    ce qui signifie que \( a_k=0\) pour tout \( k\).
\end{proof}

Avant de continuer nous devons définir comment nous calculons des sommes sur des ensembles quelconques. Si \( I\) est un ensemble et si pour chaque \( i\in I\) nous avons un nombre réel positif \( a_i\), alors nous définissons
\begin{equation}
    \sum_{i\in I}a_i=\sup_{\substack{J\subset I\\\text{ \( J\) fini }}}\sum_{j\in J} a_j.
\end{equation}

\begin{proposition}[Inégalités de Bessel]\index{Bessel!inégalités}\index{inégalités!Bessel} \label{PropHKqVHj}
    Soit \( \hH\) un préhilbert. Si \( (u_i)_{i\in I}\) est un système orthonormé et si \( x\in \hH\), alors
    \begin{equation}
        \sum_{i\in I}\big| \langle x, u_i\rangle  \big|^2\leq \| x \|^2.
    \end{equation}
\end{proposition}

\begin{proof}
    Les éléments de la somme étant des réels positifs, la notion de somme à utiliser est celle de la définition \ref{DefHYgkkA}.

    Posons \( c_i(x)=\langle x, u_i\rangle \). Pour toute partie finie \( J\subset I\) nous avons
    \begin{equation}
        0\leq \big\| x-\sum_{j\in J}c_j(x)u_j \big\|^2=\| x \|^2-2\Reel\sum_{j}\langle x, c_j(x)u_j\rangle +\sum_j| c_j(x) |^2.
    \end{equation}
    Mais en tenant compte du fait que 
    \begin{equation}
        \langle x, c_j(x)u_j\rangle =\overline{  c_j(x)}\langle x, u_j\rangle =| c_j(x) |^2,
    \end{equation}
    nous restons avec
    \begin{equation}    \label{EqvwXWEA}
        \| x-\sum_{j\in J}c_j(x)u_j \|=\| x \|^2-\sum_{j\in J}| c_j(x) |^2.
    \end{equation}
    Finalement,
    \begin{equation}
        \sum_{j\in J}| c_j(x) |^2\leq \| x \|^2.
    \end{equation}
    Ayant cette inégalité pour toute partie finie de \( I\), nous l'avons encore pour le supremum.
\end{proof}

\begin{remark}  \label{RemfdJcQF}
    À mon avis il doit exister un théorème de complétion de base hilbertienne disant que si on a une famille orthonormée, alors elle se prolonge en base. Utilisant cela, nous trouvons une nouvelle démonstration de la proposition \ref{PropHKqVHj} en disant que la somme sur la «partie de base» est plus petite que la somme sur la «base complète».
\end{remark}

%TODO : trouver un tel théorème.

Dans la suite les systèmes orthonormés dénombrables vont jouer un rôle important. Nous allons conventionnellement les indicer par \( \eN\).
\begin{proposition}
    Soit \( \hH\) un préhilbert et un système orthonormé dénombrable \( (u_n)_{n\in \eN}\). Si 
    \begin{equation}        \label{EqitrzXi}
        x=\sum_{n=1}^{\infty}\xi_nu_n,
    \end{equation}
    alors \( \xi_n=\langle x, u_n\rangle \).
\end{proposition}

\begin{proof}
    Nous appliquons \( \Phi_{u_k}\) à l'équation \eqref{EqitrzXi}. Étant donné que le produit scalaire est continu, nous pouvons permuter avec la somme infinie et obtenir
    \begin{equation}
        \langle x, u_k\rangle =\sum_{n=0}^{\infty}\langle u_n, u_k\rangle =\xi_k.
    \end{equation}
\end{proof}

La proposition suivante explique que la notion de projection est compatible avec la décomposition d'un vecteur dans un système orthonormé.
\begin{proposition}
    Soit \( (u_k)_{k\geq 1}\) un système orthonormé d'un préhilbert \( \hH\). Soient 
    \begin{equation}
        x=\sum_{k=1}^{\infty}\xi_ku_k
    \end{equation}
    et
    \begin{equation}
        F=\Span\{ u_1,\ldots, u_n \}.
    \end{equation}
    Alors 
    \begin{equation}
        \pr_F(x)=\sum_{k=1}^n\xi_ku_k.
    \end{equation}
\end{proposition}

\begin{proof}
    Nous allons dans un premier temps montrer que 
    \begin{equation}
        y=x-\sum_{k=1}^n\xi_ku_k
    \end{equation}
    est dans \( F^{\perp}\). Pour cela nous calculons
    \begin{equation}
        \langle x-\sum_{k=1}^n\xi_ku_k,u_j, \rangle =\langle \sum_{k=n+1}^{\infty}\xi_ku_k, u_j\rangle =0
    \end{equation}
    où nous avons utilisé la continuité du produit scalaire pour permuter la somme (infinie) et le produit. Étant donné que \( y\in F^{\perp}\) nous avons \( \pr_Fy=0\) par le point \ref{ItemThowZyaizii} du théorème \ref{ThowZyaiz}.

    D'autre part \( \pr_Fy\) peut être calculé selon
    \begin{equation}
        \pr_Fy=\pr_Fx-\sum_{k=1}^n\xi_k\pr_Fu_k
    \end{equation}
    tandis que \( \pr_Fu_k=u_k\) lorsque \( 1\leq k\leq n\). Par conséquent l'annulation de \( \pr_Fy\) donne
    \begin{equation}
        \pr_Fx=\sum_{k=1}^n\xi_ku_k,
    \end{equation}
    donc le résultat.
\end{proof}

\begin{theorem}[Meillere approximation] \label{ThooRArDp}
    Soit \( \{ u_i \}_{i\in I}\) une famille orthonormé de l'espace de Hilbert \( \hH\). Alors pour tout \( x\in \hH\) et pour tout \( J\) fini dans \( I\) et pour toute famille de nombres complexes  \( (a_j)_{j\in J}\) nous avons
    \begin{equation}
        \| \sum_{j\in J}\langle x, u_j\rangle u_j-x \|\leq \| \sum_{j\in J}a_ju_j-x \|.
    \end{equation}
\end{theorem}
Ce théorème exprime le fait que les nombres \( \langle x, u_i\rangle \) sont les meilleurs coefficients à mettre devant les \( u_i\) pour approximer \( x\).

\begin{corollary}
    Soit \( \{ u_i \}_{i\in I}\) une famille orthonormé de \( \hH\). Pour tout \( x\in \hH\) et pour toutes parties finies \( J,K\) de \( I\) avec \( J\subset K\) nous avons
    \begin{equation}
        \| \sum_{j\in K}\langle x, u_j\rangle u_j-x \|\leq \| \sum_{j\in J}\langle x, u_j\rangle u_j-x \|.
    \end{equation}
\end{corollary}
Ce corollaire exprime le fait que plus on prend de termes de la forme \( \langle , u_i\rangle u_i\), mieux c'est.

\begin{proposition}     \label{PropzaKXHq}
    Soit \( \hH\) un espace de Hilbert et \( (u_n)\) un système orthonormé dans \( \hH\). Si \( (\xi_n)_{n\in \eN}\) est une suite dans \( \ell^2\) alors la série 
    \begin{equation}
        \sum_{n=1}^{\infty}\xi_nu_n
    \end{equation}
    converge dans \( \hH\).

    Autrement dit l'application
    \begin{equation}
        \begin{aligned}
            S\colon \hH&\to \ell^2 \\
            x&\mapsto \big( \langle x, u_n\rangle  \big)_{n\geq 1} 
        \end{aligned}
    \end{equation}
    est surjective.
\end{proposition}

\begin{proof}
    Nous allons montrer que la série \( \sum_{n=1}^{\infty}\xi_nu_n\) est de Cauchy, c'est à dire que la limite
    \begin{equation}
        \lim_{n\to \infty} \big\| \sum_{k=n}^{n+p}\xi_ku_k \big\|=0
    \end{equation}
    est uniforme en \( p\). Cela est un corollaire de la formule \eqref{EqCLQbMy} parce que
    \begin{equation}
        \big\| \sum_{k=n}^{n+p}\xi_ku_k \big\|^2=\sum_{k=n}^{n+p}| x_k |^2.
    \end{equation}
    Mais si \( (\xi_n)\) est dans \( \ell^2\), pour tout \( \epsilon\), il existe \( N\) tel que si \( n>N\) alors le membre de droite est inférieur à \( \epsilon\) indépendamment de \( p\).
\end{proof}

%---------------------------------------------------------------------------------------------------------------------------
\subsection{Bases d'espaces de Hilbert}
%---------------------------------------------------------------------------------------------------------------------------
\label{SubsecDxkjut}

\begin{definition}
    Une \defe{base orthonormée}{base!espace préhilbertien} est une famille dénombrable orthonormé et totale.
\end{definition}
Notons que cela n'est pas la même notion de base qu'en algèbre. En effet pour avoir une base algébrique d'un espace vectoriel, nous demandons que les éléments soient des combinaisons linéaires \emph{finies} des éléments de la base, tandis qu'ici en demandant que la partie soit totale nous demandons simplement que les combinaisons linéaires finies soient denses.

Nous allons voir qu'un espace de Hilbert est généré par les sommes \emph{infinies} de vecteurs d'une base orthonormé avec des coefficients qui forment une suite dans \( \ell^2\).

Étant donné une base hilbertienne de \( \hH\), nous notons
\begin{equation}
    c_k(x)=\langle x, u_k\rangle .
\end{equation}
Dans le théorème suivant (et d'ailleurs partout), les sommes sur \( I\) sont prises au sens de la définition \ref{DefIkoheE}.


\begin{theorem}[Décomposition dans une base orthogonale]     \label{ThoyAjoqP}
    Soit \( \hH\) un espace de Hilbert séparable \( \{ u_i \}_{i\in I}\) une base orthonormée (\( I\) est un ensemble dénombrable quelconque). 

    \begin{enumerate}
        \item
            Pour tout \( x\in \hH\) nous avons
            \begin{equation}        \label{EqPUPTXJ}
                x=\sum_{i\in I}\langle x, u_i\rangle u_i
            \end{equation}
            où la somme est prise au sens de la définition \ref{DefIkoheE}. En particulier, la somme converge de façon commutative.

        \item

            Si \( \{ e_i \}_{i\in I}\) est une famille orthonormée qui satisfait la décomposition \eqref{EqPUPTXJ} pour tout \( x\in \hH\) alors \( \{ e_i \}\) est une base hilbertienne.
        \item
            Nous avons l'identité de \defe{Plancherel}{Plancherel}
            \begin{equation}
                \| x \|^2=\sum_{i\in I}| \langle x, u_i\rangle  |^2.
            \end{equation}
            Le point \ref{ItemQGwoIx} nous indiquera que cette égalité est en fait suffisante pour dire que nous avons une base.
        \item
            Nous avons l'identité de \defe{Parseval}{Parseval}
            \begin{equation}
                \langle x, y\rangle =\sum_{i\in I}\langle x, u_i\rangle \overline{ \langle y, u_i\rangle  }.
            \end{equation}
        \item   \label{ItemQGwoIx}
            Si \( \{ e_i \}\) est une famille de vecteurs unitaires vérifiant l'identité de Plancherel, alors c'est une base hilbertienne.
        \item
            Si \( \{ u_i \}_{i\in I}\) est une base hilbertienne, la suite \( n\mapsto| \langle x, e_n\rangle  |\) appartient à \( \ell^2(I)\).
    \end{enumerate}
\end{theorem}

\begin{proof}
    \begin{enumerate}
        \item
            Étant donné que le système \( \{ u_i \}_{i\in I}\) est total, nous pouvons considérer une suite de combinaisons linéaires finies des \( u_i\) qui converge vers \( x\). Nous écrivons
            \begin{equation}
                x_n=\sum_{x\in J_n}a_{n,k}u_j
            \end{equation}
            et \( x_n\to x\) dans \( \hH\). Les ensembles \( J_n\) sont des sous-ensembles finis de \( I\). Nous pouvons les choisir de telle sorte que \( J_n\subset J_{n+1}\) et \( \bigcup_{n\in \eN}J_n=I\). Ce choix correspond à éventuellement prendre \( a_{n,j}=0\) pour toutes les valeurs de \( j\) «en trop».

            Soit \( \epsilon>0\) et \( N\) tel que \( \| x_n-x \|<\epsilon\) pour tout \( n\geq N\). Nous allons montrer que pour tout \( J\) fini tel que \( J_N\subset J\) nous avons \( \| \sum_{j\in J}\langle x, u_j\rangle u_j-x \|<\epsilon\). Étant donné que
            \begin{equation}
                \sum_{j\in J_N}\langle x, u_j\rangle u_j=\sum_{j\in J}a_ju_j
            \end{equation}
            avec
            \begin{equation}
                a_j=\begin{cases}
                    \langle x, u_j\rangle     &   \text{si \( j\in J_N\)}\\
                    0    &    \text{sinon},
                \end{cases}
            \end{equation}
            le théorème de meilleure approximation \ref{ThooRArDp} nous enseigne que
            \begin{equation}
                \| \sum_{j\in J}\langle x, u_j\rangle u_j-x \|\leq\| \sum_{j\in J_N}\langle x, u_j\rangle u_j-x \|<\epsilon
            \end{equation}
            par conséquent la somme \( \sum_{i\in I}\langle x, u_i\rangle u_i\) converge vers \( x\) au sens général, et en particulier commutativement.

    Les sommes étant commutatives (en particulier \( x=\sum_{i\in I}\langle x, u_i\rangle u_i\)), et les bases hilbertiennes étant dénombrables, nous ne perdons aucune généralité en ne considérant que des bases indexées par \( \eN\).

        \item
            Nous devons montrer que l'ensemble des combinaisons linéaires finies est dense dans \( \hH\). Par hypothèse, pour tout \( \epsilon\), il existe un ensemble fini \( J\) tel que
            \begin{equation}
                \| \sum_{j\in J}\langle x, u_j\rangle u_j-x \|<\epsilon.
            \end{equation}
            Cela prouve la densité dont nous avions besoin.
        \item
            La norme étant une fonction continue, elle commute avec les sommes infinies, de telle sorte que l'égalité de Plancherel donne
            \begin{equation}
                \| x \|^2=\| \sum_{i\in I}\langle x, u_i\rangle u_i \|^2.
            \end{equation}
            Le système des \( u_i\) étant orthonormé,
            \begin{equation}
                \| x \|^2=\| \sum_{i\in I}\langle x, u_i\rangle u_i \|^2=\sum_{i\in I}| \langle x, u_i\rangle  |^2.
            \end{equation}
        \item
            Au tour de \wikipedia{fr}{Perceval}{Parseval}. Nous commençons par prouver que la somme du membre de droite converge. En utilisant l'inégalité \( | zz' |\leq | z |^2+| z' |^2\) (valable pour \( z,z'\in \eC\)) et Plancherel, nous avons
            \begin{equation}
                \begin{aligned}[]
                    \sum_{i\in I}| \langle x, u_i\rangle \overline{ \langle y, u_i\rangle  } |&\leq\sum_{i\in I}| \langle x, u_i\rangle  |^2+| \langle y, u_i\rangle  |^2\\
                    &\leq \| x \|^2+\| y \|^2.
                \end{aligned}
            \end{equation}
            Nous en déduisons que la famille \( \langle x, u_i\rangle \overline{ \langle y, u_i\rangle  }\) est (commutativement) sommable en utilisant la proposition \ref{PropMpBStL}. Par ailleurs nous savons que
            \begin{subequations}
                \begin{align}
                    x=\sum_{i\in I}c_i(x)u_i\\
                    y=\sum_{i\in I}c_i(y)u_i,
                \end{align}
            \end{subequations}
           et le produit scalaire étant une forme bilinéaire continue,
           \begin{subequations}
               \begin{align}
                   \langle x, y\rangle &=\sum_{i\in I}\sum_{j\in I}\langle c_i(x)u_i, \overline{ c_j(y)u_j }\rangle \\
                   &=\sum_{i\in I}\sum_{j\in J}c_i(x)\overline{ c_j(y) }\delta_{ij}\\
                   &=\sum_{i\in I}c_i(x)\overline{ c_j(y) }.
               \end{align}
           \end{subequations}
       \item
            Nous utilisons l'égalité de Plancherel avec \( x=u_j\) :
            \begin{equation}
                \| u_j \|^2=\| u_j \|+\sum_{i\in I\setminus\{ j \}}| \langle u_j, u_i\rangle  |^2.
            \end{equation}
            Par conséquent \( \langle u_j, u_i\rangle =0\) dès que \( i\neq j\). Cela prouve que le système \( \{ u_i \}_{i\in I}\) est orthonormé.

            Nous devons encore prouver que le système est total. Pour cela nous repartons de l'équation \eqref{EqvwXWEA} que nous avions déduites dans la démonstration de l'inégalité de Bessel :
    \begin{equation}
        \| x-\sum_{j\in J}c_j(x)u_j \|=\| x \|^2-\sum_{j\in J}| c_j(x) |^2.
    \end{equation}
    Par hypothèse le membre de droite peut être rendu aussi petit que l'on veut en prenant \( J\) grand (mais fini) dans \( I\). Le membre de gauche indique alors que le système \( \{ u_i \}_{i\in I}\) est total.

\item
    L'identité de Plancherel signifie entre autres que si \( x\in \hH\) alors \( \sum_{i\in I}| \langle x, u_i\rangle  |^2\) converge. Du coup la suite \( (\langle x, u_i\rangle )_{i\in I}\) est dans \( \ell^2(I)\).
    \end{enumerate}
\end{proof}

    
\begin{remark}
    Nous avons décidé d'indexer les bases hilbertiennes par \( \eN\); cela est légitime parce que les sommes sont commutatives. Il ne faut cependant pas perdre de vue qu'en pratique l'ensemble naturel avec lequel on indexe une base est parfois \( \eZ\). Un tel cas est donné par la base trigonométrique de \( L^2\). Indexer cette dernière par \( \eN\) plutôt que par \( \eZ\) serait une contorsion inutile.
\end{remark}

\begin{remark}
    L'égalité de Parseval est la raison pour laquelle les physiciens écrivent souvent
    \begin{equation}
        \id=\sum_{n=1}^{\infty}|u_n\rangle\langle u_n|
    \end{equation}
    dans les livres de mécanique quantique par exemple. Dans \cite{SchomblondQEDQCD}, nous lisons même 
    \begin{equation}
        \int_{-\infty}^{+\infty}dq|q\rangle\langle q| =\hat I.
    \end{equation}
    Notons que ces personnes travaillent avec un espace de Hilbert dont la base n'est pas dénombrable. Pour dire que la physique, ça n'utilise pas des mathématiques pour rire !
\end{remark}

% TODO : cette remarque devrait être plus haut.
\begin{remark}
    Par définition une base orthonormée est donc une partie dénombrable dont l'espace vectoriel engendré est dense. Un espace de Hilbert possédant une base orthonormée est donc séparable. C'est ce fait qui nous pousse à ne considérer que des espaces de Hilbert séparables; nous n'allons donc pas étudier ce qu'il se passerait par exemple en considérant l'espace vectoriel librement engendré par les éléments de \( \eR\).
\end{remark}

\begin{corollary}
    Soit \( \hH\) un espace de Hilbert et \( (u_n)\) une base orthonormée. L'application
    \begin{equation}
        \begin{aligned}
            S\colon \hH&\to \ell^2 \\
            x&\mapsto \big( \langle x, u_n\rangle  \big)_{n\in \eN} 
        \end{aligned}
    \end{equation}
    est un isomorphisme d'espaces de Hilbert.

    De plus l'isomorphisme réciproque est
    \begin{equation}
        \begin{aligned}
            S^{-1}\colon \ell^2&\to \hH \\
            (\xi_n)&\mapsto \sum_{n=1}^{\infty}\xi_nu_n. 
        \end{aligned}
    \end{equation}
\end{corollary}

\begin{proof}
    Nous devons prouver que l'application est bijective et qu'elle vérifie
    \begin{equation}
        S(x)\cdot S(y)=\langle x, y\rangle 
    \end{equation}
    où le point dénote le produit dans \( \ell^2\).

    Pour la surjectivité, si \( (\xi_n)\in \ell^2\) alors nous savons que la somme \( \sum_n\xi_nu_n\) converge par la proposition \ref{PropzaKXHq} et par conséquent \( (\xi_n)\) est l'image par \( S\) de ce vecteur de \( \hH\).

    Pour l'injectivité, si \( S(x)=S(y)\) alors
    \begin{equation}
        x=\sum_n\langle x, u_n\rangle u_n=\sum_n\langle y, u_n\rangle u_n=y
    \end{equation}
    en utilisant la décomposition \eqref{EqPUPTXJ}.

    Le fait que \( S\) soit une isométrie est contenu dans Parseval.
\end{proof}

\begin{proposition}
    Tout espace de Hilbert séparable possède des bases orthonormées.
\end{proposition}

\begin{proof}
    Soit \( (v_n)_{n\in \eN}\) une partie totale et séparable. Quitte à supprimer les \( v_i\) qui sont combinaisons linéaires des précédents, nous pouvons supposer que cette partie est libre. Nous considérons l'espace vectoriel
    \begin{equation}
        F_n=\Span\{ v_1,\ldots, v_n \}.
    \end{equation}
    Sur \( F_n\) nous pouvons appliquer un procédé de Gram-Schmidt pour construire une base orthonormée \( \{ u_1,\ldots, u_n \}\) de \( F_n\) au sens usuel. En considérant \( F_{n+1}\) et en recommençant, les vecteurs \( u_1,\ldots, u_n\) ne changent pas, mais nous obtenons un vecteur \( u_{n+1}\).

    Nous construisons ainsi une suite \( (u_n)\) qui est alors orthonormée au sens des espaces de Hilbert. Nous devons encore prouver qu'il s'agit d'un ensemble total. Cela est simplement dû au fait que tout élément de \( \Span\{ v_n \}\) est contenu dans \( \Span\{ u_n \}\) parce que \( \Span\) ne considère que des combinaisons linéaires finies.
\end{proof}

%---------------------------------------------------------------------------------------------------------------------------
\subsection{Digression sur les normes opérateurs}
%---------------------------------------------------------------------------------------------------------------------------
\label{subsecaeSywF}

Le théorème \ref{ThoyAjoqP} nous indique que si \( \{ u_i \}_{i\in I}\) est une base hilbertienne, alors pour tout \( x\in \hH\) nous avons
\begin{equation}
    \pr_{u_i}(x)=\langle x, u_i\rangle u_i,
\end{equation}
et donc
\begin{equation}
    \sum_{i\in I}\pr_{u_i}x=x.
\end{equation}
Nous ne pouvons cependant pas conclure que
\begin{equation}    \label{EqvKDzNl}
    \sum_{i\in I}\pr_{u_i}=\id
\end{equation}
au sens de la norme opérateur de la section \ref{SeckwyQjK}. En effet en prenant \( I=\eN\), l'égalité \eqref{EqvKDzNl} demanderait d'avoir
\begin{equation}
    \lim_{N\to \infty} \left\| \sum_{i=1}^N\pr_{u_i}-\id \right\|_{\infty}=0,
\end{equation}
or pour tout \( N\), le vecteur \( u_{N+1}\) réalise
\begin{equation}
    \sum_{i=1}^N\pr_{u_i}(u_{N+1})-u_{N+1}=-u_{N+1}.
\end{equation}
Par conséquent pour tout \( N\) nous avons
\begin{equation}
    \sup_{\| x \|=1}\left\| \sum_{i=1}^N\pr_{u_i}x-x \right\|\geq 1.
\end{equation}
Nous ne pouvons donc pas dire
\begin{equation}
    \sum_{n=1}^{\infty}\pr_{u_i}=\id
\end{equation}
au sens de la norme opérateur \eqref{EqThUCEJ}. 

Nous avons cependant la convergence au sens faible.
\begin{proposition}
    Soit \( \hH\) un espace de Hilbert et \( \{ u_i \}_{i\in \eN}\) une base hilbertienne de \( \hH\). Au sens de la topologie faible sur l'espace des opérateurs nous avons
    \begin{equation}
        \sum_{i=1}^{\infty}\pr_{u_i}=\id.
    \end{equation}
\end{proposition}

\begin{proof}
    Pour chaque \( N\in \eN\) et \( x\in\hH\), en vertu de la décomposition \eqref{EqPUPTXJ} nous avons
    \begin{equation}
        \sum_{i=1}^N\pr_{u_i}(x)-x=\sim_{i=N+1}^{\infty}\langle x, u_i\rangle u_i.
    \end{equation}
    Par l'orthonormalité de la base nous avons
    \begin{equation}
        \sum_{i=N+1}^{\infty}\| \langle x, u_i\rangle u_i \|=\sum_{i=N+1}^{\infty}| \langle x, u_i\rangle  |,
    \end{equation}
    dont la limite \( N\to \infty\) est zéro étant donné que la suite \( i\mapsto| \langle x, u_i\rangle  |\) est dans \( \ell^2(\eR)\) par le théorème \ref{ThoyAjoqP}.
\end{proof}

%+++++++++++++++++++++++++++++++++++++++++++++++++++++++++++++++++++++++++++++++++++++++++++++++++++++++++++++++++++++++++++
\section{Théorème de Kochen-Specker}
%+++++++++++++++++++++++++++++++++++++++++++++++++++++++++++++++++++++++++++++++++++++++++++++++++++++++++++++++++++++++++++

Le théorème suivant est \wikipedia{fr}{Variable_cachée}{central en mécanique quantique}. La démonstration provient de \cite{BainKochen} et de \wikipedia{en}{Kochen-Specker_theorem}{Wikipédia}. Nous allons démontrer complètement le théorème seulement pour les espaces de Hilbert de dimension plus grande ou égale à \( 4\).

\begin{theorem}[Kochen-Specker]
    Soit \( \hH\) un espace de Hilbert de dimension plus grande ou égale à \( 3\). Une fonction \( v\) sur l'ensemble des opérateurs de \( \hH\) ne peut pas satisfaire aux conditions suivantes :
    \begin{enumerate}
        \item
            Si \( A\) et \( B\) sont compatibles, alors \( v(A+B)=v(A)+v(B)\),
        \item
            Si \( A\) et \( B\) sont compatibles, alors \( v(AB)=v(A)v(B)\).
    \end{enumerate}
    Ici nous disons que deux opérateurs sont \defe{compatibles}{opérateurs!compatibles} lorsqu'ils possèdent une base hilbertienne commune de vecteurs propres.
\end{theorem}

\begin{proof}
    Soit \( \{ u_n \}_{n\in \eN}\) une base hilbertienne de \( \hH\). Nous notons \( \pr_i\) l'opérateur de projection sur l'espace (fermé) engendré par \( u_i\). Ce sont des opérateurs compatibles deux à deux parce que la base \( \{ u_n \}_{n\in \eN}\) est une base commune de vecteurs propres\footnote{Pour les besoins de la physique, nous remarquons que ces opérateurs sont des opérateurs hermitiens qui commutent, mais ça ne joue pas ici.}.

    D'abord nous devons avoir \( v(\mtu)=1\). En effet pour tout opérateur \( A\), nous avons
    \begin{equation}
        v(A)=v(A\mtu)=v(A)v(\mtu).
    \end{equation}
    Pour peu que \( v(A)\neq 0\), cela nous fait \( v(\mtu)=1\).

    En vertu du théorème \ref{ThoyAjoqP}, un vecteur \( x\in\hH\) se décompose en \( x=\sum_n\langle x, u_n\rangle u_n\), et nous avons
    \begin{equation}    \label{EqFFuXvq}
        \pr_ix=\langle x, u_i\rangle u_i.
    \end{equation}
    En effet le théorème de la projection orthogonale \ref{ThoMXwOjb} nous enseigne que \( \pr_ix\) serait l'unique vecteur de la forme \( \lambda u_i\) tel que \( \lambda u_i-x\perp u_i\). Il est facile de vérifier que le vecteur proposé par \eqref{EqFFuXvq} vérifie cette propriété.

    Une conséquence est que
    \begin{equation}
        \left( \sum_{i=0}^{\infty}\pr_i \right)(x)=x.
    \end{equation}
    Par conséquent, par hypothèse du théorème nous devons avoir
    \begin{equation}    \label{EqUeZHJH}
        \sum_iv(\pr_i)=v(\mtu)=1.
    \end{equation}
    Étant donné que les projections sont idempotentes,
    \begin{equation}
        v(\pr_i)=v(\pr_i^2)=v(\pr_i)^2
    \end{equation}
    et donc \( v(\pr_i)\) doit valoir zéro ou un. Mais la relation \eqref{EqUeZHJH} donne une forte contrainte sur le choix de \( 0\) et de \( 1\). En effet, parmi les \( v(\pr_i)\), un et un seul doit valoir \( 1\), les autres doivent valoir \( 0\).

    Refaisant le raisonnement pour une autre base orthonormale hilbertienne, nous trouvons que les valeurs de \( v\) sur les opérateurs de projection sur les différentes directions doivent être choisies de telle façon à ce que tout choix de base hilbertienne orthogonale contienne exactement un \( 1\) et les reste de zéros.

    Nous voudrions maintenant insister sur un point. Le problème de déterminer de façon cohérente les valeurs \( 0\) ou \( 1\) pour tous les \( v(\pr_i)\) revient à attacher \( 0\) ou \( 1\) à tous les \emph{rayons} de \( \hH\) de façon à ce que toute base orthogonale de \( \hH\) contienne exactement un \( 1\). Un rayon est une direction, c'est à dire une classe d'équivalence \( x\sim \lambda x\). Si nous décidons de nommer «blanc»  les rayons attachés à la valeur \( 0\) et «noirs»  ceux attachés à la valeur \( 1\), le problème se réduit à colorer la boule unité de façon compatible.

    Soit \( \{ u_i \}_{i\in \eN}\) une base orthogonale de \( \hH\) numérotée de telle sorte que \( v(u_0)=1\) et \( v(u_k)=0\) pour \( k\neq 0\). Nous allons maintenant nous particulariser au cas de dimension supérieure ou égale à \( 4\). Si \( R\) est une rotation dans le plan \( \Span\{ u_0,u_1,u_2,u_3 \}\), alors l'ensemble
    \begin{equation}
        \{ Ru_0,Ru_1,Ru_2,Ru_3,u_k \}_{k\geq 3}
    \end{equation}
    est encore une base orthogonale de \( \hH\) et nous avons encore \( v(u_k)=0\) pour \( k\geq 4\). Par conséquent un et un seul des vecteurs \( Ru_0\), \( Ru_1\), \( Ru_2\) ou \( Ru_3\) est colorié en noir; les trois autres étant blancs. Le problème est maintenant complètement réduit à la dimension \( 4\). Note : pour réduire à la dimension \( 3\), on procède de même, mais pour conclure, il faut travailler plus.

    Nous allons construire \( 9\) base orthogonales de \( \eR^4\) à partir de \( 18\) vecteurs, chacun arrivant dans exactement deux des bases. Ils sont donnés dans le tableau suivant\ :


    \begin{tabular}{c|c|c|c|c|c|c|c|c}
    $\begin{matrix}
        0    &   1    \\ 
        0    &   0    
    \end{matrix}$
    &
    $\begin{matrix}
        0    &   0    \\ 
        1    &   0
    \end{matrix}$
    &
    $\begin{matrix}
        0    &   0    \\ 
        0    &   1    
    \end{matrix}$
    &
    $\begin{matrix}
        1    &   1    \\ 
        1    &   1    
    \end{matrix}$
    &
    $\begin{matrix}
        1    &   -1    \\ 
        1    &   -1    
    \end{matrix}$
    &
    $\begin{matrix}
        1    &   -1    \\ 
        -1    &   1    
    \end{matrix}$
    &
    $\begin{matrix}
        -1    &   1    \\ 
        1    &   1    
    \end{matrix}$
    &
    $\begin{matrix}
        1    &   1    \\ 
        -1    &   1    
    \end{matrix}$
    &
    $\begin{matrix}
        1    &   1    \\ 
        1    &   -1    
    \end{matrix}$\\
    \hline      
$\begin{matrix}
        0    &   0    \\ 
        0    &   1    
    \end{matrix}$
    &
    $\begin{matrix}
        0    &   1    \\ 
        0    &   0    
    \end{matrix}$
    &
    $\begin{matrix}
        0    &   0    \\ 
        1    &   0    
    \end{matrix}$
    &
    $\begin{matrix}
        1    &   -1    \\ 
        1    &   -1    
    \end{matrix}$
    &
    $\begin{matrix}
        1    &   -1    \\ 
        -1    &   1
    \end{matrix}$
    &
    $\begin{matrix}
        1    &   1    \\ 
        1    &   1
    \end{matrix}$
    &
    $\begin{matrix}
          1  &   1    \\ 
        1    &   -1    
    \end{matrix}$
    &
    $\begin{matrix}
        -1    &   1    \\ 
        1    &   1    
    \end{matrix}$
    &
    $\begin{matrix}
        1    &   1    \\ 
        -1    &   1    
    \end{matrix}$\\
    \hline     
$\begin{matrix}
        1    &   0    \\ 
        1    &   0    
    \end{matrix}$
    &
    $\begin{matrix}
        1    &   0    \\ 
        0    &   1    
    \end{matrix}$
    &
    $\begin{matrix}
        1    &   1    \\ 
         0   &   0    
    \end{matrix}$
    &
    $\begin{matrix}
        1    &   0    \\ 
        -1    &   0    
    \end{matrix}$
    &
    $\begin{matrix}
        1    &   1    \\ 
        0    &   0    
    \end{matrix}$
    &
    $\begin{matrix}
        1    &   0    \\ 
        0    &   -1    
    \end{matrix}$
    &
    $\begin{matrix}
        1    &   0    \\ 
        0    &   1    
    \end{matrix}$
    &
    $\begin{matrix}
        1    &   0    \\ 
        1    &   0    
    \end{matrix}$
    &
    $\begin{matrix}
        1    &   -1    \\ 
        0    &   0    
    \end{matrix}$\\
    \hline     
$\begin{matrix}
        1    &   0    \\ 
        -1    &   0    
    \end{matrix}$
    &
    $\begin{matrix}
        1    &   0    \\ 
        0    &   -1    
    \end{matrix}$
    &
    $\begin{matrix}
        1    &   -1    \\ 
        0    &   0    
    \end{matrix}$
    &
    $\begin{matrix}
        0    &   1    \\ 
        0    &   -1    
    \end{matrix}$
    &
    $\begin{matrix}
        0    &   0    \\ 
        1    &   1    
    \end{matrix}$
    &
    $\begin{matrix}
        0    &   1    \\ 
       -1    &   0    
    \end{matrix}$
    &
    $\begin{matrix}
        0    &   1    \\ 
        -1    &   0    
    \end{matrix}$
    &
    $\begin{matrix}
        0    &   1    \\ 
        0    &   -1    
    \end{matrix}$
    &
    $\begin{matrix}
        0    &   0    \\ 
        1    &   1    
    \end{matrix}$\\
    \hline     
    \end{tabular}

    Chaque case de ce table représente un rayon de \( \eR^4\); il y en a \( 18\) différents, chacun écris deux fois. Une simple vérification montre que chaque colonne est un système orthogonal. La preuve du théorème de Kochen-Specker revient à montrer que nous ne pouvons pas colorier ce tableau de façon cohérente. En effet, étant donné que chaque vecteur est écrit deux fois, le tableau doit contenir un nombre pair de cases blanches et un nombre pair de cases noires.

    Par ailleurs chaque colonne étant un système orthogonal, chaque colonne contient exactement une case noire; il y a donc exactement neuf cases noires dans le tableau, ce qui est impossible.

\end{proof}

% Tiens, tu lis mes sources ? Je croyais que personne ne lirait jamais ;)
% Bon, alors c'est que tu es un gars bien. Je vais te confier mon plan secret.
% Mon plan est que le théorème de Kochen-Specker et le Karmack's reverse deviennent
% des classiques de l'agreg. Ce sont certainement les deux meilleurs trucs 
% inconnus de niveau agreg.

%+++++++++++++++++++++++++++++++++++++++++++++++++++++++++++++++++++++++++++++++++++++++++++++++++++++++++++++++++++++++++++
\section{Espaces de Bergman}
%+++++++++++++++++++++++++++++++++++++++++++++++++++++++++++++++++++++++++++++++++++++++++++++++++++++++++++++++++++++++++++

Source : \cite{ytMOpe}.

Soit \( \Omega\) un borné dans \( \eC\) et \( D\) le disque unité ouvert de \( \eC\).

\begin{definition}
    L'\defe{espace de Bergman}{espace!de Bergman}\index{Bergman (espace)} sur \( \Omega\), noté \( A^2(\Omega)\)\nomenclature[Y]{\( A^2(\Omega)\)}{espace de Bergman} est l'espace des fonctions holomorphes sur \( \Omega\) qui sont en même temps dans \( L^2(\Omega)\).
\end{definition}
Nous mettons sur \( A^2(\Omega)\) le produit scalaire usuel hérité de \( L^2\) :
\begin{equation}
    \langle f, g\rangle =\int_{\Omega}f(z)\overline{ g(z) }dz.
\end{equation}

\begin{lemma}   \label{LemIZxKfB}
    Soit \( K\subset \Omega\) un compact et \( f\in A^2(\Omega)\). Alors
    \begin{equation}
        \max_{z\in K}| f(z) |\leq \frac{1}{ \sqrt{\pi} }\frac{1}{ d(K,\partial \Omega) }\| f \|_2.
    \end{equation}
\end{lemma}

\begin{proof}
    Soient \( a\in \Omega\) et \( r>0\) tels que \( B(a,r)\subset\Omega\). Nous considérons aussi \( \rho\leq r\). La formule de Cauchy \eqref{EqPzUABM} nous donne   
    \begin{equation}
        f(a)=\frac{1}{ 2\pi i }\int_{B(a,\rho)}\frac{ f(\xi) }{ \xi-a }f\xi=\frac{1}{ 2\pi }\int_0^{2\pi}f(a+\rho e^{i\theta})d\theta
    \end{equation}
    où nous avons utilisé le chemin \( \gamma(\theta)=a+\rho e^{i\theta}\), \( \gamma'(\theta)=i\rho e^{i\theta}\) et \( \rho=| \xi-a |\). Maintenant une astuce est d'écrire
    \begin{equation}
        \frac{ r^2 }{2}f(a)=\int_0^rf(a)\rho d\rho,
    \end{equation}
    et d'y substituer la valeur de \( f(a)\) que nous venons de calculer :
    \begin{subequations}
        \begin{align}
            \frac{ r^2 }{2}f(a)&=\int_0^r\frac{1}{ 2\pi }\int_0^{2\pi}f(a+\rho e^{i\theta})d\theta\rho d\rho\\
            &=\frac{1}{ 2\pi }\int_{B(a,r)}f(z)dz   &   \text{passage aux polaires}\\
            &=\frac{1}{ 2\pi }\langle 1, f\rangle_B   &   \text{produit scalaire sur \( B(a,r)\)}\\
            &\leq\frac{1}{ 2\pi }\sqrt{\langle 1, 1\rangle_B\langle f, f\rangle_B }\\
        \end{align}
    \end{subequations}
    Nous avons donc
    \begin{equation}
        r^2f(a)\leq \frac{1}{ \pi }\sqrt{\langle 1, 1\rangle_B\langle f, f\rangle_B},
    \end{equation}
    et donc
    \begin{equation}
        \pi r^2 f(a)\leq \sqrt{\pi r^2}\| f \|_2,
    \end{equation}
    parce que \( \langle f, f\rangle_B\leq \| f \|_2^2\). En effet le produit scalaire \( \| . \|_2\) est donné par une intégrale sur \( \Omega\) alors que \( B(a,r)\subset \Omega\) et que la fonction qu'on y intègre est positive (c'est \( | f(z) |^2\)). En simplifiant,
    \begin{equation}
        f(a)\leq \frac{1}{ \sqrt{\pi}r }\| f \|_2.
    \end{equation}
    Mais \( r\) a été choisit pour avoir \( B(a,r)\subset\Omega\), donc \( r\leq d(a,\partial \Omega)\) et
    \begin{equation}
        | f(a) |\leq \frac{1}{ d(a,\partial\Omega)\sqrt{\pi} }\| f \|_2.
    \end{equation}
    
    Maintenant si nous prenons \( a\in K\), nous avons encore la minoration \( d(a,\partial K)\leq d(a,\partial \Omega)\) et donc
    \begin{equation}
        | f(a) |\leq\frac{1}{ d(a,\partial K)\sqrt{\pi} }\| f \|_2.
    \end{equation}

\end{proof}

\begin{theorem}
    Soit \( \Omega\) un ouvert de \( \eC\).
    \begin{enumerate}
        \item
            L'espace \( A^2(\Omega)\) est un espace de Hilbert.
        \item
            Si \( D\) est la boule unité dans \( \eC\), une base hilbertienne de \( A^2(D)\) est donnée par les fonctions
            \begin{equation}
                e_n(z)=\sqrt{\frac{ n+1 }{ \pi }}z^n
            \end{equation}
            pour \( n\geq 0\).
    \end{enumerate}
\end{theorem}

\begin{proof}
    Nous commençons par montrer que \( A^2(\Omega)\) est complet. Pour cela nous considérons une suite de Cauchy \( (f_n)\) dans \( A^2(\Omega)\) et un compact \( K\subset \Omega\). Nous savons par le lemme \ref{LemIZxKfB} que
    \begin{equation}
        \max_{z\in K}\big| f_n(z)-f_m(z) \big|\leq \frac{1}{ \sqrt{\pi}d(K,\partial\Omega) }\| f_n-f_m \|_2.
    \end{equation}
    Donc \( f_n\) converge uniformément sur \( K\). Par le théorème de Weierstrass \ref{ThoArYtQO}, la fonction \( f\) est holomorphe. Il existe donc une fonction holomorphe \( f\) qui est limite uniforme sur tout compact de \( \Omega\) de la suite \( (f_n)\).

    Mais \( L^2(\Omega)\) étant complet, la suite \( (f_n)\) a une limite \( g\in L^2(\Omega)\). Ce que nous voudrions faire est prouver que \( f=g\). Notons que tel quel, ce n'est pas vrai parce que \( f\) est une vraie fonction alors que \( g\) est une classe. Ce que nous enseigne la proposition \ref{PropWoywYG} est qu'il existe une sous-suite (qu'on note \( (g_n)\)) qui converge vers \( g\) presque partout. Dans cette dernière phrase, \( g_n\) et \( g\) sont de vraies fonctions, des représentants des classes dans \( L^2\).

    Nous déduisons que \( f=g\) presque partout (ici \( f\) et \( g\) sont les fonctions) parce que la sous-suite converge uniformément vers \( f\) en même temps que presque partout vers \( g\). Donc \( f=g\) dans \( L^2(\Omega)\) (ici \( f\) et \( g\) sont les classes). Donc \( f\in L^2(\Omega)\) et l'espace \( A^2(\Omega)\) est de Hilbert.

    Il nous faut encore prouver que \( (e_n)_{n\geq 0}\) est une base orthonormale. En ce qui concerne les produits scalaires,
    \begin{subequations}
        \begin{align}
            \langle e_m, e_n\rangle &=\sqrt{\frac{ (m+1)(n+1) }{ \pi }}\int_Dz^n\overline{ z^m }dz\\
            &=\sqrt{\frac{ (m+1)(n+1) }{ \pi^2 }}\int_0^1\rho\,d\rho\int_0^{2\pi}d\theta \rho^{m+n} e^{i\theta(n-m)}\\
            &=\sqrt{\frac{ (m+1)(n+1) }{ \pi^2 }}\frac{1}{ m+n+2 }\underbrace{\int_{0}^{2\pi} e^{i\theta(n-m)}d\theta}_{2\pi \delta_{mn}}\\
            &=\sqrt{\frac{ (n+1)^2 }{ \pi^2 }}\frac{1}{ 2n+2 }2\pi \delta_{nm}\\
            &=\delta_{nm}.
        \end{align}
    \end{subequations}
    Donc les fonctions données sont bien orthonormales. Nous devons montrer qu'elles sont denses dans \( A^2(D)\). Soit \( f\in A^2(D)\) et \( c_n(f)=\langle f, e_n\rangle \); nous allons montrer que
    \begin{equation}
        \| f \|_2^2=\sum_{n=0}^{\infty}| \langle f, e_n\rangle  |^2,
    \end{equation}
    parce que le point \ref{ItemQGwoIx} du théorème \ref{ThoyAjoqP} nous indique que ce sera suffisant pour avoir une base hilbertienne.

    Étant donné que \( f\) est holomorphe sur \( D\), le théorème \ref{ThoUHztQe} nous développe \( f\) en série entière :
    \begin{equation}    \label{EqObkbPK}
        f(z)=\sum_{k=0}^{\infty}a_kz^k.
    \end{equation}
    En permutant la somme avec le produit scalaire,
    \begin{equation}
        c_n(f)=\int_Df(z)\bar e_n(z)=\sqrt{\frac{ n+1 }{ \pi }}\int_Df(z)\bar z^ndz.
    \end{equation}
    Afin de profiter de la convergence uniforme de la série \eqref{EqObkbPK} à l'intérieur de \( D\), nous allons exprimer l'intégrale sur \( D\) comme une intégrale sur \( | z |<r\) en faisant tendre \( r\) vers \( 1\) (par le bas). Pour ce faire nous considérons les fonctions
    \begin{equation}
        g_k(z)=\begin{cases}
            f(z)\bar z^n    &   \text{si \( | z |<1-1/k\)}\\
            0    &    \text{sinon.}
        \end{cases}
    \end{equation}
    Ces fonctions sont intégrables sur \( D\) et dominées par \( f(z)\bar z^n\) qui est intégrable sans dépendre de \( k\). Mais nous avons évidemment \( g_k(z)\to f(z)\bar z^n\). Le théorème de la convergence dominée permet alors de permuter l'intégrale et la limite \( k\to \infty\). Cela nous permet d'écrire
    \begin{equation}
        c_n(f)=\sqrt{\frac{ n+1 }{ \pi }}\lim_{r\to 1^-}\int_{| z |<r}\bar z^nf(z)dz=\sqrt{\frac{ n+1 }{ \pi }}\lim_{r\to 1^-}\int_{| z |<r}\sum_{k=0}^{\infty}a_kz^k\bar z^n.
    \end{equation}
    Par la convergence uniforme de la série entière \emph{à l'intérieur} du disque \( D\) nous pouvons permuter l'intégrale et la somme (proposition \ref{PropfeFQWr}) :
    \begin{equation}
        c_n(f)=\sqrt{\frac{ n+1 }{ \pi }}\lim_{r\to 1^-}\sum_{k=0}^{\infty}a_k\int_{| z |<r}z^k\bar z^ndz.
    \end{equation}
    L'intégrale proprement dite est vite calculée et vaut
    \begin{equation}
        \int_{| z |<1}\bar z^nz^kdz=\frac{ \pi r^{2n+2} }{ n+1 }\delta_{kn}.
    \end{equation}
    Nous pouvons donc continuer le calcul de \( c_n(f)\) en effectuant la somme sur \( k\) qui se réduit à changer \( k\) en \( n\) puis en effectuant la limite :
    \begin{equation}
        c_n(f)=\sqrt{\frac{ n+1 }{ \pi }}\lim_{r\to 1^-}\sum_ka_k\frac{ \pi r^{2n+2} }{ n+1 }\delta_{kn}=\sqrt{\frac{ \pi }{ n+1 }}a_n.
    \end{equation}
    
    Nous effectuons le même genre de calculs pour évaluer \( \| f \|^2_2\) :
    \begin{subequations}
        \begin{align}
            \| f \|_2^2&=\int_D| f(z) |^2dz\\
            &=\lim_{r\to 1^-}\int_{| z |<r}f(z)\sum_{k=0}^{\infty}\bar a_k\bar z_kdz\\
            &=\lim_{r\to 1^-}\sum_{k=0}^{\infty}\bar a_k\int_{| z |<r}f(z)\bar z^kdz&\text{permuter \( \sum\) et \( \int\)}\\
            &=\lim_{r\to 1^-}\sum_{k=0}^{\infty}\bar a_ka_k\frac{ \pi r^{2k+2} }{ k+1 }&\text{intégrale déjà faite}.
        \end{align}
    \end{subequations}
    Mais nous savons déjà que \( c_n(f)=\sqrt{\pi/(n+1)}\), donc ce qui est dans la somme est \( \pi\bar a_ka_k/(n+1)=| c_k(f) |^2\). Nous avons donc
    \begin{equation}
        \| f \|^2_2=\lim_{r\to 1^-}\sum_{k=0}^{\infty}| c_k(f) |^2 r^{2k+2}.
    \end{equation}
    La fonction (de \( r\)) constante \( | c_k(f) |^2\) domine \( | c_k(f)r^{2k+2} |\) tout en ayant une somme (sur \( k\)) qui converge; en effet la proposition \ref{PropHKqVHj} nous indique que \( \sum_j| c_k(f) |^2\leq \| f \|_2^2\). Le théorème de la convergence dominée nous permet d'inverser la limite et la somme pour obtenir le résultat attendu :
    \begin{equation}
        \| f \|_2^2=\sum_{k=0}^{\infty}| c_k(f) |^2.
    \end{equation}
\end{proof}


\chapter{Théorie de la mesure}
% This is part of Agregation : modélisation
% Copyright (c) 2011
%   Laurent Claessens
% See the file fdl-1.3.txt for copying conditions.

Ce chapitre provient en grande partie de \cite{ProbaDanielLi}.

%+++++++++++++++++++++++++++++++++++++++++++++++++++++++++++++++++++++++++++++++++++++++++++++++++++++++++++++++++++++++++++
\section{Espaces mesurés}
%+++++++++++++++++++++++++++++++++++++++++++++++++++++++++++++++++++++++++++++++++++++++++++++++++++++++++++++++++++++++++++

\begin{definition}
    Si \( \Omega\) est un ensemble, un ensemble \( \tribA\) de sous-ensembles de \( \Omega\) est une \defe{tribu}{tribu} si
    \begin{enumerate}
        \item
            \( \Omega\in\tribA\);
        \item
            \( A\cup B\in\tribA\) pour tout \( A,B\in\tribA\), ce qui signifie que toutes les unions finies d'éléments de \( \tribA\) sont dans \( \tribA\);
        \item
            \( \complement A\in A\) pour tout \( A\in\tribA\);
        \item
            si \( A_n\) est une suite dénombrable d'éléments de \( \tribA\), alors \( \sup_{n\geq 1}A_n\in\tribA\).
    \end{enumerate}
    Le couple \( (\Omega,\tribA)\) est alors un \defe{espace mesuré}{espace!mesuré}.
\end{definition}

La tribu que nous utiliserons toujours dans \( \eR^d\) est la tribu des \defe{boréliens}{boréliens}, notée \( \Borelien(\eR^d)\), qui est la tribu engendrée par les ouverts de \( \eR^d\). Une fonction \( f\colon (\Omega,\tribA)\to (\eR^d,\Borelien(\eR^d))\) est \defe{borélienne}{borélienne} si pour tout \( \mO\in\Borelien\), \( f^{-1}(\mO)\in\tribA\).

\begin{definition}
    Une \defe{\wikipedia{en}{Measure_space}{mesure}}{mesure} sur l'espace mesurable \( (\Omega,\tribA)\) est une application \( \mu\colon \tribA\to \eR\cup\{ \infty \}\) telle que
    \begin{enumerate}
        \item
            \( \mu(A)\geq 0\) pour tout \( A\in\tribA\);
        \item
            \( \mu(\emptyset)=0\);
        \item
            \( \mu\left( \bigcup_{i=0}^{\infty}A_i\right)=\sum_{i=0}^{\infty}\mu(A_i)\) si les \( A_i\) sont des éléments de \( \tribA\) deux à deux disjoints.
    \end{enumerate}
\end{definition}

Si \( \mu\) est une mesure sur \( \eR^d\), une fonction \( f\colon \eR^d\to \eR\) est une \defe{densité}{densité d'une mesure} si pour tout \( A\subset\eR^d\) nous avons
\begin{equation}
    \mu(A)=\int_Af(x)dx
\end{equation}
où \( dx\) est la mesure de Lebesgue.

%---------------------------------------------------------------------------------------------------------------------------
\subsection{Intégrale par rapport à une mesure}
%---------------------------------------------------------------------------------------------------------------------------

Une fonction \( f\colon (\Omega,\tribA)\to (\Omega',\tribA')\) est \defe{mesurable}{mesurable!fonction} si 
\begin{equation}
    f^{-1}(E)\in\tribA
\end{equation}
pour tout \( E\in\tribA'\).


Une mesure \( \mu\) sur un espace mesurable \( (\Omega,\tribA)\) permet de définir une fonctionnelle linéaire sur l'ensemble des fonctions mesurables \( \Omega\to \eR\). Cette fonctionnelle linéaire est l'intégrale que nous allons définir à présent.

D'abord nous considérons les fonction \defe{simples}{simple!fonction}\index{fonction!simple}, c'est à dire les fonctions de la forme
\begin{equation}
    f=\sum_{i=1}^Na_i\caract_{E_i}
\end{equation}
où \( a_i\in\eR\) tandis que les \( E_i\) sont des ensembles \( \mu\)-mesurables. Si \( Y\in \tribA\) nous définissons
\begin{equation}
    \int_Yfd\mu=\sum_ia_i\mu(Y\cap E_i).
\end{equation}
Pour une fonction \( \mu\)-mesurable générale \( f\colon \Omega\to \mathopen[ 0 , \infty \mathclose]\) nous définissons l'intégrale de \( f\) sur \( Y\) par
\begin{equation}        \label{EqDefintYfdmu}
    \int_Yfd\mu=\sup\Big\{ \int_Yhd\mu\,\text{où \( h\) est une fonction simple et mesurable telle que \( 0\leq h\leq f\)} \Big\}.
\end{equation}
Maintenant nous définissons
\begin{equation}
    \mu(f)=\int_{\Omega}f
\end{equation}
si \( f\) est une fonction mesurable sur \( \Omega\).

\begin{remark}
    Dans \( \eR^d\), quasiment toutes les fonctions et ensembles sont mesurables. En effet la construction d'ensembles non mesurables demande obligatoirement l'utilisation de l'axiome du choix; de tels ensembles doivent être construits «exprès pour». Il y a très peu de chances pour que vous tombiez sur un ensemble non mesurable de \( \eR^d\) sans que vous ne vous en rendiez compte.

    Par exemple la variable aléatoire 
    \begin{equation}
        X(\omega)=\begin{cases}
            \frac{1}{ \omega }    &   \text{si $ \omega\neq 0$}\\
            \infty    &    \text{$\omega=0$}.
        \end{cases}
    \end{equation}
    est mesurable, mais non intégrable.
\end{remark}



%+++++++++++++++++++++++++++++++++++++++++++++++++++++++++++++++++++++++++++++++++++++++++++++++++++++++++++++++++++++++++++
\section{Espace de probabilité}
%+++++++++++++++++++++++++++++++++++++++++++++++++++++++++++++++++++++++++++++++++++++++++++++++++++++++++++++++++++++++++++

Une \defe{mesure de probabilité}{mesure!probabilité} sur un espace mesuré \( (\Omega,\tribA)\) est une mesure positive telle que \( P(\Omega)=1\). Dans ce cas, le triple \( (\Omega,\tribA,P)\) est un \defe{espace de probabilité}{espace!de probabilité}.

Un point \( \omega\in\Omega\) est une \defe{observation}{observation}, une partie mesurable \( A\in\tribA\) est un \defe{événement}{événement}. L'ensemble \( A\cup B\) représente l'événement \( A\) ou \( B\) tandis que l'ensemble \( A\cap B\) représente l'événement \( A\) et \( B\).


Si les \( A_n\) sont des événements, nous définissons la \defe{limite supérieur}{limite!supérieure} et la \defe{limite inférieure}{limite!inférieure} de la suite \( A_n\) par
\begin{equation}
    \limsup_{n\to\infty}A_n=\bigcap_{n\geq 1}\bigcup_{k\geq n}A_k
\end{equation}
et
\begin{equation}
    \liminf_{n\to\infty}A_n=\bigcup_{n\geq 1}\bigcap_{k\geq n}A_k
\end{equation}
Si \( \omega\in\liminf A_n\), alors \( \omega\) réalise tous les \( A_n\) sauf un nombre fini.

Nous avons
\begin{equation}
    \limsup A_n=\{ \omega\in\Omega\tq \omega\in A_n\text{pour une infinité de \( n\)} \}.
\end{equation}

\begin{theorem}[Borel-Cantelli]\index{théorème!Borel-Cantelli}
    Si
    \begin{equation}
        \sum_{n=1}^{\infty}P(A_n)<\infty
    \end{equation}
    alors \( P(\limsup A_n)=0\).
\end{theorem}

\begin{proof}
    La condition \( \sum_{n\geq 1}P(A_n)<\infty\) signifie que la fonction
    \begin{equation}
        \varphi=\sum_{n\geq 1}\caract_{A_n}
    \end{equation}
    est \( P\)-intégrable. Par conséquent, elle est finie presque partout (au sens de \( P\)), c'est à dire
    \begin{equation}
        P(\varphi=\infty)=0.
    \end{equation}
    Les points \( \omega\) sur lesquels \( \varphi(\omega)=\infty\) sont ceux tels que
    \begin{equation}
        \sum_{n\geq 1}\caract_{A_n}(\omega)=\infty,
    \end{equation}
    c'est à dire les \( \omega\) qui appartiennent à une infinité d'ensembles \( A_n\), ou encore les \( \omega\in\limsup A_n\). Nous avons donc montré que
    \begin{equation}
        \{ \omega\tq \varphi(\omega)=\infty \}=\{ \omega\in\Omega\tq \omega\in A_n\text{pour une infinité de \( n\)} \}=\limsup A_n.
    \end{equation}
    Or l'hypothèse signifie que la probabilité du membre de gauche est nulle.
\end{proof}

\begin{corollary}
    Si \( \sum_{n=1}^{\infty}P(\complement A_n)<\infty\), alors presque surement tous les \( B_n\) sont réalisés à l'exception d'un nombre fini.
\end{corollary}

%+++++++++++++++++++++++++++++++++++++++++++++++++++++++++++++++++++++++++++++++++++++++++++++++++++++++++++++++++++++++++++
\section{Variables aléatoires}
%+++++++++++++++++++++++++++++++++++++++++++++++++++++++++++++++++++++++++++++++++++++++++++++++++++++++++++++++++++++++++++

\begin{definition}
    Une \defe{variable aléatoire}{variable aléatoire} est une application mesurable
    \begin{equation}
        X\colon (\Omega,\tribA)\to (\eR^d,\Borelien(\eR^d)).
    \end{equation}
\end{definition}
Nous convenons que \( \eR^1=\bar\eR\), c'est à dire que dans le cas où la variable aléatoire \( X\) est réelle, nous acceptons les valeurs \( \pm\infty\).

On dit que la variable aléatoire \( X\) a un \defe{moment d'ordre \( p\)}{moment} si \( X\in L^p(\Omega,\tribA,P)\) (\( 1\leq p<\infty\)). Si \( X\) est \defe{intégrable}{variable aléatoire!intégrable} (c'est à dire si \( X\in L^1\)), alors nous définissons
\begin{equation}
    E(X)=\int_{\Omega}XdP\in\eR^d.
\end{equation}
Si \( E(X)=0\) nous disons que la variable aléatoire est \defe{centrée}{variable aléatoire!centrée}. La variable aléatoire \( X-E(X)\) est la variable aléatoire centrée associée à \( X\).

%---------------------------------------------------------------------------------------------------------------------------
\subsection{Variance}
%---------------------------------------------------------------------------------------------------------------------------

Si \( X\in L^2(\Omega,\tribA,P)\) alors nous définissons la \defe{variance}{variance} de \( X\) par
\begin{equation}
    \Var(X)=E\big( [X-E(X)]^2 \big).
\end{equation}

\begin{proposition}
    La variance de la variable aléatoire \( X\) peut être exprimée par la formule
    \begin{equation}
        \Var(X)=E(X^2)-[E(X)]^2
    \end{equation}
    où \( X^2=X\cdot X\) et \( E(X)^2=\) sont des produits scalaires dans \( \eR^d\).
\end{proposition}

\begin{proof}
    De façon explicite, nous avons
    \begin{equation}
        E\big( [X-E(X)]^2 \big)=\int_{\Omega}\big( X(\omega)-E(X) \big)\cdot\big( X(\omega)-E(X) \big)dP(\omega)
    \end{equation}
    où \( E(X)\in\eR^d\) est une constante. En développant le produit scalaire nous avons
    \begin{subequations}
        \begin{align}
            E\big( [X-E(X)]^2 \big)&=E\big( X^2-2X\cdot E(X)+E(X)^2 \big)\\
            &=E(X^2)-2E(X)^2+E(X)^2\\
            &=E(X^2)-E(X)^2.
        \end{align}
    \end{subequations}
\end{proof}

Si \( a\in\eR\) nous avons
\begin{equation}
    \Var(ax)=a^2\Var(X)
\end{equation}
et nous définissons l'\defe{écart-type}{écart-type} de \( X\) par
\begin{equation}
    \sigma_X=\sqrt{\Var(X)}.
\end{equation}
En d'autres termes,
\begin{equation}
    \sigma_X=\| X-E(X) \|_{L^2}.
\end{equation}
On définit encore la \defe{moyenne quadratique}{moyenne!quadratique} de \( X\) par
\begin{equation}
    \| X \|_{L^2}=\big[ E(X^2) \big]^{1/2}.
\end{equation}

Si \( X\) est une variable aléatoire réelle, nous définissons sa \defe{fonction de répartition}{fonction!de répartition} par
\begin{equation}
    \begin{aligned}
        F_X\colon \eR&\to \mathopen[ 0 , 1 \mathclose] \\
        F_X(x)&=P(X\leq x). 
    \end{aligned}
\end{equation}

%---------------------------------------------------------------------------------------------------------------------------
\subsection{Loi d'une variable aléatoire}
%---------------------------------------------------------------------------------------------------------------------------

La \defe{loi}{loi d'une variable aléatoire} de la variable aléatoire \( X\), notée \( P_X\) est la mesure image de \( P\) par \( X\), c'est à dire
\begin{equation}
    P_X(B)=P(X\in B)
\end{equation}
pour tout borélien \( B\subset\eR^d\). Note :
\begin{equation}
    P(X\in B)=P\big( \{ \omega\in\Omega\tq X(\omega)\in B \} \big)=P\big( X^{-1}(B) \big).
\end{equation}
En particulier \( P_X\) est une mesure de probabilité sur \( \eR^d\) parce que 
\begin{equation}
    P_X(\eR^d)=P(\Omega)=1.
\end{equation}
Si \( Q\) est une mesure de probabilité sur \( \eR^d\), nous notons \( X\sim Q\) si \( P_X=Q\). Nous disons alors que «\( X\) suit la loi \( Q\)».

La proposition suivante permet de calculer en pratique les intégrales qui définissent par exemple l'espérance mathématique d'une variable aléatoire.
\begin{proposition}
    Si \( f\colon \eR^d\to \bar\eR\) est borélienne et si \( X\) est une variable aléatoire, alors
    \begin{equation}
        \int_{\eR^d}f(x)dP_x(x)=\int_{\Omega}f\big( X(\omega) \big)dP(\omega)=E(f\circ X).
    \end{equation}
\end{proposition}

\begin{proof}
    No proof.
\end{proof}

En utilisant cette proposition nous trouvons une formule pratique pour l'espérance d'une variable aléatoire réelle:
\begin{equation}
    E(X)=\int_{\Omega}X(\omega)dP(\omega)=\int_{\eR}xdP_X(x),
\end{equation}
en vertu de la proposition appliquée à la fonction \( f(x)=x\).

\begin{proposition}
    Une variable aléatoire réelle \( X\) est intégrable si et seulement si \( P(x=\pm\infty)=0\) et
    \begin{equation}
        \int_{\eR}| x |dP_X(x)<\infty.
    \end{equation}
\end{proposition}

%+++++++++++++++++++++++++++++++++++++++++++++++++++++++++++++++++++++++++++++++++++++++++++++++++++++++++++++++++++++++++++
\section{Lois usuelles}
%+++++++++++++++++++++++++++++++++++++++++++++++++++++++++++++++++++++++++++++++++++++++++++++++++++++++++++++++++++++++++++

%---------------------------------------------------------------------------------------------------------------------------
\subsection{Loi de Bernouilli}
%---------------------------------------------------------------------------------------------------------------------------

Une variable aléatoire réelle est de \defe{Berouilli}{Bernouilli} de paramètre \( p\) (\( 0<p<1\)) si
\begin{equation}
    X\colon \Omega\to \eR
\end{equation}
avec \( P(x=1)=p\) et \( P(X=0)=1-p\). En tant que mesure sur \( \eR\), nous avons
\begin{equation}
    P_X=p\delta_1+(1-p)\delta_0.
\end{equation}
Une fonction \( h\) qui réalise le supremum de la formule \eqref{EqDefintYfdmu} est par exemple une fonction en escalier qui vaut en \( x\) le plus petit entier plus grand ou égal à \( x\). L'espérance d'une loi de Bernouilli est alors
\begin{equation}
    E(x)=p.
\end{equation}
Étant donné que la variable aléatoire \( X\) prend seulement les valeurs \( 0\) et \( 1\), nous avons pour tout ensemble mesurable \( B\)
\begin{equation}
    P_{X^2}(B)=P(X^2\in B)=P(X\in B),
\end{equation}
et par conséquent \( P_{X^2}=P_X\) et \( E(X^2)=E(X)\). Nous trouvons donc la variance
\begin{equation}
    \Var(X)=E(X^2)-E(X)^2=p-p^2=p(1-p).
\end{equation}

%---------------------------------------------------------------------------------------------------------------------------
\subsection{Loi normale}
%---------------------------------------------------------------------------------------------------------------------------






\chapter{Variables aléatoires et théorie des probabilités}
% This is part of Agregation : modélisation
% Copyright (c) 2011
%   Laurent Claessens
% See the file fdl-1.3.txt for copying conditions.


Une bonne référence pour ce chapitre est \cite{ProbaDanielLi}.

%+++++++++++++++++++++++++++++++++++++++++++++++++++++++++++++++++++++++++++++++++++++++++++++++++++++++++++++++++++++++++++
\section{Espace de probabilité}
%+++++++++++++++++++++++++++++++++++++++++++++++++++++++++++++++++++++++++++++++++++++++++++++++++++++++++++++++++++++++++++

Une \defe{mesure de probabilité}{mesure!probabilité} sur un espace mesuré \( (\Omega,\tribA)\) est une mesure positive telle que \( P(\Omega)=1\). Dans ce cas, le triple \( (\Omega,\tribA,P)\) est un \defe{espace de probabilité}{espace!de probabilité}.

Un point \( \omega\in\Omega\) est une \defe{observation}{observation}, une partie mesurable \( A\in\tribA\) est un \defe{événement}{événement}. L'ensemble \( A\cup B\) représente l'événement \( A\) ou \( B\) tandis que l'ensemble \( A\cap B\) représente l'événement \( A\) et \( B\).


Si les \( A_n\) sont des événements, nous définissons la \defe{limite supérieur}{limite!supérieure} et la \defe{limite inférieure}{limite!inférieure} de la suite \( A_n\) par
\begin{equation}
    \limsup_{n\to\infty}A_n=\bigcap_{n\geq 1}\bigcup_{k\geq n}A_k
\end{equation}
et
\begin{equation}
    \liminf_{n\to\infty}A_n=\bigcup_{n\geq 1}\bigcap_{k\geq n}A_k
\end{equation}
Si \( \omega\in\liminf A_n\), alors \( \omega\) réalise tous les \( A_n\) sauf un nombre fini.

Nous avons
\begin{equation}
    \limsup A_n=\{ \omega\in\Omega\tq \omega\in A_n\text{pour une infinité de \( n\)} \}.
\end{equation}

\begin{theorem}[Borel-Cantelli]\index{théorème!Borel-Cantelli}
    Si
    \begin{equation}
        \sum_{n=1}^{\infty}P(A_n)<\infty
    \end{equation}
    alors \( P(\limsup A_n)=0\).
\end{theorem}

\begin{proof}
    La condition \( \sum_{n\geq 1}P(A_n)<\infty\) signifie que la fonction
    \begin{equation}
        \varphi=\sum_{n\geq 1}\caract_{A_n}
    \end{equation}
    est \( P\)-intégrable. Par conséquent, elle est finie presque partout (au sens de \( P\)), c'est à dire
    \begin{equation}
        P(\varphi=\infty)=0.
    \end{equation}
    Les points \( \omega\) sur lesquels \( \varphi(\omega)=\infty\) sont ceux tels que
    \begin{equation}
        \sum_{n\geq 1}\caract_{A_n}(\omega)=\infty,
    \end{equation}
    c'est à dire les \( \omega\) qui appartiennent à une infinité d'ensembles \( A_n\), ou encore les \( \omega\in\limsup A_n\). Nous avons donc montré que
    \begin{equation}
        \{ \omega\tq \varphi(\omega)=\infty \}=\{ \omega\in\Omega\tq \omega\in A_n\text{pour une infinité de \( n\)} \}=\limsup A_n.
    \end{equation}
    Or l'hypothèse signifie que la probabilité du membre de gauche est nulle.
\end{proof}

\begin{corollary}
    Si \( \sum_{n=1}^{\infty}P(\complement A_n)<\infty\), alors presque surement tous les \( B_n\) sont réalisés à l'exception d'un nombre fini.
\end{corollary}

%+++++++++++++++++++++++++++++++++++++++++++++++++++++++++++++++++++++++++++++++++++++++++++++++++++++++++++++++++++++++++++
\section{Variables aléatoires}
%+++++++++++++++++++++++++++++++++++++++++++++++++++++++++++++++++++++++++++++++++++++++++++++++++++++++++++++++++++++++++++

\begin{definition}
    Une \defe{variable aléatoire}{variable aléatoire} est une application mesurable
    \begin{equation}
        X\colon (\Omega,\tribA)\to (\eR^d,\Borelien(\eR^d)).
    \end{equation}
\end{definition}
Nous convenons que \( \eR^1=\bar\eR\), c'est à dire que dans le cas où la variable aléatoire \( X\) est réelle, nous acceptons les valeurs \( \pm\infty\).


\begin{definition}      \label{DefAbsoluCont}
    Une fonction \( F\colon \eR\to \eR\) est \defe{absolument continue}{absolument continue} sur \( \mathopen[ a , b \mathclose]\) si il existe une fonction \( f\) sur \( \mathopen[ a , b \mathclose]\) telle que
    \begin{equation}
        F(x)=\int_a^xf(t)dt
    \end{equation}
    pour tout \( x\in\mathopen[ a , b \mathclose]\).

    Une variable aléatoire réelle \( X\) est \defe{absolument continue}{variable aléatoire!absolument continue} si il existe une fonction positive et intégrable \( f\colon \eR\to \eR\) telle que pour tout intervalle \( I\subset\eR\),
    \begin{equation}
        P(X\in I)=\int_If(t)dt.
    \end{equation}
    Nous disons alors que \( f\) est la \defe{densité}{densité} de \( X\).
\end{definition}


%---------------------------------------------------------------------------------------------------------------------------
\subsection{Indépendance}
%---------------------------------------------------------------------------------------------------------------------------

La définition suivante vient de l'instructive motivation de \cite{CourgGudRennes}.


La définition d'indépendance de deux événements se généralise à \( n\) événements de la façon suivante.
\begin{definition}
    Nous disons que les événements \( A_1,\ldots,A_n\) sont \defe{indépendants}{indépendance!événements} si pour tout choix \( \{ i_1,\ldots,i_k \}\subset\{ 1,\ldots,n \}\) nous avons
    \begin{equation}
        P(A_{i_1}\cap\ldots\cap A_{i_k})=P(A_{i_1})\ldots P(A_{i_k}).
    \end{equation}
    Les sous tribus \( \tribA_1,\ldots,\tribA_n\) sont \defe{indépendantes}{indépendance!sous tribus} si pour tout choix \( A_i\in \tribA_i\), les événements \( A_i\) sont indépendants.
\end{definition}

\begin{example}
    Soit \( \Omega=\mathopen[ 0 , 1 \mathclose]\times \mathopen[ 0 , 1 \mathclose]\) muni de la mesure de Lebesgue. Soient \( A=\mathopen[ 0 , a \mathclose]\times \mathopen[ 0 , 1 \mathclose]\) et \( B=\mathopen[ 0 , 1 \mathclose]\times \mathopen[ 0 , b \mathclose]\). Nous avons \( P(A)=a\) et \( P(B)=b\) ainsi que \( P(A\cup B)=ab\).
\end{example}

\begin{lemma}       \label{LemTribIndepProdProb}
    Les tribus \( \tribA_1,\ldots,\tribA_n\) sont indépendantes si et seulement si
    \begin{equation}
        P(A_1\cap\ldots\cap A_n)=P(A_1)\ldots P(A_n)
    \end{equation}
    pour tout \( A_i\in\tribA_i\).
\end{lemma}

\begin{proof}
    L'implication dans le sens direct découle immédiatement des définitions.

    Nous supposons avoir un choix \( (A_i)_{i=1,\ldots,n}\) avec \( A_i\in\tribA_i\) et nous devons montrer que ces événements sont indépendants, c'est à dire que si \( J\subset\{ 1,\ldots,n \}\) alors les événements \( (A_j)_{j\in J}\) sont indépendants. Sans perte de généralité, nous pouvons supposer que si \( i\notin J\), \( A_i=\Omega\). Alors nous avons
    \begin{equation}
        P\big( \bigcap_{j\in J}A_j \big)=P\big( \bigcap_{i=1}^nA_i \big)=\prod_{i=1}^nP(A_i)=\prod_{j\in J}P(A_j)
    \end{equation}
    parce que \( P(A_i)=P(\Omega)=1\) lorsque \( i\) n'est pas dans \( J\).
\end{proof}

Si \( A\) est un événement, la \defe{tribu engendrée}{tribu!engendrée!par un événement} par \( A\) est
\begin{equation}
    \sigma(A)=\{ \emptyset,A,\complement A,\Omega \}.
\end{equation}

Soit \( X\colon \Omega\to eR^d\) une variable aléatoire. La \defe{tribu engendrée}{tribu!engendrée!par une variable aléatoire} est
\begin{equation}
    \tribA_X=\{ X^{-1}(B)\tq B\in\Borelien(\eR^d) \}.
\end{equation}
Cela est la plus petite tribu sous tribu de \( \tribA\) pour laquelle \( X\) est mesurable.


Nous disons que les variables aléatoires \( X_k\colon \Omega\to \eR^d\) sont \defe{indépendantes}{indépendance!variables aléatoires} si les tribus \( \tribA_{X_1},\ldots,\tribA_{X_n}\) le sont.

\begin{proposition}
    Soient \( (X_k\colon \Omega\to \eR^{d_k})\) des variables aléatoires indépendantes et \( B_k\in \Borelien(\eR^{d_k})\). Alors
    \begin{equation}
        P(X_k\in B_k\forall k\leq n)=P(X_1\in B_1)\ldots P(X_n\in B_n).
    \end{equation}
\end{proposition}

\begin{proof}
    Lorsque nous écrivons \( X_i\in B_i\), nous parlons de l'événement
    \begin{equation}
        (X_i\in _i)=\{ \omega\in\Omega\tq X_i(\omega)\in B_i \}=X_i^{-1}(B_i)\in \tribA_{X_i}.
    \end{equation}
    Vu que par hypothèse les tribus \( (\tribA_i)\) sont indépendantes, le lemme \ref{LemTribIndepProdProb} nous montre que
    \begin{equation}
        P\big( \bigcap_{i=1}^nX_i\in B_i \big)=\prod_iP(X_i\in B_i).
    \end{equation}
    Il reste à voir que l'ensemble \( X_i^{-1}(B_i)\) fait partie de la tribu \( \tribA\) de départ. Cela est la définition du fait que l'application \( X_i\) soit une variable aléatoire : elle doit être mesurable en tant qu'application
    \begin{equation}
        X_i\colon (\Omega,\tribA)\to (\eR^d,\Borelien(\eR^d)).
    \end{equation}
\end{proof}

\begin{lemma}       \label{LemIndepEvenCompl}
    Les événements \( (A_i)_{i=0,\ldots,n}\) sont indépendants si et seulement si les événements obtenus en remplaçant certains des \( A_i\) par \( \complement A_i\).
\end{lemma}

\begin{proof}
    Sans perte de généralité, nous pouvons nous contenter de prouver que les événements \( \complement A_0,A_1,\ldots,A_n\) sont indépendants sous l'hypothèse que les événements \( A_0,A_1,\ldots,A_n\) sont indépendants. Soit \( I\) un sous-ensemble de \( \{ 1,\ldots,n \}\). Nous avons
    \begin{subequations}
        \begin{align}
            P\big( \complement A_0\bigcap_{i\in I}A_i \big)&=P\big( \bigcap_{i\in I}A_i\setminus\bigcap_{i\in I}A_i\cap A_0 \big)\\
            &=P\big( \bigcap_{i\in I}A_i \big)-P\big( \bigcap_{i\in I}A_i\cap A_0 \big)\\
            &=P\big( \bigcap_{i\in I}A_i \big)\big( 1-P(\complement A_0) \big)\\
            &=P\big( \bigcap_{i\in I}A_i \big)P(\complement A_0).
        \end{align}
    \end{subequations}
\end{proof}

\begin{proposition}
    Les événements \( (A_i)_{i=1,\ldots,n}\) sont indépendants si et seulement si les variables aléatoires \( \mtu_{A_1},\ldots,\mtu_{A_n}\) le sont.
\end{proposition}

\begin{proof}
    La tribu engendrée par la variable aléatoire \( \mtu_{A_k}\) est
    \begin{equation}    \label{EqtribAAimtu}
        \tribA_{\mtu_{A_k}}=\{ \emptyset,A_k,\complement A_k,\Omega \}.
    \end{equation}
    En effet si \( 1\in B\), alors \( A_i\subset\mtu_{A_i}^{-1}(B)\), et si \( 0\in B\), alors \( \complement A_i\subset\mtu_{A_i}^{-1}(B)\). Les éléments \( 0\) et \( 1\) sont tous deux soit dans \( B\), soit hors de \( B\). Cela donne les \( 4\) possibilités énumérées dans \eqref{EqtribAAimtu}.

    Supposons que les événements \( (A_i)\) sont indépendants. Nous devons vérifier que les tribus le soient, c'est à dire que les événements \( A_i\) et \( \complement A_j\) sont indépendants. Cela est une conséquence du lemme \ref{LemIndepEvenCompl}.
\end{proof}

\begin{theorem}     \label{ThofrestemesurablesXYYX}
    Soit \( X\colon \Omega\to \eR^d\) une variable aléatoire. Une fonction \( Y\colon \Omega\to \eR^{d'}\) est une variable aléatoire \( \tribA_X\)-mesurable si et seulement si il existe une fonction borélienne \( f\colon \eR^d\to \eR^{d'}\) telle que \( Y=f(X)\).
\end{theorem}

\begin{proposition}
    Soient des variables aléatoires \( X_k\colon \Omega\to \eR^{d_k}\) des variables aléatoires indépendantes et des fonctions boréliennes \( f_k\colon eR^{d_k}\to \eR^{d'_k}\). Alors les variables aléatoires \( f_k(X_k)\) sont indépendantes.
\end{proposition}

\begin{proof}
    Le théorème \ref{ThofrestemesurablesXYYX} assure que les applications
    \begin{equation}
        f_k\circ X_k\colon \Omega\to \eR^{d_k}
    \end{equation}
    sont \( \tribA_{X_k}\)-mesurables. En particulier pour tout borélien \( B\subset\eR^{d'_k}\), nous avons \( X^{-1}_k\circ f^{-1}_k(B)\in\tribA_{X_k}\). Nous avons donc
    \begin{equation}
        \sigma(f_k\circ X_k)\subset\sigma(X_k),
    \end{equation}
    et par conséquent les tribus \( \sigma(f_k\circ X_k)\) sont indépendantes étant donné que les tribus \( \sigma(X_k)\) le sont.
\end{proof}

%---------------------------------------------------------------------------------------------------------------------------
\subsection{Probabilité conditionnelle}
%---------------------------------------------------------------------------------------------------------------------------

Soit \( (\Omega,\tribA,P)\) un espace de probabilité et \( B\in\tribA\) avec $P(B)>0$. Nous pouvons introduire une nouvelle loi de probabilité sur \( (\Omega,\tribA)\) en définissant
\begin{equation}    \label{EqProbCond}
    P_B(A)=\frac{ P(A\cap B) }{ P(B) }=P(A|B).
\end{equation}
La première égalité est la définition de \( P_B\). La seconde est une notation. On vérifie que \( (\Omega,\tribA,P)\) est un espace de probabilité parce que \( P_B(\Omega)=1\) et 
\begin{equation}
    P_B(\bigcup_iA_i)=\sum_iP_B(A_i)
\end{equation}
si les \( A_i\) sont deux à deux disjoints.

Une conséquence immédiate de \eqref{EqProbCond} est que si \( A\) et \( B\) sont des événements indépendants alors
\begin{equation}
    P(A|B)=\frac{ P(A\cap B) }{ P(B) }=P(A).
\end{equation}

\begin{theorem}     \label{ThoBayesEtAutres}
    Soient \( (B_n)_{n\geq 1}\) une partition finie de \( \Omega\) telle que \( P(B_i)>0\). Soit \( A\in\tribA\) tel que \( P(A)>0\).
    \begin{enumerate}
        \item
            Si \( A\), \( B\) et \( C\) sont des événements, alors
            \begin{equation}
                P(A\cap B|C)=P(A|B\cap C)P(B|C).
            \end{equation}
        \item
            Si \( P(B)>0\), alors \( P(A\cap B)=P(A|B)P(B)=P(B|A)P(A)\).
        \item On a la \defe{formule des probabilités totales}{formule!probabilité totales} :
            \begin{equation}
                P(A)=\sum_{i=1}^nP(A|B_i)P(B_i)=\sum_iP(A\cap B_i).
            \end{equation}
        \item
            On a la \defe{formule de Bayes}{formule!Bayes} :
            \begin{equation}
                P(B_k|A)=\frac{ P(A|B_k)P(B_k) }{ \sum_iP(A|B_i)P(B_i) }.
            \end{equation}
    \end{enumerate}
\end{theorem}

\begin{proof}
    \begin{enumerate}
        \item
            En développant le membre de droite,
            \begin{equation}
                \begin{aligned}[]
                    P(A\cap B|C)&=\frac{ P(A\cap B\cap C) }{ P(B\cap C) }\frac{ P(B\cap C) }{ P(C) }\\
                    &=P(A\cap B|C).
                \end{aligned}
            \end{equation}
        \item
            C'est la définition de \( P(A|B)\) et \( P(B|A)\).
        \item
            Vu que les \( B_i\) forment une partition, nous avons
            \begin{equation}
                P(A)=\sum_iP(A\cap B_i)=\sum_iP(A|B_i)P(B_i).
            \end{equation}
        \item
            En utilisant les deux premiers points, nous trouvons
            \begin{equation}
                \begin{aligned}[]
                    P(A|B_k)P(B_k)&=P(A\cap B_k)\\
                    &=P(B_k|A)P(A)\\
                    &=P(B_k|A)\sum_iP(A|B_i)P(B_i).
                \end{aligned}
            \end{equation}
    \end{enumerate}
\end{proof}

%---------------------------------------------------------------------------------------------------------------------------
\subsection{Espérance}
%---------------------------------------------------------------------------------------------------------------------------

Nous dirons que la variable aléatoire \( X\) a un \defe{moment d'ordre \( p\)}{moment} si \( X\in L^p(\Omega,\tribA,P)\) (\( 1\leq p<\infty\)). Si \( X\) est \defe{intégrable}{variable aléatoire!intégrable} (c'est à dire si \( X\in L^1\)), alors nous définissons l'\defe{espérance}{espérance} de \( X\) par
\begin{equation}
    E(X)=\int_{\Omega}XdP\in\eR^d.
\end{equation}
Si \( E(X)=0\) nous disons que la variable aléatoire est \defe{centrée}{variable aléatoire!centrée}. La variable aléatoire \( X-E(X)\) est la variable aléatoire centrée associée à \( X\).

Le \defe{moment}{moment} d'ordre \( p\) de la variable aléatoire \( X\) est l'espérance
\begin{equation}
    m_n(X)=E(X^n).
\end{equation}

\begin{proposition}
    Si \( X\) et \( Y\) sont deux variables aléatoires, nous avons
    \begin{equation}
        E(X+Y)=E(X)+E(Y).
    \end{equation}
    Ce résultat est encore valide si \( X\) et \( Y\) ne sont pas indépendantes.
\end{proposition}

\begin{proof}
    Nous avons le calcul suivant :
    \begin{subequations}
        \begin{align}
            E(X+Y)&=\int_{\eR}xf_{X+Y}(x)dx\\
            &=\int_{\eR}x\int_{\eR}f_Y(x-t)f_X(t)dtdx\\
            &=\int_{\eR}f_X(t)\underbrace{\int_{\eR}xf_Y(x-t)dx}_{=E(Y)+t}\,dt\\
            &=\int_{\eR}f_X(t)\big( E(Y)+t \big)dt\\
            &=E(Y)+\int_{\eR}ff_X(t)dt\\
            &=E(Y)+E(X)
        \end{align}
    \end{subequations}
    où nous avons utilisé la proposition \ref{EqdensitesooemXYint} et le fait que l'intégrale sur \( \eR\) d'une densité vaut \( 1\).
\end{proof}

Une application de l'inégalité de Hölder est la suivante. Si \( X\) et \( Y\) sont des variables aléatoires intégrables alors
\begin{equation}
    E(XY)\leq E(X^2)^{1/2}E(Y^2)^{1/2}.
\end{equation}
En effet
\begin{equation}    \label{EqEXYleqXdYdNormHolder}
    E(XY)\leq \| XY \|_{L^1(\Omega)}\leq \| X \|_{L^2(\Omega)}\| Y \|_{L^2(\Omega)}.
\end{equation}

%---------------------------------------------------------------------------------------------------------------------------
\subsection{Espérance conditionnelle}
%---------------------------------------------------------------------------------------------------------------------------


\begin{theorem}
    Soit un espace de probabilité \( (\Omega,\tribA,P)\) et une variable aléatoire intégrable \( X\colon \Omega\to \eR\). Pour chaque sous tribu \( \tribF\) de \( \tribA\), il existe une (presque partout) unique variable aléatoire \( Y\colon \Omega\to \eR\) telle que
    \begin{enumerate}
        \item
            \( Y\) est \( \tribF\)-mesurable
        \item
            \( Y\) est \( P_{\tribF}\)-intégrable
        \item
            pour tout \( A\in\tribF\),
            \begin{equation}
                \int_{A}XdP=\int_A YdP.
            \end{equation}
    \end{enumerate}
    Cette variable aléatoire sera notée \( E(X|\tribF)\)\nomenclature[P]{\( E(X|\tribF)\)}{Espérance conditionnelle de \( X\) sachant \( \tribF\)} pour des raisons qui apparaîtront plus tard.
\end{theorem}

\begin{remark}
    Prendre \( Y=X\) ne fonctionne pas parce qu'en général si \( \mO\) est mesurable dans \( \eR\), alors \( X^{-1}(\mO)\) est dans la tribu \( \tribA\), mais n'est pas automatiquement dans la tribu \( \tribF\).
\end{remark}

\begin{proof}
    Cette preuve provient en bonne partie de \cite{ProbCOndutetz}.

    \begin{description}
        \item[Unicité] Si \( Y_1\) et \( Y_2\) vérifient tous les deux les conditions, l'ensemble \( \{ Y_1<Y_2 \}\) est un élément de \( \tribF\) et nous avons
            \begin{equation}
                \int_{\{ Y_1<Y_2 \}}X=\int_{Y_1<Y_2}Y_1=\int_{Y_1<Y_2}Y_2.
            \end{equation}
            En particulier nous avons \( \int_{\{ Y_1<Y_2 \}}(Y_1-Y_2)=0\) et donc
            \begin{equation}
                (Y_1-Y_2)\mtu_{Y_1-Y_2}=0
            \end{equation}
            presque partout. Le corollaire \ref{CorjLYiSm} montre alors que \( Y_1-Y_2\geq 0\) presque partout. De la même manière, l'ensemble \( \{ Y_2<Y_1 \}\) est dans \( \tribF\) et nous trouvons que \( Y_2-Y_1\geq 0\) presque partout. Par conséquent \( Y_1=Y_2\) presque partout.
        \item[Existence dans le cas de carré intégrable]

            Nous supposons que \( X\in L^2(\Omega,\tribA,P)\) et nous considérons \( K\), le sous ensemble de \( L^2(\Omega,\tribA,P)\) des fonctions \( \tribF\)-mesurables. Le théorème des projections \ref{ThoProjOrthuzcYkz} nous indique que
            \begin{equation}
                L^2(\Omega,\tribA,P)=K\oplus K^{\perp}
            \end{equation}
            par la décomposition \( X=\pr_{K}X+(X-\pr_KX)\). La variable aléatoire \( Y=\pr_KX\) a les propriétés d'être \( \tribF\)-mesurable et \( \langle Y-X, Z\rangle =0\) pour tout \( Z\in K\). Soit \( A\in\tribF\), si nous considérons \( Z=\mtu_A\), la dernière condition signifie que
            \begin{equation}
                \int_{\Omega}X\mtu_A=\int_{\Omega}Y\mtu_A,
            \end{equation}
            ou encore
            \begin{equation}
                \int_AY=\int_AX.
            \end{equation}
            La variable aléatoire \( Y=\pr_KX\) répond donc à la question dans le cas où \( X\in L^2(\Omega,\tribF,P)\).

        \item[Existence en général] 

            Nous considérons maintenant que \( X\in L^1(\Omega,\tribA,P)\). Quitte à décomposer \( X\) en deux fonctions positives \( X_+\) et \( X_-\) telles que \( X=X_++X_-\), nous pouvons supposer que \( X\) est positive. Par hypothèse \( X\in L^1(\Omega,\tribA,P)\); pour chaque \( n\in\eN\) nous posons
            \begin{equation}
                X_n(\omega)=\min\{ X(\omega),n \}.
            \end{equation}
            Étant donné que la mesure \( P\) est une mesure de probabilité, les constantes sont intégrables et \( X_n\in L^2(\Omega,\tribA,P)\). De plus la suite \( (X_n)\) est croissante et
            \begin{equation}
                \lim_{n\to \infty} X_n(\omega)=X(\omega).
            \end{equation}

            Si nous notons encore \( K\) l'ensemble des variables aléatoires dans \( L^2(\Omega,\tribA,P)\) qui sont \( \tribF\)-mesurables, pour chaque \( n\) nous avons donc la variable aléatoire \( Y_n=\pr_KX_n\) qui est \( \tribF\)-mesurable et telle que
            \begin{equation}
                \int_A X_n=\int_AY_n
            \end{equation}
            pour tout \( A\in\tribF\). Nous voudrions prouver que la variable aléatoire \( Y=\lim_nY_n\) existe et est la solution au problème. D'abord la suite \( (X_n)\) étant croissante, la suite \( Y_n\) est également croissante.

    \end{description}
    <++>
\end{proof}
<++>


Soit \( \tribA\) une sous tribu de la tribu \( \tribF\) et \( X\), une variable aléatoire \( \tribF\)-mesurable et intégrable. Alors nous notons 
\begin{equation}
    E(X|\tribA)
\end{equation}
l'unique variable aléatoire à être \( \tribA\)-mesurable telle que pour toute variable aléatoire \( \tribA\)-mesurable \( Y\),
\begin{equation}
    E\big( E(X|\tribA)Y \big)=E(XY).
\end{equation}


%---------------------------------------------------------------------------------------------------------------------------
\subsection{Variance}
%---------------------------------------------------------------------------------------------------------------------------

Si \( X\in L^2(\Omega,\tribA,P)\) alors nous définissons la \defe{variance}{variance} de \( X\) par
\begin{equation}
    \Var(X)=E\big( [X-E(X)]^2 \big).
\end{equation}

\begin{proposition}     \label{PrropVarAlterfrom}
    La variance de la variable aléatoire \( X\) peut être exprimée par la formule
    \begin{equation}
        \Var(X)=E(X^2)-[E(X)]^2
    \end{equation}
    où \( X^2=X\cdot X\) et \( E(X)^2=\) sont des produits scalaires dans \( \eR^d\).
\end{proposition}

\begin{proof}
    De façon explicite, nous avons
    \begin{equation}
        E\big( [X-E(X)]^2 \big)=\int_{\Omega}\big( X(\omega)-E(X) \big)\cdot\big( X(\omega)-E(X) \big)dP(\omega)
    \end{equation}
    où \( E(X)\in\eR^d\) est une constante. En développant le produit scalaire nous avons
    \begin{subequations}
        \begin{align}
            E\big( [X-E(X)]^2 \big)&=E\big( X^2-2X\cdot E(X)+E(X)^2 \big)\\
            &=E(X^2)-2E(X)^2+E(X)^2\\
            &=E(X^2)-E(X)^2.
        \end{align}
    \end{subequations}
\end{proof}


Nous définissons l'\defe{écart-type}{écart-type} de \( X\) par
\begin{equation}
    \sigma_X=\sqrt{\Var(X)}.
\end{equation}
En d'autres termes,
\begin{equation}
    \sigma_X=\| X-E(X) \|_{L^2}.
\end{equation}
On définit encore la \defe{moyenne quadratique}{moyenne!quadratique} de \( X\) par
\begin{equation}
    \| X \|_{L^2}=\big[ E(X^2) \big]^{1/2}.
\end{equation}

La variable aléatoire 
\begin{equation}
    \bar V_n=\frac{1}{ n }\sum_i(X_i-\bar X_n)^2
\end{equation}
est la \defe{variance empirique}{variance!empirique} de l'échantillon \( (X_i)\).

%---------------------------------------------------------------------------------------------------------------------------
\subsection{Covariance}
%---------------------------------------------------------------------------------------------------------------------------

Soient \( X\) et \( Y\), deux variables aléatoires réelles. Leur \defe{covariance}{covariance} est définie par
\begin{equation}
    \Cov(X,Y)=E\Big[ \big( X-E(X) \big)\big( Y-E(Y) \big) \Big].
\end{equation}
L'idée est que la covariance devient grande si \( X\) et \( Y\) s'écartent de leurs moyennes dans le même sens. Il existe une formule alternative :
\begin{equation}
    \Cov(X,Y)=E(XY)-E(X)E(Y)
\end{equation}

En ce qui concerne les dimensions plus hautes, si \( X\colon \Omega\to \eR^d\) est un vecteur aléatoire de carré intégrable, nous définissons
\begin{equation}
    \Cov(X)=E\Big[ \big( X-E(X) \big)^t(X-E(X)) \Big].
\end{equation}

\begin{proposition}     \label{PropoVarXpYCov}
    Si \( X\) et \( Y\) sont deux variables aléatoires non spécialement indépendantes, nous avons
    \begin{equation}
        \Var(X+Y)=\Var(X)+\Var(Y)+2\Cov(X+Y).
    \end{equation}
\end{proposition}

\begin{proof}
    Il s'agit d'un calcul en partant de
    \begin{equation}
        \begin{aligned}[]
            \Var(X+Y)&=E\big( (X+Y)^2 \big)-E(X+Y)^2\\
            &=E(X^2)+E(Y^2)+2E(XY)\\
            &\quad+\big( E(X)+E(Y) \big)^2-2E(X)^2-2E(X)E(Y)\\
            &\quad-2 E(Y)E(X)-2E(Y)^2.
        \end{aligned}
    \end{equation}
    À partir d'ici il s'agit de recombiner tous les termes pour former la formule annoncée.
\end{proof}

Plus généralement nous avons la formule
\begin{equation}
    \Var(\sum_i X_i)=\sum_i\Var(X_i)+2\sum_{1\leq i< j\leq n}\Cov(X_i,X_j).
\end{equation}


%---------------------------------------------------------------------------------------------------------------------------
\subsection{Fonction de répartition et fonction caractéristique}
%---------------------------------------------------------------------------------------------------------------------------

Si \( X\) est une variable aléatoire réelle, nous définissons sa \defe{fonction de répartition}{fonction!de répartition} par
\begin{equation}
    \begin{aligned}
        F_X\colon \eR&\to \mathopen[ 0 , 1 \mathclose] \\
        F_X(x)&=P(X\leq x). 
    \end{aligned}
\end{equation}


La \defe{fonction caractéristique}{fonction!caractéristique} de la variable aléatoire \( X\colon \Omega\to \eR\) est la fonction réelle définie par
\begin{equation}
    \Phi_X(t)=E( e^{itX}).
\end{equation}
Une autre façon d'écrire la définition est
\begin{equation}
    \Phi_X(t)=\int_{\eR} e^{itX}dP_X(x),
\end{equation}
ou encore, si \( X\) a une densité \( f_X\),
\begin{equation}        \label{EqFnCaractfncadens}
    \Phi_X(t)=\int_{\eR} e^{itx}f_X(x)dx
\end{equation}
Nous reconnaissons la transformée de Fourier :
\begin{equation}
    \Phi_X(t)=(TF\,f_X)(-t/2\pi).
\end{equation}

La proposition suivante se déduit en utilisant le théorème de dérivation sous l'intégrale \ref{ThoDerSousIntegrale}.
\begin{proposition}     \label{PropDerFnCaract}
    Soit $X$ une variable aléatoire qui accepte un moment d'ordre \( r\geq 1\). Alors la fonction de répartition \( \Phi_X\) est \( r\) fois continument dérivable et
    \begin{equation}
        \Phi_X^{(r)}(t)=E\big( (iX)^r e^{itX} \big).
    \end{equation}
\end{proposition}

\begin{proof}
    Nous étudions la fonction
    \begin{equation}
        \Phi(t)=\int_{\Omega} e^{itX(\omega)}dP(\omega).
    \end{equation}
    Nous considérons la fonction
    \begin{equation}
        \begin{aligned}
            f\colon \eR\times\Omega&\to \eR \\
            (t,\omega)&\mapsto  e^{itX(\omega)}. 
        \end{aligned}
    \end{equation}
    et nous regardons si ce contexte vérifie les hypothèses du théorème \ref{ThoDerSousIntegrale}.
    \begin{enumerate}
        \item
            Étant donné que \( X\) est mesurable, \( f\) sera mesurable.
        \item
            La fonction \( t\mapsto e^{itX(\omega)}\) est absolument continue pour chaque \( \omega\).
        \item
            Note : par rapport aux notations du théorème \ref{ThoDerSousIntegrale}, nous avons ici \( A=\eR\). Prenons donc un intervalle (compact) \( \mathopen[ a , b \mathclose]\subset\eR\) et calculons
            \begin{equation}        \label{EqfpfpttoiXieitXo}
                \frac{ \partial f }{ \partial t }(t,\omega)=iX(\omega) e^{itX},
            \end{equation}
            et
            \begin{equation}
                \int_a^b\int_{\Omega}\left| iX(\omega) e^{itX(\omega)} \right| d\omega\,dt=\int_a^b\int_{\Omega}| X(\omega) |d\omega\,dt.
            \end{equation}
            Par hypothèse \( X\) accepte un moment d'ordre \( 1\), de sorte que l'intégrale par rapport à \( \omega\) converge vers un nombre qui ne dépend pas de \( t\). L'intégrale sur \( t\) ne pose alors aucun problèmes.
    \end{enumerate}
    Par conséquent nous pouvons effectuer la première dérivation :
    \begin{equation}
        \Phi'(t)=\frac{ d\Phi }{ dt }(t)=\int_{\Omega}iX(\omega) e^{itX(\omega)}d\omega=E(iX e^{itX})
    \end{equation}
    et la fonction \( \Phi'\) est absolument continue. Ce dernier point est important parce que c'est lui qui permet de faire la récurrence et passer à l'ordre deux.

    Le résultat ressort alors en dérivant successivement l'expression \eqref{EqfpfpttoiXieitXo}.
\end{proof}

\begin{example}
    Sachant la fonction caractéristique de \( X\), nous pouvons calculer les moments. Par exemple
    \begin{equation}
        E(X^2)=\Phi''_X(0).
    \end{equation}
\end{example}

\begin{theorem}
    Si \( \Phi_X=\Phi_Y\), alors \( P_X=P_Y\).
\end{theorem}
\begin{proof}
    No proof.
\end{proof}
Notons que cela n'implique pas que \( X=Y\). En effet \( X\) et \( Y\) peuvent même être définis sur des espaces probabilisés différents.

Dans le cas d'une variable aléatoire vectorielle, nous définissons \( \Phi_X\colon \eR^d\to \eR\) par
\begin{equation}
    \Phi_X(v)=E\big(  e^{i\langle v, X\rangle } \big)
\end{equation}
%---------------------------------------------------------------------------------------------------------------------------
\subsection{Loi d'une variable aléatoire}
%---------------------------------------------------------------------------------------------------------------------------

La \defe{loi}{loi d'une variable aléatoire} de la variable aléatoire \( X\), notée \( P_X\) est la mesure image de \( P\) par \( X\), c'est à dire
\begin{equation}
    P_X(B)=P(X\in B)
\end{equation}
pour tout borélien \( B\subset\eR^d\). Note :
\begin{equation}
    P(X\in B)=P\big( \{ \omega\in\Omega\tq X(\omega)\in B \} \big)=P\big( X^{-1}(B) \big).
\end{equation}
En particulier \( P_X\) est une mesure de probabilité sur \( \eR^d\) parce que 
\begin{equation}
    P_X(\eR^d)=P(\Omega)=1.
\end{equation}
Si \( Q\) est une mesure de probabilité sur \( \eR^d\), nous notons \( X\sim Q\) si \( P_X=Q\). Nous disons alors que «\( X\) suit la loi \( Q\)».

La proposition suivante permet de calculer en pratique les intégrales qui définissent par exemple l'espérance mathématique d'une variable aléatoire.
\begin{proposition}     \label{PropintdPintdPXeR}
    Si \( f\colon \eR^d\to \bar\eR\) est borélienne et si \( X\) est une variable aléatoire, alors
    \begin{equation}
        \int_{\eR^d}f(x)dP_X(x)=\int_{\Omega}f\big( X(\omega) \big)dP(\omega)=E(f\circ X).
    \end{equation}
\end{proposition}

\begin{proof}
    No proof.
\end{proof}

En utilisant cette proposition nous trouvons une formule pratique pour l'espérance d'une variable aléatoire réelle:
\begin{equation}
    E(X)=\int_{\Omega}X(\omega)dP(\omega)=\int_{\eR}xdP_X(x),
\end{equation}
en vertu de la proposition \ref{PropintdPintdPXeR} appliquée à la fonction \( f(x)=x\).

\begin{proposition}
    Une variable aléatoire réelle \( X\) est intégrable si et seulement si \( P(x=\pm\infty)=0\) et
    \begin{equation}
        \int_{\eR}| x |dP_X(x)<\infty.
    \end{equation}
\end{proposition}

Le lien entre la densité \( f_X\) de la variable aléatoire \( X\) et sa loi est
\begin{equation}
    P_X(A)=\int_Af_X(x)dx
\end{equation}
pour tout ensemble mesurable \( A\subset\eR\). Le lien entre la mesure de Lebesgue et celle de la loi de \( X\) est alors donné par
\begin{equation}
    dP_X(x)=f_X(x)dx.
\end{equation}
En particulier l'espérance de \( X\) peut être calculée à partir de sa densité via la formule
\begin{equation}        \label{EqEspDensform}
    E(X)=\int_{\eR}xdP_X(x)=\int_{\eR}x f_X(x)dx.
\end{equation}




%---------------------------------------------------------------------------------------------------------------------------
\subsection{Lois conjointes et indépendance}
%---------------------------------------------------------------------------------------------------------------------------

\begin{definition}
    Deux événements \( A\) et \( B\) sont dits \defe{indépendants}{indépendance} si
    \begin{equation}
        P(A\cap B)=P(A)P(B).
    \end{equation}
\end{definition}
Si nous considérons \( n\) variables aléatoires réelles \( X_1,\ldots,X_n\colon\Omega\to\eR\), la loi du \( n\)-uplet \( X=(X_1,\ldots,X_n)\) est une variable aléatoire \( X\colon \Omega\to \eR^n\) appelée la \defe{loi conjointe}{loi!conjointe} des lois \( X_i\). Dans ce cas, les variables aléatoires \( X_i\) elles-mêmes sont dites lois \defe{marginales}{loi!marginale} de \( X\).

\begin{proposition}     \label{PropPXXXPXPXPX}
    Les variables aléatoires \( \{ X_i \}\) sont indépendantes si et seulement si
    \begin{equation}
        P_{(X_1,\ldots,X_n)}=P_{X_1}\otimes\ldots\otimes P_{X_n}.
    \end{equation}
\end{proposition}

\begin{definition}      \label{DefFonrepConj}
    Soient \( \{ X_i \}_{1\leq i\leq n}\) des variables aléatoires réelles (pas spécialement indépendantes). La \defe{densité conjointe}{densité!conjointe} de \( X_1\),\ldots,\( X_n\) est la fonction \( f\colon \eR^n\to \eR\) qui satisfait
    \begin{enumerate}
        \item
            \( f(x_1,\ldots,x_n)\geq 0\) pour tout \( (x_1,\ldots,x_n)\in\eR^n\),
        \item
            \( \int_{\eR^n}f=1\),
        \item       \label{ItemDefFonrepConjiii}
            pour tout \( A_i\subset\eR \) nous avons
            \begin{equation}
                P(\bigcap_{i=1}^n X_i\in A_i)=\int_{\prod_i A_i}f(x_1,\ldots,x_n)dx_1\ldots dx_n.
            \end{equation}
    \end{enumerate}
\end{definition}

\begin{proposition}     \label{PropDensiteConjIndep}
    Si les variables aléatoires \( X_1\),\ldots \( X_n\) sont indépendantes et ont des densités \( f_{X_1}\),\ldots,\( f_{X_n}\), alors la variable aléatoire conjointe \( X=(X_1,\ldots,X_n)\) a pour densité conjointe la fonction
    \begin{equation}
        f_X(x_1,\ldots,x_n)=f_{X_1}(x_1)\ldots f_{X_n}(x_n).
    \end{equation}
\end{proposition}

\begin{proof}
    En partant de la définition de l'indépendance et de la fonction de densité conjointe, ainsi qu'en utilisant le théorème de Fubini,
    \begin{equation}
        \begin{aligned}[]
            \int_{A_1\times \ldots\times A_n}f_X(x_1,\ldots,x_n)dx_1\ldots dx_n&=
            P(X_1\in A_1,\ldots,X_n\in A_n)\\
            &=P(X_1\in A_1)\ldots P(X_n\in A_n)\\
            &=\left( \int_{A_1}f_{X_1}(x_1)dx_1 \right)\ldots\left( \int_{A_n}f_{X_n}(x_n)dx_n \right)\\
            &=\int_{A_1\times\ldots\times A_n}f_{X_1}(x_1)\ldots f_{X_n}(x_n)dx_1\ldots dx_n.
        \end{aligned}
    \end{equation}
    La fonction \( (x_1,\ldots,x_n)\mapsto f_{X_1}(x_1)\ldots f_{X_n}(x_n)\) vérifie donc la condition \ref{ItemDefFonrepConjiii} de la définition \ref{DefFonrepConj}. La vérification des autres conditions est immédiate.
\end{proof}


La proposition suivante\cite{ProbaDanielLi} provient du fait que la mesure d'une loi conjointe est le produit des mesures lorsque les variables aléatoires sont indépendantes (proposition \ref{PropPXXXPXPXPX}).
\begin{proposition}
    Si les variables aléatoires réelles \( X_1\),\ldots,\( X_n\) sont intégrables et indépendantes, alors leur produit est intégrable et l'espérance du produit est égal au produit des espérances :
    \begin{equation}
        E(X_1\cdots X_n)=E(X_1)\ldots E(X_n).
    \end{equation}
\end{proposition}

%---------------------------------------------------------------------------------------------------------------------------
\subsection{Somme et produit de variables aléatoires indépendantes}
%---------------------------------------------------------------------------------------------------------------------------
\label{subsecscnvommevariablsindep}

Le cas plus général de variable aléatoires non indépendantes peut être trouvé dans \cite{Marazzi}.

Soient \( X\) et \( Y\), deux variables aléatoires réelles indépendantes. Nous voudrions étudier la loi de la variable aléatoire \( S=X+Y\). Nous commençons par calculer la fonction de répartition en utilisant le résultat de la proposition \ref{PropDensiteConjIndep} :
\begin{subequations}
    \begin{align}
        F_{X+Y}(z)=P(X+Y\leq z)&=\int_{x+y\leq z}f_{X,Y}(x,y)dx\,dy\\
        &=\int_{-\infty}^{\infty}dx\int_{-\infty}^{z-x}dyf_X(x)f_Y(y)\\
        &=\int_{\eR}\left( \int_{-\infty}^{z-x}f_Y(y)dy \right)f_X(x)dx\\
        &=\int_{\eR}F_Y(z-x)f_X(x)dx.
    \end{align}
\end{subequations}
Pour calculer la fonction de densité de \( S\), nous dérivons la fonction de répartition :
\begin{subequations}
    \begin{align}
        f_{X+Y}(z)&=\frac{ d F_{X+Y} }{ d z }(z)\\
        &=\int_{\eR}f_Y(z-x)f_X(x)dx,
    \end{align}
\end{subequations}
ce qui nous amène à dire que la densité de la somme est le produit de convolution\index{convolution} des densités :
\begin{equation}        \label{EqdensitesooemXYint}
    f_{X+Y}(x)=\int_{\eR}f_Y(x-t)f_X(t)dt,
\end{equation}
ou encore \( f_{X+Y}=f_X\star f_Y\).

Notez que nous avons passé sous le silence la difficulté d'inverser la dérivée et l'intégrale. Un exemple sera donné au point \ref{subsecPoissonetexpo}.


\begin{lemma}       \label{LemEXYEXEYprodindep}
    Soient \( X\) et \( Y\), deux variables aléatoires indépendantes et identiquement distribuées. Alors
    \begin{equation}
        E(XY)=E(X)E(Y).
    \end{equation}
\end{lemma}

\begin{proof}
    Par indépendance, fonction de densité conjointe de \( X\) et \( Y\) vaut \( f_{X,Y}=f_Xf_Y\). Par conséquent l'utilisation de Fubini entraine
    \begin{equation}
        E(XY)=\int_{\eR\times\eR}xyf_{X,Y}(x,y)dxdy=E(X)E(Y).
    \end{equation}
\end{proof}

\begin{lemma}   \label{LemVarXpYsmindep}
    Soit \( X\) et \( Y\) deux variables aléatoires indépendantes et identiquement distribuées. Alors
    \begin{equation}
        \Var(X+Y)=\Var(X)+\Var(Y).
    \end{equation}
\end{lemma}

\begin{proof}
    Par définition, \( \Var(X+Y)=E\big( [X+Y-E(X)-E(Y)]^2 \big)\). En développant le carré et en utilisant le lemme \ref{LemEXYEXEYprodindep},
    \begin{equation}
        \Var(X+Y)=E(X^2)-E(X)^2+E(Y^2)-E(Y)^2=\Var(X)+\Var(Y).
    \end{equation}
\end{proof}

%+++++++++++++++++++++++++++++++++++++++++++++++++++++++++++++++++++++++++++++++++++++++++++++++++++++++++++++++++++++++++++
\section{Événements}
%+++++++++++++++++++++++++++++++++++++++++++++++++++++++++++++++++++++++++++++++++++++++++++++++++++++++++++++++++++++++++++

\begin{lemma}       \label{LemEXYEXEYindep}\label{PropVarPropnnlin}
    Si \( X\) est une variable aléatoire,
    \begin{enumerate}
        \item
            $\Var(ax)=a^2\Var(X)$ pour tout \( a\in\eR\);
        \item
            Si de plus \( Y\) est une variable aléatoire indépendante de \( X\), alors $\Var(X+Y)=\Var(X)+\Var(Y)$.
    \end{enumerate}
\end{lemma}

\begin{proof}
    Nous avons
    \begin{subequations}
        \begin{align}
            \Var(X+Y)&=E(X^2+Y^2+2XY)-\big( E(X)+E(Y) \big)^2\\
            &=E(X^2)+E(Y^2)+2E(XY)-E(X)^2-E(Y)^2E(X)E(Y).
        \end{align}
    \end{subequations}
    Étant donné que \( X\) et \( Y\) sont indépendantes nous avons \( E(XY)=E(X)E(Y)\).
\end{proof}


Si les \( X_1,\ldots,X_n\) sont des variables aléatoires on considère la \defe{moyenne empirique}{moyenne!empirique}
\begin{equation}
    \bar X_n=\frac{ X_1+\ldots+X_n }{ n }.
\end{equation}

%+++++++++++++++++++++++++++++++++++++++++++++++++++++++++++++++++++++++++++++++++++++++++++++++++++++++++++++++++++++++++++
\section{Convergence}
%+++++++++++++++++++++++++++++++++++++++++++++++++++++++++++++++++++++++++++++++++++++++++++++++++++++++++++++++++++++++++++

Soient \( X_i\) des variables aléatoires réelles définies sur le même espace de probabilité \( (\Omega,\tribA,P)\). Nous disons que \( X_i\) converge \defe{presque surement}{convergence!presque surement} vers la variable aléatoire \( X\) et nous notons 
\begin{equation}
    X_n\stackrel{p.s.}{\longrightarrow}X
\end{equation}
si
\begin{equation}
    P\big( \{ \omega\in\Omega\tq\,X_n(\omega)\to X(\omega) \} \big)=1
\end{equation}
où la convergence \( X_n(\omega)\to X(\omega)\) est la convergence usuelle dans \( \eR\).

\begin{lemma}
    Nous avons \( X_n\stackrel{p.s.}{\longrightarrow}X\) si et seulement si il existe un événement \( A\in\tribA\) tel que \( P(A)=1\) et tel que \( X_n(\omega)\to X(\omega)\) pour tout \( \omega\in A\).
\end{lemma}

Nous disons que les variables aléatoires réelles \( X_n\) convergent \defe{en probabilité}{convergence!en probabilité} vers la variable aléatoire \( X\) si pour tout \( \eta>0\), on a
\begin{equation}
    P\big( | X_n-X |\geq \eta \big)\to 0,
\end{equation}
et on note
\begin{equation}
    X_n\stackrel{P}{\longrightarrow}X.
\end{equation}

Nous disons que \( X_n\) converge vers \( X\) \defe{en loi}{convergence!en loi} vers la variable aléatoire \( X\) et nous notons
\begin{equation}
    X_n\stackrel{\hL}{\longrightarrow}X
\end{equation}
si pour toute fonction continue et bornée \( g\) nous avons
\begin{equation}
    E\big( g(X_n) \big)\to E\big( g(X) \big)=\int gdP_X.
\end{equation}


\begin{proposition}     \label{PrpopCaractCvL}
    Deux autres caractérisations de la convergence en loi.
    \begin{enumerate}
        \item
            Nous avons \( X_n\stackrel{\hL}{\longrightarrow}X\) si et seulement si
            \begin{equation}
                \Phi_{X_n}(v)\to\Phi_X(v)
            \end{equation}
            pour tout \( v\in\eR^d\). Ici \( \Phi_X\) est la fonction caractéristique de \( X\).
        \item
            Dans la définition de la convergence en loi nous pouvons indifféremment utiliser les fonctions continues et bornées, les fonctions continues à support compact ou les fonctions bornées uniformément continues.
    \end{enumerate}
\end{proposition}

\begin{proposition}
    La convergence presque sure entraine à la fois la convergence en probabilité et celle en loi. La convergence en probabilité entraine la convergence en loi.
\end{proposition}
La convergence en loi n'implique pas la convergence en probabilité, et par conséquent pas non plus la convergence presque certaine.

Dans le cas particulier \( d=1\) nous avons quelque critères supplémentaires. 

\begin{proposition}     \label{PropoFnrepCvL}
    Supposons que les variables aléatoires \( X_n\) soient réelles, et notons \( F_n\) la fonction de répartition de \( X_n\). Si \( F_n(x)\to F(x)\) pour tout \( x\) dans l'ensemble des points de continuité de \( F\), alors \( X_n\stackrel{\hL}{\longrightarrow}X\).
\end{proposition}

\begin{proposition}
    Si les \( X_n\) sont des variables aléatoires réelles positives, et si \( X\) est une variable aléatoire positive, alors \( X_n\stackrel{\hL}{\longrightarrow}X\) si les transformées de Laplace des fonctions de répartition convergent ponctuellement, c'est à dire si
    \begin{equation}
        E\big(  e^{-\alpha X_n} \big)\to E\big(  e^{-\alpha X} \big)
    \end{equation}
    pour tout \( \alpha\geq 0\).
\end{proposition}

\begin{proposition}
    Si les \( X_n\) et \( X\) sont des variables aléatoires réelles discrètes à valeurs dans \( \{ x_0,x_1,\ldots \}\) alors \( X_n\stackrel{\hL}{\longrightarrow} X\) si et seulement si
    \begin{equation}
        P(X_n=x_k)\to P(X=x_k)
    \end{equation}
    pour tout \( k\in\eN\).
\end{proposition}

\begin{proposition}     \label{PropXncvXFXcvFxt}
    Soient \( X_n\) et \( X\) des variables aléatoires réelles. Nous avons
    \begin{equation}
        X_n\stackrel{\hL}{\longrightarrow}X
    \end{equation}
    si et seulement si pour tout \( t\) où \( F_X\) est continue,
    \begin{equation}
        \lim_{n\to \infty} F_{X_n}(t)=F_X(t).
    \end{equation}
\end{proposition}
Pour une preuve, voir \cite{CourgGudRennes}.

\begin{proposition}     \label{PropCvLfcvPsicst}
    Soit \( X_n\) une suite de variables aléatoires \( \Omega\to \eR^d\) et \( a\in\eR^d\). Si \( X_n\stackrel{\hL}{\longrightarrow}\), alors
    \begin{equation}
        X_n\stackrel{P}{\longrightarrow}a.
    \end{equation}
\end{proposition}

\begin{proof}
    Cette preuve provient de \cite{CourgGudRennes}. Quitte à passer aux composantes, nous pouvons supposer que \( d=1\). Nous avons
    \begin{subequations}
        \begin{align}
           P\big( | X_n-a |>\eta \big)&=P(X_n-a>\eta)+P(-X_n-a>\eta)\\
           &=P(X_n>\eta+a)+1-P(X_n>a-\eta)  \label{EqPXngeqetaapumP}\\
           &=1-F_{X_n}(\eta+a)+F_{X_n}(a-\eta).
        \end{align}
    \end{subequations}
    Nous allons utiliser la proposition \ref{PropXncvXFXcvFxt}. La fonction de répartition de la variable aléatoire constante \( X=a\) est donnée par
    \begin{equation}
        F_a(t)=P(a<t)=\caract{\eR^+}(t-a).
    \end{equation}
    Par conséquent, la convergence en loi \( X_n\stackrel{\hL}{\longrightarrow}a\) nous montre que
    \begin{equation}
        F_{X_n}(t)\to \caract_{\eR^+}(t-a)
    \end{equation}
    pour tout \( t\neq a\) parce que \( t=0\) est un point de discontinuité de \( \caract_{\eR^+}\). Nous avons par conséquent
    \begin{equation}
        P\big( | X_n-a |>\eta \big)=1-\caract_{\eR^+}(\eta)+\caract_{\eR^+}(-\eta)=0
    \end{equation}
    parce que \( \eta>0\).
\end{proof}

\begin{probleme}
    Dans \cite{CourgGudRennes}, l'équation \eqref{EqPXngeqetaapumP} arrive avec une inégalité. Pourquoi ?
\end{probleme}

\begin{lemma}[Slutsky]\index{lemme!de Slutsky}
    Soient \( X_n\) et \( Y_n\) des suites de variables aléatoires réelles telles que
    \begin{equation}
        \begin{aligned}[]
            X_n&\stackrel{\hL}{\longrightarrow} X&\text{et}&&Y_n&\stackrel{\hL}{\longrightarrow}a\in\eR.
        \end{aligned}
    \end{equation}
    Alors \( (X_n,Y_n)\stackrel{\hL}{\longrightarrow}(X,a)\).
\end{lemma}

\begin{proof}
    La preuve qui suit provient de \cite{MPSmaitrise}. Étant donné que \( Y_n\stackrel{\hL}{\longrightarrow} a\), nous avons \( Y_n\stackrel{P}{\longrightarrow} a\) par la proposition \ref{PropCvLfcvPsicst}. Soit une fonction \(f\colon \eR^2\to \eR^2 \); nous devons prouver que
    \begin{equation}
        E\big( f(X_n,Y_n) \big)\to E\big( f(X,a) \big).
    \end{equation}
    Soit \( \epsilon>0\). Nous avons
    \begin{equation}    \label{EqEparXnYnfXa}
        E\big( \| f(X_n,Y_n)-f(X,a) \| \big)\leq E\big(  \| f(X_n,Y_n)-f(X_n,a) \|  \big)+E\big(   \| f(X_n,a)-f(X,a) \|  \big).
    \end{equation}
    La fonction\( g(t)=f(t,a)\) étant continue et bornée, la convergence en loi \( X_n\stackrel{\hL}{\longrightarrow}X\) donne
    \begin{equation}
        E\big( \| f(X_n,a)-f(X,a) \| \big)\to 0.
    \end{equation}
    Étudions à présent le premier terme du membre de droite de \eqref{EqEparXnYnfXa}. Pour tout \( \eta> 0\) et toute variables aléatoires \( Z\) et \( Z'\) nous avons
    \begin{equation}
        E(Z)=E(Z\caract_{| Z' |<\eta})+E(Z\caract_{| Z' |\geq \eta}).
    \end{equation}
    Nous décomposons donc le premier terme de \eqref{EqEparXnYnfXa} en
    \begin{equation}    \label{EqEXnADecomsecvolta}
        \begin{aligned}[]
            E\big( \| f(X_n,Y_n)-f(X_n,a) \| \big)&=E\big( \| f(X_n,Y_n)-f(X_n,a) \|\caract_{| Y_n-a |<\eta} \big)\\
            &\quad+E\big( \| f(X_n,Y_n)-f(X_n,a) \|\caract_{| Y_n-a |\geq\eta} \big).
        \end{aligned}
    \end{equation}
    Choisissons maintenant une valeur de \( \eta\) telle que
    \begin{equation}
        | (x,y)-(x',y') |<\eta\Rightarrow| f(x,y)-f(x',y') |\leq \epsilon.
    \end{equation}
    Un tel \( \eta\) existe par l'uniforme continuité de \( f\). Dans le premier terme, \( | Y_n-a |<\eta\), par conséquent
    \begin{equation}
        \| (X_n,Y_n)-(X_n,a) \|=| Y_n-a |<\eta
    \end{equation}
    et donc
    \begin{equation}
        \| f(X_n,Y_n)-f(X_n,a) \|\leq \epsilon.
    \end{equation}
    Le premier terme devient donc
    \begin{equation}
        E\big( \| f(X_n,Y_n)-f(X_n,a) \|\caract_{| Y_n-a |<\eta} \big)\leq \epsilon E(\caract_{| Y_n-a |<\eta})\leq \epsilon
    \end{equation}
    parce que \( E(\caract_A)=P(A)\leq 1\). Pour le second terme de \eqref{EqEXnADecomsecvolta} nous effectuons la majoration
    \begin{equation}
        \| f(X_n,Y_n)-f(X_n,a) \|\leq 2\| f \|_{\infty}
    \end{equation}
    tandis que la convergence \( Y_n\stackrel{P}{\longrightarrow} a\) entraine 
    \begin{equation}
        P\big( | Y_n-a |\geq \eta \big).
    \end{equation}
\end{proof}

\begin{proposition}     \label{PropcvLsousfonc}
    Soient \( X_n\) des variables aléatoires telles que
    \begin{equation}
        X_n\stackrel{\hL}{\longrightarrow} X
    \end{equation}
    et \( h\), une fonction mesurable sur l'espace d'arrivée de \( X_n\). Soit \( C\) l'ensemble des points de continuité de \( h\). Alors si \( P(X\in C)=1\), nous avons 
    \begin{equation}
        h(X_n)\stackrel{\hL}{\longrightarrow}h(X).
    \end{equation}
\end{proposition}

\begin{lemma}[Borel-Cantelli]\index{lemme!de Borel-Cantelli}
    Soit \( (A_n)\) une suite d'événements (avec \( A_n\in\tribA\) pour tout \( n\)).
    \begin{enumerate}
        \item
            Si \( \sum_{n=0}^{\infty}P(A_n)\) converge, alors
            \begin{equation}
                P(A_n\,\text{i.s.})=0.
            \end{equation}
        \item
            Si la somme \( \sum_nP(A_n)\) diverge, et si de plus les \( A_i\) sont indépendants, alors
            \begin{equation}
                P(A_n\,\text{i.s.})=1.
            \end{equation}
    \end{enumerate}
\end{lemma}
La notation \( P(A_n\,\text{i.s.})\) signifie «infiniment souvent», c'est à dire
\begin{equation}
    P(A_n\,\text{i.s.})=P\big( \bigcap_{N\in\eN}\bigcup_{k\geq N}A_k \big)=P(\limsup A_n)
\end{equation}

Une façon de paraphraser le lemme de Borel-Cantelli est que nous avons l'alternative
\begin{equation}    \label{EqparaphrCantelli}
    P(\limsup A_n)=\begin{cases}
        0    &   \text{si $\sum_{n\geq 0}P(A_n)<\infty$}\\
        1    &    \text{sinon}.
    \end{cases}
\end{equation}

\begin{proposition}
    Soit \( X_n\), une suite de variables aléatoires et \( X\) une variable aléatoires. Si
    \begin{equation}
        \sum_nP\big( \| X_n-X \|>\eta \big)<\infty
    \end{equation}
    pour tout \( \epsilon\), alors \( X_n\stackrel{p.s.}{\longrightarrow}X\).
\end{proposition}

\begin{proof}
    Fixons \( \epsilon\) et considérons les événements \( A_{n}=\| X_n-X \|>\epsilon\). L'hypothèse dit que
    \begin{equation}
        \sum_nP(A_{n,\epsilon})<\infty
    \end{equation}
    et le lemme de Borel-Cantelli implique que
    \begin{equation}
        P(\limsup\| X_n-X \|>\epsilon)=0.
    \end{equation}
    Un élément \( \omega\) est dans \( \limsup A_n\) si il est contenu dans tous les \( A_n\), par conséquent, pour chaque \( \epsilon\) nous avons l'inclusion
    \begin{equation}
        \{ \omega\in\Omega\tq X_n(\omega)\to X(\omega) \}\subset\complement\limsup A_n.
    \end{equation}
    Nous pouvons aller plus loin et écrire
    \begin{equation}        \label{EqProbomOmXnmXforalleps}
        \{ \omega\in\Omega\tq X_n(\omega)\to X(\omega) \}=\complement\{ \omega\in\Omega\tq \| X_n-X \|>\epsilon,\forall\epsilon>0 \}.
    \end{equation}
    Or la probabilité de l'ensemble
    \begin{equation}
        \{ \omega\in\Omega\tq\| X_n-X \|>\epsilon \}
    \end{equation}
    est \( 0\) pour chaque \( \epsilon\), et par conséquent la probabilité du membre de droite de \eqref{EqProbomOmXnmXforalleps} est \( 1\).
\end{proof}

\begin{example}
    Considérons une suite de \( 0\) et de \( 1\) dans laquelle le \( 1\) arrive avec une probabilité \( p\) et le \( 0\) avec une probabilité \( 1-p\). Une telle suite est modélisée par une suite de variables aléatoires de Bernoulli \( (X_n)_{n\in\eN}\) indépendantes de paramètre \( p\).

    Question : une telle suite contient elle une infinité de \( 1\) ? Considérons les événements indépendants \( A_n=\{ X_n=1 \}\). Nous avons
    \begin{equation}
        \sum_n P(A_n)=\sum_nP(X_n=1)=\sum_np=\infty.
    \end{equation}
    Par Borel-Cantelli et son expression \eqref{EqparaphrCantelli}, nous avons alors
    \begin{equation}
        P(\limsup A_n)=1.
    \end{equation}
    Donc une infinité d'événements \( A_n\) se produisent, et nous avons bien une infinité de \( 1\) dans la suite.

    Remarque : dans ce raisonnement nous pouvons considérer une probabilité non constante \( p_n\) tant que la série \( \sum_np_n\) diverge.
\end{example}

%+++++++++++++++++++++++++++++++++++++++++++++++++++++++++++++++++++++++++++++++++++++++++++++++++++++++++++++++++++++++++++
\section{Loi des grands nombres, théorème central limite}
%+++++++++++++++++++++++++++++++++++++++++++++++++++++++++++++++++++++++++++++++++++++++++++++++++++++++++++++++++++++++++++

%---------------------------------------------------------------------------------------------------------------------------
\subsection{Loi des grands nombres}
%---------------------------------------------------------------------------------------------------------------------------

\begin{lemma}[Inégalité de Markov]\index{Markov (inégalité)}
    Soit une variable aléatoire \( X\in L^p\) et \( \epsilon>0\). Nous avons
    \begin{equation}
        P(| X |\geq \epsilon)\leq \frac{1}{ \epsilon^r }E(| X |^r).
    \end{equation}
\end{lemma}

\begin{proof}
    Nous avons
    \begin{equation}
        E(| X |^r)\geq\int_{| X |\geq \epsilon}X(\omega)dP(\omega)\geq \epsilon^r\int_{| X |\geq\epsilon}dP=\epsilon^rP(| Y |\geq\epsilon).
    \end{equation}
\end{proof}

\begin{theorem}[Loi forte des grands nombres]\index{loi!des grands nombres (forte)}
    Soit \( (X_n)\) une suite de variables aléatoires réelles
    \begin{enumerate}
        \item
            indépendantes et identiquement distribuées,
        \item
            intégrables (c'est à dire dans \( L^1\)),
    \end{enumerate}
    et \( \bar X_n=\frac{1}{ n }\sum_{i=1}^nX_i\). Alors
    \begin{equation}
        \bar X_n\stackrel{p.s.}{\longrightarrow} E(X_1).
    \end{equation}
\end{theorem}
Note : étant donné que les variables aléatoires sont identiquement distribuées, nous avons évidemment \( E(X_1)=E(X_2)=\ldots\)

\begin{probleme}
    Est-ce que les variables aléatoires doivent vraiment être réelles ?
\end{probleme}

\begin{corollary}
    Si les variables aléatoires réelles \( X_n\) sont
    \begin{enumerate}
        \item
            indépendantes et identiquement distribuées,
        \item
            dans \( L^2\)
    \end{enumerate}
    alors
    \begin{equation}
        \bar X_n\stackrel{P}{\longrightarrow}E(X_1).
    \end{equation}
\end{corollary}

\begin{proof}
    Nous voulons prouver que pour tout \( \eta>0\),
    \begin{equation}
        P\big( | \bar X_n-E(X_1) |>\eta \big)\to 0.
    \end{equation}
    Remarquons d'abord que les variables aléatoires \( X_n\) étant identiquement distribuées, \( E(\bar X_n)=E(X_1)\) parce que \( E(X_i)=E(X_1)\) pour tout \( i\). L'inégalité de Markov avec \( r=2\) nous donne
    \begin{equation}
        P\big( | \bar X_n-E(\bar X_n) |>\eta \big)\leq\frac{1}{ \eta^2 }E\big( | \bar X_n-E(\bar X_n) |^2 \big)
    \end{equation}
    où nous reconnaissons \( E\big( | \bar X_n-E(\bar X_n) |^2 \big)=\Var(\bar X_n)\). Par la proposition \ref{LemEXYEXEYindep} nous avons \( \Var(\bar X_n)=\Var(X_1)/n\) et par conséquent
    \begin{equation}
        P\big( | \bar X_n-E(\bar X_n) |>\frac{1}{ n\eta^2 }\Var(X_1),
    \end{equation}
    qui tend vers zéro lorsque \( n\to\infty\).
\end{proof}

\begin{proposition}
    Soient \( X_n\) des variables aléatoires indépendantes et identiquement distribuées avec \( X_n\geq 0\). Nous acceptons \( E(X_1)=\infty\), c'est à dire que nous relaxons la condition \( X_n\in L^1\) par rapport à la loi des grands nombres.

    Alors
    \begin{equation}
        \bar X_n\stackrel{p.s.}{\longrightarrow} E(X_1)\in\mathopen[ 0 , \infty \mathclose].
    \end{equation}
\end{proposition}

\begin{proof}
    Si \( E(X_1)<\infty\), nous sommes dans le cas de la loi des grands nombres. Pour chaque \( N\in\eN\) nous considérons la suite de variables aléatoires
    \begin{equation}
        X_n^{(N)}=\min(X_n,N).
    \end{equation}
    Nous avons évidement \( \bar X^{(N)}_n\leq \bar X_n\). Les variables aléatoires \( X^{(N)}_n\) étant bornées par \( N\), elles vérifient la loi des grands nombres pour chaque \( N\) séparément. Par conséquent nous avons pour chaque \( N\) la limite
    \begin{equation}        \label{EqbarXNbtoXnubus}
        \bar X^{(N)}_n\to E(X^{(N)}_1)
    \end{equation}
    Nous supposons que \( E(X_1)=\infty\), par conséquent  \( \lim_{N\to\infty}E(X_1^{(N)})=\infty\). Soit \( \eta>0\) et choisissons \( N\) de telle manière à avoir
    \begin{equation}
        E(X_1^{(N)})>\eta+1.
    \end{equation}
    La limite \eqref{EqbarXNbtoXnubus} nous permet de trouver \( n_0\) tel que pour tout \( n>n_0\) nous ayons \( \bar X^{(N)}_n>\eta\). Au final,
    \begin{equation}
        \eta<\bar X^{(N)}_n\leq \bar X_n,
    \end{equation}
    ce qui montre que \( \bar X_n\to\infty\).
\end{proof}

\begin{example}
    La loi des grands nombres justifie la pratique courante d'approximer une grandeur physique par la moyenne empirique d'un grand nombre de mesures.
\end{example}

%---------------------------------------------------------------------------------------------------------------------------
\subsection{Marche aléatoire}
%---------------------------------------------------------------------------------------------------------------------------

Nous considérons un mobile qui se déplace sur l'axe \( \eZ\). À chaque pas de temps, nous supposons qu'il va faire un pas à gauche avec une probabilité \( p\) et un pas à droite avec une probabilité \( (1-p)\). Nous nous demandons quel est le mouvement du mobile sur le long terme.

La position \( S_n\) du mobile à l'instant \( n\) est donnée par
\begin{equation}
    S_n=\sum_{i=1}^nX_i
\end{equation}
où \( X_i\) est le pas effectué à l'instant \( i\). Ce sont des variables de Bernoulli indépendantes avec
\begin{equation}
    X_i\stackrel{\hL}{=}p\delta_{-1}+(1-p)\delta_1
\end{equation}
c'est à dire
\begin{subequations}
    \begin{align}
        P(X_i=-1)&=p\\
        P(X_i=1)&=1-p.
    \end{align}
\end{subequations}
Ces variables vérifient les hypothèses de la loi des grands nombres :
\begin{enumerate}
    \item
        elles sont indépendantes et identiquement distribuées,
    \item
        elles sont intégrables.
\end{enumerate}
Pour le second point, le calcul est
\begin{equation}
    \int_{\Omega}| X_i |dP=\int_{\eR}| x |dP_X=\int_{\eR}| x |(p\delta_{-1}+(1-p)\delta_{1})=| 1-p |+| p |=1.
\end{equation}
Nous avons par conséquent
\begin{equation}
    \bar X_n=\frac{1}{ n }\sum_{i=1}^nX_i\stackrel{p.s.}{\longrightarrow} E(X_1)=(1-2p)
\end{equation}
et
\begin{equation}
    \frac{ S_n }{ n }\to(1-2p).
\end{equation}
Si \( p\neq 1/2\) nous pouvons conclure que
\begin{enumerate}
    \item
        si \( p>1/2\), alors \( S_n\stackrel{p.s.}{\longrightarrow}-\infty\)
    \item
        si \( p<1/2\), alors \( S_n\stackrel{p.s.}{\longrightarrow}\infty\).
\end{enumerate}
De plus nous connaissons la vitesse de divergence : elle est linéaire. Le mobile suit essentiellement l'équation
\( p(n)=(1-2p)n\).

\begin{remark}
    Cela ne traite pas le cas \( p=1/2\). Dans ce cas, nous pouvons simplement dire que \( S_n=o(n)\).
\end{remark}

%---------------------------------------------------------------------------------------------------------------------------
\subsection{Théorème central limite}
%---------------------------------------------------------------------------------------------------------------------------

\begin{lemma}       \label{Lemexpznznsurnton}
    Soit \( z_n\to z\) une suite convergente dans \( \eC\). Alors
    \begin{equation}
        \left( 1+\frac{ z_n }{ n } \right)^n\to e^z.
    \end{equation}
\end{lemma}

\begin{theorem}
    Si les variables aléatoires \( X_n\) sont
    \begin{enumerate}
        \item
            indépendantes et identiquement distribuées de loi parente \( X\),
        \item
            \( X_1\in L^2(\Omega,\tribA)\),
    \end{enumerate}
    alors nous notons \( \bar X_n=\frac{1}{ n }\sum_{i=1}^{n}X_i\), \( m=E(X_1)\) et \( \sigma^2=\Var(X_1)\) et nous avons
    \begin{equation}
        \frac{ \bar X_n-E(X) }{ \sqrt{\Var(X)/n} }=\frac{ \bar X_n-E(X) }{ \sigma/\sqrt{n} }\stackrel{\hL}{\longrightarrow}\dN(0,1).
    \end{equation}
    
\end{theorem}

\begin{proof}
    Nous allons écrire la démonstration dans le cas de variables aléatoires réelles. La proposition \ref{PrpopCaractCvL} dit que la suite \( X_n\) converge en loi vers \( X\) si et seulement si les fonctions caractéristiques convergent ponctuellement. Nous devons donc prouver, pour chaque\footnote{Chuck Norris peut \emph{vraiment} le faire pour \emph{chaque} $t\in\eR$} \( t\in\eR\) que
    \begin{equation}        \label{EqPhitophidNtznu}
        \Phi_{\frac{ S_n-nm }{ \sigma\sqrt{n} }}(t)\to\Phi_{\dN(0,1)}(t).
    \end{equation}
    Supposons dans un premier temps que \( E(X_i)=0\) et \( \sigma(X_i)=1\). Dans ce cas nous considérons la fonction 
    \begin{subequations}
        \begin{align}
            \Phi_{\frac{ S_n }{ \sqrt{n} }}&=E\big(  e^{i\frac{ t }{ \sqrt{n} }\sum_{k=1}^nX_k} \big)\\
            &=\prod_{k=1}^nE\big(  e^{i\frac{ t }{ \sqrt{n} }X_1} \big)\\
            &=\prod_{k=1}^n\Phi_{X_1}\left( \frac{ t }{ \sqrt{n} } \right)\\
            &=\Phi_{X_1}\left( \frac{ t }{ \sqrt{n} } \right)^n.
        \end{align}
    \end{subequations}
    Cette quantité est a priori complexe; nous ne pouvons donc pas immédiatement passer au logarithme. Nous pouvons par contre utiliser un développement en puissances de \( t\) en nous servant de la proposition \ref{PropDerFnCaract} et de l'hypothèse comme quoi \( X_1\in L^2\) :
    \begin{equation}
        \Phi_{X_1}(t)=\Phi_{X_1}(0)+\Phi_{X_1}'(0)t+\Phi_{X_1}''(0)\frac{ t^2 }{2}+\alpha(t)t^2
    \end{equation}
    où \( \alpha\) est une fonction qui a la propriété \( \lim_{x\to 0} \alpha(x)=0\).

    En utilisant les hypothèses et la formule de dérivation de la fonction caractéristique,
    \begin{subequations}
        \begin{align}
            \Phi_{X_1}(0)&=1\\
            \Phi'_{X_1}(0)&=E\big( (iX) \big)=0\\
            \Phi''_{X_1}(0)&=E(-X^2)=-\Var(X_1)=-1.
        \end{align}
    \end{subequations}
    Nous avons donc
    \begin{equation}
        \Phi_{X_1}\left( \frac{ t }{ \sqrt{n} } \right)=\underbrace{1-\frac{ 1 }{2}\frac{ t^2 }{ n }}_{\in\eR}+\underbrace{\frac{ t^2 }{ n }\alpha\left( \frac{ t }{ \sqrt{n} } \right)}_{\in\eC},
    \end{equation}
    de telle sorte que, en considérant une valeur fixée de \( t\),
    \begin{equation}        \label{EqPhifracfacbetanrigh}
        \Phi_{\frac{ S_n }{ \sqrt{n} }}(t)=\left( 1-\frac{ \frac{ t^2 }{ 2 }+\beta_n }{ n } \right)^n
    \end{equation}
    où \( \beta_n=t^2\alpha(t/\sqrt{n})\). Nous avons bien entendu \( \lim_{n\to \infty} \beta_n=0\).
    
    Nous pouvons appliquer le lemme \ref{Lemexpznznsurnton} pour obtenir la limite
    \begin{equation}        \label{EqlimninfySnsqrtntdsnd}
        \lim_{n\to \infty} \Phi_{\frac{ S_n }{ \sqrt{n} }}(t)= e^{-t^2/2}.
    \end{equation}
    La convergence \eqref{EqPhitophidNtznu} est par conséquent prouvée dans le cas où \( E(X_i)=0\) et $\Var(X_i)=1$.

    Considérons maintenant des variables aléatoires avec \( E(X_i)=m\) et \( \Var(X_i)=\sigma^2\). Elles peuvent être écrites sous la forme
    \begin{equation}
        X_i=\sigma X_i+m
    \end{equation}
    où \( X'_i\) est d'espérance nulle et de variance un. Nous avons alors
    \begin{equation}
        S_n=\sigma\sum_{i=1}^nX_i'+nm,
    \end{equation}
    et
    \begin{equation}
        \frac{ S_n-nm }{ \sigma\sqrt{n} }=\frac{ S'_n }{ \sqrt{n} }
    \end{equation}
    où \( S'_n=\sum_iX'_i\). L'étude de la variable aléatoire 
    \begin{equation}
        \frac{ S_n-nm }{ \sigma\sqrt{n} }
    \end{equation}
    revient donc à celle de \( S'_n/\sqrt{n}\) qui vient d'être effectuée.
\end{proof}

\begin{remark}
    Nous pouvons obtenir la limite \eqref{EqlimninfySnsqrtntdsnd} d'une façon alternative. Nous considérons la détermination du logarithme complexe sur \( \eC\setminus\eR^-\); cela est une fonction analytique vérifiant l'équation
    \begin{equation}
        e^{\ln(z)}=z
    \end{equation}
    pour tout \( z\in\eC\setminus\eR^-\) et le développement
    \begin{equation}
        \ln(1+z)=\sum_{n=1}^{\infty}(-1)^{n+1}\frac{ z^n }{ n }.
    \end{equation}
    En particulier, \( \ln(1+z)=z+z\alpha(z)\) où \( \lim_{z\to 0}\alpha(z)=0\). Nous reprenons à l'équation \eqref{EqPhifracfacbetanrigh} en fixant \( t\). Nous avons
    \begin{subequations}
        \begin{align}
            \Phi_{\frac{ S_n }{ \sqrt{n} }}(t)&=\exp\left[ \ln\big(\Phi_{\frac{ S_n }{ \sqrt{n} }}(t)\big) \right]\\
            &=\exp\left[ n\ln\left( 1+\frac{ -\frac{ t^2 }{2}-\beta_n }{ n } \right) \right]\\
            &=\exp\left[ -\frac{ t^2 }{2}-\beta_n+ \big( -\frac{ t^2 }{2}-\beta_n \big)\alpha\left( \frac{ -\frac{ t^2 }{2}-\beta_n }{ n } \right) \right].
        \end{align}
    \end{subequations}
    À la limite \( n\to 0\) nous tombons sur \(  e^{-t^2/2}\).
    
\end{remark}

\begin{remark}
    Étant donné que la variable aléatoire 
    \begin{equation}
        \frac{ S_n-nm }{ \sigma\sqrt{n} }
    \end{equation}
    converge en loi vers \( \dN(0,1)\), nous avons la convergence des fonctions de répartition partout où la fonction de répartition de la normale est continue (donc sur tout \( \eR\)). En particulier,
    \begin{equation}
        \left| P\left( \frac{ S_n-nm }{ \sigma\sqrt{n} }\leq x \right)-\int_{-\infty}^x\frac{1}{ \sqrt{2\pi} } e^{-y^2/2}dy \right| \to 0.
    \end{equation}
    Nous avons la borne de \defe{Berry-Esséen}{Berry-Esséen (borne)} qui donne une estimation de la vitesse de convergence: si \( X\in L^3\), alors  il existe une constante \( C\), indépendante de \( x\), des \( X_i\) et de \( n\) telle que
    \begin{equation}
        \left| P\left( \frac{ S_n-nm }{ \sigma\sqrt{n} }\leq x \right)-\int_{-\infty}^x\frac{1}{ \sqrt{2\pi} } e^{-y^2/2}dy \right| \leq\frac{ X\mu_3 }{ \sigma^3\sqrt{n} }
    \end{equation}
    où \( \mu_3=E\big( | X_1-m |^3 \big)\) est le moment d'ordre \( 3\) de \( X\). La chose à retenir est que la convergence est à la vitesse de \( 1/\sqrt{n}\).
\end{remark}

En dimension \( d>1\), nous avons encore un théorème central limite.
\begin{theorem}
    Si \( d>1\), et si nous avons des variables aléatoires \( X_n\) à valeurs dans \( \eR^d\) avec
    \begin{enumerate}
        \item
            les \( X_n\) sont indépendantes et identiquement distribuées
        \item
            les \( X_n\) sont dans \( L^2\).
    \end{enumerate}
    Alors nous notons \( X_1=(X_1^{(1)},\ldots,X_1^{(d)})\). Nous avons
    \begin{equation}
        \frac{ S_n-nm }{ \sqrt{n} }\stackrel{\hL}{\longrightarrow}\dN(0,\Sigma)
    \end{equation}
    où \( \Sigma\) est ma matrice de covariance du vecteur aléatoire \( X_1\) :
    \begin{equation}
        \Sigma=\big( \Cov(X_1^{(1),\ldots,X_1^{(d)}}) \big)_{i,j=1,\ldots,d}.
    \end{equation}
\end{theorem}




\chapter{Lois de probabilités}
% This is part of Agregation : modélisation
% Copyright (c) 2011
%   Laurent Claessens
% See the file fdl-1.3.txt for copying conditions.

%+++++++++++++++++++++++++++++++++++++++++++++++++++++++++++++++++++++++++++++++++++++++++++++++++++++++++++++++++++++++++++
\section{Les lois usuelles}
%+++++++++++++++++++++++++++++++++++++++++++++++++++++++++++++++++++++++++++++++++++++++++++++++++++++++++++++++++++++++++++

%---------------------------------------------------------------------------------------------------------------------------
\subsection{Loi de Bernoulli}
%---------------------------------------------------------------------------------------------------------------------------

Une expérience de Bernoulli consiste à tirer au hasard un \( 0\) ou un \( 1\) avec une probabilité \( p\) de tomber sur \( 1\) et \( 1-p\) de tomber sur zéro. Il s'agit donc d'une expérience qui réussi ou qui rate.

Le cas typique est une urne avec des boules indiscernables blanches ou noires. La probabilité \( p\) est la proportion de blanches dans l'urne (avec remise entre les tirages). Dans ce cas, nous avons l'espace de probabilité \( (\Omega,\tribA,P)\) où \( \Omega\) représente l'ensemble des boules, \( \tribA\) est l'ensemble des parties de \( \Omega\) et \( P\) est l'équiprobabilité sur \( \Omega\). Une variable aléatoire est une application
\begin{equation}
    \begin{aligned}
        X\colon \Omega&\to \{ 0,1 \} \\
        \omega&\mapsto \text{couleur de la boule \( \omega\)}.
    \end{aligned}
\end{equation}
Nous notons \( \dB(1,p)\) la loi de Bernoulli. Elle a une expression très simple :
\begin{subequations}
    \begin{align}
        \dB(0,1)\big( \{ 1 \} \big)&=p\\
        \dB(0,1)\big( \{ 0 \} \big)&=1-p
    \end{align}
\end{subequations}

%---------------------------------------------------------------------------------------------------------------------------
\subsection{Loi binomiale}
%---------------------------------------------------------------------------------------------------------------------------

Une expérience binomiale consiste à répéter \( n\) expériences de Bernoulli de paramètre \( p\) et de compter le nombre de réussites. Une telle expérience peut être réalisée selon la procédure suivante.

Soit une urne contenant \( N\) boules dont une proportion \( p\) de \( 1\) et \( 1-p\) de \( 0\). Une expérience binomiale de paramètres \( n\) et \( p\) consistera à prendre \( n\) boules \emph{avec remise} et à compter le nombre de \( 1\) obtenus.

En termes d'espaces probabilisé, nous avons \( \Omega\) qui est l'ensemble des tuple de taille \( n\) à valeurs dans \( \{ 0,1 \}\), la tribu \( \tribA\) est l'ensemble des parties de \( \Omega\), et la probabilité \( P\) est l'équiprobabilité :
\begin{equation}
    P(\omega)=\frac{1}{ N^n }
\end{equation}
si il y a \( N\) boules dans l'urne. Nous construisons alors la variable aléatoire
\begin{equation}
    X(\omega)=\sum_{i=1}^n\omega_i
\end{equation}
où \( \omega\) est une suite de taille \( n\) de $0$ et de \( 1\).

Calculons \( P(x=k)\). Il s'agit de considérer tous les sous-ensembles de taille \( n\) de \( \Omega\) contenant exactement \( k\) fois \( 1\). Il y a \( {n\choose k}\) manière de décider lesquelles des \( n\) boules seront blanches. Ensuite, chaque boule blanche peut être choisie parmi les \( m\) boules disponibles, et chaque boule noire peut être choisie parmi les  \( (N-m)\) disponibles. Nous avons donc
\begin{equation}        \label{EqformunPxkBin}
    P(X=k)={n\choose k}\frac{ m^k(N-m)^{n-k} }{ N^k }.
\end{equation}
En effet la mesure de probabilité sur \( \Omega\) est la mesure de comptage renormalisée par le cardinal de \( \Omega\) qui vaut \( N^m\). Étant donné que \( p=m/N\), nous transformons facilement \eqref{EqformunPxkBin} en
\begin{equation}
    P(X=k)={n\choose k}p^k(1-p)^{n-k}.
\end{equation}

%---------------------------------------------------------------------------------------------------------------------------
\subsection{Loi géométrique}
%---------------------------------------------------------------------------------------------------------------------------

Soit \( (X_n)\) une suite indépendante et identiquement distribuée de lois de Bernoulli de paramètre \( p\). Alors la variable aléatoire
\begin{equation}
    Z=\inf\{ n\geq 1\tq X_n=1 \}
\end{equation}
est une loi géométrique de paramètre \( p\).

La loi géométrique compte donc le nombre d'expériences de Bernoulli à effectuer avant que le premier succès soit au rendez-vous. Nous avons
\begin{equation}
    P(Z=k)=P(X_k=1)P(X_1,\ldots,X_{k-1}=0)=p(1-p)^{k-1}
\end{equation}

%---------------------------------------------------------------------------------------------------------------------------
\subsection{Loi de Poisson}
%---------------------------------------------------------------------------------------------------------------------------

Une variable aléatoire \( Z\) suit une \defe{loi de Poisson}{loi!de Poisson} de paramètre \( \lambda\), notée \( \dP(\lambda)\) si
\begin{equation}
    P(Z=k)= e^{-\lambda}\frac{ \lambda^k }{ k! }
\end{equation}
pour tout \( k\in\eN\).

%---------------------------------------------------------------------------------------------------------------------------
\subsection{Approximation de la binomiale par une Poisson}
%---------------------------------------------------------------------------------------------------------------------------

\begin{proposition}
    Soit \( (X_n)\) une suite de variables aléatoires avec \( X_n\sim\dB(n,p_n)\) telle que \( np_n\) converge vers une constante \( \lambda>0\). Alors \( X_n\stackrel{\hL}{\longrightarrow}\dP(\lambda)\).
\end{proposition}

\begin{proof}
    Commençons par écrire la loi binomiale sous une forme plus adaptée au passage à la limite :
    \begin{equation}
        P(X=k)={n\choose k}p^k(1-p)^{n_k}=\frac{ n(n-1)\ldots (n-k+1) }{ k! }p^k(1-p)^{n-k}.
    \end{equation}
    Le produit au numérateur contient \( k\) termes dans lesquels nous mettons \( n\) en évidence. Nous trouvons
    \begin{equation}
        P(X=k)=\frac{ (np)^k\left( 1-\frac{1}{ n } \right)\left( 1-\frac{ 2 }{ n } \right)\ldots\left( 1-\frac{ k-1 }{ n } \right) }{ k! }p^k(n-p)^{n-k}.
    \end{equation}
    Lorsque nous passons à la limite, tous les facteurs du type \( 1-l/n\) tendent vers \( 1\) ainsi que \( (1-p_n)^{-k}\). Les facteurs sont la limite n'est pas \( 1\) sont donc
    \begin{equation}
        P(X_n=k)\simeq\frac{ (np_n)^k }{ k! }(1-p_n)^k.
    \end{equation}
    Nous avons
    \begin{equation}
        \lim_{n\to \infty} (1-p_n)^n=\lim_{n\to \infty} \left( 1-\frac{ np_n }{ n } \right)^n= e^{-\lambda}.
    \end{equation}
    La thèse est alors obtenue en remettant les morceaux ensemble.
\end{proof}

\begin{example}
    Considérons un serveur informatique qui reçoit des requêtes. Toutes les \( \unit{10^{-3}}{\second}\) il reçoit une requête avec une probabilité \( p=0.05\). La variable aléatoire qui consiste à donner le nombre de requêtes effectivement effectuées en une seconde suit une loi binomiale \( \dP(1000,p)\).

    Déterminons la probabilité que le serveur reçoive \( 20\) requêtes en une seconde. Nous approximons \( \dB(1000,0.05)\) par \( \dP(50)\), et la réponse est
    \begin{equation}
        e^{-50}\frac{ 50^{20} }{ 20! }\simeq 7\cdot 10^{-7}.
    \end{equation}
\end{example}

%---------------------------------------------------------------------------------------------------------------------------
\subsection{Loi exponentielle}
%---------------------------------------------------------------------------------------------------------------------------

La loi exponentielle de paramètre \( \lambda\), notée \( \dE(\lambda)\) est la loi de densité
\begin{equation}
    x\mapsto\begin{cases}
        \lambda e^{-\lambda x}    &   \text{si \( x\geq 0\)}\\
        0    &    \text{sinon}.
    \end{cases}
\end{equation}
Si \( X\sim\dE(\lambda)\), alors la fonction de répartition de \( X\) est donnée par
\begin{equation}
    P(X\leq x)=\begin{cases}
        1- e^{-\lambda x}    &   \text{si \( x\geq 0\)}\\
        0    &    \text{sinon}.
    \end{cases}
\end{equation}
En ce qui concerne l'espérance nous faisons le calcul suivant :
\begin{equation}
    E(X)=\int_{\eR}xf_X(x)dx=\lambda\int_{\eR^+}x e^{-\lambda x}dx=\frac{1}{ \lambda }.
\end{equation}
La loi exponentielle est une loi \emph{sans mémoire} en ce sens que
\begin{equation}
    P(X>x+y|X>y)=P(X>x).
\end{equation}
En effet nous utilisons la règle de la probabilité conditionnelle
\begin{equation}
    P(A|B)=\frac{ P(A\cap B) }{ P(B) }.
\end{equation}
Ici,
\begin{equation}
    P(X>x+y|X>y)=\frac{ P(X>x+y) }{ P(X>y) }= e^{-\lambda x}.
\end{equation}

\begin{example}
    Une machine a une durée de vie représentée par une variable aléatoire suivant une loi de Poisson de paramètre \( \lambda\). Soit \( T_y\) la variable aléatoire qui représente la temps de vie restant sachant que la machine a déjà vécu un temps \( y\). Nous voulons trouver la fonction de répartition de \( T_y\). Nous avons
    \begin{equation}
        P(T_y>x)=P(X>x+y|X>y)=P(X>x)= e^{-\lambda x}.
    \end{equation}
    Dans ce cas, la loi de \( T_y\) ne dépend pas de \( y\). Cela signifie que la machine ne vieilli pas et surtout que le modèle n'est pas réaliste.
\end{example}

%---------------------------------------------------------------------------------------------------------------------------
\subsection{Loi de Poisson et loi exponentielle}
%---------------------------------------------------------------------------------------------------------------------------






\chapter{Statistiques}
%+++++++++++++++++++++++++++++++++++++++++++++++++++++++++++++++++++++++++++++++++++++++++++++++++++++++++++++++++++++++++++
\section{Notations et hypothèses}
%+++++++++++++++++++++++++++++++++++++++++++++++++++++++++++++++++++++++++++++++++++++++++++++++++++++++++++++++++++++++++++

Nous notons \( X\) le caractère à étudier, et \( \Omega\) l'ensemble des individus. Le caractère à étudier est vu comme une fonction sur \( \Omega\) :
\begin{equation}
    X\colon \Omega\to \eR,\eN,\{ 0,1 \},\ldots
\end{equation}
Les \defe{statistiques descriptives}{statistiques!descriptives} sont les techniques pour présenter et résumer les données : diagrammes, graphiques, indicateurs numériques : moyenne, écart-type, médiane, \ldots

Nous faisons les hypothèses suivantes :
\begin{enumerate}
    \item
        Chaque observation \( x_i\) est la réalisation de la variable aléatoire \( X\) qui sera de loi inconnue \( \mu\).
    \item
        Le \( n\)-uple \( (x_1,\ldots,x_n)\) est la réalisation de \( (X_1,\ldots,X_n)\) qui est l'échantillon de taille \( n\).
    \item
        Les variables aléatoires \( X_i\) sont indépendantes et identiquement distribuées, de loi commune \( \mu\). La loi \( \mu\) est la \defe{loi parente}{loi!parente} de l'échantillon.
\end{enumerate}

\begin{example}
    Un échantillon de taille \( 1\) consisterait à tirer au sort une personne dans une population et mesurer sa taille.
\end{example}

\begin{example}
    Une échantillon de taille \( n\) consisterait à tirer au sort \( n\) personnes dans une population et de mesurer leurs tailles.
\end{example}

L'\defe{inférence statistique}{inférence statistique} est l'art de dégager des informations sur la population à partir d'informations partielles : intervalles de confiance, estimateurs, test d'hypothèses, \ldots

En théorie des probabilités, nous connaissons la loi de la variable aléatoire \( X\) et nous en déduisons des informations sur le réalisations de \( X\) : valeur la plus probable, moyenne, intervalle dans lequel \( X(\omega)\) a le plus de chance d'appartenir. En statistique, au contraire, la loi est inconnues et nous cherchons des informations sur la loi à partir d'un échantillon de données numériques observées.

%+++++++++++++++++++++++++++++++++++++++++++++++++++++++++++++++++++++++++++++++++++++++++++++++++++++++++++++++++++++++++++
\section{Modèle statistique}
%+++++++++++++++++++++++++++++++++++++++++++++++++++++++++++++++++++++++++++++++++++++++++++++++++++++++++++++++++++++++++++

Une \defe{modèle statistique}{modèle statistique} est un triplet
\begin{equation}
    \statS=\Big[ (\Omega,\tribF,P),(X_{\theta})_{\theta\in\Theta},(\mu_{\theta})_{\theta\in\Theta} \Big]
\end{equation}
où \( (\Omega,\tribF,P)\) est un espace probabilisé, \( (X_{\theta})\) est une famille de variables aléatoires définies sur \( \Omega\) et telles que pour tout \( \theta\in\Theta\), la variable aléatoire \( X_{\theta}\) suive la loi \( \mu_{\theta}\). Les $\mu_{\theta}$ sont des mesures sur les boréliens de \( \eR\) et pour tout \( B\in\Borelien(\eR)\) nous avons
\begin{equation}
    P(X_{\theta}\in B)=\mu_{\theta}(B).
\end{equation}
\begin{remark}
    D'une certaine manière, l'introduction de \( \mu_{\theta}\) dans la définition est redondante parce que ces mesures sont déjà contenues dans la données des variables aléatoires \( X_{\theta}\).
\end{remark}

\begin{example}[Modèle statistique gaussien]
    Si nous savons que les variables aléatoires \( X_i\) suivent une loi gaussienne, alors nous considérons \( \Theta=\eR\times\eR^+\) et \( \theta=(m,\sigma^2j)\). Dans ce cas, \( \mu_{\theta}=\dN(m,\sigma^2)\) et le but de la statistique est de déterminer la valeur de \( \theta\) qui correspond à une population en partant de l'observation d'un échantillon.
\end{example}

\begin{definition}
    Si \( \Theta\subset\eR^k\), nous disons que le modèle statistique est un modèle \defe{paramétrique}{modèle!paramétrique}.
\end{definition}
Le modèle gaussien est un modèle paramétrique : dès que \( m\) et \( \sigma^2\) sont déterminés, la loi du phénomène \( X\) est connue.

\begin{definition}
    Pour chaque \( \theta\in\theta\), un \defe{échantillon}{échantillon} de taille \( n\) associé à un modèle statistique \( \big[ (\Omega,\tribF,P),(X_{\theta}),(\mu_{\theta}) \big]\) est un vecteur \( \big( X_{\theta,1},\ldots,X_{\theta,n} \big)\) de taille \( n\) de variables aléatoires indépendantes et identiquement distribuées de la même loi que la variable aléatoire \( X_{\theta}\). La loi \( \mu_{\theta}\) est la \defe{loi parente}{loi!parente} de l'échantillon.
\end{definition}

\begin{definition}
    Un \defe{modèle d'échantillonnage}{modèle!échantillonnage} sur le modèle statistique \( \statS\) est une famille \( (X_{\theta,1},\ldots, X_{\theta,n})_{\theta\in\Theta}\) d'échantillons de taille \( n\geq 1\).
\end{definition}

Nous noterons souvent \( (X_1,\ldots,X_n)\) à la place de \( (X_{\theta,1},\ldots,X_{\theta,n})\) un échantillon, mais il faut se souvenir que les \( X_i\) suivent toujours la même loi donnée par \( \theta\). La loi du vecteur \( (X_1,\ldots,X_n)\) est \( \mu_{\theta}\otimes\ldots\otimes\mu_{\theta}\) et est définie sur l'espace \( (\Omega^n,\tribF\otimes\ldots\otimes\tribF,P^{\otimes n})\).

\begin{remark}
    Le travail du statisticien est de proposer un modèle statistique \( \statS\) a priori. Si nous étudions la taille d'une population, nous allons choisir un modèle gaussien. Plus le modèle est précis, plus l'espace \( \Theta\) est petit mais plus il y a de risques que le vérité soit hors de l'ensemble considéré.
\end{remark}

\begin{example}
    Soit \( X\) une variable aléatoire de carré intégrable que l'on sait simuler. Affin d'évaluer la moyenne \( \mu\) de \( X\), nous pouvons considérer la moyenne empirique des simulations : \( \bar X_n=\frac{1}{ n }\sum_{i=1}^nX_i\) où les variables aléatoires \( X_i\) sont indépendantes, identiquement distribuées et de même loi que \( X\). 

    La loi des grands nombres nous enseigne que \( \bar X_n\to\mu\). De plus,
    \begin{equation}
        \lim_{n\to \infty} P\left( \mu\in\left[ \bar X_n-\frac{ a\sigma }{ \sqrt{n} },\bar X_n+\frac{ a\sigma }{ \sqrt{n} } \right] \right)=\int_{-a}^a e^{-x^2/2}\frac{ dx }{ \sqrt{2\pi} }.
    \end{equation}
    En effet, la condition sur \( \mu\) est équivalente à
    \begin{equation}
        -a\leq \frac{ \bar X_n-\mu }{ \sigma/\sqrt{n} }\leq a,
    \end{equation}
    tandis que le théorème central limite nous enseigne que la variable aléatoire \( \frac{ \bar X_n-\mu }{ \sigma/\sqrt{n} }\) se comporte comme \( \dN(0,1)\) lorsque \( n\) est grand. Dans ce cas, nous avons que
    \begin{equation}
        P\left( \frac{ \bar X_n-\mu }{ \sigma/\sqrt{n} }\in\mathopen[ -a , a \mathclose] \right)=\frac{1}{ \sqrt{2\pi} }\int_{-a}^a e^{-x^2/2}dx.
    \end{equation}
    Notons que dans ce calcul nous avons utilisé le fait que \( \mu=E(X_1)\).

    Montrons que la suite
    \begin{equation}
        \sigma_n^2=\frac{1}{ n }\sum_{i=1}^nX_i^2-\left( \frac{1}{ n }\sum_{i=1}^nX_i \right)^2
    \end{equation}
    converge presque sûrement vers \( \sigma^2\). Le théorème central limite implique que
    \begin{equation}
        \frac{1}{ n }\sum_{i=1}^nX_i^2\stackrel{p.s.}{\longrightarrow} E(X_1^2)
    \end{equation}
    et que
    \begin{equation}
        \left( \frac{1}{ n }\sum_{i=1}^nX_i \right)^2\stackrel{p.s.}{\longrightarrow}E(X_1)^2.
    \end{equation}
    La différence converge donc presque sûrement vers \( \sigma^2\) en vertu de la proposition \ref{PrropVarAlterfrom}.

    Nous avons également \( E(\sigma_n^2)=\sigma^2\). En effet, sachant que \( E(X_i)=E(X_1)=\mu\) et que \( E(X_i^2)=E(X_1^2)=\mu^2+\sigma^2\),
    \begin{subequations}
        \begin{align}
            E(\sigma_n^2)&=\frac{1}{ n }\sum_{i=1}^nE(X_i^2)-\frac{1}{ n^2 }\left( \sum_{i=1}^nE(X_i^2) \right)+\sum_{i\neq j}E(X_iX_j))\\
            &=\sigma^2+\mu^2-\frac{1}{ n^2 }\big( n(\sigma^2+\mu^2)+(n^2-n)\mu^2 \big)\\
            &=\sigma^2-\frac{1}{ n }\sigma^2,
        \end{align}
    \end{subequations}
    dont la limite \( n\to\infty\) donne bien \( \sigma^2\).

    Nous voudrions à présent montrer que 
    \begin{equation}
        \frac{ \bar X_n-\mu }{ \sigma/\sqrt{n} }\stackrel{\hL}{\longrightarrow}\dN(0,1).
    \end{equation}
    Vu que le théorème central limite donne une convergence en loi, nous pouvons utiliser le lemme de Slutsky pour montrer que
    \begin{equation}
        \left( \frac{ \bar X_n-\mu }{ 1/\sqrt{n} },\sigma_n^2 \right)\stackrel{\hL}{\longrightarrow}(\sigma Z,\sigma^2)
    \end{equation}
    où \( Z\sim\dN(0,1)\). En vertu de la proposition \ref{PropcvLsousfonc} appliqué à la fonction \( f\colon \eR^2\to \eR\),
    \begin{equation}
        f(x,y)=\begin{cases}
            \frac{ x }{ \sqrt{y} }    &   \text{si \( y\neq 0\)}\\
            0    &    \text{sinon}
        \end{cases}
    \end{equation}
    nous avons la convergence en loi
    \begin{equation}
        f\left( \frac{ \bar X_n-\mu }{ 1/\sqrt{n} },\sigma_n^2 \right)\stackrel{\hL}{\longrightarrow}f(\sigma Z,\sigma^2),
    \end{equation}
    c'est à dire 
    \begin{equation}
        \frac{ \bar X_n-\mu }{ \sigma_n/\sqrt{n} }\stackrel{\hL}{\longrightarrow}Z.
    \end{equation}
    Affin d'être complet, précisons que 
    \begin{equation}
        P\big( (\sigma Z,\sigma)\in\eR\times \{ 0 \} \big)=0.
    \end{equation}
\end{example}

%+++++++++++++++++++++++++++++++++++++++++++++++++++++++++++++++++++++++++++++++++++++++++++++++++++++++++++++++++++++++++++
\section{Statistiques descriptives}
%+++++++++++++++++++++++++++++++++++++++++++++++++++++++++++++++++++++++++++++++++++++++++++++++++++++++++++++++++++++++++++

%+++++++++++++++++++++++++++++++++++++++++++++++++++++++++++++++++++++++++++++++++++++++++++++++++++++++++++++++++++++++++++
\section{Modèles d'échantillonnages}
%+++++++++++++++++++++++++++++++++++++++++++++++++++++++++++++++++++++++++++++++++++++++++++++++++++++++++++++++++++++++++++

Soit \( X\), une variable aléatoire sur \( (\Omega,\tribF,P)\). Un \defe{échantillon}{échantillon} de taille \( n\) pour \( X\) est une suite de \( n\) variables aléatoires \( (X_1,\ldots,X_n)\) définies sur \( (\Omega,\tribF,P)\) indépendantes et de même loi que \( X\). Nous disons que la loi de \( X\) est la \defe{loi parente}{loi!parente d'un échantillon} de la suite \( X_i\).

\begin{definition}
    Soit
    \begin{equation}
        \statS=\Big[ (\Omega,\tribF,P),(X_{\theta}),(\mu_{\theta}) \Big]_{\theta\in\Theta},
    \end{equation}
    un modèle statistique. Un \defe{modèle d'échantillonnage}{modèle!échantillonnage} de taille \( n\) associée au modèle statistique \( \statS\) est la donnée d'une famille de \( n\)-échantillons \( (X_{\theta,1}),\ldots,X_{\theta,n}\) telle que pour tout \( \theta\in\Theta\), l'échantillon \( (X_{\theta,i})\) soit de variable parente \( X_{\theta}\).
\end{definition}

La \defe{moyenne empirique}{moyenne!empirique d'un échantillon} du \( n\)-échantillon \( (X_i)\) est la variable aléatoire
\begin{equation}
    \bar X_n=\frac{1}{ n }\sum_{i=1}^{n}X_i.
\end{equation}

La proposition suivante signifie que la moyenne empirique est une «bonne» façon d'approcher la variable aléatoire.
\begin{proposition}
    Soit \( X\), une variable aléatoire dans \( L^2(\Omega)\) (c'est à dire \( E(X^2)<\infty\)) d'espérance \( m\) et de variance \( \sigma^2\). Alors
    \begin{enumerate}
        \item
            \( E(\bar X_n)=m\) et \( \Var(\bar X_n)=\frac{ \sigma^2 }{ n }\).
        \item
            Nous avons les convergences
            \begin{subequations}
                \begin{align}
                    \bar X_n\stackrel{p.s.}{\longrightarrow} m\\
                    \frac{ \bar X_n-m }{ \sigma/\sqrt{n} }\stackrel{\hL}{\longrightarrow} Z\sim\dN(0,1).
                \end{align}
            \end{subequations}
        \item
            Si \( X\) est de loi \( \dN(m,\sigma^2)\), alors \( \bar X_n\sim\dN(m,\frac{ \sigma^2 }{ n })\), c'est à dire
            \begin{equation}
                \frac{ \bar X_n -m}{ \sigma/\sqrt{n} }\sim\dN(0,1).
            \end{equation}
            <++>
            
    \end{enumerate}
    <++>
\end{proposition}
<++>




\chapter{Chaînes de Markov à temps discret}
Les chaînes de Markov interviennent pour la description des systèmes dont l'évolution future ne dépend que de l'état présent.

\begin{definition}
    Soit \( E\) un ensemble au plus dénombrable et \( (\Omega,\tribF,P)\) un espace probabilisé. Une \defe{chaîne de Markov}{chaîne de Markov} à valeurs dans \( E\) est une famille \( (X_n)_{n\in\eN}\) de variables aléatoires telles que pour tout \( x_0,\ldots,x_{n+1}\in E\),
    \begin{equation}
        P(X_{n+1}=x_{n=1}|X_n=x_n,\ldots,X_0=x_0)=P(X_{n+1}=x_{n+1}|X_n=x_n).
    \end{equation}
\end{definition}
Pour une chaîne de Markov, il n'est pas important de savoir l'historique pour prédire la futur : \( X_{n+1}\) est seulement déterminé par \( X_n\).

\begin{remark}
    Il existe une théorie des chaînes de Markov à temps continu ou avec \( E\) non dénombrable, mais ce n'est pas au programme.
\end{remark}

Nous notons
\begin{equation}
    p_n(x,y)=P(X_{n+1}=y|X_n=x)
\end{equation}
la \defe{probabilité de transition}{transition!probabilité} de la chaîne à l'instant \( n\). Si cette probabilité ne dépend pas de \( n\), nous disons que la chaîne de Markov est \defe{homogène}{homogène!chaîne de Markov}, et nous notons \( p(x,y)\) au lieu de \( p_n(x,y)\). Nous notons \( Q\) la matrice (éventuellement infinie) de transition\index{matrice!de transition}
\begin{equation}
    Q_{xy}=p(x,y).
\end{equation}
Nous avons
\begin{equation}
    \sum_{y\in E}p(x,y)=1
\end{equation}
parce que c'est la somme de toutes les transitions possibles en partant de \( x\). Notons aussi que \( p(x,y)\geq 0\).

\begin{definition}
    Une matrice dont tous les éléments sont positifs ou nuls et donc la somme de toutes les lignes sont \( 1\) est une \defe{matrice stochastique}{matrice!stochastique}.
\end{definition}

\begin{lemma}
    Si \( U\) est une matrice stochastique, alors il existe une chaîne de Markov dont la matrice de transition soit \( U\).
\end{lemma}

\begin{remark}
    La somme \( \sum_{x\in E}p(x,y)\) ne vaut pas spécialement \( 1\). Si les états \( x_1\) et \( x_2\) arrivent tous les deux en \( y\) de façon certaine, alors nous avons \( \sum_xp(x,y)\geq 2\). Il n'y a donc pas de limites aux sommes des colonnes.
\end{remark}

\begin{example}
    Nous considérons une fourmi qui se déplace dans un appartement à trois pièces \( A\), \( B\), \( C\). Supposons qu'à chaque minute, elle a une probabilité \( 1/3\) de rester dans la pièce et une probabilité \( 2/3\) de se déplacer. Le plan de l'appartement est
    \begin{equation}
        \xymatrix{%
        A \ar[r]      &  B\ar[r]&C
           }
    \end{equation}
    De la pièce \( A\) est est donc uniquement possible d'aller vers la pièce \( B\); de la \( B\) il est possible d'aller en \( A\) et en \( C\) et de la \( C\) il est uniquement possible d'aller en \( B\).

    La matrice de transition de cette chaîne de Markov est 
    \begin{equation}
        Q=\begin{pmatrix}
            1/3    &   2/3    &   0    \\
            1/3    &   1/3    &   1/3    \\
            0    &   2/3    &   1/3
        \end{pmatrix}
    \end{equation}
\end{example}

%+++++++++++++++++++++++++++++++++++++++++++++++++++++++++++++++++++++++++++++++++++++++++++++++++++++++++++++++++++++++++++
\section{Marche aléatoire sur \texorpdfstring{$\eZ$}{\( \eZ\)}}
%+++++++++++++++++++++++++++++++++++++++++++++++++++++++++++++++++++++++++++++++++++++++++++++++++++++++++++++++++++++++++++

Soit \( (Y_n)\) une suite de variables aléatoires indépendantes et identiquement distribuées valant \( -1\) avec une probabilité \( p\) et \( 1\) avec une probabilité \( (1-p)\). La loi est
\begin{equation}
    Y_n\sim p\delta_{-1}+(1-p)\delta_{1}.
\end{equation}
Nous considérons la variable aléatoire
\begin{equation}
    X_n=X_0+\sum_{i=1}^nY_i
\end{equation}
où \( X_0\) est une variable aléatoire indépendante des \( Y_i\) à valeurs dans \( \eZ\). Nous vérifions à présent que \( X_n\) est une chaîne de Markov avec comme espace d'états \( E=\eZ\). Nous devons montrer que
\begin{equation}        \label{EqAVoirMarkovMAZ}
    P\big( X_{n+1}=x_{n+1}| X_n=x_n,\ldots,X_0=x_0 \big)=P\big( X_{n+1}=x_{n+1}|X_n=x_n \big).
\end{equation}
Étant donné que \( X_{n+1}=X_n+Y_{n+1}\), nous pouvons remplacer \( X_{n+1}=x_{n+1}\) par \( x_n+Y_{n+1}=x_{n+1}\) dans le membre de gauche de \ref{EqAVoirMarkovMAZ}.



\part{Calcul différentiel et intégral 1}
Cette partie sera progressivement digérée dans les autres chapitres. C'est ici que sont les contributions de Nicolas Richard et Ivik Swan.
Cette partie concerne des rappels et des exercices des cours de calcul différentiel et intégral donnés à Bruxelles en première et deuxième année de physique en 2008-2009. Ces cours n'existent plus aujourd'hui et toute la matière de cette partie sera progressivement fondue dans les autres parties.

\chapter{Rappels de CdI}
% This is part of Exercices et corrigés de CdI-1
% Copyright (c) 2011
%   Laurent Claessens
% See the file fdl-1.3.txt for copying conditions.

\section{Séries}
\label{secseries}

\subsection{Rappels et définitions}
La notion de série formalise le concept de somme infinie. L'absence de
certaines propriétés de ces objets (problèmes de commutativité et même
d'associativité) incitent à la prudence et montrent à quel point une
définition précise est importante.

Soit ${(a_k)}_{k \geq k_0}$ une suite réelle. La \Defn{série de terme
  général $(a_k)$}, notée
\begin{math}
  \sum_{i=k_0}^\infty a_i,
\end{math}
est la suite ${(s_n)}_{n \geq k_0}$ dont les termes sont donnés par
\begin{equation*}
  s_n \pardef \sum_{i=k_0}^k a_i
\end{equation*}
et sont appelés les \Defn{sommes partielles} de la série.

La série $\sum_{i=k_0}^\infty a_i$ \Defn{converge} si la suite $(s_n)$
converge vers un réel $s$. Sa limite est appelée la \Defn{somme de la
  série} et on note
\begin{equation}		\label{EqDefSommeLim}
  \sum_{i=k_0}^\infty a_i = \lim_{n\to\infty}\sum_{i=k_0}^na_i.
\end{equation}
Si la série ne converge pas, elle \Defn{diverge} et peut alors avoir
une limite infinie (uniquement si le terme général est réel, on note
alors $+\infty$ ou $-\infty$ sa limite) ou pas de limite.  La série
$\sum_{i=k_0}^\infty a_i$ \Defn{converge absolument} si la série
$\sum_{i=k_0}^\infty \abs{a_i}$ converge.

\begin{proposition}\label{propnseries_propdebase}
Les principales propriétés de la somme définie par la limite \eqref{EqDefSommeLim} sont
  \begin{enumerate}
  \item Si une série converge absolument, alors elle converge
    (simplement).
  \item Si la série est à termes positifs --c'est-à-dire pour tout
    indice $k$, $a_k \in \eR$ et $a_k \geq 0$-- il n'y a aucune
    différence entre convergence absolue et convergence simple.
  \item\label{point3-seriepropdebase} Si une série converge, son terme général doit tendre vers $0$.
\item 
Si la série converge alors la somme est associative
\item
Si la série converge absolument, alors la somme est commutative.
  \end{enumerate}
\end{proposition}

\begin{remark}Vue comme somme infinie, l'associativité et la
  commutativité dans une série sont perdues. Néanmoins, il subsiste
  que
  \begin{enumerate}
  \item si la série converge, on peut regrouper ses termes sans
    modifier la convergence ni la somme (associativité),
  \item si la série converge absolument, on peut modifier l'ordre des
    termes sans modifier la convergence ni la somme (commutativité).
  \end{enumerate}
\end{remark}

\begin{example}\label{exemplesseries}
\begin{enumerate}

\item
La \Defn{série harmonique} est
\begin{equation*}
\sum_{i=1}^\infty \frac1i
\end{equation*}
et diverge (possède une limite $+\infty$).

\item
La \Defn{série géométrique de raison $q \in \CC$} est
\begin{equation*}
\sum_{i=0}^\infty q^i
\end{equation*}
et converge absolument si $\abs{q} < 1$ et diverge si $\abs{q}
\geq 1$. De plus, pour tout naturel $n$,
\begin{equation*}
\sum_{i=0}^n q^i =  \frac{1-q^{n+1}}{1-q} \quad\text{et}\quad \sum_{i=0}^\infty q^i = \frac{1}{1-q}.
\end{equation*}
La seconde égalité est vraie si et seulement si $| q |<1$.

\item
Pour $\alpha \in \RR$, la série de Riemann (ou Dirichlet)
\begin{equation}		\label{EqSerRiem}
\sum_{i=1}^\infty \frac1{i^\alpha}
\end{equation}
converge (absolument, puisque réelle et positive) si et seulement
si $\alpha > 1$, et diverge sinon.
\end{enumerate}
\end{example}

\subsection{Critères de convergence}
\begin{remark}
  Étant donné le terme général d'une série, il est souvent --dans les cas qui nous intéressent-- difficile de déterminer la somme de la série. L'exemple de la série géométrique est particulier, puisqu'on connait une formule pour chaque somme partielle, mais pour l'exemple des séries de Riemann il n'y a aucune formule simple pour un $\alpha$ général. D'où l'intérêt d'avoir des critères de convergence ne nécessitant aucune connaissance de l'éventuelle limite de la série.
\end{remark}

\subsubsection{Critère de comparaison} 

Soient $\sum_i a_i$ et $\sum_j
b_j$ deux séries à termes positifs vérifiant
\begin{equation*}
  0 \leq a_i \leq b_i
\end{equation*}
alors
\begin{enumerate}
\item si $\sum_i a_i$ diverge, alors $\sum_j b_j$ diverge,
\item si $\sum_j b_j$ converge, alors $\sum_i a_i$ converge
  (absolument).
\end{enumerate}

\subsubsection{Critère d'équivalence}
\label{PgCritEquiv}

 Soient $\sum_i a_i$ et $\sum_j
b_j$ deux séries à termes positifs. Supposons l'existence de la limite
(éventuellement infinie) suivante
\begin{equation*}
  \limite i \infty \frac{a_i}{b_i} = \alpha \in \RR \text{ ou $\alpha =
    \infty$.}
\end{equation*}
Dans ce cas, nous avons
\begin{enumerate}
\item si $\alpha \neq 0$ et $\alpha\neq \infty$, alors
  \begin{equation*}
    \sum_i a_i \text{~converge} \ssi \sum_j b_j\text{~converge,}
  \end{equation*}
\item si $\alpha = 0$ et $\sum_j b_j$ converge, alors $\sum_i a_i$
  converge (absolument),
\item si $\alpha = +\infty$ et $\sum_j b_j$ diverge, alors $\sum_i
  a_i$ diverge.
\end{enumerate}

Une preuve est donnée à la page \pageref{PgPreuveCritEquiv}.

\subsubsection{Critère du quotient} Soit $\sum_i a_i$ une
série. Supposons l'existence de la limite (éventuellement infinie)
suivante
\begin{equation*}
  \limite i \infty \abs{\frac{a_{i+1}}{a_i}} = L\in \RR \text{ ou $L =
    \infty$.}
\end{equation*}
Alors
\begin{enumerate}
\item si $L < 1$, la série converge absolument,
\item si $L > 1$, la série diverge,
\item si $L = 1$ le critère échoue : il existe des exemple de convergence et des exemples de divergence.
\end{enumerate}

Une preuve est donnée à la page \pageref{PgPreuveCritQuotient}.

\subsubsection{Critère de la racine} Soit $\sum_i a_i$ une série, et
considérons
\begin{equation*}
  \limsup_{i \rightarrow \infty} \sqrt[i]{\abs{a_i}} = L \in \RR
  \text{ ou $L =
    \infty$.}
\end{equation*}
Alors
\begin{enumerate}
\item si $L < 1$, la série converge absolument,
\item si $L> 1$, la série diverge,
\item si $L = 1$ le critère échoue.
\end{enumerate}

Une preuve est donnée à la page \pageref{PgPreuveCritRacine}.

Les critères de comparaison, d'équivalence, du quotient et de la racine sont des critères de convergence absolue. Pour conclure à une convergence simple qui n'est pas une convergence absolue, le critère d'Abel sera notre outil principal.  

\subsubsection{Critère d'Abel}


\begin{proposition}[Proposition 5, page 122]

Soit la série $\sum_i c_iz_i$ avec
\begin{enumerate}
\item $(c_i)$ est une suite réelle décroissante qui tend vers zéro,
\item $(z_i)$ est une suite dans $\eC$ dont la suite des sommes partielles est bornée dans $\eC$, c'est à dire qu'il existe un $M>0$ tel que pour tout $n$,
\begin{equation}
	\left| \sum_{i=1}^nz_i \right| \leq M.
\end{equation}
Alors la série $\sum_ic_iz_i$ est convergente.

\end{enumerate}
\end{proposition}
Remarquons que ce critère ne donne pas de convergence absolue.


\begin{example}
  Le \Defn{critère des séries alternées} est un cas particulier de ce
  critère : si $c_i$ est une suite décroissante à limite nulle, alors
  la série
  \begin{equation*}
    \sum_{i=0}^\infty {(-1)}^i c_i
  \end{equation*}
  converge simplement.
\end{example}

%---------------------------------------------------------------------------------------------------------------------------
					\subsection{Série de puissances}
%---------------------------------------------------------------------------------------------------------------------------

Une \defe{série de puissance}{Série!de puissance} est une série de la forme
\begin{equation}		\label{eqseriepuissance}
	\sum_{k=0}^{\infty}c_k(z-z_0)^k
\end{equation}
où $z_0\in \eC$ est fixé, $(c_k)$ est une suite complexe fixée, et $z$ est un paramètre complexe. Nous disons que cette série est \emph{centrée} en $z_0$.

\begin{theorem}		\label{ThoSerPuissRap}
Considérons la série de puissances donnée par le terme général $c_k(z-z_0)^k$. Si nous posons
\begin{equation}		\label{EqRayCOnvSer}
	\alpha=\frac{1}{ R } =\limsup\sqrt[k]{| c_k |}
\end{equation}
alors la série converge absolument si $| z-z_0 |<R$ et diverge si $| z-z_0 |>R$.

De plus, ce rayon de convergence peut être calculé par la formule alternative 
\begin{equation}		\label{EqAlphaSerPuissAtern}
	\alpha=\frac{1}{ R }=\limite k \infty \abs{\frac{c_{k+1}}{c_k}}
\end{equation}
lorsque $c_k$ est non nul à partir d'un certain $k$.
\end{theorem}
Le disque $| z-z_0 |\leq R$ est le \defe{disque de convergence}{Disque de convergence} de la série \eqref{eqseriepuissance}. Notez que ce théorème ne dit rien pour les points tels que $| z-z_0 |=R$. Il faut traiter ces points au cas par cas. Et le pire, c'est qu'une série donnée peut converger pour certain des points sur le bord du disque, et diverger en d'autres.  Il y a un dessin à la figure \ref{LabelFigDisqueConv}.
\newcommand{\CaptionFigDisqueConv}{À l'intérieur du disque de convergence, la convergence est absolue. En dehors, la série diverge. Sur le cercle proprement dit, tout peut arriver.}
\input{Fig_DisqueConv.pstricks}

Le disque de centre $z_0$ et de rayon $R$ est appelé \Defn{disque de convergence}. Pour un complexe $z$ sur le bord de ce disque, c'est-à-dire tel que $\abs{z-z_0} = R$, le comportement peut-être très varié (convergence absolue, convergence simple ou divergence) et n'est éventuellement pas le même sur tout le bord.

L'étude de ce qu'il se passe sur le bord du disque de convergence commence par y étudier la convergence absolue, c'est à dire étudier la série
\begin{equation}
	\sum_k| c_k(z-z_0)^k |=\sum_k| c_k |R^k
\end{equation}
parce que sur le bord, $| z-z_0 |=R$. L'étude du terme général $| c_k |R^k$ a deux utilités :
\begin{enumerate}
\item Si la somme $\sum_{k}| c_k |R^k$ converge, alors la série converge uniformément sur le bord,
\item si la suite $| c_k |R^k$ ne tend pas vers zéro, alors la série ne converge même pas simplement sur le bord.
\end{enumerate}


\section{Continuité et dérivabilité}

\label{seccontetderiv}
\subsection{Rappels et outils}
On considère dans la suite une fonction $f : A \to \eR$, où $a \in A \subset \eR$ ; cependant, les notions de continuité et de dérivabilité se généralisent immédiatement au cas de fonctions à valeurs vectorielles ; la notion de continuité se généralise au cas des fonctions à plusieurs variables (la notion de dérivabilité est remplacée par celle de différentiabilité dans ce cadre).

\begin{defn}La fonction $f$ est \Defn{continue en $a$} si
  \begin{equation*}
    \limite x a f(x)
  \end{equation*}
  existe (et vaut alors $f(a)$).
\end{defn}

\begin{defn}La fonction $f$ est \Defn{dérivable en $a$} si $a \in
  \operatorname{int} A$ et si
  \begin{equation*}
    \lim_{\substack{x\rightarrow a\\x\neq a}} \frac{f(x)-f(a)}{x-a}
  \end{equation*}
  existe. On note alors cette quantité $f^\prime(a)$, c'est le nombre
  dérivé de $f$ en $a$. La \Defn{fonction dérivée} de $f$ est
  \begin{equation*}
    f^\prime : A^\prime \to \eR : a \mapsto f^\prime(a)
  \end{equation*}
  définie sur l'ensemble noté $A^\prime$ des points $a$ où $f$ est
  dérivable.
\end{defn}

\begin{defn}
  Une fonction est \Defn{continue} (resp. \Defn{dérivable}) si elle
  est continue (resp. dérivable) en tout point $a \in A$ de son
  domaine.
\end{defn}

\begin{example}
  Montrons que la fonction $f : \eR \to \eR : x\mapsto x$ est
  continue et dérivable.
  \begin{description}
  \item[Continuité] Soit $a \in \eR$, et montrons
    $\limite x a x = a$. En effet, soit $\epsilon > 0$, et choisissons
    $\delta = \epsilon$. Si $\abs{x-a} < \delta$ alors
    \begin{equation*}
      \abs{f(x) - f(a)} = \abs{x-a} < \delta = \epsilon
    \end{equation*}
    ce qui prouve la continuité.

  \item[Dérivabilité] Soit $a \in \eR$. Calculons la
    limite du quotient différentiel
    \begin{equation*}
      \limite[x\neq a]{x}{a} \frac{x-a}{x-a} = \limite[x\neq a]x a 1 = 1
    \end{equation*}
    ce qui prouve que $f$ est dérivable et que sa dérivée vaut $1$ en
    tout point $a$ de $\eR$.
  \end{description}

% Exceptionellement (bien qu'on sache que la dérivabilité implique la
% continuité), montrons ces deux assertions séparément.

% Il faut d'abord montrer que pour tout $a \in \eR$,
% \begin{equation*}
%   \limite x a x = a
% \end{equation*}
% c'est-à-dire
% \begin{equation*}
%   \forall \epsilon > 0, \exists \delta > 0 :  \forall x \in \eR \abs{x-a} <
%   \delta \Rightarrow \abs{x-a} < \epsilon
% \end{equation*}
% ce qui est clair en prenant $\delta = \epsilon$.

% Il faut ensuite montrer que pour tout $a \in \eR$,
% \begin{equation*}
%   \limite[x\neq a] x a \frac{x-a}{x-a} \qquad\text{ existe.}
% \end{equation*}
% ce qui, là encore, est clair, puisque la fonction $\frac{x-a}{x-a} =
% 1$ pour tout $x \neq a$, et il s'agit donc de la limite de la
% fonction constante $1$, qui vaut évidemment $1$.

% On a donc montré que la fonction $x \mapsto x$ est continue,
% dérivable, et que sa dérivée vaut $1$ en tout point $a$ de son
% domaine.

\end{example}


%+++++++++++++++++++++++++++++++++++++++++++++++++++++++++++++++++++++++++++++++++++++++++++++++++++++++++++++++++++++++++++
					\section{Développements de Taylor et Maclaurin}
%+++++++++++++++++++++++++++++++++++++++++++++++++++++++++++++++++++++++++++++++++++++++++++++++++++++++++++++++++++++++++++

Nous voulons formaliser l'idée d'une fonction qui tend vers zéro \og plus vite\fg{} qu'une autre. Nous disons que $f\in o\big(\varphi(x)\big)$ si
\begin{equation}
	\lim_{x\to 0} \frac{ f(x) }{ \varphi(x) }=0.
\end{equation}
En particulier, nous disons que $f\in o(x)$ lorsque $\lim_{x\to 0} f(x)/x=0$.

La proposition suivante se trouve à la page 216.
\begin{proposition}		\label{PropDevTaylorPol}
Soient $f\colon I\subset\eR\to \eR$ et $a\in\Int(I)$. Soit un entier $k\geq 1$. Si $f$ est $k$ fois dérivable en $a$, alors il existe un et un seul polynôme $P$ de degré $\leq k$ tel que
\begin{equation}
	f(x)-P(x-a)\in o\big( | x-a |^k \big)
\end{equation}
lorsque $x\to a$, $x\neq a$. Ce polynôme  est donné par
\begin{equation}
	P(h)=f(a)+f'(a)h+\frac{ f''(a) }{ 2! }h^2+\ldots+\frac{ f^{(k)}(a) }{ k! }h^k.
\end{equation}
\end{proposition}

La proposition suivant, de la page 218, donne une intéressante façon de trouver le reste d'un développement de Taylor.
\begin{proposition}		\label{PropResteTaylorc}
Soient $I$, un intervalle dans $\eR$ et $f\colon I\to \eR$ une fonction de classe $C^k$ sur $I$ telle que $f^{(k+1)}$ existe sur $I$. Soient $a\in\Int(I)$ et $x\in I$. Alors il existe $c$ strictement compris entre $x$ et $a$ tel que 
\begin{equation}
	R_{f,a,k}(x)=\frac{ f^{(k+1)}(c) }{ (k+1)! }(x-a)^{k+1}.
\end{equation}
\end{proposition}

%---------------------------------------------------------------------------------------------------------------------------
					\subsection{Exemple : un calcul heuristique de limite}
%---------------------------------------------------------------------------------------------------------------------------
\label{SubSecCalcLimHeuris}

Au cours de la résolution de l'exercice \ref{exoEqsDiff0002}\ref{ItemfEqsDiff00002}, nous devons calculer la limite suivante :
\begin{equation}
	\lim_{x\to 0} \frac{  e^{-2\cos(x)+2}\sin(x) }{ \sqrt{ e^{2\cos(x)+2}}-1 }.
\end{equation}
La stratégie que nous allons suivre pour calculer cette limite est de développer certaines parties de l'expression en série de Taylor, afin de simplifier l'expression. La première chose à faire est de remplacer $ e^{y(x)}$ par $1+y(x)$ lorsque $y(x)\to 0$. La limite devient
\begin{equation}
	\lim_{x\to 0} \frac{ \big( -2\cos(x)+3 \big)\sin(x) }{ \sqrt{-2\cos(x)+2} }.
\end{equation}
Nous allons maintenant remplacer $\cos(x)$ par $1$ au numérateur et par $1-x^2/2$ au dénominateur. Pourquoi ? Parce que le cosinus du dénominateur est dans une racine, donc nous nous attendons à ce que le terme de degré deux du cosinus donne un degré un en dehors de la racine, alors que du degré un est exactement ce que nous avons au numérateur : le développement du sinus commence par $x$.

Nous calculons donc
\begin{equation}
	\begin{aligned}[]
		\lim_{x\to 0} \frac{ \sin(x) }{ \sqrt{-2\left( 1-\frac{ x^2 }{ 2 } \right)+2} }=\lim_{x\to 0} \frac{ \sin(x) }{ x }=1.
	\end{aligned}
\end{equation}
Tout ceci n'est évidement pas très rigoureux, mais en principe vous avez tous les éléments en main pour justifier les étapes.
%+++++++++++++++++++++++++++++++++++++++++++++++++++++++++++++++++++++++++++++++++++++++++++++++++++++++++++++++++++++++++++
					\section{Équations différentielles}
%+++++++++++++++++++++++++++++++++++++++++++++++++++++++++++++++++++++++++++++++++++++++++++++++++++++++++++++++++++++++++++

%---------------------------------------------------------------------------------------------------------------------------
					\subsection{Équations à variables séparées}
%---------------------------------------------------------------------------------------------------------------------------

Une \defe{équation à variables séparées}{Équation!différentielle!à variables séparées} est une équation  de la forme
\begin{equation}		\label{EqDiffSeparee}
	y'=u(t)f(y).
\end{equation}
Les deux résultats qui résolvent ce cas sont donnés aux pages 311 et 313 du cours de première année.
\begin{proposition}
Nous considérons l'équation \eqref{EqDiffSeparee} avec $u(t)$ continue sur $I$ et $f$ continue sur $J$ avec $f(\eta)\neq 0$ pour tout $\eta\in J$. Soit $U$, une primitive de $u$ sur $I$, et $G$, une primitive de $1/f$ sur $J$.

Si $y\colon Y'\to J$ est une fonction sur un intervalle $I'\subset I$, alors $y$ est solution de l'équation \eqref{EqDiffSeparee} si et seulement si il existe $C\in\eR$ tel que
\begin{equation}		\label{EqSoluceEqDiffSep}
	G\big( y(t) \big)=U(t)+C.
\end{equation}
\end{proposition}
Cette proposition dit que toutes les solutions qui ne s'annulent jamais sur un intervalle ont la forme $G\big( y(t) \big)=U(t)+C$ et peuvent donc être trouvées en calculant des primitives.

La formule \eqref{EqSoluceEqDiffSep} peut être obtenue de la façon heuristique suivante, en écrivant $y'=dy/dt$, et en passant le $dt$ à droite. Nous trouvons successivement
\begin{equation}
	\begin{aligned}[]
		y'&=u(t)f(y)\\
		dy&=u(t)f(y)dt\\
		\frac{ dy }{ f(y) }&=u(t)dt\\
		\int\frac{ dy }{ f(y) }&=\int u(t)dt\\
		G(y)&=U(t)+C.
	\end{aligned}
\end{equation}

\begin{proposition}
Soient $u$ continue sur $I$ et $f$ continue sur $J$, et $f(\eta)\neq 0$ sur $J$. Soient $t_0\in I$ et $y_0\in J$. Alors il existe $I'\subset I$ avec $t_0\in I'$ et $f\in C^1(I'\to J)$ tels que
\begin{enumerate}

\item
$y$ est solution de \eqref{EqDiffSeparee} sur $I'$ et vérifie $y(t_0)=y_0$,
\item
si $z$ est une solution de \eqref{EqDiffSeparee} sur $I''\subset I'$ avec $t_0\in I''$ et $z(t_0)=y_0$, alors $I''\subset I'$ et $z(t)=y(t)$ pour tout $t\in I''$.

\end{enumerate}
\end{proposition}

%---------------------------------------------------------------------------------------------------------------------------
					\subsection{Équations linéaires}
%---------------------------------------------------------------------------------------------------------------------------

Une \defe{équation différentielle linéaire}{Équation différentielle!linéaire} est une équation de la forme
\begin{equation}
	y'+u(t)y=v(t).
\end{equation}
La technique pour résoudre cette équation est de commencer par résoudre l'équation homogène associée. Si $U(t)$ est une primitive de $u(t)$, nous avons
\begin{equation}
	\begin{aligned}[]
		y'_H(t)+u(t)y_H(t)&=0\\
		\frac{ y'_H }{ y_H }&=-u(t)\\
		\ln(y_H)&=-U(t)+C\\
		y_H(t)&= e^{-U(t)+C}=K e^{-U(t)}
	\end{aligned}
\end{equation}
où $K= e^{C}$.

Cela fournit la solution générale de l'équation homogène. Il existe un truc génial qui permet d'en tirer la solution générale du système non homogène. Lorsque nous avons trouvé $y_H(t)=K e^{-U(t)}$, le symbole $K$ désigne une constante. La méthode de \defe{variation des constantes}{variation des constantes} consiste à essayer la solution
\begin{equation}		\label{EqEssayVarSctr}
	y(t)=K(t) e^{-U(t)},
\end{equation}
c'est à dire à dire que la constante est en réalité une fonction. Afin de trouver quelle fonction $K(t)$ fait en sorte que l'essai \eqref{EqEssayVarSctr} soit une solution, nous la remplaçons dans l'équation de départ $y'+uy=v$. Maintenant,
\begin{equation}
	y'(t)=K'(t) e^{-U(t)}-K(t)u(t) e^{-U(t)}.
\end{equation}
En remettant dans l'équation,
\begin{equation}
	y'+uy=K' e^{-U}-Ku e^{-U}+uK e^{-U}=K' e^{-U}=v.
\end{equation}
Notez que les termes en $K$ se sont miraculeusement simplifiés. Cela est directement dû au fait que $ e^{-U}$ est solution de l'équation homogène. Nous restons avec l'équation
\begin{equation}
	K'=\frac{ v }{  e^{-U} }
\end{equation}
pour $K(t)$. La solution générale du problème non homogène est donc finalement donnée par
\begin{equation}
	y(t)=\big( W(t)+C \big) e^{-U(t)}
\end{equation}
si $W(t)$ est une primitive de $v(t)e^{U(t)}$.

Tout ceci est un peu heuristique. La proposition suivante (page 308) dit dans quels cas ça fonctionne.
\begin{proposition}
Soient $u$ et $v$ continues sur $I$ et $U$, une primitive de $u$ sur $I$ et $W$ une primitive de $v e^{-U}$ sur $I$. Une fonction $y\colon I\to \eR$ est solution de $y'+u(t)y=v(t)$ si et seulement si il existe une constante $C\in \eR$ telle que
\begin{equation}
	y(t)=\big( W(t)+C \big) e^{U(t)}
\end{equation}
pour tout $t\in I$.
\end{proposition}
Le corollaire de la page $310$ lui règle son compte au problème de Cauchy. Si $t_0\in I$ et si $y_0\in\eR$, alors il existe une unique solution $y$ sur $I$ à l'équation linéaire telle que $y(t_0)=y_0$.

%+++++++++++++++++++++++++++++++++++++++++++++++++++++++++++++++++++++++++++++++++++++++++++++++++++++++++++++++++++++++++++
					\section{Topologie en général}
%+++++++++++++++++++++++++++++++++++++++++++++++++++++++++++++++++++++++++++++++++++++++++++++++++++++++++++++++++++++++++++

Cette section sur la topologie en général peut être sautée dans un premier temps. Celles et ceux qui ont dans l'idée de faire de la mathématique plus avancée dans les années à venir devraient toutefois y jeter un œil.

\begin{definition}		\label{DefTopologieGene}

Soit $E$, un ensemble et $\mT$, une partie de l'ensemble de ses parties qui vérifie les propriétés suivantes
\begin{enumerate}

\item
les ensembles $\emptyset$ et $E$ sont dans $\mT$,

\item
Si $I$ est n'importe quel ensemble et si pour tout $i\in I$, nous avons un élément $\mO_i\in\mT$, alors $\cup_{i\in I}\mO_i\in\mT$,

\item
Si $J$ est un ensemble fini et si pour tout $j\in J$, nous avons un élément $\mO_j\in\mT$, alors $\cap_{j\in J}\mO_j\in\mT$.

\end{enumerate}
Les deux dernières propriétés s'énoncent en disant que toute réunions d'éléments de $\mT$ est un élément de $\mT$ et que toute intersection \emph{finie} d'éléments de $\mT$ est un élément de $\mT$.

Un tel choix $\mT$ de sous-ensembles de $E$ est une  \defe{\href{http://fr.wikipedia.org/wiki/Espace_topologique}{topologie}}{Topologie} sur $E$, et les éléments de $\mT$ sont appelés des \defe{ouverts}{ouvert}. Nous disons que un sous ensemble $A$ de $E$ est \defe{fermé}{fermé} si son complémentaire, $A^c$ est ouvert.
\end{definition}

Dès que nous avons une topologie, nous avons une notion de convergence de suite : nous disons qu'une suite $x_n$ d'éléments de $E$ \defe{converge}{convergence!en topologie} vers l'élément $x$ de $E$ si pour tout ouvert $\mO$ contenant $x$, il existe un $K$ tel que $k>K$ implique $x_k\in\mO$. Cette définition est exactement celle donnée pour la convergence de suites dans $\eR^n$, à part que nous avons remplacé le mot \og boule\fg{} par \og ouvert\fg.

Dans un espace topologique, nous avons une caractérisation très importante des ouverts.
\begin{theorem}		\label{ThoPartieOUvpartouv}
Une partie $A$ de $E$ est ouverte si et seulement si pour tout $x\in A$, il existe un ouvert autour de $x$ contenu dans $A$.
\end{theorem}

\begin{proof}
Le sens direct est évident : $A$ lui-même est un ouvert autour de $x\in A$, qui est inclus à $A$.

Pour le sens inverse, pour chaque $x\in A$, nous considérons l'ensemble $\mO_x\subset A$, un ouvert autour de $x$. Nous avons que
\begin{equation}	\label{EqAUniondesOx}
	A=\bigcup_{x\in A}\mO_x.
\end{equation}
En effet $A\subset\cup_{x\in A}\mO_x$ parce que tous les éléments de $A$ sont dans un des $\mO_x$, par construction. D'autre part, $\cup_{x\in A}\mO_x\subset A$ parce que chacun des $\mO_x$ est compris dans $A$.

L'union du membre de droite de \eqref{EqAUniondesOx} est une union d'ouverts et est donc un ouvert. Cela prouve que $A$ est un ouvert.

\end{proof}

Une utilisation typique de ce théorème est faite à l'exercice \ref{exo0083}.

%%%%%%%%%%%%%%%%%%%%%%%%%%
%
   \section{Topologie dans \texorpdfstring{$\eR^n$}{Rn}}
%
%%%%%%%%%%%%%%%%%%%%%%%%



Dans cette section, nous travaillons dans l'espace $\eR^n$ pour un certain naturel $n$. Nous y définissons la notion d'ouvert et de fermé, qui sont la base de la topologie générale. Notons que ces définitions n'ont de sens que relativement à l'espace ambiant, aussi un ouvert de $\eR$ ne sera en général pas un ouvert de $\eR^2$~: d'une part, il n'y a pas d'inclusion canonique de $\eR$ dans $\eR^2$ (les ouverts du second ne sont même pas des sous-ensembles du premier) et, d'autre part, les définitions se basent sur la notion de boule de $\eR^n$ qui dépend évidemment de la valeur de $n$ (une boule dans $\eR$ est un intervalle, dans $\eR^2$ c'est un disque, etc.)

%---------------------------------------------------------------------------------------------------------------------------
					\subsection{Ouverts et fermés}
%---------------------------------------------------------------------------------------------------------------------------

\begin{definition}
	La \defe{boule ouverte}{Boule!ouverte} de centre $x_0 \in \eR^n$ et de rayon $r \in
	\eR^+$ est définie par
	\begin{equation}
		B(x_0,r) = \{ x \in \eR^n \tq \norme{x - x_0} < r \},
	\end{equation}
	tandis que la \defe{boule fermée}{Boule!fermée} de centre $x_0$ et de rayon $r$ est
	\begin{equation}
		\bar B(x_0,r) = \{ x \in \eR^n \tq \norme{x - x_0} \leq r \};
	\end{equation}
	la différence est que l'inégalité dans la première est stricte.
\end{definition}

%---------------------------------------------------------------------------------------------------------------------------
					\subsection{Intérieur, adhérence et frontière}
%---------------------------------------------------------------------------------------------------------------------------

\begin{definition}
  Soit $A \subset \eR^n$ et $x \in \eR^n$. Le point $x$ est \defe{intérieur}{intérieur} à $A$ si il existe une boule autour de $x$ complètement contenue dans $A$. L'ensemble des points intérieurs à $A$ est noté $\interieur A$ ou $\mathring A$, de sorte qu'on a précisément
  \begin{equation*}
    x \in \interieur A \iffdefn  \exists \epsilon > 0 \tq
    B(x,\epsilon) \subset A.
  \end{equation*}
\end{definition}


\begin{definition}
Le point $x$ est dans l'\defe{adhérence}{adhérence} de $A$ si toute boule autour de $x$ intersecte $A$. L'ensemble de ces points est noté $\adh A$ ou $\bar A$, et on a donc de manière plus précise
\begin{equation}
	x \in \adh A \iffdefn \forall \epsilon > 0, B(x,\epsilon) \cap A \neq \emptyset
\end{equation}
\end{definition}

\begin{proposition}
Pour $A \subset \eR^n$, nous avons
\begin{equation*}
	\interieur A \subseteq A  \subseteq \adh A
\end{equation*}
\end{proposition}

\begin{definition}
  La \defe{frontière}{frontière} ou le \defe{bord}{bord} de $A$ est défini par $\partial A = \adh A \setminus \interieur A$. L'ensemble $A$ est un \defe{ouvert}{ouvert} si $A = \interieur A$, et c'est un \defe{fermé}{fermé} si $A = \adh A$.
\end{definition}

On vérifiera que les notations et les dénominations sont cohérentes en
prouvant la proposition suivante.
\begin{proposition}Pour $\epsilon > 0$,
  \begin{enumerate}
  \item l'adhérence de $B(x,\epsilon)$ est $\bar B(x,\epsilon)$,
  \item l'intérieur de $\bar B(x,\epsilon)$ est $B(x,\epsilon)$,
  \item la boule ouverte $B(x,\epsilon)$ est un ouvert,
  \item la boule fermée $\bar B(x,\epsilon)$ est un fermé.
  \end{enumerate}
\end{proposition}

Nous avons également les liens suivants entre intérieur, adhérence,
ouvert, fermé et passage au complémentaire (noté ${}^c$)~:
\begin{proposition}
Si $A \subset \eR^n$ et $A^c = \eR^n\setminus A$, nous
  avons
  \begin{enumerate}
  \item $(\interieur A)^c = \adh (A^c)$ et $(\adh A)^c = \interieur
    (A^c)$,
  \item $A$ est ouvert si et seulement si $A^c$ est fermé,
  \item $\interieur A$ est le plus grand ouvert contenu dans $A$,
  \item $\adh A$ est le plus petit fermé contenant $A$,
    % \item
  \end{enumerate}
\end{proposition}

%---------------------------------------------------------------------------------------------------------------------------
					\subsection{Bornés et compacts}
%---------------------------------------------------------------------------------------------------------------------------


\begin{definition}
  Un sous ensemble $A \subset \eR^n$ est \defe{borné}{borné} si il existe une boule de $\eR^n$ contenant $A$.
\end{definition}

\begin{proposition}
  Toute réunion finie d'ensembles bornés est un ensemble borné. Toute
  partie d'un ensemble borné est un ensemble borné.
\end{proposition}

\begin{definition}
  La partie $A \subset \eR^n$ est \defe{compacte}{compact} si et seulement si, pour tout
  recouvrement de $A$ par des ouverts (c'est-à-dire une collection
  d'ouverts dont la réunion contient $A$) on peut tirer un
  recouvrement fini.
\end{definition}

% En particulier, si on recouvre $A$ par l'ensemble des boules
% $B(x,1)$ où $x$ parcourt $A$ (de sorte que tout point de $A$ est
% dans \og sa\fg{} boule, et donc la réunion des boules recouvre bien
% $A$), on doit pouvoir en tirer un recouvrement fini, c'est-à-dire
% des boules $B(x_1,1), B(x_2,1), \ldots, B(x_k,1)$ (avec $k$ un
% naturel) dont la réunion contient $A$.

\begin{proposition}
Une partie de $\eR^n$ est compacte si et seulement si elle est fermée et bornée.
\end{proposition}

%---------------------------------------------------------------------------------------------------------------------------
					\subsection{Connexité}
%---------------------------------------------------------------------------------------------------------------------------

\begin{definition}
  Le sous ensemble $A \subset \eR^n$ est \defe{connexe par arcs}{Connexe!par arc} si pour tout $x, y \in
  A$, il existe un chemin\footnote{Attention : ici quand on dit \emph{chemin}, on demande que l'application soit continue. Dans de nombreux cours de géométrie différentielle, on demande $ C^{\infty}$. Il faut s'adapter au contexte.} contenu dans $A$ les reliant, c'est-à-dire
  une application continue
  \begin{equation*}
    \gamma : [0,1] \to \eR^n \tq \gamma(0) = x~\text{et}~\gamma(1) = y
  \end{equation*}
  avec $\gamma(t) \in A$ pour tout $t\in [0,1]$.
\end{definition}

%+++++++++++++++++++++++++++++++++++++++++++++++++++++++++++++++++++++++++++++++++++++++++++++++++++++++++++++++++++++++++++
					\section{Topologie des espaces métriques}
%+++++++++++++++++++++++++++++++++++++++++++++++++++++++++++++++++++++++++++++++++++++++++++++++++++++++++++++++++++++++++++

Si $E$ est un ensemble, une \defe{distance}{distance} sur $E$ est une application $d\colon E\times E\to \eR$ telle que pour tout $x,y\in E$,
\begin{enumerate}

\item
$d(x,y)\geq 0$

\item
$d(x,y)=0$ si et seulement si $x=y$,

\item
$d(x,y)=d(y,x)$

\item
$d(x,y)\leq d(x,z)+d(z,y)$.

\end{enumerate}
La dernière condition est l'\defe{inégalité triangulaire}{Inégalité!triangulaire}. Le couple $(E,d)$ d'un ensemble et d'une métrique est un \defe{espace métrique}{Espace!métrique}.

Dès que l'ensemble $E$ est muni d'une distance, nous définissons une topologie en disant que les boules
\begin{equation}
	B(x,r)=\{ y\in E\tq d(x,y)<r \}
\end{equation}
sont ouvertes.


%+++++++++++++++++++++++++++++++++++++++++++++++++++++++++++++++++++++++++++++++++++++++++++++++++++++++++++++++++++++++++++
					\section{Uniforme continuité}
%+++++++++++++++++++++++++++++++++++++++++++++++++++++++++++++++++++++++++++++++++++++++++++++++++++++++++++++++++++++++++++

Nous avons la proposition de la page $181$.
\begin{proposition}	\label{PropoInvCompCont}
Soit $f\colon A\subset\eR^n\to B\subset\eR^m$ une bijection continue. Si $A$ est compact, alors $f^{-1}\colon B\to A$ est continue.
\end{proposition}

La proposition suivante est la proposition 3, à la page $170$.
\begin{proposition}		\label{PropIntContMOnIvCont}
Soient $I$ un intervalle dans $\eR$ et $f\colon I\to \eR$ une fonction continue strictement monotone. Alors la fonction réciproque $f^{-1}\colon f(I)\to \eR$ est continue sur l'intervalle $f(I)$.
\end{proposition}



%+++++++++++++++++++++++++++++++++++++++++++++++++++++++++++++++++++++++++++++++++++++++++++++++++++++++++++++++++++++++++++
					\section{Maximisation sans contraintes}
%+++++++++++++++++++++++++++++++++++++++++++++++++++++++++++++++++++++++++++++++++++++++++++++++++++++++++++++++++++++++++++

%---------------------------------------------------------------------------------------------------------------------------
					\subsection{Fonctions à une seule variable}
%---------------------------------------------------------------------------------------------------------------------------

\begin{definition}
Soit $f\colon A\subset \eR\to \eR$ et $a\in A$. Le point $a$ est un \defe{maximum local}{Maximum!local} de $f$ si il existe un voisinage $\mU$ de $a$ tel que $f(a)\geq f(x)$ pour tout $x\in\mU\cap A$. Le point $a$ est un \defe{maximum global}{Maximum!global} si $f(a)\geq g(x)$ pour tout $x\in A$.
\end{definition}

La proposition basique à utiliser lors de la recherche d'extrema est la suivante :
\begin{proposition}
Soit $f\colon A\subset\eR\to \eR$ et $a\in\Int(A)$. Supposons que $f$ est dérivable en $a$. Si $a$ est un \href{http://fr.wikipedia.org/wiki/Extremum}{extremum} local, alors $f'(a)=0$.
\end{proposition}

La réciproque n'est pas vraie, comme le montre l'exemple de la fonction $x\mapsto x^3$ en $x=0$ : sa dérivée est nulle et pourtant $x=0$ n'est ni un maximum ni un minimum local. 

Cette proposition ne sert donc qu'à sélectionner des \emph{candidats} extremum. Afin de savoir si ces candidats sont des extrema, il y a la proposition suivante (proposition 1, page 227).

\begin{proposition}
Soit $f\colon I\subset \eR\to \eR$, une fonction de classe $C^k$ au voisinage d'un point $a\in\Int I$. Supposons que
\begin{equation}
	f'(a)=f''(a)=\ldots=f^{(k-1)}(a)=0,
\end{equation}
et que
\begin{equation}
	f^{(k)}(a)\neq 0.
\end{equation}
Dans ce cas,
\begin{enumerate}

\item
Si $k$ est pair, alors $a$ est un point d'extremum local de $f$, c'est un minimum si $f^{(k)}(a)>0$, et un maximum si $f^{(k)}(a)<0$,
\item
Si $k$ est impair, alors $a$ n'est pas un extremum local de $f$.

\end{enumerate}
\end{proposition}

Note : jusqu'à présent nous n'avons rien dit des extrema \emph{globaux} de $f$. Il n'y a pas grand chose à en dire. Si un point d'extremum global est situé dans l'intérieur du domaine de $f$, alors il sera extremum local (a fortiori). Ou alors, le maximum global peut être sur le bord du domaine. C'est ce qui arrive à des fonctions strictement croissantes sur un domaine compact.

Une seule certitude : si une fonction est continue sur un compact, elle possède une minimum et un maximum global.

%---------------------------------------------------------------------------------------------------------------------------
					\subsection{Fonctions à plusieurs variables}
%---------------------------------------------------------------------------------------------------------------------------


%---------------------------------------------------------------------------------------------------------------------------
					\subsection{Quelque mots à propos de matrices}
%---------------------------------------------------------------------------------------------------------------------------

Les notions qui suivent seront vues au cours de géométrie lorsqu'il sera temps. 

Si $g$ est une application bilinéaire sur $\eR^2$, nous disons qu'elle est
\begin{enumerate}

\item
\defe{Définie positive}{Application!définie positive} si $g(u,u)\geq 0$ pour tout $u\in\eR^2$ et $g(u,u)=0$ si et seulement si $u=0$.

\item
\defe{semi-définie positive}{Application!semi-définie positive} si $g(u,u)\geq 0$ pour tout $u\in\eR^2$. 

\end{enumerate}

Une matrice $M$ est définie positive si $v^tMv>0$ pour tout $v\neq 0$, en particulier si $v$ est un vecteur propre de valeur propre $\lambda$ (c'est à dire si $Mv=\lambda v$), alors $\lambda v^tv>0$, et donc $\lambda>0$. Donc $M$ sera définie positive si toutes ses valeurs propres sont positives.

\begin{proposition}
	Soit $M$, une matrice $2\times 2$ symétrique\footnote{la matrice $d^2f(a)$ est toujours symétrique quand $f$ est de classe $C^2$.}. Nous avons
	\begin{enumerate}
		\item
		$\det M>0$ et $\tr(M)>0$ implique $M$ définie positive,
		\item
		$\det M>0$ et $\tr(M)<0$ implique $M$ définie négative,
		\item
		$\det M<0$ implique ni semi définie positive, ni définie négative (et donc pas un extrema dans le cas où $M=d^2f(a)$ par le point \ref{ItemPropoExtreRn} de la proposition \ref{PropoExtreRn}),
		\item
		$\det M=0$ implique $M$ semi-définie positive ou semi-définie négative.
	\end{enumerate}
\end{proposition}
 

%---------------------------------------------------------------------------------------------------------------------------
					\subsection{Les théorèmes}
%---------------------------------------------------------------------------------------------------------------------------



Un point $a$ à l'intérieur du domaine d'une fonction $f\colon A\subset\eR^n\to \eR$ est une \defe{point critique}{Critique!point} de $f$ lorsque $df(a)=0$. Ces points sont analogues aux points où la dérivée d'une fonction sur $\eR$ s'annule. Les points critiques de $f$ sont dons les candidats à être des points d'extremum.

\begin{proposition}		\label{PropoExtreRn}
Soit $f\colon A\subset\eR^n\to \eR$ une fonction de classe $C^2$ au voisinage de $a\in\Int(A)$.
\begin{enumerate}

\item
Si $a$ est un point critique de $f$, et si $d^2f(a)$ est \href{http://fr.wikipedia.org/wiki/Matrice_définie_positive}{définie positive}, alors $a$ est un minimum local strict de $f$,
\item		\label{ItemPropoExtreRn}
Si $a$ est un minimum local, alors $a$ est un point critique et $d^2f(a)$ est semi-définie positive.

\end{enumerate}
\end{proposition}

La seconde partie de l'énoncé est tout à fait comparable au fait bien connu que, pour une fonction $f\colon \eR\to \eR$, si le point $a$ est minimum local, alors $f'(a)=0$ et $f''(a)\geq 0$. Notez le fait que l'inégalité n'est pas stricte, ce qui correspond à $d^2f(a)$ \emph{semi}-definie positive.

Pour rappel, dans le cas d'une fonction à deux variables, $d^2f(a)$ est la matrice (et donc l'application linéaire)
\begin{equation}
	d^2f(a)=\begin{pmatrix}
	\frac{ d^2f  }{ dx^2 }(a)	&	\frac{ d^2f  }{ dx\,dy }(a)	\\ 
	\frac{ d^2f  }{ dy\,dx }(a) 	&	\frac{ d^2f  }{ dy^2 }(a)
\end{pmatrix}.
\end{equation}
Dans le cas d'une fonction $C2$, cette matrice est symétrique.

La méthode pour chercher les extrema de $f$ est donc de suivre le points suivants :
\begin{enumerate}

\item
Trouver les candidats extrema en résolvant $\nabla f=(0,0)$,
\item
écrire $d^2f(a)$ pour chacun des candidats
\item
calculer les valeurs propres de $d^2f(a)$, déterminer si la matrice est définie positive ou négative,
\item
conclure.

\end{enumerate}


%+++++++++++++++++++++++++++++++++++++++++++++++++++++++++++++++++++++++++++++++++++++++++++++++++++++++++++++++++++++++++++
					\section{Limites à plusieurs variables}
%+++++++++++++++++++++++++++++++++++++++++++++++++++++++++++++++++++++++++++++++++++++++++++++++++++++++++++++++++++++++++++
		\label{SecLimVarsPlus}

Prenons une fonction $f\colon \eR^n\to \eR$. Nous disons que
\begin{equation}
	\lim_{x\to x_0}f(x)=l\in\eR
\end{equation}
lorsque $\forall \epsilon>0$, $\exists\delta$ tel que $\| x-x_0 \|\leq\delta$ implique $| f(x)-l |\leq \epsilon$. 

Remarquez qu'ici, $x\in\eR^n$, et sachez distinguer $\| . \|$, la norme dans $\eR^n$ de $| . |$ qui est la valeur absolue dans $\eR$. Une autre façon d'exprimer cette définition est que l'ensemble des valeurs atteintes par $f$ dans une boule de rayon $\delta$ autour de $x_0$ n'est pas très loin de $l$. Nous définissons donc
\begin{equation}
	E_{\delta}=\{ f(x)\tq x\in B(x_0,\delta) \}.
\end{equation}
Notez que si $f$ n'est pas définie en $x_0$, il n'y a pas de valeurs correspondantes au centre de la boule dans $E_{\delta}$. Ceci est évidement la situation générique lorsqu'il y a une indétermination à lever dans le calcul de la limite. Nous avons alors que
\begin{equation}
	\lim_{x\to x_0}f(x)=l
\end{equation}
lorsque $\forall\epsilon>0$, $\exists\delta$ tel que 
\begin{equation}		\label{Eqvmoinsrapplimdeux}
	\sup\{ | v-l |\tq v\in E_{\delta} \}\leq\epsilon.
\end{equation}
Une façon classique de montrer qu'une limite n'existe pas, est de prouver que, pour tout $\delta$, l'ensemble $E_{\delta}$ contient deux valeurs constantes. Si par exemple $0\in E_{\delta}$ et $1\in E_{\delta}$ pour tout $\delta$, alors aucune valeur de $l$ (même pas $l=\pm\infty$) ne peut satisfaire à la condition \eqref{Eqvmoinsrapplimdeux} pour toute valeur de $\epsilon$.

Nous laissons à la sagacité de l'étudiant le soin d'adapter tout ceci pour le cas $\lim_{x\to x_0}f(x)=\pm\infty$.

La proposition suivante semble évidente, mais nous sera tellement
utile qu'il est préférable de l'expliciter~:
\begin{proposition}
Soit $f : D \to \eR$ une fonction dont le domaine
  s'écrit comme une réunion \emph{finie}
  \begin{equation*}
    D = \bigcup_{i=1}^k A_i
  \end{equation*}  
  où $k$ est un entier. Soit $a \in \adh D$ tel que $a \in \adh A_i$
  pour tout $i \leq k$, et soit $b \in \eR$. Alors, la limite
  \begin{equation*}
    \limite x a f(x)
  \end{equation*}
  existe et vaut $b$ si et seulement si chacune des limites
  \begin{equation*}
    \limite[x \in A_i] x a f(x)
  \end{equation*}
  existe et vaut $b$.
\end{proposition}

\begin{proof}On sait déjà que si la limite de $f : D \to \eR$
  existe, alors toute restriction à $A_i$ admet la même limite. Il
  suffit donc de prouver la réciproque.

  Par hypothèse, pour tout $i = 1 \ldots k$, nous savons que
  \begin{equation*}
    \forall \epsilon > 0\, \exists \delta_i > 0 \tq (x \in A_i)
    \text{ et }
    (\norme{x-a} < \delta_i) \Rightarrow \norme{f(x) - b} < \epsilon
  \end{equation*}

  Si $\epsilon$ est fixé, posons $\delta = \min_i\{\delta_i\}$. Nous
  savons alors que
  \begin{enumerate}
  \item pour tout $x \in D$, il existe $i$ tel que $x \in A_i$, et
  \item si $x$ vérifie $\norme{x-a} < \delta$, alors pour tout $i$,
    $\norme{x-a} < \delta_i$ par définition de $\delta$.
  \end{enumerate}
  
  On en déduit que si $x \in D$ et $\norme{x-a} < \delta$, alors il
  existe $i$ tel que $x \in A_i$ et $\norme{x-a} < \delta_i$, ce qui
  implique $\abs{f(x) - b} < \epsilon$ et prouve la continuité.
\end{proof}

\begin{example}
  \begin{enumerate}
  \item Pour qu'une fonction $f : \eR \to \eR$ admette une limite en
    $a \in \eR$, il faut et il suffit qu'elle y admette une limite à
    droite et une limite à gauche qui soient égales.

  \item Une suite $(x_k)$ admet une limite si et seulement si les
    sous suites $(x_{2k})$ et $(x_{2k+1})$ convergent vers la même
    limite. Ceci n'est pas une application directe de la proposition,
    mais la teneur est la même.
  \end{enumerate}
\end{example}
%+++++++++++++++++++++++++++++++++++++++++++++++++++++++++++++++++++++++++++++++++++++++++++++++++++++++++++++++++++++++++++
					\section{Différentiabilité}
%+++++++++++++++++++++++++++++++++++++++++++++++++++++++++++++++++++++++++++++++++++++++++++++++++++++++++++++++++++++++++++


%---------------------------------------------------------------------------------------------------------------------------
					\subsection{Le pourquoi et le comment de la dérivée}
%---------------------------------------------------------------------------------------------------------------------------

La notion de dérivée est associée à la recherche de la droite tangente à une courbe. Reprenons rapidement le cheminement. La dérivée de $f\colon \eR\to \eR$ au point $a$ est un nombre $f'(a)$, qui définit donc une application linéaire dont le coefficients angulaire est $f'(a)$, et que nous notons $df_a$ :
\begin{equation}
	\begin{aligned}
		df_a\colon \eR&\to \eR \\
		u&\mapsto f'(a)u. 
	\end{aligned}
\end{equation}
La droite donnée par l'équation
\begin{equation}
	y(a+u)=f'(a)u
\end{equation}
est parallèle à la tangente en $a$. Pour trouver la tangente, il suffit de la décaler de la hauteur qu'il faut. L'équation de la droite tangente au graphe de $f$ au point $\big( a,f(a) \big)$ devient
\begin{equation}		\label{EqDiffRapTgDer}
	y(x)=f(a)+f'(a)(x-a)=f(a)+df_a(x-a).
\end{equation}
Nous nous proposons de généraliser cette formule au cas de la recherche du plan tangent à une surface.
 
%---------------------------------------------------------------------------------------------------------------------------
					\subsection{Dérivée partielle et directionnelles}
%---------------------------------------------------------------------------------------------------------------------------

Soit une fonction $f:A\subset \mathbb{R}^n \rightarrow \mathbb{R}^m$. Si $n\neq 1$, la notion de \emph{dérivée} de la fonction $f$ n'a plus de sens puisqu'on ne peut plus parler de pente de \emph{la} tangente au graphe de $f$ en un point. On introduit alors quelque notions qui feront, en dimension quelconque, le même travail que la dérivée en dimension un : les dérivées directionnelles et la différentielle. Nous allons voir qu'en dimension un, la différentielle coïncide avec la dérivée.


\begin{definition} 
	Soit un point $a \in int\,A$ et un vecteur $u \in \mathbb{R}^n$ avec $\| u \| =1$. La dérivée de $f$ au point $a$ dans la direction $u$ est donnée par la limite suivante, si elle existe 
	\begin{equation}
		\frac{\partial f}{\partial u}(a) = \lim_{t\rightarrow 0}\frac{f(a+tu) - f(a)}{t}
	\end{equation}
\end{definition}

Géométriquement, il s'agit du taux de variation instantané de $f$ en $a$ dans la direction du vecteur $u$, c'est-à-dire de la pente de la tangente dans la direction du vecteur $u$ au graphe de $f$ au point $(a, f(a))$.

\begin{remark}
On peut reformuler la définition en écrivant $x = a + u$, on obtient~:
\begin{equation}
	\limite[u\neq 0]{u}{0} \frac{f(a+u)-f(a)-T(u)}{\norme{u}} = 0.
\end{equation}
\end{remark}

\begin{remark}
Pourquoi avons-nous posé la condition $\| u \|=1$ ? Le but de la dérivée directionnelle dans la direction $u$ est de savoir à quelle vitesse la fonction monte lorsque l'on se déplace en suivant la direction $u$. Cette information n'aura un caractère \og objectif\fg{} que si l'on avance à une vitesse donnée. En effet, si on se déplace deux fois plus vite, la fonction montera deux fois plus vite. Par convention, nous demandons donc d'avancer à vitesse $1$.
\end{remark}

\subsubsection*{Cas particulier où $n=2$:} $a = (a_1, a_2)$, $u =
(u_1,u_2)$ et
$$\frac{\partial f}{\partial u}(a_1, a_2) = \lim_{t\rightarrow
0}\frac{f(a_1+tu_1,a_2+tu_2) - f(a_1, a_2)}{t}$$

Un cas particulier des dérivées directionnelles est la dérivée partielle. Si nous considérons la base canonique $e_i$ de $\eR^n$, nous notons
\begin{equation}
	\frac{ \partial f }{ \partial x_i }=\frac{ \partial f }{ \partial e_i }.
\end{equation}
Dans le cas d'une fonction à deux variables, nous avons donc les deux dérivées partielles
\begin{equation}
	\begin{aligned}[]
		\frac{ \partial f }{ \partial x }(a)&&\text{et}&&\frac{ \partial f }{ \partial y }(a)
	\end{aligned}
\end{equation}
qui correspondent aux dérivées directionnelles dans les directions des axes. Ces deux nombres représentent de combien la fonction $f$ monte lorsqu'on part de $a$ en se déplaçant dans le sens des axes $X$ et $Y$.

%///////////////////////////////////////////////////////////////////////////////////////////////////////////////////////////
					\subsubsection{Quelque propriétés et notations}
%///////////////////////////////////////////////////////////////////////////////////////////////////////////////////////////

\begin{enumerate}
\item
 $\forall \alpha \in \mathbb{R}$,
si $v = \alpha\,u$, alors $\frac{\partial f}{\partial v}(a) =
\alpha\,\frac{\partial f}{\partial u}(a)$.
\item Si on prend $u=e_j$ le $j$ème vecteur de la base canonique de
$\mathbb{R}^n$, alors
$$\frac{\partial f}{\partial e_j}(a) = \frac{\partial f}{\partial
x_j}(a)$$ c'est-à-dire que la dérivée de $f$ au point $a$ dans la
direction $e_j$ est la dérivée partielle de $f$ par rapport à sa
$j$ème variable.

\item 
Une fonction peut être dérivable dans certaines directions
mais pas dans d'autres (rappelez vous que si la limite à droite est
différente de la limite à gauche, la limite n'existe pas). 

\item
Même si une fonction est dérivable en un point dans toutes les
directions, on n'est pas sûr qu'elle soit continue en ce point. La
dérivabilité directionnelle n'est donc pas une notion suffisante
pour assurer la continuité. C'est pourquoi on introduit le concept
de \emph{différentiabilité}. 
\end{enumerate}

%---------------------------------------------------------------------------------------------------------------------------
					\subsection{Différentielle}
%---------------------------------------------------------------------------------------------------------------------------

\begin{definition}		\label{DefDifferentiablFnRn}
Soit un point $a \in int\,A$. La fonction $f$ est \defe{différentiable}{différentiable} au point $a$ si il existe une application linéaire $df_a\colon \eR^n\to \eR^m$ telle que 
\begin{equation}		\label{EqDefDiffableT}
	\lim_{x\to a} \frac{f(x) - f(a) - df_a (x-a)}{\|x-a\|}=0.
\end{equation}
\end{definition}

Si $f$ est différentiable en $a$, l'application $df_a$ est appelée la différentielle de $f$ en $a$. Voyons comment cette application linéaire agit sur les vecteurs de $\mathbb{R}^n$.

Le théorème suivant reprend pas principales propriétés d'une fonction différentiable.
\begin{theorem}		\label{ThoRapPropDiffSi}
Si $f$ est différentiable en $a\in\eR^n$, alors
\begin{enumerate}
\item $f$ est continue en $a$.

\item  Toute les dérivées directionnelles $\partial_uf(a)$ existent et nous avons l'égalité
\begin{equation}		\label{EqDiffPartRap}
	\begin{aligned}
		df_a\colon \eR^n&\to \eR^m \\
		u&\mapsto df_a(u)=\frac{ \partial f }{ \partial u }(a)=\sum_i \frac{ \partial f }{ \partial x_i }u^i,
	\end{aligned}
\end{equation}
si les $u^i$ sont les composantes de $u$ dans la base canonique de $\eR^n$.

La différentielle de $f$ en $a$ envoie donc un vecteur $u$ sur la dérivée directionnelle de $f$ au point $a$ dans la direction $u$. 

\item\label{ItemThoDiffSiLin} L'application $df_a$ est une application linéaire.
\end{enumerate}
\end{theorem}
Le point \ref{ItemThoDiffSiLin} est évidement contenu dans la définition de la différentielle, mais c'est bien de la remettre en toute lettre. En regard avec la formule \eqref{EqDiffPartRap}, elle dit que $\partial_uf(a)$ est linéaire par rapport à $u$.

\subsubsection*{Cas particuliers} \begin{description} \item $n=1$:
$f: \mathbb{R}\rightarrow \mathbb{R}$ est dérivable en $a$ si et
seulement si $f$ est différentiable en $a$ et
$$df_a:\mathbb{R}\rightarrow \mathbb{R}: x \mapsto df_a(x) =
f'(a)\,.\,x$$ \item $n=2$: $f$ est différentiable en $a =(a_1, a_2)$
si et seulement si
$$\lim_{(v_1,v_2)\rightarrow (0,0)} \frac{f(a_1+v_1, a_2+v_2) - f(a_1,a_2) - [ \frac{\partial f}{\partial x}(a)\,v_1+
\frac{\partial f}{\partial y}(a)\,v_2]}{\sqrt{v_1^2+v_2^2}} = 0
$$\end{description}\vspace{0.3cm}


Parmi les vecteurs $u \in \mathbb{R}^n$, un vecteur d'origine $(a,
f(a))$ se distingue des autres: le vecteur gradient de $f$ en $a$
donnant la direction de plus grande pente de $f$ en
$a$.\vspace{0.3cm}


\begin{definition}
La courbe de niveau de $f$ associée à a est donnée par
$$ S_a = f^{-1}\,(f(a)) = \{(x_1, \ldots, x_n)\in \mathbb{R}^n : f(x_1, \ldots,
x_n)=f(a) \}$$
\end{definition}

Nous avons maintenant en main les concepts utiles pour trouver l'équation du plan tangent à une surface.

%---------------------------------------------------------------------------------------------------------------------------
					\subsection{Gradient et recherche du plan tangent}
%---------------------------------------------------------------------------------------------------------------------------

De la même manière que la tangente à une courbe était la droite de coefficient angulaire donné par la dérivée, maintenant, le plan tangent à une surface est le plan dont les vecteurs directeurs sont les dérivées partielles :

La généralisation de l'équation \eqref{EqDiffRapTgDer} est 
\begin{equation}		\label{EqDefPlanTag}
	T_a(x)=f(a)+\sum_i\frac{ \partial f }{ \partial x_i }(a)(x-a)^i
\end{equation}

Nous introduisons aussi souvent l'opérateur différentiel abstrait \defe{nabla}{nabla}, noté $\nabla$ et qui est donné par le vecteur
\begin{equation}
	\nabla=\left( \frac{ \partial  }{ \partial x_1 },\ldots,\frac{ \partial  }{ \partial x_n } \right).
\end{equation}
Les égalités suivantes sont juste des notations, sommes toutes logiques, liées à $\nabla$ :
\begin{equation}
	\nabla f=\left( \frac{ \partial f }{ \partial x_1 },\ldots,\frac{ \partial f }{ \partial x_n } \right),
\end{equation}
et
\begin{equation}		\label{EqDefGradient}
	\nabla f(a) = \left(\frac{\partial f}{\partial x_1}(a), \frac{\partial f}{\partial x_2}(a), \ldots, \frac{\partial f}{\partial x_n}(a)\right).
\end{equation}
Ce dernier est un élément de $\eR^n$ : chaque entrée est un nombre réel.

\begin{definition} 
Le vecteur gradient de $f$ au point $a$ est le vecteur donné par la formule \eqref{EqDefGradient}.
\end{definition}
La notation $\nabla$ permet d'écrire la différentielle sous forme un peu plus compacte. En effet, la formule \eqref{EqDiffPartRap} peut être notée
\begin{equation}
	df_a(u)=\scal{\nabla f(a)}{u}.
\end{equation}

En utilisant ce produit scalaire, l'équation \eqref{EqDefPlanTag} peut se récrire
\begin{equation}
	T_a(x)=f(a)+\sum_i\frac{ \partial f }{ \partial x_i }(a)(x-a)^i=f(a)+\scal{\nabla f(a)}{x-a}.
\end{equation}

Affin d'éviter les confusions, il est parfois souhaitable de bien mettre les parenthèses et noter $(\nabla f)(a)$ au lieu de $\nabla f(a)$.

\begin{proposition}
$\nabla f(a)\,\bot \,S_a$
\end{proposition}


\begin{equation}		\label{EqPlanTgSansNabla}
	z=f(a)+\sum_i\frac{ \partial f }{ \partial f }(a)(x-a)^i.
\end{equation}

\subsubsection*{Cas particulier où $n=2$:} 
Le plan $T_a$ avec $a=(a_1,a_2)$ a pour équation dans $\eR^3$:
\begin{equation}		\label{EqPlanTgEnDimDeux}
	z = f(a_1,a_2) + \frac{\partial f}{\partial x}(a_1,a_2)\,(x-a_1)+ \frac{\partial f}{\partial y}(a_1,a_2)\,(y-a_2).
\end{equation}




%---------------------------------------------------------------------------------------------------------------------------
					\subsection{Différentielle comme élément de l'espace dual}
%---------------------------------------------------------------------------------------------------------------------------


Si nous considérons la base canonique $\{ e_i \}_{i=1,\ldots,n}$ de $\eR^n$. À partir d'elle, nous considérons la \defe{base duale}{Base duale}. En termes pratiques, nous définissons $dx_i$ comme la forme sur $\eR^n$ qui à un vecteur $u$ fait correspondre sa composante $i$ :
\begin{equation}
	dx_i\begin{pmatrix}
	u^1	\\ 
	\vdots	\\ 
	u^n	
\end{pmatrix}=u^i.
\end{equation}
En termes savants, $dx_i$ est le dual de $e_i$. Si tu ne l'as pas encore compris, Jean Doyen va te le faire comprendre !


Maintenant, dans la formule \eqref{EqDiffPartRap}, nous pouvons remplacer $u^i$ par $dx_i(u)$, et écrire
\begin{equation}
	df_a(u)=\sum_i\frac{ \partial f }{ \partial x_i }(a)u^i=\sum_i\frac{ \partial f }{ \partial x_i }(a)dx_i(u).
\end{equation}
Ce qui arrive tout à droite est explicitement vu comme une forme sur $\eR$, dont les composantes dans la base duale sont les dérivées partielles de $f$ au point $a$, agissant sur $u$. En faisant un pas en arrière, nous omettons le $u$, et nous écrivons
\begin{equation}
	df_a=\sum_{i=1}^n\frac{ \partial f }{ \partial x_i }(a)dx^i
\end{equation}

Cette notation $dx_i$ pour la forme duale de $e_i$ est en réalité parfaitement logique parce que $dx^i$ est la différentielle de la projection
\begin{equation}
	\begin{aligned}
		x^i\colon \eR^n&\to \eR \\
		(x^1,\ldots,x^n)&\mapsto x^i. 
	\end{aligned}
\end{equation}
Je te laisse un peu méditer sur cette différentielle de la projection. L'important est que tu aies compris cela d'ici la fin de ta deuxième année.


%---------------------------------------------------------------------------------------------------------------------------
					\subsection{Prouver qu'un fonction n'est pas différentiable}
%---------------------------------------------------------------------------------------------------------------------------

Chacun des point du théorème \ref{ThoRapPropDiffSi} est en soi un critère pour montrer qu'une fonction n'est pas différentiable en un point.

%///////////////////////////////////////////////////////////////////////////////////////////////////////////////////////////
					\subsubsection{Continuité}
%///////////////////////////////////////////////////////////////////////////////////////////////////////////////////////////


Le premier critère à vérifier est donc la continuité. Si une fonction n'est pas continue en un point, alors elle n'y sera pas différentiable. Pour rappel, la continuité en $a$ se teste en vérifiant si $\lim_{x\to a}f(x)=f(a)$.

%///////////////////////////////////////////////////////////////////////////////////////////////////////////////////////////
					\subsubsection{Linéarité}
%///////////////////////////////////////////////////////////////////////////////////////////////////////////////////////////

Un second test est la linéarité de la dérivée directionnelle par rapport à la direction : l'application $u\mapsto\frac{ \partial f }{ \partial u }(a)$ doit être linéaire, sinon $df_a$ n'existe pas.
\begin{example}		\label{Exemple0046Diff}
Examinons la fonction
\begin{equation}
	\begin{aligned}
		f\colon \eR^2&\to \eR \\
		(x,y)&\mapsto \begin{cases}
	\frac{ xy^2 }{ x^2+y^4 }	&	\text{si $(x,y)\neq (0,0)$}\\
	0	&	 \text{sinon}.
\end{cases}
	\end{aligned}
\end{equation}
Prenons $u=(u_1,u_2)$ et calculons la dérivée de $f$ dans la direction de $u$ au point~$(0,0)$ :
\begin{equation}
	\begin{aligned}[]
		\frac{ \partial f }{ \partial u }(0,0)	
			&=\lim_{t\to 0}\frac{ f(tu_1,tu_2)-f(0,0) }{ t }\\
			&=\lim_{t\to 0}\frac{1}{ t }\left( \frac{ tu_1t^2u_2 }{ t^2u_1^2+t^4u_2^4 } \right)\\
			&=\lim_{t\to 0}\left( \frac{ u_1u_2^2 }{ u_1^2+t^2u_2^4 } \right)\\
			&=\begin{cases}
	\frac{ u_2^2 }{ u_1 }	&	\text{si $u_1\neq 0$}\\
	0	&	 \text{si $u_1=0$}.
\end{cases}
	\end{aligned}
\end{equation}
Cette application n'est pas linéaire par rapport à $u$. En effet, notons
\begin{equation}
	\begin{aligned}
		A\colon \eR^n&\to \eR \\
		u&\mapsto \frac{ \partial f }{ \partial u }(0,0), 
	\end{aligned}
\end{equation}
et vérifions que pour tout $u$ et $v$ dans $\eR^n$ et $\lambda\in\eR$, nous ayons $A(\lambda u)=\lambda A(u)$ et $A(u+v)=A(u)+A(v)$. Le premier fonctionne parce que
\begin{equation}
	A(\lambda u)=A(\lambda u_1,\lambda u_2)=\frac{ \lambda^2 u_2^2 }{ \lambda u_1 }=\lambda\frac{ u_2^2 }{ u_1 }=\lambda A(u).
\end{equation}
Mais nous avons par exemple
\begin{equation}
	A\big( (0,1)+(2,3) \big)=A(2,4)=\frac{ 16 }{ 2 }=8,
\end{equation}
tandis que
\begin{equation}
	A(0,1)+A(2,3)=0+\frac{ 9 }{ 2 }\neq 8.
\end{equation}
La fonction $f$ n'est donc pas différentiable en $(0,0)$, parce que la candidate différentielle, $df_{(0,0)}(u)=\frac{ \partial f }{ \partial u }(0,0)$, n'est même pas linéaire.

Nous verrons, dans l'exercice \ref{exo0046}, une autre manière de traiter cette fonction.
\end{example}

%///////////////////////////////////////////////////////////////////////////////////////////////////////////////////////////
					\subsubsection{Cohérence des dérivées partielles et directionnelle}
%///////////////////////////////////////////////////////////////////////////////////////////////////////////////////////////

Dans la pratique, nous pouvons calculer $\partial_uf(a)$ pour une direction $u$ générale, et puis en déduire $\partial_xf$ et $\partial_yf$ comme cas particuliers en posant $u=(1,0)$ et $u=(0,1)$. Une chose incroyable, mais pourtant possible est qu'il peut arriver que
\begin{equation}
	\frac{ \partial f }{ \partial u }(a)\neq \sum_i\frac{ \partial f }{ \partial x_i }(a)u^i.
\end{equation}
Ceci se produit lorsque $f$ n'est pas différentiable en $a$. En voici un exemple.

%///////////////////////////////////////////////////////////////////////////////////////////////////////////////////////////
					\subsubsection{Un candidat dans la définition (marche toujours)}
%///////////////////////////////////////////////////////////////////////////////////////////////////////////////////////////

Lorsqu'une fonction est donné, un candidat différentielle au point $(a_1,a_2)$ est souvent assez simple à trouver en un point :
\begin{equation}
	T(u_1,u_2)=\frac{ \partial f }{ \partial x }(a_1,a_2)u_1+\frac{ \partial f }{ \partial y }(a_1,a_2)u_2.
\end{equation}
L'application $T$ est la candidate différentielle en ce sens que si la différentielle existe, alors elle est égale à $T$. Ensuite, il faut vérifier si
\begin{equation}		\label{EqLimDefDiff}
	\lim_{(x,y)\to (a_1,a_2)} \frac{f(x,y) - f(a_1,a_2) - T\big( (x,y)-(a_1,a_2) \big)}{\| (x,y)-(a_1,a_2) \|}=0
\end{equation}
ou non. Si oui, alors la différentielle existe et $df_{(a,b)}(u)=T(u)$, sinon\footnote{y compris si la limite \eqref{EqLimDefDiff} n'existe même pas.}, la différentielle n'existe pas.

Attention : dans la ZAP, les dérivées partielles $\partial_xf$ et $\partial_yf$ ne peuvent en général pas être calculées en utilisant les règles de calcul (c'est bien pour ça que la ZAP est une zone à problèmes). Il faut d'office utiliser la définition
\begin{equation}
	\frac{ \partial f }{ \partial x }(a_1,a_2)=\lim_{t\to 0}\frac{ f(a_1+t,a_2)-f(a_1,a_2) }{ t },
\end{equation}
et la définition correspondante pour $\partial_yf$.


\subsubsection*{Conclusion}
Soient $f:A\subset \eR^n \rightarrow \eR^m$, et $a\in int\,A$. Si $f$ est différentiable en $a$, $$ (df_a (e_j))_i = d(f_i)_a(e_j) =\frac{\partial f_i}{\partial x_j}(a)= [Jac(f)_{|a}]_{ij}$$ et la matrice de l'application linéaire $df_a$ est la matrice jacobienne $m\times n$ de $f$ en $a$ notée $Jac(f)_{|a}$.


%---------------------------------------------------------------------------------------------------------------------------
					\subsection{Calcul de différentielles}
%---------------------------------------------------------------------------------------------------------------------------


\begin{remark}		\label{deriveepartielles}
  En pratique, ayant une formule pour la fonction $f$, on dérive --grâce aux règles usuelles de dérivation-- par rapport à la variable $x_i$ en considérant que les autres ($x_j$ avec $j \neq i$) sont des constantes.
\end{remark}

\begin{example}Pour $f(x,y) = xy + x^2$, les dérivées partielles
  s'écrivent
  \begin{equation*}
    \frac{\partial f}{\partial x} = y + 2x \quad\text{et}\quad \frac{\partial f}{\partial y} = x
  \end{equation*}
\end{example}


Des \emph{règles de calcul} sont d'application. En particulier, quand
ces opérations existent, les sommes, différences, produits, quotients
et compositions d'applications différentiables sont différentiables.

Toute application linéaire est différentiable, et sa différentielle en
tout point est égale à l'application elle-même. En particulier, les
\Defn{projections canoniques}, c'est-à-dire les applications du type
$(x,y,z) \mapsto y$, sont linéaires donc différentiables.

\begin{example}
Les cas suivants sont faciles :
  \begin{enumerate}
  \item En combinant les projections canoniques avec les règles de
    calculs, on obtient que toute fonction polynômiale à $n$ variables
    est différentiable comme application de $\eR^n$ dans $\eR$.

  \item Toute fonction rationnelle, du type $f(x) \pardef
    \frac{P(x)}{Q(x)}$ où $P$ et $Q$ sont des polynômes, est
    différentiable en tout point $a$ tel que $Q(a) \neq 0$.

  \item Pour une fonction d'une variable $f : D \subset \eR \to
    \eR$, le caractère différentiable et le caractère dérivable
    coïncident. De plus, on a
    \begin{equation*}
      d f_a(u) = f'(a) u.
    \end{equation*}
  \end{enumerate}
\end{example}

%---------------------------------------------------------------------------------------------------------------------------
					\subsection{Notes idéologiques sur le concept de plan tangent}
%---------------------------------------------------------------------------------------------------------------------------
\label{ssecConceptPlanTag}

La page 94 du syllabus dit des choses très intéressantes que tu n'auras sans doute pas comprises \ldots et que tu ne comprendra sans doute pas complètement avant d'avoir un vrai cours de géométrie différentielle. Notons $G$, le graphe d'une fonction $f$, c'est à dire
\begin{equation}
	G=\{ (x,y,z)\in\eR^3\tq z=f(x,y) \}.
\end{equation}
Première affirmation : si $\gamma\colon \eR\to G$ est une courbe telle que $\gamma(0)=\big( a,f(a) \big)$, alors $\gamma'(0)\in\eR^n$ est dans le plan tangent à $G$ au point $\big( a,f(a) \big)$.

Plus fort : tous les éléments du plan tangent sont de cette forme.

Le plan tangent à $G$ en un point $x\in G$ est donc constitué des vecteurs vitesse de tous les chemins qui passent par $x$.

Prenons maintenant $S$, une courbe de niveau de $G$, c'est à dire
\begin{equation}
	S=\{ (x,y)\in\eR^2\tq f(x,y)=C \}.
\end{equation}
Si nous prenons un chemin dans $G$ qui est, de plus, contraint à $S$, c'est à dire tel que $\gamma(t)\in S$, alors $\gamma'(0)$ sera tangent à $G$ (ça, on le savait déjà), mais en plus, $\gamma'(0)$ sera tangent à $S$, ce qui est logique.

La morale est que si vous prenez un chemin qui se ballade dans n'importe quoi, alors la dérivée du chemin sera un vecteur tangent à ce n'importe quoi.

En outre, si $\gamma(t)\in S$ et $\gamma(0)=a$, alors
\begin{equation}
	\scal{\nabla f(a)}{\gamma'(0)}=0,
\end{equation}
c'est à dire que le vecteur tangent à la courbe de niveau est perpendiculaire au gradient. Cela est intuitivement logique parce que la tangente à la courbe de niveau correspond à la direction de \emph{moins} grande pente.

%+++++++++++++++++++++++++++++++++++++++++++++++++++++++++++++++++++++++++++++++++++++++++++++++++++++++++++++++++++++++++++
					\section{Jacobienne et calcul de différentielles}
%+++++++++++++++++++++++++++++++++++++++++++++++++++++++++++++++++++++++++++++++++++++++++++++++++++++++++++++++++++++++++++

\subsection{Rappels et définitions}

Dans cette section nous considérons des fonctions $f : D \to \eR^m$
où $D \subset \eR^n$, et un point $a \in \interieur D$ où $f$ est
différentiable.
\begin{remark}
  La définition de continuité (resp. différentiabilité) pour une
  fonction à valeurs vectorielles est celle introduite précédemment,
  et on remarque que pour avoir la continuité
  (resp. différentiabilité) de $f$ en un point, il faut et il suffit
  de chacune des composantes de $f = (f_1,\ldots, f_m)$, vues
  séparément comme fonctions à $n$ variables et à valeurs réelles,
  soit continue (resp. différentiable) en ce point.
\end{remark}

\begin{defn}La \Defn{jacobienne} de $f$ en $a$ est la matrice
  \begin{equation*}
    (\Jac f)_a \begin{pmatrix}
      \pder {f_1} {x_1}(a) & \ldots &\pder {f_1} {x_n}(a)\\
      \vdots& & \vdots\\
      \pder {f_m} {x_1}(a) & \ldots &\pder {f_m} {x_n}(a)
    \end{pmatrix}
  \end{equation*}
  composée de l'ensemble des dérivées partielles de $f$.

  Si $m = 1$, cette matrice ne contient qu'une ligne ; c'est donc un
  vecteur appelé le \Defn{gradient de $f$ en $a$} et noté $\nabla f(a)$.
\end{defn}

\begin{remark}
  \begin{enumerate}
  \item Si la fonction est supposée différentiable, calculer la
    jacobienne revient à connaître la différentielle. En effet, par
    linéarité de la différentielle et par définition des dérivées
    partielles, nous avons
    \begin{equation*}
      d f_a (u) =%
      \begin{pmatrix}
        \pder {f_1} {x_1}(a) & \ldots &\pder {f_1} {x_n}(a)\\
        \vdots& & \vdots\\
        \pder {f_m} {x_1}(a) & \ldots &\pder {f_m} {x_n}(a)
      \end{pmatrix}
      \begin{pmatrix}u_1\\\vdots\\u_n\end{pmatrix}
    \end{equation*}
    où $u = (u_1, \ldots, u_n)$ et où le membre de droite est un
    produit matriciel

  \item Remarquons que la jacobienne peut exister en un point donné
    sans que la fonction soit différentiable en ce point !
  \end{enumerate}
\end{remark}

\begin{proposition}[Règles de calculs] Soient $f$ et $g$ des fonctions
  différentiables en $g(a)$ et $a$ respectivement, alors la composée
  $f\circ g$ est différentiable en $a$ et
  \begin{equation*}
    d (f\circ g)_a = d f_{g(a)} \circ d g_a
  \end{equation*}
  et de plus les jacobiennes correspondantes vérifient
  \begin{equation*}
    \Jac (f\circ g)_a = \Jac f_{g(a)} \Jac g_a
  \end{equation*}
  où le membre de droite est le produit (non-commutatif !) des deux matrices.
\end{proposition}

\begin{corollary}[Chain rule] Si $f : \eR^p \to \eR$ et $g : \eR \to
  \eR^p$, alors
  \begin{equation*}
    (f\circ g)^\prime(t) = \sum_{i=1}^p \pder f {x_i}(g(t)) g_i^\prime(t).
  \end{equation*}
\end{corollary}

\begin{remark}
  \begin{enumerate}
  \item Si $p = 1$, on retrouve la règle usuelle de dérivation de
    fonctions composées.

  \item 
	  Si $g$ est à plusieurs variables, cette règle permet de déterminer les dérivées partielles de $f \circ g$, puisqu'une dérivée partielle peut être vue comme dérivée usuelle par rapport à une seule variable (voir remarque page \pageref{deriveepartielles}).

  \item Si $f$ est à valeurs vectorielles, cette formule permet de
    retrouver la jacobienne de $f \circ g$ puisqu'il suffit de traiter
    chaque composante de $f$ séparément.
  \end{enumerate}
\end{remark}

\begin{defn}
  Soit $f : \eR^n \to\eR$ une fonction différentiable en un point
  $a$. Le \emph{plan tangent} au graphe de $f$ en $(a,f(a))$ est
  l'ensemble des points
  \begin{equation*}
    \begin{split}
      T_af &= \{ (x,z) \in \eR^n \times \eR \tq z = f(a) + d f_a (x-a)\}\\
      &= \{ (x,z) \in \eR^n \times \eR \tq z = f(a) + \scalprod{\nabla f(a)}{x-a}\}
    \end{split}
  \end{equation*}
\end{defn}

%+++++++++++++++++++++++++++++++++++++++++++++++++++++++++++++++++++++++++++++++++++++++++++++++++++++++++++++++++++++++++++
					\section{Intégrales multiples}
%+++++++++++++++++++++++++++++++++++++++++++++++++++++++++++++++++++++++++++++++++++++++++++++++++++++++++++++++++++++++++++

%%%%%%%%%%%%%%%%%%%%%%%%%%%%%%%%%%%
%
%  NOTE : toute cette partie a été reprise dans OutilsMath le 3 avril 2011.
%       

Il est expliqué dans le cours théorique comment on définit le nombre
\begin{equation}
	\int_Ef
\end{equation}
lorsque $E\subset\eR^n$ et $f\colon \eR^n\to \eR$. Nous allons maintenant montrer comment calculer des intégrales en pratique.

%\label{PgRapIntMultFubiniRect}
% ceci était le label d'une sous section ``Intégrale sur un rectangle''.


%---------------------------------------------------------------------------------------------------------------------------
					\subsection{Intégrales sur d'autres domaines}
%---------------------------------------------------------------------------------------------------------------------------
Voir la sous section \ref{PgRapIntMultFubiniTri}.


%---------------------------------------------------------------------------------------------------------------------------
					\subsection{Changement de variables}
%---------------------------------------------------------------------------------------------------------------------------


Le domaine $E=\{ (x,y)\in\eR^2\tq x^2+y^2<1 \}$ s'écrit plus facilement $E=\{ (r,\theta)\tq r<1 \}$ en coordonnées polaires. Le passage aux coordonnées polaire permet de transformer une intégration sur un domaine rond à une intégration sur le domaine rectangulaire $\mathopen]0,2\pi\mathclose[\times\mathopen]0,1\mathclose[$. La question est évidement de savoir si nous pouvons écrire
\begin{equation}
	\int_Ef=\int_{0}^{2\pi}\int_0^1f(r\cos\theta,r\sin\theta)drd\theta.
\end{equation}
Hélas, non; la vie n'est pas aussi simple. Le théorème est celui de la page 480.

\begin{theorem}
Soit $g\colon A\to B$ un difféomorphisme. Soient $F\subset B$ un ensemble mesurable et borné et $f\colon F\to \eR$ une fonction bornée et intégrable. Supposons que $g^{-1}(F)$ soit borné et que $Jg$ soit borné sur $g^{-1}(F)$. Alors
\begin{equation}
	\int_Ff(x)dy=\int_{g^{-1}(F)f\big( g(x) \big)}| Jg(x) |dx
\end{equation}
\end{theorem}
Pour rappel, $Jg$ est le déterminant de la matrice \href{http://fr.wikipedia.org/wiki/Matrice_jacobienne}{jacobienne} (aucun lien de \href{http://fr.wikipedia.org/wiki/Jacob}{parenté}) donnée par
\begin{equation}
	Jg=\det\begin{pmatrix}
	\partial_xg_1	&	\partial_yg_1	\\ 
	\partial_xg_2	&	\partial_tg_2	
\end{pmatrix}.
\end{equation}
Un \defe{difféomorphisme}{difféomorphisme} est une application $g\colon A\to B$ telle que $g$ et $g^{-1}\colon B\to A$ soient de classe $C^1$.

%///////////////////////////////////////////////////////////////////////////////////////////////////////////////////////////
					\subsubsection{Coordonnées polaires}
%///////////////////////////////////////////////////////////////////////////////////////////////////////////////////////////

Les coordonnées polaires sont données par le difféomorphisme
\begin{equation}
	\begin{aligned}
		g\colon \mathopen]0,\infty\mathclose[\times\mathopen]0,2\pi\mathclose[ &\to\eR^2\setminus D\\
		(r,\theta)&\mapsto \big( r\cos(\theta),r\sin(\theta) \big)
	\end{aligned}
\end{equation}
où $D$ est la demi droite $y=0$, $x\geq 0$. Le fait que les coordonnées polaires ne soient pas un difféomorphisme sur tout $\eR^2$ n'est pas un problème pour l'intégration parce que le manque de difféomorphisme est de mesure nulle dans $\eR^2$. Le jacobien est donné par
\begin{equation}
	Jg=\det\begin{pmatrix}
	\partial_rx	&	\partial_{\theta}x	\\ 
	\partial_ry	&	\partial_{\theta}y
\end{pmatrix}=\det\begin{pmatrix}
	\cos(\theta)	&	-r\sin(\theta)	\\ 
	\sin(\theta)	&	r\cos(\theta)	
\end{pmatrix}=r.
\end{equation}

Voir l'exemple \ref{ExpmfDtAtV}.

%///////////////////////////////////////////////////////////////////////////////////////////////////////////////////////////
\subsubsection{Coordonnées sphériques}
%///////////////////////////////////////////////////////////////////////////////////////////////////////////////////////////

Voir le point \ref{SubSubCoordSpJxhMwm}
 
%+++++++++++++++++++++++++++++++++++++++++++++++++++++++++++++++++++++++++++++++++++++++++++++++++++++++++++++++++++++++++++
					\section{Théorème de la fonction implicite}
%+++++++++++++++++++++++++++++++++++++++++++++++++++++++++++++++++++++++++++++++++++++++++++++++++++++++++++++++++++++++++++

%---------------------------------------------------------------------------------------------------------------------------
\subsection{Mise en situation}
%---------------------------------------------------------------------------------------------------------------------------

Dans un certain nombre de situation, il n'est pas possible de trouver des solutions explicites aux équations qui apparaissent. Néanmoins, l'existence \og théorique\fg{} d'une telle solution est déjà suffisante pour certaines situations, et c'est l'objet du théorème de la fonction implicite.

Prenons par exemple la fonction sur $\eR^2$ donnée par 
\begin{equation}
	F(x,y)=x^2+y^2-1.
\end{equation}
Nous pouvons bien entendu regarder l'ensemble des points donnés par $F(x,y)=0$. C'est le cercle dessiné à la figure \ref{LabelFigCercleImplicite}.
\newcommand{\CaptionFigCercleImplicite}{Un cercle pour montrer l'intérêt de la fonction implicite.}
\input{Fig_CercleImplicite.pstricks}
Nous ne pouvons pas donner le cercle sous la forme $y=y(x)$ à cause du $\pm$ qui arrive quand on prend la racine carrée. Mais si on se donne le point $P$, nous pouvons dire que \emph{autour de $P$}, le cercle est la fonction
\begin{equation}
	y(x)=\sqrt{1-x^2}.
\end{equation}
Tandis que autour du point $P'$, le cercle est la fonction
\begin{equation}
	y(x)=-\sqrt{1-x^2}.
\end{equation}
Autour de ces deux point, donc, le cercle est donné par une fonction. Il n'est par contre pas possible de donner le cercle autour du point $Q$ sous la forme d'une fonction.

Ce que nous voulons faire, en général, est de voir si l'ensemble des points tels que
\begin{equation}
	F(x_1,\ldots,x_n,y)=0
\end{equation}
peut être donné par une fonction $y=y(x_1,\ldots,x_n)$. En d'autre termes, est-ce qu'il existe une fonction $y(x_1,\ldots,x_n)$ telle que
\begin{equation}
	F\big( x_1,\ldots,x_n,y(x_1,\ldots,x_n)\big)=0.
\end{equation}



\subsection{Définitions et rappels}
Soit
\begin{equation*}
  F : D \subset (\RR^n \times \RR^m) \to \RR^m : (x,y) \mapsto F(x,y) =
  (F_1(x,y),\ldots,F_m(x,y))
\end{equation*}
avec $x = (x_1,\ldots, x_n)$ et $y = (y_1,\ldots,y_m)$.
% Pour chaque $x$ fixé, on s'intéresse aux solutions du système de $m$
% équations $F(x,y) = 0$ pour les inconnues $y$ ; en particulier, on
% voudrait pouvoir écrire $y = \varphi(x)$ vérifiant $F(x,\varphi(x))
% = 0$.

Pour $(x,y) \in \interieur D$, la matrice
\begin{equation*}
\begin{pmatrix}
\pder {F_1}{y_1}(x,y)& \ldots& \pder {F_1}{y_m}(x,y)\\
\vdots& \ddots & \vdots\\
\pder {F_m}{y_1}(x,y)& \ldots& \pder {F_m}{y_m}(x,y)\\
\end{pmatrix}
\end{equation*}
est la \defe{matrice jacobienne}{jacobienne!matrice} de $F$ par rapport à $y$ (au point
$(x,y)$; son déterminant est appelé le \defe{jacobien}{jacobien} de F par
    rapport à y et se note
  $\pder{(F_1,\ldots,F_m)}{(y_1,\ldots,y_m)}(x,y)$.


\begin{theorem}
	Soit $(\bar x,\bar y)$ tel que $F(\bar x,\bar y) = 0$ et
  $\pder{(F_1,\ldots,F_m)}{(y_1,\ldots,y_m)}(\bar x,\bar y) \neq
  0$. Alors il existe un voisinage $U$ de $x$ dans $\RR^n$, un
  voisinage $V$ de $y$ dans $\RR^m$ et une unique application $\varphi
  : U \to V$ tels que
  \begin{enumerate}
  \item $\varphi(\bar x) = \bar y$ ; 
  \item $F(x,\varphi(x)) = 0$ pour tout $x \in U$.
  \end{enumerate}
\end{theorem}
	
Le théorème de la fonction implicite a pour objet de donner l'existence de la fonction $\varphi$. Maintenant nous pouvons dire beaucoup de choses sur les dérivées de $\varphi$ en considérant la fonction
\begin{equation}
	x\mapsto F\big( x,\varphi(x) \big).
\end{equation}
Par définition de $\varphi$, cette fonction est toujours nulle. En particulier, nous pouvons dériver l'équation
\begin{equation}
	F\big( x,\varphi(x) \big)=0,
\end{equation}
et nous trouvons plein de choses.

%---------------------------------------------------------------------------------------------------------------------------
\subsection{Exemple}
%---------------------------------------------------------------------------------------------------------------------------

Prenons par exemple la fonction
\begin{equation}
	F\big( (x,y),z \big)=ze^z-x-y,
\end{equation}
et demandons nous ce que nous pouvons dire sur la fonction $z(x,y)$ telle que
\begin{equation}
	F\big( x,y,z(x,y) \big)=0,
\end{equation}
c'est à dire telle que
\begin{equation}		\label{EqDefZImplExemple}
	z(x,y) e^{z(x,y)}-x-y=0.
\end{equation}
pour tout $x$ et $y\in\eR$. Nous pouvons facilement trouver $z(0,0)$ parce que
\begin{equation}
	z(0,0) e^{z(0,0)}=0,
\end{equation}
donc $z(0,0)=0$.

Nous pouvons dire des choses sur les dérivées de $z(x,y)$. Voyons par exemple $(\partial_xz)(x,y)$. Pour trouver cette dérivée, nous dérivons la relation \eqref{EqDefZImplExemple} par rapport à $x$. Ce que nous trouvons est
\begin{equation}
	(\partial_xz)e^z+ze^z(\partial_xz)-1=0.
\end{equation}
Cette équation peut être résolue par rapport à $\partial_xz$~:
\begin{equation}
	\frac{ \partial z }{ \partial x }(x,y)=\frac{1}{ e^z(1+z) }.
\end{equation}
Remarquez que cette équation ne donne pas tout à fait la dérivée de $z$ en fonction de $x$ et $y$, parce que $z$ apparaît dans l'expression, alors que $z$ est justement la fonction inconnue. En général, c'est la vie, nous ne pouvons pas faire mieux.

Dans certains cas, on peut aller plus loin. Par exemple, nous pouvons calculer cette dérivée au point $(x,y)=(0,0)$ parce que $z(0,0)$ est connu :
\begin{equation}
	\frac{ \partial z }{ \partial x }(0,0)=1.
\end{equation}
Cela est pratique pour calculer, par exemple, le développement en Taylor de $z$ autour de $(0,0)$.


%+++++++++++++++++++++++++++++++++++++++++++++++++++++++++++++++++++++++++++++++++++++++++++++++++++++++++++++++++++++++++++
\section{Formes différentielles et son intégrale sur un chemin}
%+++++++++++++++++++++++++++++++++++++++++++++++++++++++++++++++++++++++++++++++++++++++++++++++++++++++++++++++++++++++++++

%---------------------------------------------------------------------------------------------------------------------------
\subsection{Forme différentielle}
%---------------------------------------------------------------------------------------------------------------------------

La formule d'intégration d'un champ de vecteur,
\begin{equation}
	\int_{\gamma}G=\int_{[a,b]}\langle G (\gamma(t)), \gamma'(t)\rangle dt,
\end{equation}
contient quelque chose d'intéressant : la combinaison $\langle G( \gamma(t) ), \gamma'(t)\rangle$. Cette combinaison sert à transformer le vecteur tangent $\gamma'(t)$ en un nombre en utilisant le produit scalaire avec le vecteur $G( \gamma(t) )$.

Si $G$ est un champ de vecteur sur $\eR^n$, et si $x\in\eR^n$, nous pouvons considérer, de façon un peu plus abstraite, l'application
\begin{equation}		\label{EqDefBemol}
	\begin{aligned}[]
		G^{\flat}_x\colon \eR^n&\to \eR \\
			v&\mapsto \langle G(x), v\rangle . 
	\end{aligned}
\end{equation}
Cela permet de compactifier la notation et écrire
\begin{equation}
	\int_{\gamma}G=\int_{[a,b]} G^{\flat}_{\gamma(t)}\big( \gamma'(t)\big) dt.
\end{equation}

Nous nous proposons maintenant d'étudier plus en détail ce qu'est l'objet $G^{\flat}$. La règle \eqref{EqDefBemol} dit que pour chaque $x$, l'application $G_x^{\flat}$ est une forme sur $\eR^n$, c'est à dire une application linéaire de $\eR^n$ vers $\eR$. Nous écrivons que
\begin{equation}
	G_x^{\flat}\in\big( \eR^n \big)^*.
\end{equation}
Nous connaissons la \defe{base duale}{Base duale} de $(\eR^n)^*$, ce sont les formes $e^*_i$ définies par $e^*_i(e_j)=\delta_{ij}$. Dans le cadre du cours d'analyse, nous allons noter ces formes\footnote{Parce que ce sont les différentielles des fonctions (projections)
\begin{equation}
	\begin{aligned}
			x_i\colon \eR^n&\to \eR \\
			x&\mapsto x_i 
		\end{aligned}
	\end{equation}
}
par $dx_i$ :
\begin{equation}
	\begin{aligned}[]
		e^*_1&=dx_1\colon v\mapsto v_1	\\
			&\vdots			\\
		e^*_n&=dx_n\colon v\mapsto v_n
	\end{aligned}
\end{equation}
Étant donné que ces $dx_i$ forment une base de l'espace vectoriel $(\eR^n)^*$, toute application linéaire $L\colon \eR^n\to \eR$ s'écrit
\begin{equation}
	\begin{aligned}[]
		Lv&=a_1v_1+\ldots+a_nv_n\\
			&=a_1dx_1(v)+\ldots+a_ndx_n(v).
	\end{aligned}
\end{equation}
Plus abstraitement, nous notons
\begin{equation}
	\begin{aligned}[]
		L&=a_1dx_1+\ldots+a_ndx_n\\
		&=\sum_{i=1}^na_idx_i.
	\end{aligned}
\end{equation}
L'application $L$ est une combinaison linéaire des $dx_i$ au sens de l'espace vectoriel $(\eR^n)^*$.

L'objet $G^{\flat}$ est la donnée, en chaque point de $D$, d'une telle forme sur $\eR^n$. Nous donnons alors la définition suivante.
\begin{definition}
	Soit $D$, un domaine dans $\eR^n$. Une $1$-\defe{forme différentielle}{Forme différentielle} $\omega$ sur $D$ est une application
	\begin{equation}
		\begin{aligned}
				\omega\colon D&\to (\eR^n)^* \\
				x&\mapsto \omega_x. 
			\end{aligned}
		\end{equation}
\end{definition}
Étant donné que $\{ dx_i \}$ est une base de $(\eR^n)^*$, pour chaque $x\in D$, il existe des uniques réels $a_i(x)$ tels que
\begin{equation}
	\omega_x=a_1(x)dx_1+\ldots+a_n(x)dx_n.
\end{equation}
Nous disons qu'une $1$-forme différentielle est \defe{continue}{continue!forme différentielle} si les fonctions $a_i$ sont continues. La forme sera $C^k$ quand les $a_i$ seront $C^k$.

\begin{remark}
	L'ensemble des $1$-formes différentielles forment un espace vectoriel avec les définitions
	\begin{equation}
		\begin{aligned}[]
			(\lambda\omega)_x(v)&=\lambda\omega_x(v)\\
			(\omega+\mu)_x(v)&=\omega_x(v)+\mu_x(v).
		\end{aligned}
	\end{equation}
\end{remark}

Lorsque une $1$-forme différentielle s'écrit toujours sous la forme
\begin{equation}
	\omega=\sum_i a_idx_i
\end{equation}
pour certaines fonctions $a_i$. Évidemment, ces fonctions $a_i$ peuvent être trouvées en appliquant $\omega$ aux éléments de la base canonique de $\eR^n$ :
\begin{equation}
	a_j(x)=\omega_x(e_j)
\end{equation}
parce que $\omega_x(e_j)=\sum_ia_i(x)dx_i(e_i)=\sum_ia_i(x)\delta_{ij}=a_j(x)$.

Nous pouvons ainsi déterminer le développement de $G^{\flat}$ dans la base des $dx_i$ en faisant le calcul
\begin{equation}
	G_x^{\flat}(e_i)=\langle G(x), e_i\rangle =G_i(x),
\end{equation}
donc les composantes de $G^{\flat}$ dans la base $dx_i$ sont exactement les composantes de $G$ dans la base $e_i$ :
\begin{equation}
	G^{\flat}_x=G_1(x)dx_1+\ldots+G_n(x)dx_n.
\end{equation}


%///////////////////////////////////////////////////////////////////////////////////////////////////////////////////////////
\subsubsection{Une petite note pour titiller monsieur Jean Doyen}
%///////////////////////////////////////////////////////////////////////////////////////////////////////////////////////////

Pensons pendant quelque minutes aux fonctions de $\eR$ dans $\eR$. Monsieur Jean Doyen dit toujours que quand le sage demande la fonction $f$, le simple dit \og $f(x)$\fg. Or $f(x)$ n'est pas une fonction; c'est $f$, la fonction. Avec un tout petit peu de mauvaise foi, nous pouvons prétendre que $x$ désigne la fonction identité qui à chaque $x$ fait correspondre $x$ lui-même. Il est d'ailleurs un peu normal de désigner comme ça cette fonction. Dans ce cas, $f(x)$ désigne la fonction composée de la fonction $f$ avec la fonction $x$, et tout le monde est content.

Avouons que cela est un petit peu de mauvaise foi\footnote{J'offre un pot à qui ose écrire que $f(x)$ est bien la \emph{fonction} composée de $f$ avec $x$ sur sa feuille d'examen du cours d'algèbre linéaire.}. Vraiment ?

La fonction $x$ est une fonction de $\eR$ vers $\eR$. Sa différentielle en un point est donc une application de $\eR$ vers $\eR$. Devinez ce qu'elle vaut ? Ben oui : la différentielle de la fonction $x$ est \emph{vraiment} le $dx$ qu'on écrit tout le temps, la forme différentielle, la base de l'espace dual !

%///////////////////////////////////////////////////////////////////////////////////////////////////////////////////////////
\subsubsection{L'isomorphisme musical}
%///////////////////////////////////////////////////////////////////////////////////////////////////////////////////////////

Nous savons qu'un champ de vecteur $G$ produit la forme différentielle $G^{\flat}$. La construction inverse existe également. Si $\omega$ est une $1$-forme différentielle, nous pouvons définir le champ de vecteur $\omega^{\sharp}$ par la formule (implicite)
\begin{equation}
	\omega_x(v)=\langle \omega^{\sharp}(x), v\rangle 
\end{equation}
pour tout $v\in\eR^n$. Par définition, $(\omega^{\sharp})^{\flat}=\omega$. 

\begin{exercice}
	Prouver que, en composantes, 
	\begin{equation}
		\omega^{\sharp}(x)=\big( a_1(x),\ldots,a_n(x) \big),
	\end{equation}
	et vérifier que si $G$ est un champ de vecteurs, alors $(G^{\flat})^{\sharp}=G$.
\end{exercice}


%///////////////////////////////////////////////////////////////////////////////////////////////////////////////////////////
\subsubsection{Formes différentielles exactes et fermées}
%///////////////////////////////////////////////////////////////////////////////////////////////////////////////////////////

Considérons une fonction différentiable $f\colon D\to \eR$. Pour chaque $x\in D$, nous avons l'application différentielle
\begin{equation}
	\begin{aligned}
		df(x)\colon \eR^n&\to \eR \\
		v&\mapsto \sum_{i=1}^n\frac{ \partial f }{ \partial x_i }(x)v_i, 
	\end{aligned}
\end{equation}
c'est à dire que $df$ est une $1$-forme différentielle dont les composantes sont
\begin{equation}
	df(x)=\frac{ \partial f }{ \partial x_1 }(x)dx_1+\ldots+\frac{ \partial f }{ \partial x_n }(x)dx_n.
\end{equation}

Il est naturel de se demander si toutes les formes différentielles sont des différentielles de fonctions. Une réponse complète est délicate à établir, mais a d'innombrables conséquences en physique, notamment en ce qui concerne l'existence d'un potentiel vecteur pour le champ magnétique dans les équations de Maxwell.
\begin{definition}
	Deux classes importantes de formes différentielles sont à mettre en évidence
	\begin{enumerate}
		\item
			Une forme différentielle $\omega$ sur un ouvert $D\subset\eR^n$ est \defe{exacte}{forme différentielle!exacte} si il existe une fonction différentiable $f\colon D\to \eR$ telle que
			\begin{equation}
				 \omega_x=df(x)
			\end{equation}
			pour tout $x\in D$.
		\item
			Une $1$-forme de classe $C^1$ sur l'ouvert $D$ est \defe{fermée}{forme différentielle!fermée} si pour tout $i,j=1,\ldots n$, nous avons
			\begin{equation}
				\frac{ \partial a_i }{ \partial x_j }=\frac{ \partial a_j }{ \partial x_i }.
			\end{equation}
	\end{enumerate}
\end{definition}

\begin{proposition}
	Si $\omega$ est une $1$-forme exacte de classe $C^1$, alors $\omega$ est fermée.
\end{proposition}

\begin{proof}
	Le fait que $\omega$ soit exacte implique l'existence d'une fonction $f$ telle que $\omega=df$, c'est à dire
	\begin{equation}
		\omega_x=\sum_i a_i(x)dx_i=\sum_i\frac{ \partial f }{ \partial x_i }(x)dx_i,
	\end{equation}
	c'est à dire que $a_i(x)=\frac{ \partial f }{ \partial x_i }(x)$. L'hypothèse que $\omega$ est $C^1$ implique que $f$ est $C^2$, et donc que nous pouvons inverser l'ordre de dérivation pour les dérivées secondes $\partial^2_{ij}f=\partial^2_{ji}f$. Nous pouvons donc faire le calcul suivant :
	\begin{equation}
		\frac{ \partial a_i }{ \partial x_j }=\frac{ \partial  }{ \partial x_j }\frac{ \partial f }{ \partial x_i }=\frac{ \partial  }{ \partial x_i }\frac{ \partial f }{ \partial x_j }=\frac{ \partial a_j }{ \partial x_i },
	\end{equation}
	ce qu'il fallait démontrer.
\end{proof}

Ceci est assez pour les formes différentielles pour cette année.

%---------------------------------------------------------------------------------------------------------------------------
\subsection{Intégration d'une forme différentielle sur un chemin}
%---------------------------------------------------------------------------------------------------------------------------

Les formes intégrales que nous avons déjà vues sont celles de fonctions et de champs de vecteur sur des chemins. Si $\gamma\colon [a,b]\to \eR^n$ est le chemin, les formules sont
\begin{equation}
	\begin{aligned}[]
		\int_{\gamma}f&=\int_{[a,b]}f\big( \gamma(t) \big)\| \gamma'(t) \|dt\\
		\int_{\gamma}G&=\int_{[a,b]}\langle G\big( \gamma(t) \big), \gamma'(t)\rangle dt.
	\end{aligned}
\end{equation}
Dans les deux cas, le principe est que nous disposons de quelque chose (la fonction $f$ ou le vecteur $G$), et du vecteur tangent $\gamma'(t)$, et nous essayons d'en tirer un nombre que nous intégrons. Lorsque nous avons une $1$-forme, la façon de l'utiliser pour produire un nombre avec le vecteur tangent est évidement d'appliquer la $1$-forme au vecteur tangent. La définition suivante est donc naturelle.

\begin{definition}
	Soit $\gamma\colon [a,b]\to \eR^n$, un chemin de classe $C^1$ tel que son image est contenue dans le domaine $D$. Si $\omega$ es une $1$-forme différentielle sur $D$, nous définissons l'\defe{intégrale de $\omega$ le long de $\gamma$}{Intégrale d'une forme différentielle} le nombre
	\begin{equation}
		\begin{aligned}[]
			\int_{\gamma}\omega&=\int_a^b\omega_{\gamma(t)}\big( \gamma'(t) \big)dt\\
				&=\int_a^b\Big[ a_1\big( \gamma(t) \big)\gamma'_1(t)+\ldots +  a_n\big( \gamma(t) \big)\gamma'_n(t) \Big]dt.
		\end{aligned}
	\end{equation}
\end{definition}

Cette définition est une bonne définition parce que si on change la paramétrisation du chemin, on ne change pas la valeur de l'intégrale, c'est la proposition suivante.
\begin{proposition}
	Si $\gamma$ et $\beta$ sont des chemins équivalents, alors
	\begin{equation}
		\int_{\gamma}\omega=\int_{\beta}\omega,
	\end{equation}
	c'est à dire que l'intégrale est invariante sous les reparamétrisation du chemin.
\end{proposition}
\begin{proof}
	Deux chemins sont équivalents quand il existe un difféomorphisme $C^1$ $h\colon [a,b]\to [c,d]$ tel que $\gamma(t)=(\beta\circ h)(t)$. En remplaçant $\gamma$ par $(\beta\circ h)$ dans la définition de $\int_{\gamma}\omega$, nous trouvons
	\begin{equation}
		\int_a^b\omega_{\gamma(t)}\big( \gamma'(t) \big)dt=\int_a^b\omega_{(\beta\circ h)(t)}\big( (\beta\circ h)'(t) \big)dt.
	\end{equation}
	Un changement de variable $u=h(t)$ transforme cette dernière intégrale en $\int_{\beta}\omega$, ce qui prouve la proposition.
\end{proof}

\begin{remark}
	Si $\gamma$ est une somme de chemins, $\gamma=\gamma^{(1)}+\ldots+\gamma^{(n)}$, où chacun des $\gamma^{(i)}$ est un chemin, alors
	\begin{equation}
		\int_{\gamma}\omega=\sum_{i=1}^n\int_{\gamma_i}\omega
	\end{equation}
	parce que $\omega$ est linéaire.
\end{remark}

\begin{remark}
	Si $-\gamma$ est le chemin
	\begin{equation}
		\begin{aligned}
			- \gamma\colon [a,b]&\to \eR^n \\
			t&\mapsto \gamma\big( b-(t-a) \big),
		\end{aligned}
	\end{equation}
	alors
	\begin{equation}
		\int_{-\gamma}\omega=-\int_{\gamma}\omega,
	\end{equation}
	c'est à dire que si l'on parcours le chemin en sens inverse, alors on change le signe de l'intégrale.
\end{remark}

L'intégrale d'une forme différentielle sur un chemin est compatible avec l'intégrale déjà connue d'un champ de vecteur sur le chemin parce que si $G$ est un champ de vecteurs,
\begin{equation}
	\int_{\gamma}G^{\flat}=\int_{\gamma}G.
\end{equation}
En effet,
\begin{equation}
	\begin{aligned}[]
		\int_{\gamma G^{\flat}}	&=\int_a^b G_{\gamma(t)}^{\flat}(\gamma'(t))\\
					&=\int_a^b\big[ G_1( \gamma(t) )dx_1+\ldots G_n(\gamma(t))dx_n \big]\big( \gamma'_1(t),\ldots,\gamma'_n(t) \big)\\
					&=\int_{a}^b\langle G(\gamma(t)), \gamma'(t)\rangle \\
					&=\int_{\gamma}G.
	\end{aligned}
\end{equation}


\begin{proposition}
	Soit $\omega=df$, une $1$-forme exacte et continue sur le domaine $D$. Alors la valeur de $\int_{\gamma}df$ ne dépend que des valeurs de $f$ aux extrémités de $\gamma$.
\end{proposition}

\begin{proof}
	Nous avons
	\begin{equation}
		\begin{aligned}[]
			\int_{\gamma}\omega=\int_{\gamma}df&=\int_{a}^b\sum_{i=1}n\frac{ \partial f }{ \partial x_i }\big( \gamma(t) \big)\gamma'_i(t)dt\\
				&=\int_a^b\frac{ d }{ dt }\Big( (f\circ\gamma)(t) \Big)dt\\
				&=(f\circ\gamma)(b)-(f\circ\gamma(a)).
		\end{aligned}
	\end{equation}
\end{proof}

%---------------------------------------------------------------------------------------------------------------------------
\subsection{Interprétation physique : travail}
%---------------------------------------------------------------------------------------------------------------------------

\begin{definition}
	Une force $F\colon D\subset\eR^n\to \eR^n$ est \defe{\href{http://fr.wikipedia.org/wiki/Force_conservative}{conservative}}{Conservative} si elle dérive d'un potentiel, c'est à dire si il existe une fonction $V\in C^1(D,\eR)$ telle que 
	\begin{equation}
		F(x)=(\nabla V)(x).
	\end{equation}
\end{definition}
Étant donné que $F$ est un champ de vecteurs, nous avons une forme différentielle associée $F^{\flat}$,
\begin{equation}
	F^{\flat}_x\colon x\mapsto \langle F(x), v\rangle .
\end{equation}

\begin{lemma}
	Le champ $F$ est conservatif si et seulement si la $1$-forme différentielle $F^{\flat}$ est exacte.
\end{lemma}

\begin{proof}
	Supposons que la force $F$ soit conservative, c'est à dire qu'il existe une fonction $V$ telle que $F=\nabla V$. Dans ce cas, il est facile de prouver que $F^{\flat}$ est exacte et est donnée par $F_x^{\flat}=dV(x)$. En effet,
	\begin{equation}
		\begin{aligned}[]
			F_x^{\flat}(v)	&=\langle F(x), v\rangle \\
					&=F_1(x)v_1+\ldots+F_n(x)v_n\\
					&=\frac{ \partial V }{ \partial x_1 }(x)v_1+\ldots\frac{ \partial V }{ \partial x_n }(x)v_n\\
					&=dV(x)v.
		\end{aligned}
	\end{equation}
	
	Pour le sens inverse, supposons que $F^{\flat}$ soit exacte. Dans ce cas, nous avons une fonction $V$ telle que $F^{\flat}=dV$. Il est facile de prouver qu'alors, $F=\nabla V$.
\end{proof}
En résumé, nous avons deux façons équivalentes d'exprimer que la force $F$ dérive du potentiel $V$ :  soit nous disons $F=\nabla V$, soit nous disons $F^{\flat}=dV$.

\begin{proposition}
	Si $F$ est une force conservative, alors le \href{http://fr.wikipedia.org/wiki/Travail_d'une_force}{travail} de $F$ lors d'un déplacement ne dépend pas du chemin suivit.
\end{proposition}

\begin{proof}
	Le travail d'une force le long d'un chemin n'est autre que l'intégrale de la force le long du chemin, et le calcul est facile :
	\begin{equation}
		W_{\gamma}(F)=\int_{\gamma}F=\int_{\gamma}dV=V\big( \gamma(b) \big)-V\big( \gamma(a) \big).
	\end{equation}
	Donc si $\beta$ est un autre chemin tel que $\beta(a)=\gamma(a)$ et $\beta(b)=\gamma(b)$, nous avons $W_{\beta}(F)=W_{\gamma}(F)$.
\end{proof}

%+++++++++++++++++++++++++++++++++++++++++++++++++++++++++++++++++++++++++++++++++++++++++++++++++++++++++++++++++++++++++++
\section{Intégrale sur une variété}
%+++++++++++++++++++++++++++++++++++++++++++++++++++++++++++++++++++++++++++++++++++++++++++++++++++++++++++++++++++++++++++

%---------------------------------------------------------------------------------------------------------------------------
\subsection{Mesure sur une carte}
%---------------------------------------------------------------------------------------------------------------------------

Nous considérons dans cette section uniquement des variétés $M$ de dimension $2$ dans $\eR^3$.  Une particularité de $\eR^3$ (par rapport aux autres $\eR^n$) est qu'il existe le produit vectoriel. 

Si $v$, $w\in\eR^3$, alors le vecteur $v\times w$ est une vecteur normal au plan décrit par $v$ et $w$ qui jouit de l'importante propriété suivante :
\begin{equation}
	\text{aire du parallélogramme}=\| v\times w \|.
\end{equation}
L'aire du parallélogramme construit sur $v$ et $w$ est donnée par la norme du produit vectoriel. Afin de donner une mesure infinitésimale en un point $p\in M$, nous voudrions prendre deux vecteurs tangents à $M$ en $p$, et puis considérer la norme de leur produit vectoriel. Cette idée se heurte à la question du choix des vecteurs tangents à considérer.

Dans $\eR^2$, le choix est évident : nous choisissons $e_x$ et $e_y$, et nous avons $\|e_x\times e_y\|=1$. L'idée est donc de choisir une carte $F\colon W\to F(w)$ autour du point $p=F(w)$, et de choisir les vecteurs tangents qui correspondent à $e_x$ et $e_y$ via la carte, c'est à dire les vecteurs
\begin{equation}
	\begin{aligned}[]
		\frac{ \partial F }{ \partial x }(w),&&\text{et}&&\frac{ \partial F }{ \partial y }(w).
	\end{aligned}
\end{equation}
L'\defe{élément infinitésimal de surface}{Élément de surface} sur $M$ au point $p=F(w)$ est alors défini par
\begin{equation}
	d\sigma_F=\|  \frac{ \partial F }{ \partial x }(w)\times\frac{ \partial F }{ \partial y }(w) \|dw,
\end{equation}
et si la partie $A\subset M$ est entièrement contenue dans $F(W)$, nous définissons la \defe{mesure}{mesure!dans une carte} de $A$ par
\begin{equation}		\label{EqDefMuDeuxDF}
	\mu_2(A)=\int_{F^{-1}(A)}d\sigma_F=\int_{F^{-1}(A)}\| \frac{ \partial F }{ \partial x }(w)\times\frac{ \partial F }{ \partial y }(w) \|dw.
\end{equation}
\begin{remark}
	Afin que cette définition ait un sens, nous devons prouver qu'elle ne dépend pas du choix de la carte $F$. En effet, les vecteurs $\partial_xF$ et $\partial_yF$ dépendent de la carte $F$, donc leur produit vectoriel (et sa norme) dépendent également de la carte $F$ choisie. Il faudrait donc un petit miracle pour que le nombre $\mu_2(A)$ donné par \eqref{EqDefMuDeuxDF} soit indépendant du choix de $F$.  Nous allons bientôt voir comme cas particulier du théorème \ref{ThoIntIndepF} que c'est en fait le cas. C'est à dire que si $F$ et $\tilde F$ sont deux cartes qui contiennent $A$, alors
	\begin{equation}
		\int_{F^{-1}(A)}d\sigma_F=\int_{\tilde F^{-1}(A)}d\sigma_{\tilde F}.
	\end{equation}
\end{remark}

%///////////////////////////////////////////////////////////////////////////////////////////////////////////////////////////
\subsubsection{Exemple : la mesure de la sphère}
%///////////////////////////////////////////////////////////////////////////////////////////////////////////////////////////

Nous nous proposons maintenant de calculer la surface de la sphère $S^2=x^2+y^2+z^2=R^2$. L'application $F\colon B( (0,0),R)\to R^3$ donnée par
\begin{equation}
	F(x,y)=\begin{pmatrix}
		x	\\ 
		y	\\ 
		\sqrt{R^2-x^2-y^2}	
	\end{pmatrix}
\end{equation}
est une carte pour une demi sphère. Ses dérivées partielles sont
\begin{equation}
	\begin{aligned}[]
		\frac{ \partial F }{ \partial x }&=\begin{pmatrix}
			1	\\ 
			0	\\ 
			-\frac{ x }{ \sqrt{R^2-x^2-y^2} }	
		\end{pmatrix},
		&\frac{ \partial F }{ \partial y }&=\begin{pmatrix}
			0	\\ 
			1	\\ 
			-\frac{ y }{ \sqrt{R^2-x^2-y^2} }	
		\end{pmatrix}.
	\end{aligned}
\end{equation}
Le produit vectoriel de ces deux vecteurs tangents donne
\begin{equation}
	\frac{ \partial F }{ \partial x }(x,y)\times\frac{ \partial F }{ \partial y }(x,y)=\frac{ x }{ \alpha }e_1+\frac{ y }{ \alpha }e_2+e_3
\end{equation}
où $\alpha=\sqrt{R^2-x^2-y^2}$. En calculant la norme, nous trouvons
\begin{equation}
	\| \frac{ \partial F }{ \partial x }(x,y)\times\frac{ \partial F }{ \partial y }(x,y)\| =\sqrt{  \frac{ R^2 }{ R^2-x^2-y^2 } },
\end{equation}
et en passant aux coordonnées polaires, nous écrivons l'intégrale \eqref{EqDefMuDeuxDF} sous la forme
\begin{equation}
	\int_B\| \partial_xF\times\partial_yF \|=\int_0^{2\pi}d\theta\int_0^R r\sqrt{  \frac{ R^2 }{ R^2-x^2-y^2 } }dr=2\pi R^2,
\end{equation}
qui est bien la mesure de la demi sphère.

%---------------------------------------------------------------------------------------------------------------------------
\subsection{Intégrale sur une carte}
%---------------------------------------------------------------------------------------------------------------------------

Nous pouvons maintenant définir l'intégrale d'une fonction sur une carte de la variété $M$.
\begin{definition}
	Soit $F\colon W\subset \eR^2\to \eR^3$, une carte pour une variété $M$. Soit $A$, une partie de $F(W)$ telle que $A=F(B)$ où $B\subset W$ est mesurable.  Soit encore $f\colon A\to \eR$, une fonction continue. L'\defe{intégrale}{intégrale!d'une fonction sur une carte} de $f$ sur $A$ est le nombre
	\begin{equation}	\label{EqDefIntDeuxDF}
		\int_Af=\int_Afd\sigma_F=\int_{F^{-1}(A)}(f\circ F)(w)\|  \frac{ \partial F }{ \partial x }(w)\times\frac{ \partial F }{ \partial y }(w) \| dw
	\end{equation}
\end{definition}

\begin{remark}
	L'intégrale \eqref{EqDefIntDeuxDF} n'est pas toujours bien définie. Étant donné que $F$ est $C^1$ et que $f$ est continue, l'intégrante est continue. L'intégrale sera donc bien définie par exemple lorsque $B$ est borné et si la fermeture $\bar A$ est un compact contenu dans $F(w)$.
\end{remark}

Le théorème suivant montre que le travail que nous avons fait jusqu'à présent ne dépend en fait pas du choix de carte $F$ effectué.

\begin{theorem}\label{ThoIntIndepF}
	Soient $F\colon W\to F(w)$ et $\tilde F\colon \tilde W\to \tilde F(\tilde W)$, deux cartes de la variété $M$. Soit une partie $A\subset F(W)\cap\tilde F(\tilde W)$ telle que $A=F(B)$ avec $B\subset W$ mesurable.  Alors $A=\tilde F(\tilde B)$ avec $\tilde B\subset\tilde W$ mesurable.

	Si $f$ est une fonction continue, et si $\int_Afd\sigma_F$ existe, alors $\int_Afd\sigma_{\tilde F}$ existe et
	\begin{equation}
		\int_Afd\sigma_F=\int_Afd\sigma_{\tilde F}.
	\end{equation}
\end{theorem}


%---------------------------------------------------------------------------------------------------------------------------
\subsection{Exemples}
%---------------------------------------------------------------------------------------------------------------------------

Intégrons la fonction $f(x,y,z)$ sur le carré $K=\mathopen] 0 , 1 \mathclose[\times \mathopen] 0 , 2 \mathclose[\times\{ 1 \}$. La première carte que nous pouvons utiliser est
\begin{equation}
	\begin{aligned}
		F\colon \mathopen] 0 , 1 \mathclose[\times\mathopen] 0 , 2 \mathclose[&\to K \\
		(x,y)&\mapsto (x,y,1). 
	\end{aligned}
\end{equation}
Nous trouvons aisément les vecteurs tangents qui forment l'élément de surface:
\begin{equation}
	\begin{aligned}[]
		\frac{ \partial F }{ \partial x }&=\begin{pmatrix}
			1	\\ 
			0	\\ 
			0	
		\end{pmatrix},
		&\frac{ \partial F }{ \partial y }&=\begin{pmatrix}
			0	\\ 
			1	\\ 
			0	
		\end{pmatrix},
	\end{aligned}
\end{equation}
donc $d\sigma_F=1\cdot dxdy$, et
\begin{equation}		\label{IntKSurcarrUn}
	\int_Kfd\sigma_F=\int_{\mathopen] 0 , 1 \mathclose[\times\mathopen] 0 , 2 \mathclose[}f(x,y,1)\cdot 1\cdot dxdy.
\end{equation}

Nous pouvons également utiliser la carte
\begin{equation}
	\begin{aligned}
		\tilde F\colon \mathopen] 0 , \frac{ 1 }{2} \mathclose[\times\mathopen] 0 , 6 \mathclose[&\to K \\
		(\tilde x,\tilde y)&\mapsto (2\tilde x,\frac{ \tilde y }{ 3 },1). 
	\end{aligned}
\end{equation}
Les vecteurs tangents sont maintenant
\begin{equation}
	\begin{aligned}[]
		\frac{ \partial \tilde F }{ \partial \tilde x }&=\begin{pmatrix}
			2	\\ 
			0	\\ 
			0	
		\end{pmatrix},
		&\frac{ \partial \tilde F }{ \partial \tilde y }&=\begin{pmatrix}
			0	\\ 
			1/3	\\ 
			0	
		\end{pmatrix},
	\end{aligned}
\end{equation}
de telle façon à ce que $d\sigma_{\tilde F}=\| \frac{ 2 }{ 3 }e_3 \|=\frac{ 2 }{ 3 }$. Cette fois, l'intégrale de $f$ sur $K$ s'écrit
\begin{equation}
	\int_Kfd\sigma_{\tilde F}=\int_{\mathopen] 0 , \frac{ 1 }{2} \mathclose[\times\mathopen] 0 , 6 \mathclose[}f\big( 2\tilde x,\frac{ \tilde y }{ 3 },1 \big)\cdot\frac{ 2 }{ 3 }\cdot d\tilde xs\tilde y.
\end{equation}
Conformément au théorème \ref{ThoIntIndepF}, cette dernière intégrale est égale à l'intégrale \eqref{IntKSurcarrUn} parce qu'il s'agit juste d'un changement de variable.


%---------------------------------------------------------------------------------------------------------------------------
\subsection{Orientation}
%---------------------------------------------------------------------------------------------------------------------------


Soient $F\colon W\to F(w)$ et $\tilde F\colon \tilde W\to \tilde F(\tilde W)$, deux cartes de la variété $M$. Nous pouvons considérer la fonction $h=\tilde F^{-1}\circ F$, définie uniquement sur l'intersection des cartes :
\begin{equation}
	h\colon F^{-1}\big( F(W)\cap\tilde F(\tilde W) \big)\to \tilde F^{-1}\big( F(W)\cap\tilde F(\tilde W) \big).
\end{equation}
Nous disons que $F$ et $\tilde F$ ont même \defe{orientation}{orientation} si
\begin{equation}
	J_h(w)>0
\end{equation}
pour tout $w\in  F^{-1}\big( F(W)\cap\tilde F(\tilde W) \big)$.

Considérons les deux carte suivantes pour le même carré:
\begin{equation}
	\begin{aligned}
		F\colon\mathopen] 0 , 1 \mathclose[\times \mathopen] 0 , 1 \mathclose[ &\to \eR^3 \\
		(x,y)&\mapsto (x,y,0) 
	\end{aligned}
\end{equation}
et
\begin{equation}
	\begin{aligned}
		\tilde F\colon\mathopen] 0 , \frac{ 1 }{2} \mathclose[\times\mathopen] 0 , \frac{1}{ 3 } \mathclose[ &\to \eR^3 \\
		(x,y)&\mapsto (2x,3y,0) 
	\end{aligned}
\end{equation}
Ici, $h(x,y)=\left( \frac{ x }{ 2 },\frac{ y }{ 3 } \right)$ et nous avons $J_h=\frac{1}{ 6 }>0$. Ces deux cartes ont même orientation. Notez que
\begin{equation}
	\frac{ \partial F }{ \partial x }\times\frac{ \partial F }{ \partial y }=e_3,
\end{equation}
tandis que
\begin{equation}
	\frac{ \partial \tilde F }{ \partial x }\times\frac{ \partial \tilde F }{ \partial y }=6e_3.
\end{equation}
Les vecteurs normaux à la paramétrisation pointent dans le même sens.

Si par contre nous prenons la paramétrisation
\begin{equation}
	\begin{aligned}
		G\colon \mathopen] 0,1 \mathclose[\times\mathopen] 0,1 ,  \mathclose[&\to \eR^2 \\
		(x,y)&\mapsto (x,(1-y),0), 
	\end{aligned}
\end{equation}
nous avons
\begin{equation}
	\frac{ \partial G }{ \partial x }\times\frac{ \partial G }{ \partial y }=-e_3,
\end{equation}
et si $g=G^{-1}\circ F$, alors $J_g=-1$.

L'orientation d'une carte montre donc si le vecteur normal à la surface pointe d'un côté ou de l'autre de la surface.

\begin{definition}
	Une variété $M$ est \defe{orientable}{orientable!variété} si il existe un atlas de $M$ tel que deux cartes quelconques ont toujours même orientation. Une variété est \defe{orientée}{variété !orientée} lorsque qu'un tel choix d'atlas est fait.
\end{definition}

\begin{proposition}
	Soit $M$, une variété orientable et un atlas orienté $\{ F_i\colon W_i\to \eR^3 \}$. Alors le vecteur unitaire
	\begin{equation}
		\frac{   \frac{ \partial F }{ \partial x }(x,y)\times\frac{ \partial F }{ \partial y }(x,y)   }{ \| \frac{ \partial F }{ \partial x }(x,y)\times\frac{ \partial F }{ \partial y }(x,y)\| }
	\end{equation}
	ne dépend pas du choix de $F$ parmi les $F_i$.
\end{proposition}


\begin{proof}
	Considérons deux cartes $F_1$ et $F_2$, ainsi que l'application $h=F_2^{-1}\circ F_1$. Écrivons le vecteur $\partial_x F_1\times\partial_yF_1$ en utilisant $F_1=F_2\circ h$. D'abord, par la règle de dérivation de fonctions composées,
	\begin{equation}
		\frac{ \partial (F_2\circ h) }{ \partial x }=\frac{ \partial F_2 }{ \partial x }\frac{ \partial h_1 }{ \partial x }+\frac{ \partial F_2 }{ \partial y }\frac{ \partial h_2 }{ \partial x }.
	\end{equation}
	Après avoir fait le même calcul pour $\frac{ \partial (F_2\circ h) }{ \partial y }$, nous pouvons écrire
	\begin{equation}
		\partial_x(F_2\circ h)\times\partial_y(F_2\circ h)=(\partial_xh_1\partial_xF_2+\partial_xh_2\partial_yF_2)\times(\partial_yh_1\partial_xF_2+\partial_yh_2\partial_yF_2).
	\end{equation}
	Dans cette expression, les facteurs $\partial_ih_j$ sont des nombres, donc ils se factorisent dans les produits vectoriels. En tenant compte du fait que $\partial_xF_2\times\partial_xF_2=0$ et $\partial_yF_2\times\partial_yF_2=0$, ainsi que de l'antisymétrie du produit vectoriel, l'expression se réduit à
	\begin{equation}
		\left( \frac{ \partial F_2 }{ \partial x }\times\frac{ \partial F_2 }{ \partial y } \right)(\partial_xh_1\partial_yh_2-\partial_xh_2\partial_yh_2).
	\end{equation}
	Par conséquent,
	\begin{equation}
		\frac{ \partial F_1 }{ \partial x }\times\frac{ \partial F_1 }{ \partial y } =\frac{ \partial (F_2\circ h) }{ \partial x }\times\frac{ \partial (F_2\circ h) }{ \partial y } =\left( \frac{ \partial F_2 }{ \partial x }\times\frac{ \partial F_2 }{ \partial y } \right)\det J_h.
	\end{equation}
	Donc, tant que $J_h$ est positif, les vecteurs unitaires correspondants au membre de gauche et de droite sont égaux.
\end{proof}

\begin{corollary}
	Si nous avons choisit un atlas orienté pour la variété $M$, nous avons une fonction continue $G\colon M\to \eR^3$ telle que $\| G(p) \|=1$ pour tout $p\in M$. Cette fonction est donnée par
	\begin{equation}		\label{DefCarteGOritn}
		G(F(x,y))=\frac{   \frac{ \partial F }{ \partial x }(x,y)\times\frac{ \partial F }{ \partial y }(x,y)   }{ \| \frac{ \partial F }{ \partial x }(x,y)\times\frac{ \partial F }{ \partial y }(x,y)\| }
	\end{equation}
	sur l'image de la carte $F$.
\end{corollary}

\begin{proof}
	La fonction $G$ est construite indépendamment sur chaque carte $F(W)$ en utilisant la formule \eqref{DefCarteGOritn}. Cette fonction est une fonction bien définie sur tout $M$ parce que nous venons de démontrer que sur $F_1(W_1)\cap F_2(W_2)$, les fonctions construites à partir de $F_1$ et à partir de $F_2$ sont égales.
\end{proof}

Il est possible que prouver, bien que cela soit plus compliqué, que la réciproque est également vraie.
\begin{proposition}
	Une variété $M$ de dimension $2$ dans $\eR^3$ est orientable si et seulement si il existe une fonction continue $G\colon M\to \eR^3$ telle que pour tout $p\in M$, le vecteur $G(p)$ soit de norme $1$ et normal à $M$ au point $p$.
\end{proposition}

%---------------------------------------------------------------------------------------------------------------------------
\subsection{Intégrale d'une fonction sur une variété}
%---------------------------------------------------------------------------------------------------------------------------

Nous supposons à présent que $M$ est une variété compacte de dimension $2$ dans $\eR^3$. La compacité fait que $M$ possède un atlas contenant un nombre fini de cartes $F_i\colon W_i\to F_i(W_i)$. 

Si $A\subset M$ est tel que pour chaque $i$, $A\cap F_i(W_i)=F_i(V_i)$ pour une ensemble $V_i$ mesurable dans $\eR^2$, alors nous considérons
\begin{equation}
	A_1=A\cap F_1(W_2)=F_1(V_1).
\end{equation}
Ensuite, nous construisons $A_2$ en considérant $F_A(W_2)$ et en lui retranchant $A_1$ :
\begin{equation}
	A_2=\big( A\cap F_2(W_2) \big)\cap F_1(V_1).
\end{equation}
En continuant de la sorte, nous construisons la décomposition
\begin{equation}
	A=A_1\cup\ldots\cup A_p
\end{equation}
de $A$ en ouverts disjoints, chacun de ouverts $A_p$ étant compris dans une carte.

Il est possible de prouver que dans ce cas, la définition suivante a un sens et ne dépend pas du choix de l'atlas effectué.
\begin{definition}
	Si $f\colon A\to \eR$ est une fonction continue, alors l'\defe{intégrale}{intégrale!d'une fonction sur une variété} est le nombre
	\begin{equation}
		\int_Af=\sum_{i=1}^p\int_{A_i}fd\sigma_{F_i}.
	\end{equation}
\end{definition}

%+++++++++++++++++++++++++++++++++++++++++++++++++++++++++++++++++++++++++++++++++++++++++++++++++++++++++++++++++++++++++++
\section{Variétés et extrema liés}
%+++++++++++++++++++++++++++++++++++++++++++++++++++++++++++++++++++++++++++++++++++++++++++++++++++++++++++++++++++++++++++

\subsection{Introduction}
Soit $f : S^2 \to \eR$ une fonction définie sur la sphère usuelle
$S^2 \subset \eR^3$. Une question naturelle est d'estimer la
régularité de $f$ ; est-elle continue, dérivable, différentiable ? Il
n'existe pas de dérivée directionnelle étant donné que le quotient
différentiel
\begin{equation*}
  \frac{f(x + \epsilon u_1 ,y + \epsilon u_2) - f(x,y)}{\epsilon}
\end{equation*}
n'a pas de sens pour un point $(x + \epsilon u_1 ,y + \epsilon u_2)$
qui n'est pas --sauf valeurs particulières-- dans la surface. Pour la
même raison il n'est pas possible de parler de différentiabilité de
cette manière. Comment faire, sans devoir étendre le domaine de
définition de $f$ à un voisinage de la sphère ? Une solution possible
est de parler de la notion de variété.

Une variété est un objet qui ressemble, vu de près, à $\eR^m$ pour un
certain $m$. En d'autres termes, on imagine une variété comme un
recollement de morceaux de $\eR^m$ vivant dans un espace plus grand
$\eR^n$. Ces morceaux sont appelés des ouverts de carte, et
l'application qui exprime la ressemblance à $\eR^m$ est l'application
de carte.

%---------------------------------------------------------------------------------------------------------------------------
\subsection{Définition et propriétés}
%---------------------------------------------------------------------------------------------------------------------------

\begin{definition}
  Soit $\emptyset \neq M \subset \eR^n$, $1 \leq m < n$ et $k \geq
  1$. $M$ est une \Defn{variété de classe $C^k$ de dimension $m$} si
  pour tout $a \in \var M$, il existe un voisinage ouvert $U$ de $a$
  dans $\eR^n$, et un ouvert $V$ de $\eR^m$ tel que $U \cap \var M$
  soit le graphe d'une fonction $f : V \subset \eR^m \to \eR^{n-m}$
  de classe $C^1$, c'est-à-dire qu'il existe un réagencement des
  coordonnées $(x_{i_1}, \ldots, x_{i_m}, x_{i_{m+1}}, \ldots,
  x_{i_n})$ avec
  \begin{equation*}
    M \cap U = \left\{ (x_1, \ldots, x_n) \in \eR^n \tq
%      \begin{array}{l} % deux conditions
      (x_{i_1}, \ldots, x_{i_m}) \in V \quad \left\{\begin{array}{c!{=}l} % 1: equations
        x_{i_{m+1}} & f_1(x_{i_1}, \ldots, x_{i_m})\\
        \vdots & \vdots \\
        x_{i_n} & f_{n-m}(x_{i_1}, \ldots, x_{i_m})
      \end{array}\right.
%    \end{array}
    \right\}
  \end{equation*}
  où $V$ est un voisinage ouvert de $(a_{i_1}, \ldots, a_{i_m}) \in \eR^m$.
\end{definition}


Le cours regorge de théorèmes qui proposent des conditions équivalentes à la définition d'une variété. Celle que nous allons le plus utiliser est la suivante, de la page 268.

\begin{proposition}
	Soit $M\subset\eR^n$ et $1\leq m\leq n-1$. L'ensemble $M$ est une variété si et seulement si $\forall a\in M$, il existe un voisinage ouvert $\mU$ de $a$ dans $\eR^n$ et une application $F\colon W\subset\eR^m\to \eR^n$ où $W$ est un ouvert tels que
	\begin{enumerate}
		\item
			$F$ est un homéomorphisme de $W$ vers $M\cap\mU$,
		\item
			$F\in C^1(W,\eR^n)$,
		\item
			Le rang de $dF(w)\in L(\eR^m,\eR^n)$ est de rang maximum (c'est à dire $m$) en tout point $w\in W$.
	\end{enumerate}
\end{proposition}
Pour rappel, si $T\colon \eR^m\to \eR^n$ est une application linéaire, son \defe{rang}{rang} est la dimension de son image. On peut prouver que si $A$ est la matrice d'une application linéaire, alors le rang de cette application linéaire est égal à la taille de la plus grande matrice carré de déterminant non nul contenue dans $A$.

La condition de rang maximum sert à éviter le genre de cas de la figure \ref{LabelFigExempleNonRang} qui représente l'image de l'ouvert $\mathopen] -1 , 1 \mathclose[$ par l'application $F(t)=(t^2,t^3)$.
\newcommand{\CaptionFigExempleNonRang}{Quelque chose qui n'est pas de rang maximum et qui n'est pas une variété.}
\input{Fig_ExempleNonRang.pstricks}
La différentielle a pour matrice
\begin{equation}
	dF(t)=(2t,3t^2).
\end{equation}
Le rang maximum est $1$, mais en $t=0$, la matrice vaut $(0,0)$ et son rang est zéro. Pour toute autre valeur de $t$, c'est bon.

Une autre caractérisation des variétés est donnée par la proposition suivante (proposition 3, page 274).
\begin{proposition}		\label{PropCarVarZerFonc}
	Soit $M\in \eR^n$ et $1\leq m\leq n-1$. L'ensemble $M$ est une variété si et seulement si $\forall a\in M$, il existe un voisinage ouvert $\mU$ de $a$ dans $\eR^n$ tel et une application $G\in C^1(\mU,\eR^{n-m})$ tel que
	\begin{enumerate}

		\item
			le rang de $dG(a)\in L(\eR^n,\eR^{n-m})$ soit maximum (c'est à dire $n-m$) en tout $a\in M$,
		\item
			$M\cap\mU=\{ x\in\mU\tq G(x)=0 \}$.

	\end{enumerate}
\end{proposition}


%---------------------------------------------------------------------------------------------------------------------------
\subsection{Espace tangent}
%---------------------------------------------------------------------------------------------------------------------------

Soit $M$, une variété dans $\eR^n$, et considérons un chemin $\gamma\colon I\to \eR^n$ tel que $\gamma(t)\in M$ pour tout $t\in I$ et tel que $\gamma(0)=a$ et que $\gamma$ est dérivable en $0$. La \defe{tangente}{tangente à un chemin} au chemin $\gamma$ au point $a\in M$ est la droite
\begin{equation}
	s\mapsto a+s\gamma'(0).
\end{equation}
L'\defe{espace tangent}{Espace tangent} de $M$ au point $a$ est l'ensemble décrit par toutes les tangentes en $a$ pour tous les chemins $\gamma$ possibles.

\begin{proposition}			\label{PropDimEspTanVarConst}
	Une variété de dimension $m$ dans $\eR^n$ a un espace tangent de dimension $m$ en chacun de ses points.
\end{proposition}

%---------------------------------------------------------------------------------------------------------------------------
\subsection{Extrema liés}
%---------------------------------------------------------------------------------------------------------------------------
 
Soit $f$, une fonction sur $\eR^n$, et $M\subset \eR^n$ une variété de dimension $m$. Nous voulons savoir quelle sont les plus grandes et plus petites valeurs atteintes par $f$ sur $M$.

Pour ce faire, nous avons un théorème qui permet de trouver des extrema \emph{locaux} de $f$ sur la variété. Pour rappel, $a\in M$ est une \defe{extrema local de $f$ relativement}{Extrema local!relatif} à l'ensemble $M$ si il existe une boule $B(a,\epsilon)$ telle que $f(a)\leq f(x)$ pour tout $x\in B(a,\epsilon)\cap M$.

\begin{theorem}
	Soit $A\subset\eR^n$ et
	\begin{enumerate}

		\item
			une fonction (celle à minimiser) $f\in C^1(A,\eR)$,
		\item 
			des fonctions (les contraintes) $G_1,\ldots,G_k\in C^1(A,\eR)$,
		\item
			$M=\{ x\in A\tq G_i(x)=0\,\forall i\}$,
		\item
			un extrema local $a\in M$ de $f$ relativement à $M$.
	\end{enumerate}
	Supposons que les gradients $\nabla G_1(a)$, \ldots,$\nabla G_k(a)$ soient linéairement indépendants.

	Alors $a=(x_1,\ldots,x_n)$ est dans les solutions de $\nabla L=0$ où
	\begin{equation}
		L(x_1,\ldots,x_n,\lambda_1,\ldots,\lambda_k)=f(x_1,\ldots,x_n)+\sum_{i=1}^k\lambda_iG_i(x_1,\ldots,x_n)
	\end{equation}
\end{theorem}
La fonction $L$ est le \defe{lagrangien}{lagrangien} du problème.

\href{http://www.sagenb.org/home/pub/353/}{En pratique}, les candidats extrema locaux sont tous les points où les gradients ne sont pas linéairement indépendants, plus tous les points donnés par l'équation $\nabla L=0$. Parmi ces candidats, il faut trouver lesquels sont maxima ou minima, locaux ou globaux.

L'existence d'extrema locaux se prouve généralement en invoquant de la compacité, et en invoquant le lemme suivant qui permet de réduire le problème à un compact.

\begin{lemma}		\label{LemmeMinSCimpliqueS}
	Soit $S$, un ensemble dans $\eR^n$ et $C$, un ouvert de $\eR^n$. Si $a\in\Int S$ est un minimum local relatif à $S\cap C$, alors il est un minimum local par rapport à $S$.
\end{lemma}

\begin{proof}
	Nous avons que $\forall x\in B(a,\epsilon_1)\cap S\cap C$, $f(x)\geq f(x)$. Mais étant donné que $C$ est ouvert, et que $a\in C$, il existe un $\epsilon_2$ tel que $B(a,\epsilon_2)\subset C$. En prenant $\epsilon=\min\{ \epsilon_1,\epsilon_2 \}$, nous trouvons que $f(x)\geq f(a)$ pour tout $x\in B(a,\epsilon)\cap(S\cap C)=B(a,\epsilon)\cap S$.
\end{proof}

\section{Intégrales curvilignes}
\label{secintcurvi}

\subsection{Chemins de classe \texorpdfstring{$C^1$}{C1}}

Soit $p, q\in \eR^n$. Un \defe{chemin}{chemin} $C^1$ par morceaux joignant $p$ à $q$ est une application continue
\begin{equation}
  \gamma : [a,b] \to \eR^n \qquad \gamma(a) = p, \gamma(b) = q
\end{equation}
pour laquelle il existe une subdivision $a = t_0 < t_1 < \ldots < t_{r-1} < t_r = b$ telle que :
\begin{enumerate}
\item la restriction de $\gamma$ sur chaque ouvert $\mathopen]t_i,
  t_{i+1}\mathclose[$ est de classe $C^1$~;
\item pour tout $0 \leq i \leq r$, $\gamma^\prime$ possède une limite
  à gauche (sauf pour $i = 0$) et une limite à droite (sauf pour $i =
  r$) en $t_i$.
\end{enumerate}
Le \defe{chemin $\gamma$ est (globalement) de classe $C^1$}{Chemin!classe $C^2$} si la
subdivision peut être choisie de \og longueur\fg{} $r = 1$.

\begin{remark}
	Si $a$ et $b$ sont des points de
  $\eR^n$, on peut créer le chemin particulier
  \begin{equation*}
    \gamma : [0,1] \to \eR^n : t \mapsto (1-t)a + tb
  \end{equation*}
  qui relie ces points par un segment de droite.
\end{remark}

\subsection{Intégrer une fonction}
Soit $f : D \subset \eR^n \to \eR$ une fonction continue, et $\gamma
: [a,b] \to D$ un chemin $C^1$. On définit \Defn{l'intégrale de $f$
  sur $\gamma$} par
\begin{equation*}
  \int_\gamma f d s = \int_\gamma f = \int_a^b f(\gamma(t)) \norme{\gamma^\prime(t)} d t.
\end{equation*}

\begin{remark}
  Cette définition ne dépend pas de la paramétrisation choisie, ni du
  sens du chemin (échange entre point de départ et point d'arrivée).
\end{remark}

La formule qui donne la longueur d'un chemin est évidement l'intégrale de la fonction $1$ sur le chemin, c'est à dire
\begin{equation}
	L=\int_a^b\| \gamma'(t) \|dt.
\end{equation}
Si on veut savoir la longueur d'une courbe donnée sous la forme d'une fonction $y=y(x)$, un chemin qui trace la courbe est évidement donné par
\begin{equation}
	\gamma(t)=(t,y(t)),
\end{equation}
et le vecteur tangent au chemin est $\gamma'(t)=(1,y'(t))$. Donc
\begin{equation}
	\| \gamma'(t) \|=\sqrt{1+y'(t)^2},
\end{equation}
et 
\begin{equation}			\label{EqLongFonction}
	L=\int_a^b\sqrt{1+y'(t)^2}.
\end{equation}

\subsection{Intégrer un champ de vecteurs}
Un \Defn{champ de vecteur} est une application $G : \eR^n \to
\eR^n$. On définit l'intégrale de $G$ sur un chemin $\gamma : [a,b]
\to \eR^n$ par
\begin{equation*}
  \int_\gamma G \pardef \int_a^b \scalprod {G(\gamma(t))}{\gamma^\prime(t)} d t.
\end{equation*}

\begin{remark}
  Cette définition ne dépend pas de la paramétrisation choisie, mais
  le signe change selon le sens du chemin.
\end{remark}

%---------------------------------------------------------------------------------------------------------------------------
\subsection{Intégrer une forme différentielle sur un chemin}
%---------------------------------------------------------------------------------------------------------------------------

Une \defe{forme différentielle}{Forme!différentielle} sur $\eR^n$ est une application
\begin{equation}
	\begin{aligned}
		\omega\colon \eR^n&\to (\eR^n)^* \\
		x&\mapsto \omega_x 
	\end{aligned}
\end{equation}
qui à chaque point $x$ de $\eR^n$ associe une forme linéaire $\omega_x: \eR^n \to \eR$.

On sait que $\{ d x_i \}_{1\leq i\leq n}$ est une base de
${(\eR^{n})}^{*}$, donc toute forme différentielle s'écrit
\begin{equation*}
  \omega_x = \sum_{i=0}^n a_i(x) d x_i
\end{equation*}
où $a_1,\ldots,a_n$ sont les \Defn{composantes de $\omega$} dans la
base usuelle, et sont des fonctions à valeurs réelles. Pour un vecteur
$v = (v_1,\ldots,v_n)$ on a donc par définition de $d x_i$
\begin{equation*}
  \omega_x v = \sum_{i=0}^n a_i(x) v_i.
\end{equation*}

L'intégrale d'une forme différentielle sur un chemin est définie par
\begin{equation*}
  \int_\gamma \omega = \int_a^b \omega_{\gamma(t)}\gamma^\prime(t) d t
\end{equation*}

\begin{remark}
  Cette définition ne dépend pas de la paramétrisation choisie, mais
  le signe change selon le sens du chemin.
\end{remark}

\subsection{Lien entre forme différentielle et champ vectoriel}
Si $G$ est un champ de vecteurs, on peut définir la forme différentielle
\begin{equation*}
  \omega^G : \eR^n \to {(\eR^n)}^\ast : x \mapsto \left\lbrack \omega^G_x :
  \eR^n \to \eR : v \mapsto \omega^G_x v = \scalprod {G(x)}v \right\rbrack
\end{equation*}
et réciproquement, si $\omega_x = \sum_i a_i(x)d x_i$ est une forme
différentielle on définit le champ de vecteurs
\begin{equation*}
  G^\omega(x) = (a_1(x),\ldots,a_n(x)).
\end{equation*}

Avec ces définitions, pour un chemin $\gamma$ donné on a
\begin{equation*}
  \int_\gamma \omega^G = \int_\gamma G^\omega
\end{equation*}

%---------------------------------------------------------------------------------------------------------------------------
\subsection{Intégrer un champs de vecteurs sur un bord en $2D$}
%---------------------------------------------------------------------------------------------------------------------------

Si $D\subset\eR^2$ est tel que $\partial D$ est une variété de dimension $1$ et tel que $D$ accepte un champ de vecteur normal extérieur unitaire $\nu$. Si nous voulons définir 
\begin{equation}
	\int_{\partial D}G,
\end{equation}
le mieux est de prendre une paramétrisation $\gamma\colon \mathopen[ 0 , 1 \mathclose]\to \eR^2$ et de calculer
\begin{equation}
	\int_0^1 \langle G_{\gamma(t)}, \frac{ \dot\gamma(t) }{ \| \dot\gamma(t) \| }\rangle dt.
\end{equation}
Hélas, cette définition ne fonctionne pas parce que son signe dépend du sens de la paramétrisation $\gamma$. Si la paramétrisation tourne dans l'autre sens, il y a un signe de différence.

Nous allons définir
\begin{equation}		\label{EqIntVectbordDeux}
	\int_{\partial D}G=\int_0^1\langle G_{\gamma(t)}, T(t)\rangle dt
\end{equation}
où $T(t)=\dot\gamma(t)/\| \dot\gamma(t) \|$ et où $\gamma$ est choisit de telle façon à ce que la rotation d'angle $\frac{ \pi }{ 2 }$ amène $\nu$ sur $T$. Cela fixe le choix de sens.

Ce choix de sens aura des répercussions dans l'application de la formule de Green et du théorème de Stokes.

%---------------------------------------------------------------------------------------------------------------------------
\subsection{Intégrer une forme différentielle sur un bord en $2D$}
%---------------------------------------------------------------------------------------------------------------------------

Nous n'allons pas chercher très loin :
\begin{equation}
	\int_{\partial D}\omega=\int_{\partial D}\omega^{\sharp},
\end{equation}
c'est à dire que l'intégrale de la forme différentielle est celle du champ de vecteur associé. Le membre de droite est définit par \eqref{EqIntVectbordDeux}, avec le choix d'orientation qui va avec.

%---------------------------------------------------------------------------------------------------------------------------
\subsection{Intégrer une forme différentielle sur un bord en $3D$}
%---------------------------------------------------------------------------------------------------------------------------

Nous allons maintenant intégrer une forme différentielle sur certains chemins fermés dans $\eR^3$. Soit $F(D)\subset\eR^3$, une variété de dimension $2$ dans $\eR^3$ où $F\colon D\subset\eR^2\to \eR^3$ est la carte. Nous supposons que $D$ vérifie les hypothèses de la formule de Green. Alors nous définissons
\begin{equation}		\label{EqDefIntTroisForBord}
	\int_{F(\partial D)}\omega = \int_{\partial D} F^*\omega
\end{equation}
où $F^*\omega$ est la forme différentielle définie sur $\partial D$ par $(F^*\omega)(v)=\omega\big( dF(v) \big)$.

Cette définition est très abstraite, mais nous n'allons, en pratique, jamais l'utiliser, grâce au théorème de Stokes.

%---------------------------------------------------------------------------------------------------------------------------
\subsection{Intégrer d'un champ de vecteurs sur un bord en $3D$}
%---------------------------------------------------------------------------------------------------------------------------

Encore une fois, nous n'allons pas chercher bien loin :
\begin{equation}
	\int_{F(\partial D)G}=\int_{F(\partial D)}G^{\flat}
\end{equation}
où $G^{\flat}$ est la forme différentielle associée au champ de vecteur. Le membre de droite est définit par l'équation \eqref{EqDefIntTroisForBord}.

%+++++++++++++++++++++++++++++++++++++++++++++++++++++++++++++++++++++++++++++++++++++++++++++++++++++++++++++++++++++++++++
\section{Intégrales de surface}
%+++++++++++++++++++++++++++++++++++++++++++++++++++++++++++++++++++++++++++++++++++++++++++++++++++++++++++++++++++++++++++

%---------------------------------------------------------------------------------------------------------------------------
\subsection{Intégrale d'une fonction}
%---------------------------------------------------------------------------------------------------------------------------
\label{secintsurfaciques}
Soit $M$ une variété de dimension $n$ dans $\eR^m$. Soit $F : W \subset \eR^n \to M$ une paramétrisation d'un ouvert relatif de $M$.  

Si $f$ est une fonction définie sur un sous-ensemble $A \subset F(W)$ tel que $F^{-1}(A)$ est mesurable, l'\Defn{intégrale de $f$ sur $A$} est définie par
\begin{equation*}
  \int_A f = \int_{F^{-1}(A)} f(F(w)) \sqrt{\det(\transpose{J_F(w)} {J_F(w)})} dw
\end{equation*}
où l'intégrale est l'intégration usuelle (de Lebesgue) sur $F^{-1}(A) \subset \eR^n$. On écrit parfois cette intégrale $\int_{F^{-1}(A)} f(F(w)) d\sigma$ où
\begin{equation*}
  d\sigma = \sqrt{\det(\transpose{J_F(w)} {J_F(w)})} dw
\end{equation*}
est l'\Defn{élément infinitésimal de volume} de la variété. 

Si $m = 3$ et $n = 2$, l'élément infinitésimal de volume vaut
\begin{equation*}
  d \sigma = \norme{\pder F {w_1} \wedge \pder F {w_2}} dw
\end{equation*}
où $\wedge$ représente le produit vectoriel dans $\eR^3$, et $(w_1,w_2)$ sont les coordonnées sur $W \subset \eR^2$. Dans la suite, nous ne regarderons plus que ce cas.

\subsection{Intégrale d'un champ de vecteurs}
Dans l'intégration curviligne, on a noté que si l'intégrale d'une fonction ne dépendait pas de l'orientation du chemin, l'intégrale d'un champ de vecteurs ou d'une forme différentielle en dépendait. Ce problème d'orientation apparait également dans l'intégration sur des surfaces de l'espace.

%% Page 530, exemple 4
Une \Defn{orientation} sur une surface $S \subset \eR^3$ est le choix
d'un champ de vecteurs continu $\nu : S \to \eR^3$ dont la norme en
tout point de $S$ vaut $1$. On remarque qu'ayant fait un tel choix
d'orientation $\nu(x)$ en un point $x$, le seul autre choix possible
en $x$ est $-\nu(x)$.
%% Page 
Si $S$ est le bord d'un ouvert $D \subset \eR^3$, l'\Defn{orientation
  induite par $D$ sur $S$} est, si elle existe, l'orientation qui
pointe hors de $D$ en tout point de $S$. Plus précisément, il faut que
pour tout $x \in D$ il existe $\epsilon > 0$ vérifiant, pour tout $0 <
t < \epsilon$, la relation $t \nu(x) \notin D$. Dans ce cas, le champ
de vecteurs $\nu$ est appelé le \Defn{vecteur normal unitaire
  extérieur} à $D$ et il est forcément unique.

Soit $G$ un champ de vecteurs défini sur une surface orientée par un
champ $\nu$. L'intégrale de $G$ sur $S$, aussi appelée le \Defn{flux
de $G$ à travers $S$}, est
\begin{equation}\label{eqflux-star}
  \iint_S G \cdot d S \pardef \iint_S \scalprod{G}{\nu} d \sigma.
\end{equation}
Si on suppose que la surface est paramétrisée par une application
\begin{equation*}
  F : W \subset \eR^2 \to \eR^3 : (u,v) \mapsto (F_1(u,v),F_2(u,v),F_3(u,v))
\end{equation*}
alors un vecteur unitaire $\nu$ peut s'écrire sous la forme
\begin{equation*}
  \nu = \frac{\pder F u \wedge \pder F v}{\norme{\pder F u \wedge \pder F v}}
\end{equation*}
et grâce à cette paramétrisation l'intégrale \eqref{eqflux-star}
devient
\begin{equation*}
  \iint_S G \cdot d S = \iint_W \scalprod{G(F(u,v))}{\pder F u \wedge \pder F v} d u
  d v.
\end{equation*}
où on utilise l'expression de $d \sigma$ obtenue précédemment dans le
cas qui nous intéresse (surface dans l'espace).

%+++++++++++++++++++++++++++++++++++++++++++++++++++++++++++++++++++++++++++++++++++++++++++++++++++++++++++++++++++++++++++
\section{Divergence, Green, Stokes}
%+++++++++++++++++++++++++++++++++++++++++++++++++++++++++++++++++++++++++++++++++++++++++++++++++++++++++++++++++++++++++++


Le théorème de Stokes (et ses variations) peut se voir comme une généralisation du théorème fondamental du calcul différentiel et intégral qui stipule que
\begin{equation*}
	\int_a^b f^\prime(x) d x = f(b) - f(a)
\end{equation*}
c'est-à-dire qui relie l'intégrale de $f^\prime$ sur $I = [a,b]$ aux valeurs de $f$ sur le bord $\partial I = \{a,b\}$. Le signe $-$ qui apparait vient de l'orientation ; celle-ci requiert de la prudence dans l'utilisation des théorèmes.

Voici, pour votre culture générale, un énoncé général :
\begin{theorem}
	Si $M$ est une variété orientable de dimension $n$ avec un bord noté $\partial  M$, alors pour toute forme différentielle $\omega$ de degré $n-1$ on a 
	\begin{equation*}
		\int_{ M} d \omega = \int_{\partial  M} \omega.
	\end{equation*}
	où $d \omega$ désigne la différentielle extérieure de $\omega$.
\end{theorem}
Nous allons maintenant voir quelque cas particuliers. 


\subsection{Théorème de la divergence}

Si nous considérons une surface dans $\eR^n$ et un champ de vecteurs, il est bon de se demander quelle \og quantité de vecteurs\fg{} traverse la surface. Soit $D$, un ouvert borné de $\eR^n$ telle que $\partial D$ soit une variété de dimension $n-1$, et $G$, un champ de vecteurs défini sur $\bar D$. Afin de compter combien de $G$ traverse $\partial D$, il faudra faire en sorte de ne considérer que la composante de $G$ normale à $\partial D$ : pas question d'intégrer par exemple la norme de $G$ sur $\partial D$.

Comme nous le savons, la composante du vecteur $v$ dans la direction $w$ est le produit scalaire $v\cdot 1_w$ où $1_w$ est le vecteur de norme $1$ dans la direction $w$. Nous allons donc introduire le concept de vecteur normal extérieur. Soit $x\in\partial D$ et $\nu\in\eR^n$, nous disons que $\nu$ est un \defe{vecteur normal extérieur}{Normal extérieur!vecteur} de $\partial D$ si
\begin{enumerate}

	\item
		$\langle \nu, v\rangle =0$ pour tout vecteur tangent $v$ à $\partial D$ au point $x$. Pour rappel, $\partial D$ étant une variété de dimension $n-1$, il y a $n-1$ tels vecteurs $v$ linéairement indépendants.
	
	\item
		Il existe un $\delta>0$ tel que $\forall t\in\mathopen] 0 , \delta \mathclose[$, nous avons $c+t\nu\notin \bar D$ et $x-t\nu\in D$.
 
\end{enumerate}

Nous pouvons maintenant définir le concept de flux. Soit $D\subset \eR^n$ tel que $\partial D$ soit une variété de dimension $n-1$ qui admette un vecteur normal extérieur $\nu(x)$ en chaque point. Soit aussi $G\colon \bar D\to \eR^n$, un champ de vecteur de classe $C^1$. Le \defe{flux}{flux!d'un champ de vecteur} de $G$ au travers de $\partial D$ est le nombre
\begin{equation}
	\int_{\partial D}\langle G(x), \nu(x)\rangle d\sigma(x).
\end{equation}

Cette intégrale est en général très compliquée à calculer parce qu'il faut trouver le champ de vecteur normal, puis une paramétrisation de la surface $\partial D$ et ensuite appliquer la méthode décrite au point \ref{secintsurfaciques}. 

Heureusement, il y a un théorème qui nous permet de calculer plus facilement : sans devoir trouver de vecteurs normaux.


Il n'est pas plus contraignant d'énoncer ce théorème dans le cadre d'une hypersurface de $\eR^n$, ce que nous faisons donc~:
\begin{theorem}[Formule de la divergence]
	Soit $D$ un ouvert borné de $\eR^n$ dont le bord est \og assez régulier par morceaux\fg{}, c'est-à-dire~:
	\begin{equation}
		\partial D = A_1 \cup \ldots A_p \cup N
	\end{equation} 
	où
	\begin{enumerate}
		\item $A_1, \ldots, A_p, N$ sont deux à deux disjoints,
		\item pour tout $i \leq p$, $A_i$ est un ouvert relatif d'une certaine variété $M_i$ de dimension $(n-1)$
		\item $\bar A_i \subset M_i$
		\item $N$ est un compact contenu dans une réunion finie de variétés de dimensions $(n-2)$.
	\end{enumerate}
	Supposons également qu'en chaque point de $A_1 \cup \ldots \cup A_p$ il existe un vecteur normal extérieur $\nu$.
	
	Si $G$ est un champ de vecteurs de classe $C^1$ sur $\bar D$ alors
	\begin{equation}
		\int_D \nabla\cdot G = \sum_{i=1}^p \int_{A_i} \scalprod{G}{\nu}.
	\end{equation}
	L'intégrale du membre de gauche est l'intégrale sur un ouvert d'une simple fonction.
\end{theorem}

\subsection{Formule de Green}
La formule de Green est un cas particulier du théorème de la divergence dans
le cas $n = 2$, légèrement reformulé~:
\begin{theorem}
	Soit $D \subset \eR^2$ ouvert borné tel que son bord est est la réunion finie d'un certain nombre de chemins de classe $C^1$ de Jordan réguliers.  Supposons qu'en chaque point de son bord, $D$ possède un vecteur normal unitaire extérieur $\nu$. Soient $P$ et $Q$ deux fonctions réelles de classe $C^1$ sur $\bar D$. Alors
  \begin{equation*}
    \iint_D (\partial_xQ - \partial_yP)dx\,dy = \oint_{\partial D}
    Pd x + Q d y
  \end{equation*}
  où chaque chemin $\gamma$ formant le bord de $D$ est orienté de
  sorte que $T \nu = \frac{\dot\gamma}{\norme{\dot\gamma}}$ où $T$
  représente la rotation d'angle $+\frac\pi2$.
\end{theorem}

Pour rappel, une chemin $\gamma\colon \mathopen[ 0 , 1 \mathclose]\to \eR^n$ est \defe{régulier}{régulier!chemin} si il est $C^1$ et si $\gamma(t)\neq 0$ pour tout $r$. Le chemin est de \defe{Jordan}{Jordan!chemin} si $\gamma(1)=\gamma(0)$ et si $\gamma\colon \mathopen[ a , b [\to \eR^n$ est injective.

\subsection{Formule de Stokes}
\label{secstokesusuel}
La formule de Stokes est la version classique, qui permet d'exprimer la circulation d'un champ de vecteur le long d'une courbe de $\eR^3$ comme le flux de son rotationnel à travers n'importe quel surface dont le bord est la courbe. La version présentée ici suppose que la surface peut se paramétrer en un seul morceau~:
\begin{theorem}
  Soit $F : W\subset \eR^2 \to \eR^3$ une paramétrisation (carte) d'une surface dans $\eR^3$, supposée de classe $C^2$. Soit $D$ un ouvert de $\eR^2$ vérifiant les hypothèses de la formule de Green, et tel que $\bar D \subset W$. Soit $G$ un champ de vecteurs de classe $C^1$ défini sur $F(\bar D)$, et soit $N$ le champ normal unitaire donné par la paramétrisation
  \begin{equation}		
	N = \frac{\pder F u \wedge \pder F v}{\norme{\pder F u \wedge \pder F v}}
\end{equation}
  alors
  \begin{equation}\label{EqStokesTho}
    \iint_{F(D)} \scalprod{\rot G}{N} d\sigma_F = \int_{F(\partial D)} G
  \end{equation}
  où les chemins formant le bord $\partial D$ sont orientés comme dans le théorème de Green.
\end{theorem}
Notons, juste pour avoir une bonne nouvelle de temps en temps, que 
\begin{equation}
	d\sigma_F=\left\| \frac{ \partial F }{ \partial u }\times\frac{ \partial F }{ \partial v }  \right\|dudv,
\end{equation}
mais cette norme apparaît exactement au dénominateur de $N$. Il ne faut donc pas la calculer parce qu'elle se simplifie.

Sous forme un peu plus physicienne\footnote{et surtout plus explicite.}, la formule \eqref{EqStokesTho} s'écrit
\begin{equation}
	\int_{F(D)}\langle \nabla\times G, N(x)\rangle\, d\sigma_F(x)=\int_{F(\gamma)}\langle G, T\rangle\, ds
\end{equation}
où $T$ est le vecteur unitaire tangent à $F(\gamma)$.

%///////////////////////////////////////////////////////////////////////////////////////////////////////////////////////////
\subsubsection{Quelle est la bonne orientation ?}
%///////////////////////////////////////////////////////////////////////////////////////////////////////////////////////////

Le signe du vecteur normal $N$ dépend du choix de l'ordre des coordonnées dans la carte. Supposons que je veuille paramétrer la surface $x^2+y^2=1$, $z=1$. Nous prenons naturellement comme carte le cercle $C$ de rayon $1$ dans $\eR^2$ et la carte
\begin{equation}
	F(r,\theta)=\begin{pmatrix}
		r\cos\theta	\\ 
		r\sin\theta	\\ 
		1	
	\end{pmatrix}.
\end{equation}
Mais nous aurions aussi pu mettre les coordonnées $r$ et $\theta$ dans l'autre ordre :
\begin{equation}
	\tilde F(\theta,r)=\begin{pmatrix}
		r\cos\theta	\\ 
		r\sin\theta	\\ 
		1	
	\end{pmatrix}.
\end{equation}
Les vecteurs normaux ne sont pas les même : la carte $F$ donnera $\partial_rF\times\partial_{\theta}F$, tandis que l'autre donnera $\partial_{\theta}\tilde F\times\partial_r\tilde F$. Le signe change !

Il faut savoir laquelle choisir. Le cercle $C\subset \eR^2$ a une orientation donnée par le théorème de Green. Nous choisissons l'ordre des coordonnées pour que $1_{\theta}$ et $1_{r}$ soient dans la même orientation que les vecteurs $\nu$ et $T$ tels que donnés par le théorème de Green, et tels que dessinés sur la figure \ref{LabelFigCercleTnu}.
\newcommand{\CaptionFigCercleTnu}{L'orientation sur le cercle. Si nous les prenons dans l'ordre, les vecteurs $(1_r,1_{\theta})$ ont la même orientation que celle donnée par les vecteurs $(\nu,T)$ donnés par la convention de Green.}
\input{Fig_CercleTnu.pstricks}

Plus généralement, nous choisissons l'ordre des coordonnées $u$ et $v$ pour que la base $(1_u,1_v)$ ait la même orientation que $(\nu,T)$ où $T$ a le sens convenu dans le théorème de Green.

%+++++++++++++++++++++++++++++++++++++++++++++++++++++++++++++++++++++++++++++++++++++++++++++++++++++++++++++++++++++++++++
					\section{Un petit extra}
%+++++++++++++++++++++++++++++++++++++++++++++++++++++++++++++++++++++++++++++++++++++++++++++++++++++++++++++++++++++++++++

Soit $f$ une fonction de $\eR$ dans $\eR$. Supposons que 
\begin{enumerate}

\item		\label{ItemExtrai}
$f(1)=1$,

\item		\label{ItemExtraii}
$f(x+y)=f(x)+f(y)$ pour tout réels $x$ et $y$.

\end{enumerate}
Nous pouvons montrer\footnote{et toi, tu le peux ?} que la seule fonction {\it continue} qui possède ces propriétés est la fonction identité $f(x)=x$ pour tout $x\in\eR$.

De la même manière, il est aisé de voir que les seules applications linéaires de $\Rn$ dans $\Rn$ sont de la forme 
\begin{equation}
	f(x)=Ax
\end{equation}
pour une constante réelle $A$. Une question naturelle qu'on peut alors se poser est la suivante: 
\begin{quote} 
	Est-il possible de définir une fonction non continue ayant les propriétés \ref{ItemExtrai} et \ref{ItemExtraii} ?
\end{quote}
En fait, il est possible de démontrer que si $E$ est un espace vectoriel de dimension finie, alors toute application linéaire $f:E\rightarrow  F$ (où $F$ est un espace vectoriel) sera continue. Ceci ne reste plus vrai si l'espace vectoriel $E$ est de dimension infinie. Donc une manière de trouver une réponse positive à la question posée plus haut, serait de voir $\Rn$ comme espace vectoriel de dimension infinie. Après un peu de réflexion, la réponse est venue à nous (merci à Nicolas et à Samuel). 



Si nous admettons l'\href{http://fr.wikipedia.org/wiki/Axiome_du_choix}{axiome du choix}, alors nous pouvons appliquer le théorème de Zorn et nous savons que tout espace vectoriel admet une base. En particulier, l'ensemble des réels vu comme espace vectoriel sur $\Qn$ admet une base, i.e. $\exists (e_i)_{i\in I}$  des éléments de $\Rn$ tels que tout réel s'écrit comme combinaison linéaire à coefficients rationnels  de ces $e_i$, i.e.
\begin{equation}
	\forall r \in \Rn, \exists (\lambda_i)_{i\in I} \text{ des éléments de } \Qn \text{ tels que  } r = \sum_{i\in I} \lambda_i e_i.
\end{equation}
Utilisons cette base pour définir une fonction $h$ de la manière suivante.
\begin{equation}
\forall i \in I, \mbox{ on définit } h(e_i) = \alpha_i
\end{equation}
 où les $\alpha_i$ doivent être bien choisis dans $\Rn$. Pour satisfaire la propriété \ref{ItemExtrai}, choisissons sans perte de généralité $e_1 = 1$ et $h(e_1) = 1$.  Ajoutons à cette propriété la linéarité en imposant que 
\begin{equation}
h(\sum \lambda_i e_i) = \sum \lambda_i \alpha_i.
\end{equation}
Les équations (1) et (2) nous permettent de voir que, moyennant le choix des $\alpha_i$, la fonction $h$ est bien définie sur $\Rn$ et linéaire. Il est clair que si nous prenons par exemple
$$\alpha_i=e_i\;\forall i \in I$$ 
nous obtenons que la fonction $h$ est en fait la fonction identité sur $\Rn$. Par contre, si nous définissons la fonction $h$ comme satisfaisant la propriété (2) et si nous choisissons les $\alpha_i$ dans (1) de la manière suivante 
\begin{equation}
	\begin{aligned}[]
		h(e_1)	&= e_2\\
		h(e_2)	&= e_1\\
		h(e_i)	&= e_i	&&\forall i\in I\setminus\{ 1,2 \}
	\end{aligned}
\end{equation}
alors la fonction ainsi obtenue est  linéaire et bien définie mais n'est plus l'identité. Donc nous avons trouvé une application linéaire de $\Rn$ dans $\Rn$ qui n'est pas continue.  

\begin{exercice}
 Trouver d'autres exemples d'applications linéaires non continues (pas nécessairement des transformations de $\Rn$). 		
\end{exercice}

%+++++++++++++++++++++++++++++++++++++++++++++++++++++++++++++++++++++++++++++++++++++++++++++++++++++++++++++++++++++++++++
\section{Preuves de certains résultats cités}
%+++++++++++++++++++++++++++++++++++++++++++++++++++++++++++++++++++++++++++++++++++++++++++++++++++++++++++++++++++++++++++

%---------------------------------------------------------------------------------------------------------------------------
\subsection{Preuve du critère d'équivalence}
%---------------------------------------------------------------------------------------------------------------------------
\label{PgPreuveCritEquiv}

Prouvons le critère d'équivalence mentionné à la page \pageref{PgCritEquiv}. La preuve est tirée de celle donnée dans \cite{TrenchRealAnalisys}.

\begin{enumerate}
	\item
		Le fait que la suite $a_n/b_n$ converge vers $\alpha$ signifie que tant sa limite supérieure que sa limite inférieure convergent vers $\alpha$. En particulier la suite $\frac{ a_n }{ b_n }$ est bornée vers le haut et vers le bas. À partir d'un certain rang $N$, il existe $M$ tel que 
		\begin{equation}
			\frac{ a_n }{ b_n }<M
		\end{equation}
		et il existe $m$ tel que
		\begin{equation}
			\frac{ a_n }{ b_n }>m.
		\end{equation}
		Nous avons donc $a_n<Mb_n$ et $a_n>mb_n$. La série de $(a_n)$ converge donc si et seulement si la série de $(b_n)$ converge.
	\item
		Si $\alpha=0$, cela signifie que pour tout $\epsilon$, il existe un rang tel que $\frac{ a_n }{ b_n }<\epsilon$, et donc tel que $a_n<\epsilon b_k$. La suite de $(a_i)$ converge donc dès que la suite de $(b_i)$ converge.
	\item
		Pour tout $M$, il existe un rang dans la suite à partir duquel on a $\frac{ a_i }{ b_i }>M$, et donc $a_k>Mb_k$. Si la série de $(b_k)$ diverge, la série de $(a_k)$ doit également diverger.
\end{enumerate}

%---------------------------------------------------------------------------------------------------------------------------
\subsection{Preuve du critère du quotient}
%---------------------------------------------------------------------------------------------------------------------------
\label{PgPreuveCritQuotient}

Cette preuve est tirée de \cite{KeislerElemCalculus}.

\begin{enumerate}
	\item
		Soit $b$ tel que $L<b<1$. À partir d'un certain rang $K$, on a $\left| \frac{ a_{i+1} }{ a_i } \right| <b$. En particulier,
		\begin{equation}
			| a_{K+1} |<b| a_K |,
		\end{equation}
		et pour $a_{K+2}$ nous avons
		\begin{equation}
			| a_{K+2} |<b| a_{K+1} |<b^2| a_K |.
		\end{equation}
		Au final,
		\begin{equation}
			| a_{K+n} |<b^n| a_K |.
		\end{equation}
		Étant donné que la série $\sum_{n\geq K}b^n$ converge (parce que $b<1$), la queue de suite $\sum_{i\geq K}a_i$ converge, et par conséquent la suite au complet converge.
	\item
		Si $L>1$, on a
		\begin{equation}
			| a_K |<| a_{K+1} |<| a_{K+2} |<\ldots
		\end{equation}
		Il est donc impossible que la suite $(a_i)$ converge vers zéro. La série ne peut donc pas converger.
	\item
		Par exemple la suite harmonique $a_n=\frac{1}{ n }$ vérifie $L=1$, mais la série ne converge pas. Par contre, la suite $a_n=\frac{ 1 }{ n^2 }$ vérifie aussi le critère avec $L=1$ tandis que la série $\sum_n\frac{1}{ n^2 }$ converge.
\end{enumerate}

%---------------------------------------------------------------------------------------------------------------------------
\subsection{Preuve du critère de la racine}
%---------------------------------------------------------------------------------------------------------------------------
\label{PgPreuveCritRacine}
La preuve est tirée de \cite{TrenchRealAnalisys}.

\begin{enumerate}
	\item
		Si $L<1$, il existe un $r\in \mathopen] 0 , 1 \mathclose[$ tel que $| a_n |^{1/n}<r$ pour les grands $n$. Dans ce cas, $| a_n |<r^{n}$, et la série converge absolument parce que la série $\sum_nr^n$ converge du fait que $r<1$.
	\item
		Si $L>1$, il existe un $r>1$ tel que $| a_n |^{1/n}>r>1$. Cela fait que $| a_n |$ prend des valeurs plus grandes que $n$ pour une infinité de termes. Le terme général $a_n$ ne peut donc pas être une suite convergente. Par conséquent la suite diverge au sens où elle ne converge pas.

\end{enumerate}



\chapter{Travaux personnels}
% This is part of Exercices et corrigés de CdI-1
% Copyright (c) 2011
%   Laurent Claessens
% See the file fdl-1.3.txt for copying conditions.

%TODO : lorsque GeomAnal aura été fourchetté, ces exercices devront aller là-bas.
% Ces trois exercices étaient du travail personnel de CdI1 à l'ULB en 2008-2009.
\Exo{TP20090001}
\Exo{TP20090002}
\Exo{TP20090003}






\chapter{Convergences et espaces fonctionnels}
	% This is part of the Exercices et corrigés de CdI-2.
% Copyright (C) 2008, 2009, 2012
%   Laurent Claessens
% See the file fdl-1.3.txt for copying conditions.


%+++++++++++++++++++++++++++++++++++++++++++++++++++++++++++++++++++++++++++++++++++++++++++++++++++++++++++++++++++++++++++
					\section[Intégrales de fonctions et domaines non bornées]{Intégrales de fonctions non bornées sur des domaines non bornés}
%+++++++++++++++++++++++++++++++++++++++++++++++++++++++++++++++++++++++++++++++++++++++++++++++++++++++++++++++++++++++++++

%---------------------------------------------------------------------------------------------------------------------------
					\subsection[Fonctions et ensembles non bornés]{Intégrales de fonctions non bornées sur des ensembles non bornés}
%---------------------------------------------------------------------------------------------------------------------------

Soit $f\colon \eR^n\to \overline{ \eR }$, une fonction positive. On dit qu'elle est \defe{intégrable}{intégrable (fonction)} sur $E\subset\eR^n$ si
\begin{enumerate}
    \item $\forall r>0$, la fonction $f_r(x)=f(x)\mtu_{f<r}$ est intégrable sur $E_r$;
\item la limite $\lim_{r\to\infty}\int_{E_r}f_r$ est finie.
\end{enumerate}
Dans ce cas, on pose 
\begin{equation}
	\int_Ef=\lim_{r\to\infty}\int_{E_r}f_r.
\end{equation}

\begin{theorem}[Page I.38]		\label{ThoFnTestIntnnBorn}
Soit $E$ mesurable dans $\eR^n$ et $f\colon E\to \overline{ \eR }$. Si $f$ est mesurable et si il existe $g\colon E\to \overline{ \eR }$ intégrable sur $E$ telle que $| f(x) |\leq g(x)$ pour tout $x\in E$, alors $f$ est intégrable sur~$E$.

Réciproquement, si $f$ est intégrable sur $E$, alors $f$ est mesurable.
\end{theorem}

%---------------------------------------------------------------------------------------------------------------------------
					\subsection{Passage à la limite sous le signe intégral}
%---------------------------------------------------------------------------------------------------------------------------

Un autre résultat très important pour l'étude de l'intégrabilité est le théorème de la \defe{convergence dominée de Lebesgue}{}:
\begin{theorem}
	Soit $E\subset \eR^n$ un ensemble mesurable et $\{ f_k \}$, une suite de fonctions intégrables sur $E$ qui converge simplement vers une fonction $f\colon E\to \overline{ \eR }$. Supposons qu'il existe une fonction $g$ intégrable sur $E$ telle que pour tout $k$,
\begin{equation}
	| f(x) |\leq g(x)
\end{equation}
pour tout $x\in E$. Alors $f$ est intégrable sur $E$ et 
\begin{equation}
	\int_Ef=\lim_{k\to\infty}\int_Ef_k.
\end{equation}
\end{theorem}

%---------------------------------------------------------------------------------------------------------------------------
					\subsection{Théorème de Fubini et changement de variables}
%---------------------------------------------------------------------------------------------------------------------------

\begin{theorem}[Fubini]\index{théorème!Fubini}\index{Fubini!théorème}		\label{ThoFubini}
Soit $(x,t)\mapsto f(x,y)\in\bar \eR$ une fonction intégrable sur $B_n\times B_m\subset\eR^{n+m}$ où $B_n$ et $B_m$ sont des ensembles mesurables de $\eR^n$ et $\eR^m$. Alors :
\begin{enumerate}
\item pour tout $x\in B_n$, sauf éventuellement en les points d'un ensemble $G\subset B_n$ de mesure nulle, la fonction $y\in B_m\mapsto f(x,y)\in\bar\eR$ est intégrable sur $B_m$
\item
la fonction
\begin{equation}
	x\in B_n\setminus G\mapsto\int_{B_m}f(x,y)dy\in\eR
\end{equation}
est intégrable sur $B_n\setminus G$

\item 
On a
\begin{equation}
	\int_{B_n\times B_m}f(x,y)dxdy=\int_{B_n}\left( \int_{B_m}f(x,y)dy \right)dx.
\end{equation}

\end{enumerate}
\end{theorem}

Notons en particulier que si $f(x,y)=\varphi(x)\phi(y)$, alors $\int_{B_m}\varphi(y)dy$ est une constante qui peut sortir de l'intégrale sur $B_n$, et donc
\begin{equation}		\label{EqFubiniFactori}
	\int_{B_n\times B_m}\varphi(x)\phi(y)dxdy=\int_{B_n}\varphi(x)dx\int_{B_m}\phi(y)dy.
\end{equation}

%---------------------------------------------------------------------------------------------------------------------------
					\subsection{Intégrale en dimension un}
%---------------------------------------------------------------------------------------------------------------------------

\begin{proposition}[Critère de comparaison]
Soit $f$ mesurable sur $]a,\infty[$ et bornée sur tout $]a,b]$, et supposons qu'il existe un $X_0\geq a$, tel que sur $]X_0,\infty[$,
\begin{equation}
	| f(x) |\leq g(x)
\end{equation}
où $g(x)$ est intégrable. Alors $f(x)$ est intégrable sur $]a,\infty[$.
\end{proposition}

\begin{corollary}[Critère d'équivalence]
Soient $f$ et $g$ des fonctions mesurables et positives ou nulles sur $]a,\infty[$, bornées sur tout $]a,b]$, telles que 
\begin{equation}
	\lim_{x\to\infty}\frac{ f(x) }{ g(x) }=L
\end{equation}
existe dans $\bar\eR$.
\begin{enumerate}
\item Si $L\neq\infty$ et $\int_{a}^{\infty}g(x)$ existe, alors $\int_a^{\infty}f(x)dx$ existe,
\item Si $L\neq 0$ et si $\int_a^{\infty}f(x)dx$ existe, alors $\int_a^{\infty}g(x)dx$ existe,
\end{enumerate}
\end{corollary}

\begin{corollary}[Critère des fonctions test]			\label{CorCritFonsTest}
Soit $f(x)$ une fonction mesurable et positive ou nulle sur $]a,\infty[$ et bornée pour tout $]a,b]$. Nous posons
\begin{equation}
	L(\alpha)=\lim_{x\to\infty}x^{\alpha}f(x),
\end{equation}
et nous supposons qu'elle existe.
\begin{enumerate}
\item Si il existe $\alpha>1$ tel que $L(\alpha)\neq\infty$, alors $\int_a^{\infty}f(x)dx$ existe,
\item Si il existe $\alpha\leq1$ et $L(\alpha)\neq 0$, alors $\int_a^{\infty}f(x)dx$ n'existe pas.
\end{enumerate}
\end{corollary}

\begin{corollary}		\label{CorAlphaLCasInteabf}
	Soit $f\colon ]a,b]\to \eR$ une fonction mesurable, positive ou nulle, et bornée sur $[a+\epsilon,b]$ $\forall\epsilon>0$. Si $\lim_{x\to a}(x-a)^{\alpha}f(x)=L$ existe, alors
	\begin{enumerate}
		\item Si $\alpha<1$ et $L\neq\infty$, alors $\int_a^bf(x)dx$ existe,
		\item Si $\alpha\geq 1$ et $L\neq 0$, alors $\int_a^bf(x)dx$ n'existe pas.
	\end{enumerate}
\end{corollary}

%---------------------------------------------------------------------------------------------------------------------------
					\subsection{Intégrales convergentes}
%---------------------------------------------------------------------------------------------------------------------------

\begin{definition}

Soit $f$, une fonction mesurable sur $[a,\infty[$, bornée sur tout intervalle $[a,b]$. On dit que l'intégrale
\begin{equation}
	\int_a^{\infty}f(x)dx
\end{equation}
\defe{converge}{intégrale!convergente} si la limite
\begin{equation}		\label{EqDEfConvergeZeroInftX}
	\lim_{X\to\infty}\int_a^{X}f
\end{equation}
existe et est finie.
\end{definition}


%+++++++++++++++++++++++++++++++++++++++++++++++++++++++++++++++++++++++++++++++++++++++++++++++++++++++++++++++++++++++++++
					\section{Fonctions définies par des intégrales et régularisation}
%+++++++++++++++++++++++++++++++++++++++++++++++++++++++++++++++++++++++++++++++++++++++++++++++++++++++++++++++++++++++++++
%+++++++++++++++++++++++++++++++++++++++++++++++++++++++++++++++++++++++++++++++++++++++++++++++++++++++++++++++++++++++++++
%					\section[Fonctions définies par des intégrales]{Fonctions définies par des intégrales et régularisation}

%+++++++++++++++++++++++++++++++++++++++++++++++++++++++++++++++++++++++++++++++++++++++++++++++++++++++++++++++++++++++++++

Supposons $A\subset\eR^m$  et $B\subset\eR^n$ compact. Nous considérons $f(x,t)\colon A\times B\to \eR$ une fonction bornée sur $A\times B$ et intégrable par rapport à $t$ pour tout $x\in A$. Soit $F\colon A\to \eR$ définie par
\begin{equation}
	F(x)=\int_Bf(x,t)dt.
\end{equation}

%---------------------------------------------------------------------------------------------------------------------------
					\subsection{Fonction définies par une intégrale sur un compact}
%---------------------------------------------------------------------------------------------------------------------------

\begin{proposition}		\label{PropDerrSSIntegraleDSD}
Supposons $A\subset\eR^m$ ouvert et $B\subset\eR^n$ compact. Si pour un $i\in\{ i,\ldots,n \}$, la dérivée partielle $\frac{ \partial f }{ \partial x_i }$ existe dans $A\times B$ et est continue, alors $\frac{ \partial F }{ \partial x_i }$ existe dans $A$, est continue et
\begin{equation}
	\frac{ \partial F }{ \partial x_i }=\int_B\frac{ \partial f }{ \partial x_i }(x,t)dt,
\end{equation}
l'égalité signifie que l'on peut \og dériver sous le signe intégral\fg.
\end{proposition}

%---------------------------------------------------------------------------------------------------------------------------
					\subsection{Intégrale sur un segment variable}
%---------------------------------------------------------------------------------------------------------------------------


\begin{proposition}		\label{PropDerrFnAvecBornesFonctions}
Soit $f(x,t)$ une fonction continue sur $[\alpha,\beta]\times[a,b]$, telle que $\frac{ \partial f }{ \partial x }$ existe et soit continue sur $]\alpha,\beta[\times[a,b]$. Soient $\varphi(x)$ et $\psi(x)$, des fonctions continues de $[\alpha,\beta]$ dans $\eR$ et admettant une dérivée continue sur $]\alpha,\beta [$. Alors la fonction
\begin{equation}
	F(x)=\int_{\varphi(x)}^{\psi(x)}f(x,t)dt
\end{equation}
admet une dérivée continue sur $]\alpha,\beta[$ et
\begin{equation}	\label{EqFormDerrFnAvecBorneNInt}
	\frac{ dF }{ dx }=\int_{\varphi(x)}^{\psi(x)}\frac{ \partial f }{ \partial x }(x,t)dt+f\big( x,\psi(x) \big)\cdot\frac{ d\psi }{ dx }- f\big( x,\varphi(x) \big)\cdot\frac{ d\varphi }{ dx }.
\end{equation}
\end{proposition}

%---------------------------------------------------------------------------------------------------------------------------
					\subsection{Intégrales convergentes}
%---------------------------------------------------------------------------------------------------------------------------

\begin{theorem}		\label{ThoInDerrtCvUnifFContinue}
	Uniforme convergence, continuité et dérivation.
	\begin{enumerate}
		\item 
		Soit $f(x,t)$ continue sur $[\alpha,\beta]\times[a,\alpha[$ et $F(x)=\int_a^{\infty}f(x,t)dt$; cette intégrale étant supposée uniformément convergente. Alors $F(x)$ est continue.
		\item
		Supposons $f$ continue et sa dérivée partielle $\frac{ \partial f }{ \partial x }$ continue sur $[\alpha,\beta]\times[a,\alpha[$. Supposons que $F(x)=\int_a^{\infty}f(x,t)dt$ converge et que $\int_a^{\infty}\frac{ \partial f }{ \partial x }dt$ converge uniformément. Alors $F$ est $C^1$ sur $[\alpha,\beta]$ et 
		\begin{equation}
			\frac{ dF }{ dx }=\int_a^{\infty}\frac{ \partial f }{ \partial x }dt.
		\end{equation}
	\end{enumerate}
\end{theorem}
%---------------------------------------------------------------------------------------------------------------------------
					\subsection{Critères de convergence uniforme}
%---------------------------------------------------------------------------------------------------------------------------


Affin de tester l'uniforme convergence d'une intégrale, nous avons le \defe{critère de Weierstrass}{Critère!Weierstrass}:
\begin{theorem}		\label{ThoCritWeiIntUnifCv}
Soit $f(x,t)\colon [\alpha,\beta]\times[a,\infty[ \to \eR$, une fonction dont la restriction à toute demi-droite $x=cst$ est mesurable. Si $| f(x,t) |< \varphi(t)$ et $\int_a^{\infty}\varphi(t)dt$ existe, alors l'intégrale
\begin{equation}
	\int_0^{\infty}f(x,t)dt
\end{equation}
est uniformément convergente.
\end{theorem}


Le théorème suivant est le \defe{critère d'Abel}{Critère!Abel pour intégrales} :
\begin{theorem}		\label{ThoAbelIntUnif}
	Supposons que $f(x,t)=\varphi(x,t)\psi(x,t)$ où $\varphi$ et $\psi$ sont bornée et intégrables en $t$ au sens de Riemann sur tout compact $[a,b]$, $b\geq a$. Supposons que :
	\begin{enumerate}
		\item $| \int_a^{T}\varphi(x,t)dt |\leq M$ où $M$ est indépendant de $T$ et de $x$,
		\item $\psi(x,t)\geq 0$,
		\item pour tout $x\in[\alpha,\beta]$, $\psi(x,t)$ est une fonction décroissante de $t$,
		\item les fonctions $x\mapsto \psi(x,t)$ convergent uniformément vers $0$ lorsque $t\to\infty$.
	\end{enumerate}
	Alors l'intégrale
	\begin{equation}
		\int_a^{\infty}f(x,t)dt
	\end{equation}
	est uniformément convergente.
\end{theorem}


\chapter{Foire aux questions}
	% This is part of the Exercices et corrigés de CdI-2.
% Copyright (C) 2008-2009, 2012-2013
%   Laurent Claessens
% See the file fdl-1.3.txt for copying conditions.

Le but de ce chapitre est de répondre de façon plus détaillé aux questions qui ont été souvent posées durant les séances d'exercices.

%+++++++++++++++++++++++++++++++++++++++++++++++++++++++++++++++++++++++++++++++++++++++++++++++++++++++++++++++++++++++++++
                    \section{Le coup du compact}
%+++++++++++++++++++++++++++++++++++++++++++++++++++++++++++++++++++++++++++++++++++++++++++++++++++++++++++++++++++++++++++

Dans les sections qui suivent, nous verrons des fonctions définies par toute une série de processus de limite (suites, séries, intégrales). Une des questions centrales est de savoir si la fonction limite est continue, dérivable, intégrale, etc. étant donné que les fonctions sont continues.

Pour cela, nous inventons le concept de \emph{convergence uniforme}. Si la limite (série, intégrale) est uniforme, alors la fonction limite sera continue. Il arrive qu'une limite ne soit pas uniforme sur un intervalle ouvert $]0,1]$, et que nous voulions quand même prouver la continuité sur cet intervalle. C'est à cela que sert la notion de convergence uniforme \emph{sur tout compact}. En effet, la notion de continuité est une notion locale : savoir ce qu'il se passe dans un petit voisinage autour de $x$ est suffisant pour savoir la continuité en $x$ (idem pour sa dérivée).

Si nous avons uniforme convergence sur tout compact de $]0,1]$, mais pas uniforme convergence sur cet intervalle, la limite sera quand même continue sur $\mathopen] 0 , 1 \mathclose]$. En effet, si $x\in]0,1]$, il existe un ouvert autour de $x$ contenu dans un compact contenu dans $]0,1]$. L'uniforme convergence sur ce compact suffit à prouver la continuité en $x$.

Déduire la continuité sur un ouvert à partir de l'uniforme convergence sur tout compact de l'ouvert est appelé faire le \defe{coup du compact}{Compact!le coup du}.

%+++++++++++++++++++++++++++++++++++++++++++++++++++++++++++++++++++++++++++++++++++++++++++++++++++++++++++++++++++++++++++
\section{Vitesses de $x^{\alpha}$, de l'exponentielle et du logarithme}
%+++++++++++++++++++++++++++++++++++++++++++++++++++++++++++++++++++++++++++++++++++++++++++++++++++++++++++++++++++++++++++

Une astuce\label{PgAstuceLnPoly} usuelle qu'il faut savoir est que $\forall \alpha>0$, $\exists N$ tel que $n>N\Rightarrow \ln(n)\leq n^{\alpha}$. En effet, nous avons
\begin{equation}
    \lim_{x\to\infty} \frac{ x^{\alpha} }{ \ln(x) }=\lim_{x\to\infty} \frac{ \alpha x^{\alpha-1} }{ 1/x }=\lim_{x\to\infty} \alpha x^{\alpha}=\infty
\end{equation}
quand $\alpha>0$. Cela tient également lorsque nous considérons $\ln(x)^p$ au lieu de $\ln(x)$. De cela, nous disons que le logarithme croit moins vite que n'importe quel polynôme. 

Dans le même ordre d'idées, l'exponentielle croit plus vite que tout polynôme, et plus vite que que logarithme :
\begin{equation}        \label{EqExpDecrtPlusVite}
    \lim_{t\to\infty} e^{-t}(\ln t)^{n}t^{\alpha}=0
\end{equation}
pour tout $n$ et pour tout $\alpha$.

%+++++++++++++++++++++++++++++++++++++++++++++++++++++++++++++++++++++++++++++++++++++++++++++++++++++++++++++++++++++++++++
\section{Remarque : Abel et \texorpdfstring{$\sin(xt)$}{sin(xt)}}
%+++++++++++++++++++++++++++++++++++++++++++++++++++++++++++++++++++++++++++++++++++++++++++++++++++++++++++++++++++++++++++

Étant donné que la fonction sinus est bornée, il est tentant de l'utiliser comme $\varphi$ dans le critère d'Abel (théorème \ref{ThoAbelIntUnif}). Hélas,
\begin{equation}
    \int_0^T\sin(xt)=-\frac{ 1 }{ x }\big( \cos(xT)-\cos(x) \big),
\end{equation}
qui n'est pas bornée pour tout $x$ ! Poser $\varphi(x,t)=\sin(xt)$ \emph{ne fonctionne pas} pour assurer la convergence uniforme sur un intervalle qui contient des $x$ arbitrairement proches de $0$. Le critère d'Abel avec $\varphi(x,t)=\sin(xt)$ ne permet que de conclure à l'uniforme convergence \emph{sur tout compact} ne contenant pas $0$. Cela est toutefois souvent suffisant pour étudier la continuité ou la dérivabilité en se servant du fameux coup du compact.

%+++++++++++++++++++++++++++++++++++++++++++++++++++++++++++++++++++++++++++++++++++++++++++++++++++++++++++++++++++++++++++
                    \section{Que faire avec \texorpdfstring{$f(z)dz=g(t)dt$}{fzdz} ?}
%+++++++++++++++++++++++++++++++++++++++++++++++++++++++++++++++++++++++++++++++++++++++++++++++++++++++++++++++++++++++++++
\label{SecFairedzdt}

Dans de nombreux exercices d'équations différentielles, nous tombons sur $u'=f(t)$, et nous faisons formellement
\begin{equation}
    \begin{aligned}[]
        \frac{ du }{ dt }&=f(t) &\Rightarrow    &&du=f(t)dt,
    \end{aligned}
\end{equation}
et ensuite, il y a la formule un peu magique
\begin{equation}
    u-u_0=\int_{t_0}^tf(t)dt.
\end{equation}
Voyons ce qu'il en est. Tout d'abord, il faut comprendre ce que signifie la formule
\begin{equation}        \label{EqDiffAstufzdz}
    f(z)dz=g(t)dt.
\end{equation}
Il s'agit d'une égalité entre deux formes différentielles sur $\eR$ où $z$ est une fonction de $t$.  Étant donné que $z$ est une fonction de $t$, il faut voir $dz$ comme la différentielle de cette fonction. La différentielle d'une fonction à une variable est donné par la dérivée :
\begin{equation}
    dz_t=z'(t)dt
\end{equation}
Écrire l'équation \eqref{EqDiffAstufzdz} pour chaque $t$ revient donc à écrire
\begin{equation}
    f\big( z(t) \big)z'(t)dt=g(t)dt
\end{equation}
Cela est une égalité entre deux formes différentielles. Nous avons donc égalité entre les intégrales des formes sur un chemin. Prenons un chemin tout simple de $t_0$ vers $t$ :
\begin{equation}
    \int_{t_0}^tf\big( z(t) \big)z'(t)dt=\int_{t_0}^tg(t)dt.
\end{equation}
Dans le premier membre, nous faisons un changement de variable $\xi=z(t)$, $d\xi=z'(t)dt$, et nous obtenons
\begin{equation}        \label{EqIntDiffAstuztz}
    \int_{z_0}^{z(t)}f(\xi)d\xi=\int_{t_0}^tg(t)dt.
\end{equation}
où nous avons remplacé la constante $z(t_0)$ par $z_0$ dans la borne d'intégration.  Si $F$ est une primitive de $f$ et $G$ une primitive de $g$, nous avons
\begin{equation}
    F(z)-F(z_0)=G(t)-G(t_0).
\end{equation}
Si aucun problème de Cauchy n'est donné, les constantes $F(z_0)$ et $G(t_0)$ sont mises en une seule et nous écrivons la solution
\begin{equation}
    F\big( z(t) \big)=G(t)+C,
\end{equation}
qui est une équation implicite pour $z(t)$. 

Nous trouvons assez souvent le cas simple
\begin{equation}    \label{EqAstfzdzdt}
    f(z)dz=dt.
\end{equation}
En remplaçant $g(t)=1$ dans \eqref{EqIntDiffAstuztz}, nous trouvons la fameuse
\begin{equation}        \label{Eqttzint}
    t-t_0=\int_{z_0}^zf(z)dz,
\end{equation}
dans laquelle il y a un abus de notation terrible entre le $z$ de la borne (que les étudiant(e)s oublient souvent) et la variable d'intégration $z$ !!

Le passage de \eqref{EqAstfzdzdt} à \eqref{Eqttzint} sera très souvent utilisé dans le cours de mécanique par exemple.


%+++++++++++++++++++++++++++++++++++++++++++++++++++++++++++++++++++++++++++++++++++++++++++++++++++++++++++++++++++++++++++
                    \section{Pourquoi la variation des constantes fonctionne toujours ?}
%+++++++++++++++++++++++++++++++++++++++++++++++++++++++++++++++++++++++++++++++++++++++++++++++++++++++++++++++++++++++++++

Prenons une équation non homogène 
\begin{equation}        \label{EqAstNNHomo}
    z'(t)=f(t)z(t)+g(t),
\end{equation}
et supposons avoir une solution de l'homogène associée sous la forme $z_H(t)=Ch(t)$. Le coup de la variation des constates consiste à essayer une solution pour l'équation non homogène sous la forme\footnote{Je ne sais plus qui a eu l'idée de changer le nom de la constante de $C$ vers $K$ au moment de la transformer en fonction, mais c'est une bonne idée.}
\begin{equation}
    z(t)=K(t)h(t).
\end{equation}
Nous injectons cette solution dans l'équation de départ en utilisant le fait que $z'(t)=K'(t)h(t)+K(t)h'(t)$ :
\begin{equation}
    K'(t)h(t)+K(t)h'(t)=f(t)K(t)h(t)+g(t).
\end{equation}
Le terme $K(t)h'(t)$ se récrit en utilisant la propriété de définition de $h$, c'est à dire que $h'(t)=f(t)h(t)$. Nous voyons que les termes ne contenant pas de $K'$ se simplifient; il reste
\begin{equation}
    K'h=g.
\end{equation}
Cette équation a comme solution
\begin{equation}
    K=\int \frac{ f }{ h }+C.
\end{equation}
J'insiste sur la constante d'intégration ! En réalité, celles et ceux qui auront compris l'équation \eqref{Eqttzint} sauront que $K$ est donné par
\begin{equation}
    K(t)=\int_{\xi_0}^{t}\frac{ f(\xi) }{ g(\xi) }d\xi
\end{equation}
où $\xi_0$ joue le rôle de la constante d'intégration.

Quoi qu'il en soit, la solution générale de l'équation non homogène est
\begin{equation}        \label{EqSolVarCosntCool}
    z(t)=K(t)h(t)=\left( \int\frac{ g }{ h }+C \right)h.
\end{equation}
Cette solution comprend deux termes : $Ch$ qui est solution de l'homogène, et $\left( \int \frac{ g }{ h } \right)h$ qui est une particulière de l'équation non homogène.

Quelque conclusions :

\begin{enumerate}
\item
Si vous avez encore du $K$ (et pas que du $K'$) dans votre équation qui donne $K$, c'est que vous n'être pas dans le cadre d'une équation de type \eqref{EqAstNNHomo}. Le plus souvent, c'est que vous avez fait une faute de calcul quelque part.

\item
La méthode des variations des constantes n'est pas en contradiction avec le principe de \og SGEH+SPENH\fg. En effet, la SGEP et la SPENH sont toutes deux dans la solution \eqref{EqSolVarCosntCool}.

\item
La variation des constantes peut être vue comme une façon cool de trouver une solution particulière de l'équation non homogène.

\item
La simplification ne se fait que après avoir remplacé $Kh'$ par $Kfh$, c'est à dire après avoir utilisé le fait que $z_H$ est solution de l'homogène. Sinon, la simplification n'est pas du tout évidente a priori. Il se peut même que, visuellement, les termes $Kh'$ et $Kfh$ ne se ressemblent pas du tout. C'est ce qui arrive par exemple à l'exercice \ref{exoreserve0004}\ref{ItemCII107}, pour arriver à l'équation \eqref{EqFracII107exoVVprb}.

\end{enumerate}

%+++++++++++++++++++++++++++++++++++++++++++++++++++++++++++++++++++++++++++++++++++++++++++++++++++++++++++++++++++++++++++
\section{Trucs et astuces de calculs d'intégrales}
%+++++++++++++++++++++++++++++++++++++++++++++++++++++++++++++++++++++++++++++++++++++++++++++++++++++++++++++++++++++++++++

Afin d'alléger le texte de calculs parfois un peu longs, nous regroupons ici les intégrales à une variable que nous devons utiliser dans les autres parties du cours.

\begin{enumerate}
	\item	\label{ItemIntegrali}
		L'intégrale
		\begin{equation}
			\boxed{I=\int x\ln(x)dx=\frac{ x^2 }{2}\big( \ln(x)-\frac{ 1 }{2} \big)}
		\end{equation}
		se fait par partie en posant
		\begin{equation}
			\begin{aligned}[]
				u&=\ln(x),		& dv&=x\,dx\\
				du&=\frac{1}{ x }\,dx,	& v&=\frac{ x^2 }{2},
			\end{aligned}
		\end{equation}
		et ensuite
		\begin{equation}
			I=\ln(x)\frac{ x^2 }{2}-\int\frac{ x }{2}=\frac{ x^2 }{2}\big( \ln(x)-\frac{ 1 }{2} \big).
		\end{equation}
		
	\item	
		L'intégrale
		\begin{equation}
			\boxed{I=\int x\ln(x^2)dx=x^2\ln(x)-\frac{ x^2 }{2}.}
		\end{equation}
		En utilisant le fait que $\ln(u^2)=2\ln(u)$, nous retombons sur une intégrale du type \ref{ItemIntegrali} :
		\begin{equation}
			I=x^2\ln(x)-\frac{ x^2 }{2}.
		\end{equation}
	\item
		L'intégrale
		\begin{equation}		\label{EqTrucIntxlnxsqpun}
			\boxed{I=\int x\ln(1+x^2)dx=\frac{ 1 }{2}\ln(x^2+1)(x^2+1)-x^2-\frac{ 1 }{2}}
		\end{equation}
		se traite en posant $v=1+x^2$ de telle sorte à avoir $dx=\frac{ dv }{ 2x }$ et donc
		\begin{equation}
			I=\frac{ 1 }{2}\ln(x^2+1)(x^2+1)-x^2-\frac{ 1 }{2}.
		\end{equation}
		
	\item
		L'intégrale
		\begin{equation}
			I=\int \cos(\theta)\sin(\theta)\ln\left( 1+\frac{1}{ \cos^2(\theta) } \right)\,d\theta
		\end{equation}
		demande le changement de variable $u=\cos(\theta)$, $d\theta=-\frac{ du }{ \sin(\theta) }$. Nous tombons sur l'intégrale
		\begin{equation}
			I=-\int u\ln\left( \frac{ 1+u^2 }{ u^2 } \right)=-\int u\ln(1+u^2)+\int u\ln(u^2),
		\end{equation}
		qui sont deux intégrales déjà faites. Nous trouvons
		\begin{equation}
			I=-\frac{ 1 }{2}\ln\left( \frac{ \sin^2(\theta)-1 }{ \sin^2(\theta)-2 } \right)\sin^2(\theta)-\ln\big( \sin^2(\theta)-2 \big)+\frac{ 1 }{2}\ln\big( \sin^2(\theta)-1 \big)
		\end{equation}
	
	\item
		L'intégrale
		\begin{equation}
			\boxed{\int \frac{ r^3 }{ 1+r^2 }dr=\frac{ r^2 }{2}-\frac{ 1 }{2}\ln(r^2+1).}
		\end{equation}
		commence par faire la division euclidienne de $r^3$ par $r^2+1$; ce que nous trouvons est $r^3=(r^2+1)r-r$. Il reste à intégrer
		\begin{equation}
			\int \frac{ r^3 }{ 1+r^2 }dr=\int r\,dr-\int\frac{ r }{ 1+r^2 }dr.
		\end{equation}
		La fonction dans la seconde intégrale est $\frac{ r }{ 1+r^2 }=\frac{ 1 }{2}\frac{ f'(r) }{ f(r) }$ où $f(r)=1+r^2$, et donc $\int \frac{ r }{ 1+r^2 }=\frac{ 1 }{2}\ln(1+r^2)$. Au final,
		\begin{equation}
			I=\frac{ 1 }{2}r^2-\frac{ 1 }{2}\ln(r^2+1).
		\end{equation}


	\item	
		L'intégrale
		\begin{equation}	\label{EqTrucIntsxcxdx}
			\boxed{I=\int \cos(\theta)\sin(\theta)d\theta=\frac{ \sin^2(\theta) }{ 2 }}
		\end{equation}
		se traite par le changement de variable $u=\sin(\theta)$, $du=\cos(\theta)d\theta$, et donc
		\begin{equation}
			\int\cos(\theta)\sin(\theta)d\theta=\int udu=\frac{ u^2 }{2}=\frac{ \sin^2(\theta) }{ 2 }.
		\end{equation}
	\item
		L'intégrale
		\begin{equation}	\label{EqTrucsIntsqrtAplusu}
			\boxed{\int\sqrt{1+x^2}dx=\frac{ x }{2}\sqrt{1+x^2}+\frac{ 1 }{2}\arcsinh(x)}
		\end{equation}
		s'obtient en effectuant le changement de variable $u=\sinh(\xi)$.

    \item
        L'intégrale
        \begin{equation}        \label{EqTrucIntcossqsinsq}
            \boxed{ \int\cos^2(x)\sin^2(x)dx=\frac{ x }{ 8 }-\frac{ \sin(4x) }{ 32 } }
        \end{equation}
        s'obtient à coups de formules de trigonométrie. D'abord, $\sin(t)\cos(t)=\frac{ 1 }{2}\sin^2(2t)$ fait en sorte que la fonction à intégrer devient 
        \begin{equation}
            f(x)=\frac{1}{ 4 }\sin^2(x).
        \end{equation}
        Ensuite nous utilisons le fait que $\sin^2(t)=(1-\cos(2t))/2$ pour transformer la formule à intégrer en
        \begin{equation}
            f(x)=\frac{ 1-\cos(4x) }{ 8 }.
        \end{equation}
        Cela s'intègre facilement en posant $u=4x$, et le résultat est
        \begin{equation}
            \int f(x)dx=\frac{ x }{ 8 }-\frac{ \sin(4x) }{ 32 }.
        \end{equation}

\end{enumerate}

\chapter{Appendices}
% This is part of Exercices et corrigés de CdI-1
% Copyright (c) 2011,2012
%   Laurent Claessens
% See the file fdl-1.3.txt for copying conditions.

%+++++++++++++++++++++++++++++++++++++++++++++++++++++++++++++++++++++++++++++++++++++++++++++++++++++++++++++++++++++++++++
\section{Quelque fautes usuelles}
%+++++++++++++++++++++++++++++++++++++++++++++++++++++++++++++++++++++++++++++++++++++++++++++++++++++++++++++++++++++++++++

Cette section a pour but de vous montrer un certain nombre de fautes que l'on rencontre régulièrement sur les feuilles d'examen. Les affirmations contenues dans cette sections sont donc fausses.

\begin{exercice}

    Soit la fonction $f : \mathbb{R}^2\to \mathbb{R}$,
    \begin{equation}
        f(x,y)=\begin{cases}
            \frac{ \sin(x^2+y^2+y^3) }{ x^2+y^2 }	&	\text{si $(x,y)\neq (0,0)$}\\
                1				&	 \text{si $(x,y)=(0,0)$}
        \end{cases}
    \end{equation}
    Étudier la continuité et la différentiabilité de $f$
\end{exercice}

Voici différente fautes qui ont été commises pour cette question.
\begin{enumerate}

\item
Nous avons $\frac{ \sin(y^3) }{ x^2+y^2 }\leq\sin(y^3)$, or $\lim_{y\to 0}\sin(y^3)=0$, donc $\lim_{(x,y)\to(0,0)}\frac{ \sin(y^3) }{ x^2+y^2 }=0 $. La fonction est donc continue en $(0,0)$.

\item
Étudions la limite lorsque $(x,y)$ tend vers $(0,0)$ :
\begin{equation}
	\begin{aligned}[]
		\lim_{ \begin{subarray}{l} (x,y)\to(0,0)\\x=\sqrt{y^3-y^2} \end{subarray} }\frac{ \sin(y^3) }{ x^2+y^2 }=\lim_{y\to 0} \frac{ \sin(y^3) }{ y^3-y^2+y^2 }=\lim_{y\to 0} \frac{ \sin(y^3) }{ y^3 }=1,
	\end{aligned}
\end{equation}
donc si la limite existe, elle est égale à un (parce qu'il existe un chemin sur lequel la limite vaut $1$).

\item
Affin de vérifier la différentiabilité en $(0,0)$, nous étudions l'existence de la limite suivante :
\begin{equation}
	\begin{aligned}[]
		\lim_{t\to 0} \frac{ f(0,0)-f\big( (0,0)+t(u,v) \big) }{ t }=\lim_{t\to 0} \frac{ \sin(t^3v^3) }{ t^2(v^2+u^2)t }=\lim_{t\to 0} \frac{ u^3 }{ (v^2+u^2) },
	\end{aligned}
\end{equation}
donc la limite existe et la fonction est différentiable en $(0,0)$.

\item
Petit calcul :
\begin{equation}
	\lim_{(x,y)\to (0,0)} \frac{ \frac{ \sin(y^3) }{ x^2+y^2 } }{ \sqrt{x^2+y^2} }=\lim_{(x,y)\to (0,0)} \frac{ \sin(y^3) }{ (x^2+y^2)\sqrt{x^2+y^2} }=\frac{ 0 }{ 0 },
\end{equation}
donc la limite n'existe pas.

\item
Soit $u=(u_1,u_2)$. Nous avons
\begin{equation}
	\frac{ \partial f }{ \partial u }(0,0)=
	\begin{cases}
	\lim_{t\to 0} \frac{ f(0,0)-f(0,0) }{ t }=0	&	\text{si $u=(0,0)$}\\
	\lim_{t\to 0} \frac{ f(tu_1,tu_2)-f(0,0) }{ t }	&	 \text{si $u\neq(0,0)$.}
\end{cases}
\end{equation}
Dans le second cas, nous trouvons
\begin{equation}
	\lim_{t\to 0} \frac{ \sin(t^3u_2^3) }{ t^2u_1^2+t^2u_2^2 }\cdot\frac{1}{ t }=\lim_{t\to 0} \left( \frac{ \sin(t^3u_1^3) }{ t^3u_1^3 }\cdot\frac{ u_1^3 }{ u_1^2+u_2^2 } \right)=\frac{ u_1^3 }{ u_1^2+u_2^2 }.
\end{equation}
Donc, la dérivée directionnelle existe dans toutes les directions, et donc $f$ est différentiable en $(0,0)$.


    \end{enumerate}



\part{Calcul différentiel et intégral 2}
Cette partie sera progressivement digérée dans les autres chapitres. 
% This is part of the Exercices et corrigés de CdI-2.
% Copyright (C) 2008-2011
%   Laurent Claessens
% See the file fdl-1.3.txt for copying conditions.


\chapter*{Les séances d'exercices}


\chapter{Convergences et espaces fonctionnels}
	% This is part of the Exercices et corrigés de CdI-2.
% Copyright (C) 2008, 2009, 2012
%   Laurent Claessens
% See the file fdl-1.3.txt for copying conditions.


%+++++++++++++++++++++++++++++++++++++++++++++++++++++++++++++++++++++++++++++++++++++++++++++++++++++++++++++++++++++++++++
					\section[Intégrales de fonctions et domaines non bornées]{Intégrales de fonctions non bornées sur des domaines non bornés}
%+++++++++++++++++++++++++++++++++++++++++++++++++++++++++++++++++++++++++++++++++++++++++++++++++++++++++++++++++++++++++++

%---------------------------------------------------------------------------------------------------------------------------
					\subsection[Fonctions et ensembles non bornés]{Intégrales de fonctions non bornées sur des ensembles non bornés}
%---------------------------------------------------------------------------------------------------------------------------

Soit $f\colon \eR^n\to \overline{ \eR }$, une fonction positive. On dit qu'elle est \defe{intégrable}{intégrable (fonction)} sur $E\subset\eR^n$ si
\begin{enumerate}
    \item $\forall r>0$, la fonction $f_r(x)=f(x)\mtu_{f<r}$ est intégrable sur $E_r$;
\item la limite $\lim_{r\to\infty}\int_{E_r}f_r$ est finie.
\end{enumerate}
Dans ce cas, on pose 
\begin{equation}
	\int_Ef=\lim_{r\to\infty}\int_{E_r}f_r.
\end{equation}

\begin{theorem}[Page I.38]		\label{ThoFnTestIntnnBorn}
Soit $E$ mesurable dans $\eR^n$ et $f\colon E\to \overline{ \eR }$. Si $f$ est mesurable et si il existe $g\colon E\to \overline{ \eR }$ intégrable sur $E$ telle que $| f(x) |\leq g(x)$ pour tout $x\in E$, alors $f$ est intégrable sur~$E$.

Réciproquement, si $f$ est intégrable sur $E$, alors $f$ est mesurable.
\end{theorem}

%---------------------------------------------------------------------------------------------------------------------------
					\subsection{Passage à la limite sous le signe intégral}
%---------------------------------------------------------------------------------------------------------------------------

Un autre résultat très important pour l'étude de l'intégrabilité est le théorème de la \defe{convergence dominée de Lebesgue}{}:
\begin{theorem}
	Soit $E\subset \eR^n$ un ensemble mesurable et $\{ f_k \}$, une suite de fonctions intégrables sur $E$ qui converge simplement vers une fonction $f\colon E\to \overline{ \eR }$. Supposons qu'il existe une fonction $g$ intégrable sur $E$ telle que pour tout $k$,
\begin{equation}
	| f(x) |\leq g(x)
\end{equation}
pour tout $x\in E$. Alors $f$ est intégrable sur $E$ et 
\begin{equation}
	\int_Ef=\lim_{k\to\infty}\int_Ef_k.
\end{equation}
\end{theorem}

%---------------------------------------------------------------------------------------------------------------------------
					\subsection{Théorème de Fubini et changement de variables}
%---------------------------------------------------------------------------------------------------------------------------

\begin{theorem}[Fubini]\index{théorème!Fubini}\index{Fubini!théorème}		\label{ThoFubini}
Soit $(x,t)\mapsto f(x,y)\in\bar \eR$ une fonction intégrable sur $B_n\times B_m\subset\eR^{n+m}$ où $B_n$ et $B_m$ sont des ensembles mesurables de $\eR^n$ et $\eR^m$. Alors :
\begin{enumerate}
\item pour tout $x\in B_n$, sauf éventuellement en les points d'un ensemble $G\subset B_n$ de mesure nulle, la fonction $y\in B_m\mapsto f(x,y)\in\bar\eR$ est intégrable sur $B_m$
\item
la fonction
\begin{equation}
	x\in B_n\setminus G\mapsto\int_{B_m}f(x,y)dy\in\eR
\end{equation}
est intégrable sur $B_n\setminus G$

\item 
On a
\begin{equation}
	\int_{B_n\times B_m}f(x,y)dxdy=\int_{B_n}\left( \int_{B_m}f(x,y)dy \right)dx.
\end{equation}

\end{enumerate}
\end{theorem}

Notons en particulier que si $f(x,y)=\varphi(x)\phi(y)$, alors $\int_{B_m}\varphi(y)dy$ est une constante qui peut sortir de l'intégrale sur $B_n$, et donc
\begin{equation}		\label{EqFubiniFactori}
	\int_{B_n\times B_m}\varphi(x)\phi(y)dxdy=\int_{B_n}\varphi(x)dx\int_{B_m}\phi(y)dy.
\end{equation}

%---------------------------------------------------------------------------------------------------------------------------
					\subsection{Intégrale en dimension un}
%---------------------------------------------------------------------------------------------------------------------------

\begin{proposition}[Critère de comparaison]
Soit $f$ mesurable sur $]a,\infty[$ et bornée sur tout $]a,b]$, et supposons qu'il existe un $X_0\geq a$, tel que sur $]X_0,\infty[$,
\begin{equation}
	| f(x) |\leq g(x)
\end{equation}
où $g(x)$ est intégrable. Alors $f(x)$ est intégrable sur $]a,\infty[$.
\end{proposition}

\begin{corollary}[Critère d'équivalence]
Soient $f$ et $g$ des fonctions mesurables et positives ou nulles sur $]a,\infty[$, bornées sur tout $]a,b]$, telles que 
\begin{equation}
	\lim_{x\to\infty}\frac{ f(x) }{ g(x) }=L
\end{equation}
existe dans $\bar\eR$.
\begin{enumerate}
\item Si $L\neq\infty$ et $\int_{a}^{\infty}g(x)$ existe, alors $\int_a^{\infty}f(x)dx$ existe,
\item Si $L\neq 0$ et si $\int_a^{\infty}f(x)dx$ existe, alors $\int_a^{\infty}g(x)dx$ existe,
\end{enumerate}
\end{corollary}

\begin{corollary}[Critère des fonctions test]			\label{CorCritFonsTest}
Soit $f(x)$ une fonction mesurable et positive ou nulle sur $]a,\infty[$ et bornée pour tout $]a,b]$. Nous posons
\begin{equation}
	L(\alpha)=\lim_{x\to\infty}x^{\alpha}f(x),
\end{equation}
et nous supposons qu'elle existe.
\begin{enumerate}
\item Si il existe $\alpha>1$ tel que $L(\alpha)\neq\infty$, alors $\int_a^{\infty}f(x)dx$ existe,
\item Si il existe $\alpha\leq1$ et $L(\alpha)\neq 0$, alors $\int_a^{\infty}f(x)dx$ n'existe pas.
\end{enumerate}
\end{corollary}

\begin{corollary}		\label{CorAlphaLCasInteabf}
	Soit $f\colon ]a,b]\to \eR$ une fonction mesurable, positive ou nulle, et bornée sur $[a+\epsilon,b]$ $\forall\epsilon>0$. Si $\lim_{x\to a}(x-a)^{\alpha}f(x)=L$ existe, alors
	\begin{enumerate}
		\item Si $\alpha<1$ et $L\neq\infty$, alors $\int_a^bf(x)dx$ existe,
		\item Si $\alpha\geq 1$ et $L\neq 0$, alors $\int_a^bf(x)dx$ n'existe pas.
	\end{enumerate}
\end{corollary}

%---------------------------------------------------------------------------------------------------------------------------
					\subsection{Intégrales convergentes}
%---------------------------------------------------------------------------------------------------------------------------

\begin{definition}

Soit $f$, une fonction mesurable sur $[a,\infty[$, bornée sur tout intervalle $[a,b]$. On dit que l'intégrale
\begin{equation}
	\int_a^{\infty}f(x)dx
\end{equation}
\defe{converge}{intégrale!convergente} si la limite
\begin{equation}		\label{EqDEfConvergeZeroInftX}
	\lim_{X\to\infty}\int_a^{X}f
\end{equation}
existe et est finie.
\end{definition}


%+++++++++++++++++++++++++++++++++++++++++++++++++++++++++++++++++++++++++++++++++++++++++++++++++++++++++++++++++++++++++++
					\section{Fonctions définies par des intégrales et régularisation}
%+++++++++++++++++++++++++++++++++++++++++++++++++++++++++++++++++++++++++++++++++++++++++++++++++++++++++++++++++++++++++++
%+++++++++++++++++++++++++++++++++++++++++++++++++++++++++++++++++++++++++++++++++++++++++++++++++++++++++++++++++++++++++++
%					\section[Fonctions définies par des intégrales]{Fonctions définies par des intégrales et régularisation}

%+++++++++++++++++++++++++++++++++++++++++++++++++++++++++++++++++++++++++++++++++++++++++++++++++++++++++++++++++++++++++++

Supposons $A\subset\eR^m$  et $B\subset\eR^n$ compact. Nous considérons $f(x,t)\colon A\times B\to \eR$ une fonction bornée sur $A\times B$ et intégrable par rapport à $t$ pour tout $x\in A$. Soit $F\colon A\to \eR$ définie par
\begin{equation}
	F(x)=\int_Bf(x,t)dt.
\end{equation}

%---------------------------------------------------------------------------------------------------------------------------
					\subsection{Fonction définies par une intégrale sur un compact}
%---------------------------------------------------------------------------------------------------------------------------

\begin{proposition}		\label{PropDerrSSIntegraleDSD}
Supposons $A\subset\eR^m$ ouvert et $B\subset\eR^n$ compact. Si pour un $i\in\{ i,\ldots,n \}$, la dérivée partielle $\frac{ \partial f }{ \partial x_i }$ existe dans $A\times B$ et est continue, alors $\frac{ \partial F }{ \partial x_i }$ existe dans $A$, est continue et
\begin{equation}
	\frac{ \partial F }{ \partial x_i }=\int_B\frac{ \partial f }{ \partial x_i }(x,t)dt,
\end{equation}
l'égalité signifie que l'on peut \og dériver sous le signe intégral\fg.
\end{proposition}

%---------------------------------------------------------------------------------------------------------------------------
					\subsection{Intégrale sur un segment variable}
%---------------------------------------------------------------------------------------------------------------------------


\begin{proposition}		\label{PropDerrFnAvecBornesFonctions}
Soit $f(x,t)$ une fonction continue sur $[\alpha,\beta]\times[a,b]$, telle que $\frac{ \partial f }{ \partial x }$ existe et soit continue sur $]\alpha,\beta[\times[a,b]$. Soient $\varphi(x)$ et $\psi(x)$, des fonctions continues de $[\alpha,\beta]$ dans $\eR$ et admettant une dérivée continue sur $]\alpha,\beta [$. Alors la fonction
\begin{equation}
	F(x)=\int_{\varphi(x)}^{\psi(x)}f(x,t)dt
\end{equation}
admet une dérivée continue sur $]\alpha,\beta[$ et
\begin{equation}	\label{EqFormDerrFnAvecBorneNInt}
	\frac{ dF }{ dx }=\int_{\varphi(x)}^{\psi(x)}\frac{ \partial f }{ \partial x }(x,t)dt+f\big( x,\psi(x) \big)\cdot\frac{ d\psi }{ dx }- f\big( x,\varphi(x) \big)\cdot\frac{ d\varphi }{ dx }.
\end{equation}
\end{proposition}

%---------------------------------------------------------------------------------------------------------------------------
					\subsection{Intégrales convergentes}
%---------------------------------------------------------------------------------------------------------------------------

\begin{theorem}		\label{ThoInDerrtCvUnifFContinue}
	Uniforme convergence, continuité et dérivation.
	\begin{enumerate}
		\item 
		Soit $f(x,t)$ continue sur $[\alpha,\beta]\times[a,\alpha[$ et $F(x)=\int_a^{\infty}f(x,t)dt$; cette intégrale étant supposée uniformément convergente. Alors $F(x)$ est continue.
		\item
		Supposons $f$ continue et sa dérivée partielle $\frac{ \partial f }{ \partial x }$ continue sur $[\alpha,\beta]\times[a,\alpha[$. Supposons que $F(x)=\int_a^{\infty}f(x,t)dt$ converge et que $\int_a^{\infty}\frac{ \partial f }{ \partial x }dt$ converge uniformément. Alors $F$ est $C^1$ sur $[\alpha,\beta]$ et 
		\begin{equation}
			\frac{ dF }{ dx }=\int_a^{\infty}\frac{ \partial f }{ \partial x }dt.
		\end{equation}
	\end{enumerate}
\end{theorem}
%---------------------------------------------------------------------------------------------------------------------------
					\subsection{Critères de convergence uniforme}
%---------------------------------------------------------------------------------------------------------------------------


Affin de tester l'uniforme convergence d'une intégrale, nous avons le \defe{critère de Weierstrass}{Critère!Weierstrass}:
\begin{theorem}		\label{ThoCritWeiIntUnifCv}
Soit $f(x,t)\colon [\alpha,\beta]\times[a,\infty[ \to \eR$, une fonction dont la restriction à toute demi-droite $x=cst$ est mesurable. Si $| f(x,t) |< \varphi(t)$ et $\int_a^{\infty}\varphi(t)dt$ existe, alors l'intégrale
\begin{equation}
	\int_0^{\infty}f(x,t)dt
\end{equation}
est uniformément convergente.
\end{theorem}


Le théorème suivant est le \defe{critère d'Abel}{Critère!Abel pour intégrales} :
\begin{theorem}		\label{ThoAbelIntUnif}
	Supposons que $f(x,t)=\varphi(x,t)\psi(x,t)$ où $\varphi$ et $\psi$ sont bornée et intégrables en $t$ au sens de Riemann sur tout compact $[a,b]$, $b\geq a$. Supposons que :
	\begin{enumerate}
		\item $| \int_a^{T}\varphi(x,t)dt |\leq M$ où $M$ est indépendant de $T$ et de $x$,
		\item $\psi(x,t)\geq 0$,
		\item pour tout $x\in[\alpha,\beta]$, $\psi(x,t)$ est une fonction décroissante de $t$,
		\item les fonctions $x\mapsto \psi(x,t)$ convergent uniformément vers $0$ lorsque $t\to\infty$.
	\end{enumerate}
	Alors l'intégrale
	\begin{equation}
		\int_a^{\infty}f(x,t)dt
	\end{equation}
	est uniformément convergente.
\end{theorem}


\chapter{Équations différentielles}
	% This is part of the Exercices et corrigés de CdI-2.
% Copyright (C) 2008-2011
%   Laurent Claessens
% See the file fdl-1.3.txt for copying conditions.


Ici se trouvent les rappels sur les équations différentielles pour les premières et seconde BAC. Tout ce qui se trouve dans cette section n'est donc pas à lire pour les premières.

\label{sec:equadiff1erordre}
Une équation différentielle ordinaire est la recherche de toutes les fonctions définie sur une partie de $\eR$ satisfaisant à une certaine égalité, faisant intervenir les dérivées de la fonction recherchée.


Dans la suite, $I$ désignera un intervalle de $\eR$. Une fonction sera \defe{dérivable sur $I$}{Dérivalble!fonction} si elle est dérivable au sens usuel sur l'intérieur de $I$, et si elle est dérivable à droite (resp. à gauche) sur l'éventuel bord gauche (resp. droit) de $I$.

\begin{definition}
  Une \defe{équation différentielle ordinaire d'ordre $n$ sur $I$}{Équation différentielle!ordinaire!ordre 1} est la recherche d'une fonction $y : I \to \eR$ dérivable $n$ fois, satisfaisant à une équation du type
  \begin{equation}\label{eq:equadiff}
    F(t, y(t), y^\prime(t), \ldots, y^{n\prime}(t)) = 0 \quad \text{pour tout $t \in I$}
  \end{equation}
  où $I$ est un intervalle de $\eR$ et \begin{math}F : (I \times D) \subset (\eR\times\eR^{n+1})\to \eR\end{math} est une fonction donnée.
\end{definition}

\begin{remark}
L'équation différentielle~(\ref{eq:equadiff}) sera raccourcie sous la forme
  \begin{equation}
    F(t, y, y^\prime, \ldots, y^{n\prime}) = 0
  \end{equation}
  où la dépendance en $t$ est sous-entendue.
\end{remark}

\begin{example}
	Soit $f : I \to \eR$ une fonction continue fixée. L'équation différentielle
	\begin{equation}
		y^\prime = f(t)
	\end{equation}
	se ramène à la recherche des primitives de $f$ sur l'intervalle $I$.
\end{example}

%+++++++++++++++++++++++++++++++++++++++++++++++++++++++++++++++++++++++++++++++++++++++++++++++++++++++++++++++++++++++++++
					\section{Théorèmes d'existence et d'unicité}
%+++++++++++++++++++++++++++++++++++++++++++++++++++++++++++++++++++++++++++++++++++++++++++++++++++++++++++++++++++++++++++

%---------------------------------------------------------------------------------------------------------------------------
					\subsection{Équations différentielles sous forme normale}
%---------------------------------------------------------------------------------------------------------------------------

%+++++++++++++++++++++++++++++++++++++++++++++++++++++++++++++++++++++++++++++++++++++++++++++++++++++++++++++++++++++++++++
					\section{Systèmes différentiels linéaires}
%+++++++++++++++++++++++++++++++++++++++++++++++++++++++++++++++++++++++++++++++++++++++++++++++++++++++++++++++++++++++++++

Rapide aperçu de ce qui est dit entre les pages II.30 et II.49.

%---------------------------------------------------------------------------------------------------------------------------
					\subsection{La magie de l'exponentielle\ldots}
%---------------------------------------------------------------------------------------------------------------------------

Prenons l'équation différentielle très simple
\begin{equation}
	y'=ay.
\end{equation}
La solution est $y(t)=A e^{at}$. Et si on a la donnée que Cauchy $y(t_0)=y_0$, alors
\begin{equation}		\label{EqytexposimpleProp}
	y(t)=A e^{at} e^{-at_0} e^{at_0}= e^{a(t-t_0)}y(t_0).
\end{equation}
Donc on a le facteur multiplicatif $ e^{a(t-t_0)}$ qui sert à faire passer de $y(0)$ à $y(t)$. C'est un peu un opérateur d'évolution. Ce qui fait la magie  de l'exponentielle, c'est son développement en série
\begin{equation}		\label{EqDevExpoMag}
	e^x=1+x+\frac{ x^2 }{ 2 }+\frac{ x^3 }{ 3! }+\frac{ x^4 }{ 4! }+\ldots
\end{equation}
qui est tel que chaque terme est la dérivée du terme suivant.

Maintenant, si on a un système
\begin{equation}
	\bar y'=A\bar y,
\end{equation}
il n'est pas du tout étonnant d'avoir comme solution $\bar y(t)= e^{At}$ où l'exponentielle de la matrice est définie exactement par la série \eqref{EqDevExpoMag}. C'est un peu longuet, mais dans le cours, c'est effectivement ce qui est prouvé. La matrice résolvante $R(t,t_0)\colon \bar y_0\to \bar y(t;t_0,y_0)$ est donné par
\begin{equation}
	R(t,t_0)= e^{(t-t_0)A},
\end{equation}
exactement comme dans l'équation \eqref{EqytexposimpleProp}.

%---------------------------------------------------------------------------------------------------------------------------
					\subsection{\ldots mais la difficulté}
%---------------------------------------------------------------------------------------------------------------------------

Maintenant, il est suffisant de calculer des exponentielles de matrices pour résoudre des systèmes. Hélas, il est en général très difficile de calculer des exponentielles. Tu peux essayer de prouver les deux suivantes :
\begin{equation}
	\begin{aligned}[]
		A=\begin{pmatrix}
	0	&	a	\\ 
	-a	&	0	
\end{pmatrix}	&\leadsto  e^{A}=\begin{pmatrix}
	\cos(a)	&	\sin(a)	\\ 
	-\sin(a)	&	\cos(a)	
\end{pmatrix}\\
		S=\begin{pmatrix}
	0	&	a	\\ 
	a	&	0	
\end{pmatrix}	&\leadsto  e^{S}=\begin{pmatrix}
	\cosh(a)	&	\sinh(a)	\\ 
	\sinh(a)	&	\cosh(a)	
\end{pmatrix}.
	\end{aligned}
\end{equation}
La première, tu vas la revoir si tu fais de la géométrie différentielle ou de la mécanique quantique : l'algèbre de Lie du groupe des matrices orthogonales de déterminant $1$ est l'algèbre des matrices antisymétriques.

La seconde se retrouve en relativité parce que $e^S$ est la matrice qui préserve $x^2-y^2$, tout comme $e^A$ préserve $x^2+y^2$. Si tu as besoin d'une rafraîchissure de mémoire sur les fonctions hyperboliques, ou bien si leur utilité en relativité t'intéresse, tu peux aller voir la partie sur la relativité \href{http://student.ulb.ac.be/~lclaesse/physique-math.pdf}{ici}\footnote{ \url{http://student.ulb.ac.be/~lclaesse/physique-math.pdf}}.

%---------------------------------------------------------------------------------------------------------------------------
					\subsection{La recette}
%---------------------------------------------------------------------------------------------------------------------------

Afin d'éviter de devoir calculer explicitement des exponentielles de matrices, nous faisons appel à toutes sortes de trucs, dont la forme de Jordan. Le résultat final est la méthode suivante. Soit le système homogène
\begin{equation}
	\bar y'=A\bar y.
\end{equation}

\let\oldTheEnumi\theenumi
\renewcommand{\theenumi}{\arabic{enumi}}
\begin{enumerate}

\item 
D'abord, nous calculons les valeurs propres de $A$.

\item
Ensuite les vecteurs propres.

\item\label{ItemRapSystDc}
Une bonne valeur propre, c'est une valeur propre dont l'espace propre a une dimension égale à sa multiplicité. C'est à dire que si $\lambda$ est de multiplicité $m$, alors on a, dans les bon cas,  $m$ vecteur propres linéairement indépendants.

Dans ce cas, si $v_1,\ldots,v_m$ sont les vecteurs, alors on a les solutions linéairement indépendantes suivantes :
\begin{equation}
	\begin{pmatrix}
	\vdots	\\ 
	v_1	\\ 
		\vdots	
\end{pmatrix} e^{\lambda t},\ldots,
\begin{pmatrix}
	\vdots		\\ 
	v_m	\\ 
	\vdots		
\end{pmatrix} e^{\lambda t}.
\end{equation}
Pour chaque bonne valeur propre, ça nous fait un tel paquet de solutions linéairement indépendantes.

\item
Si $\lambda$ n'est pas une bonne valeur propre, alors les choses se compliquent. Mettons que $\lambda$ ait $k$ vecteurs propres en moins que sa multiplicité. Dans ce cas, il faut chercher des solutions sous la forme
\begin{equation}		\label{EqEqRapAsTestPolk}
	 \begin{pmatrix}
	a^{(k)}_1t^k+\ldots+a_1^{(0)}	\\ 
	\vdots	\\ 
	a^{(k)}_1t^k+\ldots+a_n^{(0)}		
\end{pmatrix} e^{\lambda t}.
\end{equation}
C'est à dire qu'on prend comme coefficient de $ e^{\lambda t}$, un vecteur de polynômes de degré $k$. Il faut mettre cela dans l'équation de départ pour voir quelles sont les contraintes sur les constantes $a_i^{(j)}$ introduites.

\item\label{ItemRapSystDe}
Nous avons un cas particulier du cas précédent. Si $\lambda$ est une valeur propre de multiplicité $m$ qui n'a que un seul vecteur propre $v$, alors il faut chercher des polynômes de degré $m-1$, et on peut directement fixer le coefficient de $t^{m-1}$, ce sera l'unique vecteur propres :
\begin{equation}
\left[
	\begin{pmatrix}
	\vdots	\\ 
	v	\\ 
	\vdots	
\end{pmatrix}+
\begin{pmatrix}
	a_1^{(m-2)}	\\ 
	\vdots	\\ 
	a_n^{(m-2)}	
\end{pmatrix}t^{m-2}+\ldots
\right] e^{\lambda t}.
\end{equation}
Cela économise quelque calculs par rapport à poser brutalement \eqref{EqEqRapAsTestPolk}.

\end{enumerate}
\let\theenumi\oldTheEnumi

%+++++++++++++++++++++++++++++++++++++++++++++++++++++++++++++++++++++++++++++++++++++++++++++++++++++++++++++++++++++++++++
					\section{Continuité et dérivabilité des solutions}
%+++++++++++++++++++++++++++++++++++++++++++++++++++++++++++++++++++++++++++++++++++++++++++++++++++++++++++++++++++++++++++

%+++++++++++++++++++++++++++++++++++++++++++++++++++++++++++++++++++++++++++++++++++++++++++++++++++++++++++++++++++++++++++
					\section{Équations résolubles}
%+++++++++++++++++++++++++++++++++++++++++++++++++++++++++++++++++++++++++++++++++++++++++++++++++++++++++++++++++++++++++++

\subsection{Équation à variables séparées}
\label{sec:varsep}


Nous suivons les pages 311--312 du cours. Une équation différentielle à variables séparées est une équation différentielle d'ordre $1$ qui peut s'écrire sous la forme
\begin{equation}		\label{EqGeneTheSep}
	y'=u(t)f(y)
\end{equation}
où $u\colon I\to \eR$ et $f\colon J\to \eR$ sont deux fonctions continues données. Cette équation est résolue par le théorème 1 de la page 311 du cours. Voyons rapidement comment l'appliquer en pratique. 

Rappelons nous que le problème est de trouver les fonction $y\colon I'\to \eR$ telles que pour tout $t \in I'$,
\begin{equation}
	y'(t)=u(t)f\big( y(t) \big).
\end{equation}
Nous considérons $U$, une primitive de $u$ sut $I$ et $G$, une primitive de $1/f$ sur $J$. Si $I'\subseteq I$ et $y\colon I'\to J$, alors $y$ est solution de \eqref{EqGeneTheSep} si et seulement si il existe une constante $C$ telle que
\begin{equation}		\label{EqSolSepThe}
	G\big( y(t) \big)=U(t)+C.
\end{equation}
La recherche des solutions de l'équation différentielle se ramène donc à la recherche de primitives et de solutions d'une équation algébrique (il faut isoler $y(t)$ dans \eqref{EqSolSepThe}). Réciproquement toute solution régulière de cette dernière relation est solution de l'équation différentielle.

Remarque : lorsque nous cherchons $U$ et $G$, nous ne cherchons que \emph{une} primitive. Il ne faut pas considérer des constantes d'intégration à ce niveau.


%---------------------------------------------------------------------------------------------------------------------------
					\subsection{Équation linéaire du premier ordre}
%---------------------------------------------------------------------------------------------------------------------------


%---------------------------------------------------------------------------------------------------------------------------
					\subsection{Équations et systèmes linéaire à coefficients constants}
%---------------------------------------------------------------------------------------------------------------------------

Nous regardons l'équation
\begin{equation}	\label{EqLinConstantRappels}
	y^{(n)} + a_1 y^{(n-1)} + \ldots + a_{n-1} y^\prime + a_n y = v(t)
\end{equation}
où les coefficients $a_k$ sont maintenant des constantes. La méthode est donnée à la page $311$ du cours. Il faut commencer par résoudre le polynôme caractéristique
\begin{equation}
	r^n+a_1 r^{n-1}+\ldots +a_n=0.
\end{equation}
Si $\lambda_1,\ldots,\lambda_k$ sont les solutions avec multiplicité $\mu_1,\ldots,\mu_k$, alors le \defe{système fondamental}{Système fondamental} de solutions linéairement indépendantes est l'ensemble suivant de solutions à l'équation homogène :
\begin{equation}
	\begin{aligned}[]
		 e^{\lambda_1 t},t e^{\lambda_1 t},	&	\ldots,t^{\mu_1-1} e^{\lambda_1  t}\\
							&\vdots\\
		 e^{\lambda_k t},t e^{\lambda_k t},	&\ldots,t^{\mu_k-1} e^{\lambda_k  t}.
	\end{aligned}
\end{equation}
Nous notons $y_i$ ces solutions. La solution générale de l'équation homogène est donc donnée par
\begin{equation}
	y_H=\sum_i c_i y_i.
\end{equation}
Afin de trouver la solution générale de l'équation non homogène, nous appliquons la méthode de variation des constantes, en imposant les $n-1$ conditions
\begin{equation}		\label{EqVarCstSubtil}
	\sum_{i=1}^n c'_i(t)y_i^{(l)}(t)=0
\end{equation}
avec $l=0,\ldots,n-2$. Ces condition plus l'équation de départ \eqref{EqLinConstantRappels} forment un système de $n$ équations différentielles pour les $n$ fonctions inconnues $c_i(t)$.

Cette condition peut paraître mystérieuse. Elle est posée à la page 337 du cours de première année. Il est cependant encore possible de travailler sans poser la condition \eqref{EqVarCstSubtil} en suivant la recette donnée à la page 338, en calculant des déterminants de Wronskien. Des exemples sont donnés dans les exercices sur le second ordre en première\footnote{\url{http://student.ulb.ac.be/~lclaesse/CdI1.pdf}}.

\paragraph{Si les coefficients ne sont pas constants ?}

Une équation différentielle linéaire d'ordre $n$ sur $I$ est une
équation de la forme
\begin{equation}	\label{EqLinRappels}
	y^{(n)} + u_1(t) y^{(n-1)} + \ldots + u_{n-1}(t) y^\prime + u_n(t) y = v(t)
\end{equation}
où $v$ et $u_k$ sont des fonctions continues fixées de $I$ vers $\eR$.

Pour résoudre cette équation, il faut commencer par résoudre l'équation homogène correspondante (c'est à dire celle que l'on obtient en posant $v(t)=0$). Ensuite, nous trouvons la solution de l'équation \eqref{EqLinRappels} en appliquant la méthode de la \defe{variation des constantes}{Variations des constantes}.

Donnons un exemple du pourquoi la méthode de variations des constantes est efficace. Soit l'équation 
\begin{equation}		\label{EqDiffExempleVarCst}
	u'+f(t)u=g(t),
\end{equation}
 et disons que $u_H$ est une solution de l'équation homogène. La méthode de variations des constantes consiste à poser $u(t)=K(t)u_H(t)$, et donc $u'(t)=K'u_H+Ku_H'$. En remettant dans l'équation de départ,
\begin{equation}
	K'u_H+Ku_H'+fKu_H=g.
\end{equation}
La somme $Ku_H'+fKu_H$ est nulle, par définition de $u_H$. Par conséquent, il ne reste que
\begin{equation}
	K'=\frac{ g(t) }{ u_H }.
\end{equation}
Lorsqu'on utilise la méthode de variation des constantes, nous trouvons toujours une simplification \og miraculeuse\fg.

En première année, seul le cas où les $u_i$ sont des constantes est vu. C'est en deuxième année que le cas avec des coefficients fonctions de $t$ est étudié\quext{Merci de me signaler si je me trompe.}.


%---------------------------------------------------------------------------------------------------------------------------
					\subsection{Équation homogène}
%---------------------------------------------------------------------------------------------------------------------------
\label{SubSecEqDiffHomo}

Page II.71. C'est une équation de la forme
\begin{equation}
	y'=f(t,y)
\end{equation}
où $f(\lambda t,\lambda y)=f(t,y)$ pour tout $\lambda\neq 0$.
\begin{lemma}[Lemme page II.71]
L'équation $y'=f(t,y)$ est homogène si et seulement si $f(t,y)$ est une fonction de $y/t$ seulement.
\end{lemma}
Pour résoudre l'équation homogène, on pose
\begin{equation}		\label{EqDiffHomoPoser}
	z(t)=\frac{ y(t) }{ t },
\end{equation}
donc $tz=y$, et 
\begin{equation}
	y'(t)=tv'(t)+v(t),
\end{equation}
à remettre dans l'équation de départ.

%---------------------------------------------------------------------------------------------------------------------------
					\subsection{L'équation \texorpdfstring{$y'=f\left( \frac{ at+by+c }{ a't+b'y+c' } \right)$}{yat}}
%---------------------------------------------------------------------------------------------------------------------------

La théorie relative à cette équation se trouve à la page II.72 du cours.

%---------------------------------------------------------------------------------------------------------------------------
					\subsection{Équation de Bernoulli}
%---------------------------------------------------------------------------------------------------------------------------
\label{SubSecBernh}

Voir page II.73. C'est une équation du type
\begin{equation}	\label{EqBerNDiffalp}
	y'=a(t)y+b(t)y^{\alpha}
\end{equation}
où $\alpha\neq 0$ ou $1$. Pour la résoudre, on divise l'équation par $y^{\alpha}$, et on pose $u=y^{1-\alpha}$, et on tombe sur une équation linéaire
\begin{equation}
	u'=(1-\alpha)\big( a(t)u+b(t) \big).
\end{equation}


%---------------------------------------------------------------------------------------------------------------------------
					\subsection{Équation de Riccati}
%---------------------------------------------------------------------------------------------------------------------------
\label{SubSecRicatti}

Contrairement à ce qui est écrit dans le cours, l'orthographe est \emph{\href{http://fr.wikipedia.org/wiki/Jacopo_Riccati}{Riccati}}, et non \emph{Ricatti}.

C'est une équation de la forme 
\begin{equation}		\label{EqDiffGFeneRicatti}
	y'=a(t)y^2+b(t)y+c(t).
\end{equation}

En général, on ne peut pas la résoudre, mais si on en connaît \emph{a priori} des solutions particulières, alors on peut s'en sortir.

\begin{enumerate}

\item 
Si on sait que $y_1(t)$ est une solution, alors on pose
\begin{equation}
	y(t)=y_1(t)+\frac{1}{ u(t) },
\end{equation}
et on obtient une équation linéaire
\begin{equation}
	u'=-\big( 2y_1(t)a(t)+b(t) \big)u-a(t).
\end{equation}

\item
Si $y_1$ et $y_2$ sont solutions, alors nous avons $y$ sous forme implicite
\begin{equation}
	\frac{ y-y_1 }{ y-y_2 }=K e^{\int a(t)\big( y_1(t)-y_2(t) \big)dt}.
\end{equation}
\end{enumerate}

Pour résoudre une équation de Ricatti, il faut donc d'abord deviner une ou deux solutions.

%---------------------------------------------------------------------------------------------------------------------------
					\subsection{Équation différentielle exacte}
%---------------------------------------------------------------------------------------------------------------------------
\label{SubSecEqDiffExacte}

%///////////////////////////////////////////////////////////////////////////////////////////////////////////////////////////
					\subsubsection{Résolution lorsque tout va bien}
%///////////////////////////////////////////////////////////////////////////////////////////////////////////////////////////

Avant de vous lancer dans les équations différentielles exacte, vous devez lire la section sur les formes différentielles \ref{SecFormDiffRappel}. Une équation différentielle exacte est de la forme $P(t,y)+Q(t,y)y'=0$ que nous allons écrire sous la forme
\begin{equation}		\label{EqExacteDiff}
	P(t,y)dt+Q(t,y)dy=0.
\end{equation}
Nous savons que si $\partial_yP=\partial_tQ$, alors il existe une fonction $f(t,y)$ telle que $Pdt+Qdy=df$. Pour trouver une telle fonction, nous pouvons simplement intégrer la forme $Pdt+Qdy$. En effet, si $\gamma\colon [0,1]\to \eR^2$ est un chemin tel que $\gamma(0)=(0,0)$ et $\gamma(1)=(t,y)$, alors en définissant
\begin{equation}
	f(t,y)=\int_{\gamma}[Pdt+Qdt]=\int_{0}^1\big[ (P\circ\gamma)(u)dt+(Q\circ\gamma)(u) \big]\big( \gamma'(u) \big)du,
\end{equation}
nous avons $df=Pdt+Qdy$. N'importe quel chemin fait l'affaire. Calculons avec $\gamma(u)=(tu,yu)$. La dérivée de ce chemin est donnée par
\begin{equation}
	\gamma'(u)=t\begin{pmatrix}
	1	\\ 
	0	
\end{pmatrix}+y\begin{pmatrix}
	0	\\ 
	1	
\end{pmatrix}.
\end{equation}
Étant donné que $dt\begin{pmatrix}
	a	\\ 
	b	
\end{pmatrix}=a$ et $dy\begin{pmatrix}
	a	\\ 
	b	
\end{pmatrix}=b$, nous avons
\begin{equation}
	\begin{aligned}[]
	f(t,y)&=\int_0^1[Pdt+Qdy]\big( \gamma(u) \big)\left( t\begin{pmatrix}
	1	\\ 
	0	
\end{pmatrix}+y\begin{pmatrix}
	0	\\ 
	1	
\end{pmatrix} \right)du\\
		&=\int_0^1P\big( \gamma(t) \big)tdu+\int_0^1Q\big( \gamma(t) \big)ydu\\
		&=\int_0^1\big[ tP(tu,uy)+yQ(tu,yu) \big]du.
	\end{aligned}
\end{equation}
Nous retrouvons exactement la formule \eqref{EqFormI33Fffdd} de l'exercice \ref{exo_I-3-3}. Si ça t'étonne, c'est que tu n'as pas compris ;) Dans le cas où nous avons la fonction $f$ qui vérifie $P=\partial_tf$ et $Q=\partial_yf$, l'équation \eqref{EqExacteDiff} devient
\begin{equation}
	\frac{ \partial f }{ \partial t }+\frac{ \partial f }{ \partial y }\frac{ dy }{ dt }=0,
\end{equation}
c'est à dire 
\begin{equation}
	\frac{ d }{ dt }\Big[ f\big( t,y(t) \big) \Big]=0,
\end{equation}
dont la solution
\begin{equation}
	f\big( t,y(t) \big)=C
\end{equation}
donne la solution $y(t)$ sous forme implicite.

%///////////////////////////////////////////////////////////////////////////////////////////////////////////////////////////
					\subsubsection{Facteur intégrant (quand tout ne va pas bien)}
%///////////////////////////////////////////////////////////////////////////////////////////////////////////////////////////

Si la forme $Pdt+Qdy$ n'est pas exacte, il n'existe pas de fonction $f$ qui résolve l'affaire. Nous pouvons toutefois essayer de trouver un \defe{facteur itnégrant}{Facteur!intégrant}. Nous cherchons une fonction $M$ telle que
\begin{equation}
	(MP)dt+(MQ)dy
\end{equation}
soit exacte. Nous cherchons donc $M(t,y)$ telle que $\partial_y(MP)=\partial_t(MQ)$. En utilisant la règle de Leibnitz, nous trouvons l'équation suivante pour $M$ :
\begin{equation}		\label{EqDuFacteurIntegrant}
	M(\partial_yP-\partial_tQ)=Q(\partial_tM)-P(\partial_yM).
\end{equation}
Cette équation est en générale extrêmement difficile à résoudre, mais dans certains cas particuliers, il est possible d'en trouver une solution à tâtons.


%---------------------------------------------------------------------------------------------------------------------------
					\subsection{Équation d'Euler}
%---------------------------------------------------------------------------------------------------------------------------

La théorie se trouve à la page II.77.


%---------------------------------------------------------------------------------------------------------------------------
					\subsection{Réduction de l'ordre}
%---------------------------------------------------------------------------------------------------------------------------

Afin de diminuer l'ordre d'une équation dans laquelle le paramètre n'apparaît pas, il y a deux changements de variables très utiles. Le premier, le plus simple, est simplement de poser $z(t)=y'(t)$, ce qui donne $z'(t)=y''(t)$. Le second, \emph{qui n'est pas le même}, est $z\big( y(t) \big)=y'(t)$, qui entraîne $y''(t)=z'\big( y(t) \big)z(t)$. Dans ce second cas, il faut également changer de variable, et utiliser $y(t)$ comme variable au lieu de $t$.

%---------------------------------------------------------------------------------------------------------------------------
					\subsection{L'équation \texorpdfstring{$y''=f(y)$}{yfy}}
%---------------------------------------------------------------------------------------------------------------------------

%---------------------------------------------------------------------------------------------------------------------------
					\subsection{Équation ne dépendant pas de \texorpdfstring{$t$}{t}}
%---------------------------------------------------------------------------------------------------------------------------

La théorie se trouve à la page II.79 du syllabus.

%---------------------------------------------------------------------------------------------------------------------------
					\subsection{\ldots et les autres}
%---------------------------------------------------------------------------------------------------------------------------








\chapter{Foire aux questions}
	% This is part of the Exercices et corrigés de CdI-2.
% Copyright (C) 2008-2009, 2012-2013
%   Laurent Claessens
% See the file fdl-1.3.txt for copying conditions.

Le but de ce chapitre est de répondre de façon plus détaillé aux questions qui ont été souvent posées durant les séances d'exercices.

%+++++++++++++++++++++++++++++++++++++++++++++++++++++++++++++++++++++++++++++++++++++++++++++++++++++++++++++++++++++++++++
                    \section{Le coup du compact}
%+++++++++++++++++++++++++++++++++++++++++++++++++++++++++++++++++++++++++++++++++++++++++++++++++++++++++++++++++++++++++++

Dans les sections qui suivent, nous verrons des fonctions définies par toute une série de processus de limite (suites, séries, intégrales). Une des questions centrales est de savoir si la fonction limite est continue, dérivable, intégrale, etc. étant donné que les fonctions sont continues.

Pour cela, nous inventons le concept de \emph{convergence uniforme}. Si la limite (série, intégrale) est uniforme, alors la fonction limite sera continue. Il arrive qu'une limite ne soit pas uniforme sur un intervalle ouvert $]0,1]$, et que nous voulions quand même prouver la continuité sur cet intervalle. C'est à cela que sert la notion de convergence uniforme \emph{sur tout compact}. En effet, la notion de continuité est une notion locale : savoir ce qu'il se passe dans un petit voisinage autour de $x$ est suffisant pour savoir la continuité en $x$ (idem pour sa dérivée).

Si nous avons uniforme convergence sur tout compact de $]0,1]$, mais pas uniforme convergence sur cet intervalle, la limite sera quand même continue sur $\mathopen] 0 , 1 \mathclose]$. En effet, si $x\in]0,1]$, il existe un ouvert autour de $x$ contenu dans un compact contenu dans $]0,1]$. L'uniforme convergence sur ce compact suffit à prouver la continuité en $x$.

Déduire la continuité sur un ouvert à partir de l'uniforme convergence sur tout compact de l'ouvert est appelé faire le \defe{coup du compact}{Compact!le coup du}.

%+++++++++++++++++++++++++++++++++++++++++++++++++++++++++++++++++++++++++++++++++++++++++++++++++++++++++++++++++++++++++++
\section{Vitesses de $x^{\alpha}$, de l'exponentielle et du logarithme}
%+++++++++++++++++++++++++++++++++++++++++++++++++++++++++++++++++++++++++++++++++++++++++++++++++++++++++++++++++++++++++++

Une astuce\label{PgAstuceLnPoly} usuelle qu'il faut savoir est que $\forall \alpha>0$, $\exists N$ tel que $n>N\Rightarrow \ln(n)\leq n^{\alpha}$. En effet, nous avons
\begin{equation}
    \lim_{x\to\infty} \frac{ x^{\alpha} }{ \ln(x) }=\lim_{x\to\infty} \frac{ \alpha x^{\alpha-1} }{ 1/x }=\lim_{x\to\infty} \alpha x^{\alpha}=\infty
\end{equation}
quand $\alpha>0$. Cela tient également lorsque nous considérons $\ln(x)^p$ au lieu de $\ln(x)$. De cela, nous disons que le logarithme croit moins vite que n'importe quel polynôme. 

Dans le même ordre d'idées, l'exponentielle croit plus vite que tout polynôme, et plus vite que que logarithme :
\begin{equation}        \label{EqExpDecrtPlusVite}
    \lim_{t\to\infty} e^{-t}(\ln t)^{n}t^{\alpha}=0
\end{equation}
pour tout $n$ et pour tout $\alpha$.

%+++++++++++++++++++++++++++++++++++++++++++++++++++++++++++++++++++++++++++++++++++++++++++++++++++++++++++++++++++++++++++
\section{Remarque : Abel et \texorpdfstring{$\sin(xt)$}{sin(xt)}}
%+++++++++++++++++++++++++++++++++++++++++++++++++++++++++++++++++++++++++++++++++++++++++++++++++++++++++++++++++++++++++++

Étant donné que la fonction sinus est bornée, il est tentant de l'utiliser comme $\varphi$ dans le critère d'Abel (théorème \ref{ThoAbelIntUnif}). Hélas,
\begin{equation}
    \int_0^T\sin(xt)=-\frac{ 1 }{ x }\big( \cos(xT)-\cos(x) \big),
\end{equation}
qui n'est pas bornée pour tout $x$ ! Poser $\varphi(x,t)=\sin(xt)$ \emph{ne fonctionne pas} pour assurer la convergence uniforme sur un intervalle qui contient des $x$ arbitrairement proches de $0$. Le critère d'Abel avec $\varphi(x,t)=\sin(xt)$ ne permet que de conclure à l'uniforme convergence \emph{sur tout compact} ne contenant pas $0$. Cela est toutefois souvent suffisant pour étudier la continuité ou la dérivabilité en se servant du fameux coup du compact.

%+++++++++++++++++++++++++++++++++++++++++++++++++++++++++++++++++++++++++++++++++++++++++++++++++++++++++++++++++++++++++++
                    \section{Que faire avec \texorpdfstring{$f(z)dz=g(t)dt$}{fzdz} ?}
%+++++++++++++++++++++++++++++++++++++++++++++++++++++++++++++++++++++++++++++++++++++++++++++++++++++++++++++++++++++++++++
\label{SecFairedzdt}

Dans de nombreux exercices d'équations différentielles, nous tombons sur $u'=f(t)$, et nous faisons formellement
\begin{equation}
    \begin{aligned}[]
        \frac{ du }{ dt }&=f(t) &\Rightarrow    &&du=f(t)dt,
    \end{aligned}
\end{equation}
et ensuite, il y a la formule un peu magique
\begin{equation}
    u-u_0=\int_{t_0}^tf(t)dt.
\end{equation}
Voyons ce qu'il en est. Tout d'abord, il faut comprendre ce que signifie la formule
\begin{equation}        \label{EqDiffAstufzdz}
    f(z)dz=g(t)dt.
\end{equation}
Il s'agit d'une égalité entre deux formes différentielles sur $\eR$ où $z$ est une fonction de $t$.  Étant donné que $z$ est une fonction de $t$, il faut voir $dz$ comme la différentielle de cette fonction. La différentielle d'une fonction à une variable est donné par la dérivée :
\begin{equation}
    dz_t=z'(t)dt
\end{equation}
Écrire l'équation \eqref{EqDiffAstufzdz} pour chaque $t$ revient donc à écrire
\begin{equation}
    f\big( z(t) \big)z'(t)dt=g(t)dt
\end{equation}
Cela est une égalité entre deux formes différentielles. Nous avons donc égalité entre les intégrales des formes sur un chemin. Prenons un chemin tout simple de $t_0$ vers $t$ :
\begin{equation}
    \int_{t_0}^tf\big( z(t) \big)z'(t)dt=\int_{t_0}^tg(t)dt.
\end{equation}
Dans le premier membre, nous faisons un changement de variable $\xi=z(t)$, $d\xi=z'(t)dt$, et nous obtenons
\begin{equation}        \label{EqIntDiffAstuztz}
    \int_{z_0}^{z(t)}f(\xi)d\xi=\int_{t_0}^tg(t)dt.
\end{equation}
où nous avons remplacé la constante $z(t_0)$ par $z_0$ dans la borne d'intégration.  Si $F$ est une primitive de $f$ et $G$ une primitive de $g$, nous avons
\begin{equation}
    F(z)-F(z_0)=G(t)-G(t_0).
\end{equation}
Si aucun problème de Cauchy n'est donné, les constantes $F(z_0)$ et $G(t_0)$ sont mises en une seule et nous écrivons la solution
\begin{equation}
    F\big( z(t) \big)=G(t)+C,
\end{equation}
qui est une équation implicite pour $z(t)$. 

Nous trouvons assez souvent le cas simple
\begin{equation}    \label{EqAstfzdzdt}
    f(z)dz=dt.
\end{equation}
En remplaçant $g(t)=1$ dans \eqref{EqIntDiffAstuztz}, nous trouvons la fameuse
\begin{equation}        \label{Eqttzint}
    t-t_0=\int_{z_0}^zf(z)dz,
\end{equation}
dans laquelle il y a un abus de notation terrible entre le $z$ de la borne (que les étudiant(e)s oublient souvent) et la variable d'intégration $z$ !!

Le passage de \eqref{EqAstfzdzdt} à \eqref{Eqttzint} sera très souvent utilisé dans le cours de mécanique par exemple.


%+++++++++++++++++++++++++++++++++++++++++++++++++++++++++++++++++++++++++++++++++++++++++++++++++++++++++++++++++++++++++++
                    \section{Pourquoi la variation des constantes fonctionne toujours ?}
%+++++++++++++++++++++++++++++++++++++++++++++++++++++++++++++++++++++++++++++++++++++++++++++++++++++++++++++++++++++++++++

Prenons une équation non homogène 
\begin{equation}        \label{EqAstNNHomo}
    z'(t)=f(t)z(t)+g(t),
\end{equation}
et supposons avoir une solution de l'homogène associée sous la forme $z_H(t)=Ch(t)$. Le coup de la variation des constates consiste à essayer une solution pour l'équation non homogène sous la forme\footnote{Je ne sais plus qui a eu l'idée de changer le nom de la constante de $C$ vers $K$ au moment de la transformer en fonction, mais c'est une bonne idée.}
\begin{equation}
    z(t)=K(t)h(t).
\end{equation}
Nous injectons cette solution dans l'équation de départ en utilisant le fait que $z'(t)=K'(t)h(t)+K(t)h'(t)$ :
\begin{equation}
    K'(t)h(t)+K(t)h'(t)=f(t)K(t)h(t)+g(t).
\end{equation}
Le terme $K(t)h'(t)$ se récrit en utilisant la propriété de définition de $h$, c'est à dire que $h'(t)=f(t)h(t)$. Nous voyons que les termes ne contenant pas de $K'$ se simplifient; il reste
\begin{equation}
    K'h=g.
\end{equation}
Cette équation a comme solution
\begin{equation}
    K=\int \frac{ f }{ h }+C.
\end{equation}
J'insiste sur la constante d'intégration ! En réalité, celles et ceux qui auront compris l'équation \eqref{Eqttzint} sauront que $K$ est donné par
\begin{equation}
    K(t)=\int_{\xi_0}^{t}\frac{ f(\xi) }{ g(\xi) }d\xi
\end{equation}
où $\xi_0$ joue le rôle de la constante d'intégration.

Quoi qu'il en soit, la solution générale de l'équation non homogène est
\begin{equation}        \label{EqSolVarCosntCool}
    z(t)=K(t)h(t)=\left( \int\frac{ g }{ h }+C \right)h.
\end{equation}
Cette solution comprend deux termes : $Ch$ qui est solution de l'homogène, et $\left( \int \frac{ g }{ h } \right)h$ qui est une particulière de l'équation non homogène.

Quelque conclusions :

\begin{enumerate}
\item
Si vous avez encore du $K$ (et pas que du $K'$) dans votre équation qui donne $K$, c'est que vous n'être pas dans le cadre d'une équation de type \eqref{EqAstNNHomo}. Le plus souvent, c'est que vous avez fait une faute de calcul quelque part.

\item
La méthode des variations des constantes n'est pas en contradiction avec le principe de \og SGEH+SPENH\fg. En effet, la SGEP et la SPENH sont toutes deux dans la solution \eqref{EqSolVarCosntCool}.

\item
La variation des constantes peut être vue comme une façon cool de trouver une solution particulière de l'équation non homogène.

\item
La simplification ne se fait que après avoir remplacé $Kh'$ par $Kfh$, c'est à dire après avoir utilisé le fait que $z_H$ est solution de l'homogène. Sinon, la simplification n'est pas du tout évidente a priori. Il se peut même que, visuellement, les termes $Kh'$ et $Kfh$ ne se ressemblent pas du tout. C'est ce qui arrive par exemple à l'exercice \ref{exoreserve0004}\ref{ItemCII107}, pour arriver à l'équation \eqref{EqFracII107exoVVprb}.

\end{enumerate}

%+++++++++++++++++++++++++++++++++++++++++++++++++++++++++++++++++++++++++++++++++++++++++++++++++++++++++++++++++++++++++++
\section{Trucs et astuces de calculs d'intégrales}
%+++++++++++++++++++++++++++++++++++++++++++++++++++++++++++++++++++++++++++++++++++++++++++++++++++++++++++++++++++++++++++

Afin d'alléger le texte de calculs parfois un peu longs, nous regroupons ici les intégrales à une variable que nous devons utiliser dans les autres parties du cours.

\begin{enumerate}
	\item	\label{ItemIntegrali}
		L'intégrale
		\begin{equation}
			\boxed{I=\int x\ln(x)dx=\frac{ x^2 }{2}\big( \ln(x)-\frac{ 1 }{2} \big)}
		\end{equation}
		se fait par partie en posant
		\begin{equation}
			\begin{aligned}[]
				u&=\ln(x),		& dv&=x\,dx\\
				du&=\frac{1}{ x }\,dx,	& v&=\frac{ x^2 }{2},
			\end{aligned}
		\end{equation}
		et ensuite
		\begin{equation}
			I=\ln(x)\frac{ x^2 }{2}-\int\frac{ x }{2}=\frac{ x^2 }{2}\big( \ln(x)-\frac{ 1 }{2} \big).
		\end{equation}
		
	\item	
		L'intégrale
		\begin{equation}
			\boxed{I=\int x\ln(x^2)dx=x^2\ln(x)-\frac{ x^2 }{2}.}
		\end{equation}
		En utilisant le fait que $\ln(u^2)=2\ln(u)$, nous retombons sur une intégrale du type \ref{ItemIntegrali} :
		\begin{equation}
			I=x^2\ln(x)-\frac{ x^2 }{2}.
		\end{equation}
	\item
		L'intégrale
		\begin{equation}		\label{EqTrucIntxlnxsqpun}
			\boxed{I=\int x\ln(1+x^2)dx=\frac{ 1 }{2}\ln(x^2+1)(x^2+1)-x^2-\frac{ 1 }{2}}
		\end{equation}
		se traite en posant $v=1+x^2$ de telle sorte à avoir $dx=\frac{ dv }{ 2x }$ et donc
		\begin{equation}
			I=\frac{ 1 }{2}\ln(x^2+1)(x^2+1)-x^2-\frac{ 1 }{2}.
		\end{equation}
		
	\item
		L'intégrale
		\begin{equation}
			I=\int \cos(\theta)\sin(\theta)\ln\left( 1+\frac{1}{ \cos^2(\theta) } \right)\,d\theta
		\end{equation}
		demande le changement de variable $u=\cos(\theta)$, $d\theta=-\frac{ du }{ \sin(\theta) }$. Nous tombons sur l'intégrale
		\begin{equation}
			I=-\int u\ln\left( \frac{ 1+u^2 }{ u^2 } \right)=-\int u\ln(1+u^2)+\int u\ln(u^2),
		\end{equation}
		qui sont deux intégrales déjà faites. Nous trouvons
		\begin{equation}
			I=-\frac{ 1 }{2}\ln\left( \frac{ \sin^2(\theta)-1 }{ \sin^2(\theta)-2 } \right)\sin^2(\theta)-\ln\big( \sin^2(\theta)-2 \big)+\frac{ 1 }{2}\ln\big( \sin^2(\theta)-1 \big)
		\end{equation}
	
	\item
		L'intégrale
		\begin{equation}
			\boxed{\int \frac{ r^3 }{ 1+r^2 }dr=\frac{ r^2 }{2}-\frac{ 1 }{2}\ln(r^2+1).}
		\end{equation}
		commence par faire la division euclidienne de $r^3$ par $r^2+1$; ce que nous trouvons est $r^3=(r^2+1)r-r$. Il reste à intégrer
		\begin{equation}
			\int \frac{ r^3 }{ 1+r^2 }dr=\int r\,dr-\int\frac{ r }{ 1+r^2 }dr.
		\end{equation}
		La fonction dans la seconde intégrale est $\frac{ r }{ 1+r^2 }=\frac{ 1 }{2}\frac{ f'(r) }{ f(r) }$ où $f(r)=1+r^2$, et donc $\int \frac{ r }{ 1+r^2 }=\frac{ 1 }{2}\ln(1+r^2)$. Au final,
		\begin{equation}
			I=\frac{ 1 }{2}r^2-\frac{ 1 }{2}\ln(r^2+1).
		\end{equation}


	\item	
		L'intégrale
		\begin{equation}	\label{EqTrucIntsxcxdx}
			\boxed{I=\int \cos(\theta)\sin(\theta)d\theta=\frac{ \sin^2(\theta) }{ 2 }}
		\end{equation}
		se traite par le changement de variable $u=\sin(\theta)$, $du=\cos(\theta)d\theta$, et donc
		\begin{equation}
			\int\cos(\theta)\sin(\theta)d\theta=\int udu=\frac{ u^2 }{2}=\frac{ \sin^2(\theta) }{ 2 }.
		\end{equation}
	\item
		L'intégrale
		\begin{equation}	\label{EqTrucsIntsqrtAplusu}
			\boxed{\int\sqrt{1+x^2}dx=\frac{ x }{2}\sqrt{1+x^2}+\frac{ 1 }{2}\arcsinh(x)}
		\end{equation}
		s'obtient en effectuant le changement de variable $u=\sinh(\xi)$.

    \item
        L'intégrale
        \begin{equation}        \label{EqTrucIntcossqsinsq}
            \boxed{ \int\cos^2(x)\sin^2(x)dx=\frac{ x }{ 8 }-\frac{ \sin(4x) }{ 32 } }
        \end{equation}
        s'obtient à coups de formules de trigonométrie. D'abord, $\sin(t)\cos(t)=\frac{ 1 }{2}\sin^2(2t)$ fait en sorte que la fonction à intégrer devient 
        \begin{equation}
            f(x)=\frac{1}{ 4 }\sin^2(x).
        \end{equation}
        Ensuite nous utilisons le fait que $\sin^2(t)=(1-\cos(2t))/2$ pour transformer la formule à intégrer en
        \begin{equation}
            f(x)=\frac{ 1-\cos(4x) }{ 8 }.
        \end{equation}
        Cela s'intègre facilement en posant $u=4x$, et le résultat est
        \begin{equation}
            \int f(x)dx=\frac{ x }{ 8 }-\frac{ \sin(4x) }{ 32 }.
        \end{equation}

\end{enumerate}


\chapter{Exercices}
	% This is part of the Exercices et corrigés de CdI-2.
% Copyright (C) 2008, 2009
%   Laurent Claessens
% See the file fdl-1.3.txt for copying conditions.

%+++++++++++++++++++++++++++++++++++++++++++++++++++++++++++++++++++++++++++++++++++++++++++++++++++++++++++++++++++++++++++
					\section{Suites de fonctions}
%+++++++++++++++++++++++++++++++++++++++++++++++++++++++++++++++++++++++++++++++++++++++++++++++++++++++++++++++++++++++++++

	\Exo{_I-1-1}
	\Exo{_I-1-2}
	\Exo{_I-1-3}
	\Exo{_I-1-4}
	\Exo{_I-1-5}
	\Exo{_I-1-6}
	\Exo{_I-1-7}
	\Exo{_I-1-8}
	\Exo{_I-1-8b}
	\Exo{_I-1-9}

%+++++++++++++++++++++++++++++++++++++++++++++++++++++++++++++++++++++++++++++++++++++++++++++++++++++++++++++++++++++++++++
					\section{Séries de fonctions}
%+++++++++++++++++++++++++++++++++++++++++++++++++++++++++++++++++++++++++++++++++++++++++++++++++++++++++++++++++++++++++++

	\Exo{_I-1-10}
	\Exo{_I-1-11}
	\Exo{_I-1-12}
	\Exo{_I-1-13}
	\Exo{_I-1-14}
	\Exo{_I-1-15}
	\Exo{_I-1-16}

Après avoir fait le théorème de Borel, essayons de voir ce qu'il dit. Le théorème de la page I.29 nous donne facilement, sous forme de séries de puissances, une fonction $ C^{\infty}$ sur un intervalle dont les dérivées en zéro sont données à l'avance. L'exercice que nous venons de faire nous permet de trouver une fonction $ C^{\infty}$ sur tout $\eR$ dont les dérivées sont données à l'avance.

% TODO : La ligne suivante m'a l'air un peu trop simple pour être vraie ... y réfléchir encore un peu avant de décommenter :
%De plus, la fonction $u$ ainsi définie est à support compact, contenu dans le support de $f$, parce que $t_n\leq 1$. Nous obtenons ainsi des fonctions $ C^{\infty}$ à support compact dont les dérivées en un point sont données à l'avance.

%+++++++++++++++++++++++++++++++++++++++++++++++++++++++++++++++++++++++++++++++++++++++++++++++++++++++++++++++++++++++++++
					\section{Existence d'intégrales}
%+++++++++++++++++++++++++++++++++++++++++++++++++++++++++++++++++++++++++++++++++++++++++++++++++++++++++++++++++++++++++++

En vertu des différents théorèmes, l'étude de la convergence et de l'existence d'une intégrale d'une fonction sur $[a,\infty[$ se fait dans l'ordre suivant :
\begin{itemize}
\item si l'intégrale de $f$ existe, alors elle converge,
\item si la fonction est positive et que son intégrale n'existe pas, alors elle ne converge pas,
\item si le signe de la fonction est variable et que l'intégrale n'existe pas, alors elle peut converger ou non, selon les cas.
\end{itemize}

	\Exo{_I-2-1}
	\Exo{_I-2-2}
	\Exo{_I-2-3}
	\Exo{_I-2-4}
	
%+++++++++++++++++++++++++++++++++++++++++++++++++++++++++++++++++++++++++++++++++++++++++++++++++++++++++++++++++++++++++++
					\section{Fonctions définies par des intégrales}
%+++++++++++++++++++++++++++++++++++++++++++++++++++++++++++++++++++++++++++++++++++++++++++++++++++++++++++++++++++++++++++


	\Exo{_I-3-1}
	\Exo{_I-3-2}

%+++++++++++++++++++++++++++++++++++++++++++++++++++++++++++++++++++++++++++++++++++++++++++++++++++++++++++++++++++++++++++
					\section{Convergence, continuité et dérivation sous le signe intégral}
%+++++++++++++++++++++++++++++++++++++++++++++++++++++++++++++++++++++++++++++++++++++++++++++++++++++++++++++++++++++++++++

	\Exo{_I-3-4}
	\Exo{_I-3-5}
	\Exo{_I-3-6}
	\Exo{_I-3-7}
	\Exo{_I-3-8}
	\Exo{_I-3-9}
	\Exo{_I-3-10}
	\Exo{_I-3-11}
%+++++++++++++++++++++++++++++++++++++++++++++++++++++++++++++++++++++++++++++++++++++++++++++++++++++++++++++++++++++++++++
					\section{Quelques propriétés des espaces fonctionnels}
%+++++++++++++++++++++++++++++++++++++++++++++++++++++++++++++++++++++++++++++++++++++++++++++++++++++++++++++++++++++++++++

	\Exo{_I-4-1}
	\Exo{_I-4-2}
	\Exo{_I-4-3}


	% This is part of the Exercices et corrigés de CdI-2.
% Copyright (C) 2008, 2009
%   Laurent Claessens
% See the file fdl-1.3.txt for copying conditions.


%+++++++++++++++++++++++++++++++++++++++++++++++++++++++++++++++++++++++++++++++++++++++++++++++++++++++++++++++++++++++++++
					\section{Équations différentielles}
%+++++++++++++++++++++++++++++++++++++++++++++++++++++++++++++++++++++++++++++++++++++++++++++++++++++++++++++++++++++++++++

%---------------------------------------------------------------------------------------------------------------------------
					\subsection{Équations différentielles résolubles}
%---------------------------------------------------------------------------------------------------------------------------

\Exo{_II-1-01}
\Exo{_II-1-02}
\Exo{_II-1-03}
\Exo{_II-1-04}
\Exo{_II-1-05}
\Exo{_II-1-06}

%---------------------------------------------------------------------------------------------------------------------------
					\subsection{Équation de Bernoulli}
%---------------------------------------------------------------------------------------------------------------------------

\Exo{reserve0004}
\Exo{_II-1-08}

%TODO : changer tous les noms de fichiers contenant «_».

%---------------------------------------------------------------------------------------------------------------------------
					\subsection{Équations de Ricatti}
%---------------------------------------------------------------------------------------------------------------------------


\Exo{_II-1-09}

%---------------------------------------------------------------------------------------------------------------------------
					\subsection{Équations homogènes}
%---------------------------------------------------------------------------------------------------------------------------


\Exo{_II-1-10}

%---------------------------------------------------------------------------------------------------------------------------
					\subsection{Équations différentielles exactes. Facteurs intégrants}
%---------------------------------------------------------------------------------------------------------------------------


\Exo{_II-1-11}
\Exo{_II-1-12}


%---------------------------------------------------------------------------------------------------------------------------
					\subsection{L'équation \texorpdfstring{$y'=f\left( \frac{ at+by+c }{ a't+b'y+c' } \right)$}{yat}}
%---------------------------------------------------------------------------------------------------------------------------

\Exo{_II-1-13}

%---------------------------------------------------------------------------------------------------------------------------
					\subsection{Équation d'Euler}
%---------------------------------------------------------------------------------------------------------------------------

\Exo{_II-1-14}

%---------------------------------------------------------------------------------------------------------------------------
					\subsection{Équation dont on peut réduire l'ordre}
%---------------------------------------------------------------------------------------------------------------------------

\Exo{_II-1-15}
\Exo{_II-1-16}

%---------------------------------------------------------------------------------------------------------------------------
					\subsection{Quelques exemples d'équations résolubles}
%---------------------------------------------------------------------------------------------------------------------------

\Exo{_II-1-17}

%---------------------------------------------------------------------------------------------------------------------------
					\subsection{L'équation différentielle linéaire du second ordre}
%---------------------------------------------------------------------------------------------------------------------------

\Exo{_II-1-20}
\Exo{_II-1-21}
\Exo{_II-1-22}
\Exo{_II-1-23}
\Exo{_II-1-24}

%---------------------------------------------------------------------------------------------------------------------------
					\subsection{Équation différentielle implicite du premier ordre}
%---------------------------------------------------------------------------------------------------------------------------

\Exo{_II-1-18}
\Exo{_II-1-19}
\Exo{_II-1-28}
\Exo{_II-1-25}
\Exo{_II-1-26}
\Exo{_II-1-27}

%+++++++++++++++++++++++++++++++++++++++++++++++++++++++++++++++++++++++++++++++++++++++++++++++++++++++++++++++++++++++++++
					\section{Systèmes différentiels linéaires à coefficients constants}
%+++++++++++++++++++++++++++++++++++++++++++++++++++++++++++++++++++++++++++++++++++++++++++++++++++++++++++++++++++++++++++


\Exo{_II-2-01}
\Exo{_II-2-02}
\Exo{_II-2-03}
\Exo{_II-2-04}
\Exo{_II-2-05}
\Exo{_II-2-06}
\Exo{_II-2-07}






\input{fdl-1.3}



\corrChapitre{Corrigés systématiques}

\bibliographystyle{unsrt}           % unsrt fait que la biblio arrive dans l'ordre de citation au lieu de l'ordre alphabétique.
\bibliography{mazhe}

    \addcontentsline{toc}{chapter}{Liste des notations}
    \printnomenclature
    \printindex


\end{document}


\Exo{Model-0007}
\Exo{Model-0008}
\Exo{Model-0009}
\Exo{Model-0010}
\Exo{Model-0011}
\Exo{Model-0012}
\Exo{Model-0013}
\Exo{Model-0014}
\Exo{Model-0015}
\Exo{Model-0016}
\Exo{Model-0017}
\Exo{Model-0018}
\Exo{Model-0019}
\Exo{Model-0020}
\Exo{Model-0021}
\Exo{Model-0022}
\Exo{Model-0023}
\Exo{Model-0024}
\Exo{Model-0025}
\Exo{Model-0026}
\Exo{Model-0027}
\Exo{Model-0028}
\Exo{Model-0029}
\Exo{Model-0030}
